% book14 extends book13 and adds qualitative neck--cap structure,
% canonical neighborhoods, and backward neck stability.
\documentclass[10pt]{coll-l}

\usepackage{amsfonts}
\usepackage{amsmath}
\usepackage{mathrsfs}
\usepackage{amssymb}
\usepackage{amstext}
\usepackage{amsthm}
\usepackage{graphicx}
\usepackage{lineno}
\usepackage{xcolor}
\usepackage{minted}
\usemintedstyle{tango}
\usepackage{tikz}
\usetikzlibrary{arrows.meta,positioning}
\usepackage{tikz-cd}
\usepackage{hyperref}
\usepackage{xstring}
\hypersetup{
  colorlinks=true,
  linkcolor=blue,
  filecolor=magenta,
  urlcolor=cyan
}

\newcommand{\inj}{\operatorname{inj}}
\newcommand{\Isom}{\operatorname{Isom}}
\newcommand{\dist}{\operatorname{dist}}
\newcommand{\supp}{\operatorname{supp}}
\newcommand{\trace}{\operatorname{trace}}
\newcommand{\nodeid}[1]{{\scriptsize\ttfamily
  \StrSubstitute{#1}{-}{-\allowbreak}}}

\newcommand{\R}{\mathbb{R}}
\newcommand{\Sph}{\mathbb{S}}

\newcommand{\Ttwo}{\mathsf{T2Space}}
\newcommand{\ChartedSpace}{\mathsf{ChartedSpace}}
\newcommand{\IsManifold}{\mathsf{IsManifold}}
\newcommand{\SimplyConnected}{\mathsf{SimplyConnectedSpace}}
\newcommand{\Compact}{\mathsf{CompactSpace}}
\newcommand{\Nonempty}{\mathsf{Nonempty}}
\newcommand{\Rstruct}{\mathcal{R}}
\newcommand{\diffeo}{\mathrel{\cong_{\mathrm{diff}}}}

\newenvironment{roadmap}[1]{%
  \begin{quote}\small\noindent\textbf{Roadmap.} #1\par\medskip\noindent
}{\end{quote}}

\newenvironment{leancode}
  {\VerbatimEnvironment
   \begin{minted}[
     fontsize=\small,
     breaklines,
     mathescape,
     frame=lines,
     escapeinside=||
   ]{lean}}
  {\end{minted}}

\newtheorem{theorem}{Theorem}[chapter]
\newtheorem{axiom}[theorem]{Axiom}
\newtheorem{corollary}[theorem]{Corollary}
\newtheorem{conjecture}[theorem]{Conjecture}
\newtheorem{example}[theorem]{Example}
\newtheorem{exercise}{Exercise}
\newtheorem{lemma}[theorem]{Lemma}
\newtheorem{proposition}[theorem]{Proposition}
\newtheorem{question}[theorem]{Question}

\theoremstyle{definition}
\newtheorem{definition}[theorem]{Definition}
\newtheorem{deflem}[theorem]{Definition and Lemma}
\theoremstyle{remark}
\newtheorem{remark}[theorem]{Remark}

\numberwithin{section}{chapter}
\numberwithin{equation}{chapter}

\newcommand{\ASCR}{\operatorname{ASCR}}
\newcommand{\AVR}{\operatorname{AVR}}
\newcommand{\E}{\mathrm{e}}
\newcommand{\Ric}{\operatorname{Ric}}
\newcommand{\Rm}{\operatorname{Rm}}
\newcommand{\Vol}{\operatorname{Vol}}

\newcommand{\vspacesub}{$\left.\right.$\vspace{3pt}}

\makeindex
\setcounter{tocdepth}{2}

\begin{document}

\frontmatter

\title{The Hamilton--Perelman Proof of the Poincar\'e Conjecture}
\author{Bennett Chow, Yuan Liao, and Ziyang Qin}
\date{}

\subjclass[2020]{Primary 53E20; Secondary 57K30}
\keywords{Poincar\'e conjecture, Ricci flow, surgery, three-manifolds, Lean}

\maketitle

\setcounter{page}{5}
\setcounter{chapter}{0}

\tableofcontents

% Include unnumbered chapters (preface, acknowledgments, etc.) here.
\include{preface}

\mainmatter

\chapter[Poincar\'e proof strategy]
{The Poincar\'e Conjecture and the Hamilton--Perelman Strategy}
\label{ch:architecture}

\textcolor{red}{For coauthors: We will put important short Lean code in the text right after the definition or statement that is formalized. Longer Lean code of important stuff will be put in an appendix. We eventually want the Lean code in the text to be the Lean code we use in our formalization.}

\section{Introduction}

This book presents a self-contained mathematical exposition of the
Hamilton--Perelman proof of the Poincar\'e conjecture.  It is written for
subsequent formalization in Lean: definitions have explicit domains, theorem
statements include the hypotheses used downstream, and every passage among
the topological, piecewise-linear, and smooth categories is identified.
Nevertheless, the book is a work of mathematics, not a transcription of Lean
code.

The book, blueprint, Lean development, and status ledger have different
roles.  This book is the fundamental human-readable source for the intended
mathematics: it states the definitions and theorems in ordinary mathematical
language and supplies their proofs.  The companion blueprint assigns stable
semantic identifiers, records exact statements and dependencies, and gives
the book-to-Lean crosswalk.  When an approved blueprint node has been
implemented, the Lean development supplies its kernel-checked declaration.
The status ledger records the evidence for compilation, trust, and semantic
correspondence node by node.  Except for a few landmark declarations, Lean
code is therefore not displayed in the book.

This opening chapter fixes the statement of the conjecture, the category
bridge that reduces it to a smooth theorem, the macro-level proof
architecture, and the interfaces of the two final inputs: finite extinction
and topological reconstruction.  Later chapters prove the milestones
recorded here.

Throughout the book, a \emph{closed manifold} means a compact manifold
without boundary.  Closedness does not by itself include connectedness or
orientability.

We use the convention that a connected manifold is nonempty; this agrees
with mathlib's \texttt{ConnectedSpace}.  Empty manifolds and empty time
slices are always described explicitly.

The companion Lean library for differential geometry and Ricci flow is the
\href{https://github.com/qinz1yang/differential-geometry}
{\texttt{differential-geometry} repository} maintained by Ziyang Qin.  For
reproducibility, this
chapter pins the source release at
\href{https://github.com/qinz1yang/differential-geometry/commit/8bd406e35c33a200e9b88895cf11ee8429194e15}
{commit \texttt{8bd406e3}}
(\href{https://github.com/qinz1yang/differential-geometry/releases/tag/arxiv-v1-preview}
{release 0.1.0}).  The endpoint declaration
\texttt{hamilton\_positive\_ricci}, which states the positive-Ricci
three-manifold conclusion used in that formalization, appears in
\href{https://github.com/qinz1yang/differential-geometry/blob/8bd406e35c33a200e9b88895cf11ee8429194e15/DifferentialGeometry/Geometry/Flow/RicciFlow/DimensionThree/PositiveRicci/Hamilton.lean}
{\texttt{Hamilton.lean}}.

\section{The topological and smooth statements}

The Clay Mathematics Institute's
\href{https://www.claymath.org/wp-content/uploads/2022/06/poincare.pdf}
{official problem description}, written by John Milnor, asks the following
question.

\begin{question}[Poincar\'e conjecture]
If a compact three-dimensional manifold $M$ has the property that every
simple closed curve in $M$ can be continuously deformed to a point, must
$M$ be homeomorphic to the sphere $S^{3}$?
\end{question}

Here $M$ is understood to be connected and without boundary; see
Poincar\'e's 1904 paper \cite[pp.~370 and 498]{MR1401792}.  We retain this
wording as a historical quotation.  The project-facing statement uses the
standard notion of simple connectivity, which is also the interface used in
Lean.

\begin{theorem}[Topological Poincar\'e conjecture]
\label{thm:topological-poincare}
Every closed, connected, simply connected topological $3$-manifold is
homeomorphic to $S^{3}$.
\end{theorem}

At the pinned mathlib revision
\href{https://github.com/leanprover-community/mathlib4/blob/e9d4384eb66de816031af800a5ac435ca06a15da/Mathlib/Geometry/Manifold/PoincareConjecture.lean\#L46-L49}
{\texttt{e9d4384}}, mathlib records the topological three-dimensional
conclusion as the following unproved target statement.

\begin{leancode}
proof_wanted SimplyConnectedSpace.nonempty_homeomorph_sphere_three
  [T2Space M]
  [ChartedSpace |$\R^3$| M]
  [SimplyConnectedSpace M]
  [CompactSpace M] :
  Nonempty (M |$\simeq_{\mathrm{t}}$| |$\Sph^3$|)
\end{leancode}

At this revision, \texttt{proof\_wanted} checks that the displayed target is
a well-formed Lean proposition and then discards the temporary declaration.
It therefore supplies neither a proof nor a theorem that subsequent Lean code
can invoke.  The hypotheses describe a compact Hausdorff topological
$3$-manifold that is simply connected; in mathlib,
\texttt{SimplyConnectedSpace} includes the relevant nonemptiness and
path-connectedness data.  The conclusion asks for a homeomorphism with the
standard $3$-sphere.

The Ricci-flow argument is carried out in the smooth category.  Its immediate
target is therefore the following theorem.

\begin{theorem}[Smooth Poincar\'e conjecture]
\label{thm:smooth-poincare}
Every closed, connected, simply connected smooth $3$-manifold is
diffeomorphic to $S^{3}$.
\end{theorem}

The
\href{https://github.com/leanprover-community/mathlib4/blob/e9d4384eb66de816031af800a5ac435ca06a15da/Mathlib/Geometry/Manifold/PoincareConjecture.lean\#L51-L55}
{same pinned mathlib file} records the smooth three-dimensional conclusion as
the following unproved target statement.

\begin{leancode}
proof_wanted SimplyConnectedSpace.nonempty_diffeomorph_sphere_three
  [T2Space M]
  [ChartedSpace |$\R^3$| M]
  [IsManifold |$(\Rstruct\ 3)$| |$\infty$| M]
  [SimplyConnectedSpace M]
  [CompactSpace M] :
  Nonempty (M |$\simeq_{\mathrm{m}} ( \Rstruct\,3,\ \Rstruct\,3 )$| |$\Sph^{3}$|)
\end{leancode}

The displayed target is a formal specification, not a proof of the Poincar\'e
conjecture.  Its hypotheses say that $M$ is a Hausdorff smooth
$3$-manifold that is compact and simply connected; in mathlib,
\texttt{SimplyConnectedSpace} includes the relevant nonemptiness and
path-connectedness data.  The conclusion asserts the existence of a
diffeomorphism from $M$ to the standard smooth $3$-sphere.

\section{The topological-to-smooth bridge}
\label{sec:moise-bridge}

The only smoothability result needed in this book is the following closed,
three-dimensional statement.

\begin{theorem}[Smoothability of closed topological $3$-manifolds]
\label{thm:moise-smoothability}
Let $M$ be a closed topological $3$-manifold.  There exists a smooth
$3$-manifold structure on the set $M$ whose underlying topology is exactly
the given topology of $M$.
\end{theorem}

The following schematic code records the proposed project-facing Lean
interface.  It has not yet been elaborated or compiled in the pinned project,
and it is not presently an exported or proved mathlib theorem.

\begin{leancode}
proof_wanted IsManifold.exists_smooth_structure_three
  [T2Space M]
  [ChartedSpace |$\R^3$| M]
  [CompactSpace M] :
  |$\exists$| s : ChartedSpace |$\R^3$| M,
    letI : ChartedSpace |$\R^3$| M := s
    IsManifold |$(\Rstruct\ 3)$| |$\infty$| M
\end{leancode}

The ambient topological-space structure on $M$ is held fixed, while only the
charted-space structure is existentially quantified.  The proposed interface
therefore retains the given topology on the same carrier.  This is the
precise same-carrier form needed downstream.  It is a consequence of Moise's
triangulation theorem for topological $3$-manifolds and the smoothability of
piecewise-linear $3$-manifolds.  No uniqueness theorem and no noncompact
generalization are required for the Poincar\'e argument.

Theorem~\ref{thm:moise-smoothability} reduces
Theorem~\ref{thm:topological-poincare} to
Theorem~\ref{thm:smooth-poincare}: after smoothing the same underlying set, a
diffeomorphism to $S^{3}$ is, on forgetting the smooth structures, a
homeomorphism for the original topology.

\section{The macro-level proof architecture}
\label{sec:macro-architecture}

The permanent blueprint uses semantic identifiers that do not change when
chapters are renumbered or Lean declarations are moved.  The argument has the
following ten macro milestones.

\begin{enumerate}
\item \nodeid{PC-M-TOPOLOGICAL-SMOOTH-BRIDGE}:
place the given closed
topological $3$-manifold in the smooth category without changing its topology.

\item \nodeid{PC-M-SMOOTH-RICCI-FLOW}:
choose an initial metric and construct
the corresponding short-time Ricci flow.

\item \nodeid{PC-M-SINGULARITY-EXTENSION}:
prove that bounded curvature
extends a closed Ricci flow and hence that finite maximal time forces
curvature blow-up.

\item \nodeid{PC-M-HIGH-CURVATURE-CONTROL}:
establish noncollapsing,
singularity-model, and canonical-neighborhood results at high curvature.

\item \nodeid{PC-M-CUT-CAP-TOPOLOGY}:
formalize spherical cutting, capping,
discarding, and their reversal by connected sums.

\item \nodeid{PC-M-SURGERY-CONTINUATION}:
construct Perelman's controlled
surgery continuation and verify the abstract interface below.

\item \nodeid{PC-M-FINITE-EXTINCTION}:
prove finite extinction in the simply
connected case.

\item \nodeid{PC-M-EXTINCTION-RECONSTRUCTION}:
reverse the finite surgery
history and show that the initial manifold is Poincar\'e-standard.

\item \nodeid{PC-M-SIMPLY-CONNECTED-CLASSIFICATION}:
prove that a connected,
simply connected Poincar\'e-standard manifold is diffeomorphic to $S^{3}$.

\item \nodeid{PC-M-POINCARE-CONCLUSION}:
conclude the smooth theorem and then
the topological theorem by the Moise bridge.
\end{enumerate}

Figure~\ref{fig:macro-architecture} displays the two main lanes of the proof.
The analytic lane produces an extinct surgery history; the topological lane
explains how to reverse its cut--cap events.  They meet at reconstruction.

\clearpage
\begin{figure}[p]
\centering
\begin{tikzpicture}[
  box/.style={
    draw,
    rounded corners,
    align=center,
    text width=43mm,
    minimum height=10mm,
    inner sep=3pt
  },
  >={Stealth[length=2mm]},
  node distance=8mm and 14mm
]
\node[box] (bridge) {1. Topological--smooth bridge};
\node[box,below=of bridge] (flow) {2. Smooth Ricci flow};
\node[box,below=of flow] (extension) {3. Finite-time extension criterion};
\node[box,below=of extension] (high) {4. High-curvature control};
\node[box,right=of high] (cutcap) {5. Cut--cap topology};
\node[box,below=of high] (surgery) {6. Surgery continuation};
\node[box,below=of surgery] (extinction) {7. Finite extinction};
\node[box,below=of extinction] (reconstruct) {8. Extinction reconstruction};
\node[box,below=of reconstruct] (classification) {
  9. Simply connected classification};
\node[box,below=of classification] (conclusion) {10. Poincar\'e conclusion};

\draw[->] (bridge) -- (flow);
\draw[->] (flow) -- (extension);
\draw[->] (extension) -- (high);
\draw[->] (high) -- (surgery);
\draw[->] (cutcap) -- (surgery);
\draw[->] (surgery) -- (extinction);
\draw[->] (extinction) -- (reconstruct);
\draw[->] (cutcap) |- (reconstruct);
\draw[->] (reconstruct) -- (classification);
\draw[->] (classification) -- (conclusion);
\draw[->] (bridge.west) -- ++(-3mm,0) |- (conclusion.west);
\end{tikzpicture}
\caption{The macro-level dependency graph.}
\label{fig:macro-architecture}
\end{figure}
\clearpage

\section{Smooth Ricci flow}

\begin{definition}[Ricci flow]
Let $M$ be a smooth manifold and let $I\subseteq\mathbb R$ be an interval.  A
\emph{Ricci flow} on $M$ over $I$ is a jointly smooth family of Riemannian
metrics $g(t)$, $t\in I$, satisfying
\[
  \frac{\partial}{\partial t}g(t)=-2\Ric_{g(t)}
\]
for every $t\in I$.  At an interior time the derivative is the ordinary time
derivative, and at an endpoint belonging to $I$ it is the corresponding
one-sided derivative.
\end{definition}

At the pinned \texttt{differential-geometry} revision above, the reusable
interval-wise equation predicate writes the tensor equation componentwise,
after evaluation on arbitrary tangent vectors:

\begin{leancode}
def MetricConnectionFamilyVariationEquationOn
    {D : RealTimeInterval}
    (G : MetricConnectionFamilyOn (I := I) (M := M) D)
    (Ric : RicciTensorField (I := I) (M := M) Real) : Prop :=
  forall (t : RealTimeInterval.RegularTime D)
      (x : M) (X Y : TangentSpace I x),
    HasDerivWithinAt
      (fun s : Real => (G.metric s).inner x X Y)
      ((-2 : Real) * Ric (t : Real) x X Y)
      D.carrier
      (t : Real)
\end{leancode}

\medskip

The source is
\href{https://github.com/qinz1yang/differential-geometry/blob/8bd406e35c33a200e9b88895cf11ee8429194e15/DifferentialGeometry/Geometry/Flow/RicciFlow/Solution/Defs.lean\#L63-L73}
{\texttt{Solution/Defs.lean}}.  Because that reusable predicate accepts a
tensor field \texttt{Ric} as an argument, the actual solution interface
specializes it to the Ricci tensor constructed from the evolving metric.  In
the library's ambient namespace, the specialization has the following
normalized form:

\begin{leancode}
def MetricVariationEquationOn
    {D : RealTimeInterval}
    (S : SolutionOn (I := I) (M := M) D) : Prop :=
  MetricConnectionFamilyVariationEquationOn (I := I) S.family
    (RicciAtFamily.toTensorField (I := I) S.ricciAt)
\end{leancode}

\medskip

This specialization appears in
\href{https://github.com/qinz1yang/differential-geometry/blob/8bd406e35c33a200e9b88895cf11ee8429194e15/DifferentialGeometry/Geometry/Flow/RicciFlow/Solution/Basic.lean\#L320-L325}
{\texttt{Solution/Basic.lean}}.  Here \texttt{HasDerivWithinAt} expresses the
time derivative relative to the allowed time set \texttt{D.carrier}, while
\texttt{S.ricciAt} is metric-derived rather than independent data.  For a
closed--open interval $[a,b)$, the regular times in this predicate are the
interior times $(a,b)$.  The Lean development imposes the one-sided equation
at the included initial endpoint separately in its short-time-existence
theorem.  The predicate \texttt{IsSolutionOn} packages the interior-time
equation with the regularity needed later; it does not itself add the
one-sided endpoint equation.

\subsection{Short-time existence}

\begin{theorem}[Hamilton's short-time existence theorem]
\label{thm:short-time-existence}
Let $M$ be a closed smooth $3$-manifold, and let $g_{0}$ be a smooth
Riemannian metric on $M$.  Then there exist $T>0$ and a jointly smooth family
of Riemannian metrics
\[
  g(t),\qquad 0\leq t<T,
\]
such that
\[
  g(0)=g_{0},
  \qquad
  \frac{\partial}{\partial t}g(t)=-2\Ric_{g(t)}
  \quad\text{for }0\leq t<T.
\]
At $t=0$, both joint smoothness and the time derivative are understood from
the right.
\end{theorem}

\textcolor{red}{Since short-time existence is a major theorem, put the Lean code for the statement here.}

\subsection{Continuation and curvature blow-up}

\begin{theorem}[Hamilton's finite-time extension criterion]
\label{thm:bounded-curvature-extension}
Let $M$ be a closed smooth $3$-manifold, let $0<T<\infty$, and let $g(t)$,
$t\in[0,T)$, be a smooth Ricci flow.  Suppose that there exists $K\geq0$ such
that
\[
  \bigl|\Rm_{g(t)}(x)\bigr|_{g(t)}\leq K
\]
for every $x\in M$ and every $t\in[0,T)$.  Then there exist
$\varepsilon>0$ and a smooth Ricci flow
\[
  \widetilde g(t),\qquad t\in[0,T+\varepsilon),
\]
on $M$ such that $\widetilde g(t)=g(t)$ for every $t\in[0,T)$.
\end{theorem}

\textcolor{red}{Question for coauthors: The statement is true for $\varepsilon$ depending only on $K$. Have we Lean formalized a proof of this?}

\begin{corollary}[Curvature blow-up at a finite maximal time]
\label{cor:finite-time-curvature-blowup}
Let $M$ be a closed smooth $3$-manifold, let $0<T<\infty$, and let $g(t)$,
$t\in[0,T)$, be a smooth Ricci flow.  Suppose that $[0,T)$ is maximal in
forward time: there are no $\varepsilon>0$ and Ricci flow
$\widetilde g(t)$ on $[0,T+\varepsilon)$ for which
$\widetilde g(t)=g(t)$ on $[0,T)$.  Then
\[
  \sup_{0\leq t<T}\ \sup_{x\in M}
  \bigl|\Rm_{g(t)}(x)\bigr|_{g(t)}
  =+\infty.
\]
Equivalently, for every $K>0$ there exist a time $t\in[0,T)$ and a point
$x\in M$ for which
\[
  \lvert\Rm_{g(t)}(x)\rvert_{g(t)}>K.
\]
\end{corollary}

\textcolor{red}{Since this is a major theorem, put the Lean code for the statement here.}

Thus a closed Ricci flow can fail to continue through a finite time only if
its curvature becomes unbounded.  Surgery replaces selected high-curvature
regions in a controlled manner and restarts the smooth flow.

\section{A robust surgery interface}
\label{sec:surgery-interface}

\begin{remark}[Formalization priority]
This finitary, producer-independent interface is intended as one of the
first targets for Lean autoformalization.  It records exactly the data
consumed by backward topological reconstruction, while the quantitative
hypotheses and estimates that produce those data remain in the later
analytic chapters.  Formalizing the interface first will freeze the contract
between the surgery and topology developments; it will not by itself
formalize the existence of Perelman's surgery flow.
\end{remark}

All objects in the next definitions lie in the smooth category.  The time
slices are smooth oriented $3$-manifolds, the neck parametrizations and
survivor maps are smooth embeddings, the cuts and caps are smooth
constructions, and every asserted equivalence is a diffeomorphism.  The
Riemannian metrics are additional geometric data on these smooth objects.
A time slice may be disconnected, and the empty disjoint union is allowed as
the terminal time slice after extinction.  Orientations belong to the
transition data; discard predicates concern the underlying smooth manifolds
and do not remember those choices.

\subsection{Discarding}

At an event time, some components of the capped incoming slice may already
have recognized diffeomorphism types.  Once such a component has been
classified in a family sufficient for the final topological reconstruction,
continuing its metric evolution supplies no further information needed for
the Poincar\'e argument.  We therefore allow it to leave the active flow.
Here \emph{discarding} means only that the component is no longer evolved
analytically: the history still records the component and its classification
certificate, so that it reappears in the backward connected-sum
reconstruction.

\begin{definition}[Diffeomorphism-invariant discard predicate]
\label{def:discard-predicate}
A \emph{discard predicate} is a property $\mathcal D(N)$ of connected closed
smooth $3$-manifolds that is invariant under diffeomorphism: whenever $N$ and
$N'$ are such manifolds,
\[
  N\diffeo N'
  \quad\Longrightarrow\quad
  \bigl(\mathcal D(N)\Longleftrightarrow\mathcal D(N')\bigr).
\]
If $N$ carries an orientation, $\mathcal D(N)$ refers to its underlying
smooth manifold and is independent of that choice.  The structural
transition below records what is discarded without depending on a choice of
$\mathcal D$; control by a particular predicate is imposed afterward.
\end{definition}

\begin{remark}[Why a predicate rather than a category]
At a surgery event the consumer asks only whether each discarded component
$Q$ satisfies $\mathcal D(Q)$.  No morphism between two discardable
components, and no identity or composition law, is used.  A category would
therefore add structure absent from both the hypothesis and the
reconstruction proof.  A predicate also gives the direct Lean interface: a
proposition-valued classifier, together with the displayed transport law
under diffeomorphism, matches the componentwise condition in
Definition~\ref{def:controlled-transition}.  If a later development needs
maps between discardable manifolds, the corresponding full subcategory (or
groupoid of diffeomorphisms) can be recovered from the predicate without
changing this interface.
\end{remark}

\subsection{The topological cut--cap trace}

\begin{definition}[Spherical cut--cap transition]
\label{def:cut-cap-transition}
Let $M^{-}$ be a nonempty, possibly disconnected, closed oriented smooth
$3$-manifold, and let $M^{+}$ be a possibly empty, possibly disconnected,
closed oriented smooth $3$-manifold.  A \emph{spherical cut--cap transition}
from $M^{-}$ to $M^{+}$ consists of the following data.

\begin{enumerate}
\item A finite set $A$ and, for each $\alpha\in A$, a smooth tubular
embedding
\[
  \nu_{\alpha}\colon S^{2}\times[-2,2]\longrightarrow M^{-},
\]
such that the images of the embeddings are pairwise disjoint.

\item The \emph{cut-open core}
\[
  M_{\mathrm{cut}}
  =
  M^{-}\setminus
  \bigcup_{\alpha\in A}
  \nu_{\alpha}\bigl(S^{2}\times(-1,1)\bigr).
\]
We give this complement its induced smooth manifold-with-boundary structure.
It is compact, and every component of $\partial M_{\mathrm{cut}}$ is a
$2$-sphere.  For each boundary component $\Sigma$, the data include an
oriented $3$-ball $B_{\Sigma}$ and an orientation-reversing diffeomorphism
\[
  \varphi_{\Sigma}\colon\partial B_{\Sigma}\longrightarrow\Sigma
\]
with respect to the boundary orientations.  They also include a closed
oriented smooth $3$-manifold $M^{\sharp}$, a smooth embedding
\[
  j\colon M_{\mathrm{cut}}\longrightarrow M^{\sharp},
\]
and, for every $\Sigma$, a smooth embedding
$c_{\Sigma}\colon B_{\Sigma}\to M^{\sharp}$.  Their images cover
$M^{\sharp}$, their interiors are pairwise disjoint, and
\[
  c_{\Sigma}|_{\partial B_{\Sigma}}
  =
  j|_{\Sigma}\circ\varphi_{\Sigma}.
\]
Thus these data exhibit $M^{\sharp}$ as the manifold obtained by attaching
one $3$-ball cap to each boundary sphere of $M_{\mathrm{cut}}$.

\item A closed oriented smooth $3$-manifold $M^{\mathrm{disc}}$, possibly
empty.  The data also include an orientation-preserving diffeomorphism
\[
  F\colon M^{\sharp}\longrightarrow
  M^{+}\sqcup M^{\mathrm{disc}}.
\]
The two manifolds on the right record exhaustively the retained and discarded
components; no other component is created or deleted.
\end{enumerate}

Using the distinguished embedding $j$, define the \emph{retained core}
\[
  M_{\mathrm{ret}}
  =
  j^{-1}\bigl(F^{-1}(M^{+})\bigr).
\]
We give $M_{\mathrm{ret}}$ its induced structure as a compact
codimension-zero smooth submanifold with boundary of $M_{\mathrm{cut}}$.
Let $\iota^{-}$ be the composite inclusion
$M_{\mathrm{ret}}\hookrightarrow M_{\mathrm{cut}}\hookrightarrow M^{-}$, and
let $\iota^{+}$ be the restriction of $F\circ j$ to $M_{\mathrm{ret}}$.
Thus
\[
  \iota^{-}\colon M_{\mathrm{ret}}\longrightarrow M^{-},
  \qquad
  \iota^{+}\colon M_{\mathrm{ret}}\longrightarrow M^{+}.
\]
The gluing data have the following exhaustiveness properties, which are part
of the transition interface.
Every connected component of $M^{+}$ meets
$\iota^{+}(M_{\mathrm{ret}})$, and
\[
  M^{+}\setminus
  \operatorname{int}\bigl(\iota^{+}(M_{\mathrm{ret}})\bigr)
\]
is exactly the union of the sets
$F(c_{\Sigma}(B_{\Sigma}))$ that lie in $M^{+}$.  In particular,
$M_{\mathrm{ret}}=\varnothing$ if and only if $M^{+}=\varnothing$.

The transition is required to be nontrivial: either $A\neq\varnothing$ or
$M^{\mathrm{disc}}\neq\varnothing$.
\end{definition}

\begin{definition}[Controlled transition]
\label{def:controlled-transition}
Let $\mathcal D$ be a discard predicate.  A spherical cut--cap transition is
\emph{$\mathcal D$-controlled} if every connected component $Q$ of
$M^{\mathrm{disc}}$ satisfies $\mathcal D(Q)$.
\end{definition}

The tubular embeddings remember which two boundary spheres arose from the
same cut.  That finite pairing is the topological certificate needed to
reverse the transition.

Capping does not join distinct connected components of
$M_{\mathrm{cut}}$.  Consequently, every connected component of
$M^{\sharp}$ contains $j(C)$ for a unique connected component $C$ of
$M_{\mathrm{cut}}$.  A connected component $N$ of
$M^{+}\sqcup M^{\mathrm{disc}}$ is \emph{associated with} a connected
component $Q$ of $M^{-}$ if the component $C$ determined by
$F^{-1}(N)$ is contained in $Q$.

Throughout this section, finite connected sums are understood up to
diffeomorphism and parentheses are suppressed.  The connected sum of an
empty finite family is $S^{3}$, so adjoining zero copies of
$S^{2}\times S^{1}$ changes no diffeomorphism type.

\begin{proposition}[Local cut--cap reconstruction]
\label{prop:local-cut-cap-reconstruction}
For every connected component $Q$ of $M^{-}$ in a spherical cut--cap
transition, the components of $M^{+}\sqcup M^{\mathrm{disc}}$ associated
with $Q$ form a nonempty finite family.  If
$N_{1},\ldots,N_{r}$ is its complete list, then there is an integer $b\geq0$
such that
\[
  Q\diffeo
  N_{1}\mathbin{\#}\cdots\mathbin{\#}N_{r}
  \mathbin{\#}
  \mathop{\#}^{\,b}(S^{2}\times S^{1}).
\]
\end{proposition}

\begin{proof}
Fix $Q$ and form its finite incidence graph $G_Q$.  Its vertices are the
connected components of $M_{\mathrm{cut}}$ contained in $Q$.  For each
removed cylinder in $Q$, add an edge whose endpoints are the components
containing its two boundary spheres; the endpoints may coincide.  Reattaching
all the cylinders recovers $Q$, so $G_Q$ is connected.

Capping a vertex component produces precisely its associated component of
$M^{+}\sqcup M^{\mathrm{disc}}$.  Enumerate these components as
$N_1,\ldots,N_r$, where $r=\lvert V(G_Q)\rvert\geq1$.  To recover $Q$, remove
the interiors of the paired caps and restore the corresponding product
regions $S^{2}\times[-1,1]$.

Choose a spanning tree of $G_Q$.  Restoring the product regions
corresponding to its edges forms
$N_1\mathbin{\#}\cdots\mathbin{\#}N_r$.  Each remaining edge joins two
boundary spheres in an already connected oriented $3$-manifold and
contributes one $S^{2}\times S^{1}$ summand.  Hence
\[
  b=\lvert E(G_Q)\rvert-\lvert V(G_Q)\rvert+1
\]
has the required property.  The alternative twisted $S^{2}$-bundle over
$S^{1}$ is nonorientable and is excluded by the orientation of $Q$.  The
smooth extension theorem for diffeomorphisms of $S^{2}$ over $B^{3}$ shows
that no additional gluing invariant occurs.
\end{proof}

\subsection{Terminal regular regions and metric surgery events}

\begin{definition}[Terminal regular region]
\label{def:terminal-regular-region}
Let $M$ be a closed smooth $3$-manifold, let $a<s<\infty$, and let $g(t)$,
$a\leq t<s$, be a Ricci flow.  Its \emph{terminal regular region} at $s$ is
\[
\begin{split}
  \Omega(g,s)=\bigl\{x\in M:\;&\text{there exist an open set }U\ni x,\
  a'\in[a,s),\text{ and }K\geq0\\
  &\text{such that }
  \lvert\Rm_{g(t)}(y)\rvert_{g(t)}\leq K\\
  &\text{for every }y\in U\text{ and }t\in[a',s)\bigr\}.
\end{split}
\]
\end{definition}

\begin{lemma}[Openness of the terminal regular region]
\label{lem:terminal-regular-region-open}
The subset $\Omega(g,s)$ is open in $M$.
\end{lemma}

\begin{proof}
If $x\in\Omega(g,s)$ and $U,a',K$ witness membership in the definition, then
the same data witness that every point of $U$ belongs to $\Omega(g,s)$.
Thus $U\subseteq\Omega(g,s)$.
\end{proof}

We equip $\Omega(g,s)$ with its inherited smooth structure.

\begin{definition}[Terminal limit metric]
\label{def:terminal-limit-metric}
A \emph{terminal limit metric} is a smooth Riemannian metric $\bar g$ on
$\Omega(g,s)$ such that
\[
  g(t)\longrightarrow\bar g
  \quad\text{in }C^{\infty}_{\mathrm{loc}}\bigl(\Omega(g,s)\bigr)
  \quad\text{as }t\uparrow s.
\]
\end{definition}

\begin{definition}[Metric cut--cap event]
\label{def:metric-cut-cap-event}
Suppose that $g^{-}(t)$ is a Ricci flow on $M^{-}$ for $a\leq t<s$, that
$\bar g^{-}$ is a terminal limit metric on $\Omega(g^{-},s)$, and that
$h^{+}$ is a Riemannian metric on a possibly empty manifold $M^{+}$, with
the unique empty metric used when $M^{+}=\varnothing$.

A \emph{metric cut--cap event at time $s$} consists of a spherical cut--cap
transition from $M^{-}$ to $M^{+}$ and a compact codimension-zero smooth
submanifold with boundary
\[
  K_{\mathrm{old}}\subseteq M_{\mathrm{ret}},
\]
called the \emph{unchanged locus}, such that
\[
  \iota^{-}(K_{\mathrm{old}})\subseteq\Omega(g^{-},s)
\]
and
\[
  (\iota^{-}|_{K_{\mathrm{old}}})^{*}\bar g^{-}
  =
  (\iota^{+}|_{K_{\mathrm{old}}})^{*}h^{+}.
\]
The unchanged locus contains every retained point outside the certified
surgery neighborhoods:
\[
  M_{\mathrm{ret}}\setminus
  (\iota^{-})^{-1}\!\left(
    \bigcup_{\alpha\in A}
    \nu_{\alpha}\bigl(S^{2}\times(-2,2)\bigr)
  \right)
  \subseteq K_{\mathrm{old}}.
\]
If $M^{+}\neq\varnothing$, every connected component of $M^{+}$ meets
$\iota^{+}(K_{\mathrm{old}})$; if $M^{+}=\varnothing$, then
$K_{\mathrm{old}}=\varnothing$.
\end{definition}

The unchanged locus separates two logically different assertions.  The
cut--cap trace records all topological changes, while the pullback identity
records the region on which the old and new metrics agree exactly.  This does
not require the interpolated cap collars to retain the old metric.

\subsection{Finite Ricci-flow-with-surgery histories}

For the Poincar\'e proof only a finite history up to a finite extinction time
is needed.  We therefore take finiteness as part of the primary interface and
do not introduce countably infinite histories.

\begin{definition}[Finite Ricci flow with surgery]
\label{def:finite-surgery-history}
A \emph{finite Ricci-flow-with-surgery history} consists of the following
data.

\begin{enumerate}
\item An integer $N\geq1$ and break times
\[
  0=t_{0}<t_{1}<\cdots<t_{N}<\infty.
\]
The positive times $t_{1},\ldots,t_{N}$ are the event times; $t_{0}=0$ is
only the initial time.

\item Closed oriented smooth $3$-manifolds
$M_{0},\ldots,M_{N}$, each possibly disconnected, with
$M_{i}\neq\varnothing$ for $i<N$.  The final slice $M_{N}$ may be empty.
Each $M_i$ carries a Riemannian metric $h_i$; on the empty slice this means
the unique empty metric.

\item For every $0\leq i<N$, a Ricci flow $g_i(t)$ on $M_i$ for
$t_i\leq t<t_{i+1}$, jointly smooth from the right at $t_i$, such that
\[
  g_i(t_i)=h_i.
\]

\item For every $0\leq i<N$, a terminal limit metric $\bar g_i$ on
$\Omega(g_i,t_{i+1})$ and a metric cut--cap event at time $t_{i+1}$ from
the incoming slab $(M_i,g_i)$ to the outgoing slice
$(M_{i+1},h_{i+1})$.
\end{enumerate}

The history \emph{starts from} a closed oriented Riemannian $3$-manifold
$(M,g_{\mathrm{in}})$ if its data include a specified
orientation-preserving diffeomorphism
\[
  \Phi_{0}\colon M\longrightarrow M_{0}
\]
such that $\Phi_{0}^{*}h_{0}=g_{\mathrm{in}}$.
It \emph{becomes extinct at time $T$} if $T=t_N$ and
$M_N=\varnothing$.  The finite history terminates at this empty slice.  When
it is convenient to speak of a post-extinction continuation, it is defined
by declaring every slice at $t\geq T$ to be empty; no component may reappear.
\end{definition}

\begin{definition}[Controlled history]
\label{def:controlled-history}
Let $\mathcal D$ be a discard predicate.  A finite Ricci-flow-with-surgery
history is \emph{$\mathcal D$-controlled} if every one of its cut--cap
transitions is $\mathcal D$-controlled.
\end{definition}

\begin{remark}[Relation with Perelman's surgery]
\label{rem:classical-surgery}
The preceding definition is an axiomatization, not a verbatim reproduction
of Perelman's cutoff algorithm.  The Ricci-flow slabs, terminal regular
regions, limit metrics, and isometric survivor regions abstract
Perelman~\cite[\S\S3 and~4.1]{Perelman2}; the spherical cuts, caps, and discarded
components abstract \cite[\S4.4]{Perelman2}.  The closest systematic
formulations are Kleiner--Lott
\cite[Definition~68.1]{MR2460872} for the slab-and-survivor structure and
\cite[Definition~73.1]{MR2460872} for the cutoff operation.

Our explicit finite cut--cap trace records the data needed for reverse
topological reconstruction.  The intended producer is Perelman's
sufficiently precise $\delta$-cutoff construction; later chapters will define
that producer notion precisely and prove that its output satisfies this
interface.  The converse is false: the interface intentionally omits the
quantitative neck and cap quality, pinching, noncollapsing,
canonical-neighborhood estimates, and surgery-scale hierarchies that
distinguish Perelman's construction.
\end{remark}

\begin{roadmap}{After fixing an exact definition of a Perelman
$\delta$-cutoff flow, the surgery chapters will prove that its output realizes
this interface.}
The theorem will state that if such a flow starts at time zero and becomes
extinct at a finite time $T$, then its slabs and events on $[0,T]$ determine a
finite Ricci-flow-with-surgery history in the sense of
Definition~\ref{def:finite-surgery-history}.  Its spherical cut--cap trace
will be exhaustive, its discarded components will have the topological types
recorded in Perelman's surgery analysis, and its unchanged loci will be
identified isometrically across event times.  In particular, the resulting
history will be controlled by the Poincar\'e predicate defined below.  Until
the producer notion is defined, this is a roadmap obligation rather than a
frozen theorem interface; it may not be used as a substitute for the
estimates needed to construct the flow.
\end{roadmap}

\subsection{The fixed Poincar\'e control predicate}

\begin{definition}[Spherical space form]
\label{def:spherical-space-form}
A \emph{spherical space form} is a quotient $S^{3}/\Gamma$, where
\[
  \Gamma\leq\Isom^{+}(S^{3},g_{\mathrm{round}})
\]
is finite and acts freely on $S^{3}$.  Thus the action is, in particular, by
orientation-preserving isometries of the round sphere.
\end{definition}

\begin{definition}[Poincar\'e-standard manifold]
\label{def:poincare-standard}
A connected closed smooth $3$-manifold $N$ is
\emph{Poincar\'e-standard} if
\[
  N\diffeo
  \mathop{\#}_{j=1}^{a}(S^{3}/\Gamma_j)
  \mathbin{\#}
  \mathop{\#}^{\,b}(S^{2}\times S^{1})
\]
for some $a,b\in\mathbb N$, where each $\Gamma_j$ is as in
Definition~\ref{def:spherical-space-form}.  The iterated connected sum is
taken over the $a+b$ displayed factors; a one-factor sum is that factor, and
the sum with no factors is understood to be $S^{3}$.
\end{definition}

\textcolor{red}{Since this is a major definition, put the Lean code for the statement here.}

By the convention fixed in the introduction, connectedness already includes
nonemptiness, so no separate hypothesis is needed here.  Orientability is
not an additional hypothesis: every manifold on the right-hand side is
orientable, so any Poincar\'e-standard manifold is automatically orientable.
Components discarded from the oriented time slices above also inherit
orientations.  The property just defined, however, depends only on their
underlying diffeomorphism types.

\begin{definition}[Poincar\'e discard property]
\label{def:poincare-discard-property}
For a connected closed smooth $3$-manifold $N$, define
\[
  \mathcal D_{\mathrm{PC}}(N)
  \quad\Longleftrightarrow\quad
  N\text{ is Poincar\'e-standard}.
\]
\end{definition}

\begin{lemma}[Closure of Poincar\'e-standard manifolds]
\label{lem:poincare-standard-closure}
The property $\mathcal D_{\mathrm{PC}}$ is a discard predicate.  Every
spherical space form and $S^{2}\times S^{1}$ is Poincar\'e-standard.
Moreover, the connected sum of any finite family of Poincar\'e-standard
manifolds is Poincar\'e-standard, with the empty connected sum interpreted as
$S^{3}$.
\end{lemma}

\begin{proof}
Diffeomorphism invariance follows immediately from the definition.  A
spherical space form has a presentation with $(a,b)=(1,0)$,
$S^{2}\times S^{1}$ has one with $(a,b)=(0,1)$, and $S^{3}$ has the empty
presentation.  For a finite connected sum of Poincar\'e-standard manifolds,
concatenate their lists of spherical factors and add their numbers of
$S^{2}\times S^{1}$ factors.  Associativity and commutativity of connected
sum, up to diffeomorphism, put the resulting presentation in the displayed
form of Definition~\ref{def:poincare-standard}.
\end{proof}

\begin{definition}[Poincar\'e-controlled history]
\label{def:poincare-control}
A \emph{Poincar\'e-controlled Ricci-flow-with-surgery history} is a
$\mathcal D_{\mathrm{PC}}$-controlled history.
\end{definition}

The predicate $\mathcal D_{\mathrm{PC}}$ is fixed once and for all; it is not
a caller-supplied hypothesis.  In particular,
$\mathbb{RP}^{3}\#\mathbb{RP}^{3}$ is Poincar\'e-standard.  The breadth and
closure recorded in Lemma~\ref{lem:poincare-standard-closure} are exactly
what the backward reconstruction uses.

\section{Extinction and reconstruction}

The analytic chapters establish the following project-facing form of
Perelman's construction and finite-extinction theorem.

\begin{theorem}[Perelman surgery and finite extinction: Poincar\'e case]
\label{thm:finite-extinction}
Let $M$ be a closed, connected, simply connected smooth $3$-manifold, and let
$g_{\mathrm{in}}$ be a smooth Riemannian metric on $M$.  After choosing an
orientation of $M$, there exist a Poincar\'e-controlled
Ricci-flow-with-surgery history starting from $(M,g_{\mathrm{in}})$ and a
time $T>0$ at which the history becomes extinct.
\end{theorem}

\textcolor{red}{Since this is a major theorem, put the Lean code for the statement here.}

A connected simply connected manifold is orientable: otherwise its
orientation double cover would be a nontrivial connected two-sheeted cover.
Thus the choice of orientation in the theorem is always available.  The
existence portion of Theorem~\ref{thm:finite-extinction} comes from
Perelman's surgery construction \cite[\S6.1]{Perelman2}; see also
Kleiner--Lott \cite[Proposition~77.2]{MR2460872}.  Those existence statements
use a normalized initial metric.  Because $M$ is closed, multiplying
$g_{\mathrm{in}}$ by a sufficiently large constant gives such a
normalization: curvature decreases under this rescaling, and the required
unit-ball volume bound follows uniformly from the small-ball asymptotics on
the compact manifold.  Parabolic rescaling of the resulting cutoff flow
returns a history starting from $g_{\mathrm{in}}$.

The cutoff construction has no accumulation of surgery times in a bounded
time interval.  Thus extinction at a finite time produces only finitely many
slabs and events, as required by
Definition~\ref{def:finite-surgery-history}.  The extinction conclusion is
the simply connected case of \cite[Theorem~1.1]{Perelman3}; see
\cite[\S1.4]{Perelman3} for the homotopy-sphere specialization.  Perelman's
theorem is stronger: it treats closed oriented $3$-manifolds whose prime
decompositions contain no aspherical factors.  We record only the case
consumed by this book.

The topological lane has the following exact output.

\begin{theorem}[Extinction reconstruction]
\label{thm:extinction-reconstruction}
Let $M$ be a connected closed oriented smooth $3$-manifold, and let
$g_{\mathrm{in}}$ be a Riemannian metric on $M$.  Suppose that a finite Ricci
flow with surgery is Poincar\'e-controlled, starts from
$(M,g_{\mathrm{in}})$, and becomes extinct in finite time.  Then $M$ is
Poincar\'e-standard.
\end{theorem}

\textcolor{red}{Since this is a major theorem, put the Lean code for the statement here.}

\begin{proof}
There are finitely many events by
Definition~\ref{def:finite-surgery-history}.  Work backward from the empty
final slice.  Proposition~\ref{prop:local-cut-cap-reconstruction} expresses
each component before the last event as a connected sum of components after
the event, discarded components, and copies of $S^{2}\times S^{1}$.  The
outgoing list is empty, and every discarded component is
Poincar\'e-standard.  Proposition~\ref{prop:local-cut-cap-reconstruction}
and Lemma~\ref{lem:poincare-standard-closure} therefore show that every
component before the last event is Poincar\'e-standard.  Repeating the same
argument at the preceding events proves that every component of $M_0$ is
Poincar\'e-standard.  The specified initial diffeomorphism identifies the
connected manifold $M$ with the unique component of $M_0$.
\end{proof}

This is the abstract form of the reverse-surgery topology in
Kleiner--Lott \cite[Lemmas~73.4 and~81.2]{MR2460872}.

\begin{lemma}[The simply connected case]
\label{lem:simply-connected-standard}
Every connected, simply connected, Poincar\'e-standard smooth
$3$-manifold is diffeomorphic to $S^{3}$.
\end{lemma}

\textcolor{red}{Since this is a major lemma, put the Lean code for the statement here.}

\begin{proof}
For each spherical factor, the quotient map
$S^{3}\to S^{3}/\Gamma_j$ is a universal covering, and hence
$\pi_1(S^{3}/\Gamma_j)\cong\Gamma_j$.  For a presentation as in
Definition~\ref{def:poincare-standard}, the Seifert--van Kampen theorem then
gives
\[
  \pi_{1}(M)
  \cong
  \left(\mathop{\ast}_{j=1}^{a}\Gamma_j\right)
  \ast
  \left(\mathop{\ast}_{k=1}^{b}\mathbb Z\right).
\]
Here an empty free product is the trivial group.
A free product is trivial only if every factor is trivial.  Simple
connectivity therefore forces $b=0$ and every $\Gamma_j$ to be trivial.  All
remaining factors are copies of $S^{3}$, and a finite connected sum of copies
of $S^{3}$ is diffeomorphic to $S^{3}$.
\end{proof}

\section{Completion of the proof}

\begin{proof}[Proof of Theorem~\ref{thm:smooth-poincare}]
Let $M$ be a closed, connected, simply connected smooth $3$-manifold.  Choose
a smooth Riemannian metric $g_{\mathrm{in}}$ on $M$ and an orientation.
Theorem~\ref{thm:finite-extinction} supplies a Poincar\'e-controlled surgery
history starting from $(M,g_{\mathrm{in}})$ and becoming extinct in finite
time.  Theorem~\ref{thm:extinction-reconstruction} shows that $M$ is
Poincar\'e-standard.  Lemma~\ref{lem:simply-connected-standard} now gives
$M\diffeo S^{3}$.
\end{proof}

\begin{proof}[Proof of Theorem~\ref{thm:topological-poincare}]
Let $M$ be a closed, connected, simply connected topological $3$-manifold.
By Theorem~\ref{thm:moise-smoothability}, the set $M$ admits a smooth
$3$-manifold structure whose underlying topology is its original topology.
The resulting smooth manifold remains closed, connected, and simply
connected.  Theorem~\ref{thm:smooth-poincare} gives a diffeomorphism from it
to $S^{3}$.  Forgetting the smooth structures gives a homeomorphism from the
original topological manifold $M$ to $S^{3}$.
\end{proof}

\section{Guide to the remaining chapters}

The rest of the book proves the nodes used above.  The Ricci-flow chapters
establish short-time existence, the extension criterion, noncollapsing,
singularity models, and canonical neighborhoods.  The surgery chapters
define sufficiently precise cutoff flows, construct them, and prove the
producer-to-interface theorem described in
Section~\ref{sec:surgery-interface}.  The extinction chapters prove
Theorem~\ref{thm:finite-extinction}.  The topology chapters develop the
sphere-cutting and connected-sum bookkeeping underlying
Proposition~\ref{prop:local-cut-cap-reconstruction} and
Theorem~\ref{thm:extinction-reconstruction}.

The next chapter begins the analytic development from the bottom up.  It
proves Hamilton's weak maximum principle for systems on a metric vector
bundle, in a form designed both for Uhlenbeck's trick and for later Lean
formalization.  Subsequent chapters will specialize and extend that theorem
to the curvature cones and moving barriers needed in Ricci flow.

Chapters~\ref{ch:hamilton-ivey}--\ref{ch:curvature-evolution} carry out the
first specialization.  They prove the time-improved Hamilton--Ivey estimate
and then discharge its sole curvature-evolution input by deriving the exact
fixed-bundle equation in Uhlenbeck gauge.  Chapter~\ref{ch:strong-curvature}
develops the strong maximum principle needed to turn
nonnegative curvature into the positive-curvature-versus-splitting
alternative used in the later analysis of ancient solutions.
Chapter~\ref{ch:gradient-solitons} then develops gradient Ricci solitons and
their canonical self-similar flows in arbitrary dimension.  It proves the
Chow--Lu--Yang inverse-quadratic scalar-curvature lower bound for complete
noncompact non-Gaussian shrinkers and the stronger Munteanu--Wang
compactness theorem for shrinkers with nonnegative sectional curvature and
positive Ricci curvature.
Chapter~\ref{ch:three-dimensional-shrinkers} combines these ingredients to
classify every complete three-dimensional gradient shrinker, with no
curvature bound or noncollapsing hypothesis.  It separates the localized
Hamilton--Ivey input needed for that general theorem from the cheaper
Poincar\'e-project route, where nonnegative curvature passes directly from
the closed approximating flows to their smooth blow-up limits.
Chapter~\ref{ch:bando-shi} then proves Shi's local curvature-derivative
estimates in every dimension, using the exact Uhlenbeck calculus already
developed in Chapter~\ref{ch:curvature-evolution}.  Besides supplying the
regularity input for Hamilton compactness, it gives the complete proof of
the bounded-curvature extension criterion stated above.
Chapter~\ref{ch:matrix-harnack} uses that positive-time regularity and the
fixed-bundle calculus to prove Hamilton's matrix Harnack estimate in every
dimension.  Its shifted-time formulation simultaneously yields the usual
finite-origin theorem, the trace Harnack estimate, and the ancient matrix
and trace limits.
Chapter~\ref{ch:reduced-geometry} then develops Perelman's spacetime reduced
geometry from the clocked action through weak reduced-volume monotonicity.
Its corrected moving-center tightness argument produces asymptotic
shrinkers of ancient $\kappa$-solutions and records the domain and crossing
interfaces required later for Ricci flow with surgery.

This chapter fixes interfaces and dependencies; it does not claim that the
deep analytic inputs have already been proved here.  The blueprint refines
each macro milestone into stable fine nodes, and the status ledger records
which of those nodes have source-complete proofs, approved Lean interfaces,
compiled implementations, and independent audit evidence.

\chapter{Hamilton's Weak Maximum Principle for Systems}
\label{ch:weak-maximum-principle-systems}
\index{maximum principle!for systems}
\index{Hamilton!weak maximum principle}

\section{Purpose and scope}

Hamilton's weak maximum principle for systems is the bridge from an
ordinary differential equation in one fiber to a parabolic differential
equation on a manifold.  Its geometric content is especially clean after
Uhlenbeck's trick: the vector bundle and its fiber metric are fixed, whereas
the metric on the base and the metric connection on the bundle may vary with
time.  The theorem then says that diffusion cannot drive a section out of a
closed convex set which parallel transport cannot distinguish from one fiber
to another, provided the fiberwise reaction ODE cannot do so.

This chapter proves precisely that fixed-set theorem.  The base is closed,
the fiber metric is time independent, the base metric is an arbitrary smooth
one-parameter family, and the compatible bundle connection may depend on
time.  Ricci flow is therefore an application, not a hypothesis.  Neither
orientability nor connectedness is needed.

The proof is organized for later formalization.  It uses a strictly expanding
tube around the invariant set.  At a hypothetical first contact, a nearest
point and its supporting normal turn the bundle-valued equation into a scalar
inequality.  Spatial information is obtained one geodesic at a time.  This
avoids differentiating the distance to a convex set, avoids choosing a smooth
nearest-point field, and avoids asserting that a radially parallel section
has vanishing full second covariant derivative.  Only the diagonal second
derivative along a single geodesic is used.

\begin{roadmap}{Section~\ref{sec:wmp-setup} fixes the geometric data and
states the ODE-to-PDE theorem.}
Sections~\ref{sec:wmp-convex} and
\ref{sec:wmp-geodesic} isolate the finite-dimensional convex lemma and the
geodesic scalarization lemma.  Section~\ref{sec:wmp-proof} gives the strict
tube proof.  Sections~\ref{sec:wmp-variants} and
\ref{sec:wmp-uhlenbeck} record the support, associated-bundle, and Uhlenbeck
interfaces.  Section~\ref{sec:wmp-formalization} gives the dependency order
for Lean.
\end{roadmap}

\section{Geometric setup and theorem}
\label{sec:wmp-setup}

Let $M$ be a closed smooth manifold, where \emph{closed} means compact and
without boundary.  Fix a finite real number $T>0$.  A solution on
$[0,\infty)$ is covered by applying the theorem on every finite interval
$[0,T)$.  Let
\[
   \pi:E\longrightarrow M
\]
be a smooth real vector bundle of finite rank.  We fix once and for all a
smooth positive-definite fiber metric $h$ on $E$.  Thus the norm
$|\,\cdot\,|_h$ does not depend on time.

For $t\in[0,T)$, let $g(t)$ be a Riemannian metric on $M$, let
$\nabla^{E,t}$ be an $h$-compatible connection on $E$, and let $X(t)$ be a
smooth time-dependent vector field on $M$.  All three families are assumed
smooth in space and time.  We denote the Levi--Civita connection of
$g(t)$ by $\nabla^{g(t)}$.  For a smooth section $\sigma$ of $E$, its second
covariant derivative is the tensor
\[
  (\nabla^{E,t})^2\sigma
     \in\Gamma(T^*M\otimes T^*M\otimes E)
\]
defined by
\begin{align}
 \bigl((\nabla^{E,t})^2\sigma\bigr)(X,Y)
   &=
   \nabla^{E,t}_X\bigl(\nabla^{E,t}_Y\sigma\bigr)
   -\nabla^{E,t}_{\nabla^{g(t)}_X Y}\sigma,                 \label{eq:wmp-hessian}\\
 \Delta_t^E\sigma
   &=\trace_{g(t)}\bigl((\nabla^{E,t})^2\sigma\bigr).
                                                               \label{eq:wmp-laplacian}
\end{align}
Here $X$ and $Y$ may be any smooth local extensions of the tangent vectors
at the point in question; the right side of \eqref{eq:wmp-hessian} is
independent of those extensions.  Equivalently, it is the covariant
derivative of $\nabla^{E,t}\sigma\in\Gamma(T^*M\otimes E)$ for the connection
induced by $\nabla^{g(t)}$ and $\nabla^{E,t}$.
Our sign convention is the one for which $\Delta f\leq0$ at a local maximum
of a smooth real-valued function $f$.

Let $\operatorname{pr}_M:M\times[0,T)\to M$ be projection.  By a smooth
time-dependent section of $E$ we mean a smooth section of
$\operatorname{pr}_M^*E\to M\times[0,T)$, equivalently a smooth map
$u:M\times[0,T)\to E$ satisfying $\pi\circ u=\operatorname{pr}_M$.
Smoothness at $t=0$ is understood in the standard
manifold-with-boundary sense, hence locally by a smooth extension in the
time variable.  Because $E$ itself is fixed, the curve $t\mapsto u(x,t)$
lies in the one vector space $E_x$ and has an ordinary vertical derivative
$\partial_tu(x,t)\in E_x$.  No connection in the time direction is being
silently introduced.

Let
\[
  F:E\times[0,T)\longrightarrow E
\]
be a smooth map satisfying
\[
   \pi\circ F=\pi\circ\operatorname{pr}_E,
\]
where $\operatorname{pr}_E:E\times[0,T)\to E$ is projection.  Thus $F$ is
vertical.  For a fixed $x\in M$, write $F_x(v,t)=F(v,t)$ for
$v\in E_x$.  The associated fiber ODE is
\begin{equation}
   \dot a(t)=F(a(t),t),\qquad a(t)\in E_x.                 \label{eq:wmp-ode}
\end{equation}
For every $t_0\in[0,T)$ and $a_0\in E_x$, finite-dimensional
Picard--Lindel\"of theory gives a unique maximal solution on a maximal
interval $I(x,t_0,a_0)\subset[0,T)$ containing $t_0$.  Only its forward
part $I(x,t_0,a_0)\cap[t_0,T)$ will be used.

\begin{definition}[Parallel closed convex family]
\label{def:wmp-parallel-convex}
A subset $K\subset E$ is a \emph{parallel closed convex family} for
$\{\nabla^{E,t}\}$ if the following conditions hold.
\begin{enumerate}
\item For every $x\in M$, the fiber
      $K_x=K\cap E_x$ is nonempty, closed, and convex.
\item The subset $K$ is closed in the total space $E$.
\item For every $t\in[0,T)$, every piecewise smooth path
      $\gamma:[0,1]\to M$, and parallel transport
      $P_{\gamma,t}:E_{\gamma(0)}\to E_{\gamma(1)}$ for
      $\nabla^{E,t}$,
      \[
          P_{\gamma,t}(K_{\gamma(0)})=K_{\gamma(1)}.
      \]
\end{enumerate}
It would be enough in the last condition to require inclusion: transport
along the reversed path gives the reverse inclusion.
\end{definition}

\begin{remark}[The chosen representation of the family]
\label{rem:wmp-total-set}
The public object in this chapter is the single total subset $K\subset E$;
the dependent fiber sets are always defined by $K_x=K\cap E_x$.  Some
sources do not separately assume that every $K_x$ is nonempty, because the
initial condition $u(x,0)\in K_x$ implies it in the theorem.  We include
nonemptiness in the reusable definition so that the metric projection is
available before a particular solution $u$ has been introduced.
\end{remark}

\begin{theorem}[Hamilton's weak maximum principle for systems]
\label{thm:hamilton-wmp-systems}
Assume the geometric setup above, and let $K\subset E$ be a parallel closed
convex family.  Suppose that the fiber ODE has the following invariance
property:
\begin{quote}
for every $x\in M$, every starting time $t_0\in[0,T)$, and every
$a_0\in K_x$, the maximal solution of \eqref{eq:wmp-ode} with
$a(t_0)=a_0$ satisfies
\[
    a(t)\in K_x
    \qquad\text{for every }
    t\in I(x,t_0,a_0)\cap[t_0,T).
\]
\end{quote}
If $u$ is a smooth section of
$\operatorname{pr}_M^*E\to M\times[0,T)$ and solves
\begin{equation}
     \partial_tu
       =\Delta_t^E u+\nabla^{E,t}_{X(t)}u+F(u,t)            \label{eq:wmp-pde}
\end{equation}
and $u(x,0)\in K_x$ for every $x\in M$, then
\[
     u(x,t)\in K_x
     \qquad\text{for all }(x,t)\in M\times[0,T).
\]
\end{theorem}

\textcolor{red}{Since this is a major theorem, put the Lean code for the statement here.}

\begin{remark}[Regularity actually used]
\label{rem:wmp-regularity}
Smoothness is the convenient interface for this book and covers the
curvature reactions used in Ricci flow.  The proof uses only enough
regularity to make the displayed PDE classical, to have continuous
fiberwise distances, and to obtain a fiberwise Lipschitz constant for $F$ on
each bounded spacetime region under consideration.  In particular, global
Lipschitz continuity of $F$ on the generally noncompact total space $E$ is
not required.
\end{remark}

\begin{remark}[Why every starting time is stated]
\label{rem:wmp-starting-time}
For an autonomous reaction, invariance from time $0$ can be shifted to any
starting time.  That implication is not automatic for a time-dependent
reaction.  The quantifier over $t_0$ is therefore part of the theorem, rather
than an implicit convention.
\end{remark}

\section{Convex geometry in one fiber}
\label{sec:wmp-convex}

All convex arguments needed below occur in a finite-dimensional inner-product
space.  We record them in a form which can be imported without any bundle
geometry.

\begin{definition}[Normal cone]
\label{def:wmp-normal-cone}
If $C$ is a nonempty closed convex subset of a finite-dimensional
inner-product space $V$ and $p\in C$, define
\[
  N_C(p)=
  \bigl\{\nu\in V:\langle\nu,c-p\rangle\leq0
                  \text{ for every }c\in C\bigr\}.
\]
Thus $N_C(p)$ consists of the outward supporting normals at $p$.  This is a
vector cone based at the origin; it should not be confused with an affine
tangent set based at $p$.
\end{definition}

\begin{lemma}[Metric projection and support]
\label{lem:wmp-projection-support}
Let $C\subset V$ be nonempty, closed, and convex.  For every $z\in V$ there
is a unique point $p=\operatorname{proj}_C(z)\in C$ satisfying
\[
     |z-p|=\dist(z,C).
\]
Moreover,
\begin{equation}
     \langle z-p,c-p\rangle\leq0
     \qquad(c\in C).                                      \label{eq:wmp-projection-normal}
\end{equation}
If $z\notin C$, put $r=|z-p|$ and $\nu=(z-p)/r$.  Then
$|\nu|=1$, $\nu\in N_C(p)$, and for every $y\in V$,
\begin{equation}
     \langle\nu,y-p\rangle\leq\dist(y,C).                  \label{eq:wmp-support-distance}
\end{equation}
\end{lemma}

\begin{proof}
Choose $c_0\in C$.  A minimizing sequence may be restricted to the
intersection of $C$ with a closed ball centered at $z$ of radius
$|z-c_0|+1$.  That intersection is compact, so the distance is attained.

If $p,q\in C$ are both nearest points and the common distance is $r$, then
their midpoint belongs to $C$, while the parallelogram identity gives
\[
 \left|z-\frac{p+q}{2}\right|^2
   =\frac{|z-p|^2+|z-q|^2}{2}-\frac{|p-q|^2}{4}
   =r^2-\frac{|p-q|^2}{4}.
\]
Minimality forces $p=q$.

For $c\in C$ and $0\leq s\leq1$, the point $p+s(c-p)$ lies in $C$.
The function
\[
   s\longmapsto |z-p-s(c-p)|^2
\]
has a one-sided minimum at $s=0$.  Its right derivative is nonnegative,
which is exactly \eqref{eq:wmp-projection-normal}.  Finally, for arbitrary
$c\in C$,
\[
 \langle\nu,y-p\rangle
   =\langle\nu,y-c\rangle+\langle\nu,c-p\rangle
   \leq |y-c|.
\]
Taking the infimum over $c\in C$ proves
\eqref{eq:wmp-support-distance}.
\end{proof}

\begin{lemma}[ODE invariance gives the inward support inequality]
\label{lem:wmp-ode-support}
Fix $x\in M$ and $t_0\in[0,T)$.  Suppose the solution of the fiber ODE
starting at $p\in K_x$ at time $t_0$ remains in $K_x$ for a nontrivial
right-hand time interval.  Then
\begin{equation}
     h\bigl(\nu,F(p,t_0)\bigr)\leq0
     \qquad\text{for every }\nu\in N_{K_x}(p).             \label{eq:wmp-inward-support}
\end{equation}
\end{lemma}

\begin{proof}
Let $a(t)$ be the indicated solution.  For every sufficiently small $s>0$,
the inclusion $a(t_0+s)\in K_x$ and the definition of the normal cone give
\[
 h\bigl(\nu,a(t_0+s)-p\bigr)\leq0.
\]
Divide by $s$ and let $s\downarrow0$.  Since $a(t_0)=p$ and
$\dot a(t_0)=F(p,t_0)$, this yields
\eqref{eq:wmp-inward-support}.
\end{proof}

\begin{remark}[The exact ODE-to-PDE interface]
\label{rem:wmp-normal-interface}
The parabolic proof uses the ODE hypothesis only through
\eqref{eq:wmp-inward-support}.  Hamilton's formulation in terms of invariant
ODE solutions is nevertheless the most useful application interface: in
curvature problems the reaction equation is computed first, and preservation
of a cone by that finite-dimensional ODE is then checked algebraically.
\end{remark}

\section{Parallel transport and the spatial support test}
\label{sec:wmp-geodesic}

For $v\in E_x$, write
\[
      d_K(v)=\dist_h(v,K_x).
\]
The next elementary point is what makes a first-contact argument legitimate
even though the sets $K_x$ live in different vector spaces.

\begin{lemma}[Local product structure and continuity of distance]
\label{lem:wmp-distance-continuous}
If $K$ is a parallel closed convex family, then $d_K:E\to[0,\infty)$ is
continuous.  More precisely, near any $x_0\in M$ one may use parallel
transport for any one fixed connection $\nabla^{E,t_*}$ to identify
$E$ isometrically with a product in such a way that $K$ becomes the product
with the single closed convex set $K_{x_0}$.
\end{lemma}

\begin{proof}
Choose a coordinate ball $U$ centered at $x_0$ which is star-shaped in its
coordinate chart.  Let
\[
   \Gamma:U\times[0,1]\longrightarrow U
\]
be the corresponding smooth contraction, so that
$\Gamma(x,0)=x_0$ and $\Gamma(x,1)=x$.  For fixed $t_*$, set
\[
  Q_x=P_{\Gamma(x,\,\cdot\,),t_*}:E_{x_0}\longrightarrow E_x.
\]
Smooth parameter dependence for the linear parallel-transport ODE shows
that $(x,v)\mapsto Q_xv$ is smooth.  Transport along the reversed path gives
its smooth inverse.  Hence the maps $Q_x$ give a smooth trivialization
\[
       E|_U\cong U\times E_{x_0}.
\]
Compatibility with $h$ makes every identifying map an isometry.  Parallel
invariance of $K$, applied also to each reversed path, identifies every
$K_x$ exactly with $K_{x_0}$.  In this trivialization $d_K$ is the ordinary
distance to the fixed closed set $K_{x_0}$ in a Euclidean space, and is
therefore continuous.
\end{proof}

\begin{lemma}[One-dimensional geodesic support test]
\label{lem:wmp-geodesic-support}
Fix $t\in[0,T)$ and $r>0$.  Let $\sigma$ be a smooth section of $E$ such that
\[
       d_K(\sigma(y))\leq r\qquad\text{for every }y\in M.
\]
Suppose equality holds at $x\in M$.  Let
$p=\operatorname{proj}_{K_x}(\sigma(x))$ and
\[
       \nu=\frac{\sigma(x)-p}{r}.
\]
Then, for every $X\in T_xM$,
\begin{align}
 h\bigl(\nu,\nabla^{E,t}_X\sigma\bigr)&=0,                 \label{eq:wmp-first-spatial}\\
 h\bigl(\nu,((\nabla^{E,t})^2\sigma)(X,X)\bigr)&\leq0.     \label{eq:wmp-second-spatial}
\end{align}
Consequently,
\begin{equation}
       h(\nu,\Delta_t^E\sigma)\leq0.                       \label{eq:wmp-laplacian-support}
\end{equation}
\end{lemma}

\begin{proof}
Let $\gamma:(-\delta,\delta)\to M$ be the $g(t)$-geodesic with
$\gamma(0)=x$ and $\dot\gamma(0)=X$.  Parallel transport $p$ and $\nu$ along
$\gamma$ using the single, frozen connection $\nabla^{E,t}$, and denote the
resulting fields by $p(s)$ and $\nu(s)$.  Then $p(s)\in K_{\gamma(s)}$.
Moreover, $\nu(s)$ supports $K_{\gamma(s)}$ at $p(s)$: transport any point
of $K_{\gamma(s)}$ back to $K_x$ and use that parallel transport preserves
$h$.

Lemma~\ref{lem:wmp-projection-support}, applied in $E_{\gamma(s)}$, gives
\[
 \psi(s):=
 h\bigl(\nu(s),\sigma(\gamma(s))-p(s)\bigr)
 \leq d_K(\sigma(\gamma(s)))\leq r.
\]
At $s=0$ equality holds.  Hence $\psi'(0)=0$ and $\psi''(0)\leq0$.
The fields $p(s)$ and $\nu(s)$ are parallel, so the first equality is
\eqref{eq:wmp-first-spatial}.  Since $\gamma$ is a geodesic, the second
inequality is precisely \eqref{eq:wmp-second-spatial}.  Summing
\eqref{eq:wmp-second-spatial} over a $g(t)$-orthonormal basis of $T_xM$
proves \eqref{eq:wmp-laplacian-support}.
\end{proof}

\begin{remark}[Why no radial second-derivative claim is needed]
\label{rem:wmp-no-radial-hessian}
One sometimes extends $p$ and $\nu$ radially and informally declares all of
their second covariant derivatives to vanish at the center.  For a curved
bundle connection, the antisymmetric part of the second derivative detects
bundle curvature, so such a statement is too strong without qualification.
Lemma~\ref{lem:wmp-geodesic-support} needs only the second derivative in the
same direction twice.  Parallel transport along the corresponding geodesic
gives exactly that fact, with no two-dimensional extension and no curvature
commutator.
\end{remark}

\begin{lemma}[Continuity of a compact spatial maximum]
\label{lem:wmp-compact-maximum}
Let $0<\tau<T$ and let $q:M\times[0,\tau]\to\R$ be continuous.  Then
\[
       Q(t)=\max_{x\in M}q(x,t)
\]
is well-defined and continuous on $[0,\tau]$.
\end{lemma}

\begin{proof}
The maximum is attained because $M$ is compact.  The function $q$ is
uniformly continuous on the compact set $M\times[0,\tau]$, and
\[
 |Q(t)-Q(s)|
    \leq\max_{x\in M}|q(x,t)-q(x,s)|.
\]
The right side tends to zero as $s\to t$.
\end{proof}

\section{The strict expanding-tube proof}
\label{sec:wmp-proof}

\begin{proof}[Proof of Theorem~\ref{thm:hamilton-wmp-systems}]
Fix $0<\tau<T$.  All constants below may depend on $\tau$, but not on the
barrier parameter introduced later.  Write $u_0(x)=u(x,0)$.

We first obtain the only Lipschitz estimate for $F$ that will be used.
Whenever $p$ is the nearest point in $K_x$ to some $u(x,t)$ with
$0\leq t\leq\tau$, the inclusion $u_0(x)\in K_x$ gives
\[
|u(x,t)-p|_h=d_K(u(x,t))
       \leq |u(x,t)-u_0(x)|_h.
\]
More explicitly, set
\[
 A=\max_{M\times[0,\tau]}|u|_h,
 \qquad A_0=\max_M|u_0|_h,
 \qquad R=1+2A+A_0.
\]
Then $|u-p|_h\leq A+A_0$ and
\[
       |p|_h\leq |u|_h+|u-p|_h\leq2A+A_0<R.
\]
Thus all such $p$, all values $u(x,t)$, and all fiberwise line segments
joining them lie in
\[
       D_R(E)=\{v\in E:|v|_h\leq R\}.
\]
The set $D_R(E)\times[0,\tau]$ is compact.  Let
\[
 L=\max_{D_R(E)\times[0,\tau]}
      \bigl\|D^{\mathrm{vert}}F\bigr\|_{\mathrm{op}}.
\]
Applying the fundamental theorem of calculus to the restriction of $F$ to
the fiberwise line segment from $v$ to $w$ gives
\begin{equation}
   |F(v,t)-F(w,t)|_h\leq L|v-w|_h                         \label{eq:wmp-local-lipschitz}
\end{equation}
whenever $v$ and $w$ are one of the pairs just described.

Choose a constant $C>L$.  For $\varepsilon>0$, define the strictly expanding
radius
\[
       r_\varepsilon(t)=\varepsilon e^{Ct}.
\]
We claim that
\begin{equation}
       d_K(u(x,t))<r_\varepsilon(t)
       \qquad (x\in M,\ 0\leq t\leq\tau).                 \label{eq:wmp-tube-claim}
\end{equation}
At $t=0$ this is true because $d_K(u_0(x))=0$.

Suppose the claim fails.  By Lemmas~\ref{lem:wmp-distance-continuous} and
\ref{lem:wmp-compact-maximum}, the function
\[
       H(t)=\max_{x\in M}\bigl(d_K(u(x,t))-r_\varepsilon(t)\bigr)
\]
is continuous.  Choose $t_1\in(0,\tau]$ with $H(t_1)\geq0$.  Since
$H(0)=-\varepsilon$, the intermediate value theorem shows that
\[
 Z=\{t\in[0,t_1]:H(t)=0\}
\]
is a nonempty compact set.  Let $t_0=\min Z$.  Then $t_0>0$ and
$H(t)<0$ for $0\leq t<t_0$.  Choose $x_0\in M$ at which the maximum
$H(t_0)=0$ is attained, and put
\[
  w=u(x_0,t_0),\qquad
  r=r_\varepsilon(t_0),\qquad
  p=\operatorname{proj}_{K_{x_0}}(w),\qquad
  \nu=\frac{w-p}{r}.
\]
Then $r>0$, $|\nu|_h=1$, and $\nu\in N_{K_{x_0}}(p)$.

First consider time while keeping $x_0,p,\nu$ fixed.  For $t<t_0$ close to
$t_0$, Lemma~\ref{lem:wmp-projection-support} and the definition of the
first contact give
\[
 h\bigl(\nu,u(x_0,t)-p\bigr)
     \leq d_K(u(x_0,t))
     <r_\varepsilon(t),
\]
whereas equality holds at $t=t_0$.  The left derivative at $t_0$ is
therefore nonnegative:
\begin{equation}
       0\leq h(\nu,\partial_tu(x_0,t_0))-Cr.               \label{eq:wmp-time-contact}
\end{equation}

Next freeze time at $t_0$.  The inequality
$d_K(u(x,t_0))\leq r$ holds for every $x\in M$, with equality at $x_0$.
Lemma~\ref{lem:wmp-geodesic-support} yields
\begin{equation}
       h(\nu,\Delta_{t_0}^E u(x_0,t_0))\leq0.              \label{eq:wmp-diffusion-contact}
\end{equation}
The first-derivative conclusion of the same lemma, evaluated on
$X(t_0)_{x_0}$, gives
\begin{equation}
 h\bigl(\nu,\nabla^{E,t_0}_{X(t_0)}u(x_0,t_0)\bigr)=0.      \label{eq:wmp-drift-contact}
\end{equation}
The ODE hypothesis and Lemma~\ref{lem:wmp-ode-support} give
\begin{equation}
       h(\nu,F(p,t_0))\leq0.                               \label{eq:wmp-reaction-at-p}
\end{equation}
Finally, \eqref{eq:wmp-local-lipschitz} and $|w-p|_h=r$ give
\begin{equation}
 h\bigl(\nu,F(w,t_0)-F(p,t_0)\bigr)
     \leq |F(w,t_0)-F(p,t_0)|_h
     \leq Lr.                                             \label{eq:wmp-reaction-lipschitz}
\end{equation}

Insert the PDE into \eqref{eq:wmp-time-contact} and use
\eqref{eq:wmp-diffusion-contact}--\eqref{eq:wmp-reaction-lipschitz}.  We
obtain
\[
 \begin{split}
  0
  &\leq h\bigl(\nu,\Delta_{t_0}^E u
             +\nabla^{E,t_0}_{X(t_0)}u+F(w,t_0)\bigr)-Cr\\
  &\leq 0+0+0+Lr-Cr
   =(L-C)r<0,
 \end{split}
\]
a contradiction.  Thus \eqref{eq:wmp-tube-claim} holds.

The constants $L$ and $C$ were chosen independently of $\varepsilon$.
Letting $\varepsilon\downarrow0$ gives
$d_K(u(x,t))=0$ on $M\times[0,\tau]$.  Since each $K_x$ is closed,
$u(x,t)\in K_x$ there.  The choice of $\tau<T$ was arbitrary, completing
the proof.
\end{proof}

\begin{remark}[Where compactness and absence of boundary enter]
\label{rem:wmp-compactness}
Compactness of $M$ is used twice: first, to turn pointwise bounds into a
uniform fiberwise Lipschitz constant for $F$ on a finite time slab, and
second, to ensure that the continuous first-contact function attains its
spatial maximum.  The hypothesis $\partial M=\varnothing$ is used
separately in Lemma~\ref{lem:wmp-geodesic-support}, where one takes a
two-sided geodesic through a possible contact point.  On a manifold with
boundary, an appropriate boundary condition and a boundary maximum
principle would also be required.  A noncompact theorem likewise needs
additional growth hypotheses and a spatial cutoff or exhaustion argument;
neither extension is a formal consequence of
Theorem~\ref{thm:hamilton-wmp-systems}.
\end{remark}

\section{Equivalent engine and useful variants}
\label{sec:wmp-variants}

\begin{corollary}[Support-form maximum principle]
\label{cor:wmp-support-form}
In Theorem~\ref{thm:hamilton-wmp-systems}, replace the ODE-invariance
hypothesis by the pointwise inward condition
\begin{equation}
  h(\nu,F(p,t))\leq0
  \quad\text{for every }x,t,\ p\in K_x,\text{ and }
  \nu\in N_{K_x}(p).                                      \label{eq:wmp-support-hypothesis}
\end{equation}
Then the same conclusion holds.
\end{corollary}

\begin{proof}
Repeat the proof of Theorem~\ref{thm:hamilton-wmp-systems}, using
\eqref{eq:wmp-support-hypothesis} directly in place of
Lemma~\ref{lem:wmp-ode-support}.
\end{proof}

\begin{remark}[Why the associated-bundle interface is recorded]
\label{rem:wmp-associated-purpose}
The next corollary is a construction lemma, not a second maximum principle.
Its immediate use is to turn the single orthogonally invariant model region
$\mathcal K_{\mathrm{HI}}\subset\operatorname{SymEnd}(\mathbb R^3)$ into
the parallel family $K_{\mathrm{HI}}\subset\mathcal E$ in
Corollary~\ref{cor:hi-associated-fixed-family}.  Closedness, convexity,
reaction equivariance, and parallel invariance can then be checked in one
fixed representation space rather than separately in every fiber.  The same
language is useful later in frame-bundle formulations of Hamilton's matrix
Harnack estimate, but that estimate also requires additional spacetime
tensors and evolution equations and is not a consequence of this corollary.
\end{remark}

\begin{corollary}[Associated-bundle form]
\label{cor:wmp-associated-bundle}
Let $G$ be a finite-dimensional Lie group, let $P\to M$ be a smooth
principal $G$-bundle, and let
\[
       \rho:G\longrightarrow O(W)
\]
be a smooth orthogonal representation on a finite-dimensional real
inner-product space $W$.  Let $C\subset W$ be nonempty, closed, convex, and
$G$-invariant.  Give $E=P\times_GW$ a smooth time-dependent family of
connections induced from principal connections on $P$.  Then
\[
       K=P\times_GC\subset E
\]
is a well-defined parallel closed convex family.

More precisely, let $f:W\times[0,T)\to W$ be smooth and equivariant:
\begin{equation}
 f(\rho(g)w,t)=\rho(g)f(w,t)
 \qquad(g\in G,\ w\in W,\ 0\leq t<T).                     \label{eq:wmp-equivariant-reaction}
\end{equation}
Then
\[
       F([p,w],t)=[p,f(w,t)]
\]
is a well-defined smooth vertical map on $E$.  If every solution of
$\dot w=f(w,t)$ which starts in $C$ at an arbitrary starting time remains
in $C$, then the induced fiber ODE preserves every $K_x$, and
Theorem~\ref{thm:hamilton-wmp-systems} applies.
\end{corollary}

\begin{proof}
The $G$-invariance of $C$ makes $P\times_GC$ independent of the choice of
representative in $P\times W$.  A smooth local section of $P$ over an open
set $U\subset M$ identifies $E|_U$ with $U\times W$ and $K|_U$ with
$U\times C$.  Hence $K$ is closed in $E$, and every fiber is nonempty,
closed, and convex.  Orthogonality makes the fiber identifications
isometric.  Parallel transport in an associated bundle is
induced by the principal connection and changes a representative only by the
$G$-action, so it carries the corresponding copy of $C$ exactly onto the
copy in the target fiber.  Equation
\eqref{eq:wmp-equivariant-reaction} makes $F$ independent of the
representative.  In an associated-bundle trivialization, its fiber ODE is
the model ODE conjugated by an orthogonal transformation, so preservation
of $C$ gives preservation of $K_x$.
\end{proof}

\begin{remark}[Fixed sets versus moving sets]
\label{rem:wmp-moving-sets}
The set $K$ in this chapter does not depend on time.  A family $K(t)$ has an
additional temporal velocity, so condition
\eqref{eq:wmp-inward-support} must be replaced by a spacetime tangent or
avoidance condition.  This distinction matters for pinching estimates,
including the Hamilton--Ivey estimate.  The time-dependent-set theorem will
be stated and proved as a separate result; it must not be obtained by simply
writing $K(t)$ in Theorem~\ref{thm:hamilton-wmp-systems}.
\end{remark}

\section{Compatibility with Uhlenbeck's trick}
\label{sec:wmp-uhlenbeck}
\index{Uhlenbeck's trick}

We now explain why the fixed fiber metric in
Theorem~\ref{thm:hamilton-wmp-systems} is exactly the right convention for
Ricci flow.  This subsection records the interface; the curvature evolution
formula itself belongs to the later curvature chapter.

Let $g(t)$ be any smooth family of metrics on $TM$.  Choose a fixed vector
bundle $E$ isomorphic to $TM$, a metric $h$ on $E$, and an initial isometry
\[
       \iota_0:(E,h)\longrightarrow(TM,g(0)).
\]
Define the $g(t)$-self-adjoint endomorphism $B_t:TM\to TM$ by
\begin{equation}
       g(t)(B_tX,Y)=-\frac12\,\partial_tg(t)(X,Y),          \label{eq:wmp-uhlenbeck-B}
\end{equation}
and evolve the bundle isomorphism by
\begin{equation}
       \partial_t\iota_t=B_t\circ\iota_t,\qquad
       \iota_t|_{t=0}=\iota_0.                             \label{eq:wmp-uhlenbeck-iota}
\end{equation}
On every compact time subinterval, this pointwise linear ODE has a unique
solution depending smoothly on $x$ and $t$.  Fix $x\in M$ and
time-independent vectors $\xi,\eta\in E_x$.  Then
\eqref{eq:wmp-uhlenbeck-B} gives
\[
 \frac{d}{dt}g(t)(\iota_t\xi,\iota_t\eta)
   =\partial_tg(t)(\iota_t\xi,\iota_t\eta)
    +g(t)(B_t\iota_t\xi,\iota_t\eta)
    +g(t)(\iota_t\xi,B_t\iota_t\eta)
   =0.
\]
Hence
\[
       \iota_t^*g(t)=h
\]
for every $t$.  In particular, this identity keeps $\iota_t$ injective and
surjective, so the solution of \eqref{eq:wmp-uhlenbeck-iota} remains a
smooth bundle isomorphism.  Pulling back the Levi--Civita connection defines
\begin{equation}
 D^t_X\xi
   =\iota_t^{-1}\!\left(
       \nabla^{g(t)}_X(\iota_t\xi)\right).                 \label{eq:wmp-pulled-connection}
\end{equation}
The connection $D^t$ is $h$-compatible and varies with time.  Its rough
Laplacian is exactly of the form
\eqref{eq:wmp-laplacian}.

For Ricci flow,
\[
       \partial_tg=-2\Ric(g),
\]
equation \eqref{eq:wmp-uhlenbeck-B} says $B_t=\Ric(g(t))^\sharp$.
The curvature operator, transported to the appropriate tensor bundle built
from $(E,h)$, then satisfies a reaction--diffusion equation of the schematic
Hamilton form
\begin{equation}
       \partial_t\mathcal R
          =\Delta_{D^t}\mathcal R
             +\mathcal R^2+\mathcal R^\# ,                \label{eq:wmp-curvature-schematic}
\end{equation}
with the precise coefficients determined by the curvature-operator
normalization.  Equation \eqref{eq:wmp-curvature-schematic} is explanatory,
not yet a formalization statement: no downstream proof may cite it until a
later chapter fixes the exact curvature carrier, pullback identification,
normalization, and coefficients of $\mathcal R^2+\mathcal R^\#$.
Orthogonally invariant closed convex cones in the model fiber produce
associated closed fiberwise-convex subsets $K$, often called cone bundles,
as in
Corollary~\ref{cor:wmp-associated-bundle}.  Thus the evolving base metric
and the evolving connection are already permitted by the theorem, while the
fiber metric and the algebraic cone stay fixed.  This is the precise sense
in which the result is compatible with Uhlenbeck's trick.

\begin{remark}[What the theorem does not require]
\label{rem:wmp-no-ricci-hypothesis}
Neither the proof nor the statement uses the Ricci-flow equation, a formula
for $\partial_t\nabla^{E,t}$, or a sign condition on
$\partial_tg$.  At a first contact, $g(t_0)$ and $\nabla^{E,t_0}$ are simply
frozen long enough to take geodesics and trace the spatial Hessian.
\end{remark}

\section{Formalization interfaces and dependency order}
\label{sec:wmp-formalization}

The human proof above deliberately separates facts which belong to different
parts of a Lean development.  A faithful blueprint should refine it in the
following dependency order.

\begin{enumerate}
\item \emph{Fiber metric and transport.}
      Define the time-indexed metric connections, prove that their parallel
      transports are linear isometries, prove smooth parameter dependence
      for transport along the contraction $\Gamma$, and formulate invariance
      of a fiberwise set under transport and reversed transport.
\item \emph{Local product and distance.}
      Trivialize by transport along a smoothly varying family of paths and
      prove continuity of $v\mapsto\dist(v,K_{\pi(v)})$.
\item \emph{Finite-dimensional projection.}
      Prove existence and uniqueness of the nearest point to a nonempty
      closed convex set, the projection inequality
      \eqref{eq:wmp-projection-normal}, and the support bound
      \eqref{eq:wmp-support-distance}.
\item \emph{ODE-to-normal reduction.}
      Import the finite-dimensional nonautonomous Picard--Lindel\"of theorem,
      define the maximal interval $I(x,t_0,a_0)$, and, starting from an
      arbitrary time $t_0$, derive \eqref{eq:wmp-inward-support} by a right
      derivative.
\item \emph{Geodesic scalarization.}
      Parallel transport one supporting pair $(p,\nu)$ along one frozen
      $g(t_0)$-geodesic and derive the first derivative, diagonal second
      derivative, and traced inequalities in
      Lemma~\ref{lem:wmp-geodesic-support}.
\item \emph{Uniform local Lipschitz bound.}
      Use the known initial section in $K$, the nearest-point estimate, and
      compactness of $D_R(E)\times[0,\tau]$; then apply the fiberwise
      fundamental theorem of calculus to $D^{\mathrm{vert}}F$.
\item \emph{First contact with a strict tube.}
      Prove Lemma~\ref{lem:wmp-compact-maximum}, construct the least zero of
      $H$, obtain the left time derivative inequality, and combine the three
      scalar inequalities with $C>L$.
\item \emph{Vanishing tube radius.}
      Make explicit that $C$ is independent of $\varepsilon$, let
      $\varepsilon\downarrow0$, and use closedness of each $K_x$.
\item \emph{Headline and derived interfaces.}
      Package the support-form theorem first if convenient, derive Hamilton's
      ODE-to-PDE statement with drift, and then derive the associated-bundle
      corollary.
\item \emph{Uhlenbeck bridge.}
      Separately formalize smooth parameter dependence for the pointwise
      linear ODE \eqref{eq:wmp-uhlenbeck-iota}, the fixed-metric identity,
      and the pulled-back metric connection.  The exact curvature carrier
      and equation \eqref{eq:wmp-curvature-schematic} are later obligations.
\end{enumerate}

The convex-analysis portion is finite dimensional and should require no
choice of measurable or smooth projections.  The nearest point is selected
only at the single hypothetical contact point.  The analytically deepest
formalization work is therefore expected to lie in the bundle-valued
geodesic differentiation and the time-dependent rough Laplacian, rather
than in nonsmooth convex analysis.

\begin{remark}[Suggested theorem boundary]
\label{rem:wmp-formal-boundary}
For the first compiled implementation, it is reasonable to take
$F$, $g$, $\nabla^{E,t}$, $X$, and $u$ smooth and $M$ compact without boundary,
exactly as above.  Weakening regularity, adding a boundary, treating
noncompact bases, and allowing moving convex families are separate
generalizations.  They should not be folded into the first theorem statement
unless a downstream use actually requires them.
\end{remark}

\section{Source record and reconstruction notes}
\label{sec:wmp-sources}

The primary source is R.~S.~Hamilton,
\emph{Four-manifolds with positive curvature operator},
Journal of Differential Geometry \textbf{24} (1986), 153--179,
\href{https://doi.org/10.4310/jdg/1214440433}
     {online DOI record}.
Uhlenbeck's trick is developed on pp.~155--157, the upper-derivative lemmas
on pp.~158--159, the ODE tangent-cone lemma on pp.~160--161, the
trivial-bundle maximum principle in Theorem~4.2 on p.~161, and the
vector-bundle theorem and proof in Theorem~4.3 on pp.~162--163.

The statement and its role are also discussed in \emph{The Ricci Flow: An
Introduction}, Mathematical Surveys and Monographs~110, pp.~100--101, with
the Uhlenbeck construction on pp.~180--183 and the curvature evolution and
algebraic reduction continuing through p.~186; in \emph{Hamilton's Ricci
Flow}, Graduate Studies in Mathematics~77, pp.~121--122 and 128--132; and
in \emph{The Ricci Flow: Techniques and Applications, Part II: Analytic
Aspects}, Mathematical Surveys and Monographs~144, Chapter~10, especially
pp.~9--17 and 43--46.
The fixed-bundle Uhlenbeck formulation in Mathematical Surveys and
Monographs~135, Appendix~A, pp.~457--459, and the nearest-point support
lemma in Mathematical Surveys and Monographs~163, Appendix~H,
Lemma~H.8, p.~416, are useful supplementary references.  A further expanded
treatment is B.~Chow and P.~Lu, \emph{The time-dependent maximum principle
for systems of parabolic equations subject to an avoidance set}, Pacific Journal of
Mathematics \textbf{214} (2004), 201--222,
\href{https://doi.org/10.2140/pjm.2004.214.201}{online DOI record}.

The proof in this chapter is a reconstruction, not a transcription.  It
retains Hamilton's normal-support mechanism but replaces the upper-Dini
derivative of a global distance supremum by a strict exponentially expanding
tube.  It also makes explicit six hypotheses or distinctions that are easy
to leave implicit in prose:
\begin{enumerate}
\item the nonautonomous fiber ODE starts at every time $t_0$;
\item the base metric may be any smooth family and no orientation is used;
\item the fiber metric is fixed while the compatible connection may vary;
\item an associated convex set must be invariant under the structure group;
\item a time-dependent convex family is a different theorem.
\item Hamilton allows the reaction on an open subset $U\subset E$ containing
      the invariant set, whereas this first project interface uses a global
      smooth map $F:E\times[0,T)\to E$.
\end{enumerate}

There is one deliberate regularity difference from Mathematical Surveys and
Monographs~144, Theorem~10.13.  That theorem permits a reaction continuous in
all variables and locally Lipschitz in the fiber variable, whereas smoothness
is imposed here for the first implementation.  Both statements include a
drift term in the headline.  That volume proves its fixed-set theorem by
specializing the moving-set theorem, Theorem~10.16; this chapter instead
gives an independent fixed-set proof.  These are intentional interface
choices, not claims of verbatim reproduction.

This is a complete human proof of the fixed-set weak maximum principle.
It supplies the mathematical source for a later blueprint and Lean
implementation, but it makes no claim that those formal artifacts have
already been compiled or audited.

\chapter{The Hamilton--Ivey Curvature Pinching Estimate}
\label{ch:hamilton-ivey}
\index{Hamilton--Ivey estimate}
\index{curvature pinching}
\index{Uhlenbeck's trick}

\section{Purpose, convention, and main theorem}
\label{sec:hi-purpose}

The Hamilton--Ivey estimate is the dimension-three mechanism which forces
large negative sectional curvature to be accompanied by much larger
positive curvature.  Its strongest elementary form contains an improving
time term.  In the normalization used below, it says that if the least
curvature-operator eigenvalue is initially at least $-1$, then, wherever
that eigenvalue is negative,
\[
 R\geq(-\nu)\bigl(\log(-\nu)+\log(1+t)-3\bigr).
\]
The logarithm is not an auxiliary error: it is the term which makes the
negative part negligible under high-curvature blow-up.

The classical proofs apply a maximum principle to a time-dependent convex
set.  That route is mathematically sound, but it would force the project to
formalize a second maximum principle before using
Theorem~\ref{thm:hamilton-wmp-systems}.  We instead observe that Hamilton's
moving sets are homothetic.  Multiplication of the curvature operator by
$C^{-1}+t$ turns them into one fixed closed convex set.  Chapter~2 then
applies without alteration.  This rescaling is the organizing idea of the
chapter and is the source of the time-improvement term.

We first freeze the convention which controls every factor of two.

\begin{definition}[Curvature endomorphism and covariant curvature tensor]
\label{def:hi-curvature-tensor}
Let $g$ be a Riemannian metric on a smooth manifold $M$, and let $\nabla^g$
be its Levi--Civita connection.  For smooth vector fields $X,Y,Z,W$ define
\begin{align*}
 R^g(X,Y)Z
   &=\nabla_X^g\nabla_Y^gZ-\nabla_Y^g\nabla_X^gZ
        -\nabla_{[X,Y]}^gZ,\\
 \mathbf R_g(X,Y,Z,W)
   &=g\bigl(R^g(X,Y)Z,W\bigr).
\end{align*}
This is the sign for which the round unit sphere has positive sectional
curvature.
\end{definition}

\begin{definition}[Trace-normalized curvature operator]
\label{def:hi-curvature-operator}
For each $x\in M$, equip $\Lambda^2T_xM$ with the inner product induced by
$g_x$.  The \emph{trace-normalized curvature operator} is the unique
self-adjoint endomorphism
\[
 \mathcal R_g(x):\Lambda^2T_xM\longrightarrow\Lambda^2T_xM
\]
such that, for every $X,Y,Z,W\in T_xM$,
\begin{equation}
 \left\langle\mathcal R_g(x)(X\wedge Y),Z\wedge W\right\rangle_{\Lambda^2g_x}
    =2\,\mathbf R_g(X,Y,W,Z).                              \label{eq:hi-curvature-normalization}
\end{equation}
Existence and uniqueness follow because decomposable two-vectors span
$\Lambda^2T_xM$ and the curvature symmetries make the induced bilinear
form symmetric.
\end{definition}

Thus, on a manifold of constant sectional curvature $k$,
\[
             \mathcal R_g=2kI.
\]
In dimension three let
\begin{equation}
        \lambda\geq\mu\geq\nu                             \label{eq:hi-eigenvalue-order}
\end{equation}
be the eigenvalues of $\mathcal R_g$.  They are twice the three principal
sectional curvatures, and
\begin{equation}
        R=\trace\mathcal R_g=\lambda+\mu+\nu.              \label{eq:hi-trace-is-scalar}
\end{equation}
Indeed, in an orthonormal basis the trace in
\eqref{eq:hi-curvature-normalization} is twice the sum of the three
coordinate sectional curvatures, which is the scalar curvature.  All
occurrences of $\nu$ in this chapter use this convention.

\begin{theorem}[Hamilton--Ivey estimate with time improvement]
\label{thm:hamilton-ivey-scaled}
Let $M^3$ be a closed smooth three-manifold, let $0<T\leq\infty$, and let
$g(t)$, $t\in[0,T)$, be a smooth solution of
\[
                 \partial_tg=-2\Ric(g).
\]
No connectedness or orientability assumption is required.  Fix $C>0$ and
suppose that the least eigenvalue in
\eqref{eq:hi-eigenvalue-order} satisfies
\begin{equation}
                 \nu(x,0)\geq-C
                 \qquad\text{for every }x\in M.            \label{eq:hi-initial-lower}
\end{equation}
Then, at every $(x,t)\in M\times[0,T)$ for which $\nu(x,t)<0$,
\begin{equation}
 R(x,t)\geq
 (-\nu(x,t))
 \left[
   \log\!\left((C^{-1}+t)(-\nu(x,t))\right)-3
 \right].                                                 \label{eq:hi-main-scaled}
\end{equation}
The argument of the logarithm is positive and invariant under parabolic
rescaling.
\end{theorem}

\textcolor{red}{Since the above is a major theorem, put the Lean code for the statement here.}

\begin{corollary}[Normalized Hamilton--Ivey estimate]
\label{cor:hamilton-ivey-normalized}
If $\nu(x,0)\geq-1$ for every $x$, then wherever $\nu(x,t)<0$,
\begin{equation}
 R(x,t)\geq
 (-\nu(x,t))
 \bigl(\log(-\nu(x,t))+\log(1+t)-3\bigr).                 \label{eq:hi-normalized}
\end{equation}
\end{corollary}

\begin{remark}[Translation to the sectional-curvature normalization]
\label{rem:hi-sectional-normalization}
If $\kappa_{\min}=\nu/2$ denotes the least eigenvalue of the curvature
operator whose eigenvalues are the sectional curvatures themselves, then
the hypothesis $\kappa_{\min}(\cdot,0)\geq-K$, $K>0$, gives
\begin{equation}
 R\geq2(-\kappa_{\min})
 \left[
   \log\!\left((-\kappa_{\min})(K^{-1}+2t)\right)-3
 \right]                                                 \label{eq:hi-sectional-version}
\end{equation}
wherever $\kappa_{\min}<0$.  Formula
\eqref{eq:hi-sectional-version} is the same theorem, not a competing
estimate.  Recording this conversion prevents a factor-of-two change in
the eigenvalue ODE from silently changing the time term.
\end{remark}

\begin{roadmap}{Section~\ref{sec:hi-uhlenbeck} puts the curvature operator
on a bundle with a fixed fiber metric and records the exact
dimension-three reaction.}
Section~\ref{sec:hi-fixed-region} proves that the one fixed
Hamilton--Ivey region is closed, convex, and parallel invariant.
Section~\ref{sec:hi-ode} proves its ODE invariance.  Section~\ref{sec:hi-pde}
rescales the curvature section, applies Chapter~2, and proves
Theorem~\ref{thm:hamilton-ivey-scaled}.  The final sections record
project-facing consequences, formalization interfaces, and source
reconstruction notes.
\end{roadmap}

\section{The fixed-bundle curvature equation}
\label{sec:hi-uhlenbeck}

The form of Uhlenbeck's trick needed here has three separate layers:
a fixed Euclidean replacement for $TM$, its induced curvature-operator
bundle, and one exact evolution equation.  Keeping these layers separate
avoids a global frame, a choice of orientation, and a time-dependent fiber
inner product.

Choose a rank-three real vector bundle $V\to M$, a positive-definite metric
$h$ on $V$, and a bundle isometry
\[
       \iota_0:(V,h)\longrightarrow(TM,g(0)).
\]
One may take $V=TM$, $h=g(0)$, and $\iota_0$ equal to the identity.  Evolve
$\iota_t:V\to TM$ by the pointwise linear ODE
\begin{equation}
 \partial_t\iota_t=\Ric_{g(t)}^\sharp\circ\iota_t,
 \qquad \iota_t|_{t=0}=\iota_0.                           \label{eq:hi-uhlenbeck-ode}
\end{equation}
For time-independent $v,w\in V_x$,
\begin{align*}
 \frac{d}{dt}g(t)(\iota_tv,\iota_tw)
 &= -2\Ric(\iota_tv,\iota_tw)\\
 &\quad
   +g(\Ric^\sharp\iota_tv,\iota_tw)
   +g(\iota_tv,\Ric^\sharp\iota_tw)=0.
\end{align*}
Consequently
\begin{equation}
                  \iota_t^*g(t)=h                         \label{eq:hi-fixed-metric}
\end{equation}
for every $t$.  In particular, $\iota_t$ remains a bundle isomorphism.

Define a connection on $(V,h)$ by
\begin{equation}
 D_X^t\xi
   =\iota_t^{-1}\!\left(
       \nabla_X^{g(t)}(\iota_t\xi)\right).                \label{eq:hi-pulled-connection}
\end{equation}
The inverse $\iota_t^{-1}$ in \eqref{eq:hi-pulled-connection} is essential:
without it the right side is a tangent vector rather than an element of
$V$.  Equations \eqref{eq:hi-fixed-metric} and metric compatibility of the
Levi--Civita connection imply
\[
 Xh(\xi,\eta)=h(D_X^t\xi,\eta)+h(\xi,D_X^t\eta).
\]
Thus $D^t$ is compatible with the fixed fiber metric $h$, although it
depends on time.

Let
\[
 \mathcal E=\operatorname{SymEnd}_h(\Lambda^2V)
\]
be the bundle of self-adjoint endomorphisms of $\Lambda^2V$, equipped with
its fixed Hilbert--Schmidt fiber metric.  The connection $D^t$ induces a
metric connection, again denoted $D^t$, on $\mathcal E$.  For each $x\in M$
and $t\in[0,T)$, define
\[
 (\Lambda^2\iota_t)_x:\Lambda^2V_x\longrightarrow\Lambda^2T_xM,
 \qquad
 (\Lambda^2\iota_t)_x(v\wedge w)=\iota_t(v)\wedge\iota_t(w),
\]
extended linearly.  This is the exterior square of $\iota_t$; because
$\iota_t$ is a fiberwise isometry, $\Lambda^2\iota_t$ is an isometric
bundle isomorphism for the induced exterior-power metrics.  Pull
$\mathcal R_{g(t)}$ back by $\Lambda^2\iota_t$ and denote the resulting
section of $\mathcal E$ by
\begin{equation}
 \mathcal M(t)
   =(\Lambda^2\iota_t)^{-1}\circ
      \mathcal R_{g(t)}\circ(\Lambda^2\iota_t).            \label{eq:hi-pulled-curvature}
\end{equation}
Because $\Lambda^2\iota_t$ is an isometry, $\mathcal M$ has the same ordered
eigenvalues $\lambda\geq\mu\geq\nu$ and the same trace $R$ as the geometric
operator.

We now define the algebraic reaction without using a Hodge star or an
orientation.  If $A$ is a self-adjoint endomorphism of a three-dimensional
Euclidean space, set
\begin{align}
 \sigma_2(A)
   &=\frac12\left((\trace A)^2-\trace(A^2)\right),          \label{eq:hi-sigma-two}\\
 A^\#
   &=A^2-(\trace A)A+\sigma_2(A)I,                         \label{eq:hi-adjugate}\\
 Q(A)&=A^2+A^\#.                                          \label{eq:hi-reaction}
\end{align}
Thus $A^\#$ is the degree-two adjugate polynomial, not the
Hilbert-space adjoint.  The formulas are polynomial and commute with
orthogonal conjugation.  If an orthonormal basis diagonalizes $A$ with
eigenvalues $\lambda,\mu,\nu$, then
\begin{equation}
 Q(A)=\operatorname{diag}
 \bigl(\lambda^2+\mu\nu,\,
       \mu^2+\lambda\nu,\,
       \nu^2+\lambda\mu\bigr).                            \label{eq:hi-Q-diagonal}
\end{equation}

\begin{proposition}[Uhlenbeck-gauge curvature evolution]
\label{prop:hi-uhlenbeck-evolution}
With the normalization \eqref{eq:hi-curvature-normalization}, the
pulled-back curvature operator satisfies
\begin{equation}
 \partial_t\mathcal M
   =\Delta_{D^t,g(t)}\mathcal M+Q(\mathcal M).             \label{eq:hi-fixed-bundle-pde}
\end{equation}
Here the rough Laplacian is formed from the evolving base metric $g(t)$ and
the evolving metric connection $D^t$, while the vector bundle
$\mathcal E$ and its fiber metric are fixed.
\end{proposition}

\textcolor{red}{Since the above is a major proposition, put the Lean code for the statement here.}

\begin{proof}
This is Theorem~\ref{thm:ue-operator-evolution} below.  The intervening
chapter is allowed to use the displayed equation as a named input;
Chapter~\ref{ch:curvature-evolution} proves it from the covariant
curvature-tensor evolution, including the four-slot Ricci cancellation,
the rough-Laplacian naturality statement, and the dimension-three
tensor-to-operator calculation.
\end{proof}

\begin{remark}[The minimal black-box interface]
\label{rem:hi-evolution-blackbox}
A later Lean development need not begin with the full coordinate derivation
of Proposition~\ref{prop:hi-uhlenbeck-evolution}.  It may first import that
proposition with the exact normalization above.  The pinching proof uses
only: the fixed fiber metric, metric compatibility of $D^t$,
orthogonal equivariance and quadratic homogeneity of $Q$, and the diagonal
formula \eqref{eq:hi-Q-diagonal}.  Those facts form a smaller and more stable
interface than an index-expanded formula for $\partial_t\Rm$.
\end{remark}

\section{The universal fixed Hamilton--Ivey region}
\label{sec:hi-fixed-region}

All geometry now disappears temporarily.  Let $W$ be a
three-dimensional real inner-product space, and write
\[
                   \operatorname{SymEnd}(W)
\]
for its self-adjoint endomorphisms.  For
$A\in\operatorname{SymEnd}(W)$ let $\nu(A)$ be its least eigenvalue and let
$s(A)=\trace A$.  Define
$\mathcal K_{\mathrm{HI}}\subset\operatorname{SymEnd}(W)$ by
\begin{equation}
 \mathcal K_{\mathrm{HI}}
   =
   \left\{
   A\in\operatorname{SymEnd}(W):
   \begin{array}{l}
      s(A)\geq-3,\quad\text{and}\\[1mm]
      \nu(A)\leq-1\ \Longrightarrow\\
      s(A)\geq
      (-\nu(A))\bigl(\log(-\nu(A))-3\bigr)
   \end{array}
   \right\}.                                              \label{eq:hi-fixed-set}
\end{equation}
The logarithm in the definition is evaluated only when
$-\nu(A)\geq1$.  Notice that $\mathcal K_{\mathrm{HI}}$ is not a cone.

We will need more than an assertion that
\eqref{eq:hi-fixed-set} is convex.  The following proof rewrites it as the
sublevel set of one everywhere-defined convex function, a form suited to
reuse.

\begin{lemma}[Least-eigenvalue interface]
\label{lem:hi-least-eigenvalue}
For $A,B\in\operatorname{SymEnd}(W)$ and $0\leq\theta\leq1$:
\begin{align}
 \nu(A)
   &=\min_{\lvert w\rvert=1}\langle Aw,w\rangle,           \label{eq:hi-rayleigh}\\
 -\nu(\theta A+(1-\theta)B)
   &\leq
     \theta(-\nu(A))+(1-\theta)(-\nu(B)),                  \label{eq:hi-minus-min-convex}\\
 |\nu(A)-\nu(B)|
   &\leq\|A-B\|_{\mathrm{op}},                             \label{eq:hi-min-lipschitz}\\
 \trace A&\geq3\nu(A).                                    \label{eq:hi-trace-min}
\end{align}
In particular, $A\mapsto-\nu(A)$ is convex and continuous.
\end{lemma}

\begin{proof}
The unit sphere of $W$ is compact, so the Rayleigh quotient attains its
minimum; the spectral theorem identifies that minimum with the least
eigenvalue.  For any unit $w$,
\[
 -\langle(\theta A+(1-\theta)B)w,w\rangle
 \leq\theta(-\nu(A))+(1-\theta)(-\nu(B)).
\]
Taking the maximum over $w$ proves
\eqref{eq:hi-minus-min-convex}.  Applying
\[
 |\langle(A-B)w,w\rangle|\leq\|A-B\|_{\mathrm{op}}
\]
in both directions proves \eqref{eq:hi-min-lipschitz}.  Finally, all three
eigenvalues of $A$ are at least $\nu(A)$, and their sum is $\trace A$,
which proves \eqref{eq:hi-trace-min}.
\end{proof}

\begin{lemma}[Closed convexity of the fixed region]
\label{lem:hi-fixed-region-convex}
The set $\mathcal K_{\mathrm{HI}}$ is nonempty, closed, convex, and invariant
under conjugation by every orthogonal transformation of $W$.
\end{lemma}

\textcolor{red}{Since the above is a major lemma, put the Lean code for the statement here.}

\begin{proof}
For $r\geq\mathrm e^2$ put
\[
                 f(r)=r(\log r-3).
\]
On this interval,
\[
        f'(r)=\log r-2\geq0,
        \qquad f''(r)=\frac1r>0.
\]
Define a total function $G:\mathbb R\to\mathbb R$ by
\begin{equation}
 G(u)=
 \max\left\{
     -3,\,
     m(u)\bigl(\log m(u)-3\bigr)
 \right\},
 \qquad
 m(u)=\max\{u,\mathrm e^2\}.                              \label{eq:hi-clipped-log}
\end{equation}
The map $m$ is convex.  Since $f$ is convex and nondecreasing on its
domain, $f\circ m$ is convex.  The maximum of the two convex functions
$-3$ and $f\circ m$ is convex.  It is also continuous and nondecreasing.

We claim that
\begin{equation}
 \mathcal K_{\mathrm{HI}}
   =
 \left\{
    A:G(-\nu(A))-\trace A\leq0
 \right\}.                                                \label{eq:hi-fixed-sublevel}
\end{equation}
If $u\leq1$, then $m(u)=\mathrm e^2$ and
$m(u)(\log m(u)-3)=-\mathrm e^2<-3$, so $G(u)=-3$.
If $1\leq u\leq\mathrm e^2$, then
$u(\log u-3)\leq-3$ because its derivative is nonpositive on that
interval.  If $u\geq\mathrm e^2$, the right side of
\eqref{eq:hi-clipped-log} is exactly
\[
                 \max\{-3,u(\log u-3)\}.
\]
These three cases prove \eqref{eq:hi-fixed-sublevel}.

By Lemma~\ref{lem:hi-least-eigenvalue}, the map
$A\mapsto-\nu(A)$ is convex.  Because $G$ is convex and nondecreasing,
$A\mapsto G(-\nu(A))$ is convex.  Subtracting the linear function
$\trace A$ preserves convexity.  Hence the sublevel set in
\eqref{eq:hi-fixed-sublevel} is convex.  Continuity gives closedness, and
the zero operator proves nonemptiness.  Trace and least eigenvalue are
unchanged by orthogonal conjugation, proving the final assertion.
\end{proof}

\begin{corollary}[The associated fixed family]
\label{cor:hi-associated-fixed-family}
In the bundle
\[
          \mathcal E=\operatorname{SymEnd}_h(\Lambda^2V)
\]
of Section~\ref{sec:hi-uhlenbeck}, impose
\eqref{eq:hi-fixed-set} in every fiber.  The resulting subset
$K_{\mathrm{HI}}\subset\mathcal E$ is a parallel closed convex family for
every induced connection $D^t$.
\end{corollary}

\textcolor{red}{Since the above is a major corollary, put the Lean code for the statement here.}

\begin{proof}
An $h$-orthonormal frame identifies a fiber with
$\operatorname{SymEnd}(\mathbb R^3)$, and a different frame changes that
identification by orthogonal conjugation.  Lemma
\ref{lem:hi-fixed-region-convex} and
Corollary~\ref{cor:wmp-associated-bundle} therefore give a well-defined
closed total subset with nonempty closed convex fibers.  Since $D^t$ is
metric, its parallel transports act by orthogonal conjugation and preserve
the family.
\end{proof}

\begin{remark}[Why the signed variable is indispensable]
\label{rem:hi-signed-variable}
In the global convexity argument the correct spectral variable is
$-\nu(A)$, which may be negative when $A$ is positive definite.  It must
not be replaced by $|\nu(A)|$.  Absolute value agrees with $-\nu$ only
after the hypothesis $\nu<0$ has been imposed.  The clipped function
\eqref{eq:hi-clipped-log} also keeps every logarithm away from zero and
negative inputs.
\end{remark}

\section{The invariant reaction ODE}
\label{sec:hi-ode}

Define the autonomous polynomial vector field
\begin{equation}
                 H(A)=A+Q(A).                             \label{eq:hi-normalized-reaction}
\end{equation}
We now prove that its ODE preserves
$\mathcal K_{\mathrm{HI}}$.  The proof differentiates scalar eigenvalue
coordinates only along a fiber ODE.  It never differentiates the least
eigenvalue of the curvature section in space or through a parabolic
eigenvalue multiplicity.

\begin{lemma}[Diagonal and ordering reduction]
\label{lem:hi-ordering}
Let $A(\tau)$ solve
\[
                       A'=H(A).
\]
Choose an orthonormal eigenbasis for $A(\tau_0)$ and label its eigenvalues
$\lambda\geq\mu\geq\nu$.  In that fixed basis the solution remains
diagonal, and its diagonal entries solve
\begin{align}
 \lambda'&=\lambda+\lambda^2+\mu\nu,                       \label{eq:hi-lambda-ode}\\
 \mu'&=\mu+\mu^2+\lambda\nu,                               \label{eq:hi-mu-ode}\\
 \nu'&=\nu+\nu^2+\lambda\mu.                               \label{eq:hi-nu-ode}
\end{align}
The ordering $\lambda\geq\mu\geq\nu$ is preserved for as long as the
solution exists.
\end{lemma}

\textcolor{red}{Since the above is a major lemma, put the Lean code for the statement here.}

\begin{proof}
The polynomial formulas \eqref{eq:hi-adjugate}--\eqref{eq:hi-reaction}
show that $H(A)$ is diagonal in every basis which diagonalizes $A$.
The diagonal subspace is therefore invariant; uniqueness for a smooth
finite-dimensional ODE keeps $A(\tau)$ in that subspace.  Formula
\eqref{eq:hi-Q-diagonal} gives
\eqref{eq:hi-lambda-ode}--\eqref{eq:hi-nu-ode}.  Subtraction gives
\begin{align}
 (\lambda-\mu)'
   &=(\lambda-\mu)(1+\lambda+\mu-\nu),                    \label{eq:hi-order-one}\\
 (\mu-\nu)'
   &=(\mu-\nu)(1+\mu+\nu-\lambda).                        \label{eq:hi-order-two}
\end{align}
Each difference solves a scalar linear ODE.  Its initial nonnegativity is
therefore preserved, including when two eigenvalues coincide.
\end{proof}

\begin{lemma}[Trace barrier]
\label{lem:hi-trace-barrier}
Along the ODE $A'=H(A)$, if $s=\trace A$, then
\begin{equation}
 s'\geq s+\frac23s^2.                                     \label{eq:hi-trace-derivative}
\end{equation}
Consequently the half-space $s\geq-3$ is invariant.
\end{lemma}

\begin{proof}
In the ordered diagonal coordinates of
Lemma~\ref{lem:hi-ordering},
\begin{align}
 \trace Q(A)
 &=\lambda^2+\mu^2+\nu^2+\lambda\mu+\lambda\nu+\mu\nu\\
 &=\frac12\left[
     (\lambda+\mu)^2+(\lambda+\nu)^2+(\mu+\nu)^2
   \right].                                               \label{eq:hi-trace-Q}
\end{align}
The three quantities inside the squares have sum $2s$.  Cauchy--Schwarz
therefore gives $\trace Q(A)\geq\frac23s^2$, proving
\eqref{eq:hi-trace-derivative}.  At $s=-3$ the right side is $3>0$.
Thus a first crossing from $s\geq-3$ to $s<-3$ is impossible.
\end{proof}

\begin{lemma}[Hamilton--Ivey algebraic inequality]
\label{lem:hi-algebra}
Suppose $\lambda\geq\mu\geq\nu$ and $\nu<0$.  Then
\begin{equation}
 \Pi(\lambda,\mu,\nu)
   =
   -\nu(\lambda^2+\mu^2+\lambda\mu)
     +(\lambda+\mu)\lambda\mu
   \geq0.                                                  \label{eq:hi-Pi}
\end{equation}
\end{lemma}

\begin{proof}
If $\mu<0$, then
\begin{equation}
 \Pi
  =(\mu-\nu)(\lambda^2+\mu^2+\lambda\mu)-\mu^3.            \label{eq:hi-Pi-negative-mu}
\end{equation}
Here $\mu-\nu\geq0$,
\[
 \lambda^2+\mu^2+\lambda\mu
   =\left(\lambda+\frac{\mu}{2}\right)^2+\frac34\mu^2
   \geq0,
\]
and $-\mu^3>0$.  If $\mu\geq0$, then $\lambda\geq\mu\geq0$
and
\begin{equation}
 \Pi
   =\lambda^2(\mu-\nu)-\nu\mu^2
      +\lambda\mu(\mu-\nu)\geq0.                          \label{eq:hi-Pi-nonnegative-mu}
\end{equation}
The identities \eqref{eq:hi-Pi-negative-mu} and
\eqref{eq:hi-Pi-nonnegative-mu} follow by expansion, and the two cases
exhaust all possibilities.
\end{proof}

\begin{lemma}[Pinching-functional derivative]
\label{lem:hi-pinching-derivative}
For $\nu<0$, put
\begin{equation}
 q=-\nu>0,\qquad
 \Omega=\frac{\lambda+\mu+\nu}{q}-\log q.                 \label{eq:hi-Omega}
\end{equation}
Along the quadratic ODE $A'=Q(A)$,
\begin{equation}
                  \Omega'\geq q.                          \label{eq:hi-Omega-Q}
\end{equation}
Along the normalized ODE $A'=H(A)=A+Q(A)$,
\begin{equation}
                  \Omega'\geq q-1.                        \label{eq:hi-Omega-H}
\end{equation}
\end{lemma}

\begin{proof}
Write $s=\lambda+\mu+\nu$.  From the quadratic parts of
\eqref{eq:hi-lambda-ode}--\eqref{eq:hi-nu-ode}, direct differentiation of
\eqref{eq:hi-Omega} gives the exact identity
\begin{equation}
 \nu^2\Omega'_Q
    =-\nu^3+\Pi(\lambda,\mu,\nu).                          \label{eq:hi-Omega-exact}
\end{equation}
Since $\nu^2=q^2$, $-\nu^3=q^3$, and
$\Pi\geq0$ by Lemma~\ref{lem:hi-algebra},
\eqref{eq:hi-Omega-Q} follows.  Under the radial ODE $A'=A$, both $s$ and
$q$ have derivative equal to themselves.  Thus $s/q$ has derivative zero
and $-\log q$ has derivative $-1$.  Adding this radial contribution proves
\eqref{eq:hi-Omega-H}.
\end{proof}

\begin{proposition}[ODE invariance of the fixed region]
\label{prop:hi-fixed-region-ode}
Every maximal solution of
\[
                         A'=H(A)
\]
which starts in $\mathcal K_{\mathrm{HI}}$ at an arbitrary starting time
remains in $\mathcal K_{\mathrm{HI}}$ for all later times on which it
exists.
\end{proposition}

\textcolor{red}{Since the above is a major proposition, put the Lean code for the statement here.}

\begin{proof}
By Lemma~\ref{lem:hi-trace-barrier}, $s=\trace A$ remains at least $-3$.
On an interval where $q=-\nu>1$, the second condition in
\eqref{eq:hi-fixed-set} is precisely
\begin{equation}
                         \Omega\geq-3.                     \label{eq:hi-Omega-boundary}
\end{equation}
Lemma~\ref{lem:hi-pinching-derivative} gives
$\Omega'\geq q-1>0$ there.  Hence \eqref{eq:hi-Omega-boundary} cannot be
lost on such an interval.  If the solution enters this interval through
$q=1$, then
\[
                  \Omega=s\geq-3
\]
at the entry time.  If the ODE starts with $q>1$, the same inequality is
part of the hypothesis $A\in\mathcal K_{\mathrm{HI}}$ at that starting
time.

It remains only to check the seam $q=1$.  There the logarithmic constraint
is exactly $s\geq-3$, because
$q(\log q-3)=-3$.  If equality holds, the trace derivative is strictly
positive by Lemma~\ref{lem:hi-trace-barrier}.  Thus neither the seam nor a
simultaneous contact of the two boundary pieces creates an exit.  This
covers every possible first exit from the closed set
$\mathcal K_{\mathrm{HI}}$.
\end{proof}

\begin{corollary}[Positive time-dependent multiples]
\label{cor:hi-time-multiple-ode}
Let $a(t)>0$ be smooth.  The nonautonomous ODE
\[
                         A'=a(t)H(A)
\]
preserves $\mathcal K_{\mathrm{HI}}$ from every starting time.
\end{corollary}

\begin{proof}
Starting at $t_0$, introduce
\[
                    \tau(t)=\int_{t_0}^t a(r)\,dr.
\]
This is strictly increasing.  The image trajectory, parametrized by
$\tau$, solves $dA/d\tau=H(A)$.  Apply
Proposition~\ref{prop:hi-fixed-region-ode}.
\end{proof}

\section{From the fixed ODE region to the curvature PDE}
\label{sec:hi-pde}

We now prove the main theorem.  Fix $C>0$ as in
\eqref{eq:hi-initial-lower} and define the positive affine scale
\begin{equation}
                  b(t)=C^{-1}+t.                          \label{eq:hi-b}
\end{equation}
On the fixed curvature bundle of Section~\ref{sec:hi-uhlenbeck}, set
\begin{equation}
                  \mathcal A(x,t)=b(t)\mathcal M(x,t).    \label{eq:hi-rescaled-curvature}
\end{equation}
The scale is constant in the spatial variables.  Using
Proposition~\ref{prop:hi-uhlenbeck-evolution} and quadratic homogeneity
$Q(cA)=c^2Q(A)$, we obtain
\begin{align}
 \partial_t\mathcal A
 &=\mathcal M+b\left(\Delta_{D^t,g(t)}\mathcal M
                         +Q(\mathcal M)\right) \notag\\
 &=\Delta_{D^t,g(t)}\mathcal A
     +\frac1{b(t)}\left(\mathcal A+Q(\mathcal A)\right)   \notag\\
 &=\Delta_{D^t,g(t)}\mathcal A
     +\frac1{b(t)}H(\mathcal A).                          \label{eq:hi-rescaled-pde}
\end{align}
This is precisely a system of the form treated in Chapter~2.  Its reaction
is the smooth time-dependent vertical map
\[
                 F(A,t)=b(t)^{-1}H(A).
\]
Its drift field is the identically zero field $X(t)\equiv0$.
Corollary~\ref{cor:hi-time-multiple-ode}, applied with
$a(t)=b(t)^{-1}>0$, verifies ODE invariance from every starting time, as
required by Theorem~\ref{thm:hamilton-wmp-systems}.

We check the initial condition explicitly.  Since $b(0)=C^{-1}$,
\[
 \nu(\mathcal A(x,0))=C^{-1}\nu(\mathcal M(x,0))\geq-1.
\]
The trace--least-eigenvalue inequality
\eqref{eq:hi-trace-min} gives
\[
 \trace\mathcal A(x,0)
   \geq3\nu(\mathcal A(x,0))\geq-3.
\]
If the conditional inequality in
\eqref{eq:hi-fixed-set} is active initially, then the two inequalities
$\nu(\mathcal A(x,0))\geq-1$ and
$\nu(\mathcal A(x,0))\leq-1$ force equality.  Its right side is then
\[
             1(\log1-3)=-3,
\]
so the trace bound proves the conditional inequality as well.  Hence
\begin{equation}
                  \mathcal A(x,0)\in K_{\mathrm{HI}}
                  \qquad(x\in M).                         \label{eq:hi-initial-membership}
\end{equation}

All hypotheses of Theorem~\ref{thm:hamilton-wmp-systems} are now explicit:
the base is closed; $\mathcal E$ and its fiber metric are fixed; $g(t)$ is
a smooth family of base metrics; $D^t$ is a smooth compatible family of
connections; the drift is $X\equiv0$; $K_{\mathrm{HI}}$ is a parallel closed convex family by
Corollary~\ref{cor:hi-associated-fixed-family}; the reaction is smooth; its
fiber ODE preserves the fixed set from every starting time; and
\eqref{eq:hi-initial-membership} holds.  Applying that theorem to
\eqref{eq:hi-rescaled-pde} yields
\begin{equation}
                  \mathcal A(x,t)\in K_{\mathrm{HI}}
                  \qquad(x\in M,\ 0\leq t<T).             \label{eq:hi-preserved-membership}
\end{equation}
When $T=\infty$, this conclusion is obtained on each finite interval and
hence at every finite time.

\begin{proof}[Proof of Theorem~\ref{thm:hamilton-ivey-scaled}]
Only the algebraic extraction of the estimate remains.  Fix a point where
$\nu=\nu(\mathcal M)<0$, and put
\begin{equation}
 q=-\nu(\mathcal A)=b(t)(-\nu(\mathcal M))>0.              \label{eq:hi-q-rescaled}
\end{equation}
Also $\trace\mathcal A=b(t)R$.

If $q\geq1$, membership
\eqref{eq:hi-preserved-membership} and the second condition in
\eqref{eq:hi-fixed-set} give
\[
 b(t)R\geq q(\log q-3).
\]
If $0<q<1$, the eigenvalue ordering gives the stronger elementary bound
\[
 b(t)R=\trace\mathcal A
     \geq3\nu(\mathcal A)=-3q.
\]
Since $\log q<0$,
\[
                  -3q\geq q(\log q-3),
\]
so the same conclusion holds.  Dividing by $b(t)>0$ and using
\eqref{eq:hi-q-rescaled} gives
\[
 R\geq(-\nu)
       \left[\log\!\left(b(t)(-\nu)\right)-3\right],
\]
which is \eqref{eq:hi-main-scaled}.
\end{proof}

\begin{proof}[Proof of Corollary~\ref{cor:hamilton-ivey-normalized}]
Take $C=1$ in Theorem~\ref{thm:hamilton-ivey-scaled} and use
\[
       \log((1+t)(-\nu))=\log(1+t)+\log(-\nu),
\]
whose two arguments are positive.
\end{proof}

\begin{remark}[Where the time improvement comes from]
\label{rem:hi-source-time-term}
The conventional proof uses the moving family
\[
                  K(t)=b(t)^{-1}K_{\mathrm{HI}}.
\]
Equation \eqref{eq:hi-rescaled-curvature} pulls that family back to the one
fixed set $K_{\mathrm{HI}}$.  Undoing the rescaling turns
$\log(-\nu(\mathcal A))$ into
\[
 \log\!\left(b(t)(-\nu(\mathcal M))\right)
   =\log\!\left(\frac{-\nu}{C}\right)+\log(1+Ct).
\]
Thus the term $\log(1+Ct)$ is exactly the logarithm of the similarity
scale.  It is not inserted by a later comparison estimate.

Equivalently, introduce similarity time and its inverse by
\[
 \tau=\log(1+Ct),\qquad
 t(\tau)=C^{-1}(\mathrm e^\tau-1),
\]
and set
\[
 \bar g(\tau)=b(t(\tau))^{-1}g(t(\tau)),
 \qquad D^\tau=D^{t(\tau)}.
\]
The factor $b(t(\tau))$ is spatially constant.  Hence the Levi--Civita
connections of $\bar g(\tau)$ and $g(t(\tau))$ agree, while tracing the
bundle Hessian gives
\[
 \Delta_{D^\tau,\bar g(\tau)}
    =b(t(\tau))\Delta_{D^{t(\tau)},g(t(\tau))}.
\]
Since $dt/d\tau=b(t(\tau))$, equation
\eqref{eq:hi-rescaled-pde} becomes
\[
 \partial_\tau\mathcal A
    =\Delta_{D^\tau,\bar g(\tau)}\mathcal A+H(\mathcal A).
\]
The direct proof above is preferable for a first formalization because
Chapter~2 already permits a time-dependent reaction and therefore requires
only the elementary ODE time change in
Corollary~\ref{cor:hi-time-multiple-ode}.
\end{remark}

\begin{remark}[The auxiliary scalar floor]
\label{rem:hi-scalar-floor}
Membership in $K_{\mathrm{HI}}$ also gives
\[
                    R(x,t)\geq-\frac3{C^{-1}+t}.
\]
This deliberately weak floor is used to join the two pieces of the
Hamilton--Ivey region.  It is not advertised as the sharp lower estimate
obtainable from the scalar-curvature evolution.
\end{remark}

\section{Consequences used later in the project}
\label{sec:hi-consequences}

\begin{corollary}[Positive-time estimate without a chosen initial scale]
\label{cor:hi-positive-time}
Let $g(t)$ be any Ricci flow on a closed three-manifold, defined for
$0\leq t<T$.  At every point with $t>0$ and $\nu(x,t)<0$,
\begin{equation}
        R(x,t)\geq
        (-\nu(x,t))
        \left[\log\!\left(t(-\nu(x,t))\right)-3\right].    \label{eq:hi-positive-time}
\end{equation}
\end{corollary}

\begin{proof}
Continuity of $\nu(\cdot,0)$ and compactness of $M$ give a finite $C>0$
such that $\nu(\cdot,0)\geq-C$.  Since $C^{-1}+t\geq t>0$ and the logarithm
is increasing, \eqref{eq:hi-main-scaled} implies
\eqref{eq:hi-positive-time}.  No symbol $C=\infty$ and no extended-real
convention is needed.
\end{proof}

\begin{corollary}[Large negative curvature is negligible at blow-up scale]
\label{cor:hi-negative-negligible}
Under the hypotheses of Theorem~\ref{thm:hamilton-ivey-scaled}, let
$(x_i,t_i)$ be spacetime points with
\[
 q_i=-\nu(x_i,t_i)>0,
 \qquad
 (C^{-1}+t_i)q_i\longrightarrow\infty.
\]
Then $R(x_i,t_i)>0$ for all sufficiently large $i$ and
\begin{equation}
             \frac{q_i}{R(x_i,t_i)}\longrightarrow0.      \label{eq:hi-negative-ratio}
\end{equation}
\end{corollary}

\begin{proof}
Theorem~\ref{thm:hamilton-ivey-scaled} gives
\[
 \frac{R(x_i,t_i)}{q_i}
   \geq\log\!\left((C^{-1}+t_i)q_i\right)-3
   \longrightarrow\infty.
\]
The two conclusions follow.
\end{proof}

\begin{corollary}[Closed ancient three-dimensional flows]
\label{cor:hi-ancient-nonnegative}
Every smooth ancient Ricci flow on a closed three-manifold has nonnegative
curvature operator at every spacetime point.
\end{corollary}

\textcolor{red}{Since the above is a major corollary, put the Lean code for the statement here.}

\begin{proof}
Fix a spacetime point $(x,t_0)$.  If $\nu(x,t_0)<0$, restrict the ancient
solution to $[s,t_0]$ for any $s<t_0$ and translate $s$ to time zero.
Corollary~\ref{cor:hi-positive-time}, with elapsed time $t_0-s$, gives
\[
 R(x,t_0)\geq(-\nu(x,t_0))
 \left[
   \log\!\left((t_0-s)(-\nu(x,t_0))\right)-3
 \right].
\]
The left side is fixed, while the right side tends to $+\infty$ as
$s\to-\infty$, a contradiction.  Hence $\nu(x,t_0)\geq0$.  The point was
arbitrary.
\end{proof}

\begin{remark}[Scope of the consequence]
\label{rem:hi-closed-scope}
Corollary~\ref{cor:hi-ancient-nonnegative} is stated only for closed flows
because that is the scope of Chapter~2.  Complete noncompact versions
require a noncompact maximum principle or a localized Hamilton--Ivey
argument; they are not formal consequences of the present chapter.
Chapter~\ref{ch:three-dimensional-shrinkers} keeps this boundary explicit:
its unrestricted complete-shrinker theorem invokes the localized theorem,
whereas its closed-flow blow-up corollary uses the estimate proved here
before passing to the limit.
\end{remark}

\section{Formalization interfaces and dependency order}
\label{sec:hi-formalization}

The proof was written so that geometric evolution, finite-dimensional
algebra, convex analysis, and the parabolic maximum principle can be
formalized independently.  A faithful blueprint should refine the chapter
in the following dependency order.

\begin{enumerate}
\item \emph{Curvature convention.}
      Define the trace-normalized operator by
      \eqref{eq:hi-curvature-normalization}; prove that a constant-curvature
      metric has operator $2kI$, that the eigenvalues are twice sectional
      curvatures, and that the trace is the scalar curvature.
\item \emph{Uhlenbeck isometry.}
      Solve \eqref{eq:hi-uhlenbeck-ode} as a smooth family of pointwise
      linear ODEs and prove \eqref{eq:hi-fixed-metric}.  Do not introduce a
      global frame or an orientation.
\item \emph{Pulled-back connection.}
      Define \eqref{eq:hi-pulled-connection}, including
      $\iota_t^{-1}$, and prove compatibility with the fixed metric $h$.
      Construct the induced connection on
      $\operatorname{SymEnd}_h(\Lambda^2V)$.
\item \emph{Algebraic reaction.}
      Define $\sigma_2$, $A^\#$, and $Q$ by
      \eqref{eq:hi-sigma-two}--\eqref{eq:hi-reaction}.  Prove polynomial
      smoothness, quadratic homogeneity, orthogonal equivariance, and the
      diagonal formula \eqref{eq:hi-Q-diagonal}.
\item \emph{Geometric evolution input.}
      Import or prove Proposition~\ref{prop:hi-uhlenbeck-evolution} with
      exactly the same carrier, connection, Laplacian sign, and
      normalization.  This is the only tensor-evolution black box in the
      pinching proof.
\item \emph{Spectral infrastructure.}
      Prove the finite-dimensional spectral theorem in the required form,
      the Rayleigh characterization, continuity and concavity of the least
      eigenvalue, and \eqref{eq:hi-trace-min}.
\item \emph{Clipped logarithm.}
      Define $m$ and $G$ by \eqref{eq:hi-clipped-log}.  Prove positivity of
      the logarithm's argument, convexity, monotonicity, continuity, and
      the three-case identity used in \eqref{eq:hi-fixed-sublevel}.
\item \emph{Fixed spectral region.}
      Define $\mathcal K_{\mathrm{HI}}$ by
      \eqref{eq:hi-fixed-set}, prove
      \eqref{eq:hi-fixed-sublevel}, and then prove closedness, nonemptiness,
      convexity, and invariance under the full orthogonal group.
\item \emph{Associated bundle.}
      Use Corollary~\ref{cor:wmp-associated-bundle} to obtain the total
      parallel closed convex family $K_{\mathrm{HI}}\subset\mathcal E$.
\item \emph{ODE diagonalization and ordering.}
      Use uniqueness to keep one initial eigenbasis fixed and prove
      \eqref{eq:hi-order-one}--\eqref{eq:hi-order-two}.  This removes any
      need to choose differentiable eigenvectors through multiplicities.
\item \emph{Trace boundary.}
      Prove the identity \eqref{eq:hi-trace-Q}, its Cauchy--Schwarz lower
      bound, and strict inwardness at $\trace A=-3$.
\item \emph{Atomic Hamilton--Ivey algebra.}
      Prove \eqref{eq:hi-Pi} by the two exact polynomial identities
      \eqref{eq:hi-Pi-negative-mu} and
      \eqref{eq:hi-Pi-nonnegative-mu}.
\item \emph{Pinching functional.}
      Define $\Omega$ only under the proof that $q=-\nu>0$; prove
      \eqref{eq:hi-Omega-exact}, isolate the radial contribution $-1$, and
      derive \eqref{eq:hi-Omega-H}.
\item \emph{ODE invariance.}
      Prove Proposition~\ref{prop:hi-fixed-region-ode}, with separate
      branches for entry at $q=1$, an initial point already in $q>1$, and
      simultaneous contact with the trace floor.  Then prove the positive
      time-change corollary from an arbitrary starting time.
\item \emph{Similarity scaling.}
      Define $b$ and $\mathcal A$, prove the spatial constancy and
      positivity of $b$, and derive \eqref{eq:hi-rescaled-pde} using
      quadratic homogeneity.
\item \emph{Initial inclusion.}
      Prove \eqref{eq:hi-initial-membership} explicitly, including the
      equality case $\nu(\mathcal A(0))=-1$.
\item \emph{ODE-to-PDE step.}
      Instantiate Theorem~\ref{thm:hamilton-wmp-systems} with the fixed
      curvature bundle, drift $X\equiv0$, $K_{\mathrm{HI}}$, and reaction
      $b(t)^{-1}H$.
\item \emph{Extraction and scaling.}
      Split into $q\geq1$ and $0<q<1$, using
      $\trace A\geq3\nu(A)$ in the second branch.  Divide by $b>0$ and
      derive the scale-covariant, normalized, and positive-time statements.
\end{enumerate}

The recommended public theorem is
Theorem~\ref{thm:hamilton-ivey-scaled}.  The normalized theorem is a
specialization, and the sectional-normalized formula is a translation
lemma.  They should not be implemented as three unrelated analytic proofs.
Likewise, the moving family $b(t)^{-1}K_{\mathrm{HI}}$ is explanatory data,
not an additional maximum-principle interface.

\begin{remark}[Expected formalization bottlenecks]
\label{rem:hi-formal-bottlenecks}
The main geometric bottleneck is
Proposition~\ref{prop:hi-uhlenbeck-evolution}; the main algebraic
bottleneck is reusable spectral infrastructure for self-adjoint
endomorphisms of a real three-dimensional inner-product space.  Once those
are available, the pinching proof is finite-dimensional algebra followed
by the already established Chapter~2 theorem.  No PDE eigenvector field,
measurable eigenbasis, Hodge-star identification, moving convex-set theorem,
or spacetime connection is required.
\end{remark}

\section{Source record and reconstruction notes}
\label{sec:hi-sources}

The historically precise attribution distinguishes the qualitative
pinching phenomenon from the time-improved formula.  R.~S.~Hamilton,
\emph{The formation of singularities in the Ricci flow}, Surveys in
Differential Geometry \textbf{2} (1995), 7--136,
\href{https://doi.org/10.4310/SDG.1993.v2.n1.a2}{Theorem~24.4},
and T.~A.~Ivey,
\emph{Ricci solitons on compact three-manifolds}, Differential Geometry and
its Applications \textbf{3} (1993), 301--307,
\href{https://doi.org/10.1016/0926-2245(93)90008-O}{Theorem~2},
gave earlier qualitative, time-independent forms independently.  The exact
term $\log(1+t)$ used here is due to Hamilton,
\emph{Non-singular solutions of the Ricci flow on three-manifolds},
Communications in Analysis and Geometry \textbf{7} (1999), 695--729,
\href{https://doi.org/10.4310/CAG.1999.v7.n4.a2}{Theorem~4.1},
pp.~699--701.  The conventional name Hamilton--Ivey estimate is retained,
but the time improvement is specifically sourced to Hamilton's 1999
theorem.

Hamilton's \emph{Four-manifolds with positive curvature operator}, Journal
of Differential Geometry \textbf{24} (1986), 153--179,
\href{https://doi.org/10.4310/jdg/1214440433}{online DOI record},
is a primary source for Uhlenbeck's trick and the systems maximum principle,
not for the later time-improved pinching formula.  That distinction matters
for provenance.

Among the project source books, the following are the principal
cross-checks.
\begin{enumerate}
\item \emph{The Ricci Flow: An Introduction}, Mathematical Surveys and
      Monographs~110: Uhlenbeck's trick and the curvature reaction on
      pp.~180--189; Theorem~9.4 and Lemmas~9.5--9.6 on pp.~258--260.
\item \emph{Hamilton's Ricci Flow}, Graduate Studies in Mathematics~77:
      Uhlenbeck's trick and the three-dimensional operator on pp.~121--123;
      the systems theorem and ODE on pp.~131--132; Proposition~6.43 and
      Theorem~6.44 on pp.~241--243.
\item \emph{The Ricci Flow: Techniques and Applications, Part II:
      Analytic Aspects}, Mathematical Surveys and Monographs~144:
      fixed-bundle setup and curvature ODE on pp.~9--15; the moving-set
      theorem and Hamilton--Ivey Theorems~10.16--10.18 on pp.~17--19.
\item \emph{The Ricci Flow: Techniques and Applications, Part I:
      Geometric Aspects}, Mathematical Surveys and
      Monographs~135: the corrected abstract Uhlenbeck construction in
      Appendix~A, Lemma~A.23 and Theorem~A.24, pp.~457--459, and
      Theorem~A.31 on p.~462.
\item \emph{Ricci Solitons in Dimensions Four and Higher}, Mathematical
      Surveys and Monographs~293: the scale-covariant statement in
      Theorem~1.20, pp.~12--13.
\end{enumerate}
The localized theorem in Graduate Studies in Mathematics~235,
Section~5.1, is useful for complete noncompact flows but introduces cutoff
and distance-distortion machinery unnecessary for the closed-manifold
result.  Mathematical Surveys and Monographs~163 and~206 use the estimate
downstream rather than supply the proof adopted here.

The proof in this chapter is a reconstruction, not a transcription.
The cited sources use the moving sets
$b(t)^{-1}K_{\mathrm{HI}}$ and a maximum principle for time-dependent
convex families.  The homothetic fixed-set reduction
\eqref{eq:hi-rescaled-curvature} was not found in the inspected primary or
project sources.  It is presented as a formalization-oriented
reorganization of Hamilton's region, not as a claim of mathematical
novelty.

Several source compressions have been repaired explicitly.
\begin{enumerate}
\item The curvature eigenvalues are declared to be twice sectional
      curvatures, so that $R=\lambda+\mu+\nu$ and
      \eqref{eq:hi-Q-diagonal} has the displayed coefficients.
\item The pulled-back connection includes $\iota_t^{-1}$ and is proved
      compatible with the fixed metric.
\item The convexity proof uses $-\nu(A)$ globally, not
      $|\nu(A)|$.  The latter agrees with $-\nu$ only on the negative
      branch.
\item ODE preservation starts at every time.  A trajectory which begins
      with $-\nu>1$ is initialized by membership in the fixed set; it is
      not assumed to have a preceding threshold-crossing time.
\item The branch $0<b(t)(-\nu)<1$ is proved from
      $R\geq3\nu$, rather than from the weaker scalar floor.
\item Repeated eigenvalues are handled by a fixed diagonal basis and the
      multiplicative difference equations
      \eqref{eq:hi-order-one}--\eqref{eq:hi-order-two}.
\item No global parallelization, orientation, or radial extension of a
      spatial eigenvector is used.
\item The positive-time statement chooses a finite lower bound by
      compactness; it does not encode an infinite normalization parameter.
\end{enumerate}

This chapter supplies a complete human proof of the closed-manifold,
time-improved Hamilton--Ivey estimate.  Its exact fixed-gauge curvature
input, Proposition~\ref{prop:hi-uhlenbeck-evolution}, is proved in
Chapter~\ref{ch:curvature-evolution}; its ODE-to-PDE input is the maximum
principle proved in Chapter~2.  It is the mathematical source for later
blueprint and Lean work; it makes no claim that those formal artifacts have
already been implemented or compiled.

\chapter{Curvature Evolution in Exact Uhlenbeck Gauge}
\label{ch:curvature-evolution}
\index{Uhlenbeck's trick}
\index{curvature operator!evolution}
\index{rough Laplacian!naturality}

\section{Purpose and exact scope}
\label{sec:ue-purpose}

Chapter~\ref{ch:hamilton-ivey} used one geometric evolution equation as an
explicitly isolated input: after transporting the curvature operator to a
fixed Euclidean vector bundle, it satisfies
\[
       \partial_t\mathcal M
          =\Delta_{D^t,g(t)}\mathcal M
             +\mathcal M^2+\mathcal M^\#.
\]
The present chapter proves that statement with exactly the carrier, metric,
connection, Laplacian sign, curvature normalization, and ordinary time
derivative used in Chapters~\ref{ch:weak-maximum-principle-systems} and
\ref{ch:hamilton-ivey}.  It therefore discharges
Proposition~\ref{prop:hi-uhlenbeck-evolution} rather than replacing it by a
nearby convention-dependent formula.

There are three logically separate parts.  First, an arbitrary smooth
family of base metrics is frozen by a pointwise linear ODE.  Second, the
resulting bundle isometry is used to prove naturality of spatial covariant
derivatives, Hessians, and rough Laplacians.  Third, for Ricci flow, the four
terms obtained by differentiating the four bundle-isometry factors cancel
the four Ricci contractions in the raw covariant curvature evolution.  A
finite-dimensional calculation then turns the remaining quadratic
four-tensor into the orientation-free polynomial used in Chapter~3.

The bundle calculations in this chapter are local in space.  The manifold
need not be compact, connected, or orientable.  Compactness enters only
when a later theorem invokes the global maximum principle of Chapter~2.
No global frame is chosen.  A temporary orthonormal basis is used in one
fiber to verify a polynomial identity and is discarded immediately.

\begin{roadmap}{Section~\ref{sec:ue-conventions} freezes the tensor and
operator conventions.}
Section~\ref{sec:ue-freezing} proves the metric-freezing and connection
lemmas for an arbitrary metric family.  Section~\ref{sec:ue-naturality}
proves spatial naturality.  Sections~\ref{sec:ue-raw} and
\ref{sec:ue-cancellation} derive the raw tensor identity and prove
the Uhlenbeck cancellation.  Section~\ref{sec:ue-operator} converts the
result to the exact three-dimensional curvature-operator equation.
Sections~\ref{sec:ue-covariance}--\ref{sec:ue-regression} prove covariance
and normalization checks.  The final sections record the formalization
dependency order and the source corrections.
\end{roadmap}

\section{Curvature and Laplacian conventions}
\label{sec:ue-conventions}

Let $M$ be a smooth manifold and let $g$ be a Riemannian metric.  As in
Chapter~3, our curvature endomorphism convention is
\begin{equation}
 R^g(X,Y)Z
   =\nabla_X^g\nabla_Y^gZ-\nabla_Y^g\nabla_X^gZ
      -\nabla_{[X,Y]}^gZ,                                \label{eq:ue-curvature-sign}
\end{equation}
with the sign for which the round sphere has positive sectional curvature.
For the component calculation we distinguish the covariant curvature tensor
from the curvature operator by writing
\begin{equation}
 \mathbf R_g(X,Y,Z,W)
     =g\bigl(R^g(X,Y)Z,W\bigr).                          \label{eq:ue-covariant-rm}
\end{equation}
Thus, in an orthonormal frame on a positively curved round sphere,
\[
                      R_{1221}>0.
\]

The trace-normalized curvature operator of Chapter~3 is the self-adjoint
endomorphism $\mathcal R_g$ of $\Lambda^2TM$ characterized by
\begin{equation}
 \left\langle
    \mathcal R_g(X\wedge Y),Z\wedge W
 \right\rangle_{\Lambda^2g}
   =2\,\mathbf R_g(X,Y,W,Z).                             \label{eq:ue-tensor-operator-bridge}
\end{equation}
The final two arguments on the right side of
\eqref{eq:ue-tensor-operator-bridge} are intentionally reversed relative to
\eqref{eq:ue-covariant-rm}.  Hence
\[
 \left\langle
   \mathcal R_g(e_i\wedge e_j),e_i\wedge e_j
 \right\rangle=2R_{ijji}=2K(e_i,e_j).
\]
This is the convention in which the curvature-operator eigenvalues are
twice the principal sectional curvatures and their sum is the scalar
curvature.

\begin{definition}[Hessian and Laplacian on bundles]
\label{def:ue-bundle-hessian-laplacian}
Let $E\to M$ be a vector bundle with connection $D$, and let $g$ be the
metric used to trace base indices.  For $s\in\Gamma(E)$ and smooth vector
fields $X,Y\in\Gamma(TM)$, set
\begin{align}
 (D^2s)(X,Y)
   &=D_X(D_Ys)-D_{\nabla_X^gY}s,                         \label{eq:ue-rough-hessian}\\
 \Delta_{D,g}s
   &=\trace_g(D^2s).                                    \label{eq:ue-rough-laplacian}
\end{align}
The same definition is used on all induced tensor and endomorphism bundles.
\end{definition}

\textcolor{red}{Since the above is a major definition, put the Lean code for the statement here.}

If $e_1,\ldots,e_n$ is orthonormal at a point, then
\begin{equation}
 \Delta_{D,g}s
   =\sum_{i=1}^n
       \left(D_{e_i}D_{e_i}s-D_{\nabla^g_{e_i}e_i}s\right)              \label{eq:ue-laplacian-frame}
\end{equation}
at that point.  Our sign is therefore the one for which $\Delta f\leq0$ at
a local maximum of a smooth real-valued function.

\section{Freezing an arbitrary metric family}
\label{sec:ue-freezing}

The metric-freezing part of Uhlenbeck's trick does not use the Ricci-flow
equation.  We state it at its natural level because the same construction
will be reused for other geometric tensor equations.

\begin{lemma}[Metric-freezing ODE]
\label{lem:ue-metric-freezing}
Let $I\subset\mathbb R$ be an interval containing $0$, let $g(t)$,
$t\in I$, be a smooth family of Riemannian metrics on a smooth
$n$-manifold $M$, and let $(V,h)\to M$ be a rank-$n$ real vector bundle
with a time-independent positive-definite fiber metric.  Suppose
\[
             \iota_0:(V,h)\longrightarrow(TM,g(0))
\]
is a smooth bundle isometry.  Define the $g(t)$-self-adjoint endomorphism
$S_t:TM\to TM$ by
\begin{equation}
       g(t)(S_tX,Y)=-\frac12\,\partial_tg(t)(X,Y),        \label{eq:ue-general-gauge-velocity}
\end{equation}
and solve, in every fiber,
\begin{equation}
       \partial_t\iota_t=S_t\circ\iota_t,
       \qquad \iota_t|_{t=0}=\iota_0.                   \label{eq:ue-general-iota-ode}
\end{equation}
Then $\iota_t$ is a smooth bundle isometry
\[
             \iota_t:(V,h)\longrightarrow(TM,g(t))
\]
for every $t\in I$.
\end{lemma}

\begin{proof}
In any local trivialization, \eqref{eq:ue-general-iota-ode} is a linear
matrix ODE whose coefficients depend smoothly on $(x,t)$.  Existence,
uniqueness, and smooth dependence on the parameter $x$ give a smooth family
$\iota_t$ on the full interval on which $g(t)$ is defined.

Fix $x\in M$ and time-independent vectors $v,w\in V_x$.  Equations
\eqref{eq:ue-general-gauge-velocity} and
\eqref{eq:ue-general-iota-ode} give
\begin{align*}
 \frac d{dt}g(t)(\iota_tv,\iota_tw)
 &=\partial_tg(t)(\iota_tv,\iota_tw)
   +g(t)(S_t\iota_tv,\iota_tw)
   +g(t)(\iota_tv,S_t\iota_tw)\\
 &=0.
\end{align*}
At $t=0$ the displayed inner product is $h(v,w)$, so
\begin{equation}
                         \iota_t^*g(t)=h.                \label{eq:ue-frozen-metric}
\end{equation}
Consequently $\iota_t$ is injective.  Since its domain and codomain have
the same finite rank, it is bijective in every fiber, and its inverse is
smooth.  This proves the assertion.
\end{proof}

\begin{corollary}[Ricci-flow gauge velocity]
\label{cor:ue-ricci-gauge-velocity}
If $g(t)$ solves
\[
                         \partial_tg=-2\Ric(g),
\]
then $S_t=\Ric_{g(t)}^\sharp$, and the freezing ODE is
\begin{equation}
       \partial_t\iota_t
          =\Ric_{g(t)}^\sharp\circ\iota_t.               \label{eq:ue-ricci-iota-ode}
\end{equation}
\end{corollary}

\begin{proof}
Raising one index in $\Ric$ is exactly the defining identity
\eqref{eq:ue-general-gauge-velocity} when
$\partial_tg=-2\Ric$.
\end{proof}

\begin{definition}[Pulled-back connection]
\label{def:ue-pulled-connection}
Under the hypotheses of Lemma~\ref{lem:ue-metric-freezing}, define a
connection $D^t$ on $V$ by
\begin{equation}
 D_X^t\xi
   =\iota_t^{-1}\!\left(
        \nabla_X^{g(t)}(\iota_t\xi)
     \right).                                            \label{eq:ue-pulled-connection}
\end{equation}
Here $X$ is a vector field on $M$ and $\xi$ is a section of $V$.  The
inverse $\iota_t^{-1}$ is part of the definition; without it the right side
is a tangent vector field rather than a section of $V$.
\end{definition}

\begin{lemma}[Metric compatibility and induced connections]
\label{lem:ue-connection-compatible}
For every $t\in I$, the connection $D^t$ is compatible with the fixed
fiber metric $h$.  It induces metric connections on $\Lambda^2V$, on all
tensor powers of $V$ and $V^*$, and on
\[
             \mathcal E=\operatorname{SymEnd}_h(\Lambda^2V).
\]
For $A\in\Gamma(\operatorname{End}(\Lambda^2V))$, the induced connection
is characterized by
\begin{equation}
       (D_X^tA)\omega
          =D_X^t(A\omega)-A(D_X^t\omega).                \label{eq:ue-end-connection}
\end{equation}
In particular, $D^t$ preserves the self-adjoint endomorphism subbundle
$\mathcal E$.
\end{lemma}

\begin{proof}
Using \eqref{eq:ue-frozen-metric}, the metric compatibility of the
Levi--Civita connection, and \eqref{eq:ue-pulled-connection}, we obtain
\begin{align*}
 Xh(\xi,\eta)
 &=Xg(t)(\iota_t\xi,\iota_t\eta)\\
 &=g(t)(\nabla_X^{g(t)}(\iota_t\xi),\iota_t\eta)
   +g(t)(\iota_t\xi,\nabla_X^{g(t)}(\iota_t\eta))\\
 &=h(D_X^t\xi,\eta)+h(\xi,D_X^t\eta).
\end{align*}
The induced connections are then defined by the tensor-product and dual
rules.  Formula \eqref{eq:ue-end-connection} is the usual induced
connection on an endomorphism bundle.  Differentiating the identity
$\langle A\omega,\zeta\rangle=\langle\omega,A\zeta\rangle$ and using metric
compatibility shows that $D_X^tA$ is self-adjoint whenever $A$ is.
\end{proof}

\section{Spatial naturality}
\label{sec:ue-naturality}

Fix a time $t$.  Let $F_t=\Lambda^2\iota_t$.  The definition of $D^t$ says
that $\iota_t$, and hence $F_t$, is parallel for the corresponding pair of
connections.  We record all derivative levels needed later rather than
using the phrase ``pullback commutes with the Laplacian.''

\begin{lemma}[First derivative, Hessian, and rough-Laplacian naturality]
\label{lem:ue-spatial-naturality}
Fix $t\in I$.  Let $T$ be a smooth covariant $r$-tensor on $TM$, and define
the $V$-tensor
\[
 \overline T(v_1,\ldots,v_r)
    =T(\iota_tv_1,\ldots,\iota_tv_r).
\]
Let $\nabla^{D,t}$ denote the mixed connection which uses the
Levi--Civita connection of $g(t)$ on base covariant indices and $D^t$ on
$V$-indices.  Then, for vector fields $X,Y$ on $M$,
\begin{align}
 \nabla_X^{D,t}\overline T
   &=\iota_t^*(\nabla_X^{g(t)}T),                         \label{eq:ue-first-naturality}\\
 (\nabla^{D,t})^2\overline T(X,Y)
   &=\iota_t^*\!\left((\nabla^{g(t)})^2T(X,Y)\right),   \label{eq:ue-hessian-naturality}\\
 \Delta_{D^t,g(t)}\overline T
   &=\iota_t^*(\Delta_{\nabla^{g(t)},g(t)}T).             \label{eq:ue-laplacian-naturality}
\end{align}
Likewise, if $A$ is an endomorphism of $\Lambda^2TM$ and
$\overline A=F_t^{-1}AF_t$, then
\begin{align}
 D_X^t\overline A
   &=F_t^{-1}(\nabla_X^{g(t)}A)F_t,                       \label{eq:ue-end-first-naturality}\\
 \Delta_{D^t,g(t)}\overline A
   &=F_t^{-1}(\Delta_{\nabla^{g(t)},g(t)}A)F_t.           \label{eq:ue-end-laplacian-naturality}
\end{align}
\end{lemma}

\begin{proof}
For any section $\xi$ of $V$, definition
\eqref{eq:ue-pulled-connection} is equivalent to
\begin{equation}
       \nabla_X^{g(t)}(\iota_t\xi)=\iota_t(D_X^t\xi).    \label{eq:ue-iota-parallel}
\end{equation}
Apply the definition of the covariant derivative of a covariant tensor to
$\overline T$.  Every derivative of a slot $\iota_t\xi$ produced on the
tangent-bundle side is, by \eqref{eq:ue-iota-parallel}, exactly the image
of the corresponding $D^t\xi$ term on the fixed-bundle side.  The terms
cancel in pairs and give \eqref{eq:ue-first-naturality}.

Apply \eqref{eq:ue-first-naturality} a second time and subtract the same
base correction
$\nabla^{D,t}_{\nabla_X^{g(t)}Y}$ on the left and
$\nabla^{g(t)}_{\nabla_X^{g(t)}Y}$ on the right.  This proves
\eqref{eq:ue-hessian-naturality}.  Tracing the two base indices with the
same inverse metric $g(t)^{-1}$ proves
\eqref{eq:ue-laplacian-naturality}.

The exterior-square version of \eqref{eq:ue-iota-parallel} is
\[
 \nabla_X^{\Lambda^2TM}(F_t\omega)
       =F_t(D_X^{\Lambda^2V}\omega).
\]
Expanding the induced endomorphism connection
\eqref{eq:ue-end-connection} gives
\eqref{eq:ue-end-first-naturality}; the preceding Hessian-and-trace
argument gives \eqref{eq:ue-end-laplacian-naturality}.
\end{proof}

\begin{remark}[Spatial versus temporal differentiation]
\label{rem:ue-space-time-separation}
Lemma~\ref{lem:ue-spatial-naturality} is a statement at one fixed time.
The family $D^t$ may vary with time, but no $\partial_tD^t$ occurs in the
rough Laplacian.  Conversely, after pullback, a time-dependent tensor is a
section of a fixed vector bundle, so its $\partial_t$ below is the ordinary
vertical derivative used in Chapter~2.  No spacetime connection is being
introduced implicitly.
\end{remark}

\section{Derivation of the raw covariant tensor evolution}
\label{sec:ue-raw}

We now specialize to Ricci flow.  To make the convention boundary auditable,
we derive the raw covariant formula in two stages: first the differential
variation and commutator calculation, and then the purely algebraic
$B$-tensor rewrite.  All indices in this section refer to a local frame held
fixed during time differentiation.

\begin{lemma}[Variation of the Levi--Civita connection]
\label{lem:ue-connection-variation}
For a smooth metric variation $\partial_tg=h$, the difference tensor
$A=\partial_t\nabla$ is characterized by
\begin{equation}
 2g(A_XY,Z)
  =(\nabla_Xh)(Y,Z)+(\nabla_Yh)(X,Z)-(\nabla_Zh)(X,Y).        \label{eq:ue-connection-variation}
\end{equation}
In local coordinates, under Ricci flow $h=-2\Ric$,
\begin{equation}
 \partial_t\Gamma_{ij}^{k}
  =-g^{k\ell}\bigl(
      \nabla_iR_{j\ell}+\nabla_jR_{i\ell}-\nabla_\ell R_{ij}
    \bigr).                                                \label{eq:ue-christoffel-variation}
\end{equation}
\end{lemma}

\begin{proof}
Differentiate the Koszul identity while keeping the vector fields fixed in
time.  Expanding the three covariant derivatives of $h$ cancels all
Lie-bracket terms and gives \eqref{eq:ue-connection-variation}.
Substitution of $h=-2\Ric$ and raising the final index gives
\eqref{eq:ue-christoffel-variation}.
\end{proof}

\begin{lemma}[Differential and commutator stage]
\label{lem:ue-three-quadratic-evolution}
With the convention \eqref{eq:ue-curvature-sign}, Ricci flow satisfies
\begin{align}
 \partial_tR_{ijk\ell}
 &=\Delta R_{ijk\ell}
   +g^{pq}\bigl(
       R_{ijp}{}^{r}R_{rqk\ell}
       -2R_{pik}{}^{r}R_{jqr\ell}
       +2R_{pir\ell}R_{jqk}{}^{r}
     \bigr)                                                \label{eq:ue-three-quadratic}\\
 &\quad-
  \bigl(
    \Ric_i{}^pR_{pjk\ell}
   +\Ric_j{}^pR_{ipk\ell}
   +\Ric_k{}^pR_{ijp\ell}
   +\Ric_\ell{}^pR_{ijkp}
  \bigr).\notag
\end{align}
\end{lemma}

\begin{proof}
For the mixed tensor $R_{ijk}{}^\ell$, curvature variation gives
\begin{equation}
 \partial_tR_{ijk}{}^\ell
   =\nabla_i(\partial_t\Gamma_{jk}^{\ell})
      -\nabla_j(\partial_t\Gamma_{ik}^{\ell}).             \label{eq:ue-curvature-variation}
\end{equation}
Insert \eqref{eq:ue-christoffel-variation}.  The contracted second Bianchi
identity converts each divergence of $\Rm$ into a difference of first
derivatives of $\Ric$.  Apply the second Bianchi identity once more to put
the two derivatives in the order occurring in the rough Laplacian.  Expand
every interchange of covariant derivatives by the Ricci commutation rule,
which contributes one curvature action on each tensor slot.  Collecting the
four slot contributions gives
\begin{align}
 \partial_tR_{ijk}{}^\ell
 &=\Delta R_{ijk}{}^\ell
   +g^{pq}\bigl(
       R_{ijp}{}^{r}R_{rqk}{}^\ell
       -2R_{pik}{}^{r}R_{jqr}{}^\ell
       +2R_{pir}{}^\ell R_{jqk}{}^{r}
     \bigr)                                                \label{eq:ue-mixed-rm-evolution}\\
 &\quad-
  \Ric_i{}^pR_{pjk}{}^\ell
  -\Ric_j{}^pR_{ipk}{}^\ell
  -\Ric_k{}^pR_{ijp}{}^\ell
  +\Ric_p{}^\ell R_{ijk}{}^p .\notag
\end{align}
Equation \eqref{eq:ue-mixed-rm-evolution} is a useful intermediate
formalization target: the signs of its last four terms are exactly the
curvature action on three covariant slots and one contravariant slot.

Finally, $R_{ijk\ell}=g_{\ell m}R_{ijk}{}^m$.  Differentiating the lowering
metric contributes $-2\Ric_{\ell m}R_{ijk}{}^m$.  One copy combines with
the last positive term of \eqref{eq:ue-mixed-rm-evolution}, leaving
$-\Ric_\ell{}^pR_{ijkp}$.  Metric compatibility gives
$g_{\ell m}\Delta R_{ijk}{}^m=\Delta R_{ijk\ell}$.  The result is
\eqref{eq:ue-three-quadratic}.
\end{proof}

\begin{lemma}[$B$-tensor rewrite]
\label{lem:ue-B-rewrite}
Define
\[
 B_{ijk\ell}=-g^{pr}g^{qs}R_{ipjq}R_{kr\ell s}.
\]
Then the three quadratic terms in
\eqref{eq:ue-three-quadratic} equal
\begin{equation}
 2\bigl(B_{ijk\ell}-B_{ij\ell k}
             +B_{ikj\ell}-B_{i\ell jk}\bigr).             \label{eq:ue-B-rewrite}
\end{equation}
\end{lemma}

\begin{proof}
The curvature symmetries give
$B_{ijk\ell}=B_{ji\ell k}=B_{k\ell ij}$.  They immediately rewrite the
last two terms in \eqref{eq:ue-three-quadratic} as
\[
 -2g^{pq}R_{pik}{}^rR_{jqr\ell}=2B_{ikj\ell},
 \qquad
 2g^{pq}R_{pir\ell}R_{jqk}{}^r=-2B_{i\ell jk}.
\]
For the first term, apply the first Bianchi identity in both factors:
\begin{align*}
 g^{pq}R_{ijp}{}^rR_{rqk\ell}
  &=-B_{jik\ell}+B_{ji\ell k}+B_{ijk\ell}-B_{ij\ell k}\\
  &=2(B_{ijk\ell}-B_{ij\ell k}).
\end{align*}
Adding the three identities proves \eqref{eq:ue-B-rewrite}.
\end{proof}

\begin{proposition}[Raw covariant curvature evolution]
\label{prop:ue-raw-curvature-evolution}
Let $g(t)$ be a smooth solution of $\partial_tg=-2\Ric(g)$ on an
$n$-manifold.  Use \eqref{eq:ue-curvature-sign} and define
$R_{ijk\ell}=\mathbf R_g(e_i,e_j,e_k,e_\ell)$ in a local frame
$(e_i)$ which is held fixed during time differentiation.  Equivalently,
the equation below is an equality of time-dependent covariant
four-tensors, so no frame-motion terms are included.  Set
\begin{equation}
 B_{ijk\ell}
   =-g^{pr}g^{qs}R_{ipjq}R_{kr\ell s}.                  \label{eq:ue-B-unpulled}
\end{equation}
Then
\begin{align}
 \partial_tR_{ijk\ell}
 &=\Delta R_{ijk\ell}
   +2\bigl(
       B_{ijk\ell}-B_{ij\ell k}+B_{ikj\ell}-B_{i\ell jk}
     \bigr)                                             \label{eq:ue-raw-rm-evolution}\\
 &\quad-
  \bigl(
    \Ric_i{}^pR_{pjk\ell}
   +\Ric_j{}^pR_{ipk\ell}
   +\Ric_k{}^pR_{ijp\ell}
   +\Ric_\ell{}^pR_{ijkp}
  \bigr).\notag
\end{align}
Here $\Delta$ is the rough Laplacian formed from $g(t)$ and its
Levi--Civita connection, with the sign in
\eqref{eq:ue-rough-laplacian}.
\end{proposition}

\textcolor{red}{Since the above is a major proposition, put the Lean code for the statement here.}

\begin{proof}
Lemma~\ref{lem:ue-three-quadratic-evolution} supplies the Laplacian, the
four Ricci-slot terms, and the three unreduced quadratic terms.  Substitute
the algebraic identity of Lemma~\ref{lem:ue-B-rewrite}.  The result is
exactly \eqref{eq:ue-raw-rm-evolution}.
\end{proof}

\section{The four-slot cancellation}
\label{sec:ue-cancellation}

Let $(V,h)$ and $\iota_t$ be as in
Lemma~\ref{lem:ue-metric-freezing}, with the Ricci-flow ODE
\eqref{eq:ue-ricci-iota-ode}.  Define the pulled covariant curvature tensor
\begin{equation}
 \overline{\mathbf R}_t(a,b,c,d)
   =\mathbf R_{g(t)}
      (\iota_ta,\iota_tb,\iota_tc,\iota_td).             \label{eq:ue-pulled-covariant-rm}
\end{equation}

\begin{lemma}[Four-slot time derivative]
\label{lem:ue-four-slot-derivative}
Let $e_a,e_b,e_c,e_d$ be local sections of $V$ which do not depend on
time, and write $\overline R_{abcd}$ for the components in
\eqref{eq:ue-pulled-covariant-rm}.  Let
\[
 \overline\Ric_{ab}
    =\Ric_{g(t)}(\iota_te_a,\iota_te_b),
 \qquad
 \overline\Ric_a{}^e=h^{ef}\overline\Ric_{af}.
\]
Then
\begin{align}
 \partial_t\overline R_{abcd}
 &=\bigl(\iota_t^*(\partial_t\mathbf R)\bigr)_{abcd}
   +\overline\Ric_a{}^e\overline R_{ebcd}
   +\overline\Ric_b{}^e\overline R_{aecd}                \label{eq:ue-four-slot-derivative}\\
 &\quad
   +\overline\Ric_c{}^e\overline R_{abed}
   +\overline\Ric_d{}^e\overline R_{abce}.\notag
\end{align}
\end{lemma}

\textcolor{red}{Since the above is a major lemma, put the Lean code for the statement here.}

\begin{proof}
Differentiate the five time-dependent factors in
\[
 \overline R_{abcd}
   =R_{ijk\ell}(\iota_t)^i{}_a(\iota_t)^j{}_b
                   (\iota_t)^k{}_c(\iota_t)^\ell{}_d.
\]
The derivative of $R_{ijk\ell}$ gives the first term on the right side of
\eqref{eq:ue-four-slot-derivative}.  Equation
\eqref{eq:ue-ricci-iota-ode} says
\[
       \partial_t(\iota_t)^i{}_a
          =\Ric^i{}_p(\iota_t)^p{}_a.
\]
Applying this once in each of the four slots gives, in order, the four
remaining terms in \eqref{eq:ue-four-slot-derivative}.  Raising the pulled
Ricci index with $h$ is legitimate because $\iota_t^*g(t)=h$.
\end{proof}

Define the quadratic tensor on $V$ by
\begin{equation}
 \overline B_{abcd}
   =-h^{eg}h^{fh}
       \overline R_{aebf}\overline R_{cgdh}.             \label{eq:ue-B-pulled}
\end{equation}

\begin{proposition}[Exact pulled covariant curvature evolution]
\label{prop:ue-pulled-tensor-evolution}
Under Ricci flow, the tensor \eqref{eq:ue-pulled-covariant-rm} satisfies
\begin{equation}
 \partial_t\overline R_{abcd}
  =\Delta_{D^t,g(t)}\overline R_{abcd}
   +2\bigl(
       \overline B_{abcd}-\overline B_{abdc}
       +\overline B_{acbd}-\overline B_{adbc}
     \bigr).                                             \label{eq:ue-pulled-tensor-pde}
\end{equation}
The derivative on the left is the ordinary time derivative in the fixed
bundle $(V^*)^{\otimes4}$.
\end{proposition}

\textcolor{red}{Since the above is a major proposition, put the Lean code for the statement here.}

\begin{proof}
Pull back equation \eqref{eq:ue-raw-rm-evolution} by the four copies of
$\iota_t$.  Because $\iota_t$ is an isometry, contraction by two copies of
$g(t)^{-1}$ becomes contraction by two copies of $h^{-1}$; therefore the
pullback of \eqref{eq:ue-B-unpulled} is exactly
\eqref{eq:ue-B-pulled}.  Lemma~\ref{lem:ue-spatial-naturality} changes the
pulled rough Laplacian into $\Delta_{D^t,g(t)}$.

It remains to inspect the zero-order terms.  The pullback of the first
negative Ricci contraction in \eqref{eq:ue-raw-rm-evolution} is
\[
             -\overline\Ric_a{}^e\overline R_{ebcd},
\]
which cancels the first slot term in
\eqref{eq:ue-four-slot-derivative}.  The second, third, and fourth negative
Ricci contractions cancel, respectively,
\[
 \overline\Ric_b{}^e\overline R_{aecd},\qquad
 \overline\Ric_c{}^e\overline R_{abed},\qquad
 \overline\Ric_d{}^e\overline R_{abce}.
\]
No terms remain from the motion of the bundle isometry.  The diffusion and
the four displayed quadratic terms are therefore exactly those in
\eqref{eq:ue-pulled-tensor-pde}.
\end{proof}

\section{The exact three-dimensional operator equation}
\label{sec:ue-operator}

Assume from now through Theorem~\ref{thm:ue-operator-evolution} that
$\dim M=3$.  Let
\[
                    F_t=\Lambda^2\iota_t.
\]
Since $\iota_t$ is an isometry, so is
\[
       F_t:(\Lambda^2V,\Lambda^2h)
              \longrightarrow
           (\Lambda^2TM,\Lambda^2g(t)).
\]

\begin{definition}[Fixed-bundle curvature operator]
\label{def:ue-fixed-curvature-operator}
Define
\begin{equation}
 \mathcal M(t)
    =F_t^{-1}\circ\mathcal R_{g(t)}\circ F_t
       \in\Gamma\bigl(\operatorname{SymEnd}_h(\Lambda^2V)\bigr).       \label{eq:ue-fixed-operator}
\end{equation}
Equivalently, for $a,b,c,d\in V_x$,
\begin{equation}
 \left\langle
    \mathcal M(a\wedge b),c\wedge d
 \right\rangle_h
    =2\overline{\mathbf R}(a,b,d,c).                    \label{eq:ue-fixed-operator-tensor}
\end{equation}
\end{definition}

The equivalence follows by applying the isometry $F_t$ to both entries of
the inner product and then using
\eqref{eq:ue-tensor-operator-bridge}.  It also proves directly that
$\mathcal M$ is self-adjoint.

For a self-adjoint endomorphism $A$ of a three-dimensional Euclidean
space, retain the Chapter~3 definitions
\begin{align}
 \sigma_2(A)
   &=\frac12\left((\trace A)^2-\trace(A^2)\right),        \label{eq:ue-sigma-two}\\
 A^\#
   &=A^2-(\trace A)A+\sigma_2(A)I,                       \label{eq:ue-sharp-polynomial}\\
 Q(A)&=A^2+A^\#.                                        \label{eq:ue-Q-polynomial}
\end{align}
This is an everywhere-defined polynomial; $A^\#$ is the degree-two
adjugate, not the Hilbert-space adjoint.

\begin{lemma}[Three-dimensional quadratic conversion]
\label{lem:ue-three-dimensional-reaction}
Let $(V,h)$ be a three-dimensional Euclidean vector space.  Let
$\overline{\mathbf R}$ be an algebraic curvature tensor and let $A$ be its
trace-normalized curvature operator, related by
\[
 \langle A(a\wedge b),c\wedge d\rangle_h
       =2\overline{\mathbf R}(a,b,d,c).
\]
Form $\overline B$ by \eqref{eq:ue-B-pulled}, and set
\begin{equation}
 P_{abcd}
   =2\bigl(
       \overline B_{abcd}-\overline B_{abdc}
       +\overline B_{acbd}-\overline B_{adbc}
     \bigr).                                             \label{eq:ue-P-quadratic-tensor}
\end{equation}
Let $\Psi(P)$ be the unique self-adjoint endomorphism characterized by
\begin{equation}
 \langle\Psi(P)(a\wedge b),c\wedge d\rangle_h
       =2P(a,b,d,c).                                     \label{eq:ue-P-operator-bridge}
\end{equation}
Then
\begin{equation}
                            \Psi(P)=Q(A).                \label{eq:ue-reaction-conversion}
\end{equation}
\end{lemma}

\begin{proof}
The curvature symmetries give
\[
       \overline B_{abcd}=\overline B_{badc}
          =\overline B_{cdab}.
\]
Together with the first Bianchi identity, these relations show directly
that the combination $P$ in \eqref{eq:ue-P-quadratic-tensor} is
antisymmetric in each pair, symmetric under interchange of the pairs, and
satisfies the first Bianchi identity.  Thus
\eqref{eq:ue-P-operator-bridge} defines a unique self-adjoint endomorphism
$\Psi(P)$.

Both sides of \eqref{eq:ue-reaction-conversion} are quadratic polynomial
functions of $A$ and commute with orthogonal changes of the underlying
basis.  We nevertheless give the full diagonal calculation, since it fixes
the factor of two.

By the spectral theorem choose an orthonormal eigenbasis
$\omega_1,\omega_2,\omega_3$ of $\Lambda^2V$.  The elementary
three-dimensional exterior-algebra basis lemma supplies an orthonormal
basis $e_1,e_2,e_3$ of $V$, after harmlessly changing signs of the
$\omega_i$, such that
\[
       \omega_1=e_2\wedge e_3,\qquad
       \omega_2=e_3\wedge e_1,\qquad
       \omega_3=e_1\wedge e_2.
\]
An orientation may be selected to name this basis, but no orientation is
retained in the statement or in the resulting polynomial.  Write
\[
              A=\operatorname{diag}(\lambda,\mu,\nu)
\]
in this basis.  Equation defining $A$ gives the three independent
sectional components
\begin{equation}
 \overline R_{2332}=\frac\lambda2,\qquad
 \overline R_{3113}=\frac\mu2,\qquad
 \overline R_{1221}=\frac\nu2,                           \label{eq:ue-diagonal-r-components}
\end{equation}
with all remaining components determined by the curvature symmetries.

Substitute \eqref{eq:ue-diagonal-r-components} into
\eqref{eq:ue-B-pulled} and then into
\eqref{eq:ue-P-quadratic-tensor}.  The finite contractions give
\begin{align}
 2P_{2332}&=\lambda^2+\mu\nu,                            \label{eq:ue-P-lambda}\\
 2P_{3113}&=\mu^2+\lambda\nu,                            \label{eq:ue-P-mu}\\
 2P_{1221}&=\nu^2+\lambda\mu,                            \label{eq:ue-P-nu}
\end{align}
and all off-diagonal entries of $\Psi(P)$ vanish.  Hence
\begin{equation}
 \Psi(P)=\operatorname{diag}
    (\lambda^2+\mu\nu,\,
     \mu^2+\lambda\nu,\,
     \nu^2+\lambda\mu).                                \label{eq:ue-diagonal-reaction}
\end{equation}
On the other hand, \eqref{eq:ue-sharp-polynomial} gives
\[
        A^\#=\operatorname{diag}(\mu\nu,\lambda\nu,\lambda\mu),
\]
so the right side of \eqref{eq:ue-diagonal-reaction} is $A^2+A^\#$.
This proves \eqref{eq:ue-reaction-conversion}.  Since the final expression
is the orthogonally equivariant polynomial
\eqref{eq:ue-Q-polynomial}, it is independent of every temporary basis and
orientation choice.
\end{proof}

\begin{theorem}[Exact Uhlenbeck-gauge curvature-operator evolution]
\label{thm:ue-operator-evolution}
Let $M^3$ be a smooth three-manifold, with no compactness, connectedness,
or orientability assumption.  Let $I\subset\mathbb R$ be an interval
containing $0$, and let $g(t)$, $t\in I$, be a smooth solution of
\[
                         \partial_tg=-2\Ric(g).
\]
Let $(V,h)$ be a rank-three real metric vector bundle, let
$\iota_0:(V,h)\to(TM,g(0))$ be a bundle isometry, let $\iota_t$ solve
\eqref{eq:ue-ricci-iota-ode}, and let $D^t$ and $\mathcal M(t)$ be defined
by \eqref{eq:ue-pulled-connection} and
\eqref{eq:ue-fixed-operator}.  Then $h$ is independent of time, $D^t$ is
$h$-compatible, and
\begin{equation}
 \boxed{
 \partial_t\mathcal M
     =\Delta_{D^t,g(t)}\mathcal M
       +\mathcal M^2+\mathcal M^\# .}
                                                               \label{eq:ue-main-operator-pde}
\end{equation}
Here $\partial_t\mathcal M$ is the ordinary derivative of a section of the
fixed bundle $\operatorname{SymEnd}_h(\Lambda^2V)$, the rough Laplacian has
the sign \eqref{eq:ue-rough-laplacian}, and $\mathcal M^\#$ is the
orientation-free polynomial \eqref{eq:ue-sharp-polynomial}.
\end{theorem}

\textcolor{red}{Since this is a major theorem, put the Lean code for the statement here.}

\begin{proof}
Lemma~\ref{lem:ue-metric-freezing} and
Lemma~\ref{lem:ue-connection-compatible} prove the assertions about the
fiber metric and connection.  Proposition~\ref{prop:ue-pulled-tensor-evolution}
gives the precise equation for $\overline{\mathbf R}$.  Because $h$ is fixed
and $D^th=0$, the linear tensor-to-operator correspondence
\eqref{eq:ue-fixed-operator-tensor} is independent of time and parallel in
space.  It therefore takes $\partial_t\overline{\mathbf R}$ to
$\partial_t\mathcal M$ and, by
Lemma~\ref{lem:ue-spatial-naturality}, takes the tensor rough Laplacian to
$\Delta_{D^t,g(t)}\mathcal M$.  Finally,
Lemma~\ref{lem:ue-three-dimensional-reaction} takes the quadratic tensor
in \eqref{eq:ue-pulled-tensor-pde} to
$Q(\mathcal M)=\mathcal M^2+\mathcal M^\#$.  These are exactly the three
terms in \eqref{eq:ue-main-operator-pde}.
\end{proof}

\begin{corollary}[Discharge of the Hamilton--Ivey evolution input]
\label{cor:ue-discharge-hamilton-ivey}
Proposition~\ref{prop:hi-uhlenbeck-evolution} holds with all of the
conventions stated in Chapter~3.
\end{corollary}

\begin{proof}
Chapter~3 uses the same bundle $(V,h)$, the same isometry ODE
\eqref{eq:ue-ricci-iota-ode}, the same pulled connection
\eqref{eq:ue-pulled-connection}, the same conjugation definition
\eqref{eq:ue-fixed-operator}, the rough-Laplacian convention
\eqref{eq:ue-rough-laplacian}, and the polynomial
\eqref{eq:ue-Q-polynomial}.  Thus
Theorem~\ref{thm:ue-operator-evolution} is exactly, and not merely up to a
normalization, the equation asserted in that proposition.
\end{proof}

\section{Gauge, scaling, and nullspace covariance}
\label{sec:ue-covariance}

The next statements ensure that the construction exposes geometric data
rather than an artifact of the initial bundle identification.

\begin{proposition}[Gauge covariance]
\label{prop:ue-gauge-covariance}
Suppose $(\widehat V,\widehat h,\widehat\iota_0)$ is a second choice in
Theorem~\ref{thm:ue-operator-evolution}.  Define the time-independent
bundle isometry
\[
       U=\iota_0^{-1}\circ\widehat\iota_0:
          (\widehat V,\widehat h)\longrightarrow(V,h).
\]
Then, for all $t\in I$,
\begin{align}
 \widehat\iota_t&=\iota_t\circ U,                        \label{eq:ue-gauge-iota}\\
 \widehat D_X^t\xi
   &=U^{-1}\!\left(D_X^t(U\xi)\right),                   \label{eq:ue-gauge-connection}\\
 \widehat{\mathcal M}
   &=(\Lambda^2U)^{-1}\mathcal M(\Lambda^2U).            \label{eq:ue-gauge-operator}
\end{align}
Consequently equation \eqref{eq:ue-main-operator-pde} for one choice is
carried to the same equation for the other choice.
\end{proposition}

\begin{proof}
Both $\widehat\iota_t$ and $\iota_tU$ solve the same fiberwise linear ODE
\eqref{eq:ue-ricci-iota-ode} and have the same initial value.  Uniqueness
gives \eqref{eq:ue-gauge-iota}.  Substitution into
\eqref{eq:ue-pulled-connection} gives
\eqref{eq:ue-gauge-connection}, including the spatial derivative of $U$
inside $D_X^t(U\xi)$.  Conjugating the geometric curvature operator gives
\eqref{eq:ue-gauge-operator}.  The induced connection Laplacian and the
orthogonally equivariant polynomial $Q$ obey the same conjugation law, so
the full equation is gauge covariant.
\end{proof}

\begin{proposition}[Parabolic-scaling covariance]
\label{prop:ue-parabolic-scaling}
Let the data be as in Theorem~\ref{thm:ue-operator-evolution}, fix
$t_0\in I$ and $\rho>0$, and define
\[
 \widetilde I
    =\{s\in\mathbb R:t_0+s/\rho\in I\},
 \qquad
 \widetilde g(s)=\rho\,g(t_0+s/\rho).
\]
On the same fixed metric bundle $(V,h)$, set
\begin{equation}
 \widetilde\iota_s
    =\rho^{-1/2}\iota_{t_0+s/\rho}.                      \label{eq:ue-scaled-iota}
\end{equation}
Then $\widetilde\iota_s$ is the Uhlenbeck isometry for
$\widetilde g(s)$, and, writing $t=t_0+s/\rho$,
\begin{align}
 \widetilde D^{\,s}&=D^t,                               \label{eq:ue-scaled-connection}\\
 \widetilde{\mathcal M}(s)&=\rho^{-1}\mathcal M(t),      \label{eq:ue-scaled-operator}\\
 \partial_s\widetilde{\mathcal M}
   &=\rho^{-2}\partial_t\mathcal M,                      \label{eq:ue-scaled-time}\\
 \Delta_{\widetilde D^{\,s},\widetilde g(s)}
        \widetilde{\mathcal M}
   &=\rho^{-2}\Delta_{D^t,g(t)}\mathcal M,               \label{eq:ue-scaled-laplacian}\\
 Q(\widetilde{\mathcal M})&=\rho^{-2}Q(\mathcal M).     \label{eq:ue-scaled-reaction}
\end{align}
Thus every term in \eqref{eq:ue-main-operator-pde} has parabolic scaling
weight $\rho^{-2}$.
\end{proposition}

\begin{proof}
Constant spatial rescaling does not change the Levi--Civita connection or
the Ricci tensor as a covariant two-tensor.  It changes the Ricci
endomorphism by
\[
       \widetilde\Ric^{\,\sharp}=\rho^{-1}\Ric^\sharp.
\]
Equation \eqref{eq:ue-scaled-iota} is an isometry because
\[
 \widetilde g(s)(\widetilde\iota_sv,\widetilde\iota_sw)
 =\rho\,g(t)(\rho^{-1/2}\iota_tv,\rho^{-1/2}\iota_tw)
 =h(v,w).
\]
Differentiating in $s$ proves that it solves the corresponding isometry
ODE.  Direct substitution into \eqref{eq:ue-pulled-connection} proves
\eqref{eq:ue-scaled-connection}.

The covariant curvature tensor scales by $\rho$, whereas each of the four
copies of \eqref{eq:ue-scaled-iota} contributes $\rho^{-1/2}$.
Consequently the pulled covariant tensor, and equivalently the
trace-normalized operator, scales by $\rho^{-1}$, proving
\eqref{eq:ue-scaled-operator}.  The chain rule gives
\eqref{eq:ue-scaled-time}.  The inverse base metric in the trace contributes
one factor $\rho^{-1}$ and the section contributes the other, proving
\eqref{eq:ue-scaled-laplacian}.  Finally $Q$ is homogeneous of degree two,
which proves \eqref{eq:ue-scaled-reaction}.
\end{proof}

\begin{proposition}[Nullspace transfer]
\label{prop:ue-nullspace-transfer}
At every $(x,t)$,
\begin{equation}
 F_t\bigl(\ker\mathcal M(x,t)\bigr)
       =\ker\mathcal R_{g(t)}(x).                        \label{eq:ue-kernel-transfer}
\end{equation}
For every piecewise smooth path $\gamma$ in $M$, parallel transports satisfy
\begin{equation}
 F_{t,\gamma(1)}\circ P^{D^t}_\gamma
   =P^{\nabla^{g(t)}}_\gamma\circ F_{t,\gamma(0)}.       \label{eq:ue-parallel-transport-intertwining}
\end{equation}
Hence a subbundle of $\ker\mathcal M(\cdot,t)$ is $D^t$-parallel if and
only if its image under $F_t$ is a Levi--Civita-parallel subbundle of the
geometric curvature nullspace.
\end{proposition}

\begin{proof}
Equation \eqref{eq:ue-kernel-transfer} follows immediately from the
conjugation identity \eqref{eq:ue-fixed-operator}.  Equation
\eqref{eq:ue-iota-parallel}, applied on $\Lambda^2$, says that $F_t$ carries
$D^t$-parallel sections to Levi--Civita-parallel sections.  Uniqueness for
the parallel-transport ODE gives
\eqref{eq:ue-parallel-transport-intertwining}.  Applying it also to the
reversed path proves the stated equivalence.
\end{proof}

\section{Regression tests for signs and factors}
\label{sec:ue-regression}

The following tests are consequences of the proof, but recording them
protects the public interface against future normalization drift.

\begin{proposition}[Scalar trace check]
\label{prop:ue-scalar-trace-check}
If $\lambda,\mu,\nu$ are the eigenvalues of $\mathcal M$, then
\begin{align}
 \trace Q(\mathcal M)
 &=\lambda^2+\mu^2+\nu^2
     +\lambda\mu+\lambda\nu+\mu\nu                    \label{eq:ue-trace-Q}\\
 &=2|\Ric|^2.                                            \label{eq:ue-trace-Q-ricci}
\end{align}
Consequently the trace of \eqref{eq:ue-main-operator-pde} is
\[
                   \partial_tR=\Delta R+2|\Ric|^2.
\]
\end{proposition}

\begin{proof}
The first identity is the trace of
\eqref{eq:ue-diagonal-reaction}.  In a corresponding principal frame, the
Ricci eigenvalues are
\[
        \frac{\mu+\nu}{2},\qquad
        \frac{\lambda+\nu}{2},\qquad
        \frac{\lambda+\mu}{2}.
\]
Squaring and adding proves \eqref{eq:ue-trace-Q-ricci}.  Since
$R=\trace\mathcal M$, tracing the operator PDE gives the standard scalar
evolution.
\end{proof}

\begin{remark}[Canonical model checks]
\label{rem:ue-model-checks}
The exact equation passes the following independent tests.
\begin{enumerate}
\item For a flat metric, $\mathcal M=0$ and $Q(0)=0$.
\item On a round three-manifold of sectional curvature $k(t)$,
      $\mathcal M=2kI$ and $Q(\mathcal M)=8k^2I$.  Ricci flow gives
      $k'=4k^2$, and hence $\mathcal M'=8k^2I$.
\item On $S^2(k)\times\mathbb R$, the eigenvalues are $(2k,0,0)$ and
      $Q(\mathcal M)=(4k^2,0,0)$.  The surface factor satisfies
      $k'=2k^2$, so the nonzero curvature-operator eigenvalue has derivative
      $4k^2$.
\end{enumerate}
The round and product tests simultaneously detect a wrong sign and the most
common factors of two.
\end{remark}

\begin{remark}[Translation to sectional normalization]
\label{rem:ue-sectional-normalization}
If $\mathcal K=\frac12\mathcal M$ has eigenvalues equal to the principal
sectional curvatures, quadratic homogeneity gives
\begin{equation}
       \partial_t\mathcal K
          =\Delta_{D^t,g(t)}\mathcal K
             +2\bigl(\mathcal K^2+\mathcal K^\#\bigr).   \label{eq:ue-sectional-normalized-pde}
\end{equation}
Thus the diagonal sectional-normalized eigenvalue ODE contains an extra
factor $2$.  Formula \eqref{eq:ue-sectional-normalized-pde} is a translation
of the theorem, not an alternative convention which may be substituted into
Chapter~3 without also translating its time scale.
\end{remark}

\section{Formalization interfaces and stable dependency order}
\label{sec:ue-formalization}

The following semantic identifiers name mathematical proof units, not Lean
declarations.  They are stable under chapter renumbering.  Proposed Lean
module and declaration names should be assigned later in the blueprint and
must be back-translated to these exact statements before the interfaces are
frozen.

\begin{enumerate}
\item \nodeid{PC-B-CURVATURE-TENSOR-OPERATOR-NORMALIZATION}: equations
      \eqref{eq:ue-curvature-sign}--\eqref{eq:ue-tensor-operator-bridge},
      including the last-slot swap, the factor $2$, and the scalar-trace
      normalization.
\item \nodeid{PC-F-UHLENBECK-METRIC-FREEZING}: the general metric-family
      Lemma~\ref{lem:ue-metric-freezing}.  Its dependencies are smooth
      linear ODEs with parameters and finite-dimensional linear algebra.
\item \nodeid{PC-F-UHLENBECK-RICCI-GAUGE}: Corollary
      \ref{cor:ue-ricci-gauge-velocity}, specializing the gauge velocity to
      $+\Ric^\sharp$.
\item \nodeid{PC-F-PULLBACK-CONNECTION-NATURALITY}: Definition
      \ref{def:ue-pulled-connection} and Lemma
      \ref{lem:ue-connection-compatible}, including the inverse
      $\iota_t^{-1}$ and the induced connection on self-adjoint
      endomorphisms.
\item \nodeid{PC-F-PULLBACK-ROUGH-LAPLACIAN-NATURALITY}: Lemma
      \ref{lem:ue-spatial-naturality}.  The Hessian statement must precede
      the rough-Laplacian statement so the base correction is not lost.
\item \nodeid{PC-F-RICCI-FLOW-COVARIANT-CURVATURE-EVOLUTION}: Proposition
      \ref{prop:ue-raw-curvature-evolution}, derived through
      Lemmas~\ref{lem:ue-connection-variation}--\ref{lem:ue-B-rewrite}.
      Its statement must retain its sign, four Ricci terms, and Laplacian
      convention.
\item \nodeid{PC-F-TIME-DERIVATIVE-OF-BUNDLE-PULLBACK}: Lemma
      \ref{lem:ue-four-slot-derivative}, an ordinary derivative in a fixed
      tensor bundle.
\item \nodeid{PC-F-UHLENBECK-RICCI-TERM-CANCELLATION}: Proposition
      \ref{prop:ue-pulled-tensor-evolution}; it depends exactly on the raw
      identity, spatial naturality, and the four-slot derivative.
\item \nodeid{PC-B-UHLENBECK-FIXED-CURVATURE-OPERATOR}: Definition
      \ref{def:ue-fixed-curvature-operator} and the proof that conjugation
      agrees with the four-tensor pairing.
\item \nodeid{PC-B-GENERAL-CURVATURE-REACTION-TO-THREE-DIMENSIONAL-ADJUGATE}: Lemma
      \ref{lem:ue-three-dimensional-reaction}, including the exterior-algebra
      basis lemma and the finite contractions
      \eqref{eq:ue-P-lambda}--\eqref{eq:ue-P-nu}.
\item \nodeid{PC-F-UHLENBECK-CURVATURE-EVOLUTION}: Theorem
      \ref{thm:ue-operator-evolution}.  Its direct dependencies are the
      metric and connection nodes, the pulled tensor evolution, and the
      three-dimensional reaction node.
\item \nodeid{PC-B-HAMILTON-IVEY-UHLENBECK-EVOLUTION}: Corollary
      \ref{cor:ue-discharge-hamilton-ivey}, recording exact correspondence
      with Proposition~\ref{prop:hi-uhlenbeck-evolution}.
\item \nodeid{PC-F-UHLENBECK-GAUGE-COVARIANCE}: Proposition
      \ref{prop:ue-gauge-covariance}.
\item \nodeid{PC-F-UHLENBECK-PARABOLIC-SCALING}: Proposition
      \ref{prop:ue-parabolic-scaling}.
\item \nodeid{PC-B-UHLENBECK-NULLSPACE-TRANSFER}: Proposition
      \ref{prop:ue-nullspace-transfer}, the bridge consumed by the strong
      maximum principle and splitting theorem.
\end{enumerate}

The deepest narrow backbone has the following ordered stages:
\begin{enumerate}
\item \nodeid{PC-B-CURVATURE-TENSOR-OPERATOR-NORMALIZATION};
\item \nodeid{PC-F-UHLENBECK-METRIC-FREEZING};
\item \nodeid{PC-F-PULLBACK-CONNECTION-NATURALITY};
\item \nodeid{PC-F-PULLBACK-ROUGH-LAPLACIAN-NATURALITY}.
\end{enumerate}
The derived node
\nodeid{PC-F-RICCI-FLOW-COVARIANT-CURVATURE-EVOLUTION} joins that chain at
\nodeid{PC-F-UHLENBECK-RICCI-TERM-CANCELLATION}.  The proof then passes
through
\nodeid{PC-B-GENERAL-CURVATURE-REACTION-TO-THREE-DIMENSIONAL-ADJUGATE} and
ends at \nodeid{PC-F-UHLENBECK-CURVATURE-EVOLUTION}.  Gauge covariance,
scaling, and nullspace transfer are downstream siblings and should not be
made hidden fields of the main theorem.

\section{Source ledger and convention repairs}
\label{sec:ue-sources}

The primary source is R.~S.~Hamilton,
\emph{Four-manifolds with positive curvature operator}, Journal of
Differential Geometry \textbf{24} (1986), 153--179,
\href{https://doi.org/10.4310/jdg/1214440433}{pp.~155--157}.  Those pages
introduce the abstract fixed metric bundle, evolve the isometry, pull back
the connection and curvature tensor, prove the Ricci-term cancellation, and
rewrite the remaining reaction as $M^2+M^\#$.  Hamilton's covariant
curvature-component sign differs from
\eqref{eq:ue-covariant-rm}; therefore his displayed sign in the auxiliary
$B$-tensor must be translated rather than copied.

The project sources were used as follows.
\begin{enumerate}
\item \emph{The Ricci Flow: An Introduction}, Mathematical Surveys and
      Monographs~110: equations (6.15)--(6.17), pp.~179--180, give the raw
      tensor equation; pp.~180--183 give the Uhlenbeck construction and
      cancellation; Theorem~6.26, pp.~183--186, gives the operator equation;
      equation (6.32), p.~188, gives the three-dimensional reaction.
\item \emph{Hamilton's Ricci Flow}, Graduate Studies in Mathematics~77:
      Lemma~2.58 and equations (2.70)--(2.74), pp.~121--123, give a concise
      independent cross-check of the four-slot cancellation and the
      three-dimensional coefficients.
\item \emph{The Ricci Flow: Techniques and Applications, Part I:
      Geometric Aspects}, Mathematical Surveys and Monographs~135:
      Lemma~A.23 and
      Theorem~A.24, pp.~457--459, give the corrected abstract construction.
      In particular, p.~458 explicitly includes $\iota^{-1}$ in the pulled
      connection and the minus sign in the pulled $B$-tensor.
\item \emph{The Ricci Flow: Techniques and Applications, Part II},
      Mathematical Surveys and Monographs~144: equations (10.29)--(10.32),
      pp.~14--15, record the exact PDE and ODE as inputs to the systems
      maximum principle.
\end{enumerate}

Two source hazards are repaired rather than propagated.
\begin{enumerate}
\item In the inspected copy of Mathematical Surveys and Monographs~110,
      the displayed pulled $B_{abcd}$ in Lemma~6.22 on p.~182 omits the
      minus sign.  This conflicts with its own equation (6.15) on p.~179
      and with the proof of Theorem~6.26.  Equation (2.73) in Graduate
      Studies in Mathematics~77 and the corrected Appendix~A formula in
      Mathematical Surveys and Monographs~135 both give
      \[
          \overline B_{abcd}
             =-h^{eg}h^{fh}
                 \overline R_{aebf}\overline R_{cgdh},
      \]
      which is the formula used here.
\item Any shorthand for the pulled connection which omits
      $\iota_t^{-1}$ is not type-correct: $\nabla_X(\iota_t\xi)$ lies in
      $TM$, not in $V$.  Definition~\ref{def:ue-pulled-connection} uses the
      corrected formula stated explicitly on p.~458 of Mathematical Surveys
      and Monographs~135.
\end{enumerate}

The result proved in this chapter is the complete raw tensor,
Uhlenbeck-gauge, and three-dimensional operator derivation.  The raw
identity \eqref{eq:ue-raw-rm-evolution} is derived above from the connection
variation, Bianchi, commutator, and $B$-tensor lemmas.  No maximum principle, curvature positivity,
orientation, compactness, or global trivialization is used.  The chapter is
the human mathematical source for the corresponding blueprint nodes; it
makes no claim that their Lean implementations have already been written,
compiled, or audited.

\chapter{The Strong Maximum Principle and Three-Dimensional Splitting}
\label{ch:strong-curvature}
\index{maximum principle!strong, for systems}
\index{curvature operator!rank rigidity}
\index{splitting theorem!parallel line}

\section{Purpose and main outputs}
\label{sec:strong-curvature-purpose}

Hamilton's weak maximum principle preserves a convex curvature cone.  The
strong principle asks what equality with the boundary of that cone forces
after positive time.  For a nonnegative three-dimensional curvature
operator the answer is rigid: its rank is spatially constant, rank two is
impossible, and rank one produces a parallel tangent line.  At a complete
time slice that line splits the universal cover as a positively curved
surface times $\R$.

This chapter separates the proof into interfaces suitable for later Lean
formalization.  The abstract parabolic theorem is proved with lower Ky Fan
spectral sums, so no eigenvector field is differentiated through an
eigenvalue multiplicity.  Spatial parallelism of the kernel is then proved
by differentiating the identity $Aw=0$.  Ordinary time-independence of the
kernel is stated only after a separate kernel-annihilation hypothesis has
been verified.  This distinction repairs a false general formulation while
retaining the curvature theorem: the three-dimensional reaction
$Q(A)=A^2+A^\#$ is nonnegative and commutes with every nonnegative $A$.

The local strong theorem assumes that the solution is already nonnegative;
it uses neither completeness nor a global curvature bound.  Completeness
enters only in the global parallel-line splitting.  A single product for an
entire noncompact flow additionally uses complete bounded-curvature
Ricci-flow uniqueness, and that input is displayed rather than hidden.

\begin{roadmap}{Section~\ref{sec:smp-abstract-systems} develops the
strong maximum principle and its exact kernel-motion law.}
Sections~\ref{sec:curvature-three-dimensional-reaction}--
\ref{sec:curvature-rank-one-line} specialize it to the three-dimensional
curvature reaction and construct the parallel tangent line without an
orientation.  Sections~\ref{sec:parallel-line-global-splitting}--
\ref{sec:curvature-fixed-time-splitting} prove the complete fixed-time
trichotomy.  The final sections isolate the extra uniqueness input for a
whole-flow product, record the ancient-flow consequence, and give the
formalization and source ledgers.
\end{roadmap}

% Abstract strong maximum principle engine.

\section{The strong maximum principle for positive systems}
\label{sec:smp-abstract-systems}

Hamilton's strong maximum principle has two logically different parts.
The first is genuinely parabolic: positivity at one point spreads through a
connected manifold and forces the rank of a nonnegative system to be
spatially constant at positive time.  The second is geometric: once the
rank is constant, the nullspace is preserved by spatial parallel transport.
There is then a third assertion which must be kept separate from both:
ordinary constancy of the nullspace in time, viewed inside a fixed vector
bundle, requires the reaction to annihilate the nullspace.  The usual
null-eigenvector inequality by itself does not imply that annihilation.

This section states and proves the precise abstract theorem.  Its bundle
conventions agree exactly with Chapter~\ref{ch:weak-maximum-principle-systems}.
The base metric and the compatible bundle connection may depend on time,
whereas the vector bundle and its fiber metric are fixed.  The proof is
local in space once nonnegativity of the solution is known: neither
completeness nor a global curvature bound is needed.

The macro formalization milestone for this section is
\nodeid{PC-M-STRONG-MAXIMUM-PRINCIPLE-SYSTEMS}.

\subsection{Geometric and analytic setup}
\label{subsec:smp-setup}

Let $M$ be a nonempty connected smooth manifold without boundary.  It is
not assumed compact or complete.  Fix $T>0$, a smooth family of Riemannian metrics
$g(t)$ on $M$ for $0\leq t\leq T$, and a smooth time-dependent vector field
$X(t)$ on $M$.

Let $\pi:W\to M$ be a smooth real vector bundle of finite rank $r\geq1$.
Fix a positive-definite fiber metric $h$ on $W$, independent of time, and
let $D^t$ be a smooth family of $h$-compatible connections.  The connection
$D^t$ induces a metric connection, denoted by the same symbol, on
\[
              \mathcal S=\operatorname{SymEnd}_h(W),
\]
the bundle of $h$-self-adjoint endomorphisms of $W$.  We use the rough
Laplacian
\[
 \Delta_t A
 =\trace_{g(t)}(D^t)^2A
 =\sum_{i=1}^{n}
   \left(D^t_{e_i}D^t_{e_i}A
        -D^t_{\nabla^{g(t)}_{e_i}e_i}A\right)
\]
in a $g(t)$-orthonormal frame.  Its sign is such that $\Delta f\geq0$ at a
local minimum of a smooth scalar function.

Let
\[
 \beta:\mathcal S\times[0,T]\longrightarrow\mathcal S
\]
be a continuous fiber-preserving map which is locally Lipschitz in its
endomorphism argument, locally uniformly in space and time.  Explicitly,
for every compact $C\subset M$, compact interval $J\subset[0,T]$, and
$R<\infty$, there is $L<\infty$ such that
\begin{equation}
 \|\beta_x(B,t)-\beta_x(C',t)\|_{\mathrm{op}}
 \leq L\|B-C'\|_{\mathrm{op}}                         \label{eq:smp-local-lipschitz}
\end{equation}
whenever $x\in C$, $t\in J$, and $B,C'\in\mathcal S_x$ have operator norm
at most $R$.  Smooth reactions automatically have this property.

\paragraph{Why local Lipschitzness is not cosmetic.}
On a one-point base, define the scalar reaction on all of $\mathbb R$ by
\[
 \beta(a)=-2\operatorname{sgn}(a)\sqrt{|a|},
 \qquad \beta(0)=0.
\]
Its nonnegative branch is $\beta(a)=-2\sqrt a$; it satisfies the null
condition at zero, but the stopped solution of
\[
             a'=-2\sqrt a,\qquad a(0)=1,
\]
is positive for $0\leq t<1$ and zero for $t\geq1$.  Thus rank can be lost
in finite forward time when uniqueness at the cone boundary fails.  This
example is $C^1$ in time, the natural time regularity for a classical
spatially constant solution.  If a smooth-in-time example is desired, take
\[
 \beta(a)=
 \begin{cases}
   -a(\log a)^2,&a>0,\\
   0,&a\leq0,
 \end{cases}
 \qquad
 a(t)=
 \begin{cases}
   \exp\!\bigl(-1/(1-t)\bigr),&t<1,\\
   0,&t\geq1.
 \end{cases}
\]
Again $\beta$ is continuous, satisfies the null condition, and is not
locally Lipschitz at zero, while $a$ is smooth and loses rank at $t=1$.
Only the displayed nonnegative solution branch is used in either example.
This necessity test has stable node
\nodeid{PC-B-REACTION-LOCAL-LIPSCHITZ-NECESSITY}.

We impose the following tangent condition on the cone of nonnegative
endomorphisms.

\begin{definition}[Null-eigenvector condition]
\label{def:smp-null-eigenvector}
The reaction $\beta$ satisfies the \emph{null-eigenvector condition} if,
for every $x\in M$, $t\in[0,T]$, $B\in\mathcal S_x$, and $v\in W_x$,
\begin{equation}
 B\geq0,\quad Bv=0
 \quad\Longrightarrow\quad
 h\bigl(\beta_x(B,t)v,v\bigr)\geq0.                     \label{eq:smp-null-eigenvector}
\end{equation}
Equivalently, if
$\beta(B,t)(v,w):=h(\beta(B,t)v,w)$ denotes the associated symmetric
bilinear form, then \eqref{eq:smp-null-eigenvector} says
$\beta(B,t)(v,v)\geq0$ on every null vector of $B$.
\end{definition}

\textcolor{red}{Since the above is a major definition, put the Lean code for the statement here.}

Let $A$ be a smooth section of the pullback of $\mathcal S$ to
$M\times[0,T]$ satisfying
\begin{equation}
 \partial_tA
   =\Delta_tA+D^t_{X(t)}A+\beta(A,t).                    \label{eq:smp-system}
\end{equation}
Because $W$ is fixed, $\partial_tA(x,t)$ is the ordinary derivative of a
curve in the single vector space $\mathcal S_x$.  No time connection is
being introduced.  Throughout this section we assume
\begin{equation}
                  A(x,t)\geq0
       \qquad\text{for every }(x,t)\in M\times[0,T].      \label{eq:smp-solution-nonnegative}
\end{equation}
In the intended applications, this hypothesis is first supplied by the weak
maximum principle of Chapter~\ref{ch:weak-maximum-principle-systems}.  The
strong theorem below does not silently extend that weak maximum principle
to a noncompact base.

\begin{theorem}[Strong maximum principle for systems]
\label{thm:smp-systems}
Assume the setup above, including
\eqref{eq:smp-null-eigenvector} and
\eqref{eq:smp-solution-nonnegative}.
Then the following statements hold.

\begin{enumerate}
\item \emph{Point-to-point rank spreading.}
      If $0\leq s<t\leq T$ and $x,y\in M$, then
      \begin{equation}
          \operatorname{rank}A(y,t)
          \geq\operatorname{rank}A(x,s).                 \label{eq:smp-point-rank-spreading}
      \end{equation}
      Equivalently,
      \begin{equation}
       \inf_{y\in M}\operatorname{rank}A(y,t)
       \geq
       \sup_{x\in M}\operatorname{rank}A(x,s).          \label{eq:smp-global-rank-spreading}
      \end{equation}
      The infimum and supremum in
      \eqref{eq:smp-global-rank-spreading} are integer values; compactness
      is not needed for their attainment.

\item \emph{Positive-time spatial constancy.}
      For every $t\in(0,T]$, the integer
      $\operatorname{rank}A(x,t)$ is independent of
      $x$.  Denote its common value by $\rho(t)$.  The function $\rho$ is
      nondecreasing.  Every $t_0\in(0,T]$ has a left neighborhood
      $(t_0-\varepsilon,t_0]$ on which $\rho$ is constant.  There is also
      $\delta>0$ for which $\rho$ is constant on $(0,\delta]$.  Consequently
      every compact positive-time interval admits a finite partition into
      right-closed intervals on which the rank is constant in space and
      time.

\item \emph{Smoothness and spatial parallelism of the kernel.}
      On any open time interval $J\subset(0,T]$ on which $\rho$ is
      constant, the sets
      \[
            K_{(x,t)}=\ker A(x,t)
      \]
      form a smooth vector subbundle
      $K\subset\operatorname{pr}_M^*W|_{M\times J}$.
      For every $t\in J$, this subbundle is preserved by $D^t$-parallel
      transport in space.  Infinitesimally, if $w$ is a smooth local
      section of $K$, then
      \begin{equation}
                  D^t_Yw\in K_{(x,t)}                   \label{eq:smp-kernel-spatial-parallel}
      \end{equation}
      for every tangent vector $Y\in T_xM$.

\item \emph{The exact kernel-motion law.}
      On a constant-rank interval $J$, every smooth local section $w$ of
      $K$ satisfies
      \begin{align}
       h\bigl(\beta(A,t)w,w\bigr)&=0,                    \label{eq:smp-reaction-null-quadratic}\\
       A(\partial_tw)&=-\beta(A,t)w.                     \label{eq:smp-kernel-motion}
      \end{align}
      Equation \eqref{eq:smp-kernel-motion}, rather than ordinary time
      constancy, is the conclusion furnished by the null-eigenvector
      condition alone.

\item \emph{Temporal constancy under kernel annihilation.}
      If, in addition,
      \begin{equation}
          w\in\ker A(x,t)
          \quad\Longrightarrow\quad
          \beta(A(x,t),t)w=0,                            \label{eq:smp-kernel-annihilation}
      \end{equation}
      then $K(t)$ is independent of $t\in J$ as an actual subbundle of the
      fixed bundle $W$.  The same is true of
      $K(t)^{\perp_h}=\operatorname{im}A(t)$.
\end{enumerate}
\end{theorem}

The permanent identity of Theorem~\ref{thm:smp-systems} is
\nodeid{PC-M-STRONG-MAXIMUM-PRINCIPLE-SYSTEMS}.  The individual conclusions
are refined below into fine nodes so that later formalization does not treat
the theorem as one indivisible black box.

\subsection{Ky Fan sums and their rank content}
\label{subsec:smp-ky-fan}

Let $E$ be an $r$-dimensional real inner-product space.  For
$B\in\operatorname{SymEnd}(E)$, write its eigenvalues in increasing order:
\[
                 \lambda_1(B)\leq\cdots\leq\lambda_r(B).
\]

\begin{definition}[Lower Ky Fan sum]
\label{def:smp-ky-fan}
For $1\leq k\leq r$, define
\begin{equation}
                 \Phi_k(B)=\sum_{i=1}^k\lambda_i(B).      \label{eq:smp-ky-fan-definition}
\end{equation}
\end{definition}

\begin{lemma}[Variational, continuity, and rank formulas]
\label{lem:smp-ky-fan}
For every self-adjoint $B$ and every $1\leq k\leq r$,
\begin{align}
 \Phi_k(B)
   &=\min_{\substack{e_1,\ldots,e_k\in E\\
                     \langle e_i,e_j\rangle=\delta_{ij}}}
      \sum_{i=1}^k\langle Be_i,e_i\rangle,               \label{eq:smp-ky-fan-frame}\\
   &=\min_{P\in\mathcal P_k(E)}\trace(PB),               \label{eq:smp-ky-fan-projection}\\
 |\Phi_k(B)-\Phi_k(C)|
   &\leq k\|B-C\|_{\mathrm{op}},                         \label{eq:smp-ky-fan-lipschitz}\\
 \Phi_k(B+cI)&=\Phi_k(B)+kc.                             \label{eq:smp-ky-fan-shift}
\end{align}
Here $\mathcal P_k(E)$ is the set of rank-$k$ orthogonal projections.
If $B\geq0$, then
\begin{equation}
 \Phi_k(B)>0
 \quad\Longleftrightarrow\quad
 \dim\ker B<k
 \quad\Longleftrightarrow\quad
 \operatorname{rank}B\geq r-k+1.                        \label{eq:smp-ky-fan-rank}
\end{equation}
Moreover, if $P$ realizes the minimum in
\eqref{eq:smp-ky-fan-projection}, then
\begin{equation}
                 \trace(PC)\geq\Phi_k(C)                 \label{eq:smp-ky-fan-support}
\end{equation}
for every self-adjoint $C$.
\end{lemma}

\begin{proof}
Choose an orthonormal eigenbasis $u_1,\ldots,u_r$ for $B$.  If
$e_1,\ldots,e_k$ is any orthonormal $k$-frame and
$a_j=\sum_{i=1}^k\langle e_i,u_j\rangle^2$, then
\[
       0\leq a_j\leq1,
       \qquad
       \sum_{j=1}^r a_j=k,
\]
and
\[
 \sum_{i=1}^k\langle Be_i,e_i\rangle
 =\sum_{j=1}^r\lambda_j(B)a_j
 \geq\sum_{j=1}^k\lambda_j(B).
\]
The last inequality is the elementary mass-transfer lemma: among weights
$0\leq a_j\leq1$ of total mass $k$, an increasing sequence has least
weighted sum when the first $k$ weights equal one and the rest equal zero.
Equality is obtained from $e_i=u_i$, proving
\eqref{eq:smp-ky-fan-frame}.  A rank-$k$ orthogonal projection is the
projection onto the span of an orthonormal $k$-frame, so
\eqref{eq:smp-ky-fan-projection} is equivalent.

If $P$ minimizes for $B$, then
\[
 \Phi_k(C)-\Phi_k(B)
 \leq\trace(P(C-B))
 \leq k\|C-B\|_{\mathrm{op}}.
\]
Interchanging $B$ and $C$ proves
\eqref{eq:smp-ky-fan-lipschitz}.  Formula
\eqref{eq:smp-ky-fan-shift} follows because scalar translation shifts every
eigenvalue by $c$.  For $B\geq0$, the first eigenvalues are zero precisely
on the kernel, which proves \eqref{eq:smp-ky-fan-rank}.  Finally,
\eqref{eq:smp-ky-fan-support} is the defining minimum property of
$\Phi_k(C)$.
\end{proof}

The stable node for Definition~\ref{def:smp-ky-fan} and
Lemma~\ref{lem:smp-ky-fan} is
\nodeid{PC-F-KY-FAN-LOWER-SPECTRAL-SUM}.  In particular, later arguments
will use a minimizing projection, not a differentiable choice of ordered
eigenvectors.

\subsection{The scalar local input}
\label{subsec:smp-scalar-input}

The sole scalar strong-maximum-principle input is the following classical
Dirichlet statement.  We state its scope precisely because it is what makes
the systems theorem local on a noncompact manifold.

\begin{proposition}[Local scalar Dirichlet strong maximum principle]
\label{prop:smp-local-scalar-input}
Let $\Omega\Subset M$ be a connected open set with smooth boundary, let
$0\leq s<t_1\leq T$, and let $a\geq0$.  If
$f_0\in C_c^\infty(\Omega)$ is nonnegative and not identically zero, then
the Dirichlet problem
\begin{align}
 \partial_tf&=\Delta_{g(t)}f+X(t)f-af
       &&\text{in }\Omega\times(s,t_1],                  \label{eq:smp-scalar-dirichlet-pde}\\
 f&=0
       &&\text{on }\partial\Omega\times[s,t_1],          \label{eq:smp-scalar-dirichlet-boundary}\\
 f(\cdot,s)&=f_0
       &&\text{on }\Omega                                \label{eq:smp-scalar-dirichlet-initial}
\end{align}
has a classical solution, continuous up to the parabolic boundary and
smooth in the interior for positive elapsed time.  It satisfies
\begin{equation}
            f\geq0\quad\text{on }\overline\Omega\times[s,t_1],
 \qquad
            f>0\quad\text{on }\Omega\times(s,t_1].       \label{eq:smp-scalar-dirichlet-positive}
\end{equation}
\end{proposition}

\textcolor{red}{Since the above is a major proposition, put the Lean code for the statement here.}

The proposition combines standard linear Dirichlet existence, the scalar
weak maximum principle, and the scalar strong maximum principle.  No
completeness assumption appears because $\overline\Omega$ is compact.
When $M$ is closed one may instead take $\Omega=M$ and omit the boundary
condition.  An exact project-source locator is MSM~144,
Theorem~12.40 and Corollary~12.42, pp.~180--182, for the local scalar
strong principle, together with the linear initial-boundary-value problem
used in the proof of Proposition~12.47, pp.~184--185.  That proof cites
Evans's parabolic theory, Section~7.1, for Dirichlet existence and
Theorem~10.2 of MSM~144 for the scalar weak principle.  The stable node is
\nodeid{PC-F-LOCAL-SCALAR-STRONG-PARABOLIC-MAXIMUM-PRINCIPLE}.

We also use the elementary fact that any two points of a connected
boundaryless smooth manifold lie in some connected relatively compact
smooth domain.  Indeed, join the points by a smooth path, take a relatively
compact neighborhood of its compact image, choose a regular sublevel set
of a smooth local exhaustion which still contains the path, and take the
connected component containing that path.

\subsection{Strict propagation of lower spectral sums}
\label{subsec:smp-rank-spreading}

The next proposition is the parabolic engine.  Its proof is arranged so
that repeated eigenvalues cause no regularity problem.

\begin{proposition}[Rank spreading for positive systems]
\label{prop:smp-rank-spreading}
Under the hypotheses of Theorem~\ref{thm:smp-systems}, fix
$1\leq k\leq r$.  If
\begin{equation}
                   \Phi_k(A(x_0,s))>0                    \label{eq:smp-positive-phi-start}
\end{equation}
for some $x_0\in M$ and $0\leq s<T$, then
\begin{equation}
                   \Phi_k(A(y,t))>0                      \label{eq:smp-positive-phi-later}
\end{equation}
for every $y\in M$ and every $t\in(s,T]$.
Consequently, \eqref{eq:smp-point-rank-spreading} and
\eqref{eq:smp-global-rank-spreading} hold.
\end{proposition}

\begin{proof}
Fix $y\in M$ and $t_1\in(s,T]$.  Choose a connected smooth domain
$\Omega\Subset M$ containing both $x_0$ and $y$.  By continuity of
$x\mapsto\Phi_k(A(x,s))$ and
\eqref{eq:smp-positive-phi-start}, there is a nonnegative
$f_0\in C_c^\infty(\Omega)$ such that
\begin{equation}
 f_0(x_0)>0,
 \qquad
 kf_0(x)\leq\Phi_k(A(x,s))\quad(x\in\Omega).             \label{eq:smp-bump-dominated}
\end{equation}

Set
\[
              C=\max_{\overline\Omega\times[s,t_1]}
                    \|A\|_{\mathrm{op}}.
\]
Whenever $e$ is a unit eigenvector of a nonnegative $A(x,t)$ with
eigenvalue $\lambda$, the endomorphism
$A-\lambda e\otimes e$ is nonnegative and has operator norm at most $C$.
The local Lipschitz hypothesis therefore supplies a single constant $L$
valid for all pairs of this form over
$\overline\Omega\times[s,t_1]$.  Choose a constant $c>L$, and let $f$ be
the solution of Proposition~\ref{prop:smp-local-scalar-input} with
$a=c$ and initial value $f_0$.  Thus
\begin{equation}
 \mathcal L f=-cf,
 \qquad
 \mathcal L=\partial_t-\Delta_t-D^t_{X(t)}               \label{eq:smp-parabolic-operator}
\end{equation}
on scalar functions, where $D^t_{X(t)}f=X(t)f$, and
$f>0$ in $\Omega\times(s,t_1]$.

For $\varepsilon>0$, define
\begin{equation}
 H_\varepsilon
   =A+\bigl(\varepsilon e^{c(t-s)}-f\bigr)I              \label{eq:smp-barrier-tensor}
\end{equation}
and its lower spectral barrier
\begin{equation}
 \Theta_\varepsilon
   :=\Phi_k(H_\varepsilon)
    =\Phi_k(A)+k\bigl(\varepsilon e^{c(t-s)}-f\bigr).     \label{eq:smp-theta-barrier}
\end{equation}
The identity endomorphism is parallel for every $D^t$ and is independent
of time because $h$ is fixed.  Equations
\eqref{eq:smp-system} and \eqref{eq:smp-parabolic-operator} give
\begin{equation}
 \mathcal L H_\varepsilon
 =\beta(A,t)
   +c\bigl(\varepsilon e^{c(t-s)}+f\bigr)I.              \label{eq:smp-barrier-evolution}
\end{equation}
At the initial time, the shift formula and
\eqref{eq:smp-bump-dominated} imply
\[
 \Theta_\varepsilon(x,s)
 =\Phi_k(A(x,s))+k(\varepsilon-f_0(x))
 \geq k\varepsilon>0.
\]
On the side boundary, $f=0$ and $A\geq0$, so
\[
 \Theta_\varepsilon
 =\Phi_k(A)+k\varepsilon e^{c(t-s)}>0.
\]

We claim that
\begin{equation}
 \Theta_\varepsilon>0
 \quad\text{on }\overline\Omega\times[s,t_1].           \label{eq:smp-barrier-strict}
\end{equation}
Suppose otherwise.  Continuity, compactness of
$\overline\Omega$, and strict positivity on the parabolic boundary give a
first time $\tau\in(s,t_1]$ and an interior point $z\in\Omega$ such that
\begin{equation}
 \Theta_\varepsilon(z,\tau)=0,
 \qquad
 \Theta_\varepsilon>0
 \quad\text{on }\overline\Omega\times[s,\tau).          \label{eq:smp-first-contact}
\end{equation}

Choose an $h$-orthonormal eigenbasis $e_1,\ldots,e_r$ of
$H_\varepsilon(z,\tau)$, ordered by increasing eigenvalue, and let $P$ be
the orthogonal projection onto
$\operatorname{span}\{e_1,\ldots,e_k\}$.  If the $k$-th eigenvalue has
multiplicity, any $k$-dimensional minimizing subspace inside the relevant
eigenspace may be used.  This is the only choice made at the contact point;
no eigenvector is differentiated.  Since $H_\varepsilon$ differs from $A$
by a scalar multiple of $I$, the same vectors diagonalize $A(z,\tau)$ and
the chosen subspace also realizes $\Phi_k(A(z,\tau))$.

For completeness, we spell out the scalar first-contact calculus.  At the
frozen time $\tau$, extend each $e_i$ by $D^\tau$-parallel transport along
$g(\tau)$-geodesics emanating from $z$.  At $z$ this gives
\begin{equation}
 D^\tau_Ye_i=0,
 \qquad
 ((D^\tau)^2e_i)(Y,Y)=0                                 \label{eq:smp-frame-spatial-jets}
\end{equation}
for every $Y\in T_zM$.  Only the same-direction second derivative is being
asserted; bundle curvature may obstruct vanishing of mixed second
derivatives.  Extend the frame in time, preserving $h$-orthonormality, so
that
\begin{equation}
                    \partial_te_i(z,\tau)=0.             \label{eq:smp-frame-time-jet}
\end{equation}
Such a prescribed time jet is possible precisely because $W$ and $h$ are
fixed.

Define near $(z,\tau)$
\[
        \psi(x,t)=\sum_{i=1}^k
             h\bigl(H_\varepsilon(x,t)e_i(x,t),e_i(x,t)\bigr).
\]
The Ky Fan variational formula gives
$\psi\geq\Theta_\varepsilon$.  Equality holds at $(z,\tau)$.  Hence
\eqref{eq:smp-first-contact} implies
\begin{equation}
 \partial_t\psi(z,\tau)\leq0,
 \qquad
 d\psi(z,\tau)=0,
 \qquad
 \Delta_{g(\tau)}\psi(z,\tau)\geq0.                     \label{eq:smp-scalar-contact-jets}
\end{equation}
Equations \eqref{eq:smp-frame-spatial-jets} and
\eqref{eq:smp-frame-time-jet} remove all derivatives of the chosen frame
from the traced calculation at the contact point.  More invariantly, if
one uses only an orthonormal extension with $D^\tau e_i=0$ at $z$, its
possible Laplacian residual is
$2h(H_\varepsilon\Delta^\tau e_i,e_i)$; this vanishes because
$H_\varepsilon e_i$ is a scalar multiple of $e_i$ and
$h(\Delta^\tau e_i,e_i)=-|D^\tau e_i|^2=0$ at $z$.  Thus
\begin{equation}
 0\geq\mathcal L\psi(z,\tau)
   =\sum_{i=1}^k
       h\bigl((\mathcal LH_\varepsilon)e_i,e_i\bigr).     \label{eq:smp-contact-trace}
\end{equation}

Write
\[
                 A(z,\tau)e_i=\lambda_i e_i
       \qquad(1\leq i\leq k).
\]
Every $\lambda_i$ is nonnegative.  Set
\begin{equation}
                A_i=A(z,\tau)-\lambda_i e_i\otimes e_i.  \label{eq:smp-eigenvalue-deletion}
\end{equation}
Here $(e_i\otimes e_i)v=h(v,e_i)e_i$ is the rank-one orthogonal
projection onto $\mathbb Re_i$.
This definition is valid without any simplicity assumption on
$\lambda_i$: in the chosen orthonormal eigenbasis, it merely replaces one
diagonal entry by zero.  Hence $A_i\geq0$, $A_ie_i=0$, and
$\|A-A_i\|_{\mathrm{op}}=\lambda_i$.  The null-eigenvector condition and
the choice of $L$ imply
\begin{align}
 h\bigl(\beta(A,\tau)e_i,e_i\bigr)
 &\geq h\bigl(\beta(A_i,\tau)e_i,e_i\bigr)
      -\|\beta(A,\tau)-\beta(A_i,\tau)\|_{\mathrm{op}}\notag\\
 &\geq-L\lambda_i.                                      \label{eq:smp-beta-eigen-bound}
\end{align}
Summing and using that the chosen subspace minimizes for $A$ gives
\begin{equation}
 \sum_{i=1}^k h\bigl(\beta(A,\tau)e_i,e_i\bigr)
 \geq-L\Phi_k(A(z,\tau)).                                \label{eq:smp-beta-trace-bound}
\end{equation}

At the contact point, the scalar-shift identity and
\eqref{eq:smp-first-contact} say
\begin{equation}
 \Phi_k(A(z,\tau))
   =k\bigl(f(z,\tau)-\varepsilon e^{c(\tau-s)}\bigr).     \label{eq:smp-contact-shift-identity}
\end{equation}
Insert \eqref{eq:smp-barrier-evolution},
\eqref{eq:smp-beta-trace-bound}, and
\eqref{eq:smp-contact-shift-identity} into
\eqref{eq:smp-contact-trace}.  We obtain
\begin{align*}
 0
 &\geq-L\Phi_k(A)
       +kc\varepsilon e^{c(\tau-s)}+kcf\\
 &=k\left[(c-L)f+(c+L)\varepsilon e^{c(\tau-s)}\right]>0,
\end{align*}
a contradiction.  This proves \eqref{eq:smp-barrier-strict}.

Evaluate \eqref{eq:smp-barrier-strict} at $(y,t_1)$ and use the scalar
shift formula:
\[
 \Phi_k(A(y,t_1))
 >k f(y,t_1)-k\varepsilon e^{c(t_1-s)}.
\]
The constants $c$, $L$, and $f$ are independent of $\varepsilon$.  Let
$\varepsilon\downarrow0$.  Since $f(y,t_1)>0$,
\[
                 \Phi_k(A(y,t_1))\geq kf(y,t_1)>0.
\]
The point $y$ and time $t_1$ were arbitrary, proving
\eqref{eq:smp-positive-phi-later}.

Finally, suppose $\operatorname{rank}A(x,s)=\ell$.  If $\ell=0$, rank spreading is
automatic.  If $\ell\geq1$, set $k=r-\ell+1$.  Since $A(x,s)\geq0$, it has
exactly $r-\ell$ zero eigenvalues, so
$\Phi_k(A(x,s))>0$.  The result just proved and
\eqref{eq:smp-ky-fan-rank} give
$\operatorname{rank}A(y,t)\geq\ell$.  This proves
\eqref{eq:smp-point-rank-spreading}.  Taking integer-valued infima and
suprema gives \eqref{eq:smp-global-rank-spreading}.
\end{proof}

The stable node for Proposition~\ref{prop:smp-rank-spreading} is
\nodeid{PC-F-POSITIVE-SYSTEM-RANK-SPREADING}.  Its spectral propagation
sublemma has stable node
\nodeid{PC-F-KY-FAN-STRICT-POSITIVITY-PROPAGATION}.

\subsection{Positive-time rank and the smooth null bundle}
\label{subsec:smp-constant-rank}

\begin{proposition}[Spatial rank constancy at positive time]
\label{prop:smp-positive-time-rank}
For every $t_0\in(0,T]$, the rank of $A(\cdot,t_0)$ is constant on $M$.
Its common value $\rho(t_0)$ is nondecreasing in $t_0$, locally constant
from the left, and constant on some interval $(0,\delta]$.
\end{proposition}

\textcolor{red}{Since the above is a major proposition, put the Lean code for the statement here.}

\begin{proof}
Suppose that at one time $t_0>0$ there are points $x,y\in M$ with
\[
             \operatorname{rank}A(x,t_0)
             <\operatorname{rank}A(y,t_0)=\ell.
\]
If $\ell\geq1$, the smallest positive eigenvalue of $A(y,t_0)$ remains
positive at $(y,s)$ for all $s<t_0$ sufficiently close to $t_0$.
Equivalently, by continuity of the relevant Ky Fan sum,
$\operatorname{rank}A(y,s)\geq\ell$.
Proposition~\ref{prop:smp-rank-spreading}, applied
from $(y,s)$ to $(x,t_0)$, yields
$\operatorname{rank}A(x,t_0)\geq\ell$, a contradiction.  Thus the rank is
spatially constant.
Rank spreading immediately implies that the common integer-valued function
$\rho(t)$ is nondecreasing.  Fix $t_0>0$.  If $\rho(t_0)>0$, its smallest
positive eigenvalue at any fixed point persists for times $s<t_0$ close to
$t_0$, so $\rho(s)\geq\rho(t_0)$.  Monotonicity gives the reverse
inequality, hence equality.  If $\rho(t_0)=0$, monotonicity itself gives
$\rho(s)=0$ for every $s<t_0$.  Therefore $\rho$ is locally constant from
the left.

Finally, the image of $\rho:(0,T]\to\{0,1,\ldots,r\}$ is a nonempty finite
set.  Let $q$ be its least value and choose $t_*>0$ with
$\rho(t_*)=q$.  Monotonicity and minimality give
$\rho(t)=q$ for every $0<t\leq t_*$.  Taking $\delta=t_*$ proves the last
claim.  Since a nondecreasing function with values in
$\{0,1,\ldots,r\}$ has at most $r$ strict jumps, the asserted finite
rank stratification follows as well.
\end{proof}

The stable node is
\nodeid{PC-F-POSITIVE-TIME-SPATIAL-RANK-CONSTANCY}.

\begin{lemma}[Smooth kernel for a constant-rank self-adjoint family]
\label{lem:smp-smooth-kernel}
Let $J\subset(0,T]$ be an open interval on which $\rho$ is constant.  Then
\[
 K=\{(w,t)\in W\times J:A(\pi(w),t)w=0\}
\]
is a smooth vector subbundle of
$\operatorname{pr}_M^*W\to M\times J$.  Its orthogonal complement is
$\operatorname{im}A$.
\end{lemma}

\textcolor{red}{Since the above is a major lemma, put the Lean code for the statement here.}

\begin{proof}
This is the constant-rank kernel theorem, but we record its local mechanism.
Fix $(x_0,t_0)$ and split
\[
        W_{x_0}=K_0\oplus P_0,
        \qquad K_0=\ker A(x_0,t_0),
        \qquad P_0=K_0^{\perp_h}.
\]
Extend this to a smooth local orthogonal splitting.  In block form write
\[
 A=\begin{pmatrix}B&C^*\\ C&D\end{pmatrix}.
\]
At $(x_0,t_0)$ the block $D$ is positive definite, hence it remains
invertible nearby.  Elementary block elimination gives
\[
 \operatorname{rank}A
 =\operatorname{rank}D+\operatorname{rank}(B-C^*D^{-1}C).
\]
Because $\operatorname{rank}A=\dim P_0=\operatorname{rank}D$
throughout a sufficiently small
neighborhood, the Schur complement vanishes there.  Consequently
\[
 \ker A
 =\left\{\,u\oplus(-D^{-1}Cu):u\in K_0\,\right\}.
\]
This graph varies smoothly, proving that $K$ is a smooth subbundle.
For a self-adjoint endomorphism in finite dimension,
$(\ker A)^{\perp_h}=\operatorname{im}A$.
\end{proof}

The stable node is
\nodeid{PC-F-CONSTANT-RANK-KERNEL-SUBBUNDLE}.

\subsection{Spatial parallelism and the kernel-motion law}
\label{subsec:smp-kernel-rigidity}

\begin{proposition}[Kernel rigidity and motion]
\label{prop:smp-kernel-rigidity}
Let $J\subset(0,T]$ be a constant-rank interval and let $K=\ker A$.
Then $K(t)$ is $D^t$-parallel in space.  More precisely, if $w$ is a smooth
local section of $K$, then
\begin{align}
 D^t_Yw&\in K,                                            \label{eq:smp-parallel-conclusion}\\
 h\bigl(\beta(A,t)w,w\bigr)&=0,                          \label{eq:smp-beta-diagonal-zero}\\
 A(\partial_tw)&=-\beta(A,t)w.                           \label{eq:smp-motion-law-repeated}
\end{align}
\end{proposition}

\begin{proof}
Write
\[
       a(u,v)=h(Au,v),
       \qquad
       b(u,v)=h(\beta(A,t)u,v).
\]
Both are symmetric bilinear forms.  Since $A\geq0$ and $Aw=0$, one has
$a(w,z)=0$ for every local section $z$.

Fix a spacetime point and choose a $g(t)$-orthonormal frame
$e_1,\ldots,e_n$ which is Levi--Civita normal there.  Because
$a(w,w)$ vanishes identically and $a(w,\cdot)=0$, the product rule gives
\begin{equation}
 (\partial_ta)(w,w)=0,
 \qquad
 (D^t_{X(t)}a)(w,w)=0.                                  \label{eq:smp-time-drift-zero}
\end{equation}
For the Laplacian, first differentiate the identity
$a(w,D^t_{e_i}w)=0$:
\begin{equation}
 (D^t_{e_i}a)(w,D^t_{e_i}w)
   =-a(D^t_{e_i}w,D^t_{e_i}w).                           \label{eq:smp-cross-derivative}
\end{equation}
Expanding $\Delta(a(w,w))=0$ and using
\eqref{eq:smp-cross-derivative} yields
\begin{equation}
 (\Delta_ta)(w,w)
   =2\sum_{i=1}^n a(D^t_{e_i}w,D^t_{e_i}w).              \label{eq:smp-laplacian-null-identity}
\end{equation}
Indeed, the raw product rule has coefficient $4$ on the cross term and
coefficient $2$ on the last quadratic term; substituting
\eqref{eq:smp-cross-derivative} leaves the coefficient $2$ displayed in
\eqref{eq:smp-laplacian-null-identity}.

Evaluate the system \eqref{eq:smp-system} on $(w,w)$ and use
\eqref{eq:smp-time-drift-zero} and
\eqref{eq:smp-laplacian-null-identity}.  We obtain the decisive identity
\begin{equation}
 0
 =2\sum_{i=1}^n a(D^t_{e_i}w,D^t_{e_i}w)+b(w,w).         \label{eq:smp-decisive-kernel-identity}
\end{equation}
Every summand in the sum is nonnegative because $A\geq0$.  The last term
is nonnegative by the null-eigenvector condition.  Hence every term
vanishes.  For a nonnegative self-adjoint endomorphism,
$a(v,v)=0$ is equivalent to $Av=0$.  It follows that
$D^t_{e_i}w\in K$ for every $i$, and hence
$D^t_Yw\in K$ for every $Y$.  We also obtain $b(w,w)=0$.
This proves spatial parallelism and
\eqref{eq:smp-beta-diagonal-zero}.  Parallel transport preserves $K$ in
both directions because the reversed path gives the inverse transport.

It remains to identify temporal motion.  Use endomorphism notation again.
Spatial parallelism gives
\[
 (D^t_{X(t)}A)w
   =D^t_{X(t)}(Aw)-A(D^t_{X(t)}w)=0.
\]
It also gives $(\Delta_tA)w=0$.  To see the latter without suppressing a
product term, expand
\[
 0=\Delta_t(Aw)
   =(\Delta_tA)w
     +2\sum_i(D^t_{e_i}A)(D^t_{e_i}w)
     +A(\Delta_t^Ww).
\]
The last term vanishes because a parallel subbundle is preserved by the
rough Laplacian.  For the middle term, $D^t_{e_i}w$ is $K$-valued, and
parallelism also makes
$D^t_{e_i}(D^t_{e_i}w)$ kernel-valued at the chosen normal-frame point.
Therefore differentiating $A(D^t_{e_i}w)=0$ gives
\[
 (D^t_{e_i}A)(D^t_{e_i}w)
   +A\bigl(D^t_{e_i}D^t_{e_i}w\bigr)=0,
\]
and both terms vanish.  Hence $(\Delta_tA)w=0$.

Equation \eqref{eq:smp-system} now gives
\begin{equation}
                 (\partial_tA)w=\beta(A,t)w.             \label{eq:smp-time-A-on-kernel}
\end{equation}
Differentiate $Aw=0$ using the ordinary time derivative in the fixed fiber:
\[
 0=\partial_t(Aw)
   =(\partial_tA)w+A(\partial_tw).
\]
Substitution of \eqref{eq:smp-time-A-on-kernel} proves
\eqref{eq:smp-motion-law-repeated}.
\end{proof}

The spatial conclusion has stable node
\nodeid{PC-F-POSITIVE-SYSTEM-KERNEL-SPATIAL-PARALLEL}; the exact motion
law has stable node
\nodeid{PC-F-POSITIVE-SYSTEM-KERNEL-MOTION}.

\begin{corollary}[Sufficient hypotheses for temporal constancy]
\label{cor:smp-kernel-time-constant}
On a constant-rank interval, each of the following conditions is sufficient
for \eqref{eq:smp-kernel-annihilation}, and hence for $K(t)$ and
$\operatorname{im}A(t)$ to be independent of time as subbundles of the
fixed bundle $W$.

\begin{enumerate}
\item Direct kernel annihilation:
      $\beta(A,t)\ker A=0$.
\item Positivity of the reaction along the nonnegative cone:
      \begin{equation}
                B\geq0\quad\Longrightarrow\quad
                \beta(B,t)\geq0.                         \label{eq:smp-positive-reaction}
      \end{equation}
\item Commutation along the solution:
      \begin{equation}
                       [A,\beta(A,t)]=0.                 \label{eq:smp-commuting-reaction}
      \end{equation}
\end{enumerate}
\end{corollary}

\begin{proof}
Only the last two conditions need proof.  Under
\eqref{eq:smp-positive-reaction}, the nonnegative symmetric form
$\beta(A,t)$ satisfies
$h(\beta(A,t)w,w)=0$ for $w\in\ker A$ by
\eqref{eq:smp-beta-diagonal-zero}.  The Cauchy--Schwarz inequality for a
nonnegative symmetric form gives, for every $z$,
\[
 |h(\beta(A,t)w,z)|^2
 \leq h(\beta(A,t)w,w)h(\beta(A,t)z,z)=0.
\]
Thus $\beta(A,t)w=0$.

Under \eqref{eq:smp-commuting-reaction}, the self-adjoint endomorphism
$\beta(A,t)$ preserves $\ker A$.  Equation
\eqref{eq:smp-beta-diagonal-zero}, applied to every vector of $\ker A$ and
polarized, says that the restriction of $\beta(A,t)$ to $\ker A$ is zero.
Hence it again annihilates the kernel.

In any of the three cases, the motion law becomes
$A(\partial_tw)=0$, so $\partial_tw\in K$.  Choose a smooth local frame
$w_1,\ldots,w_p$ of $K$.  Then
\[
               \partial_tw_i=\sum_{j=1}^p c_i{}^j w_j.
\]
Solving the inverse finite-dimensional coefficient ODE produces a frame
whose ordinary time derivative is zero.  Thus the span $K_x(t)\subset W_x$
is independent of $t$ for every $x$, and these fixed fibers form the same
subbundle of $W$.  Orthogonal complementation with respect to the fixed
metric $h$ proves the assertion for $\operatorname{im}A$.
\end{proof}

The direct annihilation and temporal conclusions have stable nodes
\nodeid{PC-F-REACTION-KERNEL-ANNIHILATION} and
\nodeid{PC-F-KERNEL-TIME-CONSTANCY}; the commutation bridge has stable node
\nodeid{PC-F-COMMUTING-REACTION-KERNEL-ANNIHILATION}.

\begin{proof}[Completion of the proof of Theorem~\ref{thm:smp-systems}]
Item~(1) is Proposition~\ref{prop:smp-rank-spreading}.  Item~(2) is
Proposition~\ref{prop:smp-positive-time-rank}.  Lemma~\ref{lem:smp-smooth-kernel}
gives the smooth subbundle in item~(3), and
Proposition~\ref{prop:smp-kernel-rigidity} gives spatial parallelism,
the null-reaction identity, and the motion law.  Finally,
Corollary~\ref{cor:smp-kernel-time-constant} proves item~(5) under the
stated kernel-annihilation hypothesis.
\end{proof}

\subsection{Why the null-eigenvector condition does not freeze time}
\label{subsec:smp-source-correction}

The distinction isolated above repairs an overgeneralized statement in the
literature.  The following example satisfies every hypothesis needed for
rank spreading and spatial parallelism, but its kernel rotates in the fixed
fiber.

\begin{example}[A rotating kernel]
\label{ex:smp-rotating-kernel}
Take any connected manifold, for example $M=S^1$, with a fixed metric.
Let $W=M\times\mathbb R^2$ have its standard metric and flat connection,
and take $X=0$.  Define
\[
 J=\begin{pmatrix}0&-1\\1&0\end{pmatrix},
 \qquad
 R(t)=e^{tJ},
 \qquad
 P=\begin{pmatrix}0&0\\0&1\end{pmatrix},
\]
and set
\begin{equation}
 A(x,t)=R(t)PR(t)^{-1},
 \qquad
 \beta(B)=JB-BJ.                                         \label{eq:smp-rotating-example}
\end{equation}
The map $\beta$ sends self-adjoint endomorphisms to self-adjoint
endomorphisms and is linear, hence globally Lipschitz.  If $B\geq0$ and
$Bv=0$, then
\[
 \langle\beta(B)v,v\rangle
 =\langle JBv,v\rangle-\langle BJv,v\rangle=0.
\]
Thus the null-eigenvector condition holds with equality.  Since $A$ is
spatially constant,
\[
 \partial_tA=JA-AJ=\Delta A+\beta(A).
\]
For every $t$, $A(t)\geq0$ and $\operatorname{rank}A(t)=1$.  Nevertheless,
\begin{equation}
                  \ker A(t)=R(t)\mathbb Re_1             \label{eq:smp-rotating-kernel}
\end{equation}
rotates with time.  In particular, the null-eigenvector condition alone
does not imply that $\ker A(t)$ is an ordinary time-independent subspace of
the fixed fiber.
\end{example}

There are two independent defects to separate in the later printed
argument.  First, the proof of Proposition~12.49 on pp.~186--188 of
MSM~144 applies the null-eigenvector condition to a shifted tensor
$\bar A=A+(\varepsilon e^{ct}-f)I$ at a contact where only
$\Phi_k(\bar A)=0$ is known.  For $k>1$, that equality does not imply
$\bar A\geq0$, so the printed null-condition step is unavailable.  The
deletion tensors $A_i$ in
\eqref{eq:smp-eigenvalue-deletion} are genuinely nonnegative and repair
this rank-propagation proof without applying the null condition to the
shifted tensor.

Hamilton's original curvature-oriented strong maximum principle,
Lemma~8.2 on pp.~174--176 of \emph{Four-manifolds with positive curvature
operator}, Journal of Differential Geometry \textbf{24} (1986), 153--179,
assumes that the reaction is nonnegative whenever the solution endomorphism
is nonnegative.  That stronger hypothesis is sufficient by
Corollary~\ref{cor:smp-kernel-time-constant}.  By contrast, the temporal
constancy conclusion in Theorem~12.50(iii) of \emph{The Ricci Flow:
Techniques and Applications, Part II: Analytic Aspects}, Mathematical
Surveys and Monographs~144, pp.~188--192, is printed under only the
null-eigenvector hypothesis.  Its proof obtains the scalar equality
$\beta(A)(w,w)=0$ on p.~190 and uses the stronger vector equality
$\beta(A)w=0$ on p.~191.  Example~\ref{ex:smp-rotating-kernel} shows that
this passage needs an additional hypothesis.

Thus the $A_i$ argument reconstructs the rank-spreading proof, while the
rotating example diagnoses the separate temporal claim.  The spatial
parallelism argument remains valid after rank spreading has produced a
smooth constant-rank kernel.  The resulting public interface is
Theorem~\ref{thm:smp-systems}, whose motion law
\eqref{eq:smp-kernel-motion} records exactly what is true before a later
application verifies kernel annihilation.

The stable bridge node for this reconciliation is
\nodeid{PC-B-NULL-EIGENVECTOR-TIME-CONSTANCY-CORRECTION}.

\subsection{Formalization dependency ledger for the abstract engine}
\label{subsec:smp-formalization-ledger}

The fine nodes below all belong to the macro milestone
\nodeid{PC-M-STRONG-MAXIMUM-PRINCIPLE-SYSTEMS}.  An edge means that the
consumer actually invokes the prerequisite.

\begin{sloppypar}
\begin{enumerate}
\item \nodeid{PC-F-KY-FAN-LOWER-SPECTRAL-SUM}: Definition
      \ref{def:smp-ky-fan}, the variational formulas
      \eqref{eq:smp-ky-fan-frame}--\eqref{eq:smp-ky-fan-projection},
      continuity, scalar shifts, and the rank criterion.
\item \nodeid{PC-F-LOCAL-SCALAR-STRONG-PARABOLIC-MAXIMUM-PRINCIPLE}:
      Proposition~\ref{prop:smp-local-scalar-input}, including Dirichlet
      existence and interior strict positivity on a relatively compact
      smooth domain.
\item \nodeid{PC-F-KY-FAN-STRICT-POSITIVITY-PROPAGATION}: the
      $H_\varepsilon$ first-contact argument, depending on the first two
      nodes, the null-eigenvector condition, and local fiberwise
      Lipschitz continuity of $\beta$.
\item \nodeid{PC-F-POSITIVE-SYSTEM-RANK-SPREADING}:
      Proposition~\ref{prop:smp-rank-spreading}, depending on Ky Fan
      propagation and the rank criterion.
\item \nodeid{PC-F-POSITIVE-TIME-SPATIAL-RANK-CONSTANCY}:
      Proposition~\ref{prop:smp-positive-time-rank}, depending on rank
      spreading and continuity of the finite-dimensional spectrum.
\item \nodeid{PC-F-CONSTANT-RANK-KERNEL-SUBBUNDLE}:
      Lemma~\ref{lem:smp-smooth-kernel}, depending on the smooth
      constant-rank bundle-map theorem, here exposed through the Schur
      complement calculation.
\item \nodeid{PC-F-POSITIVE-SYSTEM-KERNEL-SPATIAL-PARALLEL}:
      the identity \eqref{eq:smp-decisive-kernel-identity}, depending on
      smooth local kernel sections, metric compatibility, nonnegativity of
      $A$, and the null-eigenvector condition.
\item \nodeid{PC-F-POSITIVE-SYSTEM-KERNEL-MOTION}: the exact law
      \eqref{eq:smp-motion-law-repeated}, depending on spatial parallelism
      and the product rule for the time-dependent rough Laplacian.
\item \nodeid{PC-F-REACTION-KERNEL-ANNIHILATION} and
      \nodeid{PC-F-COMMUTING-REACTION-KERNEL-ANNIHILATION}: the two
      reusable sufficient conditions in
      Corollary~\ref{cor:smp-kernel-time-constant}.
\item \nodeid{PC-F-KERNEL-TIME-CONSTANCY}: constancy of the subspace in
      the fixed fiber, depending on the motion law, kernel annihilation,
      and a finite-dimensional linear ODE for a kernel frame.
\item \nodeid{PC-B-NULL-EIGENVECTOR-TIME-CONSTANCY-CORRECTION}:
      Example~\ref{ex:smp-rotating-kernel} and the reconciliation of the
      original positive-reaction theorem with the overgeneralized later
      formulation.
\item \nodeid{PC-B-REACTION-LOCAL-LIPSCHITZ-NECESSITY}: the scalar
      finite-time rank-loss examples showing that a mere continuous
      null-condition reaction does not support rank spreading.
\end{enumerate}
\end{sloppypar}

The most delicate future Lean interfaces are the local scalar Dirichlet
theorem, the construction of a smooth orthonormal frame with the precise
contact-point jets \eqref{eq:smp-frame-spatial-jets}--
\eqref{eq:smp-frame-time-jet}, and the rough-Laplacian product rule for a
time-dependent metric connection.  No step requires a smooth eigenbasis
through an eigenvalue multiplicity, a global trivialization of $W$, a
complete base metric, or a global bound on the reaction.

\section{The three-dimensional curvature reaction}
\label{sec:curvature-three-dimensional-reaction}

We now specialize the strong maximum principle to the curvature operator of
a three-dimensional Ricci flow.  This specialization is deliberately
separated from the global splitting argument.  The strong maximum principle
is local and applies without completeness; the passage from a parallel line
distribution to a product on the universal cover uses completeness.

Throughout this section, $W$ denotes a three-dimensional Euclidean space and
$A\colon W\to W$ a self-adjoint endomorphism.  We use the polynomials
\begin{align}
 \sigma _2(A)
   &=\frac12\bigl((\trace A)^2-\trace(A^2)\bigr),
                                                        \label{eq:smp-sigma-two}\\
 A^\#&=A^2-(\trace A)A+\sigma _2(A)I,                  \label{eq:smp-adjugate}\\
 Q(A)&=A^2+A^\#.                                      \label{eq:smp-curvature-reaction}
\end{align}
These are the same polynomials as in
\eqref{eq:hi-sigma-two}--\eqref{eq:hi-reaction}.  We repeat the notation
because the sign and the factor of two in the curvature convention will be
used at every subsequent step.

\begin{lemma}[Positivity and commutation of the curvature reaction]
\label{lem:curvature-reaction-positive-commuting}
If $A\geq0$, then
\[
                 A^\#\geq0,
        \qquad    Q(A)\geq0,
        \qquad    A Q(A)=Q(A)A.
\]
More precisely, if an orthonormal basis diagonalizes $A$ as
$\operatorname{diag}(a,b,c)$, then
\begin{align}
 A^\#&=\operatorname{diag}(bc,ac,ab),                  \label{eq:smp-adjugate-diagonal}\\
 Q(A)&=\operatorname{diag}
       (a^2+bc,b^2+ac,c^2+ab).                         \label{eq:smp-Q-diagonal}
\end{align}
\end{lemma}

\begin{proof}
The characteristic polynomial of $A$ gives
\eqref{eq:smp-adjugate-diagonal}, and adding $A^2$ gives
\eqref{eq:smp-Q-diagonal}.  If $a,b,c\geq0$, every displayed diagonal entry
is nonnegative.  Finally,
\[
 Q(A)=2A^2-(\trace A)A+\sigma _2(A)I
\]
is a polynomial in $A$, so it commutes with $A$.  This proof uses only the
spectral theorem in dimension three and polynomial arithmetic.
\end{proof}

\begin{corollary}[Closed-flow preservation of nonnegative curvature]
\label{cor:closed-curvature-operator-preserved}
Let $M^3$ be closed and let $g(t)$ be a Ricci flow on an interval containing
$[s,T)$.  If
\[
                    \mathcal R_{g(s)}\geq0,
\]
then $\mathcal R_{g(t)}\geq0$ for every $t\in[s,T)$.
\end{corollary}

\begin{proof}
In Uhlenbeck gauge the curvature operator solves
\[
                 \partial_t\mathcal M
                  =\Delta_{D^t,g(t)}\mathcal M+Q(\mathcal M).
\]
Lemma~\ref{lem:curvature-reaction-positive-commuting} shows that, at every
$A\geq0$ and every $v\in\ker A$,
\[
                        \langle Q(A)v,v\rangle\geq0.
\]
Equivalently, the ODE $A'=Q(A)$ points into the tangent cone of the closed
convex cone of positive-semidefinite endomorphisms; its flow preserves that
cone.  Theorem~\ref{thm:hamilton-wmp-systems}, applied on the closed base
after shifting the initial time to $s$, therefore preserves
$\mathcal M\geq0$.  The Uhlenbeck isometry transfers this inequality back to
$\mathcal R_{g(t)}\geq0$.
\end{proof}

\begin{remark}[Scope of the preservation corollary]
\label{rem:closed-curvature-preservation-scope}
Corollary~\ref{cor:closed-curvature-operator-preserved} is stated on a
closed base because that is the scope of the weak systems theorem proved in
Chapter~2.  The noncompact theorems below continue to assume nonnegativity
on the whole spacetime region; they do not silently import a noncompact weak
maximum principle.
\end{remark}

The conclusion needed from Theorem~\ref{thm:smp-systems} is stronger than
mere preservation of the positive-semidefinite cone.  We spell out the
short PDE calculation, both to expose the sign and to make the
kernel-annihilation interface independently reusable.

\begin{proposition}[Kernel annihilation by the curvature reaction]
\label{prop:curvature-kernel-annihilation}
Let $\mathcal M\geq0$ solve the Uhlenbeck-gauge curvature equation
\[
                 \partial_t\mathcal M
                 =\Delta_{D^t,g(t)}\mathcal M+Q(\mathcal M).
\]
On every positive-time constant-rank slab on which
Theorem~\ref{thm:smp-systems} supplies the smooth spatially parallel bundle
$K=\ker\mathcal M$, one has
\begin{equation}
       \ker\mathcal M(x,t)\subseteq\ker Q(\mathcal M(x,t)).
                                                        \label{eq:smp-null-reaction}
\end{equation}
\end{proposition}

\begin{proof}
Let $W$ be a local smooth spacetime section of $K$.  The scalar identity
$\langle\mathcal MW,W\rangle=0$ implies, at a point and in a
$g(t)$-orthonormal normal frame $(e_i)$,
\begin{equation}
 (\Delta_{D^t,g(t)}\mathcal M)(W,W)
   =2\sum_i\langle\mathcal M D^t_{e_i}W,D^t_{e_i}W\rangle.
                                                        \label{eq:smp-kernel-laplacian-identity}
\end{equation}
Indeed, differentiate $\mathcal MW=0$ once to obtain
$(D_i\mathcal M)W=-\mathcal M D_iW$, and differentiate
$(D_i\mathcal M)(W,W)=0$ once more.  Terms containing $\partial_tW$
vanish for the same reason.  Evaluating the evolution equation on $(W,W)$
therefore gives
\[
 0=2\sum_i\langle\mathcal M D^t_{e_i}W,D^t_{e_i}W\rangle
     +\langle Q(\mathcal M)W,W\rangle.
\]
Both terms are nonnegative by $\mathcal M\geq0$ and
Lemma~\ref{lem:curvature-reaction-positive-commuting}; hence the final
quadratic form vanishes.  Since $Q(\mathcal M)\geq0$, diagonalization (or its
positive square root) implies $Q(\mathcal M)W=0$.  This is
\eqref{eq:smp-null-reaction}.  Spatial parallelism actually makes every
summand in \eqref{eq:smp-kernel-laplacian-identity} vanish separately.
\end{proof}

\begin{corollary}[Fixed-bundle constancy on a rank-stable interval]
\label{cor:curvature-fixed-bundle-rank-interval}
On every time interval $J\subset(0,T)$ on which the rank of
$\mathcal M$ is constant, the subbundles
\[
                   \ker\mathcal M(t),
       \qquad      \operatorname{im}\mathcal M(t)
\]
of the fixed Uhlenbeck bundle are independent of $t\in J$.
\end{corollary}

\begin{proof}
Proposition~\ref{prop:curvature-kernel-annihilation} verifies the
kernel-annihilation hypothesis in item~5 of
Theorem~\ref{thm:smp-systems}.  That item makes the kernel independent of
time in the fixed bundle.  Because the fiber metric is fixed and
$\mathcal M$ is self-adjoint,
$\operatorname{im}\mathcal M=(\ker\mathcal M)^\perp$ is independent of time
as well.  Equivalently, one may invoke the positive-reaction or the
commuting-reaction criterion, both of which follow from
Lemma~\ref{lem:curvature-reaction-positive-commuting}.
\end{proof}

\begin{lemma}[One-sided transfer to a rank-stable endpoint]
\label{lem:curvature-one-sided-endpoint}
Let $0<a<b<T$, and suppose that the rank of $\mathcal M$ has one constant
value on $M\times(a,b]$.  Then
\[
                     K_b=\ker\mathcal M(\cdot,b)
\]
is a smooth $D^b$-parallel subbundle of the fixed Uhlenbeck bundle and
\begin{equation}
                     K_b\subseteq\ker Q(\mathcal M(\cdot,b)).
                                                        \label{eq:smp-endpoint-null-reaction}
\end{equation}
\end{lemma}

\begin{proof}
On the open interval $(a,b)$, Corollary
\ref{cor:curvature-fixed-bundle-rank-interval} gives one smooth subbundle
$K$ of the fixed bundle such that
\[
       K=\ker\mathcal M(\cdot,t)
       \qquad(a<t<b).
\]
Fix $x\in M$ and $v\in K_x$.  Since
$\mathcal M(x,t)v=0$ for every $t<b$, continuity gives
$\mathcal M(x,b)v=0$.  Thus
$K_x\subseteq\ker\mathcal M(x,b)$.  The two spaces have the same
dimension by the rank hypothesis, so equality holds.  In particular, the
endpoint kernel is the already smooth bundle $K$.

Let $w$ be a time-independent smooth local section of $K$.  For $a<t<b$,
spatial parallelity and kernel annihilation give
\[
              D_Y^tw\in K,
       \qquad Q(\mathcal M(t))w=0.
\]
The family of connections is smooth in time, so
$D_Y^tw\to D_Y^bw$.  The reaction $Q$ is a polynomial and
$\mathcal M(t)\to\mathcal M(b)$, so
$Q(\mathcal M(t))w\to Q(\mathcal M(b))w$.  Because the finite-dimensional
fibers $K_x$ are closed, passage to the limit proves
\[
              D_Y^bw\in K,
       \qquad Q(\mathcal M(b))w=0.
\]
These are precisely $D^b$-parallelity and
\eqref{eq:smp-endpoint-null-reaction}.
\end{proof}

The stable endpoint node is
\nodeid{PC-F-CURVATURE-ONE-SIDED-ENDPOINT-TRANSFER}.

Proposition~\ref{prop:smp-rank-spreading} supplies the required
positive-time constant-rank slabs.  We now isolate the elementary algebraic
consequence which replaces the usual classification of Lie subalgebras of
$\mathfrak{so}(3)$.

\begin{lemma}[The null-reaction condition excludes rank two]
\label{lem:curvature-rank-two-excluded}
Let $A\geq0$ be self-adjoint on a three-dimensional Euclidean space and
suppose
\begin{equation}
                    \ker A\subseteq\ker Q(A).            \label{eq:smp-algebraic-null-reaction}
\end{equation}
Then $\operatorname{rank}A\in\{0,1,3\}$.  In particular,
$\operatorname{rank}A\neq2$.
\end{lemma}

\begin{proof}
Suppose $\operatorname{rank}A=2$.  Choose an orthonormal eigenbasis with
\[
                    A=\operatorname{diag}(a,b,0),
                    \qquad a,b>0.
\]
For a unit vector $v$ in the zero eigenspace,
\eqref{eq:smp-Q-diagonal} gives
\[
                       Q(A)v=ab\,v\neq0,
\]
contrary to \eqref{eq:smp-algebraic-null-reaction}.  The remaining possible
ranks are $0,1,3$.  Notice also that, since $A^2v=0$ for $v\in\ker A$,
\eqref{eq:smp-algebraic-null-reaction} implies
\[
             \ker A\subseteq\ker A^\#.
\]
Thus the full kernel-annihilation assertion is visible in the polynomial
reaction, not hidden in a Lie-algebra argument.
\end{proof}

\begin{theorem}[Positive-time curvature trichotomy]
\label{thm:curvature-time-slice-trichotomy}
Let $M^3$ be nonempty, connected, and without boundary, and let
$g(t)$, $t\in[0,T)$, be a smooth Ricci flow.  Assume that its curvature
operator is nonnegative on $M\times[0,T)$.  Fix $t_0\in(0,T)$.  Exactly one
of the following alternatives holds on the entire time slice
$(M,g(t_0))$.
\begin{enumerate}
\item \emph{Flat alternative.}
      $\mathcal R_{g(t_0)}=0$ everywhere.
\item \emph{Rank-one alternative.}
      $\mathcal R_{g(t_0)}$ has rank one everywhere.  Its ordered
      eigenvalues satisfy
      \[
                          \lambda>0,\qquad\mu=\nu=0.
      \]
      Every point has a neighborhood isometric to a product of an interval
      and a surface of positive Gaussian curvature.
\item \emph{Positive alternative.}
      $\mathcal R_{g(t_0)}$ is positive definite everywhere; equivalently,
      \[
                          \lambda\geq\mu\geq\nu>0.
      \]
\end{enumerate}
No completeness, bounded-curvature, or orientability hypothesis is used in
this theorem.
\end{theorem}

\begin{proof}
Use the Uhlenbeck isometries to regard the curvature operator as a section
$\mathcal M$ of the fixed Euclidean bundle
$\operatorname{SymEnd}(\Lambda^2V)$.  It satisfies
\[
                 \partial_t\mathcal M
                 =\Delta_{D^t,g(t)}\mathcal M+Q(\mathcal M).
\]
Lemma~\ref{lem:curvature-reaction-positive-commuting} verifies the
null-eigenvector condition required by Theorem~\ref{thm:smp-systems}.
That theorem shows that the rank is independent of $x$ at $t_0$ and has
the same value on $M\times(a,t_0]$ for some $a\in(0,t_0)$.  Lemma
\ref{lem:curvature-one-sided-endpoint} now shows that the kernel at $t_0$
is a smooth $D^{t_0}$-parallel subbundle and gives the endpoint
null-reaction condition
\eqref{eq:smp-endpoint-null-reaction}.  Lemma
\ref{lem:curvature-rank-two-excluded} leaves only ranks $0,1,3$.

Rank zero is the flat alternative.  In rank one, nonnegativity gives one
strictly positive eigenvalue and two zero eigenvalues.  In rank three all
three eigenvalues are positive.  In rank one,
Proposition~\ref{prop:ue-nullspace-transfer} carries the fixed-bundle
parallel kernel to the parallel nullspace of the geometric curvature
operator.  The asserted local product is then proved, without choosing an
orientation, in
Lemma~\ref{lem:curvature-image-to-parallel-line} and
Proposition~\ref{prop:parallel-line-local-product} below.
\end{proof}

\begin{remark}[Why the theorem starts at positive time]
\label{rem:curvature-trichotomy-positive-time}
The restriction $t_0>0$ is essential.  At $t=0$ the rank can vary with the
point, and an initially cylindrical or flat open region can instantly acquire
positive curvature when it sits inside a larger incomplete flow.  The local
nonuniqueness example in GSM~77, Example~6.59, pp.~246--247, exhibits
exactly this phenomenon.  Completeness will enter only when we ask for a
global product; uniqueness will enter only when we ask for one product that
persists through a time interval.
\end{remark}

\section{From a rank-one curvature image to a tangent line}
\label{sec:curvature-rank-one-line}

The customary three-dimensional proof applies a Hodge star to a parallel
two-form.  That is efficient on an oriented manifold, but orientation is
irrelevant to the result and would be an artificial formalization
dependency.  The following construction uses only contraction of forms.

\begin{lemma}[The orientation-free tangent-line bridge]
\label{lem:curvature-image-to-parallel-line}
Let $(U^3,g)$ be a Riemannian manifold whose nonnegative curvature operator
$\mathcal R$ has rank one.  Suppose that
$\ker\mathcal R\subset\Lambda^2TU$ is parallel.  Set
\[
                 \mathcal I=\operatorname{im}\mathcal R.
\]
For $\omega\in\Lambda^2T_xU$, let $\omega^\flat\in\Lambda^2T_x^*U$ be
obtained by lowering both indices with $g$.  Define
\begin{equation}
 L_x=\left\{X\in T_xU:\iota_X\omega^\flat=0
                  \text{ for every }\omega\in\mathcal I_x\right\}.
                                                        \label{eq:smp-line-annihilator}
\end{equation}
Then $L\subset TU$ is a smooth parallel line subbundle.  Its orthogonal
complement $P=L^\perp$ is a smooth parallel two-plane subbundle, and
\begin{equation}
                 \mathcal I_x=\Lambda^2P_x.
                                                        \label{eq:smp-curvature-image-plane}
\end{equation}
This construction requires no orientation of $U$ and is unchanged if a
local generator of $\mathcal I$ is multiplied by a nonzero function.
\end{lemma}

\begin{proof}
Self-adjointness gives
\[
       \mathcal I=(\ker\mathcal R)^\perp.
\]
The metric connection preserves orthogonal complements, so $\mathcal I$ is
a parallel line subbundle.  Every nonzero two-form on a three-dimensional
vector space has a one-dimensional contraction kernel.  Consequently
\eqref{eq:smp-line-annihilator} defines a smooth line bundle, and elementary
linear algebra gives \eqref{eq:smp-curvature-image-plane}.

It remains to check parallelism.  Locally choose a unit section $\omega$ of
$\mathcal I$.  Because $\mathcal I$ is parallel, $\nabla_Y\omega$ lies in
$\mathcal I$ for every $Y$.  Since $|\omega|=1$,
$\langle\nabla_Y\omega,\omega\rangle=0$; hence
\begin{equation}
                         \nabla\omega=0.                 \label{eq:smp-local-parallel-two-form}
\end{equation}
Metric compatibility gives $\nabla(\omega^\flat)=(\nabla\omega)^\flat=0$.
If $X$ is a local section of $L$, then $\iota_X\omega^\flat=0$.  Covariantly
differentiating and using \eqref{eq:smp-local-parallel-two-form} yields
\[
                 0=\nabla_Y(\iota_X\omega^\flat)
                   =\iota_{\nabla_YX}\omega^\flat.
\]
Thus $\nabla_YX\in L$.  Hence $L$ is parallel, and metric compatibility
makes $P=L^\perp$ parallel as well.  Replacing $\omega$ by $-\omega$ has no
effect, which is the promised orientation-free feature.
\end{proof}

We next record all normalization factors.  This avoids a common but serious
ambiguity: in this book the eigenvalues of the curvature operator are twice
the corresponding sectional curvatures.

\begin{lemma}[Rank-one normalization and Ricci kernel]
\label{lem:rank-one-normalization}
Assume the hypotheses of
Lemma~\ref{lem:curvature-image-to-parallel-line}.  At a point choose a unit
vector $e_3\in L$ and an orthonormal basis $e_1,e_2$ of $P$.  If $\lambda$
is the nonzero eigenvalue of $\mathcal R$, then
\begin{align}
 \mathcal R(e_1\wedge e_2)&=\lambda(e_1\wedge e_2),
      & K(e_1,e_2)&=\frac{\lambda}{2}>0,                 \label{eq:smp-rank-one-sectional}\\
 \mathcal R(e_i\wedge e_3)&=0,
      & K(e_i,e_3)&=0\quad(i=1,2),                       \label{eq:smp-rank-one-mixed}\\
 R&=\lambda=2K(e_1,e_2),
      & \Ric|_P&=\frac{\lambda}{2}g|_P,
      \qquad \Ric(e_3,\cdot)=0.                         \label{eq:smp-rank-one-ricci}
\end{align}
In particular,
\begin{equation}
                         L=\ker\Ric^\sharp.              \label{eq:smp-line-is-ricci-kernel}
\end{equation}
\end{lemma}

\begin{proof}
Equation \eqref{eq:smp-curvature-image-plane} says that the image is
spanned by the unit two-vector $e_1\wedge e_2$.  Since $\mathcal R$ is
self-adjoint,
\[
       \ker\mathcal R=(\operatorname{im}\mathcal R)^\perp,
\]
so both mixed two-planes lie in the kernel.  The normalization
\eqref{eq:hi-curvature-normalization} says, for orthonormal $u,v$, that
\[
        \langle\mathcal R(u\wedge v),u\wedge v\rangle
                       =2K(u,v).
\]
This proves \eqref{eq:smp-rank-one-sectional} and
\eqref{eq:smp-rank-one-mixed}.  Taking the two sectional-curvature traces
that define Ricci curvature gives
\[
 \Ric(e_1,e_1)=\Ric(e_2,e_2)=\frac{\lambda}{2},
 \qquad \Ric(e_3,e_3)=0,
 \qquad \Ric(e_i,e_j)=0\ (i\neq j).
\]
Finally, scalar curvature is twice the sum of the three sectional curvatures,
so $R=\lambda$.  This proves \eqref{eq:smp-rank-one-ricci} and
\eqref{eq:smp-line-is-ricci-kernel}.
\end{proof}

\begin{proposition}[The rank-one Uhlenbeck line is physically fixed]
\label{prop:rank-one-uhlenbeck-line-fixed}
Let $J\subset(0,T)$ be an interval on which the curvature operator has rank
one.  Use the fixed bundle $(V,h)$ and the Uhlenbeck isometries
$\iota_t\colon(V,h)\to(TM,g(t))$ satisfying
\eqref{eq:ue-ricci-iota-ode}.  Then the tangent line $L(t)$ defined by
\eqref{eq:smp-line-annihilator} is one and the same subbundle
$L\subset TM$ for all $t\in J$.

Use $h$ to identify the fixed line
$\operatorname{im}\mathcal M(t)\subset\Lambda^2V$ with a line
$(\operatorname{im}\mathcal M(t))^\flat\subset\Lambda^2V^*$.  More
locally, if $v$ is an $h$-unit, time-independent section of the contraction
annihilator of this covector line, then
\begin{equation}
                         X=\iota_tv                    \label{eq:smp-fixed-physical-line-field}
\end{equation}
is independent of $t$, is $g(t)$-unit, and satisfies
$\nabla^{g(t)}X=0$ for every $t\in J$.  The one-form
\begin{equation}
                         \alpha=g(t)(X,\cdot)            \label{eq:smp-fixed-physical-line-form}
\end{equation}
is also independent of time.  Consequently the rank-one flow has, locally
on $M\times J$, one fixed product chart
\[
                           g(t)=h(t)+ds^2,
\]
and $h(t)$ is a surface Ricci flow with positive Gaussian curvature.
\end{proposition}

\begin{proof}
Corollary~\ref{cor:curvature-fixed-bundle-rank-interval} makes
$\operatorname{im}\mathcal M(t)$ independent of time.  Its contraction
annihilator in $V$, after the displayed $h$-identification, is therefore a
fixed $h$-line.  Since $\iota_t$ is an isometry and
$\mathcal M=(\Lambda^2\iota_t)^{-1}\circ\mathcal R_{g(t)}
\circ(\Lambda^2\iota_t)$, the map $\iota_t$ carries this line exactly onto the
physical line $L(t)$ in \eqref{eq:smp-line-annihilator}.  It is
$D^t$-parallel for
each $t$, so \eqref{eq:ue-iota-parallel} makes $X=\iota_tv$ parallel for
$g(t)$.  Lemma~\ref{lem:rank-one-normalization} gives
$\Ric^\sharp_{g(t)}X=0$.  Hence the Uhlenbeck ODE gives
\[
       \partial_tX
        =\partial_t(\iota_tv)
        =\Ric^\sharp_{g(t)}(\iota_tv)=0.
\]
This proves that the physical line and the unit vector field are fixed.  As
$X$ is time independent,
\[
 \partial_t\bigl(g(t)(X,\cdot)\bigr)
       =-2\Ric_{g(t)}(X,\cdot)=0,
\]
which proves the assertion about $\alpha$.

Locally write $\alpha=ds$.  The level surfaces of $s$ and the flow lines of
$X$ are independent of time.  Proposition
\ref{prop:parallel-line-local-product}, applied in these fixed coordinates,
gives $g(t)=h(t)+ds^2$.  Restricting the Ricci-flow equation to a level
surface gives $\partial_th=-2\Ric_h$; positivity of its Gaussian curvature
is \eqref{eq:smp-rank-one-sectional}.
\end{proof}

\begin{proposition}[Local product from the curvature line]
\label{prop:parallel-line-local-product}
Under the hypotheses of
Lemma~\ref{lem:curvature-image-to-parallel-line}, every point of $U$ has a
neighborhood isometric to
\[
                      (\Sigma^2,h)\times(J,ds^2),
\]
where $J$ is an interval and $h$ has positive Gaussian curvature.  Under
this isometry $L$ is the tangent line to $J$ and $P$ is tangent to
$\Sigma$.
\end{proposition}

\begin{proof}
Because the Levi-Civita connection is torsion free, the parallel
distributions $L$ and $P$ are involutive: for sections $Y,Z$ of either
distribution,
\[
                         [Y,Z]=\nabla_YZ-\nabla_ZY
\]
lies in that distribution.  Frobenius therefore gives transverse local
leaves.  A local unit section $X$ of $L$ is parallel: parallelism gives
$\nabla_YX\in L$, while differentiation of $|X|^2=1$ makes
$\nabla_YX\perp X$.  Thus $\nabla X=0$.  Its local flow preserves $g$ because
\[
       (\mathcal L_Xg)(Y,Z)
        =g(\nabla_YX,Z)+g(Y,\nabla_ZX)=0,
\]
and it carries one $P$-leaf isometrically to the neighboring $P$-leaves.
The resulting flow-box map pulls $g$ back to $h+ds^2$.  Positivity of the
Gaussian curvature of $h$ is
\eqref{eq:smp-rank-one-sectional}.
\end{proof}

\section{A complete global splitting from one parallel line}
\label{sec:parallel-line-global-splitting}

For later formalization we prove precisely the special case of de Rham's
theorem that is needed here.  This proof avoids the full holonomy
decomposition and turns the isometry into an explicit map with an explicit
inverse.

\begin{theorem}[Complete simply connected parallel-line splitting]
\label{thm:parallel-line-splitting}
Let $(N^3,g)$ be nonempty, connected, complete, and simply connected.  If
$L\subset TN$ is a smooth $g$-parallel line subbundle, then there are a
connected, complete, simply connected Riemannian surface $(\Sigma^2,h)$ and
an isometry
\begin{equation}
             F\colon(\Sigma\times\R,h+ds^2)\longrightarrow(N,g)
                                                        \label{eq:smp-global-product-isometry}
\end{equation}
such that $dF(\R\partial_s)=L$.
\end{theorem}

\begin{proof}
Fix $p_0\in N$ and a unit vector $v_0\in L_{p_0}$.  Parallel translation in
$L$ along a path starting at $p_0$ preserves length.  The induced connection
on $L$ is flat: indeed, on any open set carrying a unit section $Y$ of $L$,
parallelism gives $\nabla_ZY\in L$, while
$0=Z|Y|^2=2g(\nabla_ZY,Y)$, and hence $\nabla Y=0$.  Flat parallel
translation is invariant under endpoint-fixed homotopy.  Concretely,
parallel transport around a loop is an element of
$O(L_{p_0})=\{+1,-1\}$; it varies continuously under a homotopy of loops and
therefore is constant during that homotopy.  Since $N$ is simply connected,
every loop contracts to the constant loop, whose transport is $+1$.
Consequently, translating $v_0$ from $p_0$ to $p$ is independent of the
chosen path.  We obtain a global unit section $X$ of $L$ satisfying
\begin{equation}
                              \nabla X=0.                 \label{eq:smp-global-parallel-vector}
\end{equation}

Let $\alpha=g(X,\cdot)$.  Then $\nabla\alpha=0$, hence $d\alpha=0$.  Define
\begin{equation}
       r(p)=\int_\gamma\alpha,                            \label{eq:smp-global-potential}
\end{equation}
where $\gamma$ is any path from $p_0$ to $p$.  Closedness of $\alpha$ and
simple connectivity make this independent of $\gamma$.  Thus
\begin{equation}
                         dr=\alpha,
             \qquad     Xr=1.                            \label{eq:smp-potential-properties}
\end{equation}
In particular, $0$ is a regular value and
\[
                         \Sigma=r^{-1}(0)
\]
is an embedded surface with $T\Sigma=X^\perp$.

The integral curve of $X$ through any point is a unit-speed geodesic by
\eqref{eq:smp-global-parallel-vector}.  Completeness of $g$ therefore makes
the flow $\Phi_s$ of $X$ defined for every $s\in\R$.  The same equation gives
$\mathcal L_Xg=0$, so each $\Phi_s$ is an isometry.  Moreover,
\eqref{eq:smp-potential-properties} gives
\begin{equation}
                       r(\Phi_s(p))=r(p)+s.               \label{eq:smp-potential-flow}
\end{equation}
Define
\[
             F(y,s)=\Phi_s(y),
             \qquad y\in\Sigma,\ s\in\R.
\]
Equation \eqref{eq:smp-potential-flow} supplies the explicit inverse
\begin{equation}
       p\longmapsto\bigl(\Phi_{-r(p)}(p),r(p)\bigr).       \label{eq:smp-global-product-inverse}
\end{equation}
Thus $F$ is a diffeomorphism.  It sends $\partial_s$ to the unit vector
$X$, sends $T\Sigma$ to $X^\perp$, and restricts on each slice to the
isometry $\Phi_s$.  Consequently
\[
                             F^*g=h+ds^2,
             \qquad h=g|_{T\Sigma}.
\]

It remains to verify the three adjectives asserted for $\Sigma$.  If
$\Sigma$ had more than one connected component, then
$\Sigma\times\R$ would be disconnected, contrary to the connectedness of
$N$ and the fact that $F$ is a diffeomorphism.  If $(y_j)$ is Cauchy in
$(\Sigma,h)$, then $(y_j,0)$ is Cauchy in the complete product
$\Sigma\times\R\cong N$; its limit has $r$-coordinate zero, so it lies in
$\Sigma$.  Thus $\Sigma$ is complete.  Finally,
$\Sigma\times\R\cong N$ is simply connected and $\R$ is contractible, so
$\Sigma$ is simply connected.
\end{proof}

\begin{remark}[No hidden line hypothesis]
\label{rem:parallel-line-versus-geodesic-line}
Theorem~\ref{thm:parallel-line-splitting} assumes a parallel
\emph{distribution}, which the curvature argument produces.  It does not
assume in advance the existence of a distance-minimizing geodesic line and
does not invoke the Cheeger--Gromoll splitting theorem.  Completeness is used
exactly once, to make the flow of the unit parallel field exist for all
$s\in\R$.
\end{remark}

\section{The fixed-time universal-cover theorem}
\label{sec:curvature-fixed-time-splitting}

\begin{theorem}[Complete three-dimensional time-slice classification]
\label{thm:curvature-complete-trichotomy}
Under the hypotheses of
Theorem~\ref{thm:curvature-time-slice-trichotomy}, fix $t_0\in(0,T)$ and
assume in addition that $(M,g(t_0))$ is complete.  Let
$\pi\colon\widetilde M\to M$ be the universal covering and put
$\widetilde g(t_0)=\pi^*g(t_0)$.  Exactly one of the following holds.
\begin{enumerate}
\item $\mathcal R_{g(t_0)}=0$, and
      $(\widetilde M,\widetilde g(t_0))$ is isometric to Euclidean
      three-space.
\item There are a complete simply connected surface $(\Sigma,h_{t_0})$
      of positive Gaussian curvature and an isometry
      \begin{equation}
       (\widetilde M,\widetilde g(t_0))
          \cong(\Sigma\times\R,h_{t_0}+ds^2).             \label{eq:smp-fixed-time-product}
      \end{equation}
      The surface $\Sigma$ is diffeomorphic to $\Sph^2$ when it is compact
      and to $\R^2$ when it is noncompact.
\item $\mathcal R_{\widetilde g(t_0)}>0$ everywhere.
\end{enumerate}
No orientability assumption is required.  If $M$ is closed, then the surface
in the second alternative is compact and hence diffeomorphic to $\Sph^2$.
\end{theorem}

\begin{proof}
The lifted metric is complete and simply connected.  In the flat
alternative, the flat case of Cartan--Hadamard identifies its exponential
map at any point with a global Euclidean isometry.

In the rank-one alternative, the Uhlenbeck connection at the fixed time is
the pullback of the Levi-Civita connection under the fiber isometry.
Consequently the parallel curvature kernel supplied by
Theorem~\ref{thm:smp-systems} transfers, by
Proposition~\ref{prop:ue-nullspace-transfer}, to a parallel kernel of the
geometric curvature operator.  Lemma~\ref{lem:curvature-image-to-parallel-line}
produces a parallel tangent line on $\widetilde M$, and
Theorem~\ref{thm:parallel-line-splitting} gives
\eqref{eq:smp-fixed-time-product}.  Lemma
\ref{lem:rank-one-normalization} says that the Gaussian curvature of
$h_{t_0}$ is $\lambda/2>0$.  The classification of simply connected
surfaces says that $\Sigma$ is diffeomorphic to $\Sph^2$ if compact and to
$\R^2$ if noncompact.

Suppose now that $M$ is closed.  The positive function $\lambda(\cdot,t_0)$
on $M$ has a positive minimum.  Hence the lifted surface satisfies
\[
                    K_{h_{t_0}}=\frac{\lambda}{2}\geq c>0.
\]
Bonnet--Myers makes $\Sigma$ compact.  A connected, compact, simply
connected surface is diffeomorphic to $\Sph^2$.
\end{proof}

\begin{corollary}[Equality at one positive-time point]
\label{cor:curvature-zero-eigenvalue-splitting}
Assume the hypotheses of
Theorem~\ref{thm:curvature-time-slice-trichotomy}, and let
$t_0\in(0,T)$.  If the least curvature-operator eigenvalue satisfies
$\nu(x_0,t_0)=0$ at one point, then precisely one of the following occurs.
\begin{enumerate}
\item $R(x_0,t_0)=0$, and the entire time slice $(M,g(t_0))$ is flat.
\item $R(x_0,t_0)>0$, and the curvature operator has rank one everywhere on
      the time slice.
\end{enumerate}
In the second case, if $(M,g(t_0))$ is complete, its universal cover has the
product form \eqref{eq:smp-fixed-time-product}; if $M$ is closed, its surface
factor is diffeomorphic to $\Sph^2$.

The same conclusion holds if one merely assumes that some sectional
curvature vanishes at $(x_0,t_0)$.
\end{corollary}

\begin{proof}
Nonnegativity and $\nu(x_0,t_0)=0$ exclude the positive alternative of
Theorem~\ref{thm:curvature-time-slice-trichotomy}.  If the trace is zero,
all three nonnegative eigenvalues vanish at $x_0$; spatial constancy of rank
then gives the flat alternative.  If the trace is positive, the rank is
nonzero and therefore equals one.  The complete conclusions follow from
Theorem~\ref{thm:curvature-complete-trichotomy}.

Finally, if $\omega$ represents a two-plane of zero sectional curvature,
then
\[
                     \langle\mathcal R\omega,\omega\rangle=0.
\]
For a positive-semidefinite self-adjoint operator, vanishing of this
quadratic form implies $\mathcal R\omega=0$ (diagonalize the operator, or
apply its positive square root).  Thus the least eigenvalue is zero, and the
first part applies.
\end{proof}

\begin{remark}[What the universal-cover product does not say]
\label{rem:curvature-product-quotients}
The base manifold need not itself be a product.  A deck transformation
preserves the curvature-determined pair $L\oplus P$, but it may translate or
reflect the $\R$ factor while acting isometrically on the surface.  Thus the
theorem includes the standard cylinder $\Sph^2\times\R$, products
$\Sph^2\times\Sph^1$, twisted quotients, and nonorientable quotients.  The
flat alternative similarly includes all complete flat quotients, not only
$\R^3$.  These are features of the statement, not exceptional cases to be
removed by an orientability assumption.
\end{remark}

\section{Persistence through time: the uniqueness input}
\label{sec:curvature-whole-flow-splitting}

The preceding results are time-slice theorems.  On an incomplete flow the
rank can jump, so a single product for all times is false without an
additional analytic input.  We state that input explicitly.

\begin{theorem}[Whole-flow trichotomy under complete-flow uniqueness]
\label{thm:curvature-spacetime-trichotomy}
Let $M^3$ be nonempty, connected, and without boundary, and let $g(t)$,
$t\in(0,T)$, be a complete Ricci flow with nonnegative curvature operator.
Assume
\begin{equation}
  \sup_{M\times[a,b]}|\Rm(g(t))|<\infty
       \qquad\text{whenever }0<a<b<T.                    \label{eq:smp-compact-slab-bound}
\end{equation}
Assume as delegated analytic inputs the following two standard theorems;
neither is proved in this chapter.
\begin{enumerate}
\item A complete bounded-curvature metric has a complete
      bounded-curvature Ricci flow for a short time.
\item If two complete Ricci flows on the same manifold have uniformly
      bounded curvature on $[a,b]$ and agree at $t=a$, then they agree on
      $[a,b]$.
\end{enumerate}
Then exactly one of the following alternatives holds for every
$t\in(0,T)$.
\begin{enumerate}
\item \emph{Stationary flat.}  Every $g(t)$ is flat, and after identifying
      times by the fixed underlying manifold, $g(t)$ is independent of $t$.
\item \emph{One product for the whole flow.}  There are a fixed complete
      simply connected surface $\Sigma$, a fixed diffeomorphism
      \[
                       F\colon\Sigma\times\R\longrightarrow\widetilde M,
      \]
      and a complete surface Ricci flow $h(t)$ of positive Gaussian
      curvature such that
      \begin{equation}
            F^*\widetilde g(t)=h(t)+ds^2
            \qquad\text{for every }t\in(0,T).             \label{eq:smp-whole-flow-product}
      \end{equation}
\item \emph{Strictly positive.}
      $\mathcal R_{g(t)}>0$ at every spacetime point.
\end{enumerate}
Here $\widetilde g(t)$ denotes the lift to the universal cover.  The
complete bounded-curvature existence and uniqueness theorems displayed
above are explicit dependencies of this whole-flow upgrade.  They are not
consequences of the strong maximum principle, and they are not used in the
local or fixed-time theorems.
\end{theorem}

The exact project-source locators for the two delegated analytic inputs are
MSM~144, Appendix~D, Theorem~D.1 for the compact case and
Theorems~D.2--D.3 for complete noncompact metrics with bounded curvature,
all on p.~357.  Theorem~D.2 records Shi's theorem; Theorem~D.3 records the
Chen--Zhu theorem and points also to Hsu.  Their stable nodes are,
respectively,
\nodeid{PC-F-COMPLETE-BOUNDED-CURVATURE-RICCI-FLOW-SHORT-TIME-EXISTENCE}
and
\nodeid{PC-F-COMPLETE-RICCI-FLOW-UNIQUENESS}.

\begin{proof}
Fix $s\in(0,T)$.  If $g(s)$ is flat, the stationary family with initial
value $g(s)$ is a complete bounded-curvature Ricci flow.  Uniqueness makes
$g(t)$ stationary and flat for $t\geq s$.

Suppose instead that the rank-one alternative holds at $s$.  By
Theorem~\ref{thm:curvature-complete-trichotomy}, choose an isometry
\[
       F_s\colon(\Sigma_s\times\R,h_s+dr^2)
                    \longrightarrow(\widetilde M,\widetilde g(s)).
\]
The curvature bound at time $s$ pulls back to a curvature bound for the
complete surface $(\Sigma_s,h_s)$.
Run the complete bounded-curvature Ricci flow $h_s(t)$ from $h_s$ and form
the product flow $h_s(t)+dr^2$.  Pulled to $\widetilde M$ by $F_s$, it has
the same initial value as $\widetilde g(t)$.  The uniqueness input makes the
two flows equal for a short forward time.  For completeness, let $b<T$ be
a putative first endpoint of this equality.  In the fixed coordinates
$F_s$, the identities ``mixed terms vanish'' and ``the $dr^2$ coefficient
is one'' pass to $t=b$ by smoothness of $\widetilde g(t)$.  Thus
$F_s^*\widetilde g(b)=h_b+dr^2$.  This product is complete, and
\eqref{eq:smp-compact-slab-bound} bounds its curvature.  Short-time
existence from $h_b$ and uniqueness at $b$ extend the equality beyond $b$,
a contradiction.  Hence the product persists for every $t\geq s$.
Thus a flat slice cannot later become
rank one or positive, and a rank-one slice cannot later become positive.
Strict positivity, once present, persists by the rank-spreading conclusion
of Theorem~\ref{thm:smp-systems}.  Comparing any two positive times therefore
shows that one of the three ranks $0,1,3$ occurs for the entire interval.

It remains to justify the word ``one'' in the product alternative.  Once
rank one is known on the whole interval, fix an arbitrary reference time
$t_*\in(0,T)$, choose a rank-three metric bundle $(V,h)$ and an isometry
$\iota_{t_*}\colon(V,h)\to(TM,g(t_*))$, and solve
\eqref{eq:ue-ricci-iota-ode} both forward and backward from $t_*$.  This
produces one Uhlenbeck gauge on all of $(0,T)$; no initial time $0$ is being
silently used.  Corollary
\ref{cor:curvature-fixed-bundle-rank-interval} and Proposition
\ref{prop:rank-one-uhlenbeck-line-fixed} apply on $(0,T)$.  Pull the fixed
Uhlenbeck line to the universal cover.  Since the cover is simply connected,
it has a global $h$-unit section $v$; because the line is $D^t$-parallel,
every such unit section is $D^t$-parallel.  Proposition
\ref{prop:rank-one-uhlenbeck-line-fixed} then makes
$X=\iota_tv$ one time-independent vector field which is unit and parallel
for every $\widetilde g(t)$.  In intrinsic terms its span is
$L=\ker\Ric^\sharp(\widetilde g(t))$ at every time.  Moreover,
\begin{equation}
 \partial_t\bigl(\widetilde g(t)(X,\cdot)\bigr)
     =-2\Ric_{\widetilde g(t)}(X,\cdot)=0.                \label{eq:smp-fixed-line-one-form}
\end{equation}
Thus the parallel one-form $\alpha=\widetilde g(t)(X,\cdot)$ is independent
of time.  The potential $r$ with $dr=\alpha$, its zero level $\Sigma$, and
the flow of $X$ used in
Theorem~\ref{thm:parallel-line-splitting} can all be chosen once and
for all.  They give the fixed map $F$ in
\eqref{eq:smp-whole-flow-product}.  Restricting
$\partial_t\widetilde g=-2\Ric_{\widetilde g}$ to $T\Sigma$ gives
\[
                              \partial_t h=-2\Ric_h,
\]
while the $ds^2$ coefficient is stationary.  Positivity of $K_h$ follows
from \eqref{eq:smp-rank-one-sectional}.  Completeness of $h(t)$ follows at
each time from completeness of the isometric product
$(\Sigma\times\R,h(t)+ds^2)$, by the factor argument in the proof of
Theorem~\ref{thm:parallel-line-splitting}.
\end{proof}

\begin{remark}[Closed flows]
\label{rem:curvature-whole-flow-closed}
For a Ricci flow on a closed manifold,
\eqref{eq:smp-compact-slab-bound}, short-time existence, and uniqueness hold
on every compact time slab.  Hence
Theorem~\ref{thm:curvature-spacetime-trichotomy} applies automatically.
For incomplete flows it need not apply; Example~6.59 of GSM~77 is the
model obstruction.  Completeness without the curvature bounds needed by the
chosen uniqueness theorem should not silently be substituted for the
stated hypothesis.
\end{remark}

\begin{corollary}[Closed ancient three-dimensional trichotomy]
\label{cor:closed-ancient-curvature-trichotomy}
Let $M^3$ be nonempty, connected, and closed, and let
$g(t)$, $t\in(-\infty,T)$, be an ancient Ricci flow.  Then precisely one of
the following holds.
\begin{enumerate}
\item The flow is stationary and flat.
\item The curvature operator is positive at every spacetime point.
\item There are a fixed diffeomorphism
      $F\colon\Sph^2\times\R\to\widetilde M$ and an ancient surface Ricci
      flow $h(t)$ on $\Sph^2$, with positive Gaussian curvature, such that
      \begin{equation}
                    F^*\widetilde g(t)=h(t)+ds^2
                    \qquad(t<T).                         \label{eq:smp-ancient-product}
      \end{equation}
\end{enumerate}
The surface flow in the third alternative is not asserted to be round.
\end{corollary}

\begin{proof}
Corollary~\ref{cor:hi-ancient-nonnegative} gives
$\mathcal R_{g(t)}\geq0$ at every spacetime point.  On each compact time
slab all curvature bounds and uniqueness hypotheses in
Theorem~\ref{thm:curvature-spacetime-trichotomy} are automatic.  Given any
two finite times, apply that theorem after translating a slightly larger
finite interval; hence the alternative cannot depend on time.  In the
rank-one case, choose one reference time $t_*$, solve the Uhlenbeck ODE in
both directions on $(-\infty,T)$, and repeat the fixed-line and fixed
one-form construction in the proof of
Theorem~\ref{thm:curvature-spacetime-trichotomy}.  This produces one
product map for the entire ancient interval, rather than unrelated product
maps on overlapping slabs.

It remains only to identify the surface.  At any fixed time in the rank-one
case, the positive curvature eigenvalue on the compact base has a positive
minimum.  The lifted factor therefore has Gaussian curvature bounded below
by a positive constant.  Bonnet--Myers makes it compact, and simple
connectivity identifies it topologically with $\Sph^2$.
\end{proof}

\begin{remark}[The surface factor need not be round]
\label{rem:ancient-surface-not-round}
The product of a nonround Rosenau ancient solution on $\Sph^2$ with a flat
circle is a closed ancient three-dimensional Ricci flow whose universal
cover has the form \eqref{eq:smp-ancient-product}.  Thus roundness does not
follow from the strong maximum principle.  It requires an additional input,
such as an appropriate noncollapsing, Type~I, or soliton-classification
hypothesis.
For a shrinking soliton, Chapter~\ref{ch:three-dimensional-shrinkers}
supplies precisely that additional rigidity: the soliton equation splits
the potential, and the resulting complete surface shrinker is round.
\end{remark}

\section{Formalization interfaces and source record}
\label{sec:curvature-splitting-formalization-record}

The following stable semantic nodes are the intended interfaces for the
formalization blueprint.  All have status ``complete human proof; Lean
implementation pending,'' relative to the explicitly delegated interfaces
for the local scalar Dirichlet/strong maximum principle and the complete
bounded-curvature Ricci-flow existence and uniqueness theorems used only in
the spacetime upgrade.
\begin{sloppypar}
\begin{enumerate}
\item \nodeid{PC-F-CURVATURE-REACTION-PSD-COMMUTES}: Lemma
      \ref{lem:curvature-reaction-positive-commuting}.
\item \nodeid{PC-F-CLOSED-NONNEGATIVE-CURVATURE-PRESERVATION}: Corollary
      \ref{cor:closed-curvature-operator-preserved}, depending on the
      Chapter~2 weak systems theorem and the Uhlenbeck evolution.
\item \nodeid{PC-F-CURVATURE-NULL-REACTION-RANKS}: Proposition
      \ref{prop:curvature-kernel-annihilation}, equation
      \eqref{eq:smp-null-reaction}, and Lemma
      \ref{lem:curvature-rank-two-excluded}.
\item \nodeid{PC-F-CURVATURE-FIXED-BUNDLE-RANK-INTERVAL}: Corollary
      \ref{cor:curvature-fixed-bundle-rank-interval}.
\item \nodeid{PC-F-CURVATURE-ONE-SIDED-ENDPOINT-TRANSFER}: Lemma
      \ref{lem:curvature-one-sided-endpoint}, which closes the right
      endpoint of the left rank-stable slab used at a prescribed positive
      time.
\item \nodeid{PC-F-CURVATURE-TIME-SLICE-TRICHOTOMY}: Theorem
      \ref{thm:curvature-time-slice-trichotomy}.
\item \nodeid{PC-B-CURVATURE-IMAGE-TANGENT-LINE}: Lemmas
      \ref{lem:curvature-image-to-parallel-line} and
      \ref{lem:rank-one-normalization}.  This is the orientation-free bridge
      from the fixed-bundle curvature kernel to tangent geometry.
\item \nodeid{PC-B-UHLENBECK-RANK-ONE-PHYSICAL-LINE}: Proposition
      \ref{prop:rank-one-uhlenbeck-line-fixed}, including the
      time-independence of the physical line, its dual one-form, and the
      local product coordinates on a rank-stable interval.
\item \nodeid{PC-F-PARALLEL-LINE-LOCAL-PRODUCT}: Proposition
      \ref{prop:parallel-line-local-product}.
\item \nodeid{PC-F-PARALLEL-LINE-GLOBAL-SPLITTING}: Theorem
      \ref{thm:parallel-line-splitting}.
\item \nodeid{PC-F-CURVATURE-FIXED-TIME-UNIVERSAL-COVER}: Theorem
      \ref{thm:curvature-complete-trichotomy} and Corollary
      \ref{cor:curvature-zero-eigenvalue-splitting}.
\item \nodeid{PC-F-CURVATURE-WHOLE-FLOW-TRICHOTOMY}: Theorem
      \ref{thm:curvature-spacetime-trichotomy}, whose interface depends on
      \nodeid{PC-F-COMPLETE-BOUNDED-CURVATURE-RICCI-FLOW-SHORT-TIME-EXISTENCE}
      and \nodeid{PC-F-COMPLETE-RICCI-FLOW-UNIQUENESS}.
\item \nodeid{PC-F-ANCIENT-THREE-DIMENSIONAL-SPLITTING}: Corollary
      \ref{cor:closed-ancient-curvature-trichotomy}.
\end{enumerate}
\end{sloppypar}

The two complete-flow nodes just named are delegated classical interfaces,
not results proved in this chapter.  Their exact source locators are
MSM~144, Appendix~D, Theorems~D.1--D.3, p.~357; Theorem~D.1 covers the
compact case and Theorems~D.2--D.3 the complete noncompact
bounded-curvature case.

The dependency chain should remain visible in Lean:
\[
\begin{aligned}
 &\text{fixed-bundle curvature evolution}
   \longrightarrow\text{strong systems theorem}\\
 &\hspace{1.5em}\longrightarrow
   \text{parallel null bundle and null reaction}\\
 &\hspace{1.5em}\longrightarrow\{\text{ranks }0,1,3\}
   \longrightarrow\text{orientation-free tangent line}\\
 &\hspace{1.5em}\longrightarrow\text{local product}.
\end{aligned}
\]
followed, only in the complete case, by the explicit global splitting
theorem.  The whole-flow node has the additional, logically separate,
existence-and-uniqueness dependency.

The primary historical source is R.~S.~Hamilton,
\emph{Four-manifolds with positive curvature operator}, Journal of
Differential Geometry \textbf{24} (1986), 153--179,
\href{https://doi.org/10.4310/jdg/1214440433}{DOI
10.4310/jdg/1214440433}.  The null-space argument and strong maximum
principle are in Section~8, especially Lemma~8.2 and Theorem~8.3,
pp.~174--176; the holonomy and splitting consequences are developed in
Section~9, pp.~176--178.

Among the supplied project sources, the closest treatments are GSM~77,
Chapter~6, Sections~6--7, pp.~245--250 (especially Example~6.59,
Theorem~6.60, Proposition~6.62, and Theorem~6.64); MSM~144,
Propositions~12.47 and~12.49, pp.~184--188, Theorem~12.50,
pp.~188--192, and Theorem~12.53 with Remark~12.54, pp.~193--194; and
MSM~135, Appendix~A, Theorems~A.53--A.54, p.~469.  MSM~110,
Chapter~4, Section~4, pp.~102--103, is a useful check of the three possible
eigenvalue patterns.  GSM~235, Proposition~5.5 and Theorem~5.7,
pp.~162--163, gives a recent statement of the image-subbundle result and its
specialization to complete three-dimensional gradient Ricci solitons.

The proof above reconstructs the sources faithfully.  It makes four
intentional changes for formalization.  First, rank two is excluded directly from
$Q(A)$ and the null-reaction condition, avoiding a classification of
subalgebras of $\mathfrak{so}(3)$.  Second, the tangent line is the
contraction annihilator \eqref{eq:smp-line-annihilator}, avoiding a Hodge
star and any orientation hypothesis.  Third, the required global splitting
is proved from one parallel line with the explicit maps
\eqref{eq:smp-global-product-isometry} and
\eqref{eq:smp-global-product-inverse}, rather than imported as the full de
Rham theorem.  Fourth, the fixed-time theorem is kept separate from the
whole-flow theorem, and complete-flow uniqueness is displayed as an
additional dependency.  Complete-flow short-time existence is displayed
beside it.  These changes strengthen traceability without
changing the mathematical content of Hamilton's argument.

\section{Combined formalization crosswalk and proof status}
\label{sec:strong-curvature-formalization}

This section records the intended source-to-formalization interfaces for
Chapters~\ref{ch:curvature-evolution} and
\ref{ch:strong-curvature}.  The identifiers below are stable semantic
crosswalk keys: they do not depend on chapter numbers or future Lean
declaration names.  They record mathematical meanings and dependencies;
they do not assert that a separate blueprint ledger or any Lean source has
already been created.

\begin{sloppypar}
\subsection*{Macro crosswalk}

\begin{enumerate}
\item \nodeid{PC-M-UHLENBECK-CURVATURE-EVOLUTION} is realized by
      Theorem~\ref{thm:ue-operator-evolution}.  Its project-facing output is
      the exact fixed-bundle equation used earlier in
      Proposition~\ref{prop:hi-uhlenbeck-evolution}.  The latter is an
      earlier consumer-facing occurrence of the same interface, not a
      second semantic node.

\item \nodeid{PC-M-STRONG-MAXIMUM-PRINCIPLE-SYSTEMS} is realized by the
      rank engine, Proposition~\ref{prop:smp-rank-spreading}, and the
      precise abstract strong principle,
      Theorem~\ref{thm:smp-systems}.

\item \nodeid{PC-M-CURVATURE-OPERATOR-RANK-RIGIDITY} is realized by
      Theorem~\ref{thm:curvature-time-slice-trichotomy}.  This is where the
      abstract system theorem and the three-dimensional reaction
      polynomial are combined to exclude curvature-operator rank two at a
      positive time.

\item \nodeid{PC-M-PARALLEL-LINE-SPLITTING} is realized by
      Theorem~\ref{thm:parallel-line-splitting}.

\item \nodeid{PC-M-THREE-DIMENSIONAL-CURVATURE-TRICHOTOMY} has two public
      outputs.  The fixed-time complete result is
      Theorem~\ref{thm:curvature-complete-trichotomy}; the spacetime result,
      with its additional persistence input, is
      Theorem~\ref{thm:curvature-spacetime-trichotomy}.
\end{enumerate}

\subsection*{Fine dependencies in proof order}

\begin{enumerate}
\item \nodeid{PC-B-CURVATURE-TENSOR-OPERATOR-NORMALIZATION} is the
      normalization bridge recorded by
      Definition~\ref{def:ue-fixed-curvature-operator} and
      Lemma~\ref{lem:ue-three-dimensional-reaction}, and used in
      Theorem~\ref{thm:ue-operator-evolution}.  With
      \[
        \mathbf R(X,Y,Z,W)=g(R(X,Y)Z,W),
      \]
      the operator convention is
      \[
        \langle\mathcal R(X\wedge Y),Z\wedge W\rangle
        =2\mathbf R(X,Y,W,Z).
      \]
      Thus a metric of constant sectional curvature $k$ has
      $\mathcal R=2kI$, and in dimension three
      $\operatorname{tr}\mathcal R=R$.  This bridge fixes the order of the
      final two tensor slots, the factor two, the tangent-versus-cotangent
      musical identification, and the Laplacian sign.  These conventions
      are part of the theorem interface, not proof-local notation.

\item \nodeid{PC-F-RICCI-FLOW-COVARIANT-CURVATURE-EVOLUTION} is the exact
      evolution equation for the covariant $(0,4)$ curvature tensor,
      derived in Proposition~\ref{prop:ue-raw-curvature-evolution} from the
      connection-variation, Bianchi, commutator, and $B$-tensor lemmas.
      Starting from it,
      Lemma~\ref{lem:ue-four-slot-derivative} proves
      \nodeid{PC-F-TIME-DERIVATIVE-OF-BUNDLE-PULLBACK};
      Lemma~\ref{lem:ue-spatial-naturality} proves
      \nodeid{PC-F-PULLBACK-ROUGH-LAPLACIAN-NATURALITY}; and
      Proposition~\ref{prop:ue-pulled-tensor-evolution} proves
      \nodeid{PC-F-UHLENBECK-RICCI-TERM-CANCELLATION}.  Together with the
      operator conversion, these declarations yield
      \nodeid{PC-F-UHLENBECK-CURVATURE-EVOLUTION}:
      \[
        \partial_t\mathcal M
        =\Delta_{D^t,g(t)}\mathcal M
         +\mathcal M^2+\mathcal M^\#.
      \]
      The conversion
      \nodeid{PC-B-GENERAL-CURVATURE-REACTION-TO-THREE-DIMENSIONAL-ADJUGATE}
      identifies
      \[
        \mathcal M^\#
        =\mathcal M^2-(\operatorname{tr}\mathcal M)\mathcal M
          +\sigma_2(\mathcal M)I
      \]
      in dimension three.  The Uhlenbeck cancellation and the
      tensor-to-operator conversion and the unpulled raw curvature evolution
      are proved here.

\item \nodeid{PC-F-LOCAL-SCALAR-STRONG-PARABOLIC-MAXIMUM-PRINCIPLE} is the
      delegated scalar-PDE input recorded in
      Proposition~\ref{prop:smp-local-scalar-input}.  The Ky Fan lower
      spectral sum \nodeid{PC-F-KY-FAN-LOWER-SPECTRAL-SUM} is
      Definition~\ref{def:smp-ky-fan}.  Using these two inputs,
      Proposition~\ref{prop:smp-rank-spreading} proves both
      \nodeid{PC-F-KY-FAN-STRICT-POSITIVITY-PROPAGATION} and
      \nodeid{PC-F-POSITIVE-SYSTEM-RANK-SPREADING}:
      \[
        \inf_x\operatorname{rank}A(x,t)
        \geq
        \sup_x\operatorname{rank}A(x,s),
        \qquad s<t.
      \]
      The finite-dimensional Ky Fan variational formula and all subsequent
      bundle arguments are proved in the chapter; only the named scalar
      strong principle is delegated.

\item Proposition~\ref{prop:smp-positive-time-rank} proves
      \nodeid{PC-F-POSITIVE-TIME-SPATIAL-RANK-CONSTANCY}, and
      Lemma~\ref{lem:smp-smooth-kernel} proves
      \nodeid{PC-F-CONSTANT-RANK-KERNEL-SUBBUNDLE}.
      Proposition~\ref{prop:smp-kernel-rigidity} proves
      \nodeid{PC-F-POSITIVE-SYSTEM-KERNEL-SPATIAL-PARALLEL}; these results
      are packaged in Theorem~\ref{thm:smp-systems}.  The bridge
      \nodeid{PC-B-NULL-EIGENVECTOR-TIME-CONSTANCY-CORRECTION} records an
      essential distinction: the null-eigenvector inequality alone gives
      rank propagation and spatial parallelity, whereas constancy of the
      kernel in the fixed bundle time variable requires the additional
      reaction hypothesis stated in that theorem; the necessity of the
      distinction is exhibited in Example~\ref{ex:smp-rotating-kernel}.
      Under the additional hypothesis,
      Corollary~\ref{cor:smp-kernel-time-constant} proves
      \nodeid{PC-F-REACTION-KERNEL-ANNIHILATION} and
      \nodeid{PC-F-KERNEL-TIME-CONSTANCY}.  No unrecorded spacetime
      connection is used.

\item Lemma~\ref{lem:curvature-reaction-positive-commuting} proves
      \nodeid{PC-F-CURVATURE-REACTION-PSD-COMMUTES}.  Proposition
      \ref{prop:curvature-kernel-annihilation}, Lemma
      \ref{lem:curvature-one-sided-endpoint}, and Lemma
      \ref{lem:curvature-rank-two-excluded} together prove
      \nodeid{PC-F-CURVATURE-NULL-REACTION-RANKS}, with the endpoint
      passage separately named
      \nodeid{PC-F-CURVATURE-ONE-SIDED-ENDPOINT-TRANSFER}.
      Theorem~\ref{thm:curvature-time-slice-trichotomy} packages these as
      \nodeid{PC-F-CURVATURE-TIME-SLICE-TRICHOTOMY}.  If
      $A=\operatorname{diag}(a,b,0)$ with $a,b>0$, then
      \[
          Q(A)=\operatorname{diag}(a^2,b^2,ab),
      \]
      so the null direction has strictly positive reaction.  This excludes
      rank two without formalizing the classification of Lie subalgebras of
      $\mathfrak{so}(3)$.

\item Lemmas~\ref{lem:rank-one-normalization} and
      \ref{lem:curvature-image-to-parallel-line} prove the algebraic bridge
      \nodeid{PC-B-CURVATURE-IMAGE-TANGENT-LINE}, using
      \[
        h(\operatorname{Ric}_A v,w)
        =\frac12\sum_i
          \langle A(v\wedge e_i),w\wedge e_i\rangle
      \]
      to turn a parallel rank-one curvature image into a parallel tangent
      line without choosing an orientation or a global Hodge star.
      Proposition~\ref{prop:parallel-line-local-product} proves
      \nodeid{PC-F-PARALLEL-LINE-LOCAL-PRODUCT}, and Theorem
      \ref{thm:parallel-line-splitting} proves
      \nodeid{PC-F-PARALLEL-LINE-GLOBAL-SPLITTING}.
      Theorem~\ref{thm:curvature-complete-trichotomy} and Corollary
      \ref{cor:curvature-zero-eigenvalue-splitting} then prove
      \nodeid{PC-F-CURVATURE-FIXED-TIME-UNIVERSAL-COVER}.

\item The delegated nodes
      \nodeid{PC-F-COMPLETE-BOUNDED-CURVATURE-RICCI-FLOW-SHORT-TIME-EXISTENCE}
      and \nodeid{PC-F-COMPLETE-RICCI-FLOW-UNIQUENESS} are used only for
      the persistence step in
      Theorem~\ref{thm:curvature-spacetime-trichotomy}.  Conditional on
      them, that theorem proves
      \nodeid{PC-F-CURVATURE-WHOLE-FLOW-TRICHOTOMY}.  Neither the
      fixed-time trichotomy nor Theorem~\ref{thm:parallel-line-splitting}
      depends on either complete-flow input.
\end{enumerate}
\end{sloppypar}

\subsection*{Status ledger}

The raw covariant curvature evolution, normalization bridge, Uhlenbeck pullback cancellation,
three-dimensional reaction conversion, Ky Fan and rank arguments,
kernel-rigidity argument, rank-two exclusion, parallel-line
splitting, and fixed-time curvature trichotomy are supplied with human
proofs in Chapters~\ref{ch:curvature-evolution} and
\ref{ch:strong-curvature}.  Their precise newly delegated dependency
boundary consists of three named inputs: the local scalar Dirichlet
existence/strong-maximum-principle package; and, only for the spacetime
persistence theorem, complete bounded-curvature short-time Ricci-flow
existence and complete-flow uniqueness.  The
Chapter~2 weak systems theorem and the Chapter~3 ancient Hamilton--Ivey
corollary are earlier proved dependencies within this book, not additional
unrecorded inputs.

All identifiers in this section currently have Lean statement and Lean
proof status \emph{not yet implemented}.  No typechecking, compilation,
kernel checking, axiom-footprint inspection, or semantic correspondence
audit is claimed.  The existence of complete human proofs or delegated
classical inputs does not by itself upgrade any Lean status.

\chapter{Gradient Ricci Solitons and Canonical Self-Similar Flows}
\label{ch:gradient-solitons}
\index{gradient Ricci soliton}
\index{Ricci soliton!gradient}
\index{Ricci flow!self-similar solution}
\index{scalar curvature!quadratic lower bound}

\section{Purpose and main results}
\label{sec:sol-purpose}

Gradient Ricci solitons are the stationary objects of Ricci flow after one
allows both diffeomorphisms and constant rescalings.  This chapter fixes the
normalization and makes that slogan into an exact global construction.  It
then proves two arbitrary-dimensional shrinker theorems needed in the
analysis of singularity models.

The first is the Chow--Lu--Yang theorem: a complete noncompact shrinker
which is not Gaussian has scalar curvature bounded below by a positive
multiple of the inverse square of distance.  After scaling the soliton
constant to $\sigma=1$ and Hamilton-normalizing the potential, the proof is
organized around the stronger potential estimate
\[
                 R\geq c\bigl(f^{-1}+nf^{-2}\bigr).
\]
Properness of $f$ turns the apparently noncompact maximum-principle argument
into an ordinary minimum argument on a compact sublevel set.

The second is the Munteanu--Wang theorem in its stronger original form: a
complete shrinker with nonnegative sectional curvature and positive Ricci
curvature is compact.  Its tensor barrier is the exact Ricci-tensor analogue
of the Chow--Lu--Yang scalar barrier.  We then use a corrected orbit
dichotomy and an elementary scalar-curvature identity on regular
$f$-sublevels.  The user-facing positive-sectional-curvature statement is a
corollary.  No bounded-curvature or noncollapsing hypothesis is used.

Unless a statement says otherwise, $M^n$ in this chapter is a nonempty,
connected, smooth $n$-manifold without boundary.  The dimension $n$ is a
positive integer.  All integrals use Riemannian densities, so orientability
is never implicit.

\begin{roadmap}{Sections~\ref{sec:sol-definition}--\ref{sec:sol-identities}
fix the soliton operations, Hamilton normalization, and weighted identities.}
Sections~\ref{sec:sol-completeness} and \ref{sec:sol-canonical} construct the
global canonical Ricci flow, and Section~\ref{sec:sol-models} records its
model checks.  Section~\ref{sec:sol-shrinker-inputs} isolates
the standard complete-shrinker inputs.  Sections~\ref{sec:sol-cly} and
\ref{sec:sol-mw} prove the Chow--Lu--Yang and Munteanu--Wang theorems.
The final sections give regression tests, the formalization dependency
order, and the source-repair ledger.
\end{roadmap}

\section{Definition, covariance, and normalization}
\label{sec:sol-definition}

\begin{definition}[Gradient Ricci soliton]
\label{def:gradient-ricci-soliton}
Let $g$ be a Riemannian metric on $M^n$, let $f\in C^\infty(M)$, and let
$\sigma\in\mathbb R$.  The quadruple $(M,g,f,\sigma)$ is a
\emph{gradient Ricci soliton} if
\begin{equation}
          \Ric_g+\nabla_g^2f=\frac{\sigma}{2}g.            \label{eq:sol-equation}
\end{equation}
It is called \emph{shrinking}, \emph{steady}, or \emph{expanding} according
as $\sigma>0$, $\sigma=0$, or $\sigma<0$.
\end{definition}

\textcolor{red}{Since the above is a major definition, put the Lean code for the statement here.}

We use $\sigma$ for the soliton constant.  The symbols
$\lambda,\mu,\nu$ already denote curvature-operator eigenvalues in the
preceding chapters.  The factor
$1/2$ makes a normalized shrinker satisfy
$\Ric+\nabla^2f=\frac12g$.

\begin{definition}[Weighted Laplacian]
\label{def:sol-weighted-laplacian}
For a smooth function $u$ on $(M,g)$, define
\begin{equation}
       \Delta_f u
          =\Delta_gu-\langle\nabla^gf,\nabla^gu\rangle_g. \label{eq:sol-weighted-laplacian}
\end{equation}
The same symbol on a covariant tensor means the rough Laplacian minus
covariant differentiation in the direction $\nabla f$.
\end{definition}

\begin{proposition}[Functorial operations]
\label{prop:sol-operations}
Let $(M,g,f,\sigma)$ satisfy \eqref{eq:sol-equation}.
\begin{enumerate}
\item If $c\in\mathbb R$, then $(M,g,f+c,\sigma)$ is a soliton.
\item If $\Phi:N\to M$ is a diffeomorphism, then
      $(N,\Phi^*g,f\circ\Phi,\sigma)$ is a soliton.
\item If $a>0$, then $(M,ag,f,\sigma/a)$ is a soliton.
\item If $(M_i,g_i,f_i,\sigma)$, $i=1,2$, are solitons with the same
      constant, then
      \[
       \bigl(M_1\times M_2,
          g_1\oplus g_2,
          f_1\circ\pi_1+f_2\circ\pi_2,
          \sigma\bigr)
      \]
      is a soliton.
\end{enumerate}
\end{proposition}

\begin{proof}
A constant shift does not change the Hessian.  Ricci curvature and Hessians
are natural under diffeomorphisms.  Under the constant scaling $g\mapsto ag$
the Levi--Civita connection, the Hessian as a covariant two-tensor, and the
Ricci tensor as a covariant two-tensor are unchanged; the right side becomes
$(\sigma/(2a))(ag)$.  The product statement follows from the direct-sum
formulas for Ricci curvature and Hessians.  The equality of the two soliton
constants in the product statement is essential.
\end{proof}

For later scaling arguments we record
\begin{equation}
 R_{ag}=a^{-1}R_g,\qquad
 |\nabla f|_{ag}^2=a^{-1}|\nabla f|_g^2,\qquad
 d_{ag}=a^{1/2}d_g.                                      \label{eq:sol-scaling-data}
\end{equation}
In particular, replacing $g$ by $\widehat g=\sigma g$ normalizes the
\emph{soliton constant} of a shrinker from $\sigma>0$ to $1$; Hamilton
normalization of the potential is a separate constant shift.

\section{Scalar and weighted soliton identities}
\label{sec:sol-identities}

\begin{theorem}[Fundamental scalar identities]
\label{thm:sol-fundamental-identities}
Let $(M^n,g,f,\sigma)$ be a gradient Ricci soliton.  Then
\begin{align}
 R+\Delta f&=\frac{n\sigma}{2},                            \label{eq:sol-trace}\\
 dR&=2\Ric(\nabla f,\mathord\cdot),                        \label{eq:sol-grad-scalar}\\
 R+|\nabla f|^2-\sigma f&=C                                \label{eq:sol-hamilton-constant}
\end{align}
for one constant $C\in\mathbb R$.  Moreover,
\begin{align}
 \Delta_ff&=\frac{n\sigma}{2}-\sigma f-C,                 \label{eq:sol-drift-potential}\\
 \Delta_fR&=\sigma R-2|\Ric|^2.                           \label{eq:sol-drift-scalar}
\end{align}
\end{theorem}

\begin{proof}
Tracing \eqref{eq:sol-equation} gives \eqref{eq:sol-trace}.  Take its
divergence.  The contracted Bianchi identity and the commutation identity
\[
 \operatorname{div}(\nabla^2f)
      =d(\Delta f)+\Ric(\nabla f,\mathord\cdot)
\]
give
\[
 0=\tfrac12dR+d(\Delta f)+\Ric(\nabla f,\mathord\cdot).
\]
Differentiating \eqref{eq:sol-trace} gives $d(\Delta f)=-dR$, and hence
\eqref{eq:sol-grad-scalar}.

The soliton equation and \eqref{eq:sol-grad-scalar} imply
\[
 d|\nabla f|^2
   =2\nabla^2f(\nabla f,\mathord\cdot)
   =\sigma\,df-2\Ric(\nabla f,\mathord\cdot)
   =\sigma\,df-dR.
\]
Thus the differential of the left side of
\eqref{eq:sol-hamilton-constant} vanishes.  Connectedness of $M$ makes it
one constant $C$.  Combining \eqref{eq:sol-trace} with
\eqref{eq:sol-hamilton-constant} proves \eqref{eq:sol-drift-potential}.

Finally, take the divergence of \eqref{eq:sol-grad-scalar}:
\begin{align*}
 \Delta R
 &=2\nabla^i(R_{ij}\nabla^jf)\\
 &=\langle\nabla R,\nabla f\rangle
      +2\langle\Ric,\nabla^2f\rangle\\
 &=\langle\nabla R,\nabla f\rangle
      +\sigma R-2|\Ric|^2.
\end{align*}
Subtracting the drift term proves \eqref{eq:sol-drift-scalar}.
\end{proof}

\begin{definition}[Hamilton normalization]
\label{def:sol-hamilton-normalization}
If $\sigma\ne0$, a soliton potential is \emph{Hamilton normalized} when
\begin{equation}
                  R+|\nabla f|^2=\sigma f.                \label{eq:sol-hamilton-normalized}
\end{equation}
Equivalently, the constant $C$ in
\eqref{eq:sol-hamilton-constant} is zero.
\end{definition}

Every potential with $\sigma\ne0$ has a unique constant shift which is
Hamilton normalized: replace $f$ by $f+C/\sigma$.  This normalization is
not the same as choosing the weighted measure $e^{-f}d\mu$ to have a
prescribed mass.  Under a shift $f\mapsto f+c$ the Hamilton constant changes
from $C$ to $C-\sigma c$; under $g\mapsto ag$ it changes to $C/a$.

\begin{proposition}[Weighted integration by parts]
\label{prop:sol-weighted-integration}
For $u,v\in C_c^\infty(M)$,
\begin{equation}
 \int_Mu\,\Delta_fv\,e^{-f}d\mu_g
   =-\int_M\langle\nabla u,\nabla v\rangle e^{-f}d\mu_g
   =\int_Mv\,\Delta_fu\,e^{-f}d\mu_g.                   \label{eq:sol-weighted-integration}
\end{equation}
\end{proposition}

\begin{proof}
The identity
$\operatorname{div}(e^{-f}u\nabla v)
=e^{-f}(u\Delta_fv+\langle\nabla u,\nabla v\rangle)$ and the compact-support
divergence theorem give the first equality.  Interchanging $u$ and $v$
gives the second.  Completeness alone would not justify omitting the compact
support.
\end{proof}

\section{Completeness of the soliton vector field}
\label{sec:sol-completeness}

The global self-similar flow requires the vector field $\nabla f$ to be
complete.  Metric completeness alone does not imply completeness of an
arbitrary smooth vector field, so we isolate the argument.

\begin{lemma}[Linear-growth vector fields are complete]
\label{lem:sol-linear-growth-complete}
Let $(M,g)$ be complete, let $X$ be a smooth vector field, and suppose that
for some $p\in M$ and $A\geq0$,
\begin{equation}
             |X|_g(x)\leq A\bigl(1+d_g(p,x)\bigr)
             \qquad(x\in M).                              \label{eq:sol-linear-growth-field}
\end{equation}
Then every maximal integral curve of $X$ is defined on all of $\mathbb R$.
\end{lemma}

\begin{proof}
Let $\gamma:(a,b)\to M$ be a maximal integral curve with $0\in(a,b)$ and
put $r(t)=d(p,\gamma(t))$.  For $0\leq t<b$, the triangle inequality and
\eqref{eq:sol-linear-growth-field} give
\[
 1+r(t)\leq1+r(0)+A\int_0^t(1+r(s))\,ds.
\]
Gr\"onwall's inequality bounds $r(t)$, and hence $|X|(\gamma(t))$, on every
finite time interval.  If $b<\infty$, the tail of $\gamma$ has length
tending to zero as its initial time tends to $b$.  Thus $\gamma(t)$ is
Cauchy as $t\uparrow b$.  Metric completeness supplies a limit, and local
ODE existence extends the curve past $b$, a contradiction.  Applying the
same argument to $-X$ shows $a=-\infty$.
\end{proof}

We use the following standard complete-soliton scalar estimate as a named
analytic input.

\begin{theorem}[Sharp scalar lower bound for complete solitons]
\label{thm:sol-sharp-scalar-lower}
If $(M^n,g,f,\sigma)$ is a complete gradient Ricci soliton, then
\begin{equation}
 R\geq L_\sigma,\qquad
 L_\sigma=
 \begin{cases}
 0,&\sigma\geq0,\\
 n\sigma/2,&\sigma<0.
 \end{cases}                                             \label{eq:sol-sharp-scalar-lower}
\end{equation}
If $\sigma>0$ and $R$ vanishes at one point, then the shrinker is Gaussian
in the sense of Definition~\ref{def:sol-gaussian} below.
\end{theorem}

The scalar estimate is the only non-elementary input in the next proof.  It
is proved by a cutoff maximum principle applied to
\eqref{eq:sol-drift-scalar}; its equality case uses the soliton equation and
complete flat-shrinker rigidity.

\begin{theorem}[Completeness of the potential field]
\label{thm:sol-potential-field-complete}
If $(M,g,f,\sigma)$ is a complete gradient Ricci soliton, then $\nabla^gf$
is a complete vector field.
\end{theorem}

\begin{proof}
Let $C$ be the Hamilton constant and let $L_\sigma$ be
\eqref{eq:sol-sharp-scalar-lower}.  Define
\[
                    H=\sigma f+C-L_\sigma.
\]
By Hamilton's identity and the scalar lower bound,
\[
       H=R-L_\sigma+|\nabla f|^2\geq|\nabla f|^2\geq0.
\]
If $\sigma\ne0$, then $\nabla H=\sigma\nabla f$, and for every
$\varepsilon>0$,
\[
 |\nabla\sqrt{H+\varepsilon}|
   =\frac{|\sigma|\,|\nabla f|}{2\sqrt{H+\varepsilon}}
   \leq\frac{|\sigma|}{2}.
\]
Integrating along a minimizing geodesic from a fixed point $p$ and letting
$\varepsilon\downarrow0$ shows
\[
 |\nabla f|(x)\leq\sqrt{H(x)}
     \leq\sqrt{H(p)}+\frac{|\sigma|}{2}d(p,x).
\]
If $\sigma=0$, the same identity gives the uniform bound
$|\nabla f|^2\leq C-L_0$.  Lemma~\ref{lem:sol-linear-growth-complete}
applies in both cases.
\end{proof}

\section{The canonical self-similar Ricci flow}
\label{sec:sol-canonical}

Let $(M,g,f,\sigma)$ be a complete gradient Ricci soliton.  Define
\begin{equation}
 \tau(t)=1-\sigma t,\qquad
 I_\sigma=\{t\in\mathbb R:\tau(t)>0\}.                   \label{eq:sol-canonical-domain}
\end{equation}
Thus
\[
 I_\sigma=
 \begin{cases}
 (-\infty,1/\sigma),&\sigma>0,\\
 \mathbb R,&\sigma=0,\\
 (1/\sigma,\infty),&\sigma<0.
 \end{cases}
\]
The set description in \eqref{eq:sol-canonical-domain}, rather than the
case split, is the preferred formal interface.

Let $\Psi_s$ be the global flow of $\nabla^gf$, supplied by
Theorem~\ref{thm:sol-potential-field-complete}, and set
\begin{equation}
 s_\sigma(t)=
 \begin{cases}
 -\sigma^{-1}\log(1-\sigma t),&\sigma\ne0,\\
 t,&\sigma=0.
 \end{cases}                                             \label{eq:sol-canonical-time-change}
\end{equation}
Then $s_\sigma(0)=0$ and $s_\sigma'(t)=\tau(t)^{-1}$.  Define
\begin{equation}
 \phi_t=\Psi_{s_\sigma(t)},\qquad
 g(t)=\tau(t)\phi_t^*g,\qquad
 f(t)=f\circ\phi_t.                                      \label{eq:sol-canonical-definition}
\end{equation}

\begin{theorem}[Canonical self-similar flow]
\label{thm:sol-canonical-flow}
Let $(M,g,f,\sigma)$ be a complete gradient Ricci soliton.  The families in
\eqref{eq:sol-canonical-definition} are defined for every $t\in I_\sigma$,
$\phi_t$ is a diffeomorphism, and
\begin{align}
 \phi_0&=\operatorname{id}_M,\qquad g(0)=g,\qquad f(0)=f,  \label{eq:sol-canonical-initial}\\
 \partial_t\phi_t(x)
   &=\tau(t)^{-1}\nabla^gf(\phi_t(x)),                  \label{eq:sol-canonical-diffeomorphism}\\
 \partial_tg(t)&=-2\Ric_{g(t)},                           \label{eq:sol-canonical-ricci-flow}\\
 \Ric_{g(t)}+\nabla_{g(t)}^2f(t)
   &=\frac{\sigma}{2\tau(t)}g(t),                        \label{eq:sol-canonical-soliton}\\
 \partial_tf(t)&=|\nabla^{g(t)}f(t)|_{g(t)}^2.            \label{eq:sol-canonical-potential-time}
\end{align}
Every $g(t)$ is complete.  If $C$ is the Hamilton constant of $(g,f)$,
then
\begin{equation}
 R_{g(t)}+|\nabla f(t)|_{g(t)}^2
   -\frac{\sigma}{\tau(t)}f(t)=\frac{C}{\tau(t)}.         \label{eq:sol-canonical-hamilton-constant}
\end{equation}
\end{theorem}

\textcolor{red}{Since the above is a major theorem, put the Lean code for the statement here.}

\begin{proof}
The first two equations follow from the definition and the time-change
identity.  For a time-dependent pullback,
\[
 \frac d{dt}\phi_t^*g
   =\phi_t^*\mathcal L_{\tau^{-1}\nabla f}g.
\]
Since $\mathcal L_{\nabla f}g=2\nabla^2f$, the soliton equation gives
\begin{align*}
 \partial_tg(t)
 &=-\sigma\phi_t^*g
      +\tau\phi_t^*\mathcal L_{\tau^{-1}\nabla f}g\\
 &=\phi_t^*(-\sigma g+2\nabla^2f)
  =-2\phi_t^*\Ric_g
  =-2\Ric_{g(t)}.
\end{align*}
The last equality uses diffeomorphism naturality and the fact that the
covariant Ricci tensor is unchanged by a positive constant scaling.

The same two facts give
\[
 \Ric_{g(t)}+\nabla_{g(t)}^2f(t)
 =\phi_t^*(\Ric_g+\nabla_g^2f)
 =\frac\sigma2\phi_t^*g
 =\frac{\sigma}{2\tau}g(t).
\]
Moreover,
\[
 \partial_tf(t)
 =\tau^{-1}|\nabla^gf|_g^2\circ\phi_t
 =|\nabla^{g(t)}f(t)|_{g(t)}^2.
\]
A diffeomorphic pullback and a positive constant scaling preserve
completeness.  Finally, scalar curvature and squared gradient norm each
acquire the factor $\tau^{-1}$ under
$g(t)=\tau\phi_t^*g$, proving
\eqref{eq:sol-canonical-hamilton-constant}.
\end{proof}

\begin{remark}[Representation domain versus maximal lifetime]
\label{rem:sol-canonical-domain}
The theorem asserts that the canonical representation is defined on
$I_\sigma$; it does not assert that this is the maximal Ricci-flow interval.
For the Gaussian shrinker the canonical metric is static and extends beyond
$t=1/\sigma$, even though the particular scale--diffeomorphism expression
has $\tau(t)\downarrow0$ there.
\end{remark}

\begin{proposition}[Tensor scaling along the canonical flow]
\label{prop:sol-canonical-tensor-scaling}
Let $(M^n,g,f,\sigma)$ be a complete gradient Ricci soliton, and let
$\tau$, $I_\sigma$, $\phi_t$, and $g(t)$ be the objects defined in
\eqref{eq:sol-canonical-domain} and
\eqref{eq:sol-canonical-definition}.  For every $t\in I_\sigma$,
\begin{align}
 \Ric_{g(t)}&=\phi_t^*\Ric_g,                             \label{eq:sol-canonical-ricci-scaling}\\
 \mathbf R_{g(t)}&=\tau(t)\phi_t^*\mathbf R_g,           \label{eq:sol-canonical-rm-scaling}\\
 R_{g(t)}&=\tau(t)^{-1}R_g\circ\phi_t,                   \label{eq:sol-canonical-scalar-scaling}\\
 d\mu_{g(t)}&=\tau(t)^{n/2}\phi_t^*d\mu_g.               \label{eq:sol-canonical-volume-scaling}
\end{align}
Consequently sectional-curvature signs, Ricci positivity, and Ricci nullity
are preserved by the canonical identifications.
\end{proposition}

\begin{proof}
These are the standard pullback laws followed by the constant-scaling laws.
The covariant Riemann tensor acquires one factor of $\tau$, whereas each
contraction by the inverse metric removes one such factor.
\end{proof}

\begin{proposition}[Conjugate-heat bridge for a normalized shrinker]
\label{prop:sol-conjugate-heat}
Let $(M^n,g,f,1)$ be a complete gradient Ricci soliton whose potential is
Hamilton normalized.  Let $\tau$, $\phi_t$, $g(t)$, and $f(t)$ be defined by
\eqref{eq:sol-canonical-domain} and
\eqref{eq:sol-canonical-definition}.  For $t\in I_1=(-\infty,1)$ put
\[
 u(x,t)=(4\pi\tau(t))^{-n/2}e^{-f(t)(x)}.
\]
Then
\begin{equation}
       \partial_tu=-\Delta_{g(t)}u+R_{g(t)}u,             \label{eq:sol-conjugate-heat}
\end{equation}
and
\begin{equation}
 u(t)d\mu_{g(t)}
   =(4\pi)^{-n/2}\phi_t^*(e^{-f}d\mu_g).                 \label{eq:sol-conjugate-measure}
\end{equation}
\end{proposition}

\begin{proof}
The trace of \eqref{eq:sol-canonical-soliton} gives
$R+\Delta f(t)=n/(2\tau)$.  Combine this with
$\partial_tf(t)=|\nabla f(t)|^2$ and
$\Delta(e^{-f})=e^{-f}(-\Delta f+|\nabla f|^2)$ to obtain
\eqref{eq:sol-conjugate-heat}.  Equation
\eqref{eq:sol-conjugate-measure} follows immediately from
\eqref{eq:sol-canonical-definition} and
\eqref{eq:sol-canonical-volume-scaling}.
\end{proof}

\begin{proposition}[Weighted Ricci equation]
\label{prop:sol-weighted-ricci}
For every gradient Ricci soliton,
\begin{equation}
 (\Delta_f\Ric)_{ij}
   =\sigma R_{ij}-2R_{kij\ell}R^{k\ell}.                 \label{eq:sol-drift-ricci}
\end{equation}
Here the contraction uses the covariant curvature convention of
Definition~\ref{def:hi-curvature-tensor}.
\end{proposition}

\begin{proof}
The calculation is local.  Fix $x\in M$.  Local ODE theory supplies a
neighborhood $U$ of $x$ and an $\varepsilon>0$ on which the flow map of
$\nabla f$ is defined for $|s|<\varepsilon$.  The construction in
\eqref{eq:sol-canonical-definition}, restricted to this spacetime
neighborhood, suffices to differentiate at $(0,x)$ even when the original
soliton is not complete.  Constant scaling does not change the covariant
Ricci tensor; hence the canonical expression gives
\[
       \left.\partial_t\Ric_{g(t)}\right|_{t=0}
          =\mathcal L_{\nabla f}\Ric.
\]
On the other hand, contraction of the raw curvature evolution from
Proposition~\ref{prop:ue-raw-curvature-evolution} gives the Ricci-flow
identity
\[
 \partial_tR_{ij}
   =\Delta R_{ij}+2R_{kij\ell}R^{k\ell}-2R_i{}^kR_{kj}.
\]
Expanding the Lie derivative and using
$\nabla^2f=(\sigma/2)g-\Ric$ gives
\[
 (\mathcal L_{\nabla f}\Ric)_{ij}
   =\nabla_{\nabla f}R_{ij}+\sigma R_{ij}-2R_i{}^kR_{kj}.
\]
Equate the last two displays and cancel the common quadratic Ricci term.
This is \eqref{eq:sol-drift-ricci}.
\end{proof}

\section{Models and normalization regression tests}
\label{sec:sol-models}

The following examples test every scaling and additive constant used later.

\begin{example}[Gaussian shrinker]
\label{ex:sol-gaussian}
On $\mathbb R^n$ with its Euclidean metric,
\[
               f(x)=\frac{|x|^2}{4},\qquad \sigma=1
\]
is Hamilton normalized.  Its canonical diffeomorphism is
$\phi_t(x)=\tau(t)^{-1/2}x$, so the factors cancel and $g(t)=g_{\rm E}$.
The potential is $f(t)(x)=|x|^2/(4\tau(t))$.
\end{example}

\begin{definition}[Gaussian shrinker]
\label{def:sol-gaussian}
A complete shrinking gradient Ricci soliton $(M,g,f,\sigma)$ is
\emph{Gaussian} if there exist a diffeomorphism
$\Psi:M\to\mathbb R^n$ and a constant $b\in\mathbb R$ such that
\begin{equation}
       \Psi^*g_{\rm E}=\sigma g,
       \qquad
       f+b=\left(\frac{|\mathord\cdot|^2}{4}\right)\circ\Psi. \label{eq:sol-gaussian-witnesses}
\end{equation}
Thus $\Psi:(M,\sigma g)\to(\mathbb R^n,g_{\rm E})$ is the isometry, and the
second equality records the potential rather than merely the metric.
\end{definition}

Allowing the isometry to include a Euclidean translation absorbs the affine
freedom in a Gaussian potential.  Thus Gaussianity is independent of the
chosen additive normalization of $f$.

\begin{example}[Einstein shrinkers and the round sphere]
If $\Ric=(\sigma/2)g$, a constant potential solves the soliton equation.
In Hamilton normalization that constant is $f=n/2$, not $0$.  In
particular, when $n\geq2$ the round sphere
$S^n(\sqrt{2(n-1)})$ with $\sigma=1$ has $f=n/2$ and canonical metric
$g(t)=(1-t)g$.
\end{example}

\begin{example}[Round cylinders]
For $2\leq k\leq n-1$,
\[
 S^k(\sqrt{2(k-1)})\times\mathbb R^{n-k},\qquad
 f(y)=\frac{k}{2}+\frac{|y|^2}{4},\qquad \sigma=1,
\]
is Hamilton normalized.  Its canonical metric is
\[
       g(t)=(1-t)g_{S^k}+g_{\mathbb R^{n-k}}.
\]
The constant $k/2$ is forced by Hamilton normalization.  The cylinder has
nonnegative sectional curvature but a nontrivial Ricci kernel; it shows
that the positive-Ricci hypothesis in Theorem~\ref{thm:sol-mw} is essential.
It is also the universal-cover model for the rank-one branch in
Chapter~\ref{ch:three-dimensional-shrinkers}.
\end{example}

\section{The complete-shrinker input package}
\label{sec:sol-shrinker-inputs}

\begin{definition}[Normalized gradient shrinking soliton]
\label{def:sol-normalized-shrinker}
A triple $(M^n,g,f)$ is a \emph{normalized shrinker} if $g$ is complete,
$f\in C^\infty(M)$, and
\begin{equation}
             \Ric+\nabla^2f=\frac12g,\qquad
             R+|\nabla f|^2=f.                           \label{eq:sol-normalized-shrinker}
\end{equation}
The standing convention at the start of this chapter supplies that $M$ is
nonempty, connected, smooth, and without boundary, with $n\geq1$.
\end{definition}

From this point through the two proof engines we use this normalization.
Thus $\Delta_ff=n/2-f$ and $\Delta_fR=R-2|\Ric|^2$.

\begin{theorem}[Potential control on a complete shrinker]
\label{thm:sol-potential-control}
Let $(M^n,g,f)$ be a normalized shrinker.  Then $R\geq0$.  The fixed
Hamilton-normalized potential $f$ is proper and attains a minimum at some
point $o$.  It satisfies
\begin{equation}
 \frac14\bigl((d(o,x)-5n)_+\bigr)^2
   \leq f(x)
   \leq\frac14\bigl(d(o,x)+\sqrt{2n}\bigr)^2.             \label{eq:sol-potential-growth}
\end{equation}
Moreover $0\leq f(o)=R(o)\leq n/2$.  If the shrinker is not Gaussian, then
\begin{equation}
                         R(x)>0                           \label{eq:sol-positive-scalar}
\end{equation}
for every $x\in M$; in particular, $f>0$ everywhere.
\end{theorem}

The scalar nonnegativity and Gaussian equality case are the shrinking part
of Theorem~\ref{thm:sol-sharp-scalar-lower}.  Properness and quadratic
growth originate in the Cao--Zhou potential theorem; the explicit universal
constants $5n$ and $\sqrt{2n}$ in
\eqref{eq:sol-potential-growth} are the sharpened Haslhofer--M\"uller form.
These are explicit imported inputs to the next two
proofs.  The upper bound can also be recovered directly without dividing at
a possible zero of $f$: for every $\varepsilon>0$,
\[
 |\nabla(2\sqrt{f+\varepsilon})|^2
   =\frac{|\nabla f|^2}{f+\varepsilon}
   \leq\frac{f}{f+\varepsilon}\leq1.
\]
Here $|\nabla f|^2\leq f$ follows from $R\geq0$ and
\eqref{eq:sol-normalized-shrinker}.  Integrate along a minimizing geodesic
from $o$, use $f(o)\leq n/2$, and let $\varepsilon\downarrow0$.

\begin{lemma}[Reciprocal-potential formulas]
\label{lem:sol-reciprocal-formulas}
On the positive set of a normalized shrinker,
\begin{align}
 \Delta_f(f^{-1})
   &=f^{-1}-\frac n2f^{-2}+2|\nabla f|^2f^{-3},           \label{eq:sol-drift-f-inverse}\\
 \Delta_f(f^{-2})
   &=2f^{-2}-nf^{-3}+6|\nabla f|^2f^{-4}.               \label{eq:sol-drift-f-inverse-square}
\end{align}
Consequently, for
\begin{equation}
                  H=f^{-1}+nf^{-2},                      \label{eq:sol-reciprocal-H}
\end{equation}
one has
\begin{equation}
                  \Delta_fH\geq H
                  \qquad\text{wherever }f\geq2n.         \label{eq:sol-H-supersolution}
\end{equation}
\end{lemma}

\begin{proof}
For a smooth one-variable function $\eta$,
\[
 \Delta_f(\eta\circ f)
   =\eta'(f)\Delta_ff+\eta''(f)|\nabla f|^2.
\]
Apply this with $\eta(s)=s^{-1}$ and $s^{-2}$ and use
$\Delta_ff=n/2-f$.  Adding the first formula to $n$ times the second gives
\[
 \Delta_fH-H
   =nf^{-3}\left(\frac f2-n\right)
      +2|\nabla f|^2f^{-3}+6n|\nabla f|^2f^{-4},
\]
which proves \eqref{eq:sol-H-supersolution}.
\end{proof}

\section{The Chow--Lu--Yang scalar lower bound}
\label{sec:sol-cly}

We first prove the potential form, which is stronger and has the cleanest
formal interface.

\begin{theorem}[Reciprocal-potential scalar barrier]
\label{thm:sol-cly-potential}
Let $(M^n,g,f)$ be a complete noncompact non-Gaussian normalized shrinker.
Then there is a constant $c>0$ such that, for every $x\in M$,
\begin{equation}
 R(x)\geq c\bigl(f(x)^{-1}+nf(x)^{-2}\bigr)
       \geq\frac{c}{f(x)}.                               \label{eq:sol-cly-potential-lower}
\end{equation}
\end{theorem}

\begin{proof}
Theorem~\ref{thm:sol-potential-control} gives $R>0$, $f>0$, and properness
of $f$.  Choose $F>2n$ for which the nonempty sublevel
\[
                         K_F=\{x:f(x)\leq F\}
\]
is compact.  Since $R/H$ is positive and continuous on $K_F$, choose
$c>0$ so small that
\begin{equation}
                 \Phi:=R-cH>0\quad\text{on }K_F.          \label{eq:sol-cly-core-positive}
\end{equation}
Using \eqref{eq:sol-drift-scalar} with $\sigma=1$ and
Lemma~\ref{lem:sol-reciprocal-formulas}, an exact calculation gives
\begin{align}
 \Delta_f\Phi
 &=\Phi-2|\Ric|^2
   -cnf^{-3}\left(\frac f2-n\right)                      \label{eq:sol-cly-barrier-identity}\\
 &\quad
   -cf^{-4}(2f+6n)|\nabla f|^2.\notag
\end{align}

Suppose that $\Phi(y)<0$ at some point.  Choose
$T>\max\{F,f(y)\}$ so large that
\[
              c(T^{-1}+nT^{-2})<-\tfrac12\Phi(y).
\]
If $f(x)>T$, scalar nonnegativity gives
\[
 \Phi(x)\geq-c(f(x)^{-1}+nf(x)^{-2})
           >\tfrac12\Phi(y)>\Phi(y).
\]
Therefore $\Phi$ attains a negative global minimum at a point $x_0$ in the
compact set $\{f\leq T\}$.  By
\eqref{eq:sol-cly-core-positive}, $f(x_0)>F>2n$.  At $x_0$,
$\nabla\Phi=0$ and $\Delta_f\Phi=\Delta\Phi\geq0$.  But
\eqref{eq:sol-cly-barrier-identity} gives
\[
 0\leq\Delta_f\Phi(x_0)
  \leq\Phi(x_0)
    -cnf(x_0)^{-3}\left(\frac{f(x_0)}2-n\right)<0,
\]
a contradiction.  Hence $\Phi\geq0$ everywhere, proving
\eqref{eq:sol-cly-potential-lower}.
\end{proof}

\begin{theorem}[Chow--Lu--Yang scalar lower bound]
\label{thm:sol-cly-distance}
Let $(M^n,g,f,\sigma)$ be a complete noncompact non-Gaussian gradient
shrinking Ricci soliton, with arbitrary $\sigma>0$.  For every basepoint
$p\in M$ there exists a constant $c_p>0$ such that
\begin{equation}
             R_g(x)\geq
             \frac{c_p}{(1+d_g(p,x))^2}
             \qquad\text{for every }x\in M.              \label{eq:sol-cly-distance-lower}
\end{equation}
The constant may depend on the soliton, $\sigma$, and $p$, but not on $x$.
\end{theorem}

\textcolor{red}{Since the above is a major theorem, put the Lean code for the statement here.}

\begin{proof}
First take $\sigma=1$ and Hamilton normalize $f$.  From
$R+|\nabla f|^2=f$ and $R\geq0$,
\[
                    |\nabla(2\sqrt f)|\leq1.
\]
Hopf--Rinow supplies a minimizing geodesic from $p$ to $x$, and integration
along it gives
\[
             f(x)\leq\frac14
                 \bigl(d(p,x)+2\sqrt{f(p)}\bigr)^2.
\]
Together with Theorem~\ref{thm:sol-cly-potential},
\[
 R(x)\geq
 \frac{4c}{(d(p,x)+2\sqrt{f(p)})^2}
 \geq\frac{c_p}{(1+d(p,x))^2}
\]
for a suitable $c_p>0$.

For general $\sigma>0$, let $C$ be the Hamilton constant in
\eqref{eq:sol-hamilton-constant} and set
\[
              \widehat g=\sigma g,
              \qquad \widehat f=f+\frac C\sigma.
\]
Then $(M,\widehat g,\widehat f)$ is a normalized shrinker.  This operation
preserves completeness, noncompactness, and Gaussianity, and satisfies
$R_{\widehat g}=\sigma^{-1}R_g$ and
$d_{\widehat g}=\sqrt{\sigma}\,d_g$.  Apply the normalized result and absorb
the fixed scaling factors into $c_p$.
\end{proof}

\begin{remark}[Sharp order and exact source statement]
\label{rem:sol-cly-sharp}
The original theorem is stated outside a compact set as
$R(x)d(x,p)^2\geq C^{-1}$ for sufficiently large $d(x,p)$.  Positivity on
the compact core makes it equivalent to
\eqref{eq:sol-cly-distance-lower}.  The inverse-quadratic order is attained
by known noncompact K\"ahler shrinkers and cannot in general be improved.
\end{remark}

\begin{proposition}[CLY bridges]
\label{prop:sol-cly-source-bridges}
Let $(M^n,g,f)$ be a normalized shrinker.
\begin{enumerate}
\item The metric $g$ is flat if and only if $(M,g,f,1)$ is Gaussian.
\item Suppose $(M,g,f,1)$ is not Gaussian.  Fix $p\in M$.  If there are
      $A>0$ and $r_0\geq0$ such that
      \[
       R(x)\geq\frac{A}{(1+d(p,x))^2}
       \qquad\text{whenever }d(p,x)\geq r_0,
      \]
      then there is $c_p>0$ for which the same inequality, with $A$
      replaced by $c_p$, holds for every $x\in M$.
\end{enumerate}
\end{proposition}

\begin{proof}
A Gaussian shrinker is flat.  Conversely, if $g$ is flat, then $R$ vanishes
at every point, so the equality case in
Theorem~\ref{thm:sol-sharp-scalar-lower} supplies the Gaussian witnesses.
For the second assertion, scalar positivity gives
\[
             Q(x):=R(x)(1+d(p,x))^2>0.
\]
Hopf--Rinow makes the closed ball $\overline{B(p,r_0)}$ compact.  Let $m>0$
be the minimum of $Q$ there and take $c_p=\min\{A,m\}$.
\end{proof}

\section{The Munteanu--Wang compactness theorem}
\label{sec:sol-mw}

For linearly independent $v,w\in T_xM$, our sectional-curvature convention
is
\begin{equation}
 K(v,w)=
 \frac{\mathbf R_g(v,w,w,v)}{|v|^2|w|^2-\langle v,w\rangle^2}.          \label{eq:sol-sectional-curvature}
\end{equation}
The hypothesis $K\geq0$ below quantifies over all tangent two-planes.  In
dimension at least four it is weaker than nonnegativity of the curvature
operator, and the two predicates must not be identified in a formalization.

\begin{definition}[Curvature--Ricci contraction]
\label{def:sol-curvature-ricci-contraction}
At a point choose an orthonormal Ricci eigenbasis
$e_1,\ldots,e_n$ with $\Ric(e_a)=\rho_ae_a$.  Define the symmetric
two-tensor $\mathcal C_{\Ric}$ by
\begin{equation}
 \mathcal C_{\Ric}(v,v)
   =\sum_{a=1}^n\rho_a\,\mathbf R_g(e_a,v,v,e_a).          \label{eq:sol-curvature-ricci-contraction}
\end{equation}
Equivalently, in indices,
$(\mathcal C_{\Ric})_{ij}=R_{kij\ell}R^{k\ell}$.
\end{definition}

The coordinate expression shows that the definition is independent of the
chosen eigenbasis, including inside repeated eigenspaces.

\begin{lemma}[Positivity of the curvature--Ricci contraction]
\label{lem:sol-curvature-ricci-positive}
If $K\geq0$, then $\Ric\geq0$ and
$\mathcal C_{\Ric}\geq0$ as symmetric bilinear forms.
\end{lemma}

\begin{proof}
For a unit vector $v$, extend it to an orthonormal basis.  Then
$\Ric(v,v)$ is the sum of the sectional curvatures of the planes containing
$v$, so it is nonnegative.  Hence every Ricci eigenvalue $\rho_a$ is
nonnegative.  Fix $a$.  If $v$ and $e_a$ are linearly dependent, then
$\mathbf R_g(e_a,v,v,e_a)=0$ by the alternating symmetry of
$\mathbf R_g$ in its first
two arguments.  If they are linearly independent, then the $a$th
summand in \eqref{eq:sol-curvature-ricci-contraction} is
\[
 \rho_a K(v,e_a)
       \bigl(|v|^2-\langle v,e_a\rangle^2\bigr)\geq0.
\]
Thus every summand is nonnegative, and so is their sum.
\end{proof}

\begin{theorem}[Reciprocal-potential Ricci barrier]
\label{thm:sol-mw-ricci-barrier}
Let $(M^n,g,f)$ be a complete noncompact normalized shrinker with
$K\geq0$ and pointwise positive-definite Ricci curvature.  Then there is
$c>0$ such that
\begin{equation}
        \Ric\geq c(f^{-1}+nf^{-2})g\geq cf^{-1}g           \label{eq:sol-mw-ricci-lower}
\end{equation}
on all of $M$.
\end{theorem}

\begin{proof}
Positive Ricci curvature implies $R>0$, hence $f>0$.  Properness follows
from Theorem~\ref{thm:sol-potential-control}.  By
Proposition~\ref{prop:sol-weighted-ricci},
Lemma~\ref{lem:sol-curvature-ricci-positive}, and $\sigma=1$,
\begin{equation}
                 \Delta_f\Ric
                    =\Ric-2\mathcal C_{\Ric}\leq\Ric.     \label{eq:sol-mw-ricci-subsolution}
\end{equation}
Let $H=f^{-1}+nf^{-2}$.  Choose $F>2n$ so that $\{f\leq F\}$ is nonempty.
It is compact, and positive definiteness of $\Ric$ gives $c>0$ for which
\begin{equation}
                 T:=\Ric-cHg\geq0
                 \quad\text{on }\{f\leq F\}.             \label{eq:sol-mw-core-tensor}
\end{equation}
Equations \eqref{eq:sol-H-supersolution} and
\eqref{eq:sol-mw-ricci-subsolution} imply
\begin{equation}
                 \Delta_fT\leq T
                 \quad\text{wherever }f>F.               \label{eq:sol-mw-tensor-inequality}
\end{equation}

Suppose $T$ has a negative direction.  Since $\Ric>0$ and $H\to0$ at
infinity, the least Rayleigh quotient of $T$ has nonnegative lower limit at
infinity.  More explicitly, after fixing one negative value, choose a
compact $f$-sublevel outside which $-cH$ is larger than half that value;
minimize $T_x(v,v)$ over the unit sphere bundle of this compact sublevel.
There is therefore a pair $(x_0,v_0)$, with $|v_0|=1$ and $f(x_0)>F$, at
which the global minimum is a negative number $\kappa$.

Vertical minimization makes $v_0$ an eigenvector of $T_{x_0}$ with
$T(v_0)=\kappa v_0$.  Choose a local unit vector field $V$ with
$V(x_0)=v_0$ and $\nabla V(x_0)=0$.  The scalar function $T(V,V)$ has a
local minimum at $x_0$.  When its drift Laplacian is expanded, all
first-derivative terms vanish.  The remaining $\Delta_fV$ terms vanish
after pairing with $T(v_0)=\kappa v_0$, because differentiating
$|V|^2=1$ twice and using $\nabla V(x_0)=0$ gives
$\langle\Delta_fV,v_0\rangle=0$.  Hence
\[
 0\leq\Delta_f[T(V,V)](x_0)
    =(\Delta_fT)(v_0,v_0)
    \leq T(v_0,v_0)=\kappa<0,
\]
a contradiction.  Thus $T\geq0$ everywhere, proving
\eqref{eq:sol-mw-ricci-lower}.
\end{proof}

\begin{lemma}[Orbit dichotomy]
\label{lem:sol-mw-orbit-dichotomy}
Under the hypotheses of Theorem~\ref{thm:sol-mw-ricci-barrier}, there is a
constant $a>0$ such that
\begin{equation}
                  R(x)\geq\min\{n,af(x)\}
                  \qquad(x\in M).                        \label{eq:sol-mw-min-dichotomy}
\end{equation}
\end{lemma}

\begin{proof}
Let $c>0$ be the constant in \eqref{eq:sol-mw-ricci-lower}.  The vector
field $\nabla f$ is complete by
Theorem~\ref{thm:sol-potential-field-complete}.  For $x\in M$, let
$\beta:\mathbb R\to M$ solve
\[
              \beta'(s)=\nabla f(\beta(s)),\qquad\beta(0)=x.
\]
Along this curve, \eqref{eq:sol-grad-scalar} and
\eqref{eq:sol-mw-ricci-lower} give
\begin{equation}
 \frac d{ds}R(\beta(s))
  =2\Ric(\nabla f,\nabla f)
  \geq2c\frac{|\nabla f|^2}{f}
  =2c\left(1-\frac Rf\right).                            \label{eq:sol-mw-orbit-R}
\end{equation}

There are two exhaustive cases on the fixed interval $[-n/c,0]$.  If
$R/f\leq1/2$ throughout the interval, then
$dR/ds\geq c$, and therefore
\[
              R(x)\geq n+R(\beta(-n/c))>n.
\]
Otherwise, at some $s\in[-n/c,0]$, with $y=\beta(s)$,
$R(y)>f(y)/2$.  Scalar curvature is nondecreasing along the forward
$\nabla f$-flow, so $R(x)\geq R(y)$.  Also
\[
 \frac d{ds}\log f(\beta(s))
   =\frac{|\nabla f|^2}{f}=1-\frac Rf\leq1.
\]
Thus $f(y)\geq e^{-n/c}f(x)$ and
\[
                  R(x)>\tfrac12e^{-n/c}f(x).
\]
Taking $a=\tfrac12e^{-n/c}$ proves
\eqref{eq:sol-mw-min-dichotomy}.  The minimum, rather than a maximum, is
the conclusion common to the two alternatives.
\end{proof}

\begin{lemma}[Scalar average on regular potential sublevels]
\label{lem:sol-sublevel-scalar-average}
Let $(M^n,g,f)$ be a complete normalized shrinker with $f$ proper.  If
$t$ is a regular value of $f$ and
$\Omega_t=\{x\in M:f(x)<t\}$, then
\begin{equation}
 \int_{\Omega_t}R\,d\mu
  =\frac n2\Vol(\Omega_t)
      -\int_{\partial\Omega_t}|\nabla f|\,d\sigma
  \leq\frac n2\Vol(\Omega_t).                            \label{eq:sol-sublevel-average}
\end{equation}
\end{lemma}

\begin{proof}
Properness makes the closure of $\Omega_t$ compact.  The trace identity is
$R+\Delta f=n/2$.  Integrate it over $\Omega_t$ and apply the divergence
theorem.  Because $t$ is regular, the outward unit normal is
$\nabla f/|\nabla f|$ and $\partial_\nu f=|\nabla f|$.  This gives the
equality and hence the inequality.  The density formulation of the
divergence theorem does not require an orientation.
\end{proof}

\begin{lemma}[Infinite volume under nonnegative Ricci curvature]
\label{lem:sol-nonnegative-ricci-infinite-volume}
If $(M^n,g)$ is complete, connected, noncompact, and has $\Ric\geq0$, then
$\Vol_g(M)=\infty$.
\end{lemma}

\begin{proof}
Hopf--Rinow supplies a unit-speed ray $\gamma:[0,\infty)\to M$ from a
point $p$.  Put $d_j=2^j$, $q_j=\gamma(d_j)$, and $r_j=d_j/4$.  The balls
$B(q_j,r_j)$ are pairwise disjoint.  Moreover
$B(p,1)\subset B(q_j,d_j+1)$.  Bishop--Gromov comparison at $q_j$ gives
\begin{align*}
 \Vol B(q_j,r_j)
 &\geq\left(\frac{r_j}{d_j+1}\right)^n
          \Vol B(q_j,d_j+1)\\
 &\geq8^{-n}\Vol B(p,1)>0.
\end{align*}
Infinitely many disjoint subsets with this uniform positive volume force
$\Vol M=\infty$.
\end{proof}

\begin{theorem}[Munteanu--Wang compactness theorem]
\label{thm:sol-mw}
Let $(M^n,g,f,\sigma)$ be a complete gradient shrinking Ricci soliton with
$\sigma>0$.  If every sectional curvature is nonnegative and $\Ric_g$ is
positive definite at every point, then $M$ is compact.
\end{theorem}

\textcolor{red}{Since the above is a major theorem, put the Lean code for the statement here.}

\begin{proof}
Suppose for contradiction that $M$ is noncompact.  Scale to
$\widehat g=\sigma g$ and Hamilton normalize the potential.  Completeness,
noncompactness, both curvature-sign assumptions, and compactness are
unchanged, so it suffices to work under
\eqref{eq:sol-normalized-shrinker}.

Lemma~\ref{lem:sol-mw-orbit-dichotomy} gives
$R\geq\min\{n,af\}$ for some $a>0$.  Let
\[
                    K=\{x:f(x)\leq n/a\}.
\]
This set is compact, and $R\geq n$ on $M\setminus K$.  By
Lemma~\ref{lem:sol-curvature-ricci-positive}, $\Ric\geq0$; hence
Lemma~\ref{lem:sol-nonnegative-ricci-infinite-volume} gives
$\Vol M=\infty$.  Properness of $f$ means that its sublevels exhaust $M$, so
their volumes tend to infinity.  Sard's theorem supplies arbitrarily large
regular values.  Choose a regular $t>n/a$ such that
\[
                   \Vol(\Omega_t)>2\Vol(K).
\]
Since $R>0$ on $K$ and $R\geq n$ outside $K$,
\[
 \int_{\Omega_t}R\,d\mu
  \geq n\bigl(\Vol(\Omega_t)-\Vol(K)\bigr)
  >\frac n2\Vol(\Omega_t).
\]
This contradicts Lemma~\ref{lem:sol-sublevel-scalar-average}.
Therefore $M$ is compact.
\end{proof}

\begin{corollary}[Positive sectional curvature forces compactness]
\label{cor:sol-positive-sectional-compact}
Let $n\geq2$, and let $(M^n,g,f,\sigma)$ be a complete gradient shrinking
Ricci soliton.  If every sectional curvature is positive, then $M$ is
compact.
\end{corollary}

\begin{proof}
Positive sectional curvature implies nonnegative sectional curvature.  For
a nonzero $v$, extend $v/|v|$ to an orthonormal basis
$e_1,\ldots,e_n$.  Since $n\geq2$,
\[
 \Ric(v,v)=|v|^2\sum_{a=2}^nK(e_1,e_a)>0.
\]
Thus Theorem~\ref{thm:sol-mw} applies.  The explicit restriction $n\geq2$
prevents a vacuous interpretation of positivity on two-planes in dimension
one.
\end{proof}

\begin{proposition}[Converse slice of a self-similar flow]
\label{prop:sol-converse-self-similar}
Let $I\subset\mathbb R$ be an open interval, let $\bar g$ be a Riemannian
metric on $M$, and let $(g(t))_{t\in I}$ be a smooth family of Riemannian
metrics satisfying $\partial_tg(t)=-2\Ric_{g(t)}$.  Suppose there are a
smooth function $c:I\to(0,\infty)$ and a smooth
map $\Phi:I\times M\to M$ such that every
$\phi_t:=\Phi(t,\mathord\cdot)$ is a diffeomorphism and
\[
                    g(t)=c(t)\phi_t^*\bar g,
\]
for every $t\in I$.  For $t\in I$ define $V_t\in\Gamma(TM)$ by the
pointwise witness-tying equation
\begin{equation}
 d\phi_t|_x\bigl(V_t(x)\bigr)=\partial_t\Phi(t,x)
       \qquad(x\in M).                                   \label{eq:sol-converse-generator}
\end{equation}
Then, for every $t\in I$, the time slice satisfies
\begin{equation}
 \Ric_{g(t)}+\frac12\mathcal L_{V_t}g(t)
    =-\frac{c'(t)}{2c(t)}g(t).                            \label{eq:sol-converse-soliton}
\end{equation}
If $V_t=\nabla^{g(t)}F_t$ for some $F_t\in C^\infty(M)$, then the slice is
a gradient Ricci soliton with soliton constant $-c'(t)/c(t)$.
\end{proposition}

\begin{proof}
Differentiate the displayed self-similar expression in domain variables:
\[
 \partial_tg(t)=\frac{c'(t)}{c(t)}g(t)+\mathcal L_{V_t}g(t).
\]
Equate this with $-2\Ric_{g(t)}$ and rearrange.
\end{proof}

\section{Formalization interfaces and dependency order}
\label{sec:sol-formalization}

\begin{sloppypar}

The public statements above are ordinary mathematical theorems.  The
following stable identifiers expose their intended decomposition before
any Lean declaration names are frozen.

\subsection*{Macro milestones}

\begin{enumerate}
\item \nodeid{PC-M-GRADIENT-RICCI-SOLITON-FOUNDATIONS}: Definition
      \ref{def:gradient-ricci-soliton}, functorial operations, Hamilton
      normalization, and the scalar identities.
\item \nodeid{PC-M-CANONICAL-SELF-SIMILAR-RICCI-FLOW}: Theorem
      \ref{thm:sol-canonical-flow} and its tensor-scaling consequences.
\item \nodeid{PC-M-SHRINKER-SCALAR-QUADRATIC-LOWER-BOUND}: Theorems
      \ref{thm:sol-cly-potential} and \ref{thm:sol-cly-distance}.
\item \nodeid{PC-M-POSITIVELY-CURVED-SHRINKER-COMPACTNESS}: Theorem
      \ref{thm:sol-mw} and Corollary
      \ref{cor:sol-positive-sectional-compact}.
\end{enumerate}

\subsection*{Fine proof order}

\begin{enumerate}
\item \nodeid{PC-F-GRADIENT-RICCI-SOLITON} is the pointwise equality
      \eqref{eq:sol-equation}, with $M,g,f,\sigma$ explicit and no
      completeness hypothesis.  \nodeid{PC-F-SOLITON-WEIGHTED-LAPLACIAN}
      is Definition~\ref{def:sol-weighted-laplacian}; it depends only on a
      metric, a smooth potential, and the rough Laplacian on the chosen
      tensor bundle.
\item \nodeid{PC-F-GAUSSIAN-SHRINKER-WITNESSES} is the existential
      interface \eqref{eq:sol-gaussian-witnesses} in
      Definition~\ref{def:sol-gaussian}.  The metric and potential witness
      equalities are both part of the node; this definition precedes every
      use of Gaussian rigidity in the dependency order.
\item The four conclusions of Proposition~\ref{prop:sol-operations} are
      separate bridge nodes:
      \nodeid{PC-B-SOLITON-POTENTIAL-SHIFT},
      \nodeid{PC-B-SOLITON-DIFFEOMORPHISM-COVARIANCE},
      \nodeid{PC-B-SOLITON-CONSTANT-METRIC-SCALING}, and
      \nodeid{PC-B-SOLITON-PRODUCT}.  Each depends directly on
      \nodeid{PC-F-GRADIENT-RICCI-SOLITON}.  The three equalities in
      \eqref{eq:sol-scaling-data} form the independent bridge
      \nodeid{PC-B-SCALAR-GRADIENT-DISTANCE-CONSTANT-SCALING}.
\item \nodeid{PC-F-SOLITON-TRACE-IDENTITY} is
      \eqref{eq:sol-trace} and depends on
      \nodeid{PC-F-GRADIENT-RICCI-SOLITON}.
      \nodeid{PC-F-SOLITON-SCALAR-GRADIENT-IDENTITY} is
      \eqref{eq:sol-grad-scalar} and additionally depends on contracted
      Bianchi and the Hessian-divergence commutator.
      \nodeid{PC-F-SOLITON-HAMILTON-CONSTANT} is
      \eqref{eq:sol-hamilton-constant} and depends directly on the preceding
      scalar-gradient identity and connectedness.
\item \nodeid{PC-F-SOLITON-WEIGHTED-POTENTIAL-EQUATION} is
      \eqref{eq:sol-drift-potential}; its direct parents are the trace and
      Hamilton-constant nodes.  \nodeid{PC-F-SOLITON-WEIGHTED-SCALAR-EQUATION}
      is \eqref{eq:sol-drift-scalar}; its parents are the scalar-gradient
      identity, the soliton equation, and contracted Bianchi.
\item \nodeid{PC-F-SOLITON-HAMILTON-NORMALIZATION} is
      Definition~\ref{def:sol-hamilton-normalization}; existence and
      uniqueness of its shift depend on
      \nodeid{PC-B-SOLITON-POTENTIAL-SHIFT} and
      \nodeid{PC-F-SOLITON-HAMILTON-CONSTANT}.
      \nodeid{PC-B-SHRINKER-TO-NORMALIZED-SHRINKER} is the simultaneous
      replacement
      $\widehat g=\sigma g$, $\widehat f=f+C/\sigma$ for $\sigma>0$; its
      direct parents are constant metric scaling, Hamilton normalization,
      and the scalar--gradient constant-scaling bridge.
      \nodeid{PC-F-SOLITON-WEIGHTED-INTEGRATION-BY-PARTS} is
      Proposition~\ref{prop:sol-weighted-integration}; it depends on the
      weighted-Laplacian definition and the compact-support density
      divergence theorem.
\item \nodeid{PC-F-LINEAR-GROWTH-VECTOR-FIELD-COMPLETE} is
      Lemma~\ref{lem:sol-linear-growth-complete}; its direct inputs are
      metric completeness, Gronwall, and local ODE continuation.
      \nodeid{PC-F-COMPLETE-SOLITON-SHARP-SCALAR-LOWER} is the imported
      Theorem~\ref{thm:sol-sharp-scalar-lower}, including the Gaussian
      equality case.  These two nodes, together with the Hamilton constant,
      imply \nodeid{PC-F-COMPLETE-SOLITON-POTENTIAL-FIELD}, namely
      Theorem~\ref{thm:sol-potential-field-complete}.
\item \nodeid{PC-F-CANONICAL-SOLITON-TIME-INTERVAL} is
      \eqref{eq:sol-canonical-domain}.
      \nodeid{PC-F-CANONICAL-SOLITON-TIME-CHANGE} is
      \eqref{eq:sol-canonical-time-change}.  The latter depends on the
      former and contains the explicit $\sigma=0$ branch.
      \nodeid{PC-F-CANONICAL-SOLITON-DIFFEOMORPHISMS} is the flow equation
      \eqref{eq:sol-canonical-diffeomorphism}; its parents are the time
      change and completeness of $\nabla f$.
\item Theorem~\ref{thm:sol-canonical-flow} is decomposed into five outputs:
      \nodeid{PC-F-CANONICAL-SOLITON-RICCI-FLOW} is
      \eqref{eq:sol-canonical-ricci-flow};
      \nodeid{PC-F-CANONICAL-TRANSPORTED-SOLITON-EQUATION} is
      \eqref{eq:sol-canonical-soliton};
      \nodeid{PC-F-CANONICAL-POTENTIAL-EVOLUTION} is
      \eqref{eq:sol-canonical-potential-time};
      \nodeid{PC-F-CANONICAL-SLICE-COMPLETE} is the completeness conclusion;
      and \nodeid{PC-F-CANONICAL-HAMILTON-CONSTANT} is
      \eqref{eq:sol-canonical-hamilton-constant}.  Their direct parents are
      the soliton equation, the canonical diffeomorphisms, and the relevant
      pullback and constant-scaling laws.
\item \nodeid{PC-F-CANONICAL-TENSOR-SCALING} is
      Proposition~\ref{prop:sol-canonical-tensor-scaling} and depends on the
      canonical metric definition plus tensor naturality and scale weights.
      \nodeid{PC-F-NORMALIZED-SHRINKER-CONJUGATE-HEAT-BRIDGE} is
      Proposition~\ref{prop:sol-conjugate-heat}; its direct parents are the
      transported soliton equation, potential evolution, and tensor-scaling
      node.
\item \nodeid{PC-F-SOLITON-WEIGHTED-RICCI-EQUATION} is Proposition
      \ref{prop:sol-weighted-ricci}.
      \nodeid{PC-B-LOCAL-CANONICAL-RICCI-DERIVATIVE} is the pointwise local
      identity
      $\left.\partial_t\Ric_{g(t)}\right|_{(0,x)}
        =\mathcal L_{\nabla f}\Ric|_x$ used in its proof; it depends only on
      local ODE existence, pullback naturality, and constant-scale invariance
      of covariant Ricci.  The weighted Ricci node depends directly on this
      bridge and the raw Ricci-flow Ricci evolution.  It does not depend on
      complete-flow existence.
\item \nodeid{PC-F-NORMALIZED-SHRINKER} is
      Definition~\ref{def:sol-normalized-shrinker}; completeness, both
      equations in \eqref{eq:sol-normalized-shrinker}, and the standing
      manifold binders are all part of its interface.
\item \nodeid{PC-F-COMPLETE-SHRINKER-POTENTIAL-GROWTH} is the properness,
      minimum, and two-sided estimate in
      Theorem~\ref{thm:sol-potential-control}.
      \nodeid{PC-F-NONGAUSSIAN-SHRINKER-SCALAR-POSITIVE} is its consequence
      from scalar nonnegativity plus Gaussian rigidity.  These nodes must
      precede every use of an inverse power of $f$.
\item \nodeid{PC-F-SHRINKER-RECIPROCAL-POTENTIAL-EQUATIONS} is
      Lemma~\ref{lem:sol-reciprocal-formulas}.  It depends on the normalized
      shrinker equations and positivity of $f$, and contains exactly the
      powers used later.
\item \nodeid{PC-F-SHRINKER-SCALAR-RECIPROCAL-POTENTIAL-BARRIER} is
      Theorem~\ref{thm:sol-cly-potential}.  Its direct parents are scalar
      positivity, potential properness, the weighted scalar equation, and
      the reciprocal formulas.  The minimum is attained on a compact
      sublevel; no Omori--Yau theorem is used.
\item \nodeid{PC-B-NONFLAT-IFF-NONGAUSSIAN-COMPLETE-SHRINKER} is the bridge
      from the primary source's ``nonflat'' hypothesis to the book's
      ``non-Gaussian'' hypothesis; it depends on Gaussian rigidity.
      \nodeid{PC-B-RECIPROCAL-POTENTIAL-TO-DISTANCE-QUADRATIC} is the
      minimizing-geodesic conversion in
      Theorem~\ref{thm:sol-cly-distance}.  Its public quantifier order is
      $\forall p\,\exists c_p>0\,\forall x$.
      \nodeid{PC-B-EXTERIOR-TO-GLOBAL-QUADRATIC-LOWER} uses scalar
      positivity and compact-core minimization to pass from the source's
      exterior estimate to \eqref{eq:sol-cly-distance-lower}.
\item \nodeid{PC-F-SECTIONAL-CURVATURE-COVARIANT-CONVENTION} is the quotient
      definition in Section~\ref{sec:sol-mw} for linearly independent
      vectors.  \nodeid{PC-F-CURVATURE-RICCI-CONTRACTION} is
      Definition~\ref{def:sol-curvature-ricci-contraction}; repeated Ricci
      eigenspaces do not affect it.
      \nodeid{PC-F-NONNEGATIVE-SECTIONAL-CURVATURE-CONTRACTION-PSD} is
      Lemma~\ref{lem:sol-curvature-ricci-positive} and depends on these two
      nodes.  Its hypothesis is sectional nonnegativity, not
      curvature-operator nonnegativity.
\item \nodeid{PC-F-SHRINKER-RICCI-RECIPROCAL-POTENTIAL-BARRIER} is
      Theorem~\ref{thm:sol-mw-ricci-barrier}.  Its parents are the weighted
      Ricci equation, curvature--Ricci positivity, potential properness, and
      reciprocal formulas.  The compact unit-sphere-bundle minimum avoids a
      global smooth least-eigenvector field.
\item \nodeid{PC-F-SHRINKER-SCALAR-FIXED-INTERVAL-DICHOTOMY} is
      Lemma~\ref{lem:sol-mw-orbit-dichotomy}; its direct parents are the Ricci
      barrier, completeness of the potential field, the scalar-gradient
      identity, and the normalized Hamilton identity
      $R+|\nabla f|^2=f$.  Its exact conclusion is
      $R\geq\min\{n,af\}$.
      \nodeid{PC-F-NORMALIZED-SHRINKER-REGULAR-SUBLEVEL-SCALAR-AVERAGE} is
      Lemma~\ref{lem:sol-sublevel-scalar-average}; its direct inputs are the
      trace identity, potential properness, a regular value, and the density
      divergence theorem.
\item \nodeid{PC-F-COMPLETE-NONCOMPACT-NONNEGATIVE-RICCI-INFINITE-VOLUME}
      is Lemma~\ref{lem:sol-nonnegative-ricci-infinite-volume}; it depends on
      a minimizing ray and Bishop--Gromov comparison.
      It and the preceding two nodes imply
      \nodeid{PC-F-NONNEGATIVE-SECTIONAL-POSITIVE-RICCI-SHRINKER-COMPACT},
      which is Theorem~\ref{thm:sol-mw} after the explicit normalization
      bridge.
\item \nodeid{PC-B-POSITIVE-SECTIONAL-TO-POSITIVE-RICCI} is the
      orthonormal-basis sum in Corollary
      \ref{cor:sol-positive-sectional-compact}; its binder $n\geq2$ is
      essential.  \nodeid{PC-B-SELF-SIMILAR-FLOW-TO-SOLITON-SLICE} is
      Proposition~\ref{prop:sol-converse-self-similar}; its direct inputs are
      the generator witness \eqref{eq:sol-converse-generator}, the derivative
      of a pullback, and the Ricci-flow equation.
\end{enumerate}

\subsection*{Module plan and status axes}

The proposed module partition is
\nolinkurl{Poincare.RicciFlow.Soliton.Basic} for definitions and identities,
\nolinkurl{Poincare.RicciFlow.Soliton.Canonical} for completeness and
self-similarity, \nolinkurl{Poincare.RicciFlow.Soliton.ShrinkerPotential} for
the imported complete-shrinker package,
\nolinkurl{Poincare.RicciFlow.Soliton.ChowLuYang}, and
\nolinkurl{Poincare.RicciFlow.Soliton.MunteanuWang}.  Candidate declaration
names are to be derived mechanically from the full semantic slug after the
\texttt{PC-F-} or \texttt{PC-B-} prefix; they are proposed, not frozen.

For every fine node above, the \emph{interface status} is ``book-stable,
blueprint approval pending,'' the \emph{Lean statement status} and
\emph{Lean proof status} are ``not started,'' and the \emph{build},
\emph{kernel-trust}, and \emph{prose--Lean correspondence} statuses are
``not yet applicable.''  The mathematical status is ``defined or proved in
this chapter'' for the nodes above except
\nodeid{PC-F-COMPLETE-SOLITON-SHARP-SCALAR-LOWER} and
\nodeid{PC-F-COMPLETE-SHRINKER-POTENTIAL-GROWTH}, whose status is ``imported
with exact source statement matched.''  The standard geometric comparison
and ODE inputs named in dependency lists also remain external.  The
editorial owners are the three book authors; a Lean implementation owner
has not yet been assigned.  The
next action is to freeze each atomic interface before declaration names are
created, never to weaken a theorem in response to missing infrastructure.

Likely infrastructure bottlenecks are the orientation-free divergence
theorem on regular sublevels, Sard's theorem, compactness of a unit sphere
bundle over a compact base, prescribed local vector-field jets, smooth
flows of complete vector fields, and Bishop--Gromov comparison.  These are
honest dependencies, not reasons to weaken either public theorem.

\end{sloppypar}

\section{Sources, reconciliations, and proof status}
\label{sec:sol-sources}

The foundational identities, completeness theorem, potential estimates,
and canonical flow were checked against B.~Chow,
\emph{Ricci Solitons in Low Dimensions}, Graduate Studies in Mathematics
\textbf{235}: Section~2.6, pp.~53--56; Theorems~2.14 and 2.27,
pp.~56--68; Theorem~4.3, pp.~107--108; and Proposition~4.19,
pp.~119--121.  The consolidated higher-dimensional identity list in
\emph{Ricci Solitons in Dimensions Four and Higher}, Mathematical Surveys
and Monographs~\textbf{293}, Section~1.6, was used as an independent sign
check.

More precisely, the source-to-node crosswalk is as follows.  The scalar
identity nodes and Hamilton constant are matched to GSM~235,
Section~2.6, pp.~53--56.  The imported sharp scalar node
\nodeid{PC-F-COMPLETE-SOLITON-SHARP-SCALAR-LOWER} is Theorem~2.14,
pp.~56--57, and the potential-field completeness node is Theorem~2.27,
pp.~67--68.  The canonical construction is matched to GSM~235,
Proposition~4.19, pp.~120--121, with the earlier interval description in
Section~2.2, pp.~40--41, repaired below.  The node
\nodeid{PC-F-COMPLETE-SHRINKER-POTENTIAL-GROWTH} is GSM~235,
Theorem~4.3, pp.~107--108, together with the
Haslhofer--M\"uller sharpening stated below.  The weighted Ricci identity is
GSM~235, equation (4.109), p.~130.  Definitions and the
elementary ODE, scaling, pullback, compact-minimum, and source-interface
bridges are internal chapter nodes, with their displayed equations or named
propositions as the exact source spans.

The properness part of Theorem~\ref{thm:sol-potential-control} is the
Cao--Zhou potential theorem.  The precise constants displayed in
\eqref{eq:sol-potential-growth} are from R.~Haslhofer and R.~M\"uller,
\emph{A compactness theorem for complete Ricci shrinkers}, Geometric and
Functional Analysis \textbf{21} (2011), 1091--1116, Lemma~2.1.

The primary source for Theorem~\ref{thm:sol-cly-distance} is B.~Chow,
P.~Lu, and B.~Yang,
\emph{Lower bounds for the scalar curvatures of noncompact gradient Ricci
solitons}, C. R. Math. Acad. Sci. Paris \textbf{349} (2011), 1265--1267,
\href{https://doi.org/10.1016/j.crma.2011.11.004}{DOI record}.
Its Theorem~1 uses ``nonflat'' and states the estimate outside a compact
set.  Complete-shrinker Gaussian rigidity identifies nonflatness with the
non-Gaussian condition used here, and compact-core positivity gives the
global form.  The compact-sublevel presentation above makes explicit the
attained negative minimum and avoids importing a noncompact maximum
principle.

\begin{sloppypar}
Thus the primary-source crosswalk sends
\nodeid{PC-F-SHRINKER-SCALAR-RECIPROCAL-POTENTIAL-BARRIER} and
\nodeid{PC-B-RECIPROCAL-POTENTIAL-TO-DISTANCE-QUADRATIC} to
Chow--Lu--Yang, Theorem~1 and its proof, pp.~1265--1267.  Two changes of
interface are Proposition~\ref{prop:sol-cly-source-bridges}, recorded as
\nodeid{PC-B-NONFLAT-IFF-NONGAUSSIAN-COMPLETE-SHRINKER} and
\nodeid{PC-B-EXTERIOR-TO-GLOBAL-QUADRATIC-LOWER}; they are proved here rather
than attributed to the primary source.  The third is
\nodeid{PC-B-SHRINKER-TO-NORMALIZED-SHRINKER}, which converts the source's
$\sigma=1$, Hamilton-normalized interface to the arbitrary-$\sigma>0$
public theorem by the scaling and potential shift given by
\eqref{eq:sol-scaling-data} and
Definition~\ref{def:sol-hamilton-normalization}.
\end{sloppypar}

The primary source for Theorem~\ref{thm:sol-mw} is O.~Munteanu and J.~Wang,
\emph{Positively curved shrinking Ricci solitons are compact}, Journal of
Differential Geometry \textbf{106} (2017), 499--505,
\href{https://doi.org/10.4310/jdg/1500084024}{DOI record}, Theorem~2.
That theorem is the stronger $K\geq0$, $\Ric>0$ statement used here and has
no bounded-curvature or noncollapsing assumption.

The primary-source crosswalk sends the Ricci reciprocal barrier, orbit
argument, and compactness conclusion to the proof of Munteanu--Wang,
Theorem~2, pp.~500--504.  The regular-sublevel average and infinite-volume
closure used here are the chapter's explicit replacement for the compressed
last step in GSM~235, Theorem~4.47, pp.~138--139.  Consequently
\nodeid{PC-F-NONNEGATIVE-SECTIONAL-POSITIVE-RICCI-SHRINKER-COMPACT} has the
same hypothesis and conclusion as the primary theorem, while its internal
dependency route is the one displayed in Section~\ref{sec:sol-formalization}.

Two source defects and one normalization hazard are repaired explicitly.
\begin{enumerate}
\item In the streamlined proof of the Munteanu--Wang theorem in Graduate
      Studies in Mathematics~235, Theorem~4.47, p.~139, the two orbit cases
      are printed as implying
      $R\geq\max\{n,\frac12e^{-n/c}f\}$.  They imply the minimum in
      \eqref{eq:sol-mw-min-dichotomy}.  That corrected minimum is exactly
      sufficient for the regular-sublevel contradiction.
\item The original Munteanu--Wang proof of Theorem~2, published p.~502,
      first unnumbered display between equations (8) and (9), contains an
      extraneous $2f^{-2}$ in the middle expression for
      $\Delta_f(f^{-1})$ immediately before an equality in which it
      disappears.  Formula \eqref{eq:sol-drift-f-inverse} is the direct
      chain-rule computation and is the formula used in the proof.
\item Hamilton normalization
      $R+|\nabla f|^2=\sigma f$ is shift-sensitive and must not be silently
      conflated with probability normalization of $e^{-f}d\mu$.  The
      Einstein and cylinder constants in Section~\ref{sec:sol-models} are
      regression tests for this distinction.
\end{enumerate}

The canonical domain is always stated as $\{t:1-\sigma t>0\}$.  This also
repairs the expander interval in Graduate Studies in Mathematics~235,
Section~2.2, p.~41, in the paragraph immediately after Proposition~2.2:
for $\sigma<0$ the interval is $(\sigma^{-1},\infty)$, not
$(-\sigma^{-1},\infty)$.  The set description also prevents the
representation domain from being mistaken for a maximal Ricci-flow
lifetime.

Complete human proofs of the two requested headline theorems, including the
barrier, orbit, volume, and scaling arguments, appear above.  The sharp
scalar lower bound for complete solitons, its Gaussian equality case,
the Cao--Zhou potential properness theorem, Hopf--Rinow, Bishop--Gromov,
Sard, and the density divergence theorem are named classical inputs.
No claim is made that any node in this chapter has yet been implemented,
compiled, kernel-checked, or audited in Lean.

\chapter[Three-Dimensional Shrinking Solitons]
{Classification of Complete Three-Dimensional Shrinking Gradient Ricci
Solitons}
\label{ch:three-dimensional-shrinkers}
\index{gradient Ricci soliton!three-dimensional classification}
\index{shrinking Ricci soliton!classification}
\index{Hamilton--Ivey estimate!localized}
\index{curvature operator!strong maximum principle}

\section{Purpose, scope, and the exact model list}
\label{sec:shrinker-classification-purpose}

This chapter classifies complete three-dimensional gradient shrinking Ricci
solitons.  The classification is valid without a curvature bound, a
$\kappa$-noncollapsing hypothesis, orientability, or prior realization as a
Ricci-flow singularity model.  The proof uses the canonical ancient flow of
Chapter~\ref{ch:gradient-solitons}, one localized curvature input, the
fixed-time strong maximum principle and splitting theorem of
Chapter~\ref{ch:strong-curvature}, and the Munteanu--Wang compactness
theorem.

There is an important dependency distinction.  An arbitrary complete
noncompact shrinker cannot be inserted into the closed-manifold
Hamilton--Ivey theorem of Chapter~\ref{ch:hamilton-ivey}.  To prove the
classification for \emph{all} complete shrinkers, one needs a localized
Hamilton--Ivey theorem, or the equivalent theorem that every complete
ancient three-dimensional Ricci flow has nonnegative curvature operator.
For the Poincar\'e-project consumer, where the shrinker is already a smooth
pointed blow-up limit of flows on closed three-manifolds, no new localized
theorem is needed: the closed Hamilton--Ivey estimate passes to the limit.
We prove these two reductions separately.

Unless otherwise stated, every manifold in this chapter is nonempty,
connected, smooth, without boundary, and finite dimensional.  Completeness
means metric completeness.  We retain the trace-normalized curvature
operator of Definition~\ref{def:hi-curvature-operator}; in dimension three
its eigenvalues are twice the principal sectional curvatures and its trace
is $R$.

For a normalized shrinker, a covering statement must remember the
potential, not only the metric.

\begin{definition}[Soliton-model covering]
\label{def:soliton-model-covering}
Let $(N,h,F)$ and $(M,g,f)$ be normalized gradient shrinking solitons.  A
\emph{soliton-model covering} from $(N,h,F)$ to $(M,g,f)$ is a smooth
surjective covering map $\pi:N\to M$ such that
\begin{equation}
                       \pi^*g=h,
             \qquad   f\circ\pi=F.                       \label{eq:soliton-model-covering}
\end{equation}
Thus $\pi$ is a Riemannian covering and the displayed potentials agree
exactly.  The second equality rules out metric quotients on which the
soliton potential does not descend.
\end{definition}

\begin{proposition}[Deck group of a simply connected soliton model]
\label{prop:soliton-model-covering-deck-group}
Let $\pi:(N,h,F)\to(M,g,f)$ be a soliton-model covering, and suppose $N$
is connected and simply connected.  Then $\pi$ is a universal covering.
Its deck group
\begin{equation}
        \Gamma_\pi=
        \{\gamma\in\operatorname{Diff}(N):\pi\circ\gamma=\pi\}             \label{eq:soliton-model-deck-group}
\end{equation}
acts freely and properly discontinuously by $h$-isometries, every
$\gamma\in\Gamma_\pi$ satisfies $F\circ\gamma=F$, and the induced map
\begin{equation}
                 N/\Gamma_\pi\longrightarrow M           \label{eq:soliton-model-quotient-identification}
\end{equation}
is an isometry carrying the descended potential to $f$.

Conversely, let $(N,h,F)$ be a connected normalized shrinker and let
\[
                       \Gamma\leq\operatorname{Diff}(N).
\]
Suppose $\Gamma$ acts freely and properly discontinuously.
If every $\gamma\in\Gamma$ satisfies $\gamma^*h=h$ and
$F\circ\gamma=F$, then the smooth quotient $M=N/\Gamma$ carries unique
data $(g,f)$ with $q^*g=h$ and $f\circ q=F$, and
$q:(N,h,F)\to(M,g,f)$ is a soliton-model covering.  If $N$ is simply
connected, its deck group is exactly $\Gamma$.
\end{proposition}

\begin{proof}
A covering with connected simply connected total space is universal.
Standard universal-cover theory says that its deck group acts freely and
properly discontinuously and that the fibers are exactly the deck orbits,
which gives the diffeomorphism in
\eqref{eq:soliton-model-quotient-identification}.  From
$\pi^*g=h$ and $\pi\circ\gamma=\pi$ we obtain
$\gamma^*h=h$.  From $F=f\circ\pi$ we obtain
$F\circ\gamma=f\circ\pi\circ\gamma=f\circ\pi=F$.  The metric and
potential identities descend through the quotient map.

For the converse, the free properly discontinuous action gives a smooth
covering $q:N\to N/\Gamma$.  Invariance of $h$ and $F$ gives unique
descended $g$ and $f$ with the asserted pullback identities.  The soliton
equation is local and is preserved by a local isometry, so it descends from
$(N,h,F)$ to $(N/\Gamma,g,f)$.  Completeness descends through a Riemannian
covering: a maximal base geodesic lifts after choosing one point over its
initial point, extends for all time in complete $N$, and projects back.
Hamilton normalization descends because its scalar identity pulls back to
the corresponding identity on $N$ and $q$ is surjective.  When $N$ is
simply connected, $q$ is a universal cover and its fibers are precisely the
$\Gamma$-orbits; hence its deck group is $\Gamma$.
\end{proof}

\begin{corollary}[Spherical model coverings are spherical space forms]
\label{cor:spherical-soliton-cover-space-form}
Let $(M^3,g,f)$ be a connected normalized shrinker.  There is a
soliton-model covering
\begin{equation}
 \left(\Sph^3(2),g_{\rm rd},\frac32\right)
       \longrightarrow(M,g,f)                           \label{eq:spherical-soliton-model-cover}
\end{equation}
if and only if there is a finite group
$\Gamma\leq\Isom(\Sph^3(2),g_{\rm rd})$ acting freely such that
$(M,g)$ is isometric to $\Sph^3(2)/\Gamma$ and $f\equiv3/2$.
Under this correspondence, $\Gamma$ is the deck group of the covering.
\end{corollary}

\begin{proof}
Since $\Sph^3$ is simply connected, Proposition
\ref{prop:soliton-model-covering-deck-group} identifies the target with the
quotient by its free properly discontinuous deck group.  This group is
finite: proper discontinuity makes the orbit of one point locally finite,
compactness makes that orbit finite, and freeness identifies the group with
the orbit.  The potential is
the descended constant $3/2$.  Conversely, a finite free isometric action
is properly discontinuous and preserves the constant potential, so the
converse part of Proposition
\ref{prop:soliton-model-covering-deck-group} supplies the asserted
quotient covering; compose it with the stated isometry
$\Sph^3(2)/\Gamma\to M$ to obtain the asserted soliton-model covering.
\end{proof}

Put
\begin{align}
 \mathcal S^2
   &=\left(\Sph^2(\sqrt2),g_{\rm rd},1\right),
                                                               \label{eq:two-shrinker-sphere-model}\\
 \mathcal G^3
   &=\left(\R^3,g_{\rm E},\frac{|x|^2}{4}\right),
                                                               \label{eq:three-shrinker-gaussian-model}\\
 \mathcal S^3
   &=\left(\Sph^3(2),g_{\rm rd},\frac32\right),
                                                               \label{eq:three-shrinker-sphere-model}\\
 \mathcal C^3
   &=\left(\Sph^2(\sqrt2)\times\R,
             g_{\rm rd}+ds^2,1+\frac{s^2}{4}\right).
                                                               \label{eq:three-shrinker-cylinder-model}
\end{align}
These are normalized shrinkers in the sense of
Definition~\ref{def:sol-normalized-shrinker}.  On the cylinder define the
free involutions
\begin{equation}
       a(y,s)=(-y,s),
       \qquad
       b(y,s)=(-y,-s).                                   \label{eq:three-shrinker-cylinder-involutions}
\end{equation}
Both preserve the metric and the potential in
\eqref{eq:three-shrinker-cylinder-model}.

\begin{theorem}[Complete three-dimensional shrinker classification]
\label{thm:complete-three-dimensional-shrinker-classification}
\leavevmode\par\noindent
Let $(M^3,g,f,\sigma)$ be a complete gradient shrinking Ricci soliton.
Assume $\sigma>0$.  Let $C\in\R$ be the unique Hamilton constant determined by
\begin{equation}
        R_g+|\nabla f|_g^2-\sigma f=C,
        \qquad
        \widehat g=\sigma g,
        \qquad
        \widehat f=f+\frac C\sigma .                    \label{eq:three-shrinker-normalization}
\end{equation}
Then $(M,\widehat g,\widehat f)$ is Hamilton normalized, and exactly one of
the following alternatives holds.
\begin{enumerate}
\item \emph{Gaussian.}  There is an isometry
      $\Phi:(M,\widehat g)\to(\R^3,g_{\rm E})$ such that
      $\widehat f=|\Phi|^2/4$.
\item \emph{Spherical.}  There is a finite group
      $\Gamma\leq\Isom(\Sph^3(2),g_{\rm rd})$ acting freely such that
      $(M,\widehat g)$ is isometric to $\Sph^3(2)/\Gamma$ and
      $\widehat f\equiv3/2$.
\item \emph{Cylindrical.}  There is a soliton-model covering
      $\mathcal C^3\to(M,\widehat g,\widehat f)$, and, up to isometry, the
      deck group is exactly one of
      \begin{equation}
          \{1\},\qquad \langle a\rangle,
                 \qquad \langle b\rangle.               \label{eq:three-shrinker-cylinder-deck-list}
      \end{equation}
      Hence the three cylindrical manifolds are
      \begin{equation}
       \Sph^2(\sqrt2)\times\R,\qquad
       \mathbb{RP}^2(\sqrt2)\times\R,\qquad
       \bigl(\Sph^2(\sqrt2)\times\R\bigr)/\langle b\rangle.
                                                               \label{eq:three-shrinker-cylinder-list}
      \end{equation}
\end{enumerate}
The middle cylinder in \eqref{eq:three-shrinker-cylinder-list} is
nonorientable.  The first and third are orientable.  In the unscaled metric
$g$, the spherical radius is $2/\sqrt\sigma$, the spherical factor of a
cylinder has radius $\sqrt{2/\sigma}$, and the lifted Hamilton-normalized
cylinder potential is $1+\sigma s^2/4$, where $s$ is $g$-arclength on the
line factor.
\end{theorem}

\noindent\emph{Formalization crosswalk.}
The public statement above is node
\nodeid{PC-F-COMPLETE-THREE-DIMENSIONAL-SHRINKER-CLASSIFICATION}; its
interface remains draft until the separate blueprint is approved.  Lean
source is not embedded in the mathematical chapter.

The alternatives are genuinely disjoint: their curvature operators have
ranks zero, three, and one, respectively.  The theorem is a classification
of smooth manifolds, not orbifolds.  In particular, a nontrivial finite
rotation quotient of the Gaussian model has a fixed point and is not an
additional smooth model.

\begin{roadmap}{Section~\ref{sec:two-dimensional-shrinker-classification}
proves the complete two-dimensional classification needed by the splitting
branch.}
Section~\ref{sec:shrinker-nonnegative-curvature} gives the general and the
closed-flow-limit routes to nonnegative curvature.  Section
\ref{sec:shrinker-product-potential} upgrades the metric product supplied by
the strong maximum principle to a product of solitons.  Section
\ref{sec:three-dimensional-shrinker-proof} treats the three curvature-rank
branches.  Section~\ref{sec:shrinker-cylinder-quotients} proves the exact
deck-group list.  The final sections record the formalization forest,
source reconciliation, and proof status.
\end{roadmap}

\section{Complete shrinking solitons in dimension two}
\label{sec:two-dimensional-shrinker-classification}

The surface theorem is short once the scalar identities of
Chapter~\ref{ch:gradient-solitons} are available.  Its only substantial
compact-surface input is the Kazdan--Warner identity.  This route is much
shorter than introducing a global circle action and solving the rotational
soliton ODE.

\begin{theorem}[Kazdan--Warner identity on the two-sphere]
\label{thm:kazdan-warner-two-sphere}
Let $(\Sph^2,h)$ be a smooth Riemannian two-sphere, let $R_h$ be its scalar
curvature, and let $X$ be a conformal vector field.  Then
\begin{equation}
               \int_{\Sph^2}\langle\nabla R_h,X\rangle_h
                    \,d\mu_h=0.                          \label{eq:kazdan-warner-two-sphere}
\end{equation}
\end{theorem}

This is a named compact-surface input.  One proof uses uniformization and
the corresponding identity for conformal vector fields of the round
sphere.  Keeping it as a separate interface makes the precise cost of the
short surface proof visible; the longer Killing-field and ODE proof avoids
this input but requires substantially more global infrastructure.

\begin{lemma}[The exponential scalar identity on a shrinking surface]
\label{lem:surface-shrinker-exponential-scalar}
\leavevmode\par\noindent
Let $(M^2,g,f)$ be a normalized shrinker.  The following identities hold:
\begin{align}
 \Ric&=\frac R2g,                                        \label{eq:surface-shrinker-ricci}\\
 \nabla^2f&=\frac{1-R}{2}g,                              \label{eq:surface-shrinker-hessian}\\
 \Delta f&=1-R,                                          \label{eq:surface-shrinker-trace}\\
 dR&=R\,df,                                               \label{eq:surface-shrinker-grad-scalar}\\
 d(Re^{-f})&=0.                                           \label{eq:surface-shrinker-exponential-closed}
\end{align}
Consequently, if $M$ is connected, there is a constant $c\in\R$ such that
\begin{equation}
                             R=ce^f.                     \label{eq:surface-shrinker-exponential-scalar}
\end{equation}
\end{lemma}

\begin{proof}
The algebraic identity \eqref{eq:surface-shrinker-ricci} holds on every
Riemannian surface.  Substitution into the normalized soliton equation and
its trace gives \eqref{eq:surface-shrinker-hessian} and
\eqref{eq:surface-shrinker-trace}.  The scalar-gradient identity
\eqref{eq:sol-grad-scalar} becomes
\[
                dR=2\Ric(\nabla f,\mathord\cdot)=R\,df,
\]
which gives the last two identities by the product rule.  A smooth function
with zero differential is constant on a connected manifold.
\end{proof}

\begin{theorem}[Complete two-dimensional shrinker classification]
\label{thm:complete-two-dimensional-shrinker-classification}
\leavevmode\par\noindent
Let $(M^2,g,f,\sigma)$ be a complete gradient shrinking Ricci soliton.
Assume $\sigma>0$.  Up to an isometry and an additive constant in the potential,
exactly one of the following holds:
\begin{enumerate}
\item the Gaussian plane $(\R^2,g_{\rm E},\sigma|x|^2/4)$;
\item the round sphere $\Sph^2(\sqrt{2/\sigma})$ with constant potential;
\item the round projective plane $\mathbb{RP}^2(\sqrt{2/\sigma})$ with
      constant potential.
\end{enumerate}
In particular, the only complete noncompact shrinking surface is Gaussian.
If $M$ is simply connected and nonflat, it is the round sphere.
\end{theorem}

\begin{sloppypar}
\noindent\emph{Formalization crosswalk.}
The public statement above is node
\nodeid{PC-F-COMPLETE-TWO-DIMENSIONAL-SHRINKER-CLASSIFICATION}; its
interface remains draft until the separate blueprint is approved.
\end{sloppypar}

\begin{proof}
Write
\[
 C_2=R_g+|\nabla f|_g^2-\sigma f,
 \qquad \widehat g=\sigma g,
 \qquad \widehat f=f+\frac{C_2}{\sigma}.
\]
Then $(M,\widehat g,\widehat f)$ is normalized.  These operations preserve completeness,
flatness, compactness, and simple connectivity, so we suppress the hat until
the final sentence.

Theorem~\ref{thm:sol-potential-control} gives $R\geq0$ and says that if $R$
vanishes anywhere, the shrinker is Gaussian.  This proves the flat branch,
including the assertion that no nontrivial complete flat quotient occurs.
Suppose henceforth that the shrinker is not Gaussian.  Then $R>0$ everywhere.
Lemma~\ref{lem:surface-shrinker-exponential-scalar} gives
$R=ce^f$ with $c>0$.  Hamilton normalization gives
\[
                         f=R+|\nabla f|^2>0,
\]
and therefore
\begin{equation}
                 R=ce^f\geq c,
        \qquad   \Ric=\frac R2g\geq\frac c2g.             \label{eq:surface-shrinker-positive-lower}
\end{equation}
Bonnet--Myers makes $M$ compact.

If $M$ is nonorientable, pass to its connected orientable double cover and
lift $g$ and $f$; otherwise use $M$ itself.  The lifted surface is closed and
has positive Gaussian curvature.  Gauss--Bonnet and the classification of
closed orientable surfaces identify it with $\Sph^2$.  On this cover
\eqref{eq:surface-shrinker-hessian} says that $X=\nabla f$ is conformal:
\[
                       \mathcal L_Xg=(1-R)g.
\]
The Kazdan--Warner identity and
\eqref{eq:surface-shrinker-grad-scalar} give
\[
 0=\int_{\Sph^2}\langle\nabla R,\nabla f\rangle\,d\mu
   =\int_{\Sph^2}R|\nabla f|^2\,d\mu.
\]
The integrand is continuous and nonnegative, and $R>0$; hence
$\nabla f=0$.  The soliton equation now gives
\[
                    \Ric=\frac12g,
             \qquad R=1,
             \qquad K=\frac12.                           \label{eq:surface-shrinker-round-data}
\]
Thus the orientable cover is the round $\Sph^2(\sqrt2)$.

It remains to identify the smooth quotient.  Its finite deck group acts
freely by round isometries on $\Sph^2$.  If its order is $d$, multiplicativity
of Euler characteristic gives
\[
                         2=d\,\chi(M).
\]
Hence $d=1$ or $2$.  A fixed-point-free isometric involution of a round
two-sphere is antipodal: as an orthogonal involution of $\R^3$ it has no
$+1$ eigenspace, and therefore equals $-I$.  The quotient is consequently
$\Sph^2$ or $\mathbb{RP}^2$.  Scaling back changes the radius from $\sqrt2$
to $\sqrt{2/\sigma}$.  In the Gaussian branch, if
$\Phi:(M,\widehat g)\to(\R^2,g_{\rm E})$ is the normalized isometry, then
\[
       \Psi=\delta_{\sigma^{-1/2}}\circ\Phi:(M,g)
             \longrightarrow(\R^2,g_{\rm E})
\]
is an isometry and
$f=\sigma|\Psi|^2/4-C_2/\sigma$.  This is the displayed Gaussian
potential up to the permitted additive constant.
\end{proof}

\begin{remark}[Why Perelman's Section~11 is not the surface proof used here]
\label{rem:perelman-not-surface-classification-proof}
Perelman's first paper, Section~11.1, works with complete nonflat ancient
solutions having bounded curvature, nonnegative curvature operator, and
$\kappa$-noncollapsing on all scales.  Its oriented two-dimensional
corollary in Section~11.3 states the classification in that restricted
singularity-model setting and points to Hamilton; it does not display the
completeness-only reduction.  The unrestricted result is Hamilton's 1995
Theorem~26.1.  Thus Perelman's brief corollary is not, by itself, a proof of
Theorem~\ref{thm:complete-two-dimensional-shrinker-classification} under
completeness alone.  The scalar-exponential, Myers, and Kazdan--Warner route
above supplies the missing reduction and is the shortest proof with the
exact hypotheses needed here.
\end{remark}

\section[Nonnegative curvature in dimension three]
{Two routes to nonnegative curvature in dimension three}
\label{sec:shrinker-nonnegative-curvature}

The strong maximum principle applies only after nonnegativity of the
curvature operator has been established.  We first state the exact
noncompact analytic input for the unrestricted theorem.

\begin{theorem}[Complete localized Hamilton--Ivey estimate]
\label{thm:complete-localized-hamilton-ivey}
Let $M^3$ be a smooth manifold without boundary, let $T>0$, and let $g(t)$,
$t\in[0,T]$, be a smooth Ricci flow such that every $(M,g(t))$ is complete.
No curvature bound, compactness, orientability, or initial lower curvature
bound is assumed.  Let $\lambda\geq\mu\geq\nu$ be the eigenvalues of the
trace-normalized curvature operator.  For every $x\in M$ and
$t\in(0,T]$ with $\nu(x,t)<0$,
\begin{equation}
 R(x,t)\geq(-\nu(x,t))
       \left[\log\!\bigl(t(-\nu(x,t))\bigr)-3\right].    \label{eq:complete-localized-hamilton-ivey}
\end{equation}
\end{theorem}

Theorem~\ref{thm:complete-localized-hamilton-ivey} is an explicitly
delegated analytic theorem of Chen--Xu--Zhang.  Its proof localizes the
same Hamilton--Ivey quantity used in Chapter~\ref{ch:hamilton-ivey}.  The
new ingredients are a compactly supported distance cutoff, a barrier for
the distance function at its cut locus, local Ricci control in a moving
ball, and removal of the cutoff radius.  These ingredients are not hidden
inside the classification proof.  For the theorem forest below, the
displayed estimate is a source-matched imported interface; its blueprint
status remains draft.

\begin{corollary}[Complete ancient three-flows have nonnegative curvature]
\label{cor:complete-ancient-three-nonnegative}
Let $g(t)$, $t\in(-\infty,T_*)$, be a smooth three-dimensional Ricci flow
such that every time slice is complete.  Then
\begin{equation}
                         \mathcal R_{g(t)}\geq0
                         \qquad(t<T_*).                  \label{eq:complete-ancient-three-nonnegative}
\end{equation}
No curvature bound is required.
\end{corollary}

\begin{proof}
Fix $(x,t_0)$ with $t_0<T_*$.  Suppose
$q=-\nu(x,t_0)>0$.  For every $L>0$, restrict the ancient flow to
$[t_0-L,t_0]$ and translate the initial time to zero.  Applying
Theorem~\ref{thm:complete-localized-hamilton-ivey} at elapsed time $L$
gives
\[
                         R(x,t_0)\geq q\,[\log(Lq)-3].
\]
The left side is one fixed finite real number, whereas the right side tends
to $+\infty$ as $L\to\infty$.  This contradiction proves $\nu(x,t_0)\geq0$.
The spacetime point was arbitrary.
\end{proof}

\begin{corollary}[A complete three-dimensional shrinker has
nonnegative curvature]
\label{cor:complete-three-shrinker-nonnegative}
Let $(M^3,g,f,\sigma)$ be a complete gradient shrinking Ricci soliton.
Then $\mathcal R_g\geq0$.
\end{corollary}

\begin{proof}
Theorem~\ref{thm:sol-potential-field-complete} and
Theorem~\ref{thm:sol-canonical-flow} produce a smooth complete canonical
Ricci flow on $(-\infty,1/\sigma)$.  Corollary
\ref{cor:complete-ancient-three-nonnegative} applies to that flow.  At
$t=0$ its metric is $g$, so $\mathcal R_g\geq0$.
\end{proof}

For the Poincar\'e critical path we use a different bridge.  The following
statement records enough data to make the limiting argument unambiguous.

\begin{proposition}[Closed Hamilton--Ivey passes to a blow-up limit]
\label{prop:closed-hamilton-ivey-blowup-limit}
Let $g(t)$, $t\in[0,T)$, be a Ricci flow on a closed three-manifold, where
$T>0$.  Let $t_i\in[0,T)$ satisfy $t_i\to T$, let $Q_i>0$ satisfy
$Q_i\to\infty$, let $x_i\in M$, and define
\begin{equation}
                       g_i(s)=Q_i g(t_i+s/Q_i)            \label{eq:closed-flow-parabolic-rescaling}
\end{equation}
whenever the right side is defined.  Let $I\subset\R$ be an open interval.
Suppose that, for every compact $J\subset I$, the flows $g_i(s)$ are
defined on $J$ for all sufficiently large $i$, and that the pointed flows
$(M,g_i(s),x_i)$ converge smoothly on compact subsets of spacetime
$M_\infty\times I$ to a pointed flow
$(M_\infty,g_\infty(s),x_\infty)$.  Concretely, this convergence includes
pointed exhaustion embeddings whose pulled-back metrics converge in
$C^\infty$ on every compact spacetime subset.  Then
\begin{equation}
                         \mathcal R_{g_\infty(s)}\geq0
                         \qquad(s\in I)                  \label{eq:closed-blowup-limit-nonnegative}
\end{equation}
at every point.
\end{proposition}

\begin{proof}
Fix a limiting spacetime point $(x_\infty,s)$ and use the convergence maps
to choose corresponding points $y_i\in M$.  Smooth convergence implies
convergence of the scalar curvature and of the ordered curvature-operator
eigenvalues.  Suppose the limiting least eigenvalue is $-q<0$.  After
discarding finitely many indices,
\[
              -\nu_{g_i}(y_i,s)\geq q/2,
       \qquad |R_{g_i}(y_i,s)|\leq A
\]
for one finite $A$.  Put $u_i=t_i+s/Q_i$.  Then $u_i\to T>0$.
Discard finitely many more indices so that $u_i>0$.

Corollary~\ref{cor:hi-positive-time}, applied to the original closed flow
at $(y_i,u_i)$ and divided by $Q_i$, gives
\begin{equation}
 R_{g_i}(y_i,s)
 \geq(-\nu_{g_i}(y_i,s))
   \left[
    \log\!\bigl(u_iQ_i(-\nu_{g_i}(y_i,s))\bigr)-3
   \right].                                              \label{eq:closed-hi-rescaled-contradiction}
\end{equation}
The right side tends to $+\infty$, while the left side is bounded by $A$,
a contradiction.  Hence the limiting least eigenvalue is nonnegative.
The point and time were arbitrary.
\end{proof}

\begin{remark}[The localization decision]
\label{rem:shrinker-localization-decision}
Corollary~\ref{cor:complete-three-shrinker-nonnegative} proves the theorem
for every complete shrinker but depends on the localized noncompact input
Theorem~\ref{thm:complete-localized-hamilton-ivey}.  Proposition
\ref{prop:closed-hamilton-ivey-blowup-limit} proves exactly the curvature
input needed for shrinkers already known to arise as singularity models of
closed flows, using only Chapter~\ref{ch:hamilton-ivey} and smooth pointed
convergence.  Neither route is silently substituted for the other.
Chen's earlier localized linear cone estimates also imply Corollary
\ref{cor:complete-ancient-three-nonnegative}; they provide a weaker but
potentially cheaper formalization interface if the logarithmic estimate is
not otherwise needed.
\end{remark}

\section{From a metric product to a product shrinker}
\label{sec:shrinker-product-potential}

The strong maximum principle produces a metric splitting of the universal
cover.  For classification we must also split the potential.  This is not
automatic from de Rham's theorem, so we record the required bridge.

\begin{proposition}[Potential splitting on a product shrinker]
\label{prop:product-shrinker-potential-splitting}
Let $(\Sigma^2,h)$ be connected, let $u\in C^\infty(\Sigma\times\R)$, and
let $\sigma>0$.  Suppose
\begin{equation}
 \Ric_{h+ds^2}+\nabla_{h+ds^2}^2u
       =\frac\sigma2(h+ds^2).                            \label{eq:product-shrinker-equation}
\end{equation}
Then there are a constant $s_0\in\R$ and a smooth function
$\psi:\Sigma\to\R$ such that
\begin{align}
 u(y,s)&=\psi(y)+\frac\sigma4(s-s_0)^2,                  \label{eq:product-shrinker-potential-split}\\
 \Ric_h+\nabla_h^2\psi&=\frac\sigma2h.                  \label{eq:surface-factor-shrinker-equation}
\end{align}
If $h+ds^2$ is complete, then $h$ is complete.
\end{proposition}

\begin{proof}
Let $Y$ be tangent to $\Sigma$.  The product connection and the mixed
component of \eqref{eq:product-shrinker-equation} give
\[
       0=\nabla^2u(Y,\partial_s)=Y(\partial_su).
\]
Because $\Sigma$ is connected, $\partial_su$ depends only on $s$.  The
$(\partial_s,\partial_s)$ component gives
\[
                         \partial_s^2u=\frac\sigma2.
\]
Thus
\[
                 u(y,s)=\psi_0(y)+\frac\sigma4s^2+as
\]
for one constant $a$.  Put $s_0=-2a/\sigma$ and absorb the constant
$-\sigma s_0^2/4$ into $\psi_0$ to obtain
\eqref{eq:product-shrinker-potential-split}.  The tangential component of
\eqref{eq:product-shrinker-equation} is exactly
\eqref{eq:surface-factor-shrinker-equation}.  Finally, a Cauchy sequence in
$(\Sigma,h)$ gives a Cauchy sequence in the complete product by adjoining
the coordinate $s=0$; its product limit has $s$-coordinate zero.  Hence
$h$ is complete.
\end{proof}

\begin{corollary}[The rank-one shrinker cover is the round cylinder]
\label{cor:rank-one-shrinker-round-cylinder}
Let $(M^3,g,f)$ be a complete normalized shrinker with nonnegative
curvature operator.  Suppose there is a point $x_0\in M$ such that
\begin{equation}
                   \operatorname{rank}\mathcal R_g(x_0)=1.                \label{eq:rank-one-shrinker-witness}
\end{equation}
Then there is a soliton-model covering
\begin{equation}
        \left(\Sph^2(\sqrt2)\times\R,
          g_{\rm rd}+ds^2,1+\frac{s^2}{4}\right)
             \longrightarrow(M,g,f).                    \label{eq:rank-one-shrinker-cylinder-cover}
\end{equation}
\end{corollary}

\begin{proof}
Restrict the canonical ancient flow to $[-1,1/2)$ and translate time by
$+1$.  Proposition~\ref{prop:sol-canonical-tensor-scaling} transports
$\mathcal R_g\geq0$ to every slice of this restricted flow.  The soliton
slice is the positive time $1$.  Theorem
\ref{thm:curvature-time-slice-trichotomy}, together with
\eqref{eq:rank-one-shrinker-witness}, places the whole slice in its
rank-one branch.  Completeness permits the use of
Theorem~\ref{thm:curvature-complete-trichotomy}.

Let $\pi:\widetilde M\to M$ be the universal covering and lift $g$ and
$f$.  Theorem~\ref{thm:curvature-complete-trichotomy} supplies an isometry
\[
       F:(\Sigma\times\R,h+ds^2)
             \longrightarrow(\widetilde M,\widetilde g),
\]
where $\Sigma$ is connected, complete, and simply connected, and $h$ has
positive Gaussian curvature.  Apply Proposition
\ref{prop:product-shrinker-potential-splitting} to
$u=\widetilde f\circ F$ and translate the line coordinate.  It makes
$(\Sigma,h,\psi)$ a complete normalized two-dimensional shrinker.
Theorem~\ref{thm:complete-two-dimensional-shrinker-classification} and
positive curvature identify it with the round $\Sph^2(\sqrt2)$; the
projective plane cannot occur because $\Sigma$ is simply connected.

The surface potential $\psi$ is constant.  Since the total potential is
Hamilton normalized, its constant is determined rather than arbitrary:
on the cylinder $R=1$ and $|\nabla(s^2/4)|^2=s^2/4$, so
\[
                    \widetilde f=1+\frac{s^2}{4}.
\]
Composing $F$ with $\pi$ gives the covering
\eqref{eq:rank-one-shrinker-cylinder-cover}, including both metric and
potential identities.
\end{proof}

\section{The three curvature-rank branches}
\label{sec:three-dimensional-shrinker-proof}

We first prove roundness in the positive branch by a compact weighted
divergence identity.  This avoids importing the long-time convergence part
of Hamilton's positive-Ricci theorem.  All normalization factors are
displayed because this calculation is especially sensitive to curvature
operator conventions.

\begin{lemma}[Three-dimensional Ricci anisotropy identity]
\label{lem:three-shrinker-anisotropy-identity}
Let $(M^3,g,f)$ be a normalized shrinker with positive sectional curvature.
Define
\begin{equation}
             \mathcal A=\frac{|\Ric|^2}{R^2}-\frac13.    \label{eq:three-shrinker-anisotropy}
\end{equation}
If $\lambda\geq\mu\geq\nu>0$ are the eigenvalues of the trace-normalized
curvature operator, define the nonnegative scalar $J$ by
\begin{equation}
 4J=\frac12\left[
       \lambda^2(\mu-\nu)^2
      +\mu^2(\lambda-\nu)^2
      +\nu^2(\lambda-\mu)^2
     \right].                                            \label{eq:three-shrinker-J-squares}
\end{equation}
Then
\begin{equation}
 \Delta_f\mathcal A+\frac{2}{R}\langle\nabla R,\nabla\mathcal A\rangle
  =\frac2{R^4}|R\nabla\Ric-dR\otimes\Ric|^2
     +\frac4{R^3}J.                                      \label{eq:three-shrinker-anisotropy-elliptic}
\end{equation}
Here $(dR\otimes\Ric)_{kij}=(\nabla_kR)R_{ij}$.
\end{lemma}

\begin{proof}
We expand the quotient calculation rather than hiding it in the phrase
``a standard computation.''  Put
\[
 S=|\Ric|^2,
 \qquad
 P=R_{kij\ell}R^{ij}R^{k\ell}.
\]
The scalar and Ricci-norm evolution equations are
\begin{equation}
 (\partial_t-\Delta)R=2S,
 \qquad
 (\partial_t-\Delta)S=-2|\nabla\Ric|^2+4P.              \label{eq:three-shrinker-R-S-evolution}
\end{equation}
For $D=\partial_t-\Delta$, the one-variable chain rule and the product
rule give
\begin{align}
 D(R^{-2})
   &=-2R^{-3}DR-6R^{-4}|\nabla R|^2,                    \label{eq:three-shrinker-inverse-square}\\
 D(SR^{-2})
   &=R^{-2}DS+S D(R^{-2})
       +4R^{-3}\langle\nabla S,\nabla R\rangle.          \label{eq:three-shrinker-quotient-rule}
\end{align}
Moreover,
\begin{align}
 &\frac2R\langle\nabla R,\nabla(SR^{-2})\rangle
   -\frac2{R^4}|R\nabla\Ric-dR\otimes\Ric|^2           \notag\\
 &\quad=-\frac2{R^2}|\nabla\Ric|^2
      +\frac4{R^3}\langle\nabla S,\nabla R\rangle
      -\frac{6S}{R^4}|\nabla R|^2,                    \label{eq:three-shrinker-gradient-square-completion}
\end{align}
because
$\langle\nabla\Ric,dR\otimes\Ric\rangle
=\frac12\langle\nabla S,\nabla R\rangle$.
Substituting \eqref{eq:three-shrinker-R-S-evolution} into
\eqref{eq:three-shrinker-quotient-rule} and using
\eqref{eq:three-shrinker-gradient-square-completion} therefore yields
\begin{equation}
 D\mathcal A
  =\frac2R\langle\nabla R,\nabla\mathcal A\rangle
   -\frac2{R^4}|R\nabla\Ric-dR\otimes\Ric|^2
   -\frac4{R^3}(S^2-RP).                                \label{eq:three-shrinker-anisotropy-before-algebra}
\end{equation}

It remains only to identify the reaction polynomial.  In an orthonormal
curvature eigenframe, if $r_1,r_2,r_3$ are the Ricci eigenvalues, then
\begin{equation}
 \begin{gathered}
 r_1=\frac{\mu+\nu}{2},\qquad
 r_2=\frac{\lambda+\nu}{2},\qquad
 r_3=\frac{\lambda+\mu}{2},\qquad
 R=\lambda+\mu+\nu,\\
 S=r_1^2+r_2^2+r_3^2,\qquad
 P=\lambda r_2r_3+\mu r_1r_3+\nu r_1r_2.
 \end{gathered}                                          \label{eq:three-shrinker-reaction-eigenframe}
\end{equation}
Direct expansion, collection by the three squared differences, and no
division by an eigenvalue give
\begin{equation}
 S^2-RP
  =\frac18\left[
       \lambda^2(\mu-\nu)^2
      +\mu^2(\lambda-\nu)^2
      +\nu^2(\lambda-\mu)^2
    \right]=J.                                          \label{eq:three-shrinker-reaction-squares}
\end{equation}
Thus the calculation is valid also where eigenvalues coincide.  Equations
\eqref{eq:three-shrinker-anisotropy-before-algebra} and
\eqref{eq:three-shrinker-reaction-squares} give
\begin{equation}
 (\partial_t-\Delta)\mathcal A
  =\frac2R\langle\nabla R,\nabla\mathcal A\rangle
   -\frac2{R^4}|R\nabla\Ric-dR\otimes\Ric|^2
   -\frac4{R^3}J.                                        \label{eq:three-shrinker-anisotropy-parabolic}
\end{equation}

The scalar $\mathcal A$ is unchanged by constant rescaling and natural
under pullback.  Along the canonical flow of the shrinker, differentiation
at $t=0$ therefore gives
\[
                  \left.\partial_t\mathcal A(g(t))\right|_{t=0}
                     =\langle\nabla f,\nabla\mathcal A\rangle.
\]
Substitution in
\eqref{eq:three-shrinker-anisotropy-parabolic} and the definition
$\Delta_f=\Delta-\nabla_{\nabla f}$ give
\eqref{eq:three-shrinker-anisotropy-elliptic}.
\end{proof}

\begin{proposition}[A compact positive three-shrinker is round]
\label{prop:compact-positive-three-shrinker-round}
\leavevmode\par\noindent
Let $(M^3,g,f,\sigma)$ be a compact gradient shrinking Ricci soliton.
Assume it has positive sectional curvature.  Then $g$ has constant sectional curvature
$\sigma/4$.  After Hamilton normalization and scaling to $\sigma=1$, the
universal-cover soliton is
\begin{equation}
                         (\Sph^3(2),g_{\rm rd},3/2).      \label{eq:positive-three-shrinker-sphere-cover}
\end{equation}
\end{proposition}

\begin{proof}
Scale and shift to a normalized shrinker.  Positive sectional curvature
gives $R>0$, so Lemma~\ref{lem:three-shrinker-anisotropy-identity} applies.
Its left side has divergence form:
\begin{equation}
 \operatorname{div}\!\left(R^2e^{-f}\nabla\mathcal A\right)
  =R^2e^{-f}\left(
      \Delta_f\mathcal A+\frac{2}{R}
             \langle\nabla R,\nabla\mathcal A\rangle
    \right).                                             \label{eq:three-shrinker-anisotropy-divergence}
\end{equation}
Integrate over the closed manifold and apply the divergence theorem.
Equations \eqref{eq:three-shrinker-anisotropy-elliptic} and
\eqref{eq:three-shrinker-anisotropy-divergence} give
\begin{equation}
 0=\int_M\left[
       \frac2{R^2}|R\nabla\Ric-dR\otimes\Ric|^2
        +\frac4R J
      \right]e^{-f}\,d\mu.                              \label{eq:three-shrinker-anisotropy-integral}
\end{equation}
Every summand is nonnegative.  Hence $J=0$ everywhere.  Positivity of
$\lambda,\mu,\nu$ and
\eqref{eq:three-shrinker-J-squares} imply
\[
                             \lambda=\mu=\nu.
\]
Thus every two-plane at a point has the same sectional curvature.  Schur's
lemma makes this curvature constant on connected $M$.  The trace identity
and compactness give
\[
       \frac1{\Vol M}\int_MR\,d\mu=\frac32;
\]
since $R$ is constant and $R=6K$, we obtain $K=1/4$.  The soliton equation
then gives $\nabla^2f=0$; a smooth function with zero Hessian on a compact
connected manifold is constant.  Hamilton normalization fixes
$f=R=3/2$.  The universal cover of a complete constant-curvature-$1/4$
three-manifold is the round sphere of radius $2$.  Scaling back gives
$K=\sigma/4$.
\end{proof}

\begin{remark}[Alternative positive-branch proof]
\label{rem:positive-three-shrinker-hamilton-alternative}
One may instead combine Munteanu--Wang compactness with Hamilton's 1982
positive-Ricci convergence theorem.  The volume-normalized canonical flow
of a normalized compact shrinker is $\phi_t^*g$ after the logarithmic time
change.  It is therefore isometric to $g$ at every time, while Hamilton's
theorem makes it converge to a round metric.  Any diffeomorphism-invariant
curvature anisotropy, such as
$\max_M|\Ric-(R/3)g|$, is constant along the former description and tends
to zero along the latter.  Hence it already vanishes on $g$.  The weighted
divergence proof above is chosen because it is local, finite, and does not
require a long-time convergence interface.
\end{remark}

\begin{theorem}[Classification assuming nonnegative curvature]
\label{thm:nonnegative-three-shrinker-classification}
\leavevmode\par\noindent
Let $(M^3,g,f,\sigma)$ be a complete gradient shrinking Ricci soliton.
Assume $\mathcal R_g\geq0$.  Let $C$ be the unique constant satisfying
$R_g+|\nabla f|_g^2-\sigma f=C$, and put
$\widehat g=\sigma g$ and $\widehat f=f+C/\sigma$.  Exactly one of the
following holds:
\begin{enumerate}
\item $(M,\widehat g,\widehat f)$ is the Gaussian model $\mathcal G^3$;
\item it admits a soliton-model covering from $\mathcal S^3$;
\item it admits a soliton-model covering from $\mathcal C^3$.
\end{enumerate}
No curvature bound, orientability, or Ricci-flow uniqueness theorem is
used.
\end{theorem}

\begin{proof}
Scale and Hamilton normalize, so $\sigma=1$.  Proposition
\ref{prop:sol-canonical-tensor-scaling} gives
\[
                  \mathcal R_{g(t)}\geq0
\]
on every slice of the canonical ancient flow.  Restrict
that flow to $[-1,1/2)$ and translate time by $+1$.  The original soliton
slice $t=0$ is now the positive time $1$.  Apply
Theorem~\ref{thm:curvature-time-slice-trichotomy} and, using completeness
of this slice, Theorem~\ref{thm:curvature-complete-trichotomy}.  This uses
the fixed-time theorem only; the bounded-curvature uniqueness assumptions
of Theorem~\ref{thm:curvature-spacetime-trichotomy} never enter.

In the flat branch, $R=0$.  The Gaussian equality case in
Theorem~\ref{thm:sol-sharp-scalar-lower} identifies the normalized soliton
with $\mathcal G^3$.

In the rank-one branch, Corollary
\ref{cor:rank-one-shrinker-round-cylinder} gives a soliton-model covering
from $\mathcal C^3$.

In the positive branch, every sectional curvature is positive.
Corollary~\ref{cor:sol-positive-sectional-compact} makes $M$ compact, with
no curvature bound.  Proposition
\ref{prop:compact-positive-three-shrinker-round} makes the metric have
constant sectional curvature $1/4$ and the Hamilton-normalized potential
equal to $3/2$.  The universal covering is therefore the soliton
$\mathcal S^3$.  The ranks zero, one, and three show that the alternatives
are pairwise disjoint.
\end{proof}

\section{Potential-preserving cylinder quotients}
\label{sec:shrinker-cylinder-quotients}

The phrase ``a quotient of the round cylinder'' is too weak for a soliton
classification.  Translations of the line produce metric quotients such as
$\Sph^2\times\Sph^1$, but the quadratic potential cannot descend to them.
We now prove the exact list in
\eqref{eq:three-shrinker-cylinder-deck-list}.

\begin{lemma}[Form of a potential-preserving cylinder isometry]
\label{lem:potential-preserving-cylinder-isometry}
Let
\[
 \gamma:\Sph^2(\sqrt2)\times\R
       \longrightarrow\Sph^2(\sqrt2)\times\R
\]
be an isometry satisfying
\begin{equation}
          \left(1+\frac{s^2}{4}\right)\circ\gamma
              =1+\frac{s^2}{4}.                          \label{eq:cylinder-isometry-preserves-potential}
\end{equation}
Then there are $A\in\Isom(\Sph^2(\sqrt2))$ and
$\varepsilon\in\{+1,-1\}$ such that
\begin{equation}
                         \gamma(y,s)=(Ay,\varepsilon s). \label{eq:cylinder-isometry-normal-form}
\end{equation}
\end{lemma}

\begin{proof}
The Ricci endomorphism of the cylinder has eigenvalue $1/2$ on the tangent
planes to the sphere and eigenvalue zero on the line distribution
$\R\partial_s$.  Isometries preserve the Ricci endomorphism, hence preserve
both distributions.  Since the cylinder is connected,
$d\gamma(\partial_s)=\varepsilon\partial_s$ for one constant sign
$\varepsilon$.  The sphere component of $\gamma$ is therefore independent
of $s$, and the line component is independent of $y$ and has derivative
$\varepsilon$.  Thus initially
\[
                         \gamma(y,s)=(Ay,\varepsilon s+d)
\]
for a sphere isometry $A$ and a constant $d$.  Substituting in
\eqref{eq:cylinder-isometry-preserves-potential} gives
$(\varepsilon s+d)^2=s^2$ for every $s$, hence $d=0$.
\end{proof}

\begin{theorem}[Exact smooth quotients of the cylinder soliton]
\label{thm:exact-cylinder-soliton-quotients}
Let
\begin{equation}
 \operatorname{Aut}_{\rm sol}(\mathcal C^3)
  =\left\{\gamma\in\Isom(\Sph^2(\sqrt2)\times\R):
       (1+s^2/4)\circ\gamma=1+s^2/4\right\}.             \label{eq:cylinder-soliton-automorphism-group}
\end{equation}
Let $\Gamma\leq\operatorname{Aut}_{\rm sol}(\mathcal C^3)$, and suppose its induced
action is free and properly discontinuous.  Then exactly one of the
following subgroup equalities holds:
\begin{equation}
          \Gamma=\{1\},\qquad
          \Gamma=\langle a\rangle,\qquad
          \Gamma=\langle b\rangle,                       \label{eq:exact-cylinder-groups}
\end{equation}
where $a$ and $b$ are defined in
\eqref{eq:three-shrinker-cylinder-involutions}.
\end{theorem}

\begin{proof}
Lemma~\ref{lem:potential-preserving-cylinder-isometry} shows that every
$\gamma\in\Gamma$ preserves the central sphere
$\Sph^2\times\{0\}$.  Its restriction to that sphere is free: if its sphere
isometry fixed $y$, then $\gamma$ would fix $(y,0)$.  The restriction is
also faithful, because an element with identity sphere component and sign
$-1$ fixes the whole central sphere.  Proper discontinuity on the invariant
compact central sphere makes $\Gamma$ finite.

The restricted action is therefore a free finite isometric action on a
round two-sphere.  Multiplicativity of Euler characteristic gives
\[
            2=|\Gamma|\,
                \chi\bigl(\Sph^2/\Gamma\bigr).
\]
The quotient has positive Gaussian curvature, so Gauss--Bonnet makes its
Euler characteristic positive.  It follows that $|\Gamma|=1$ or $2$.
In the second case let $\gamma$ be the nonidentity element.  Its sphere
component is a fixed-point-free orthogonal involution of $\R^3$, hence is
$-I$.  The line sign is independently $+1$ or $-1$, giving $a$ or $b$.
\end{proof}

\begin{corollary}[Topology and orientation of the cylinder quotients]
\label{cor:cylinder-quotient-topology}
The quotient by $a$ is the product
$\mathbb{RP}^2\times\R$ and is nonorientable.  The quotient by $b$ is the
total space of the tautological real line bundle over $\mathbb{RP}^2$; it
is orientable and is diffeomorphic to punctured $\mathbb{RP}^3$.  No
potential-preserving smooth quotient is compact, and in particular
$\Sph^2\times\Sph^1$ is not a shrinking gradient soliton quotient of
$\mathcal C^3$.
\end{corollary}

\begin{proof}
The first identification is immediate because $a$ acts only on the sphere.
For the second, the quotient relation $(y,s)\sim(-y,-s)$ is exactly the
tautological-line-bundle relation.  The antipodal map reverses the
orientation of $\Sph^2$, and $s\mapsto-s$ reverses the orientation of the
line, so their product preserves the product orientation.  The map
\[
   [(y_1,y_2,y_3),s]\longmapsto[y_1:y_2:y_3:s]
\]
identifies this line bundle with $\mathbb{RP}^3$ minus the point
$[0:0:0:1]$.  Finally, Lemma
\ref{lem:potential-preserving-cylinder-isometry} excludes every nonzero
translation in the line.  A compact cylinder quotient would require such a
translation, so none occurs.
\end{proof}

\begin{proof}[Proof of Theorem~\ref{thm:complete-three-dimensional-shrinker-classification}]
Corollary~\ref{cor:complete-three-shrinker-nonnegative} supplies
$\mathcal R_g\geq0$ without a curvature bound.  Theorem
\ref{thm:nonnegative-three-shrinker-classification} gives the Gaussian,
spherical-cover, and cylinder-cover alternatives after the exact scaling
and potential shift in \eqref{eq:three-shrinker-normalization}.

In the spherical branch Corollary
\ref{cor:spherical-soliton-cover-space-form} gives the finite free quotient
in the statement.  In the cylindrical
branch every deck transformation preserves the lifted potential exactly,
because that potential is the pullback of $\widehat f$ from $M$.
Theorem~\ref{thm:exact-cylinder-soliton-quotients} gives precisely the three
groups in \eqref{eq:three-shrinker-cylinder-deck-list}.  Conversely, each
listed free action preserves both the metric and potential, so each quotient
is a smooth normalized shrinker.  The rank distinction proves uniqueness of
the top-level alternative.  Scaling formulas
\eqref{eq:sol-scaling-data} give the unnormalized radii and line potential.
\end{proof}

\begin{corollary}[Orientable complete three-dimensional shrinkers]
\label{cor:orientable-three-shrinker-classification}
Under the hypotheses of
Theorem~\ref{thm:complete-three-dimensional-shrinker-classification}, if
$M$ is orientable, the model $\mathbb{RP}^2\times\R$ is absent.  Thus the
only cylindrical possibilities are the round product and the diagonal
quotient by $b(y,s)=(-y,-s)$.
\end{corollary}

\begin{corollary}[Project-facing singularity-model classification]
\label{cor:singularity-model-shrinker-classification}
Let $g_0(t)$, $t\in[0,T)$, be a Ricci flow on a closed
three-manifold $M_0$, where $T>0$.  Let $t_i\in[0,T)$ tend to $T$, let
$Q_i>0$ tend to infinity, let $x_i\in M_0$, and put
\[
                       g_i(s)=Q_i g_0(t_i+s/Q_i)
\]
whenever the right side is defined.  Let $I\subset\R$ be an open interval
containing zero.  Suppose that for every compact $J\subset I$, the flows
$g_i$ are defined on $J$ for all sufficiently large $i$, and that
$(M_0,g_i(s),x_i)$ converge smoothly on compact spacetime subsets, by
pointed exhaustion embeddings, to
\[
                   (M_\infty,g_\infty(s),x_\infty).
\]
Let $(N^3,h,F)$ be a complete
normalized shrinker with canonical flow $h_{\rm can}(s)$ on $I$, and fix
$z\in N$.  Suppose there is one basepoint-preserving diffeomorphism
\begin{equation}
 \Phi:(N,z)\longrightarrow(M_\infty,x_\infty)
 \quad\hbox{such that}\quad
 \Phi^*g_\infty(s)=h_{\rm can}(s)\quad(s\in I).          \label{eq:canonical-shrinker-limit-slice-identification}
\end{equation}
In particular, $h=h_{\rm can}(0)=\Phi^*g_\infty(0)$.  Then, without using
Theorem~\ref{thm:complete-localized-hamilton-ivey}, it is exactly one of:
the Gaussian $\mathcal G^3$; a finite free quotient of
$\mathcal S^3$; or a quotient of $\mathcal C^3$ by one of the three groups
in \eqref{eq:exact-cylinder-groups}.  If the shrinker is nonflat, the
Gaussian alternative is absent.  If it is also orientable, the
$\mathbb{RP}^2\times\R$ alternative is absent.
\end{corollary}

\noindent\emph{Formalization crosswalk.}
The public statement above is node
\nodeid{PC-F-SINGULARITY-MODEL-SHRINKER-CLASSIFICATION}; it belongs to the
closed-flow singularity-model lane, not to the unrestricted complete-flow
lane.

\begin{proof}
Proposition~\ref{prop:closed-hamilton-ivey-blowup-limit} gives
$\mathcal R_{g_\infty(0)}\geq0$ using the closed estimate from
Chapter~\ref{ch:hamilton-ivey}.  Pullback through $\Phi$ and
\eqref{eq:canonical-shrinker-limit-slice-identification} give
$\mathcal R_h\geq0$.  Apply Theorem
\ref{thm:nonnegative-three-shrinker-classification} and then Theorem
\ref{thm:exact-cylinder-soliton-quotients}.  A nonflat metric is not
Gaussian, and Corollary~\ref{cor:cylinder-quotient-topology} gives the
orientability assertion.
\end{proof}

\section{Normalization and boundary regression tests}
\label{sec:shrinker-classification-regression}

The following checks are part of the theorem interface.
\begin{enumerate}
\item On a normalized shrinking surface, a round solution has
      $K=1/2$, $R=1$, and Hamilton-normalized potential $f=1$.
\item On a normalized round three-shrinker, $K=1/4$, $R=3/2$, and
      $f=3/2$.
\item On the normalized cylinder, the trace-normalized curvature operator
      has eigenvalues $(1,0,0)$, so $R=1$, and
      \[
        \left|\nabla\left(1+s^2/4\right)\right|^2=s^2/4.
      \]
      Hence $R+|\nabla f|^2=f$ exactly.
\item The surface factor in the \emph{universal cover} of a rank-one
      three-shrinker is $\Sph^2$, never $\mathbb{RP}^2$.  The projective
      plane appears only after dividing the cylinder by $a$.
\item A translation $s\mapsto s+L$, $L\ne0$, is a cylinder-metric
      isometry but does not preserve $1+s^2/4$.  Thus a metric quotient
      $\Sph^2\times\Sph^1$ is not a quotient soliton in the sense of
      Definition~\ref{def:soliton-model-covering}.
\item The smooth-manifold hypothesis is essential.  Finite rotation
      quotients of a Gaussian model and the nonround teardrop or unequal
      football shrinkers are orbifold examples; they do not enlarge the
      smooth list in Theorems
      \ref{thm:complete-two-dimensional-shrinker-classification} and
      \ref{thm:complete-three-dimensional-shrinker-classification}.
\end{enumerate}

\section{Formalization theorem forest and module plan}
\label{sec:shrinker-classification-formalization}

The natural formalization separates model data, the two-dimensional
engine, nonnegative curvature, the three rank branches, and quotient
enumeration.  Most importantly, the common classification under
$\mathcal R\geq0$ is a module boundary: the unrestricted and
singularity-model routes are sibling consumers of that common core.  The
macro milestones are
\begin{sloppypar}
\begin{enumerate}
\item \nodeid{PC-M-SHRINKER-MODEL-COVERS};
\item \nodeid{PC-M-COMPLETE-TWO-DIMENSIONAL-SHRINKER-CLASSIFICATION};
\item \nodeid{PC-M-THREE-DIMENSIONAL-SHRINKER-RANK-REDUCTION};
\item \nodeid{PC-M-NONNEGATIVE-THREE-DIMENSIONAL-SHRINKER-CLASSIFICATION};
\item \nodeid{PC-M-COMPLETE-ANCIENT-THREE-DIMENSIONAL-NONNEGATIVITY};
\item \nodeid{PC-M-COMPLETE-THREE-DIMENSIONAL-SHRINKER-CLASSIFICATION};
\item \nodeid{PC-M-SINGULARITY-MODEL-SHRINKER-CLASSIFICATION}.
\end{enumerate}
\end{sloppypar}
The sixth milestone uses localized Hamilton--Ivey to feed the fourth.  The
seventh instead uses the closed-flow pointed-limit bridge to feed the same
fourth milestone.  Neither classification wrapper depends on the other.

\subsection*{Fine nodes in dependency order}

\begin{sloppypar}
\begin{enumerate}
\item \nodeid{PC-F-SOLITON-MODEL-RIEMANNIAN-COVER} is
      Definition~\ref{def:soliton-model-covering}.  Its interface contains
      the covering map, the metric pullback equality, and the potential
      equality in \eqref{eq:soliton-model-covering}.
      \nodeid{PC-B-SIMPLY-CONNECTED-SOLITON-COVER-DECK-QUOTIENT} is
      Proposition~\ref{prop:soliton-model-covering-deck-group}; it packages
      universality, the deck subgroup, free proper action, preservation of
      metric and potential, and the quotient identification.  Conjugacy in
      later quotient nodes is always inside the potential-preserving model
      automorphism group.
      The model checks
      \nodeid{PC-F-NORMALIZED-THREE-GAUSSIAN-SHRINKER-MODEL},
      \nodeid{PC-F-NORMALIZED-ROUND-TWO-SPHERE-SHRINKER-MODEL},
      \nodeid{PC-F-NORMALIZED-ROUND-THREE-SPHERE-SHRINKER-MODEL}, and
      \nodeid{PC-F-NORMALIZED-ROUND-THREE-CYLINDER-SHRINKER-MODEL}
      record the exact radii and Hamilton constants in
      \eqref{eq:two-shrinker-sphere-model}--
      \eqref{eq:three-shrinker-cylinder-model}.
\item \nodeid{PC-B-NORMALIZED-MODEL-COVER-TO-ARBITRARY-SHRINKER} is the
      scaling and potential-shift wrapper
      \eqref{eq:three-shrinker-normalization}.  It depends on
      \nodeid{PC-B-SHRINKER-TO-NORMALIZED-SHRINKER} from Chapter~6 and
      must transport both equalities in a soliton-model covering.
\item \nodeid{PC-F-TWO-DIMENSIONAL-RICCI-SCALAR-IDENTITY} is
      \eqref{eq:surface-shrinker-ricci}.
      \nodeid{PC-F-TWO-DIMENSIONAL-SOLITON-EXPONENTIAL-SCALAR} is
      Lemma~\ref{lem:surface-shrinker-exponential-scalar}; its parents are
      the surface Ricci identity and
      \eqref{eq:sol-grad-scalar}.
\item \nodeid{PC-F-NONFLAT-TWO-DIMENSIONAL-SHRINKER-COMPACT} is the
      argument \eqref{eq:surface-shrinker-positive-lower} followed by
      Bonnet--Myers.  It depends on scalar positivity, Hamilton
      normalization, connectedness, and the exponential scalar node.
\item \nodeid{PC-F-KAZDAN-WARNER-SPHERE-IDENTITY} is the imported
      Theorem~\ref{thm:kazdan-warner-two-sphere}.
      Together with the orientable-cover, Gauss--Bonnet, and round-quotient
      bridges, it yields
      \nodeid{PC-F-COMPLETE-TWO-DIMENSIONAL-SHRINKER-CLASSIFICATION},
      Theorem~\ref{thm:complete-two-dimensional-shrinker-classification}.
\item \nodeid{PC-F-COMPLETE-LOCAL-HAMILTON-IVEY-ESTIMATE} is the imported
      Theorem~\ref{thm:complete-localized-hamilton-ivey}.  The exact public
      statement has complete time slices and no bounded-curvature, scalar
      positivity, compactness, or initial lower-bound hypothesis.
      It implies
      \nodeid{PC-F-COMPLETE-ANCIENT-THREE-DIMENSIONAL-CURVATURE-NONNEGATIVE},
      Corollary~\ref{cor:complete-ancient-three-nonnegative}, by one time
      translation and a logarithmic contradiction.
\item \nodeid{PC-B-CANONICAL-SHRINKER-AS-COMPLETE-ANCIENT-FLOW} combines
      the canonical-flow and potential-field-completeness nodes of
      Chapter~6.  With the preceding ancient node it gives
      \nodeid{PC-F-COMPLETE-THREE-DIMENSIONAL-SHRINKER-CURVATURE-NONNEGATIVE},
      Corollary~\ref{cor:complete-three-shrinker-nonnegative}.
\item \nodeid{PC-B-CLOSED-BLOWUP-LIMIT-HAMILTON-IVEY-NONNEGATIVITY} is
      Proposition~\ref{prop:closed-hamilton-ivey-blowup-limit}.  Its direct
      parents are Chapter~3's closed positive-time Hamilton--Ivey estimate,
      parabolic tensor scaling, and continuity of scalar curvature and
      ordered self-adjoint eigenvalues under smooth pointed convergence.
      It is independent of the localized node.
      \nodeid{PC-B-CANONICAL-SHRINKER-LIMIT-SLICE-IDENTIFICATION} is the
      witness in Corollary
      \ref{cor:singularity-model-shrinker-classification}: a fixed pointed
      flow isometry identifies the supplied canonical shrinker flow on one
      open interval containing zero with that limit, and hence identifies
      the metric at time zero.
\item \nodeid{PC-B-ANCIENT-TIME-TRANSLATION-TO-POSITIVE-SLICE} is the
      translation of $[-1,1/2)$ used in the proof of
      Theorem~\ref{thm:nonnegative-three-shrinker-classification}.
      It permits application of the existing
      \nodeid{PC-F-CURVATURE-TIME-SLICE-TRICHOTOMY} and
      \nodeid{PC-F-CURVATURE-FIXED-TIME-UNIVERSAL-COVER} without invoking
      the bounded-curvature whole-flow theorem.  Their specialization is
      \nodeid{PC-F-NORMALIZED-THREE-DIMENSIONAL-SHRINKER-CURVATURE-TRICHOTOMY}.
\item \nodeid{PC-F-PARALLEL-LINE-SHRINKER-POTENTIAL-SPLITTING} is
      Proposition~\ref{prop:product-shrinker-potential-splitting}.
      Its mixed, line--line, and tangential outputs must remain separate
      conjuncts.  It and the two-dimensional classification imply
      \nodeid{PC-B-ROUND-SURFACE-FACTOR-TO-CYLINDER-MODEL}, Corollary
      \ref{cor:rank-one-shrinker-round-cylinder}.
\item \nodeid{PC-F-THREE-DIMENSIONAL-RICCI-ANISOTROPY-IDENTITY} is
      Lemma~\ref{lem:three-shrinker-anisotropy-identity}, including the
      invariant reaction polynomial $J=S^2-RP$ and its sum-of-squares
      identity \eqref{eq:three-shrinker-reaction-squares}.
      Compact weighted integration and Schur's lemma give
      \nodeid{PC-F-COMPACT-POSITIVE-CURVATURE-THREE-SHRINKER-ROUND},
      Proposition~\ref{prop:compact-positive-three-shrinker-round}.
      Its compactness parent is
      \nodeid{PC-F-NONNEGATIVE-SECTIONAL-POSITIVE-RICCI-SHRINKER-COMPACT}
      from Chapter~6.
\item The flat, rank-one, and positive branches assemble as
      \nodeid{PC-F-COMPLETE-NONNEGATIVE-THREE-SHRINKER-CLASSIFICATION},
      Theorem~\ref{thm:nonnegative-three-shrinker-classification}.
      \nodeid{PC-F-THREE-SHRINKER-MODEL-CURVATURE-RANK-DISJOINTNESS}
      records that the three model curvature-operator ranks are respectively
      zero, one, and three, so the resulting disjunction is exclusive.
\item \nodeid{PC-F-CYLINDER-POTENTIAL-PRESERVING-ISOMETRY} is
      Lemma~\ref{lem:potential-preserving-cylinder-isometry}.
      \nodeid{PC-F-CYLINDER-POTENTIAL-PRESERVING-FREE-ACTION-CLASSIFICATION}
      is Theorem~\ref{thm:exact-cylinder-soliton-quotients}.  These nodes
      exclude translations before any group classification is attempted.
      \nodeid{PC-B-SOLITON-MODEL-COVER-TO-SPHERICAL-SPACE-FORM} is the
      exact Corollary~\ref{cor:spherical-soliton-cover-space-form}, and
      \nodeid{PC-F-CYLINDER-QUOTIENT-TOPOLOGY-ORIENTATION} is
      Corollary~\ref{cor:cylinder-quotient-topology}.
\item The all-complete branch concludes with
      \nodeid{PC-F-COMPLETE-THREE-DIMENSIONAL-SHRINKER-CLASSIFICATION},
      Theorem~\ref{thm:complete-three-dimensional-shrinker-classification}.
      The separate project lane concludes with
      \nodeid{PC-F-SINGULARITY-MODEL-SHRINKER-CLASSIFICATION}, Corollary
      \ref{cor:singularity-model-shrinker-classification}.
\end{enumerate}
\end{sloppypar}

The coarse dependency forest is therefore
\begin{equation}
\begin{gathered}
 \text{Chapter~6 canonical flow}+\text{complete localized Hamilton--Ivey}
       \Longrightarrow \mathcal R\geq0,\\
 \mathcal R\geq0+\text{Chapter~5 fixed-time strong trichotomy}
       \Longrightarrow \{\text{flat},\text{rank one},\text{positive}\},\\
 \text{flat}\Longrightarrow\text{Gaussian},\qquad
 \text{rank one}+\text{surface classification}\Longrightarrow\text{cylinder},\\
 \text{positive}+\text{Munteanu--Wang}+\text{anisotropy identity}
       \Longrightarrow\text{sphere}.
\end{gathered}                                             \label{eq:shrinker-classification-coarse-forest}
\end{equation}
For the singularity-model lane, replace only the first line by
\[
 \text{closed Hamilton--Ivey}+\text{smooth pointed blow-up convergence}
       \Longrightarrow\mathcal R\geq0.
\]

\subsection*{Module plan and status axes}

\begin{sloppypar}
The proposed module partition is
\begin{enumerate}
\item \nolinkurl{Poincare.RicciFlow.Soliton.ModelCover} for model triples,
      soliton coverings, deck actions, and quotient reconstruction;
\item \nolinkurl{Poincare.RicciFlow.Soliton.Shrinker2D} for the surface
      identities and classification;
\item \nolinkurl{Poincare.RicciFlow.Soliton.SplitPotential} for the product
      potential bridge;
\item \nolinkurl{Poincare.RicciFlow.Soliton.Shrinker3DCore} for the
      classification under the explicit hypothesis $\mathcal R\geq0$,
      including rank reduction, roundness, and quotient enumeration;
\item \nolinkurl{Poincare.RicciFlow.DimensionThree.LocalHamiltonIvey} for
      the unrestricted complete-flow curvature input, and
      \nolinkurl{Poincare.RicciFlow.Soliton.Shrinker3DComplete} for the
      wrapper from that input to the all-complete theorem;
\item \nolinkurl{Poincare.RicciFlow.SingularityModel.ShrinkerClassification}
      for the separate wrapper using closed Hamilton--Ivey and pointed
      convergence.
\end{enumerate}
Thus \nolinkurl{Shrinker3DCore} imports neither nonnegativity route.
\nolinkurl{Shrinker3DComplete} imports the localized route, whereas
\nolinkurl{SingularityModel.ShrinkerClassification} imports the closed-flow
limit bridge.  This module DAG prevents localized Hamilton--Ivey from
entering the project-facing theorem transitively.
\end{sloppypar}

This section is a source-to-blueprint crosswalk, not an approved blueprint.
For every new node above the separate status axes are as follows:
\begin{center}
\begin{tabular}{@{}lp{0.58\textwidth}@{}}
interface & draft; approval pending,\\
informal mathematics & node-specific; recorded below,\\
Lean declaration & absent,\\
Lean proof & absent,\\
build & not run,\\
kernel and trust check & not run,\\
axiom footprint & not computed,\\
dependency closure & mapped, not closed,\\
declared-scope coverage & unverified,\\
provenance & chapter locators present; blueprint locators pending,\\
prose--Lean correspondence & unverified.
\end{tabular}
\end{center}
The explicitly stated lemmas, propositions, corollaries, and classification
theorems proved in this chapter have human-proof status ``independently
verified conditional on the named external inputs.''  The source-matched
imports are
\nodeid{PC-F-KAZDAN-WARNER-SPHERE-IDENTITY} and
\nodeid{PC-F-COMPLETE-LOCAL-HAMILTON-IVEY-ESTIMATE}.  Myers,
Gauss--Bonnet, closed-surface classification, space forms, universal
covers, and smooth pointed convergence are separately visible external
infrastructure.  Fine nodes which presently name a decomposition step but
do not have a separately boxed statement remain draft proof-extraction
obligations; the prose proof does not by itself freeze their interfaces.
The mathematical review status is therefore conditional verification, and
the formalization review status remains unverified.  Exact theorem-card
quantifiers must encode dimension by a
hypothesis such as $\dim M=2$ or $\dim M=3$; the superscripts in $M^2$ and
$M^3$ are expository abbreviations only.  Exact dependency arrays,
dependents, and macro-to-fine membership belong in the separate blueprint
before any interface is frozen.  No sentence here claims that a new
declaration has been compiled or kernel checked.

The highest-risk implementation interfaces are the localized distance
cutoff at the cut locus, Bonnet--Myers, Kazdan--Warner and its
uniformization dependencies, Riemannian universal covers and quotient
metrics, ordered eigenvalue continuity under pointed smooth convergence,
and the quotient-rule verification of
\eqref{eq:three-shrinker-anisotropy-parabolic}.  These are reasons to
modularize, not reasons to weaken the public theorem.

\section{Sources, reconciliations, and proof status}
\label{sec:shrinker-classification-sources}

The two-dimensional identities and classification were checked against
B.~Chow, \emph{Ricci Solitons in Low Dimensions}, Graduate Studies in
Mathematics~\textbf{235}, equations~(3.30)--(3.33), pp.~82--83;
Theorem~3.5, pp.~84--85; Theorem~3.12, p.~90; and the Kazdan--Warner
identity in Appendix~A, Theorem~A.1, pp.~307--309.  Hamilton's original
compact-surface input is R.~Hamilton, \emph{The Ricci flow on surfaces},
Contemporary Mathematics~\textbf{71} (1988), 237--262.  Perelman,
\emph{The entropy formula for the Ricci flow and its geometric
applications}, arXiv:math/0211159v1, Section~11.3, p.~28, was used only as
a check on the singularity-model perspective.  As explained in Remark
\ref{rem:perelman-not-surface-classification-proof}, its citation to
Hamilton does not display the completeness-only noncompact reduction.  The
primary unrestricted source is R.~Hamilton, \emph{The formation of
singularities in the Ricci flow}, Surveys in Differential
Geometry~\textbf{2} (1995), Section~26, especially Theorem~26.1; compare
GSM~235, Theorem~3.12, and MSM~293, Theorem~1.70.

The primary source for
Theorem~\ref{thm:complete-localized-hamilton-ivey} is B.-L.~Chen, G.~Xu,
and Z.-H.~Zhang, \emph{Local pinching estimates in 3-dim Ricci flow},
Mathematical Research Letters~\textbf{20} (2013), 845--855,
Corollary~1.5,
\href{https://doi.org/10.4310/MRL.2013.v20.n5.a3}{DOI record}; its
Corollary~2.3 is
Corollary~\ref{cor:complete-ancient-three-nonnegative}.  B.-L.~Chen,
\emph{Strong uniqueness of the Ricci flow}, Journal of Differential
Geometry~\textbf{82} (2009), 363--382, Corollary~2.4, gives the same ancient
nonnegativity conclusion from localized linear cone estimates.  The
streamlined account in GSM~235, Theorem~5.3 and Corollary~5.4,
pp.~155--162, assumes $R>0$ in the displayed local theorem; that weaker
form already suffices after separating the Gaussian shrinker, but the
primary Chen--Xu--Zhang statement recorded here does not require scalar
positivity.

\begin{sloppypar}
The strong maximum principle and parallel-image mechanism originate in
R.~Hamilton, \emph{Four-manifolds with positive curvature operator},
Journal of Differential Geometry~\textbf{24} (1986), Sections~8--9.  This
chapter uses the repaired, orientation-free fixed-time interfaces already
proved in Chapter~\ref{ch:strong-curvature}.  The compact positive branch
adapts the shrinker calculation of L.~Ni and N.~Wallach,
\emph{On a classification of gradient shrinking solitons}, Mathematical
Research Letters~\textbf{15} (2008), Proposition~3.2 and
equations~(3.2)--(3.6), pp.~944--946.  That calculation begins from
Hamilton's three-dimensional Ricci-anisotropy evolution formula; compare
GSM~235, equations~(5.40)--(5.52), pp.~164--167.  The compact weighted
integration here is written independently.  Compactness comes from the
Munteanu--Wang theorem proved in Chapter~\ref{ch:gradient-solitons}.
\end{sloppypar}

The three-dimensional classification and its conventional quotient form
were checked against GSM~235, Proposition~5.5, Corollary~5.6,
Theorem~5.7, and Corollary~5.8, pp.~162--164, and independently against
\emph{Ricci Solitons in Dimensions Four and Higher}, Mathematical Surveys
and Monographs~\textbf{293}, Theorem~1.74 and Corollary~1.75, pp.~52--53.
The exact three-element cylinder quotient inventory in
Theorem~\ref{thm:exact-cylinder-soliton-quotients} is the elementary
potential-preserving refinement proved here; it is not attributed verbatim
to either source.

Five source hazards are repaired explicitly.
\begin{enumerate}
\item GSM~235, Theorem~5.7, and MSM~293, Theorem~1.74, are printed for
      ``any complete three-dimensional GRS,'' but their proof and source
      chains invoke nonnegative curvature operator.  The omission is
      material for expanders.  Theorem
      \ref{thm:nonnegative-three-shrinker-classification} states the correct
      curvature hypothesis, while Corollary
      \ref{cor:complete-three-shrinker-nonnegative} supplies it for
      shrinkers.
\item Both sources say that the two-dimensional factor in the
      \emph{universal-cover} branch may be $\Sph^2$ or
      $\mathbb{RP}^2$.  The factor produced by
      Theorem~\ref{thm:parallel-line-splitting} is simply connected and is
      therefore $\Sph^2$.  The projective plane occurs only in the
      downstairs quotient by $a$.
\item A parallel rank-one curvature image need not have a global parallel
      unit generator on the base: its real line bundle may have holonomy
      $-1$.  Chapter~\ref{ch:strong-curvature} first passes to the simply
      connected universal cover and constructs the parallel tangent line
      without choosing an orientation.  The present proof reuses that
      repaired interface.
\item In the localized proof printed in GSM~235, the cutoff is called
      ``non-decreasing'' although it equals one on $(-\infty,1]$ and zero on
      $[2,\infty)$.  It is non-increasing, as its later derivative signs and
      the primary source require.  Since the localized estimate is a named
      input here, this typographical defect is recorded rather than copied.
\item Perelman's brief Section~11.3 corollary can obscure the compactness
      step needed outside the $\kappa$-solution setting.  As B.~Kleiner and
      J.~Lott observe in \emph{Notes on Perelman's papers}, Geometry \&
      Topology~\textbf{12} (2008), Section~40, in the note following
      Corollary~40.1, the compact-surface citation does not itself justify
      the noncompact case.  Hamilton's 1995 Theorem~26.1 gives the complete
      result.  Theorem
      \ref{thm:complete-two-dimensional-shrinker-classification} supplies
      the completeness-only reduction explicitly: a non-Gaussian shrinker
      has a uniform positive Ricci lower bound and is compact by
      Bonnet--Myers, while the zero-scalar branch is Gaussian.
\end{enumerate}

Complete human proofs of the two classification theorems, the potential
split, the three rank branches, the exact cylinder quotient list, and the
closed-flow blow-up bridge appear above, relative to the precisely named
classical inputs.  The full local-cutoff proof of
Theorem~\ref{thm:complete-localized-hamilton-ivey} and a proof of the
Kazdan--Warner identity are not reproduced.  No claim is made that the new
nodes in this chapter have been implemented, compiled, kernel checked, or
independently certified in Lean.

\chapter{Bando--Shi Local Curvature-Derivative Estimates}
\label{ch:bando-shi}
\index{Bando--Shi estimates}
\index{Shi estimates}
\index{curvature!derivative estimates}
\index{Bernstein method}

\section{Purpose, theorem, and proof architecture}
\label{sec:shi-purpose}

A zeroth-order curvature bound for a Ricci flow forces bounds for every
spatial curvature derivative at every positive time.  The result is local:
neither completeness of the evolving metric nor compactness of the
underlying manifold is required.  This is the regularity theorem that turns
curvature-controlled parabolic rescalings into smooth limits and that makes
the tensor coefficients in Hamilton's differential Harnack estimate
bounded on positive-time slabs.

The compact, global precursor is due to Shigetoshi Bando.  The local theorem
proved here is due to Wan-Xiong Shi.  We use the name \emph{Bando--Shi
estimates} for the resulting regularity package, but retain Shi's attribution
for the local theorem itself.

The proof is organized to separate four kinds of reasoning.
\begin{enumerate}
\item The exact tensor identities are proved in the fixed Euclidean bundle
      of Chapter~\ref{ch:curvature-evolution}.  Every curvature slot and
      every derivative slot is pulled back.  This makes the time derivative
      compatible with a time-independent norm.
\item The tensor identities are converted into scalar norm inequalities.
      Only here are universal contractions replaced by numerical bounds.
\item A cutoff built from the \emph{initial} distance localizes the first
      Bernstein quantity.  The evolving connection, rather than an evolving
      distance, is the only time-dependent error.
\item Once the first derivative is known, the cutoff errors are uniformly
      bounded.  A finite induction with radii chosen in advance proves every
      higher-order estimate on the same half-ball.
\end{enumerate}

No changing-distance estimate is used.  This is deliberate.  Initial-distance
cutoffs give a shorter dependency chain and avoid differentiating a distance
function for an evolving metric.

\begin{roadmap}{Section~\ref{sec:shi-statement-scaling} states the theorem
in scale-invariant form.}
Sections~\ref{sec:shi-full-pullback}--\ref{sec:shi-norm-evolution} construct
the exact all-order curvature evolution.  Section~\ref{sec:shi-cutoffs}
proves the two cutoff lemmas.  Section~\ref{sec:shi-first} proves the local
first-derivative estimate, and Section~\ref{sec:shi-higher} performs the
higher-order induction.  Sections~\ref{sec:shi-global}--
\ref{sec:shi-project-interfaces} derive the complete global estimates,
mixed spacetime estimates, the closed-flow extension criterion, and the
interfaces needed by Hamilton compactness and the matrix Harnack theorem.
The final sections record formalization units and the source audit.
\end{roadmap}

\section{Scale-invariant statement}
\label{sec:shi-statement-scaling}

All balls in the local theorem are measured using the initial metric.  This
is part of the interface, not a notational convenience.

\begin{theorem}[Shi's local curvature-derivative estimates]
\label{thm:shi-local-derivatives}
Fix an integer $n\geq2$, an integer $m\geq0$, and positive real numbers
$\alpha$ and $\rho$.  There exists
\[
                 C=C(n,m,\alpha,\rho)<\infty
\]
with the following property.  Let $U$ be an $n$-manifold without boundary,
let $p\in U$, let $K>0$, and let
\[
                 g(t),\qquad 0\leq t\leq T,
\]
be a smooth Ricci flow on $U$.  Assume
\begin{align}
  0<T&\leq \frac{\alpha}{K},                              \label{eq:shi-time-hypothesis}\\
  \overline{B_{g(0)}(p,\rho K^{-1/2})}&\Subset U,          \label{eq:shi-compact-ball}\\
  |\Rm_{g(t)}|_{g(t)}(x)&\leq K
     \quad\text{for }(x,t)\in
     B_{g(0)}(p,\rho K^{-1/2})\times[0,T].                \label{eq:shi-curvature-hypothesis}
\end{align}
Then
\begin{equation}
 |\nabla_{g(t)}^{m}\Rm_{g(t)}|_{g(t)}(x)
       \leq C K t^{-m/2}                                 \label{eq:shi-main-estimate}
\end{equation}
for every
\[
 (x,t)\in B_{g(0)}(p,\tfrac12\rho K^{-1/2})\times(0,T].
\]
For $m=0$, one may take $C=1$.
\end{theorem}

The theorem assumes curvature control only on the displayed ball.  It does
not assume that $g(t)$ is complete on $U$, and it does not assume bounds on
any initial curvature derivative.  Compact containment in
\eqref{eq:shi-compact-ball} is load-bearing: it keeps every localized maximum
inside the region on which the flow and the curvature bound are available.

\begin{corollary}[One-parameter curvature-scale form]
\label{cor:shi-gsm-form}
For $n,m\geq1$ and $c>0$, there is $C=C(n,m,c)$ such that the following
holds.  If $0<T\leq cK^{-1}$, the closed ball
$\overline{B_{g(0)}(p,cK^{-1/2})}$ is compactly contained in $U$, and
$|\Rm|\leq K$ on its product with $[0,T]$, then
\[
       |\nabla^m\Rm|(x,t)\leq CKt^{-m/2}
\]
on $B_{g(0)}(p,\frac c2K^{-1/2})\times(0,T]$.
\end{corollary}

\begin{proof}
Apply Theorem~\ref{thm:shi-local-derivatives} with
$\alpha=\rho=c$.
\end{proof}

\begin{lemma}[Parabolic rescaling and derivative weights]
\label{lem:shi-scaling}
Let $K>0$ and set
\[
       \widetilde g(s)=K g(s/K),\qquad s=Kt.
\]
Then $\widetilde g$ is a Ricci flow and
\begin{align}
 d_{\widetilde g(0)}&=K^{1/2}d_{g(0)},                    \label{eq:shi-distance-scaling}\\
 |\widetilde\Rm|_{\widetilde g}&=K^{-1}|\Rm|_g,          \label{eq:shi-rm-scaling}\\
 |\widetilde\nabla^{m}\widetilde\Rm|_{\widetilde g}
   &=K^{-1-m/2}|\nabla^{m}\Rm|_g.                         \label{eq:shi-derivative-scaling}
\end{align}
Consequently, it is enough to prove
Theorem~\ref{thm:shi-local-derivatives} with $K=1$, time interval
$[0,\alpha]$, and initial radius $\rho$.
\end{lemma}

\begin{proof}
A constant rescaling does not change the Levi-Civita connection.  Distances
have weight $1$, time has weight $2$, curvature has weight $-2$, and each
covariant derivative has weight $-1$.  These statements give
\eqref{eq:shi-distance-scaling}--\eqref{eq:shi-derivative-scaling} directly.
If the normalized estimate is
\[
 |\widetilde\nabla^m\widetilde\Rm|(x,s)
       \leq C s^{-m/2},
\]
then, with $s=Kt$,
\[
 |\nabla^m\Rm|(x,t)
   \leq K^{1+m/2}C(Kt)^{-m/2}=CKt^{-m/2}.
\]
The radius becomes $K^{1/2}(\rho K^{-1/2})=\rho$.
\end{proof}

If $K=0$ in a variant of the theorem, then $\Rm$ is identically zero on the
open spacetime cylinder and every spatial derivative of $\Rm$ vanishes
there.  Thus no division by $K$ is needed in the flat case.

\section{The fully pulled Uhlenbeck derivative}
\label{sec:shi-full-pullback}

Fix the normalized situation $K=1$.  Use the metric-freezing data from
Lemma~\ref{lem:ue-metric-freezing}:
\[
 (V,h)\xrightarrow{\ \iota_t\ }(TU,g(t)),
\]
where $h$ is time-independent, every $\iota_t$ is an isometry, and $D^t$ is
the pulled-back Levi-Civita connection.  We now make one refinement that is
essential for higher derivatives: all derivative directions are also pulled
back to $V$.

For $q\geq0$, write $E_q=(V^*)^{\otimes q}$.  If
$A\in\Gamma(E_q)$, define $\mathcal DA\in\Gamma(E_{q+1})$ by
\begin{equation}
 (\mathcal DA)(v_0,v_1,\ldots,v_q)
   =\bigl(D^t_{\iota_t v_0}A\bigr)(v_1,\ldots,v_q).       \label{eq:shi-full-D}
\end{equation}
Define $\mathcal D^k$ recursively and set
\begin{equation}
       \Delta^U A=\operatorname{tr}_h(\mathcal D^2A),
       \qquad L=\partial_t-\Delta^U.                     \label{eq:shi-fixed-heat-operator}
\end{equation}
The superscript $U$ means ``Uhlenbeck gauge,'' not the open manifold $U$.

Let $\mathbf R(t)$ denote the covariant $(0,4)$ curvature tensor in the
convention of Chapter~\ref{ch:curvature-evolution}, and set
\begin{equation}
 A_0=\iota_t^*\mathbf R(t),
 \qquad A_k=\mathcal D^kA_0\in\Gamma(E_{k+4}).            \label{eq:shi-Ak-definition}
\end{equation}

\begin{lemma}[Every slot is pulled]
\label{lem:shi-every-slot-pulled}
For every $k\geq0$,
\begin{equation}
 A_k(v_1,\ldots,v_{k+4})
  =(\nabla^k\mathbf R)
       (\iota_t v_1,\ldots,\iota_t v_{k+4}).             \label{eq:shi-Ak-naturality}
\end{equation}
Moreover,
\begin{equation}
 |A_k|_h=|\nabla^k\Rm|_{g(t)},
 \qquad |\mathcal DA_k|_h=|A_{k+1}|_h.                  \label{eq:shi-Ak-norm-bridge}
\end{equation}
\end{lemma}

\begin{proof}
The case $k=0$ is the definition of pullback.  If the identity holds for
$k$, apply the spatial naturality statement of
Lemma~\ref{lem:ue-spatial-naturality} in the direction $\iota_t v_0$.
Equation~\eqref{eq:shi-full-D} pulls this new derivative direction back as
the first slot.  This proves the tensor identity by induction.  Since
$\iota_t$ is an isometry in every slot, it also proves both norm identities.
\end{proof}

Pulling back only the original four curvature slots would not prove
\eqref{eq:shi-Ak-norm-bridge}: the $k$ derivative slots would still carry
the evolving inverse metric.  Definition~\eqref{eq:shi-Ak-definition}
removes that ambiguity.

There is an equivalent moving-bundle notation that will be useful later.
For a covariant $q$-tensor $T(t)$ on $U$, define
\begin{equation}
 (\nabla_tT)(X_1,\ldots,X_q)
   =\partial_tT(X_1,\ldots,X_q)
      +\sum_{r=1}^q
       T(X_1,\ldots,\Ric^\sharp X_r,\ldots,X_q),         \label{eq:shi-metric-time-derivative}
\end{equation}
where the vector fields $X_r$ are held fixed while differentiating the
components.  Then $\nabla_tg=0$ and
\begin{equation}
       \partial_t(\iota_t^*T)=\iota_t^*(\nabla_tT).       \label{eq:shi-time-pullback-bridge}
\end{equation}
Thus $\nabla_t$ is simply ordinary differentiation after full Uhlenbeck
pullback.  Formula~\eqref{eq:shi-metric-time-derivative} is used only as a
bridge; all estimates below are carried out with the ordinary derivative in
the fixed bundle.

\section{An exact one-derivative commutator}
\label{sec:shi-exact-commutator}

The principal simplification of Uhlenbeck gauge is that the spatial and
temporal connection errors cancel before any norm estimate is taken.

Fix a point and an $h$-orthonormal basis.  Let
$\rho^{(q)}:\mathfrak{so}(V,h)\to\operatorname{End}(E_q)$ be the canonical
infinitesimal representation on covariant $q$-tensors.  Its component
$\rho_{cd}^{(q)}$ at the skew generator in the $c,d$ plane is characterized
by
\begin{equation}
 (\rho_{cd}^{(q)}A)_{b_1\ldots b_q}
  =\sum_{r=1}^q
    \left(
      \delta_{c b_r}A_{b_1\ldots d\ldots b_q}
      -\delta_{d b_r}A_{b_1\ldots c\ldots b_q}
    \right).                                             \label{eq:shi-rho-action}
\end{equation}
The individual labeled component $\rho_{cd}^{(q)}$ changes equivariantly
when the orthonormal basis changes.  The representation $\rho^{(q)}$ is
intrinsic, and the fully contracted expression in
\eqref{eq:shi-rho-commutator} is basis-independent.

\begin{lemma}[Exact Uhlenbeck heat commutator]
\label{lem:shi-exact-commutator}
Let $A(t)\in\Gamma(E_q)$ be smooth.  Then
\begin{equation}
 [L,\mathcal D_a]A
    =A_{0,aedc}\,\mathcal D_e(\rho_{cd}^{(q)}A).          \label{eq:shi-rho-commutator}
\end{equation}
Equivalently,
\begin{equation}
 \bigl(L\mathcal DA-\mathcal DLA\bigr)_{a b_1\ldots b_q}
  =2\sum_{r=1}^q
       A_{0,aedb_r}
       (\mathcal D_eA)_{b_1\ldots d\ldots b_q}.          \label{eq:shi-expanded-commutator}
\end{equation}
In particular, the right side is a universal contraction of $A_0$ with
$\mathcal DA$; there is no term of the form
$\mathcal DA_0\otimes A$.
\end{lemma}

\begin{proof}
We give the cancellation rather than hiding it in star notation.  In an
orthonormal frame at the point, the connection-variation formula, the motion
of the pulled derivative direction, and the Ricci commutator give
\begin{align}
 [\partial_t,\mathcal D_a]A
   &=(\mathcal D_b\Ric_{ac})\rho_{bc}^{(q)}A
        +\Ric_{ac}\mathcal D_cA,                         \label{eq:shi-time-D-commutator}\\
 [\Delta^U,\mathcal D_a]A
   &=A_{0,eadc}\mathcal D_e(\rho_{cd}^{(q)}A)
       +\Ric_{ad}\mathcal D_dA
       +(\mathcal D_b\Ric_{ac})\rho_{bc}^{(q)}A.         \label{eq:shi-space-D-commutator}
\end{align}
For completeness, the first identity follows by differentiating the pulled
Christoffel symbols and using
\[
 \partial_t\Gamma^k_{ij}
   =-g^{k\ell}
      (\nabla_i\Ric_{j\ell}+\nabla_j\Ric_{i\ell}
                         -\nabla_\ell\Ric_{ij});
\]
the second follows by commuting the three spatial derivatives and using the
contracted second Bianchi identity.  Both are tensorial, so calculation in
the chosen frame suffices.

Subtract \eqref{eq:shi-space-D-commutator} from
\eqref{eq:shi-time-D-commutator}.  The $\mathcal D\Ric$ terms cancel, as do
the Ricci multiples of $\mathcal DA$.  Antisymmetry in the first curvature
pair changes $-A_{0,eadc}$ into $A_{0,aedc}$ and gives
\eqref{eq:shi-rho-commutator}.  Finally expand
\eqref{eq:shi-rho-action}.  The two terms belonging to each tensor slot are
equal after antisymmetry in the last curvature pair; their sum gives the
factor $2$ in \eqref{eq:shi-expanded-commutator}.
\end{proof}

Two regression tests help fix the signs.  When $q=0$, the right side is
zero, as it must be for the fully pulled differential of a scalar.  When
$q=1$,
\[
 (L\mathcal DA-\mathcal DLA)_{ab}
       =2A_{0,aedb}\mathcal D_eA_d.
\]
This is the one-form formula from which the general slot sum can also be
reconstructed.

\section{Universal contraction maps in every order}
\label{sec:shi-universal-maps}

Proposition~\ref{prop:ue-pulled-tensor-evolution} gives an exact homogeneous
quadratic map $\mathcal N:E_4\to E_4$ such that
\begin{equation}
                       LA_0=\mathcal N(A_0).             \label{eq:shi-base-evolution-N}
\end{equation}
Let $\mathcal Q:E_4\times E_4\to E_4$ be its symmetric polarization,
\begin{equation}
 \mathcal Q(P,S)
   =\frac12\bigl(\mathcal N(P+S)-\mathcal N(P)-\mathcal N(S)\bigr).  \label{eq:shi-Q-polarization}
\end{equation}
Then $\mathcal N(P)=\mathcal Q(P,P)$.  The map $\mathcal Q$ is a fixed
finite linear combination of permutations and $h$-contractions.

If $B:E_p\times E_q\to E_r$ is any parallel bilinear contraction map,
define its derivative lifts $d_1B$ and $d_2B$ by the exact Leibniz rule
\begin{equation}
 \mathcal D(B(X,Y))
     =(d_1B)(\mathcal DX,Y)+(d_2B)(X,\mathcal DY).        \label{eq:shi-derivative-lifts}
\end{equation}
More explicitly, if $I$ is the output multi-index and the bullet denotes the
remaining input slots, then
\begin{equation}
 \bigl((d_1B)(X',Y)\bigr)_{aI}=B(X'_{a\,\bullet},Y)_I,
 \qquad
 \bigl((d_2B)(X,Y')\bigr)_{aI}=B(X,Y'_{a\,\bullet})_I.   \label{eq:shi-derivative-lift-slots}
\end{equation}
Thus both lifts prepend the new derivative slot, consistently with
Definition~\eqref{eq:shi-full-D}.
For $q\geq0$, define
\begin{equation}
 \mathcal C_q(P,Z)_{a b_1\ldots b_q}
   =2\sum_{r=1}^q
       P_{aedb_r}Z_{e,b_1\ldots d\ldots b_q},            \label{eq:shi-Cq-definition}
\end{equation}
where $P\in E_4$ and $Z\in E_{q+1}$.

We now define actual maps rather than a schematic product symbol.

\begin{definition}[Recursive curvature-contraction maps]
\label{def:shi-Bki}
For $0\leq i\leq k$, define parallel bilinear maps
\[
 \mathcal B_{k,i}:E_{i+4}\times E_{k-i+4}\longrightarrow E_{k+4}
\]
recursively by
\begin{align}
 \mathcal B_{0,0}&=\mathcal Q,                            \label{eq:shi-B00}\\
 \mathcal B_{k+1,0}&=d_2\mathcal B_{k,0}+\mathcal C_{k+4},\label{eq:shi-B-left}\\
 \mathcal B_{k+1,j}&=d_1\mathcal B_{k,j-1}
                       +d_2\mathcal B_{k,j}
          \quad(1\leq j\leq k),                          \label{eq:shi-B-middle}\\
 \mathcal B_{k+1,k+1}&=d_1\mathcal B_{k,k}.              \label{eq:shi-B-right}
\end{align}
In \eqref{eq:shi-B-left}, $\mathcal C_{k+4}$ has exactly the same domain
and codomain as $d_2\mathcal B_{k,0}$.
\end{definition}

\begin{theorem}[Exact all-order curvature evolution]
\label{thm:shi-exact-all-order-evolution}
For every integer $k\geq0$,
\begin{equation}
       LA_k=\sum_{i=0}^k
          \mathcal B_{k,i}(A_i,A_{k-i}).                 \label{eq:shi-exact-Ak-evolution}
\end{equation}
Every summand in \eqref{eq:shi-exact-Ak-evolution} is a specified universal
$O(n)$-equivariant bilinear map.
\end{theorem}

\begin{proof}
For $k=0$, this is \eqref{eq:shi-base-evolution-N} and
\eqref{eq:shi-B00}.  Assume it at order $k$.  Since
$A_{k+1}=\mathcal DA_k$, Lemma~\ref{lem:shi-exact-commutator} gives
\[
 LA_{k+1}
   =\mathcal D(LA_k)+\mathcal C_{k+4}(A_0,A_{k+1}).
\]
Apply the exact Leibniz rule \eqref{eq:shi-derivative-lifts} to every term in
the inductive sum.  Terms in which the derivative lands on the first factor
give $d_1\mathcal B_{k,i}$; terms in which it lands on the second give
$d_2\mathcal B_{k,i}$.  Sorting by the order of the first factor produces
precisely \eqref{eq:shi-B-left}--\eqref{eq:shi-B-right}.
\end{proof}

The recursive definition has two advantages for formalization.  It makes
the induction literally an equality, and it postpones every estimate until
the finite-dimensional maps have been constructed.  If
\[
 \beta_{k,i}(n)=\|\mathcal B_{k,i}\|_{\mathrm{op}},
\]
then
\begin{equation}
 |LA_k|\leq\sum_{i=0}^k
       \beta_{k,i}(n)|A_i||A_{k-i}|.                     \label{eq:shi-Ak-map-bound}
\end{equation}
The constants exist because the fibers are finite-dimensional.  One may
also obtain explicit coarse constants: $\|\mathcal Q\|_{\mathrm{op}}\leq8$,
$\|\mathcal C_q\|_{\mathrm{op}}\leq2q$, and hence a bound for
$\sum_i\beta_{k,i}$ is generated by
\[
       \gamma_0=8,\qquad
       \gamma_{k+1}=2\gamma_k+2(k+4).
\]
For example, $\gamma_k=18\cdot2^k-2k-10$ is admissible.  None of the later
arguments needs this particular numerical choice.

\section{Exact norm evolution and its scalar consequence}
\label{sec:shi-norm-evolution}

Put $u_k=|A_k|_h^2$.  Because $h$ is fixed in time and $\mathcal D$ is
$h$-compatible,
\begin{align}
 \partial_t|A|_h^2&=2\langle\partial_tA,A\rangle_h,       \label{eq:shi-fixed-time-norm}\\
 \Delta^U|A|_h^2
   &=2\langle\Delta^UA,A\rangle_h+2|\mathcal DA|_h^2.    \label{eq:shi-fixed-space-norm}
\end{align}
Combining these identities with
Theorem~\ref{thm:shi-exact-all-order-evolution} gives the following equality.

\begin{proposition}[Exact all-order norm identity]
\label{prop:shi-exact-norm-identity}
For every $k\geq0$,
\begin{equation}
 Lu_k=-2u_{k+1}
    +2\sum_{i=0}^k
       \left\langle
          \mathcal B_{k,i}(A_i,A_{k-i}),A_k
       \right\rangle_h.                                 \label{eq:shi-exact-norm-evolution}
\end{equation}
Consequently, for a constant $c_k=c(n,k)$,
\begin{equation}
 Lu_k\leq-2u_{k+1}
      +c_k\sum_{i=0}^k
          u_i^{1/2}u_{k-i}^{1/2}u_k^{1/2}.               \label{eq:shi-scalar-norm-evolution}
\end{equation}
\end{proposition}

\begin{proof}
Subtract \eqref{eq:shi-fixed-space-norm} from
\eqref{eq:shi-fixed-time-norm}, use
$\mathcal DA_k=A_{k+1}$, and substitute
\eqref{eq:shi-exact-Ak-evolution}.  This proves the equality.  Apply
Cauchy--Schwarz and the operator-norm bounds $\beta_{k,i}(n)$ to obtain the
inequality.
\end{proof}

Under the isometric identification in
Lemma~\ref{lem:shi-every-slot-pulled},
\eqref{eq:shi-exact-norm-evolution} is also the evolution of
$|\nabla^k\Rm|_{g(t)}^2$.  Thus the terms caused by differentiating the
evolving inverse metrics have not been discarded: they have already been
cancelled by the Ricci-slot terms in
\eqref{eq:shi-time-pullback-bridge}.

At the first two orders, \eqref{eq:shi-scalar-norm-evolution} says that
there are constants $c_0,c_1$ depending only on $n$ such that
\begin{align}
 Lu_0&\leq-2u_1+c_0u_0^{3/2},                            \label{eq:shi-u0-evolution}\\
 Lu_1&\leq-2u_2+c_1u_0^{1/2}u_1.                        \label{eq:shi-u1-evolution}
\end{align}
The endpoint terms $i=0$ and $i=k$ in
\eqref{eq:shi-scalar-norm-evolution} are essential; they produce the
$|\Rm||\nabla^k\Rm|^2$ reaction.

\section{Fixed-initial-distance cutoffs}
\label{sec:shi-cutoffs}

Distance from a point is not smooth on its cut locus.  We state precisely
the weak sense in which the radial cutoff is used.

\begin{definition}[Lower-support Laplacian inequality]
\label{def:shi-lower-support}
Let $\eta$ be continuous near $x$ and let $g$ be smooth.  We say that
\[
                         -\Delta_g\eta(x)\leq\Lambda
\]
in the lower-support sense if for every $\varepsilon>0$ there is a smooth
function $\eta_\varepsilon$ near $x$ such that
\[
 \eta_\varepsilon\leq\eta,
 \qquad \eta_\varepsilon(x)=\eta(x),
 \qquad -\Delta_g\eta_\varepsilon(x)\leq\Lambda+\varepsilon.
\]
Any accompanying gradient bound is required of the same support functions
at $x$.
\end{definition}

This is the correct inequality direction for the maximum argument.  If
$\eta F$ has a local maximum at $x$ and $F(x)\geq0$, then
$\eta_\varepsilon F\leq\eta F$ with equality at $x$, so the classical
differential test applies to the smooth product
$\eta_\varepsilon F$.

We first construct the cutoff needed before $|\nabla\Rm|$ is known.  Work in
the normalization $K=1$.

\begin{lemma}[Self-coupled initial-distance cutoff]
\label{lem:shi-self-coupled-cutoff}
Let $g(t)$, $0\leq t\leq\tau\leq\alpha$, be a Ricci flow on $U$, and assume
\[
 |\Rm|\leq1
 \quad\text{on }B_{g(0)}(p,R)\times[0,\tau],
 \qquad \overline{B_{g(0)}(p,R)}\Subset U.
\]
Fix $0<r_1<r_2<R$ and define
\begin{equation}
             \Theta(x,t)=t|\nabla\Ric|^2(x,t),
             \qquad \Theta(x,0)=0.                       \label{eq:shi-Theta}
\end{equation}
There are constants $A,B,D<\infty$, depending only on
$n,\alpha,R,r_1,r_2$, and a time-independent locally Lipschitz function
$\eta:U\to[0,1]$ such that
\begin{align}
 \eta&=1\text{ on }B_{g(0)}(p,r_1),
 &\operatorname{supp}\eta&\Subset B_{g(0)}(p,r_2),       \label{eq:shi-cutoff-support}\\
 |\nabla\eta|_{g(t)}^2&\leq A\eta,                        \label{eq:shi-cutoff-gradient}\\
 -\Delta_{g(t)}\eta(x)&\leq
       B+D\sup_{0\leq s\leq t}
             (\eta(x)\Theta(x,s))^{1/2}.                 \label{eq:shi-cutoff-laplacian}
\end{align}
The last two inequalities hold in the lower-support sense at cut-locus
points.
\end{lemma}

\begin{proof}
Put $r_*=(r_1+r_2)/2$.  Choose a smooth nonincreasing
$\lambda:[0,\infty)\to[0,1]$ which is one on $[0,r_1]$ and zero on
$[r_*,\infty)$, and put $\kappa=\lambda^2$.  Its
constants can be chosen so that
\begin{equation}
       |\kappa'|^2\leq C\kappa,
       \qquad |\kappa''|\leq C,                           \label{eq:shi-kappa-bounds}
\end{equation}
where $C$ depends only on $r_2-r_1$.  Set
\[
              d_0(x)=d_{g(0)}(p,x),
              \qquad\eta(x)=\kappa(d_0(x)).
\]

Since $|\Ric|\leq c_n$ and $\partial_tg=-2\Ric$, integration along a fixed
tangent vector gives
\begin{equation}
 e^{-2c_n\alpha}g(0)\leq g(t)\leq e^{2c_n\alpha}g(0).     \label{eq:shi-local-metric-comparison}
\end{equation}
Together with \eqref{eq:shi-kappa-bounds}, this proves
\eqref{eq:shi-cutoff-gradient}.

On the transition annulus, $d_0\geq r_1$.  The lower sectional-curvature
bound $\sec_{g(0)}\geq-1$, Hessian comparison, and
\eqref{eq:shi-kappa-bounds} imply
\begin{equation}
                  \operatorname{Hess}_{g(0)}\eta
                       \geq-Cg(0)                         \label{eq:shi-initial-hessian-cutoff}
\end{equation}
away from the cut locus.  Let
$S(t)=\nabla^{g(t)}-\nabla^{g(0)}$.  The connection-variation formula gives,
after using \eqref{eq:shi-local-metric-comparison},
\begin{equation}
 |S(t)|(x)\leq C\int_0^t|\nabla\Ric|(x,s)\,ds.           \label{eq:shi-connection-difference}
\end{equation}
Because $|d\eta|\leq C\eta^{1/2}$,
\begin{align}
 |d\eta|\,|S(t)|
  &\leq C\int_0^t
       s^{-1/2}\bigl(\eta\Theta(x,s)\bigr)^{1/2}\,ds \\
  &\leq2C\sqrt\alpha
       \sup_{0\leq s\leq t}
          \bigl(\eta\Theta(x,s)\bigr)^{1/2}.             \label{eq:shi-connection-cutoff-error}
\end{align}
Now
\[
 \operatorname{Hess}_{g(t)}\eta
   =\operatorname{Hess}_{g(0)}\eta-S(t)*d\eta.
\]
Trace this equality with $g(t)$ and use
\eqref{eq:shi-local-metric-comparison},
\eqref{eq:shi-initial-hessian-cutoff}, and
\eqref{eq:shi-connection-cutoff-error}.  This proves
\eqref{eq:shi-cutoff-laplacian} away from the cut locus.

At a cut point $x$, choose a minimizing $g(0)$-geodesic $\gamma$ from $p$
to $x$, let $p_\varepsilon=\gamma(\varepsilon)$, and set
\[
 d_\varepsilon(y)=\varepsilon+d_{g(0)}(p_\varepsilon,y),
 \qquad \eta_\varepsilon(y)=\kappa(d_\varepsilon(y)).
\]
The triangle inequality gives $d_\varepsilon\geq d_0$, with equality at
$x$.  Since $\kappa$ is nonincreasing,
$\eta_\varepsilon\leq\eta$, with equality at $x$.  For sufficiently small
$\varepsilon$, the standard Calabi support construction makes
$d_\varepsilon$ smooth near $x$; the comparison bounds above are uniform as
$\varepsilon\downarrow0$.  This supplies the required lower supports.  Near
$p$, the cutoff is identically one and there is no singularity.
\end{proof}

\begin{corollary}[Clean cutoff after the first estimate]
\label{cor:shi-clean-cutoff}
In the setting of Lemma~\ref{lem:shi-self-coupled-cutoff}, suppose in
addition that
\[
       |\nabla\Rm|(x,t)\leq Ht^{-1/2}
\]
on a ball containing $\operatorname{supp}\eta$, for $0<t\leq\tau$.
Then the cutoff may be chosen so that
\begin{equation}
             |\nabla\eta|^2\leq A\eta,
             \qquad -\Delta\eta\leq B'                  \label{eq:shi-clean-cutoff-bounds}
\end{equation}
in the lower-support sense, where $B'$ also depends on $H$.
\end{corollary}

\begin{proof}
Contraction gives $|\nabla\Ric|\leq c_n|\nabla\Rm|$, and hence
$\Theta\leq c_n^2H^2$.  Substitute this bound into
\eqref{eq:shi-cutoff-laplacian}.
\end{proof}

The next scalar lemma isolates the maximum-principle step used at every
higher order.

\begin{lemma}[Compactly supported Bernstein maximum]
\label{lem:shi-compact-bernstein-maximum}
Let $0<\tau<\infty$, let $F\geq0$ be smooth for $t>0$ on a spacetime
neighborhood of a compact set, and suppose $F$ extends continuously by
$F(\cdot,0)=0$.  Assume
\begin{equation}
       LF\leq-\frac c tF^2+\frac Ct                       \label{eq:shi-abstract-F-evolution}
\end{equation}
for constants $c>0$ and $C\geq0$.  Let $\eta\in[0,1]$ have compact support,
be one on a prescribed inner set, and satisfy
\[
       |\nabla\eta|^2\leq A\eta,
       \qquad -\Delta\eta\leq B
\]
in the lower-support sense.  Then
\begin{equation}
        \sup \eta F\leq C'(c,C,A,B,\tau).                \label{eq:shi-abstract-local-bound}
\end{equation}
\end{lemma}

\begin{proof}
Let $Q=\eta F$.  It vanishes at $t=0$ and near the spatial boundary of the
cutoff support.  If its maximum is positive, it is attained at a point
$(x_0,t_0)$ with $t_0>0$.  Use a smooth lower support for $\eta$ if needed.
At that point
\[
 \nabla F=-\frac F\eta\nabla\eta,
 \qquad LQ\geq0.
\]
Multiplying the product evolution by $t_0\eta$ gives
\begin{align*}
 0&\leq-cQ^2+C\eta^2-t_0(\Delta\eta)Q
                  +2t_0F|\nabla\eta|^2\\
  &\leq-cQ^2+C+\tau(B+2A)Q.
\end{align*}
The positive root of this quadratic bounds $Q(x_0,t_0)$ and hence all of
$Q$.
\end{proof}

\begin{lemma}[Polynomial absorption]
\label{lem:shi-polynomial-absorption}
If $z\geq0$ and
\[
       cz^2\leq az^{3/2}+bz+d
\]
with $c>0$ and $a,b,d\geq0$, then $z$ is bounded by a number depending only
on $a,b,c,d$.
\end{lemma}

\begin{proof}
Put $y=\sqrt z$.  If
\[
 y\geq\max\left\{\frac{4a}{c},
                   \sqrt{\frac{4b}{c}},
                   \left(\frac{4d}{c}\right)^{1/4}\right\},
\]
then each term on the right of $cy^4\leq ay^3+by^2+d$ is at most
$cy^4/4$, a contradiction.  The displayed maximum therefore bounds $y$
and hence $z$.
\end{proof}

\section{The local first-derivative estimate}
\label{sec:shi-first}

We continue in the normalization $K=1$.  The proof uses a product in which
one negative Hessian term controls the gradient cross term and one negative
gradient term produces a quadratic reaction.

\begin{proposition}[Normalized local first derivative]
\label{prop:shi-normalized-first}
Assume
\[
 0<\tau\leq\alpha,
 \qquad |\Rm|\leq1
   \text{ on }B_{g(0)}(p,R)\times[0,\tau],
 \qquad \overline{B_{g(0)}(p,R)}\Subset U.
\]
For each $0<r<R$, there is
$C=C(n,\alpha,R,r)$ such that
\begin{equation}
       |\nabla\Rm|(x,t)\leq Ct^{-1/2}                    \label{eq:shi-normalized-first}
\end{equation}
on $B_{g(0)}(p,r)\times(0,\tau]$.
\end{proposition}

\begin{proof}
Write
\[
       u=|A_0|^2,
       \qquad v=|A_1|^2,
       \qquad w=|A_2|^2.
\]
Thus $0\leq u\leq1$.  Equations
\eqref{eq:shi-u0-evolution} and \eqref{eq:shi-u1-evolution} give
\begin{equation}
       Lu\leq-2v+c_0,
       \qquad Lv\leq-2w+c_1v.                            \label{eq:shi-uvw}
\end{equation}
Choose a dimensional constant $a\geq32$ and set
\begin{equation}
                  F=(a+u)v.                              \label{eq:shi-first-F}
\end{equation}
The product rule and \eqref{eq:shi-uvw} give
\begin{align}
 LF
  &\leq-2v^2+c_0v-2(a+u)w+c_1(a+u)v
        +2|\nabla u|\,|\nabla v|.                        \label{eq:shi-first-F-pre}
\end{align}
Metric compatibility gives
\[
       |\nabla u|\leq2u^{1/2}v^{1/2}\leq2v^{1/2},
       \qquad |\nabla v|\leq2v^{1/2}w^{1/2}.
\]
Therefore
\[
       2|\nabla u|\,|\nabla v|
          \leq8v\sqrt w
          \leq aw+\frac{16}{a}v^2.
\]
The remaining linear multiple of $v$ is absorbed into another half of
$v^2$ plus a dimensional constant.  Increasing constants if necessary,
\begin{equation}
               LF\leq-cv^2+C
                    \leq-c'F^2+C,                        \label{eq:shi-first-F-good}
\end{equation}
because $av\leq F\leq(a+1)v$.

Put
\begin{equation}
                         G=tF.                            \label{eq:shi-first-G}
\end{equation}
Since $0<t\leq\alpha$, \eqref{eq:shi-first-F-good} and absorption of the
linear term imply
\begin{equation}
                LG\leq\frac1t(-c_2G^2+C_2),              \label{eq:shi-first-G-good}
\end{equation}
where $c_2>0$ and $C_2<\infty$ depend only on $n$ and $\alpha$.

Put $r_2=(r+R)/2$ and take the cutoff from
Lemma~\ref{lem:shi-self-coupled-cutoff}, equal to one on $B_{g(0)}(p,r)$.
Let
\[
        Z_*=\max_{\overline{B_{g(0)}(p,r_2)}\times[0,\tau]}\eta G.
\]
The support is compactly contained in $U$, and $G(\cdot,0)=0$.  If
$Z_*>0$, a maximum occurs at $(x_0,t_0)$ with $t_0>0$.  A lower support for
$\eta$ is used if $x_0$ lies on the cut locus.

For $0\leq s\leq t_0$, maximality and contraction of the derivative of the
Ricci tensor give
\begin{align}
 \eta(x_0)\Theta(x_0,s)
   &\leq c_n\eta(x_0)s|\nabla\Rm|^2(x_0,s)\\
   &\leq \frac{c_n}{a}\eta(x_0)G(x_0,s)
    \leq \frac{c_n}{a}Z_*.                               \label{eq:shi-Theta-closes}
\end{align}
At the maximum of $Z=\eta G$,
\[
             \nabla G=-\frac G\eta\nabla\eta,
             \qquad LZ\geq0.
\]
Multiply the product evolution by $t_0\eta$.  Using
\eqref{eq:shi-first-G-good}, the cutoff estimates, and
\eqref{eq:shi-Theta-closes}, one obtains
\begin{align}
 0
 &\leq-c_2Z_*^2+C_2\eta^2
       -t_0(\Delta\eta)Z_*+2t_0G|\nabla\eta|^2\\
 &\leq-c_2Z_*^2+C_3+C_4Z_*+C_5Z_*^{3/2}.                 \label{eq:shi-first-polynomial}
\end{align}
Lemma~\ref{lem:shi-polynomial-absorption} now gives
$Z_*\leq C(n,\alpha,R,r)$.  On $B_{g(0)}(p,r)$, $\eta=1$, and therefore
\[
             at|\nabla\Rm|^2\leq G\leq Z_*.
\]
Taking square roots proves \eqref{eq:shi-normalized-first}.
\end{proof}

The proof contains no bound on $|\nabla\Rm|(\cdot,0)$.  The factor $t$ in
$G$ forces the localized quantity to vanish at the initial face, while the
quadratic term in \eqref{eq:shi-first-G-good} controls its positive-time
maximum.

\begin{corollary}[Unnormalized first derivative]
\label{cor:shi-local-first}
Under the hypotheses of Theorem~\ref{thm:shi-local-derivatives},
\[
       |\nabla\Rm|(x,t)
          \leq C(n,\alpha,\rho)Kt^{-1/2}
\]
on $B_{g(0)}(p,\frac12\rho K^{-1/2})\times(0,T]$.
\end{corollary}

\begin{proof}
Apply Proposition~\ref{prop:shi-normalized-first} with normalized outer
radius $R=\rho$, inner radius $r=\rho/2$, and then use
Lemma~\ref{lem:shi-scaling}.
\end{proof}

\section{The all-orders Bernstein induction}
\label{sec:shi-higher}

We first isolate the higher-order algebra.  Work on a normalized spacetime
region on which $|A_0|\leq1$.

\begin{lemma}[Weighted higher-order Bernstein product]
\label{lem:shi-higher-bernstein-product}
Fix $m\geq1$.  Suppose that, on the region under consideration,
\begin{equation}
             |A_j|\leq A_j^*t^{-j/2}
             \quad(1\leq j\leq m)                        \label{eq:shi-induction-hypothesis}
\end{equation}
for fixed constants $A_j^*$.  Define
\begin{equation}
 X=t^m|A_m|^2,
 \qquad Y=t^{m+1}|A_{m+1}|^2.                            \label{eq:shi-XY}
\end{equation}
Set
\begin{equation}
                         a=1+4(A_m^*)^2                  \label{eq:shi-a-choice}
\end{equation}
and define
\begin{equation}
 F_m=(a+X)Y
   =\bigl(a+t^m|A_m|^2\bigr)
       t^{m+1}|A_{m+1}|^2.                               \label{eq:shi-Fm}
\end{equation}
If $0<t\leq\alpha$, then there are constants $c_m>0$ and $C_m<\infty$,
depending only on $n,m,\alpha,A_1^*,\ldots,A_m^*$, such that
\begin{equation}
                    LF_m\leq-\frac{c_m}{t}F_m^2
                                   +\frac{C_m}{t}.        \label{eq:shi-Fm-good}
\end{equation}
\end{lemma}

\begin{proof}
Apply \eqref{eq:shi-scalar-norm-evolution} at orders $m$ and $m+1$.  In
each reaction sum, separate the two endpoint terms from the interior terms.
For every interior factor, use
\eqref{eq:shi-induction-hypothesis}; then apply Young's inequality to the
single remaining factor $|A_{m+1}|$.  Before collecting terms, the two
endpoint reactions at order $m+1$ satisfy
\begin{equation}
 t^{m+1}|A_0||A_{m+1}|^2
       \leq CY\leq C\alpha\frac Yt.                      \label{eq:shi-LY-endpoints}
\end{equation}
For $1\leq i\leq m$, every interior reaction satisfies, for a fixed
$\varepsilon>0$,
\begin{align}
 t^{m+1}|A_i||A_{m+1-i}||A_{m+1}|
   &\leq C Y^{1/2}\\
   &\leq\varepsilon\frac Yt+C_\varepsilon t.             \label{eq:shi-LY-interior}
\end{align}
At order $m$, including its endpoints,
\begin{equation}
 t^m\sum_{i=0}^m|A_i||A_{m-i}||A_m|\leq C.              \label{eq:shi-LX-reactions}
\end{equation}
The derivatives of the time weights contribute $(m+1)Y/t$ and $mX/t$.
Consequently,
\begin{align}
 LY&\leq-2t^{m+1}|A_{m+2}|^2+C\frac Yt+Ct,                \label{eq:shi-LY}\\
 LX&\leq-2\frac Yt+m\frac Xt+C.                          \label{eq:shi-LX}
\end{align}

The product rule gives
\begin{equation}
 LF_m=(a+X)LY+Y\,LX-2\langle\nabla X,\nabla Y\rangle.   \label{eq:shi-LFm-product}
\end{equation}
Since
\begin{equation}
        |\nabla|A_j|^2|
            \leq2|A_j||A_{j+1}|,                         \label{eq:shi-gradient-uj}
\end{equation}
the decisive three terms on the right of
\eqref{eq:shi-LFm-product} are at most
\begin{align}
 &-10t^{2m+1}|A_m|^2|A_{m+2}|^2
   +8t^{2m+1}|A_m||A_{m+1}|^2|A_{m+2}|\notag\\
 &\hspace{13em}-2t^{2m+1}|A_{m+1}|^4.                   \label{eq:shi-three-leading-terms}
\end{align}
Indeed, $a\geq4X$ makes $a+X\geq5X$, producing the coefficient $-10$;
\eqref{eq:shi-gradient-uj} produces the coefficient $8$; and the negative
term in \eqref{eq:shi-LX} produces the coefficient $-2$.

Put
\[
 x=t^{m+1/2}|A_m||A_{m+2}|,
 \qquad y=t^{m+1/2}|A_{m+1}|^2.
\]
The elementary identity
\begin{equation}
           -10x^2+8xy-2y^2
                \leq-\frac25y^2                         \label{eq:shi-quadratic-absorption}
\end{equation}
turns \eqref{eq:shi-three-leading-terms} into
$-(2/5)Y^2/t$.  All remaining terms from
\eqref{eq:shi-LY}--\eqref{eq:shi-LX} are bounded by
$CY/t+Ct$.  Therefore
\begin{equation}
             LF_m\leq-\frac{2}{5t}Y^2+C\frac Yt+Ct.      \label{eq:shi-LFm-Y}
\end{equation}
Finally,
\[
       a\leq a+X\leq a+(A_m^*)^2,
\]
so $F_m$ and $Y$ are comparable by fixed positive constants.  Absorb the
linear $Y/t$ term into the negative square, and use $t\leq\alpha$ to bound
$Ct$ by $C\alpha^2/t$.  This gives \eqref{eq:shi-Fm-good}.
\end{proof}

\begin{proof}[Proof of Theorem~\ref{thm:shi-local-derivatives}]
The case $m=0$ is the hypothesis.  Suppose $m=M\geq1$ is fixed.  By
Lemma~\ref{lem:shi-scaling}, it suffices to work with $K=1$, outer radius
$R=\rho$, and $0<T\leq\alpha$.

Choose in advance the finite radius sequence
\begin{equation}
       R_j=R\left(1-\frac{j}{2M}\right),
       \qquad 0\leq j\leq M.                             \label{eq:shi-predetermined-radii}
\end{equation}
Thus $R_0=R$, $R_M=R/2$, and every gap is $R/(2M)$.
Proposition~\ref{prop:shi-normalized-first}, with an initial-distance cutoff
equal to one on $B_{g(0)}(p,R_1)$ and supported strictly inside
$B_{g(0)}(p,R_0)$, gives
\begin{equation}
             |A_1|\leq C_1t^{-1/2}
       \quad\text{on }B_{g(0)}(p,R_1)\times(0,T].        \label{eq:shi-first-radius}
\end{equation}

Assume inductively that for some $1\leq j<M$ and every $1\leq\ell\leq j$,
\begin{equation}
       |A_\ell|\leq C_\ell t^{-\ell/2}
       \quad\text{on }B_{g(0)}(p,R_j)\times(0,T].        \label{eq:shi-radius-induction}
\end{equation}
Apply Lemma~\ref{lem:shi-higher-bernstein-product} with $m=j$ on this ball.
By \eqref{eq:shi-first-radius}, the first-derivative bound is already known
on every later, smaller ball.  Corollary~\ref{cor:shi-clean-cutoff} therefore
supplies a cutoff supported strictly inside $B_{g(0)}(p,R_j)$, equal to one
on $B_{g(0)}(p,R_{j+1})$, and satisfying the clean bounds
\eqref{eq:shi-clean-cutoff-bounds}.  Lemma
\ref{lem:shi-compact-bernstein-maximum}, applied to $F_j$, gives
\[
                  \eta F_j\leq C.
\]
The function $F_j$ does extend continuously by zero at $t=0$: smoothness on
the compact cutoff support makes each unweighted curvature derivative
finite there, while the factor $t^{j+1}$ tends to zero.  On the inner ball,
$\eta=1$ and
\[
       a t^{j+1}|A_{j+1}|^2\leq F_j\leq C.
\]
Hence
\[
       |A_{j+1}|\leq C_{j+1}t^{-(j+1)/2}
       \quad\text{on }B_{g(0)}(p,R_{j+1})\times(0,T].
\]
This closes the induction.  After the final step $j=M-1$ (with $M=1$
supplied by the base case), the order-$M$ estimate holds on
$B_{g(0)}(p,R_M)=B_{g(0)}(p,R/2)$.  Scaling back by
Lemma~\ref{lem:shi-scaling} proves \eqref{eq:shi-main-estimate}.
\end{proof}

The preselected radii in \eqref{eq:shi-predetermined-radii} are not merely
cosmetic.  A dyadic induction that proves order $j$ only on
$B(p,R/2^j)$ does not, by itself, prove the theorem on the fixed half-ball.
The finite sequence above repairs that quantifier gap while allowing the
constant to depend on the target order $M$, exactly as the theorem permits.

\section{Complete global estimates and positive-time slabs}
\label{sec:shi-global}

The local theorem yields the complete global estimate without any maximum
principle at infinity.

\begin{corollary}[Complete global Bando--Shi estimates]
\label{cor:shi-complete-global}
Let $g(t)$, $0\leq t\leq T$, be a complete Ricci flow on an $n$-manifold,
and suppose
\[
                 |\Rm|\leq K
       \quad\text{on }M\times[0,T]
\]
for some $K>0$.  For every integer $m\geq0$,
\begin{align}
 |\nabla^m\Rm|(x,t)
   &\leq C(n,m)K\min\{t,K^{-1}\}^{-m/2}                 \label{eq:shi-complete-global}\\
   &\leq C(n,m)K\bigl(t^{-m/2}+K^{m/2}\bigr)            \label{eq:shi-complete-global-sum}
\end{align}
for $x\in M$ and $0<t\leq T$.
\end{corollary}

\begin{proof}
Fix $(x,t)$.  If $t\leq K^{-1}$, apply
Theorem~\ref{thm:shi-local-derivatives} to the interval $[0,t]$, with
$\alpha=\rho=1$, and evaluate the conclusion at the center $x$.  If
$t>K^{-1}$, apply the same theorem to the time-translated flow on
$[t-K^{-1},t]$ and again evaluate at $x$.  Completeness of the metric at the
new initial time and Hopf--Rinow make the required closed ball compact.
The curvature hypothesis is global, so it holds on that ball.  These two
cases give \eqref{eq:shi-complete-global}; the second inequality follows
from $\max\{a,b\}\leq a+b$.  The case $K=0$ is flat and may be added
separately.
\end{proof}

The preceding time translation is safe because the flow is complete and the
curvature bound is global.  It should not be copied into a theorem whose
only hypothesis concerns one fixed initial ball: a later-time ball need not
remain inside that original local domain.

\begin{corollary}[Positive-time slab regularity]
\label{cor:shi-positive-slab}
Let $I\subset\mathbb R$ be an interval and let $g(t)$ be a complete Ricci
flow on $M\times I$.  Suppose the curvature is uniformly bounded on every
compact time subinterval.  Fix
\[
                  [a_0,b]\subset I,
                  \qquad a_0<a\leq b,
\]
and let
\[
                  K=\sup_{M\times[a_0,b]}|\Rm|<\infty.
\]
Then, for every $m\geq0$,
\begin{equation}
 \sup_{M\times[a,b]}|\nabla^m\Rm|
   \leq C(n,m)K
      \left((a-a_0)^{-1/2}+\sqrt K\right)^m.             \label{eq:shi-positive-slab}
\end{equation}
When $K=0$, the right side and every derivative vanish.
\end{corollary}

\begin{proof}
Translate $a_0$ to time zero and apply
Corollary~\ref{cor:shi-complete-global}.  Every time in $[a,b]$ is at least
$a-a_0$ from the translated initial face.  Finally use
\[
 \min\{a-a_0,K^{-1}\}^{-m/2}
   \leq\left((a-a_0)^{-1/2}+\sqrt K\right)^m.
\]
\end{proof}

\section{Mixed spacetime derivative estimates}
\label{sec:shi-mixed}

For time derivatives we use the metric-compatible operator $\nabla_t$ from
\eqref{eq:shi-metric-time-derivative}.  It is invariant under the fixed
Uhlenbeck identification and assigns parabolic weight two to one time
derivative.

\begin{proposition}[Exact mixed-derivative expansion]
\label{prop:shi-mixed-derivatives}
For every $p,q\geq0$, there is a finite family of universal contractions
such that
\begin{equation}
 \nabla_t^q\nabla^p\Rm
   =\sum_\sigma c_\sigma\,
       \operatorname{contr}_\sigma
       \left(
          \nabla^{j_{\sigma,1}}\Rm\otimes\cdots\otimes
          \nabla^{j_{\sigma,\ell_\sigma}}\Rm
       \right),                                         \label{eq:shi-mixed-expansion}
\end{equation}
where $\ell_\sigma\geq1$, every $j_{\sigma,\nu}\geq0$, and every summand
satisfies the exact weight identity
\begin{equation}
       \sum_{\nu=1}^{\ell_\sigma}(j_{\sigma,\nu}+2)
             =p+2q+2.                                   \label{eq:shi-mixed-weight}
\end{equation}
The coefficients and contraction patterns depend only on $n,p,q$ and the
fixed curvature convention.
\end{proposition}

\begin{proof}
For $q=0$, take the single factor $\nabla^p\Rm$; its weight is $p+2$.
Suppose the assertion holds at order $q$.  Because $\nabla_tg=0$, applying
$\nabla_t$ commutes with every metric contraction.  Translate
\eqref{eq:shi-exact-Ak-evolution} back through
\eqref{eq:shi-time-pullback-bridge}.  It says that
\begin{equation}
 \nabla_t\nabla^j\Rm
   =\Delta\nabla^j\Rm
      +\sum_{i=0}^j\mathcal B_{j,i}
          (\nabla^i\Rm,\nabla^{j-i}\Rm),                 \label{eq:shi-time-one-factor}
\end{equation}
with all slots interpreted through the natural pullback bridge.
The Laplacian replaces a factor of weight $j+2$ by one derivative factor of
weight $j+4$.  A quadratic reaction replaces it by two factors whose total
weight is
\[
          (i+2)+(j-i+2)=j+4.
\]
Thus every application of $\nabla_t$ increases total weight by exactly two.
The tensor Leibniz rule produces only finitely many terms, proving the
induction.
\end{proof}

\begin{corollary}[Mixed estimates on a positive slab]
\label{cor:shi-mixed-slab}
In the setting of Corollary~\ref{cor:shi-positive-slab}, for every
$p,q\geq0$,
\begin{equation}
 \sup_{M\times[a,b]}
       |\nabla_t^q\nabla^p\Rm|
 \leq C(n,p,q)K
       \left((a-a_0)^{-1/2}+\sqrt K\right)^{p+2q}.        \label{eq:shi-mixed-slab}
\end{equation}
\end{corollary}

\begin{proof}
Put $S=(a-a_0)^{-1/2}+\sqrt K$.  The spatial estimates give
$|\nabla^j\Rm|\leq CKS^j$.  A term in
\eqref{eq:shi-mixed-expansion} with $\ell$ factors is therefore bounded by
\[
 C K^\ell S^{\sum j_\nu}
   =CKS^{p+2q}\left(\frac K{S^2}\right)^{\ell-1}
   \leq CKS^{p+2q},
\]
where \eqref{eq:shi-mixed-weight} and $S^2\geq K$ were used.  Sum the
finitely many terms.
\end{proof}

\section{Fixed-background control and continuation}
\label{sec:shi-continuation}

We now supply the regularity step that was left as an interface in
Theorem~\ref{thm:bounded-curvature-extension}.  The argument is valid in
every dimension.

\begin{proposition}[Fixed-background control]
\label{prop:shi-fixed-background-control}
Let $M$ be closed, let $a<b<\infty$, and let $g(t)$ be a Ricci flow on
$[a,b)$.  Suppose
\begin{enumerate}
\item the metrics $g(t)$ are uniformly equivalent to $g(a)$; and
\item for every $j\geq0$, $|\nabla_{g(t)}^j\Rm_{g(t)}|$ is uniformly bounded
      on $M\times[a,b)$.
\end{enumerate}
Put $\widehat\nabla=\nabla^{g(a)}$.  Then, for every $r\geq0$, the tensors
\begin{equation}
 \widehat\nabla^r g(t),\quad
 \widehat\nabla^r\Ric(g(t)),\quad
 \widehat\nabla^r(\nabla^{g(t)}-\widehat\nabla)           \label{eq:shi-background-controlled-list}
\end{equation}
are uniformly bounded on $M\times[a,b)$.  Moreover $g(t)$ converges in
$C^\infty(M,g(a))$ as $t\uparrow b$ to a smooth positive-definite metric
$g(b)$.
\end{proposition}

\begin{proof}
Let $S(t)=\nabla^{g(t)}-\widehat\nabla$.  The exact connection-variation
formula is
\begin{equation}
 \partial_tS^k{}_{ij}
   =-g^{k\ell}
      \left(
       \nabla_i\Ric_{j\ell}+\nabla_j\Ric_{i\ell}
                         -\nabla_\ell\Ric_{ij}
      \right).                                           \label{eq:shi-S-evolution}
\end{equation}
Here and in this proof, unadorned $\nabla$ is $\nabla^{g(t)}$.  The assumed
first curvature-derivative bound and metric equivalence make the right side
uniformly bounded, hence $S$ is bounded.

The conversion formula is also exact.  If $T$ has type $(u,v)$, then
\begin{align}
 (\widehat\nabla_iT)^{a_1\ldots a_u}{}_{b_1\ldots b_v}
 &= (\nabla_iT)^{a_1\ldots a_u}{}_{b_1\ldots b_v}\notag\\
 &\quad-
   \sum_{\mu=1}^u S^{a_\mu}{}_{ic}
    T^{a_1\ldots c\ldots a_u}{}_{b_1\ldots b_v}\notag\\
 &\quad+
   \sum_{\nu=1}^v S^c{}_{ib_\nu}
    T^{a_1\ldots a_u}{}_{b_1\ldots c\ldots b_v}.        \label{eq:shi-background-conversion}
\end{align}
In particular, $\widehat\nabla g$ is the sum of the two lower-slot
$S$-actions on $g$.

We now state the triangular induction explicitly.  Suppose
$\widehat\nabla^sS$ is bounded for every $s<r$.  Repeatedly differentiating
$\widehat\nabla g=S*g$ first bounds $\widehat\nabla^s g$ and
$\widehat\nabla^s g^{-1}$ for $s\leq r$.  When
$\widehat\nabla^r$ is applied to the right side of
\eqref{eq:shi-S-evolution}, expand the $r$ leading fixed derivatives from
left to right by \eqref{eq:shi-background-conversion}, while retaining the
terminal factor $\nabla\Ric$.  A term containing $S$ has spent at least one
of those $r$ derivatives on that factor.  Consequently, the expansion
contains only $\widehat\nabla^sS$ with $s\leq r-1$, together with evolving
derivatives of $\Ric$ of order at most $r+1$.  All are bounded by the
induction hypothesis and the assumed curvature-derivative estimates.  Since
$\widehat\nabla$ is fixed in time and $S(a)=0$, integrate
\[
       \partial_t(\widehat\nabla^rS)
          =\widehat\nabla^r(\partial_tS)
\]
to bound $\widehat\nabla^rS$.  The conversion formula then bounds
$\widehat\nabla^{r+1}g$ and $\widehat\nabla^r\Ric$.  This closes the
induction without using the quantity currently being proved.  The inverse
metric derivatives used above follow recursively from
\[
       \widehat\nabla(g^{-1})=-g^{-1}*(\widehat\nabla g)*g^{-1}.
\]
Thus all tensors in
\eqref{eq:shi-background-controlled-list} are uniformly bounded.

Finally,
\begin{equation}
       \partial_t\widehat\nabla^r g(t)
           =-2\widehat\nabla^r\Ric(g(t))                 \label{eq:shi-background-g-time}
\end{equation}
is uniformly bounded.  Hence $g(t)$ is Cauchy in every background $C^r$
norm as $t\uparrow b$.  The compatible limits define a smooth tensor
$g(b)$.  Uniform metric equivalence makes it positive definite.
\end{proof}

\begin{lemma}[Autonomous Ricci-flow jets]
\label{lem:shi-autonomous-jets}
For every $q\geq1$ there is a natural differential operator $P_q$, of
spatial order at most $2q$, such that every smooth Ricci flow satisfies
\begin{equation}
                       \partial_t^qg=P_q(g).              \label{eq:shi-autonomous-jets}
\end{equation}
Consequently, if two one-sided Ricci flows have the same smooth metric at a
common endpoint and each converges smoothly there, then all one-sided time
jets agree at that endpoint.
\end{lemma}

\begin{proof}
Take $P_1(g)=-2\Ric(g)$ and define recursively
\begin{equation}
             P_{q+1}(g)=DP_q|_g[P_1(g)],                 \label{eq:shi-Pq-recursion}
\end{equation}
where $DP_q|_g$ is the linearization of the natural differential operator
$P_q$ at $g$.  The chain rule gives
$\partial_t^{q+1}g=P_{q+1}(g)$ whenever
$\partial_t^qg=P_q(g)$.  Since $P_1$ has spatial order two, the recursion
raises order by at most two and proves the bound $2q$ by induction.  In
fixed coordinates this is the familiar rule that differentiates the
universal expression in $g^{-1}$, $\partial g$, and $\partial^2g$, with
$\partial_tg^{-1}=-g^{-1}(\partial_tg)g^{-1}$.  Naturality is preserved by
linearization and composition.  The value at an endpoint depends only on
the spatial jet of the common endpoint metric, so both one-sided time jets
equal $P_q$ of that metric.
\end{proof}

\begin{theorem}[Closed short-time existence; imported input]
\label{thm:shi-closed-short-time-input}
Let $M^n$ be a closed smooth manifold and let $g_0$ be a smooth Riemannian
metric on $M$.  There are $\varepsilon>0$ and a smooth Ricci flow $h(s)$,
$0\leq s<\varepsilon$, such that $h(0)=g_0$.
\end{theorem}

This is the general-dimensional Hamilton--DeTurck short-time existence
theorem, recorded for example as Theorem~3.13 of MSM~110.  It is an earlier
project input and is not reproved in this chapter; Chapter
\ref{ch:architecture} states its three-dimensional specialization.

\begin{theorem}[Closed bounded-curvature extension, all dimensions]
\label{thm:shi-closed-extension}
Let $M^n$ be closed, let $0<T<\infty$, and let $g(t)$, $0\leq t<T$, be a
smooth Ricci flow.  If
\begin{equation}
             \sup_{M\times[0,T)}|\Rm|<\infty,             \label{eq:shi-extension-curvature-bound}
\end{equation}
then $g(t)$ extends smoothly as a Ricci flow to an interval
$[0,T+\varepsilon)$ for some $\varepsilon>0$.
\end{theorem}

\begin{proof}
Let $K$ be the bound in \eqref{eq:shi-extension-curvature-bound}.  The
metric evolution and $|\Ric|\leq c_nK$ give uniform equivalence of $g(t)$
and $g(0)$.  For each $b\in(T/2,T)$, apply
Corollary~\ref{cor:shi-positive-slab} on $[T/4,b]$ and restrict to
$[T/2,b]$.  Its constants depend on the fixed separation $T/4$, not on
$b$.  Letting $b\uparrow T$ therefore bounds every spatial curvature
derivative on $M\times[T/2,T)$.  Proposition
\ref{prop:shi-fixed-background-control} produces a smooth positive metric
$g(T)$ and smooth convergence $g(t)\to g(T)$.

Invoke Theorem~\ref{thm:shi-closed-short-time-input} with initial metric
$g(T)$.  Let $h(s)$, $0\leq s<\varepsilon$, be the resulting solution and
define
\[
 \widetilde g(t)=
 \begin{cases}
   g(t),&0\leq t<T,\\
   h(t-T),&T\leq t<T+\varepsilon.
 \end{cases}
\]
Both pieces solve Ricci flow and agree continuously at $T$.  Smooth spatial
convergence and Lemma~\ref{lem:shi-autonomous-jets} show that every
one-sided time jet agrees there.  Hence $\widetilde g$ is smooth and is the
required extension.
\end{proof}

In dimension three, Theorem~\ref{thm:shi-closed-extension} proves
Theorem~\ref{thm:bounded-curvature-extension}.  It asserts the existence of
some $\varepsilon>0$.  A lower bound for $\varepsilon$ depending only on
$n$ and $K$ would additionally use a quantitative complete
bounded-curvature existence theorem; that stronger input is not needed
here.

\begin{corollary}[Quantitative curvature blow-up]
\label{cor:shi-quantitative-blowup}
Let $M^n$ be closed and let $g(t)$, $0\leq t<T<\infty$, be a Ricci flow with
maximal forward time $T$.  Set
\[
                       Q(t)=\max_M|\Rm|(\cdot,t).
\]
There is $c_n>0$ such that
\begin{equation}
             Q(t)\geq\frac{c_n}{T-t}
             \quad(0\leq t<T).                           \label{eq:shi-type-I-lower-rate}
\end{equation}
In particular,
\begin{equation}
                       \lim_{t\uparrow T}Q(t)=\infty.     \label{eq:shi-curvature-limit-infinity}
\end{equation}
\end{corollary}

\begin{proof}
The order-zero inequality \eqref{eq:shi-u0-evolution} and the compact scalar
maximum principle imply the doubling-time estimate: there is $d_n>0$ such
that, if $Q(t_0)>0$, then
\begin{equation}
        Q(s)\leq2Q(t_0)
        \quad\text{whenever }t_0\leq s<T
        \text{ and }s-t_0\leq\frac{d_n}{Q(t_0)}.          \label{eq:shi-doubling-time}
\end{equation}
Indeed, apply the upper-Dini derivative inequality
$D^+Q^2\leq C_nQ^3$ and integrate the scalar ODE.

If $Q(t_0)=0$, the same ODE comparison forces $Q(s)=0$ for every
$s\in[t_0,T)$.  The terminal curvature is then bounded and
Theorem~\ref{thm:shi-closed-extension} contradicts maximality.  Hence
$Q(t)>0$ for every $t<T$.

If $Q(t_0)<d_n/(T-t_0)$, then
\eqref{eq:shi-doubling-time} bounds curvature on the whole terminal interval
$[t_0,T)$.  Curvature is already bounded on the compact earlier interval,
so Theorem~\ref{thm:shi-closed-extension} contradicts maximality.  Thus
\eqref{eq:shi-type-I-lower-rate} holds with $c_n=d_n$; its right side tends
to infinity, proving \eqref{eq:shi-curvature-limit-infinity}.
\end{proof}

\section{Interfaces for the next analytic chapters}
\label{sec:shi-project-interfaces}

The point of proving the estimates here is to remove vague regularity
phrases from later theorems.  The following two propositions state exactly
what may be imported.

\begin{proposition}[Derivative input for Hamilton's matrix Harnack estimate]
\label{prop:shi-harnack-input}
Let $I\subset[0,\infty)$ be an interval, let $g(t)$ be a complete Ricci flow
on $M\times I$, and suppose curvature is uniformly bounded on every compact
subinterval of $I$.  Fix
\[
             0\leq a_0<a\leq b,
             \qquad [a_0,b]\subset I.
\]
Then
\begin{equation}
       \sup_{M\times[a,b]}
        \bigl(|\Rm|+|\nabla\Rm|+|\nabla^2\Rm|\bigr)<\infty.             \label{eq:shi-harnack-input}
\end{equation}
Consequently all coefficients in Hamilton's tensors formed from
$\Ric$, $\nabla\Ric$, $\nabla^2\Ric$, $\Delta\Ric$, and the factor
$t^{-1}\Ric$ are bounded on $[a,b]$ when $a>0$.  If time is translated so
that $a_0$ is the initial time, the corresponding factor
$(t-a_0)^{-1}\Ric$ is bounded because $a-a_0>0$.
\end{proposition}

\begin{proof}
Use Corollary~\ref{cor:shi-positive-slab} for orders $0,1,2$ and contract
indices.  The scalar factor $t^{-1}$ is bounded on $[a,b]$ when $a>0$.
\end{proof}

This proposition supplies only analytic boundedness.  The next chapter
proves the exact Harnack evolution equations, the curvature-cone algebra,
and the appropriate compact or noncompact first-contact argument.

\begin{proposition}[Derivative input for Hamilton compactness]
\label{prop:shi-compactness-input}
Fix an integer $n\geq2$.  Let
\[
   (M_i,g_i(t),p_i),\qquad t\in(\omega_-,\omega_+),
   \qquad \omega_-<0<\omega_+,
\]
be complete pointed Ricci flows on $n$-manifolds.  Suppose that for every compact interval
$J\Subset(\omega_-,\omega_+)$ there is $K_J<\infty$ such that
\[
             \sup_i\sup_{M_i\times J}|\Rm_{g_i(t)}|\leq K_J.
\]
Then for every $p,q\geq0$ and every compact
$J\Subset(\omega_-,\omega_+)$,
\begin{equation}
       \sup_i\sup_{M_i\times J}
          |\nabla_t^q\nabla^p\Rm_{g_i(t)}|<\infty.       \label{eq:shi-compactness-derivatives}
\end{equation}
\end{proposition}

\begin{proof}
Choose a slightly larger compact interval $J'$ with
$J\Subset\operatorname{int}J'\Subset(\omega_-,\omega_+)$.  Apply
Corollary~\ref{cor:shi-mixed-slab} on $J'$, using the positive separation
between the left endpoints of $J$ and $J'$.  The resulting constant depends
only on $n,p,q,J,J'$, and $K_{J'}$, and is therefore uniform in $i$.
\end{proof}

To obtain Hamilton compactness itself one must additionally assume, for
example, a uniform lower bound
$\inj_{g_i(0)}(p_i)\geq\iota_0>0$ and prove pointed Riemannian compactness,
the construction of comparison diffeomorphisms, diagonal extraction, and
completeness of the limit.  Those steps are not part of the Bando--Shi
estimate.  The condition $\omega_-<0<\omega_+$ is also meaningful: if zero
were the left endpoint, curvature alone would not control the initial
derivative jets.  Smooth convergence through such an endpoint needs initial
derivative bounds or the modified Shi estimates.

\section{Formalization units and dependency boundary}
\label{sec:shi-formalization}

The proof should be formalized in the following dependency order.  The
identifiers record mathematical interfaces; they are not claims of existing
Lean declarations.

\begin{enumerate}
\item \nodeid{PC-F-SHI-PARABOLIC-SCALING-WEIGHTS} is
      Lemma~\ref{lem:shi-scaling}.  It records the weights of distance,
      time, curvature, and every curvature derivative.

\item \nodeid{PC-F-SHI-FULL-UHLENBECK-SPATIAL-DERIVATIVE} is
      Definition~\eqref{eq:shi-full-D} together with
      Lemma~\ref{lem:shi-every-slot-pulled}.  Its output lies in
      $(V^*)^{\otimes(k+4)}$; no derivative slot remains in $T^*M$.

\item \nodeid{PC-F-SHI-UHLENBECK-HEAT-COMMUTATOR} is
      Lemma~\ref{lem:shi-exact-commutator}.  The exact output is a finite
      slot sum $\Rm*\mathcal DA$, with no $\mathcal D\Rm*A$ term.

\item \nodeid{PC-F-SHI-UNIVERSAL-CONTRACTION-RECURSION} is
      Definition~\ref{def:shi-Bki};
      \nodeid{PC-F-SHI-EXACT-HIGHER-CURVATURE-EVOLUTION} is
      Theorem~\ref{thm:shi-exact-all-order-evolution}; and
      \nodeid{PC-F-SHI-HIGHER-CURVATURE-NORM-EVOLUTION} is
      Proposition~\ref{prop:shi-exact-norm-identity}.  These three units
      keep equality-level tensor algebra separate from operator-norm bounds.

\item \nodeid{PC-F-SHI-INITIAL-DISTANCE-CUTOFF} is
      Lemma~\ref{lem:shi-self-coupled-cutoff}, including the lower-support
      definition and Calabi replacement.  Its clean post-gradient form is
      \nodeid{PC-F-SHI-CLEAN-HIGHER-CUTOFF},
      Corollary~\ref{cor:shi-clean-cutoff}.

\item \nodeid{PC-F-SHI-FIRST-DERIVATIVE-BERNSTEIN} is
      Proposition~\ref{prop:shi-normalized-first}.  Its reusable scalar
      subunits are the product evolution
      \eqref{eq:shi-first-F-good}, the cutoff closure
      \eqref{eq:shi-Theta-closes}, and polynomial absorption in
      Lemma~\ref{lem:shi-polynomial-absorption}.

\item \nodeid{PC-F-SHI-HIGHER-BERNSTEIN-PRODUCT} is
      Lemma~\ref{lem:shi-higher-bernstein-product};
      \nodeid{PC-F-SHI-COMPACT-CUTOFF-MAXIMUM} is
      Lemma~\ref{lem:shi-compact-bernstein-maximum}; and
      \nodeid{PC-F-SHI-FIXED-HALF-BALL-INDUCTION} is the finite-radius
      induction \eqref{eq:shi-predetermined-radii}.
      Together they prove
      \nodeid{PC-F-SHI-LOCAL-DERIVATIVE-ESTIMATES},
      Theorem~\ref{thm:shi-local-derivatives}.

\item \nodeid{PC-F-SHI-COMPLETE-GLOBAL-DERIVATIVE-ESTIMATES} and
      \nodeid{PC-F-SHI-POSITIVE-TIME-SLAB-REGULARITY} are
      Corollaries~\ref{cor:shi-complete-global} and
      \ref{cor:shi-positive-slab}.  Their noncompact proof uses only the
      local theorem and Hopf--Rinow, not a maximum principle at infinity.

\item \nodeid{PC-F-RICCI-FLOW-MIXED-DERIVATIVE-EXPANSION} is
      Proposition~\ref{prop:shi-mixed-derivatives}.
      \nodeid{PC-F-BOUNDED-CURVATURE-FIXED-BACKGROUND-CONTROL} and
      \nodeid{PC-F-CLOSED-BOUNDED-CURVATURE-EXTENSION} are
      Proposition~\ref{prop:shi-fixed-background-control} and
      Theorem~\ref{thm:shi-closed-extension}.  The latter also has the
      explicit imported dependency
      \nodeid{PC-I-CLOSED-RICCI-FLOW-SHORT-TIME-EXISTENCE},
      Theorem~\ref{thm:shi-closed-short-time-input}; that node is not
      reproved here.

\item The bridge nodes \nodeid{PC-B-SHI-TO-MATRIX-HARNACK} and
      \nodeid{PC-B-SHI-TO-HAMILTON-COMPACTNESS} are exactly
      Propositions~\ref{prop:shi-harnack-input} and
      \ref{prop:shi-compactness-input}; neither claims the downstream
      theorem itself.
\end{enumerate}

The proved chapter-level dependencies are the exact pulled curvature
evolution from Chapter~\ref{ch:curvature-evolution}, elementary finite
dimensional tensor norms, the compactly supported scalar parabolic maximum
principle, metric and connection variation, Hessian comparison, the Calabi
lower-support construction, and Hopf--Rinow.  The extension theorem also
imports closed-manifold short-time existence.  No complete noncompact
existence theorem, complete-flow uniqueness theorem, noncollapsing theorem,
or injectivity-radius estimate is used in the local derivative proof.

\subsection*{Lean status}

Every new identifier in this chapter currently has Lean statement and Lean
proof status \emph{not yet implemented}.  The natural-language proofs are
decomposed for later formalization, but no typechecking, kernel checking,
axiom-footprint inspection, or source-to-Lean correspondence audit is
claimed.

\section{Sources, corrections, and deliberately deferred results}
\label{sec:shi-sources}

The primary source for the local estimates is W.-X.~Shi,
\emph{Deforming the metric on complete Riemannian manifolds}, Journal of
Differential Geometry~\textbf{30} (1989), no.~1, 223--301,
Section~7, pp.~289--297,
\href{https://doi.org/10.4310/jdg/1214443292}{doi:10.4310/jdg/1214443292}.
Shi's companion paper, \emph{Ricci deformation of the metric on complete
noncompact Riemannian manifolds}, Journal of Differential
Geometry~\textbf{30} (1989), no.~2, 303--394, supplies application and
existence context but is not the primary source for the local proof.

The compact global precursor is S.~Bando,
\emph{Real analyticity of solutions of Hamilton's equation},
Mathematische Zeitschrift~\textbf{195} (1987), no.~1, 93--97,
\href{https://doi.org/10.1007/BF01161602}{doi:10.1007/BF01161602}.
Bando proves factorial estimates leading to spatial real analyticity on
compact positive-time slices.  This chapter proves the ordinary all-orders
estimates needed by the Poincar\'e project; it does not reproduce the
factorial analyticity argument.

The detailed modern proof consulted here is Chapter~14 of MSM~144.  The
ordinary moving-connection commutators in its Chapter~13 were used as a
cross-check, while Appendix~F supplied the cleaner metric-compatible
commutator.  Chapter~7 of MSM~110 supplied the global continuation route;
Chapter~6 of GSM~77 supplied independent first- and second-derivative
checks; and Theorem~1.6 of GSM~235 supplied the clean curvature-scale
statement.  Hamilton's sharper reciprocal-barrier first-derivative estimate
is in R.~Hamilton, \emph{The formation of singularities in the Ricci flow},
Surveys in Differential Geometry~\textbf{2} (1995), 7--136, Section~13.
That sharper estimate is not needed for the all-orders induction and is
deferred because its reciprocal cutoff and first-contact bootstrap require
substantially more formal infrastructure.

The following source hazards have been repaired rather than propagated.
\begin{enumerate}
\item In the higher-derivative calculation, every one of the $k+4$ tensor
      slots is pulled to the fixed bundle.  Pulling only four curvature slots
      does not freeze the norm of the derivative slots.

\item The general formula \eqref{eq:shi-expanded-commutator} is derived
      directly from the infinitesimal orthogonal action.  In Appendix~F of
      MSM~144, the displayed special case for a covariant three-tensor omits
      the factor $2$ present in the preceding one- and two-tensor formulas.
      The general derivation fixes the factor without depending on that
      display.

\item In the proof of the local first estimate in MSM~144, the line before
      its equation (14.20) must read, in the notation used here,
      \[
          \sup(\eta\Theta)^{1/2}\leq C(\eta G)^{1/2},
      \]
      not $\eta G$.  The resulting $Z_*^{3/2}$ term in
      \eqref{eq:shi-first-polynomial} confirms the missing square root.

\item The same proof must extract the gradient estimate from
      \[
        16K^2t|\nabla\Rm|^2
          \leq t(16K^2+|\Rm|^2)|\nabla\Rm|^2.
      \]
      The printed coefficient $17$ is the upper comparison coefficient and
      is incorrect in this lower inequality.

\item The displayed dyadic induction in MSM~144 yields successively smaller
      balls and does not itself imply the stated fixed half-ball.  The
      predetermined finite sequence \eqref{eq:shi-predetermined-radii}
      proves the advertised spatial conclusion exactly.  It also removes an
      indexing mismatch in the first use of the post-gradient cutoff lemma.

\item Constants are expressed through the dimensionless radius
      $\sqrt K r$.  Thus the local estimate has the scale-correct form
      $C(n,m,\alpha,\sqrt Kr)Kt^{-m/2}$ rather than a dimensionful,
      uninterpreted dependence on $K$ and $r$ separately.

\item Time translation is used only in
      Corollary~\ref{cor:shi-complete-global}, where completeness and a
      global curvature bound make it valid.  It is not used to extend a
      fixed-initial-ball local theorem to arbitrarily long times.
\end{enumerate}

The attached surface calculation correctly emphasizes induction on the
commutator and on the squared norm.  Its schematic contraction notation is
adequate for estimates but cannot be differentiated as literal syntax.
Definition~\ref{def:shi-Bki} supplies the missing fixed contraction maps in
all dimensions.

Modified estimates that retain prescribed initial derivative bounds,
Bando's factorial real-analyticity theorem, complete noncompact
short-time existence and uniqueness, Hamilton compactness, noncollapsing,
and the Harnack-dependent instantaneous estimates are deliberately
deferred.  Hamilton's matrix Harnack estimate is proved in the next chapter,
using the bridge in Proposition~\ref{prop:shi-harnack-input}.  None of these
downstream results is used circularly in the proof of
Theorem~\ref{thm:shi-local-derivatives}.

\chapter{Hamilton's Matrix Harnack Estimate}
\label{ch:matrix-harnack}
\index{Hamilton!matrix Harnack estimate}
\index{Harnack estimate!matrix}
\index{Harnack estimate!trace}
\index{ancient solution!Harnack estimate}

\section{Purpose, conventions, and the shifted theorem}
\label{sec:harnack-purpose}

Hamilton's matrix Harnack estimate is a differential comparison theorem for
Ricci flow with nonnegative curvature operator.  It packages the curvature
tensor, its first derivative, and a particular second-derivative tensor into
one positive-semidefinite block form.  Its trace controls scalar curvature
at different points and times; its ancient limit removes the term containing
the time origin and gives the monotonicity $\partial_tR\geq0$.

The proof has four logically different layers.  The first is exact tensor
calculus.  The second is a finite-dimensional sum-of-squares identity.  The
third is the scalar first-contact mechanism underlying Chapter~2.  The
fourth is a proper barrier which prevents first contact from escaping to
spatial infinity.  We keep the layers separate.  In particular, no inverse
of the curvature operator is used, so zero curvature directions cause no
singularity.

We state the theorem with an arbitrary earlier time $\alpha$.  This is not
cosmetic: the same statement gives the usual $t^{-1}$ estimate by choosing
the origin, and gives the ancient estimate by sending $\alpha$ to
$-\infty$.  Put
\begin{equation}
                         \tau=t-\alpha>0.                 \label{eq:har-tau}
\end{equation}
Every metric contraction and every raising of an index below is taken using
$g(t)$.  We retain the curvature convention of
Definition~\ref{def:hi-curvature-tensor}; thus $R_{1221}>0$ on a round
positively curved sphere.

For covectors $A,B$, define the normalized alternation
\begin{equation}
 (A\mathbin{\wedge_{\!A}}B)_{ij}
       =\frac12(A_iB_j-A_jB_i).                          \label{eq:har-normalized-wedge}
\end{equation}
The factor $1/2$ is load-bearing in the trace calculation.

\begin{definition}[Hamilton's $P$, $M$, and block form]
\label{def:har-block-form}
Let $g(t)$ be a Ricci flow and let $\alpha<t$.  Define
\begin{align}
 P_{kij}
   &=\nabla_kR_{ij}-\nabla_iR_{kj},                       \label{eq:har-P-definition}\\
 \overline M_{ij}
   &=\Delta R_{ij}-\frac12\nabla_i\nabla_jR
      +2R_{kij\ell}R^{k\ell}-R_i{}^kR_{kj},             \label{eq:har-Mbar-definition}\\
 M^{(\alpha)}_{ij}
   &=\overline M_{ij}+\frac1{2(t-\alpha)}R_{ij}.         \label{eq:har-M-definition}
\end{align}
For $x\in M$, $U\in\Lambda^2T_x^*M$, and $W\in T_x^*M$, set
\begin{equation}
 \begin{split}
 Q^{(\alpha)}_{x,t}(U,W)
   ={}&M^{(\alpha)}_{ij}W^iW^j
      +2P_{pij}U^{pi}W^j
      +R_{pijq}U^{pi}U^{qj}.                            \label{eq:har-Q-definition}
 \end{split}
\end{equation}
Its polarization is the symmetric bilinear form on
$\Lambda^2T_x^*M\oplus T_x^*M$ denoted by the same letter.
\end{definition}

The curvature term in \eqref{eq:har-Q-definition} has the positive sign.
Indeed, after relabeling ordered indices it is
\begin{equation}
                  R_{abdc}U_{ab}U_{cd}.                 \label{eq:har-positive-block}
\end{equation}
If $u=\frac12U_{ab}e_a\wedge e_b$, then the trace-normalized curvature
operator of Definition~\ref{def:hi-curvature-operator} satisfies
\begin{equation}
 R_{abdc}U_{ab}U_{cd}
       =2\langle\mathcal R u,u\rangle.                  \label{eq:har-operator-bridge}
\end{equation}
Thus nonnegativity of the curvature operator is exactly nonnegativity of
the upper-left block.  Writing $R_{abcd}$ instead of $R_{abdc}$ in that
block would reverse the sign.

\begin{theorem}[Shifted complete matrix Harnack estimate]
\label{thm:har-shifted-matrix}
Let $n\geq2$, let $I\subset\mathbb R$ be an open interval, and let
\[
                         g(t),\qquad t\in I,
\]
be a smooth complete Ricci flow on an $n$-manifold $M$.  Assume
\begin{enumerate}
\item the curvature operator is nonnegative at every point of $M\times I$;
\item for every compact interval $J\Subset I$,
      \begin{equation}
          \sup_{M\times J}|\Rm|<\infty.                 \label{eq:har-slab-bound}
      \end{equation}
\end{enumerate}
Then for every $\alpha,t\in I$ with $\alpha<t$, every $x\in M$, and every
$U\in\Lambda^2T_x^*M$ and $W\in T_x^*M$,
\begin{equation}
                         Q^{(\alpha)}_{x,t}(U,W)\geq0.   \label{eq:har-shifted-conclusion}
\end{equation}
No connectedness, orientability, injectivity-radius, or noncollapsing
hypothesis is required.
\end{theorem}

\textcolor{red}{Since the above is a major theorem, put the Lean code for the statement here.}

\begin{corollary}[Hamilton's finite-origin matrix Harnack estimate]
\label{cor:har-finite-origin}
Let $g(t)$, $0<t<T$, be a complete Ricci flow with nonnegative curvature
operator and curvature bounded on every compact positive-time slab.  Define
$P$ by \eqref{eq:har-P-definition} and
\begin{equation}
 M_{ij}=\Delta R_{ij}-\frac12\nabla_i\nabla_jR
      +2R_{kij\ell}R^{k\ell}-R_i{}^kR_{kj}
      +\frac1{2t}R_{ij}.                                \label{eq:har-M-origin}
\end{equation}
Then
\begin{equation}
 M_{ij}W^iW^j+2P_{pij}U^{pi}W^j+R_{pijq}U^{pi}U^{qj}
       \geq0                                             \label{eq:har-origin-conclusion}
\end{equation}
at every positive time.
\end{corollary}

\begin{proof}
Fix $(x,t,U,W)$.  For every $\alpha\in(0,t)$,
Theorem~\ref{thm:har-shifted-matrix} gives
\[
 \overline Q(U,W)+\frac{\Ric(W^\sharp,W^\sharp)}{2(t-\alpha)}\geq0,
\]
where $\overline Q$ is obtained from $Q^{(\alpha)}$ by replacing
$M^{(\alpha)}$ with $\overline M$.  Let $\alpha\downarrow0$.  This is a
limit of real numbers in one fixed tangent space and yields
\eqref{eq:har-origin-conclusion}.
\end{proof}

\begin{corollary}[Closed-flow form]
\label{cor:har-closed}
Let $M^n$ be closed and let $g(t)$, $0<t<T$, be a smooth Ricci flow whose
curvature operator is nonnegative at every positive time.  Then
\eqref{eq:har-origin-conclusion} holds.  In dimension three, the same
conclusion holds for a flow on $[0,T)$ if only the curvature operator of
$g(0)$ is assumed nonnegative.
\end{corollary}

\begin{proof}
Closedness makes curvature bounded on every compact time slab, so the first
claim is Corollary~\ref{cor:har-finite-origin}.  In dimension three, the
curvature-operator cone is preserved by the weak systems maximum principle
and the exact Uhlenbeck curvature evolution of Chapters~2 and~4; this gives
the hypothesis of the first claim.
\end{proof}

\begin{roadmap}{Sections~\ref{sec:har-fixed-calculus} and
\ref{sec:har-block-evolution} prove all equality-level tensor identities in
the fixed Uhlenbeck bundle.}
Section~\ref{sec:har-algebra} proves Hamilton's reaction as a sum of
squares.  Sections~\ref{sec:har-first-contact} and
\ref{sec:har-proof-main} supply the compact first-contact argument and the
proper noncompact barrier.  Sections~\ref{sec:har-trace} and
\ref{sec:har-ancient} derive the trace, integrated, and ancient estimates.
The final sections record regression tests, formalization units, sources,
and source corrections.
\end{roadmap}

\section{Exact tensor calculus in the fixed bundle}
\label{sec:har-fixed-calculus}

Fix $\alpha<b$ with $[\alpha,b]\Subset I$.  Choose an earlier buffer time
$\alpha_-<\alpha$ with $[\alpha_-,b]\Subset I$.  Start Uhlenbeck's
metric-freezing construction at $\alpha$: take a rank-$n$ bundle $(V,h)$,
an isometry
\[
       \iota_\alpha:(V,h)\longrightarrow(TM,g(\alpha)),
\]
and solve
\begin{equation}
       \partial_t\iota_t=\Ric_{g(t)}^\sharp\circ\iota_t. \label{eq:har-iota-ode}
\end{equation}
Then $\iota_t^*g(t)=h$, and the pulled connection $D^t$ is compatible with
the time-independent metric $h$.

Use the fully pulled spatial derivative $\mathcal D$ of
\eqref{eq:shi-full-D}: every tensor slot and every derivative direction is
transported to $V$.  Put
\begin{equation}
       \Delta^U=\operatorname{tr}_h\mathcal D^2,
       \qquad L=\partial_t-\Delta^U.                    \label{eq:har-fixed-L}
\end{equation}
The time derivative in $L$ is ordinary differentiation in a fixed vector
bundle.  We suppress pullback bars from now until the geometric theorem is
recovered.  Thus $R_{abcd}$ denotes the fully pulled covariant curvature
tensor, and $R_{1221}>0$ remains the sign convention.

The safe commutator is Lemma~\ref{lem:shi-exact-commutator}.  For a
covariant $q$-tensor $A$ it says
\begin{equation}
 \bigl(L\mathcal D A-\mathcal DLA\bigr)_{a b_1\ldots b_q}
 =2\sum_{r=1}^q
      R_{aedb_r}(\mathcal D_eA)_{b_1\ldots d\ldots b_q}. \label{eq:har-master-commutator}
\end{equation}
Each covariant slot contributes its own term and its own factor $2$.

\begin{lemma}[Algebraic and divergence identities for $P$]
\label{lem:har-P-identities}
The tensor $P$ in \eqref{eq:har-P-definition}, after pullback, satisfies
\begin{align}
 P_{abc}&=-P_{bac},                                      \label{eq:har-P-skew}\\
 P_{abc}+P_{bca}+P_{cab}&=0,                             \label{eq:har-P-cyclic}\\
 h^{ac}P_{abc}&=-\frac12\mathcal D_bR,                  \label{eq:har-P-trace-one}\\
 h^{bc}P_{abc}&=\frac12\mathcal D_aR,                   \label{eq:har-P-trace-two}\\
 \mathcal D_cR_{cdab}&=P_{bad}.                          \label{eq:har-divergence-rm}
\end{align}
Moreover,
\begin{equation}
 \mathcal D_cP_{cab}
 =\Delta^UR_{ab}-\frac12\mathcal D_a\mathcal D_bR
   -R_{ad}R_{db}+R_{acdb}R_{cd}.                         \label{eq:har-divergence-P}
\end{equation}
Consequently
\begin{equation}
 M^{(\alpha)}_{ab}
 =\mathcal D_cP_{cab}+R_{acdb}R_{cd}
     +\frac1{2\tau}R_{ab}.                              \label{eq:har-M-divergence}
\end{equation}
In particular, $M^{(\alpha)}$ is symmetric.
\end{lemma}

\begin{proof}
Skew-symmetry is immediate from the definition.  Expanding the cyclic sum
and using symmetry of Ricci makes its six terms cancel in pairs.  Contracted
second Bianchi gives $\mathcal D^aR_{ab}=\frac12\mathcal D_bR$, which proves
\eqref{eq:har-P-trace-one} and \eqref{eq:har-P-trace-two}.  The differential
second Bianchi identity, contracted in its derivative and one curvature
index, gives \eqref{eq:har-divergence-rm} with the displayed order of the
last pair.

For the remaining identity, expand
\[
 \mathcal D_cP_{cab}
   =\Delta^UR_{ab}-\mathcal D_c\mathcal D_aR_{cb}.
\]
Commute the last two derivatives, then use contracted second Bianchi.  The
two curvature actions on the covariant Ricci tensor are
$R_{ad}R_{db}$ and $-R_{acdb}R_{cd}$ with the signs shown in
\eqref{eq:har-divergence-P}.  Substitution into
\eqref{eq:har-M-definition} proves \eqref{eq:har-M-divergence}.  Since the
left side of \eqref{eq:har-M-definition} is visibly symmetric except
possibly for its curvature contraction, pair symmetry and symmetry of
Ricci show that contraction is symmetric as well.
\end{proof}

Define, as in Chapter~4,
\begin{equation}
 B_{abcd}=-h^{eg}h^{fh}R_{aebf}R_{cgdh}.                 \label{eq:har-B-definition}
\end{equation}
The previously proved pulled evolutions are
\begin{align}
 LR_{abcd}
   &=2(B_{abcd}-B_{abdc}+B_{acbd}-B_{adbc}),             \label{eq:har-Rm-evolution}\\
 LR_{ab}&=2R_{acdb}R_{cd},                               \label{eq:har-Ric-evolution}\\
 LR&=2R_{ab}R_{ab}.                                      \label{eq:har-scalar-evolution}
\end{align}

\begin{proposition}[Exact evolution of $P$]
\label{prop:har-P-evolution}
In the fixed bundle,
\begin{equation}
 \begin{split}
 LP_{abc}={}&2R_{adeb}P_{dec}+2R_{adec}P_{dbe}
              +2R_{bdec}P_{ade}\\
            &-2R_{de}\mathcal D_dR_{abec}.
                                                               \label{eq:har-P-evolution}
 \end{split}
\end{equation}
\end{proposition}

\begin{proof}
Apply $L$ to $\mathcal D_aR_{bc}-\mathcal D_bR_{ac}$.  Use
\eqref{eq:har-master-commutator} once for each derivative and
\eqref{eq:har-Ric-evolution}.  Before any abbreviation, the result regroups
as
\begin{align*}
 LP_{abc}={}&2R_{de}
  (\mathcal D_aR_{bdec}-\mathcal D_bR_{adec})\\
 &+2R_{bdec}(\mathcal D_aR_{de}-\mathcal D_dR_{ae})\\
 &+2R_{adeb}(\mathcal D_dR_{ec}-\mathcal D_eR_{dc})\\
 &+2R_{adec}(\mathcal D_dR_{be}-\mathcal D_bR_{de}).
\end{align*}
The second Bianchi identity changes the first parenthesis to
$-\mathcal D_dR_{abec}$.  The other three parentheses are exactly the
three indicated components of $P$.  This proves
\eqref{eq:har-P-evolution} without schematic star notation.
\end{proof}

\begin{proposition}[Exact evolution of $M^{(\alpha)}$]
\label{prop:har-M-evolution}
For $\tau=t-\alpha$,
\begin{equation}
 \begin{split}
 LM^{(\alpha)}_{ab}={}&2R_{acdb}M^{(\alpha)}_{cd}
  +2R_{cd}(\mathcal D_cP_{dab}+\mathcal D_cP_{dba})\\
 &+P_{cda}P_{cdb}-2P_{acd}P_{bdc}
  +2R_{cd}R_{ce}R_{adeb}
  -\frac1{2\tau^2}R_{ab}.                              \label{eq:har-M-evolution}
 \end{split}
\end{equation}
\end{proposition}

\begin{proof}
We give a cancellation proof with three separately reusable identities.
Starting from Proposition~\ref{prop:har-P-evolution}, use
\eqref{eq:har-master-commutator}, the Bianchi identities, and
\eqref{eq:har-divergence-rm} to obtain
\begin{align}
 L(\mathcal D_cP_{cab})={}&
  2R_{acdb}\mathcal D_eP_{ecd}
 +2R_{cd}(\mathcal D_cP_{dab}+\mathcal D_cP_{dba})
 \notag\\
 &+P_{cda}P_{cdb}-2P_{acd}P_{bdc}
 +2R_{cd}R_{ce}R_{adeb}                                \notag\\
 &+2\mathcal D_eR_{acdb}\mathcal D_eR_{cd}              \notag\\
 &+2R_{cd}(B_{abcd}+B_{acbd}-B_{acdb}-B_{adcb}).         \label{eq:har-L-divP}
\end{align}
The product rule and \eqref{eq:har-Rm-evolution}--
\eqref{eq:har-Ric-evolution} give
\begin{align}
 L(R_{acdb}R_{cd})={}&
  2R_{acdb}R_{cefd}R_{ef}
 -2\mathcal D_eR_{acdb}\mathcal D_eR_{cd}               \notag\\
 &+2R_{cd}(B_{acdb}-B_{acbd}+B_{adcb}-B_{abcd}).         \label{eq:har-L-RmRic}
\end{align}
Finally, because $L(1/(2\tau))=-1/(2\tau^2)$ and because of
\eqref{eq:har-Ric-evolution} and \eqref{eq:har-M-divergence},
\begin{equation}
 \begin{split}
 L\!\left(\frac1{2\tau}R_{ab}\right)
  ={}&2R_{acdb}
      \left(M^{(\alpha)}_{cd}-\mathcal D_eP_{ecd}
                         -R_{cefd}R_{ef}\right)\\
    &-\frac1{2\tau^2}R_{ab}.                            \label{eq:har-L-timeRic}
 \end{split}
\end{equation}
Add \eqref{eq:har-L-divP}, \eqref{eq:har-L-RmRic}, and
\eqref{eq:har-L-timeRic}, as prescribed by
\eqref{eq:har-M-divergence}.  The two gradient products cancel.  The eight
$B$-terms cancel in pairs.  The two cubic terms
$R_{acdb}R_{cefd}R_{ef}$ cancel, and the two divergence terms combine with
the first term of \eqref{eq:har-L-divP}.  What remains is exactly
\eqref{eq:har-M-evolution}.
\end{proof}

\begin{corollary}[Trace of the time--time block]
\label{cor:har-trace-M}
One has
\begin{equation}
 h^{ab}M^{(\alpha)}_{ab}
 =\frac12\Delta^UR+|\Ric|^2+\frac{R}{2\tau}
 =\frac12\left(\partial_tR+\frac R\tau\right).          \label{eq:har-trace-M}
\end{equation}
\end{corollary}

\begin{proof}
Trace \eqref{eq:har-M-definition}.  The two curvature--Ricci contractions
combine to $2|\Ric|^2$, while the $-\Ric^2$ term subtracts
$|\Ric|^2$.  Equation~\eqref{eq:har-scalar-evolution} gives the second
equality.
\end{proof}

\section{The Harnack carrier, test jets, and block evolution}
\label{sec:har-block-evolution}

Let
\begin{equation}
       \mathcal H=\operatorname{Alt}^2(V^*)\oplus V^*,  \label{eq:har-carrier}
\end{equation}
where $\operatorname{Alt}^2(V^*)$ is realized as the subbundle of skew
covariant two-tensors.  On skew tensors use the raw fixed norm
\begin{equation}
       |U|_{\mathrm{raw}}^2=U_{ab}U_{ab}.                \label{eq:har-raw-twoform-norm}
\end{equation}
It is twice the usual exterior-power norm.  The raw norm is convenient
because every sum below is over all ordered indices; it changes no notion of
positive semidefiniteness.

Define a symmetric bilinear form $\mathscr Q^{(\alpha)}$ on $\mathcal H$ by
\begin{equation}
 \begin{split}
 \mathscr Q^{(\alpha)}(U\oplus W,V\oplus Z)
 ={}&R_{abdc}U_{ab}V_{cd}\\
    &+P_{abc}(U_{ab}Z_c+V_{ab}W_c)
      +M^{(\alpha)}_{ab}W_aZ_b.                         \label{eq:har-polarized-Q}
 \end{split}
\end{equation}
Pair symmetry of curvature and symmetry of $M^{(\alpha)}$ prove that this
form is symmetric.  Its diagonal is the pullback of
\eqref{eq:har-Q-definition}.

Hamilton's decisive choice is not to extend $U$ and $W$ by ordinary
parallel transport.  Define
\begin{equation}
 S^{(\alpha)}_{ab}=R_{ab}+\frac1{2\tau}h_{ab}            \label{eq:har-S-definition}
\end{equation}
and, for a one-form $W$, set
\begin{equation}
 (\mathcal A_aW)_{bc}
   =\frac12(S^{(\alpha)}_{ab}W_c-S^{(\alpha)}_{ac}W_b). \label{eq:har-triangular-A}
\end{equation}
The triangular spatial connection is
\begin{equation}
 \mathscr D_a(U,W)
    =(\mathcal D_aU-\mathcal A_aW,\mathcal D_aW).        \label{eq:har-triangular-connection}
\end{equation}
It is a genuine connection, but in general it is not compatible with the
fixed product metric on $\mathcal H$.  It is therefore a test-jet device;
we will not invoke Chapter~2's metric-connection theorem verbatim for it.

\begin{lemma}[Pointwise realization of Hamilton's test jets]
\label{lem:har-test-jet-realization}
Fix $(x,t)$ with $\tau>0$ and fix
$U_0\oplus W_0\in\mathcal H_x$.  There are smooth local spacetime
extensions $U,W$ taking these values at $(x,t)$ and satisfying there
\begin{align}
 \mathcal D_aW_b&=0,                                     \label{eq:har-jet-DW}\\
 \mathcal D_aU_{bc}
  &=\frac12(R_{ab}W_c-R_{ac}W_b)
     +\frac1{4\tau}(h_{ab}W_c-h_{ac}W_b),                \label{eq:har-jet-DU}\\
 LW_a&=\frac1\tau W_a,                                  \label{eq:har-jet-LW}\\
 LU_{ab}&=0.                                             \label{eq:har-jet-LU}
\end{align}
\end{lemma}

\begin{proof}
This is finite-jet interpolation, not existence for a parabolic PDE.  Choose
normal coordinates for $g(t)$ at $x$ and a local trivialization of $V$ which
is $D^t$-normal at $x$.  Prescribe the values of $U,W$.  Prescribe their
first spatial derivatives by \eqref{eq:har-jet-DW} and
\eqref{eq:har-jet-DU}.  Choose their second spatial derivatives arbitrarily,
and then prescribe their first time derivatives so that
\eqref{eq:har-jet-LW} and \eqref{eq:har-jet-LU} hold.  A polynomial of
degree two in the spatial coordinates and degree one in time realizes these
jets; multiplying by a local cutoff preserves every jet at the point.
Skew-symmetry is preserved because the prescribed right side in
\eqref{eq:har-jet-DU} is skew in $b,c$.  There is no integrability condition
for a jet prescribed at one point.
\end{proof}

For $\xi=U\oplus W$, put
\begin{align}
 \Sigma(\xi)_{ab}
   &=P_{abc}W_c+R_{abdc}U_{cd},                           \label{eq:har-Sigma}\\
 J(\xi)={}&2R_{acdb}M^{(\alpha)}_{cd}W_aW_b
            -2P_{acd}P_{bdc}W_aW_b                     \notag\\
 &\quad+8R_{adec}P_{dbe}U_{ab}W_c
            +4R_{aefc}R_{befd}U_{ab}U_{cd}.             \label{eq:har-J}
\end{align}

\begin{theorem}[Exact Harnack block evolution]
\label{thm:har-block-evolution}
If $U,W$ satisfy \eqref{eq:har-jet-DW}--\eqref{eq:har-jet-LU} at a point,
then at that point
\begin{equation}
 L\bigl(\mathscr Q^{(\alpha)}(\xi,\xi)\bigr)
      =J(\xi)+\Sigma(\xi)_{ab}\Sigma(\xi)_{ab}.          \label{eq:har-exact-Q-evolution}
\end{equation}
\end{theorem}

\textcolor{red}{Since the above is a major theorem, put the Lean code for the statement here.}

\begin{proof}
We display the computation at equality level.  The product rule, using
$\mathcal DW=0$, $LU=0$, but not yet the other prescribed jets, is
\begin{align}
 L\mathscr Q^{(\alpha)}(\xi,\xi)={}&
 (LM_{ab})W_aW_b+2M_{ab}(LW_a)W_b                       \notag\\
 &+2(LP_{abc})U_{ab}W_c+2P_{abc}U_{ab}(LW_c)
      -4\mathcal D_eP_{abc}\mathcal D_eU_{ab}W_c        \notag\\
 &+(LR_{abdc})U_{ab}U_{cd}
      -4\mathcal D_eR_{abdc}\mathcal D_eU_{ab}U_{cd}   \notag\\
 &-2R_{abdc}\mathcal D_eU_{ab}\mathcal D_eU_{cd}.       \label{eq:har-raw-Q-product}
\end{align}
Here and below $M=M^{(\alpha)}$.  Substitute
\eqref{eq:har-Rm-evolution}, \eqref{eq:har-P-evolution},
\eqref{eq:har-M-evolution}, and the remaining test jets.  Use
\eqref{eq:har-M-divergence} and \eqref{eq:har-divergence-rm} to replace the
two divergences.  Terms containing $\mathcal DP$, $\mathcal DR_{abcd}$,
$\tau^{-2}R_{ab}$, $\tau^{-1}M_{ab}$, and $\tau^{-1}P_{abc}$ cancel in
separate pairs.  The exact expression left before completing a square is
\begin{align}
 L\mathscr Q^{(\alpha)}(\xi,\xi)={}&
 2R_{acdb}M_{cd}W_aW_b-2P_{acd}P_{bdc}W_aW_b            \notag\\
 &+8R_{adec}P_{dbe}U_{ab}W_c
      +4R_{aefc}R_{befd}U_{ab}U_{cd}                    \notag\\
 &+P_{cda}P_{cdb}W_aW_b
      +4R_{adeb}P_{dec}U_{ab}W_c                        \notag\\
 &+R_{abef}R_{dcfe}U_{ab}U_{cd}.                        \label{eq:har-pre-square}
\end{align}
This line is a useful formalization checkpoint: all differential terms have
gone, but no positivity has yet been used.

Expand $\Sigma_{ab}\Sigma_{ab}$ from \eqref{eq:har-Sigma}.  Curvature pair
symmetry, the algebraic first Bianchi identity, skew-symmetry of $U$, and the
cyclic identity \eqref{eq:har-P-cyclic} identify its three summands with the
last two lines of \eqref{eq:har-pre-square}.  The first two lines are
$J(\xi)$.  This proves \eqref{eq:har-exact-Q-evolution}.
\end{proof}

\begin{remark}[Why the frame bundle is not the formal carrier]
\label{rem:har-frame-bundle}
Hamilton's original computation is elegantly organized on the orthonormal
frame bundle.  For this project the frame bundle remains a valuable source
of the formulas, but the fixed bundle $\mathcal H$ is the better theorem
carrier: its fiber metric is independent of time, its elements are ordinary
finite-dimensional tensors, and the only nonmetric object is the explicitly
isolated triangular test connection \eqref{eq:har-triangular-connection}.
\end{remark}

\section{Hamilton's finite-dimensional sum of squares}
\label{sec:har-algebra}

We now forget the manifold.  Let $E$ be an $n$-dimensional Euclidean space,
let $\mathcal A$ be its skew covariant two-tensors, and use full ordered-index
sums.  Suppose
\begin{equation}
 q(U,W)=K_{ab,cd}U_{ab}U_{cd}
        +2P_{abc}U_{ab}W_c+M_{ab}W_aW_b                 \label{eq:har-abstract-q}
\end{equation}
is a positive-semidefinite quadratic form on $\mathcal A\oplus E^*$, where
$K_{ab,cd}=-K_{ba,cd}=-K_{ab,dc}=K_{cd,ab}$,
$P_{abc}=-P_{bac}$, and $M_{ab}=M_{ba}$.  Define
\begin{align}
 J_q(U,W)={}&2K_{ac,bd}M_{cd}W_aW_b
             -2P_{acd}P_{bdc}W_aW_b                    \notag\\
 &+8K_{ad,ce}P_{dbe}U_{ab}W_c
             +4K_{ae,cf}K_{be,df}U_{ab}U_{cd}.          \label{eq:har-abstract-J}
\end{align}

\begin{lemma}[Hamilton's Gram-square lemma]
\label{lem:har-Gram-square}
If $q\geq0$, then $J_q(U,W)\geq0$ for every $(U,W)$.  More precisely, there
are skew tensors $Y^N$ and covectors $X^N$, with
$1\leq N\leq n(n+1)/2$, such that
\begin{align}
 q(U,W)&=\sum_N(Y^N_{ab}U_{ab}+X^N_aW_a)^2,              \label{eq:har-Gram-q}\\
 J_q(U,W)&=\sum_{M,N}
 \left(
  Y^M_{ac}X^N_cW_a-Y^N_{ac}X^M_cW_a
       -2Y^M_{ac}Y^N_{bc}U_{ab}
 \right)^2.                                             \label{eq:har-Gram-J}
\end{align}
\end{lemma}

\begin{proof}
The spectral theorem for the self-adjoint operator representing $q$ gives a
positive-semidefinite square root.  Expand that square root in an
orthonormal basis of $\mathcal A\oplus E^*$; pad with zero rows if necessary.
This gives \eqref{eq:har-Gram-q}.  Comparison of the three blocks gives
\begin{equation}
 K_{ab,cd}=\sum_NY^N_{ab}Y^N_{cd},\qquad
 P_{abc}=\sum_NY^N_{ab}X^N_c,\qquad
 M_{ab}=\sum_NX^N_aX^N_b.                              \label{eq:har-Gram-blocks}
\end{equation}
Substitute these three finite sums into \eqref{eq:har-abstract-J}.  Expand
the right side of \eqref{eq:har-Gram-J}.  The two pure $X$ terms give the
first two terms of \eqref{eq:har-abstract-J}; the mixed terms give its third
term; and the pure $Y$ terms give its fourth.  The remaining terms cancel
after interchanging $M,N$ and using $Y^N_{ab}=-Y^N_{ba}$.  This is a finite
polynomial identity, so no limiting argument is hidden in the proof.
Every summand in \eqref{eq:har-Gram-J} is a real square.
\end{proof}

\begin{corollary}[Reaction tangency to the PSD cone]
\label{cor:har-reaction-tangent}
If $\mathscr Q^{(\alpha)}\geq0$ in one fiber, then
\begin{equation}
 J(\xi)+|\Sigma(\xi)|_{\mathrm{raw}}^2\geq0             \label{eq:har-reaction-nonnegative}
\end{equation}
for every $\xi$ in that fiber.
\end{corollary}

\begin{proof}
Apply Lemma~\ref{lem:har-Gram-square} with
\[
        K_{ab,cd}=R_{abdc},\qquad P=P,\qquad M=M^{(\alpha)}.
\]
The polynomial \eqref{eq:har-abstract-J} is exactly
\eqref{eq:har-J}.  The remaining term is a sum of squares by definition.
\end{proof}

If $q(U,W)=0$, then every factor in \eqref{eq:har-Gram-q} vanishes.  Hence
\begin{equation}
 P_{abc}W_c+K_{ab,cd}U_{cd}
 =\sum_NY^N_{ab}(X^N_cW_c+Y^N_{cd}U_{cd})=0.            \label{eq:har-null-Sigma}
\end{equation}
This equality is useful in rigidity questions, but it is not needed to
prove nonnegativity: the square $|\Sigma|^2$ is harmless at every vector,
not only at null vectors.

\begin{remark}[Why no Schur complement is used]
\label{rem:har-no-Schur}
When the curvature block $K$ is positive definite, minimizing in $U$ gives
the Schur complement $M-P^*K^{-1}P$.  That formulation is useful for
equality analysis, but it is a poor primary proof: $K^{-1}$ is undefined at
curvature null directions and would force a separate rank analysis.  The
Gram proof is polynomial on the entire closed PSD cone and is therefore the
more stable Lean interface.
\end{remark}

\section{A first-contact principle for Harnack test jets}
\label{sec:har-first-contact}

The following lemma is the precise bridge to Chapter~2.  It isolates the
part of that proof which does not require a metric connection on the test
bundle.

\begin{lemma}[Strict rank-one support principle]
\label{lem:har-strict-first-contact}
Let $E\to M$ be a finite-rank vector bundle with a fixed positive-definite
fiber metric, and let $B(t)$ be a continuous family of symmetric bilinear
forms on $E$.  Let $g(t)$ be a smooth family of Riemannian metrics on $M$,
and let $\Delta=\Delta_{g(t)}$ on scalar functions.  Work on a time interval
$(a,b]$.  Assume:
\begin{enumerate}
\item there are $\eta>0$ and a compact set $K\Subset M$ such that
      $B>0$ on $M\times(a,a+\eta]$ and on
      $(M\setminus K)\times[a+\eta,b]$;
\item whenever $B(\cdot,t_0)\geq0$ and
      $0\ne\xi_0\in E_{x_0}$ satisfies
      $B_{x_0,t_0}(\xi_0,\xi_0)=0$, there are a smooth local extension
      $\xi$ of $\xi_0$ and a smooth scalar function $f$ near $(x_0,t_0)$
      such that
      \begin{equation}
       f\geq B(\xi,\xi),\qquad
       f(x_0,t_0)=0,\qquad
       (\partial_t-\Delta)f(x_0,t_0)>0.                \label{eq:har-strict-support-hyp}
      \end{equation}
\end{enumerate}
Then $B(t)>0$ for all $t\in(a,b]$.
\end{lemma}

\begin{proof}
Otherwise assumption 1, compactness of $K$, compactness of the unit sphere
bundle over $K$, and continuity of the least Rayleigh quotient give a first
null point $(x_0,t_0)$ and a nonzero null vector
$\xi_0$.  At time $t_0$, the scalar
$B(\xi,\xi)$ is nonnegative near $x_0$ and vanishes there.  Its smooth upper
support $f$ is therefore also nonnegative and vanishes at $x_0$, so
$\Delta f(x_0,t_0)\geq0$.  At the fixed spatial point, $f$ is nonnegative
for earlier nearby times and equals zero at the first contact, so its left
time derivative is nonpositive.  Therefore
\[
       (\partial_t-\Delta)f(x_0,t_0)\leq0,
\]
contrary to \eqref{eq:har-strict-support-hyp}.  This argument uses neither a
smooth least-eigenvector field nor differentiability of the least
eigenvalue.  It is exactly the rank-one normal-support argument of
Chapter~2.
\end{proof}

\section{The proper barrier and perturbed evolution}
\label{sec:har-proof-main}

The curvature hypothesis and Chapter~8 give all coefficient bounds needed
on the slab.  Translate $\alpha_-$ to time zero and apply
Proposition~\ref{prop:shi-harnack-input} to
$\widetilde g(s)=g(\alpha_-+s)$ with the positive buffer
$\alpha-\alpha_->0$.  It implies
\begin{equation}
 \sup_{M\times[\alpha,b]}
  (|\Rm|+|\nabla\Rm|+|\nabla^2\Rm|)<\infty.             \label{eq:har-two-derivative-slab}
\end{equation}
Hence $P$, $\overline M$, and all coefficients in the Harnack evolution are
uniformly bounded on this slab.  This is why an earlier buffer is chosen
before the Harnack clock is started.

We use an upper-support version of a fixed-distance exhaustion.  A continuous
function $h$ satisfies $\Delta_{g(t)}h\leq Ch$ in the
\emph{upper-support sense} if, at every $(x,t)$ and for every $\varepsilon>0$,
there is a smooth function $h_\varepsilon$ near $x$ such that
\begin{equation}
 h_\varepsilon\geq h,\qquad h_\varepsilon(x)=h(x),\qquad
 \Delta_{g(t)}h_\varepsilon(x)\leq Ch(x)+\varepsilon.   \label{eq:har-upper-support}
\end{equation}

\begin{lemma}[One proper exhaustion for a complete flow slab]
\label{lem:har-proper-exhaustion}
Under the hypotheses of Theorem~\ref{thm:har-shifted-matrix}, assume in
addition that $M$ is connected.  For every $[\alpha,b]\Subset I$ there are
a proper continuous function
$h:M\to[1,\infty)$ and a constant $C_h$ such that, for all
$t\in[\alpha,b]$,
\begin{equation}
        |\nabla h|_{g(t)}\leq C_hh,\qquad
        \Delta_{g(t)}h\leq C_hh                         \label{eq:har-exhaustion-bounds}
\end{equation}
in the upper-support sense.  If $M$ is closed, one may take $h=1$.
\end{lemma}

\begin{proof}
Only the noncompact case needs proof.  Fix $o\in M$ and let
$r(x)=d_{g(\alpha)}(o,x)$.  Choose a smooth nondecreasing function
$\chi:[0,\infty)\to[0,\infty)$ which is constant near zero, equals its
argument for arguments at least two, and has bounded first and second
derivatives.  Set $h=1+\chi\circ r$.  Hopf--Rinow makes $r$, and hence $h$,
proper.

Nonnegative curvature operator implies nonnegative sectional curvature.
Hessian comparison for the complete metric $g(\alpha)$ gives, away from the
cut locus, an upper quadratic-form bound for $\nabla^2_{g(\alpha)}r$.
The Calabi replacement which moves the basepoint a short distance along a
minimizing geodesic gives the same bound by smooth upper supports at a cut
point.  The cutoff $\chi$ removes the singularity at $o$.  Consequently
$h$ has $g(\alpha)$-upper supports with bounded gradient and Hessian.

The curvature bound makes all $g(t)$ uniformly equivalent on the finite
slab.  By \eqref{eq:har-two-derivative-slab}, $|\nabla\Ric|$ is bounded; the
connection-variation identity
\begin{equation}
 \partial_t\Gamma^k_{ij}
  =-g^{k\ell}(\nabla_iR_{j\ell}+\nabla_jR_{i\ell}
                          -\nabla_\ell R_{ij})           \label{eq:har-Christoffel-variation}
\end{equation}
therefore gives a uniform bound for
$\nabla^{g(t)}-\nabla^{g(\alpha)}$.  Transfer each upper support through
these metric and connection comparisons and trace it with $g(t)$.  This
proves \eqref{eq:har-exhaustion-bounds}, after enlarging $C_h$.
\end{proof}

Let $\phi:M\times(\alpha,b]\to(0,\infty)$ and
$\psi:(\alpha,b]\to(0,\infty)$.  Algebraically perturb the two diagonal
blocks by
\begin{align}
 \widehat M_{ab}&=M^{(\alpha)}_{ab}+\frac\phi\tau h_{ab},               \label{eq:har-hat-M}\\
 \widehat R_{abcd}
   &=R_{abcd}+\frac\psi2(h_{ad}h_{bc}-h_{ac}h_{bd}),                   \label{eq:har-hat-Rm}
\end{align}
and leave $P$ unchanged.  These hatted tensors are algebraic barriers; they
are not asserted to be the curvature and $M$ tensor of another Ricci flow.
With the raw norm \eqref{eq:har-raw-twoform-norm}, their block form is
\begin{equation}
 \widehat{\mathscr Q}(U,W)
   =\mathscr Q^{(\alpha)}(U,W)
      +\frac\phi\tau|W|^2+\psi|U|_{\mathrm{raw}}^2.      \label{eq:har-hat-Q}
\end{equation}

\begin{lemma}[Perturbed Harnack evolution]
\label{lem:har-perturbed-evolution}
Fix the slab $[\alpha,b]$.  There is $C<\infty$, depending only on $n$, the
slab, and the bound \eqref{eq:har-two-derivative-slab}, with the following
property.  If $0<\psi\leq1$, and $U,W$ satisfy Hamilton's test jets at a
point, then, in the upper-support sense for $\phi$,
\begin{align}
 L\widehat{\mathscr Q}(U,W)\geq{}&
 \widehat J(U,W)
 +|P_{abc}W_c+\widehat R_{abdc}U_{cd}|_{\mathrm{raw}}^2 \notag\\
 &+\frac1\tau
   \left(L\phi+\frac\phi\tau-\frac{C\psi}\tau-C\phi\right)|W|^2
   +(\psi'-C\psi)|U|_{\mathrm{raw}}^2.                 \label{eq:har-perturbed-evolution}
\end{align}
Here $\widehat J$ is the polynomial \eqref{eq:har-J} formed from
$(\widehat R,P,\widehat M)$.
\end{lemma}

\begin{proof}
Apply $L$ to the two added terms in \eqref{eq:har-hat-Q}, using the same test
jets as in Theorem~\ref{thm:har-block-evolution}.  Before estimating any
term, the exact product rule gives
\begin{align}
 L\widehat{\mathscr Q}={}&J+|\Sigma|_{\mathrm{raw}}^2
 +\left(\frac{L\phi}{\tau}+\frac\phi{\tau^2}\right)|W|^2
 +\psi'|U|_{\mathrm{raw}}^2
 -2\psi|\mathcal DU|_{\mathrm{raw}}^2.                  \label{eq:har-exact-hat-before-error}
\end{align}
Here $\mathcal DW=0$, $LW=W/\tau$, and $LU=0$ explain every added term.
Rewrite the unperturbed reaction and the last gradient square as the hatted
polynomial plus the difference.  Because $J$ has
degree two in the three blocks, that difference is a finite sum of terms
which contain at least one $\phi/\tau$ or $\psi$.  The terms with the most
singular clock weight are bounded by
\[
       \frac{C\psi}{\tau^2}|W|^2;
\]
the remaining mixed terms are bounded, using Young's inequality and
\eqref{eq:har-two-derivative-slab}, by
\[
       \frac{C\phi}{\tau}|W|^2+C\psi|U|_{\mathrm{raw}}^2.
\]
The derivative of $\tau^{-1}$ contributes $+\phi/\tau^2$ after all test-jet
terms are combined, while differentiation of the coefficients contributes
$L\phi/\tau$ and $\psi'$.  Grouping by $|W|^2$ and $|U|^2$ gives exactly
\eqref{eq:har-perturbed-evolution}.  This finite expansion is the perturbed
counterpart of the checkpoint \eqref{eq:har-pre-square}; the constant $C$
is chosen as the maximum of the finitely many contraction bounds and is
independent of the small barrier parameters.
\end{proof}

Choose constants $A,B$ so large that
\begin{equation}
        A-C_h>C,\qquad B>C.                              \label{eq:har-AB-choice}
\end{equation}
Let $S=b-\alpha$.  Choose $\kappa>0$ so small that
\begin{equation}
  C\kappa e^{(B-A)s}<1
       \qquad(0\leq s\leq S).                           \label{eq:har-kappa-choice}
\end{equation}
For $\varepsilon>0$, put
\begin{equation}
 \phi_\varepsilon(x,t)=\varepsilon e^{A\tau}h(x),
 \qquad
 \psi_\varepsilon(t)=\kappa\varepsilon e^{B\tau}.       \label{eq:har-barrier-functions}
\end{equation}
For sufficiently small $\varepsilon$, $\psi_\varepsilon\leq1$ on the slab.
Equations \eqref{eq:har-exhaustion-bounds}--
\eqref{eq:har-kappa-choice} imply
\begin{equation}
 L\phi_\varepsilon>C\phi_\varepsilon,\qquad
 \phi_\varepsilon>C\psi_\varepsilon,\qquad
 \psi_\varepsilon'>C\psi_\varepsilon.                  \label{eq:har-strict-barrier-ineq}
\end{equation}

\begin{proof}[Proof of Theorem~\ref{thm:har-shifted-matrix}]
Fix a target $(x,t,U,W)$ and choose $b>t$.  The connected component of $x$
is open and closed, is complete with the restricted metrics, and inherits
all curvature hypotheses.  Since the conclusion is pointwise, replace $M$
by this component and construct the preceding barriers there.  We first
prove that $\widehat{\mathscr Q}_\varepsilon>0$ on
$M\times(\alpha,b]$.

Choose $C_0$ so that, on the slab,
\[
 |\overline M_{ab}W_aW_b|\leq C_0|W|^2,
 \qquad
 2|P_{abc}U_{ab}W_c|\leq C_0|U|_{\mathrm{raw}}|W|.
\]
Such a constant exists by \eqref{eq:har-two-derivative-slab}.  Since the
curvature block and Ricci tensor are nonnegative, Young's inequality gives
\begin{equation}
 \widehat{\mathscr Q}_\varepsilon(U,W)
 \geq \frac{\psi_\varepsilon}{2}|U|_{\mathrm{raw}}^2
 +\left(
    \frac{\phi_\varepsilon}{\tau}-C_0
       -\frac{C_0^2}{2\psi_\varepsilon}
  \right)|W|^2.                                        \label{eq:har-initial-lower-bound}
\end{equation}
For each fixed $\varepsilon>0$, both barrier functions have a positive
uniform lower bound on the finite slab.  The coefficient of $|W|^2$ in
\eqref{eq:har-initial-lower-bound} therefore tends uniformly to $+\infty$
as $\tau\downarrow0$.  This proves positivity for all sufficiently early
times, uniformly in space.

After removing that early interval, $\tau$ is bounded below.  Properness of
$h$ makes $\phi_\varepsilon/\tau$ tend to infinity as $x$ tends to spatial
infinity.  The same lower bound shows that
$\widehat{\mathscr Q}_\varepsilon$ is positive outside one compact set.
Thus every possible first null point lies in a compact spacetime set.

Suppose a first null point $(x_0,t_0)$ exists, and choose a nonzero null
vector $U_0\oplus W_0$.  Realize Hamilton's test jets by
Lemma~\ref{lem:har-test-jet-realization}.  If $h$ is nonsmooth at $x_0$,
replace it by a smooth upper support from
\eqref{eq:har-upper-support}.  The supported hatted quadratic is no smaller
near $x_0$ and has the same value there, so it is still a valid scalar
first-contact test.  Let $\mathfrak m_0>0$ be the contribution of the last
two lines of \eqref{eq:har-perturbed-evolution} at the nonzero vector
$U_0\oplus W_0$.  Choose the support tolerance $\delta$ so small that
\begin{equation}
 \frac{\varepsilon e^{A\tau_0}\delta}{\tau_0}|W_0|^2
       <\frac12\mathfrak m_0.                            \label{eq:har-support-loss}
\end{equation}
If $W_0=0$, there is no support loss.  Thus the strict differential
inequality survives the upper-support replacement.

At this point the entire hatted block form is positive semidefinite.
Corollary~\ref{cor:har-reaction-tangent}, applied to the hatted blocks, gives
$\widehat J\geq0$; the explicit $\widehat\Sigma$ term is also nonnegative.
The last two coefficients in \eqref{eq:har-perturbed-evolution} are strictly
positive by \eqref{eq:har-strict-barrier-ineq}.  Since
$U_0\oplus W_0\ne0$, we obtain
\[
        L\widehat{\mathscr Q}_\varepsilon(U,W)(x_0,t_0)>0.
\]
This contradicts Lemma~\ref{lem:har-strict-first-contact}.  Therefore
$\widehat{\mathscr Q}_\varepsilon>0$ everywhere on the slab.

Return to the fixed $(x,t,U,W)$.  Equation~\eqref{eq:har-hat-Q} gives
\[
 \mathscr Q^{(\alpha)}(U,W)
 \geq-\frac{\varepsilon e^{A\tau}h(x)}\tau|W|^2
      -\kappa\varepsilon e^{B\tau}|U|_{\mathrm{raw}}^2.
\]
Let $\varepsilon\downarrow0$.  This proves
$\mathscr Q^{(\alpha)}(U,W)\geq0$.  Since $b>t$ was arbitrary and the
Uhlenbeck pullback is an isometry, the intrinsic conclusion
\eqref{eq:har-shifted-conclusion} follows.
\end{proof}

\begin{remark}[Where each global hypothesis is used]
\label{rem:har-hypothesis-use}
Completeness is used only for the proper fixed-distance exhaustion.
Nonnegative curvature operator initializes the two diagonal blocks and
also supplies nonnegative sectional curvature for the simple Hessian
comparison proof of Lemma~\ref{lem:har-proper-exhaustion}.  Slab curvature
bounds, through Chapter~8, control $P$, $M$, and the changing connection.
No strong maximum principle, localized Hamilton--Ivey estimate, Perelman
changing-distance estimate, injectivity-radius theorem, or noncollapsing
theorem enters the proof.
\end{remark}

\section{The trace and integrated Harnack estimates}
\label{sec:har-trace}

\begin{theorem}[Shifted trace Harnack estimate]
\label{thm:har-trace}
Under the hypotheses of Theorem~\ref{thm:har-shifted-matrix}, for every
$\alpha<t$ in $I$, every $x\in M$, and every single tangent vector
$V\in T_xM$,
\begin{equation}
 \boxed{
 \partial_tR+\frac{R}{t-\alpha}
   +2\langle\nabla R,V\rangle+2\Ric(V,V)\geq0.}          \label{eq:har-trace}
\end{equation}
The vector $V$ is pointwise; no global vector field is part of the
statement.
\end{theorem}

\textcolor{red}{Since the above is a major theorem, put the Lean code for the statement here.}

\begin{proof}
Choose a $g(t)$-orthonormal coframe $\{\theta^r\}_{r=1}^n$ at $x$.  For each
$r$ take
\begin{equation}
 W^{(r)}=\theta^r,\qquad
 U^{(r)}=V^\flat\mathbin{\wedge_{\!A}}\theta^r.          \label{eq:har-trace-test-pairs}
\end{equation}
Apply Theorem~\ref{thm:har-shifted-matrix} to every pair and sum in $r$.
There are exactly three contractions.  Corollary~\ref{cor:har-trace-M}
gives
\begin{equation}
 \sum_rM^{(\alpha)}_{ij}(W^{(r)})^i(W^{(r)})^j
    =\frac12\left(\partial_tR+\frac R\tau\right).        \label{eq:har-trace-contraction-M}
\end{equation}
Using the two contractions \eqref{eq:har-P-trace-one} and
\eqref{eq:har-P-trace-two}, and expanding the normalized alternation, gives
\begin{equation}
 \sum_r2P_{pij}(U^{(r)})^{pi}(W^{(r)})^j
       =\langle\nabla R,V\rangle.                       \label{eq:har-trace-contraction-P}
\end{equation}
Finally, the curvature symmetries and two Kronecker contractions give
\begin{equation}
 \sum_rR_{pijq}(U^{(r)})^{pi}(U^{(r)})^{qj}
       =\Ric(V,V).                                      \label{eq:har-trace-contraction-Rm}
\end{equation}
Thus the sum of the nonnegative matrix expressions is
\[
 \frac12\left(\partial_tR+\frac R\tau\right)
      +\langle\nabla R,V\rangle+\Ric(V,V).
\]
Multiply by two and use $\tau=t-\alpha$.
\end{proof}

\begin{corollary}[Fixed-point scalar monotonicity]
\label{cor:har-tR-monotone}
Under the hypotheses and notation of Theorem~\ref{thm:har-trace}, at every
fixed $x$ and every $\alpha,t\in I$ with $\alpha<t$,
\begin{equation}
       \partial_t\bigl((t-\alpha)R(x,t)\bigr)\geq0.     \label{eq:har-tR-monotone}
\end{equation}
In particular, if $\alpha,t_1,t_2\in I$ and $\alpha<t_1\leq t_2$, then
\begin{equation}
 R(x,t_2)\geq\frac{t_1-\alpha}{t_2-\alpha}R(x,t_1).     \label{eq:har-fixed-point-comparison}
\end{equation}
\end{corollary}

\begin{proof}
Set $V=0$ in \eqref{eq:har-trace}, multiply by $t-\alpha$, and integrate.
\end{proof}

The trace estimate can be optimized without pretending that Ricci is
invertible.

\begin{corollary}[Exact semidefinite optimization]
\label{cor:har-semidefinite-optimization}
Under the hypotheses and notation of Theorem~\ref{thm:har-trace}, at any
point $(x,t)$ with $\alpha,t\in I$ and $\alpha<t$, $dR$ annihilates
$\ker\Ric$.  Consequently
there is
$Y\in T_xM$ with
\begin{equation}
                         \Ric(Y,\cdot)=dR,               \label{eq:har-Ric-gradient-preimage}
\end{equation}
and every such $Y$ satisfies
\begin{equation}
 \partial_tR+\frac R{t-\alpha}-\frac12\Ric(Y,Y)\geq0.   \label{eq:har-optimized-semidefinite}
\end{equation}
The number $\Ric(Y,Y)$ does not depend on the choice of $Y$.  If
$\Ric>0$, this becomes
\begin{equation}
 \partial_tR+\frac R{t-\alpha}
 -\frac12(\Ric^{-1})^{ij}\nabla_iR\nabla_jR\geq0.      \label{eq:har-optimized-positive}
\end{equation}
\end{corollary}

\begin{proof}
Put $c=\partial_tR+R/(t-\alpha)$.  If $Z\in\ker\Ric$, substitute $V=sZ$
in \eqref{eq:har-trace} for every real $s$.  The inequality
$c+2s\,dR(Z)\geq0$ for both signs and arbitrarily large $|s|$ forces
$dR(Z)=0$.  For a finite-dimensional self-adjoint map,
$\operatorname{range}(\Ric^\sharp)=(\ker\Ric^\sharp)^\perp$, so
\eqref{eq:har-Ric-gradient-preimage} has a solution.  Put $V=-Y/2$ in
\eqref{eq:har-trace}; this gives \eqref{eq:har-optimized-semidefinite}.
Two solutions differ by a kernel vector, which proves independence.  If
Ricci is positive definite, $Y=\Ric^{-1}(dR)$ is unique.
\end{proof}

We record a path form because it is the version used when comparing
singularity models.  The proof uses Gronwall's inequality and therefore
does not divide by $R$.

\begin{corollary}[Integrated trace Harnack estimate]
\label{cor:har-integrated}
Under the hypotheses of Theorem~\ref{thm:har-shifted-matrix}, let
$\alpha,t_1,t_2\in I$ with $\alpha<t_1<t_2$, and let
$\gamma:[t_1,t_2]\to M$ be piecewise $C^1$.  Then
\begin{equation}
 \begin{split}
 &(t_2-\alpha)R(\gamma(t_2),t_2)\\
 &\quad\geq(t_1-\alpha)R(\gamma(t_1),t_1)
 \exp\!\left[-\frac14\int_{t_1}^{t_2}
                |\dot\gamma(t)|_{g(t)}^2\,dt\right].    \label{eq:har-integrated-path}
 \end{split}
\end{equation}
If $x_1,x_2$ lie in the same connected component, then
\begin{equation}
 R(x_2,t_2)\geq\frac{t_1-\alpha}{t_2-\alpha}
 \exp\!\left[-\frac{d_{g(t_1)}(x_1,x_2)^2}
                         {4(t_2-t_1)}\right]R(x_1,t_1). \label{eq:har-integrated-distance}
\end{equation}
\end{corollary}

\begin{proof}
Set $r(t)=R(\gamma(t),t)$ and use
$V=\dot\gamma(t)/2$ in \eqref{eq:har-trace}.  At differentiability times,
\begin{equation}
 r'(t)+\frac{r(t)}{t-\alpha}
       +\frac12\Ric(\dot\gamma,\dot\gamma)\geq0.        \label{eq:har-path-differential}
\end{equation}
Nonnegative curvature operator implies the sharper algebraic inequality
\begin{equation}
                         2\Ric\leq Rg.                   \label{eq:har-twoRic-less-R}
\end{equation}
To see this, complete a unit vector $e_1$ to an orthonormal basis.  Then
\[
 \Ric(e_1,e_1)=\sum_{j>1}K(e_1,e_j)
 \leq\sum_{i<j}K(e_i,e_j)=\frac R2.
\]
With $y(t)=(t-\alpha)r(t)$, equations
\eqref{eq:har-path-differential} and \eqref{eq:har-twoRic-less-R} give
\[
                    y'(t)\geq-\frac14|\dot\gamma|^2y(t).
\]
Here $y\geq0$.  Gronwall proves \eqref{eq:har-integrated-path}, including
the case where scalar curvature vanishes.

Completeness supplies a minimizing $g(t_1)$-geodesic from $x_1$ to $x_2$.
Parametrize it at constant speed on $[t_1,t_2]$.  Since $\Ric\geq0$,
$\partial_tg=-2\Ric\leq0$, so its $g(t)$-energy is at most its
$g(t_1)$-energy.  Thus
\[
 \int_{t_1}^{t_2}|\dot\gamma|_{g(t)}^2dt
       \leq\frac{d_{g(t_1)}(x_1,x_2)^2}{t_2-t_1}.
\]
Substitute this into \eqref{eq:har-integrated-path}.
\end{proof}

Hamilton's commonly quoted estimate has $2(t_2-t_1)$ rather than
$4(t_2-t_1)$ in the denominator of the exponential; it uses only the
coarser inequality $\Ric\leq Rg$.  Formula
\eqref{eq:har-integrated-distance} is the elementary sharpening available
under the actual nonnegative-curvature-operator hypothesis.  Equivalently,
where $R>0$,
\begin{equation}
       \partial_t\log R+\frac1{t-\alpha}
                    -|\nabla\log R|^2\geq0.             \label{eq:har-log-scalar}
\end{equation}

\section{The ancient matrix and trace limits}
\label{sec:har-ancient}

\begin{theorem}[Ancient matrix Harnack estimate]
\label{thm:har-ancient-matrix}
Let $n\geq2$ and let
\[
                         g(t),\qquad -\infty<t<T,
\]
be a complete ancient Ricci flow with nonnegative curvature operator.
Assume that for every finite compact interval $[a,b]\subset(-\infty,T)$,
\begin{equation}
                       \sup_{M\times[a,b]}|\Rm|<\infty. \label{eq:har-ancient-slab}
\end{equation}
With $P$ as in \eqref{eq:har-P-definition} and $\overline M$ as in
\eqref{eq:har-Mbar-definition}, define
\begin{equation}
 \overline Q(U,W)=\overline M_{ij}W^iW^j
      +2P_{pij}U^{pi}W^j+R_{pijq}U^{pi}U^{qj}.           \label{eq:har-Q-ancient}
\end{equation}
Then $\overline Q(U,W)\geq0$ for every point, time, two-form $U$, and
one-form $W$.
\end{theorem}

\textcolor{red}{Since the above is a major theorem, put the Lean code for the statement here.}

\begin{proof}
Fix $(x,t,U,W)$.  For every $\alpha<t$, apply
Theorem~\ref{thm:har-shifted-matrix} on the ancient interval.  It gives
\[
 \overline Q(U,W)+\frac{\Ric(W^\sharp,W^\sharp)}{2(t-\alpha)}\geq0.
\]
Let $\alpha\to-\infty$.  The second term tends to zero.  This is a pointwise
real-number limit in one fixed geometric tangent space; no convergence of
flows, frames, or $\alpha$-dependent Uhlenbeck bundles is involved.  The
constants in the proof for different $\alpha$ need not be uniform.
\end{proof}

\begin{corollary}[Ancient trace Harnack estimate]
\label{cor:har-ancient-trace}
Under the hypotheses of Theorem~\ref{thm:har-ancient-matrix}, for every
$t<T$, every $x\in M$, and every $V\in T_xM$,
\begin{equation}
 \boxed{
 \partial_tR+2\langle\nabla R,V\rangle+2\Ric(V,V)\geq0.}              \label{eq:har-ancient-trace}
\end{equation}
In particular,
\begin{equation}
                              \partial_tR\geq0.          \label{eq:har-ancient-R-monotone}
\end{equation}
\end{corollary}

\begin{proof}
Either repeat the three contractions
\eqref{eq:har-trace-contraction-M}--
\eqref{eq:har-trace-contraction-Rm} with $\overline M$, or let
$\alpha\to-\infty$ in \eqref{eq:har-trace}.  Set $V=0$ for the last claim.
\end{proof}

The semidefinite optimization in
Corollary~\ref{cor:har-semidefinite-optimization} loses only the clock term:
if $\Ric(Y,\cdot)=dR$, then
\begin{equation}
                       \partial_tR-\frac12\Ric(Y,Y)\geq0.             \label{eq:har-ancient-optimized}
\end{equation}
For $-\infty<t_1<t_2<T$, every piecewise $C^1$ path
$\gamma:[t_1,t_2]\to M$, and points $x_1,x_2$ in the same connected
component, the path and distance forms are
\begin{align}
 R(\gamma(t_2),t_2)
 &\geq R(\gamma(t_1),t_1)
  \exp\!\left[-\frac14\int_{t_1}^{t_2}|\dot\gamma|_{g(t)}^2dt\right], \label{eq:har-ancient-path}\\
 R(x_2,t_2)
 &\geq\exp\!\left[-\frac{d_{g(t_1)}(x_1,x_2)^2}
                    {4(t_2-t_1)}\right]R(x_1,t_1).      \label{eq:har-ancient-distance}
\end{align}
They follow exactly as in Corollary~\ref{cor:har-integrated}.

\section{Normalization and equality regression tests}
\label{sec:har-regression}

The following models test different parts of the formula.  They are not
used in the proof.

\begin{enumerate}
\item For a stationary flat flow,
\[
     \Rm=P=M^{(\alpha)}=\mathscr Q^{(\alpha)}=0.
\]
Every finite and ancient Harnack expression vanishes exactly.

\item For $2\leq m\leq n$, consider a product Ricci flow
\[
 S^m\times\mathbb R^{n-m}
\]
with spherical sectional curvature $k(t)$, and put $r=(m-1)k$.  On the
spherical distribution, $\Ric=rh$, $P=0$, and direct substitution gives
\begin{equation}
 M^{(\alpha)}\big|_{TS^m}
     =\left(r^2+\frac r{2\tau}\right)h,                 \label{eq:har-cylinder-M-test}
\end{equation}
while $M^{(\alpha)}$ vanishes on the Euclidean distribution.  Since
$r'=2r^2$, differentiation of the coefficient in
\eqref{eq:har-cylinder-M-test} gives
\[
 4r^3+\frac{r^2}{\tau}-\frac r{2\tau^2}.
\]
The right side of \eqref{eq:har-M-evolution} gives
\[
 2r\left(r^2+\frac r{2\tau}\right)+2r^3
                       -\frac r{2\tau^2},
\]
the same number.  This checks the cubic sign and the power $\tau^{-2}$.

\item On an expanding gradient soliton in canonical forward form, use the
consistent convention
\begin{equation}
       \Ric+\frac1{2\tau}g=-\nabla^2f.                  \label{eq:har-expander-equation}
\end{equation}
Differentiating \eqref{eq:har-expander-equation} and commuting the third
derivatives gives
\[
             P_{abc}=R_{abc}{}^d\nabla_df.
\]
For every one-form $W$, put
\begin{equation}
       U_{ab}=\frac12(W_a\nabla_bf-W_b\nabla_af).        \label{eq:har-expander-null-U}
\end{equation}
The divergence identity \eqref{eq:har-M-divergence} then gives
\[
 \mathscr Q^{(\alpha)}(U,W)=0,
 \qquad
 P_{abc}W_c+R_{abdc}U_{cd}=0.
\]
Thus expanders supply null directions for the finite-clock Harnack form.
This checks simultaneously the sign of $P$, the order $R_{abdc}$, the
factor $1/(2\tau)$, and the sign in
\eqref{eq:har-expander-null-U}.  In the trace formula the equality vector is
$V=-\nabla f$.
\end{enumerate}

Shrinking ancient spheres and cylinders do not make the trace Harnack
expression vanish identically.  For example, on the spherical factor of
\[
       (-2(m-1)t)g_{S^m(1)}+g_{\mathbb R^{n-m}},\qquad t<0,
\]
one has $R=-m/(2t)$ and $\partial_tR=m/(2t^2)>0$, consistent with the strict
ancient trace inequality at $V=0$.  A cylinder still has matrix-null
directions in its Euclidean factor; no contrary matrix-positivity claim is
being made.

\section{Formalization units and dependency boundary}
\label{sec:har-formalization}

The recommended formalization order is the following.  The identifiers are
stable mathematical interfaces, not claims that Lean declarations already
exist.

\begin{enumerate}
\item \nodeid{PC-F-HAR-CLOCK-DOMAIN} is the pair $(\alpha,t)$ with proof
      $\alpha<t$ and the derived positive number $\tau=t-\alpha$.
      \nodeid{PC-F-HAR-CURVATURE-BLOCK-BRIDGE} is
      \eqref{eq:har-operator-bridge}.  These nodes prevent totalizing
      $1/\tau$ at zero or silently replacing $R_{abdc}$ by $R_{abcd}$.

\item \nodeid{PC-F-HAR-P-SYMMETRIES-AND-TRACES} consists of
      \eqref{eq:har-P-skew}--\eqref{eq:har-divergence-rm};
      \nodeid{PC-F-HAR-M-DIVERGENCE} is
      \eqref{eq:har-M-divergence}.  Their parents are the differential and
      contracted second Bianchi identities.

\item \nodeid{PC-F-HAR-P-EVOLUTION} and
      \nodeid{PC-F-HAR-M-EVOLUTION} are
      Propositions~\ref{prop:har-P-evolution} and
      \ref{prop:har-M-evolution}.  The latter is split internally into the
      three exact nodes \eqref{eq:har-L-divP},
      \eqref{eq:har-L-RmRic}, and \eqref{eq:har-L-timeRic}.
      All depend on the corrected all-slot commutator
      \nodeid{PC-F-SHI-UHLENBECK-HEAT-COMMUTATOR}.

\item \nodeid{PC-F-HAR-TRIANGULAR-TEST-CONNECTION} is
      \eqref{eq:har-triangular-connection}, and
      \nodeid{PC-F-HAR-TEST-JET-REALIZATION} is
      Lemma~\ref{lem:har-test-jet-realization}.  The latter is a local
      finite-jet theorem, not a global PDE solver.

\item \nodeid{PC-F-HAR-RAW-BLOCK-PRODUCT} is
      \eqref{eq:har-raw-Q-product};
      \nodeid{PC-F-HAR-PRE-SQUARE-EVOLUTION} is
      \eqref{eq:har-pre-square}; and
      \nodeid{PC-F-HAR-EXACT-BLOCK-EVOLUTION} is
      Theorem~\ref{thm:har-block-evolution}.  These interfaces keep the
      cancellation calculation independent of positivity.

\item \nodeid{PC-F-HAR-GRAM-BLOCK-DECOMPOSITION} is
      \eqref{eq:har-Gram-q}--\eqref{eq:har-Gram-blocks};
      \nodeid{PC-F-HAR-REACTION-SUM-OF-SQUARES} is the finite polynomial
      identity \eqref{eq:har-Gram-J}; and
      \nodeid{PC-F-HAR-PSD-REACTION-TANGENCY} is
      Corollary~\ref{cor:har-reaction-tangent}.  These are pure
      finite-dimensional real linear algebra and should not import any
      differential geometry.

\item \nodeid{PC-F-HAR-STRICT-RANK-ONE-SUPPORT} is
      Lemma~\ref{lem:har-strict-first-contact};
      \nodeid{PC-F-HAR-PROPER-SLAB-EXHAUSTION} is
      Lemma~\ref{lem:har-proper-exhaustion}; and
      \nodeid{PC-F-HAR-PERTURBED-EVOLUTION} is
      Lemma~\ref{lem:har-perturbed-evolution}.  Their output is the analytic
      engine for \nodeid{PC-F-HAR-SHIFTED-COMPLETE-MATRIX},
      Theorem~\ref{thm:har-shifted-matrix}.

\item The three independent finite contractions
      \nodeid{PC-F-HAR-TRACE-M-CONTRACTION},
      \nodeid{PC-F-HAR-TRACE-P-CONTRACTION}, and
      \nodeid{PC-F-HAR-TRACE-RM-CONTRACTION} are
      \eqref{eq:har-trace-contraction-M}--
      \eqref{eq:har-trace-contraction-Rm}.  Together they prove
      \nodeid{PC-F-HAR-TRACE}, Theorem~\ref{thm:har-trace}.

\item \nodeid{PC-F-HAR-ANCIENT-MATRIX-LIMIT} is
      Theorem~\ref{thm:har-ancient-matrix};
      \nodeid{PC-F-HAR-ANCIENT-TRACE-LIMIT} is
      Corollary~\ref{cor:har-ancient-trace}.  Both limits are pointwise
      limits of real inequalities, so they do not depend on a compactness
      theorem for flows.
\end{enumerate}

The imported chapter-level inputs are: the fixed-bundle metric-freezing and
curvature evolution of Chapter~4; the exact all-slot heat commutator and
positive-slab derivative estimates of Chapter~8; the scalar first-contact
mechanism of Chapter~2; Hessian comparison, the Calabi upper-support
construction, Hopf--Rinow, finite-dimensional Gram factorization, and
Gronwall's inequality.  No result proved later in the Poincar\'e project is
used.

\subsection*{Lean status}

Every new identifier in this chapter currently has Lean statement and Lean
proof status \emph{not yet implemented}.  The displayed decompositions are
formalization plans and mathematical proofs; no typechecking, kernel
checking, axiom-footprint inspection, or source-to-Lean correspondence
claim is being made here.

\section{Sources, convention conversion, and repaired formulas}
\label{sec:har-sources}

The primary source is R.~S.~Hamilton,
\emph{The Harnack estimate for the Ricci flow}, Journal of Differential
Geometry~\textbf{37} (1993), no.~1, 225--243,
\href{https://doi.org/10.4310/jdg/1214453430}
{doi:10.4310/jdg/1214453430}.  Its main theorem is on p.~225; the trace and
integrated consequences are on p.~226; the frame-bundle calculus and
expander motivation are on pp.~227--234; the evolution calculation is on
pp.~234--238; and the perturbation and complete noncompact proof are on
pp.~238--243.

The detailed secondary source used here is Chapter~15 of MSM~144, especially
its equations (15.5), (15.9), (15.12), (15.39)--(15.65), and its
Corollaries~15.3--15.5 and 15.35.  Appendix~F of MSM~144 supplies the moving
orthonormal-frame calculus, while Chapter~8 of this book supplies its
correct fixed-bundle version.  Chapter~10 of GSM~77 gives an independent
concise cross-check of the theorem and its trace consequences.  The
inverse-curvature proof in B.~Chow,
\emph{On an alternate proof of Hamilton's matrix Harnack inequality for the
Ricci flow}, Indiana University Mathematics Journal~\textbf{52} (2003),
863--874, was consulted but not adopted because inversion is ill-suited to
the semidefinite boundary.

Hamilton's original covariant curvature tensor has the opposite sign from
the convention fixed in Chapter~3 and used in MSM~144.  Individual
four-index formulas from the original paper must therefore not be imported
without converting the complete convention.  Our invariant conversion
checkpoint is \eqref{eq:har-operator-bridge}.

The following convention-critical source hazards have been repaired rather
than propagated.
\begin{enumerate}
\item In the proof of Lemma~F.35 of MSM~144, the displayed Ricci-flow
      Christoffel variation is missing its leading minus sign.  Formula
      \eqref{eq:har-Christoffel-variation} has the correct sign.

\item Appendix~F, equation (F.27), orders the last curvature pair
      oppositely from its own derivation in (F.29).  The fixed-bundle master
      identity \eqref{eq:har-master-commutator} uses the derived slot order.

\item Appendix~F, equation (F.32), repeated as equation (15.35), omits the
      factor $2$ contributed by each covariant tensor slot.  Equation
      \eqref{eq:har-master-commutator} is derived for an arbitrary number of
      slots and has the correct factor.  The later calculation of $M$ in
      Chapter~15 in fact uses the corrected coefficients.

\item In equation (15.55), the cubic term inside the parentheses must be
      $-R_{cefd}R_{ef}$, not plus.  The next printed page silently uses the
      minus needed for cancellation.  Equation~\eqref{eq:har-L-timeRic}
      records the corrected identity.

\item Remark~15.17 and its equation (15.44) reverse both the soliton Hessian
      sign and the associated two-form sign relative to the correct
      derivation in Section~1.  The consistent pair is
      \eqref{eq:har-expander-equation} and
      \eqref{eq:har-expander-null-U}; it passes the kernel test following
      that equation.

\item In the hatted reaction displayed on p.~287, the coefficient $4$ in
      front of the quadratic curvature term is omitted.  The Gram identity
      \eqref{eq:har-Gram-J} and the perturbation proof require the
      coefficient $4$ in \eqref{eq:har-J}.

\item The complete theorem retains the spatial curvature bound on every
      compact time slab.  Removing that hypothesis is not justified by this
      proof; MSM~163, Problem~19.29 and Remark~19.30 explicitly warn about
      that issue.  Exercise~15.2 of MSM~144 does justify the positive-slab
      formulation used here, and Chapter~8 supplies the derivative bounds.
\end{enumerate}

Hamilton's source uses a smooth bounded-Hessian exhaustion in the
noncompact perturbation.  Lemma~\ref{lem:har-proper-exhaustion} replaces it
by a fixed-distance upper-support exhaustion, reusing the support calculus
already needed for the Bando--Shi chapter.  This changes no theorem
hypothesis and avoids introducing a global smoothing theorem solely for the
Harnack argument.  The sharpened coefficient $1/4$ in
\eqref{eq:har-integrated-distance} is likewise distinguished from the
classical coefficient $1/2$: it follows from the stronger elementary
inequality $2\Ric\leq Rg$ available under nonnegative curvature operator.

\chapter[Spacetime Reduced Geometry]
{Spacetime Reduced Geometry: Reduced Distance, Reduced Volume, and
Asymptotic Shrinkers}
\label{ch:reduced-geometry}
\index{reduced distance}
\index{reduced volume}
\index{asymptotic shrinker}
\index{$\mathcal L$-geodesic}

\section{Purpose and proof architecture}
\label{sec:red-purpose}

Perelman's reduced geometry assigns a cost to a spacetime path, measured
backward from a regular spacetime pole.  Its normalized value is the reduced
distance $\ell$; the corresponding density
$(4\pi\tau)^{-n/2}e^{-\ell}$ is a weak subsolution of the conjugate heat
equation; and its total mass is the reduced volume.  The monotonicity and
rigidity of this mass turn an ancient $\kappa$-solution, viewed at times
tending to minus infinity, into a nonflat gradient shrinking soliton.  That
last conclusion is the bridge from the analytic chapters of this book to the
classification of ancient three-dimensional $\kappa$-solutions and then to
Ricci flow with surgery.

The chapter is written in arbitrary dimension.  Dimension three enters only
in Corollary~\ref{cor:red-three-asymptotic-shrinker}, where we invoke the
shrinker classification of Chapter~\ref{ch:three-dimensional-shrinkers}.
There are four deliberate features.

\begin{enumerate}
\item The pole is an arbitrary regular spacetime point $(p,t_0)$, and the
      primary path functional is a \emph{clocked segment action}.  This is
      the object which is exactly additive when a path is cut at an
      intermediate time.
\item The change of variable $r=2\sqrt\tau$ is used at the level of path
      spaces, not merely as a formal calculation.  It removes the singular
      endpoint, makes the direct method an ordinary $H^1$ argument, and
      states correctly what regularity an $\mathcal L$-geodesic has at its
      pole.
\item The basic proof of reduced-volume monotonicity uses upper supports,
      semiconcavity, and a weak conjugate-heat inequality.  Thus it does not
      depend on a formalization of the full $\mathcal L$-cut locus and
      $\mathcal L$-Jacobian change-of-variables theorem.
\item In the blow-down theorem, convergence on compact pointed sets is
      separated from convergence of total mass.  A two-point coercivity
      estimate gives a uniform Gaussian tail and prevents reduced-volume
      mass from escaping.  This is the missing step in the older proof
      discussed in Section~\ref{sec:red-sources}.
\end{enumerate}

\begin{roadmap}{Sections~\ref{sec:red-clock}--\ref{sec:red-weak} build the
clocked action, minimizing curves, first and second variations, and the weak
reduced-density inequality.  Sections~\ref{sec:red-volume} and
\ref{sec:red-regular-rigidity} prove reduced-volume monotonicity and the
regular-pole Gaussian equality case.  Sections~\ref{sec:red-kappa}--
\ref{sec:red-asymptotic} prove the two-point estimate and the corrected
asymptotic-shrinker theorem.  Sections~\ref{sec:red-surgery-interface}--
\ref{sec:red-sources} record the surgery interface, formalization forest,
and source corrections.}
\end{roadmap}

Unless stated otherwise, manifolds are nonempty, connected, smooth, without
boundary, and finite dimensional.  All norms, gradients, Laplacians,
curvatures, distances, and measures at backward time $\tau$ are those of
$g(\tau)$.  An inequality between measurable functions is understood almost
everywhere; an inequality involving a Laplacian at a nonsmooth point is
given an explicit support or distributional meaning.

\section{The backward clock and the segment action}
\label{sec:red-clock}

Let $(M^n,\bar g(t))$, $t\in I$, be a smooth forward Ricci flow,
\begin{equation}
                         \partial_t\bar g=-2\Ric_{\bar g}.
                                                               \label{eq:red-forward-flow}
\end{equation}
Fix a regular spacetime pole $(p,t_0)$.  Elapsed backward time and the
backward metric are
\begin{equation}
       \tau=t_0-t,
       \qquad g(\tau)=\bar g(t_0-\tau),
       \qquad \partial_\tau g=2\Ric_g.                 \label{eq:red-backward-flow}
\end{equation}
In particular,
\begin{equation}
 \partial_\tau R=-\Delta R-2|\Ric|^2,
 \qquad
 \partial_\tau d\mu_{g(\tau)}=R\,d\mu_{g(\tau)}.       \label{eq:red-backward-basic-evolution}
\end{equation}
The weight below is $\sqrt\tau$, where $\tau$ is elapsed time from $t_0$;
it is not the square root of an absolute time coordinate.

\begin{definition}[Clocked segment action]
\label{def:red-clocked-action}
Let $0\leq a<b$ and suppose $[t_0-b,t_0-a]\subset I$.  A path
$\gamma:[a,b]\to M$ is called absolutely continuous using the distance of
the fixed metric $g(a)$ (using $g(0)$ when $a=0$).  Any smooth auxiliary
metric gives the same notion along the compact curve graph, by local metric
equivalence.  For such a path define
\begin{equation}
 \mathcal A^{t_0,g}_{a,b}(\gamma)
 :=\int_a^b\sqrt\sigma\left(
       R(\gamma(\sigma),\sigma)
       +|\dot\gamma(\sigma)|_{g(\sigma)}^2\right)d\sigma .
                                                               \label{eq:red-clocked-action}
\end{equation}
as an element of $\mathbb R\cup\{+\infty\}$.  The path has
\emph{finite segment energy} when
\begin{equation}
           \int_a^b\sqrt\sigma\,|\dot\gamma(\sigma)|^2_{g(\sigma)}
                    d\sigma<\infty.                    \label{eq:red-finite-segment-energy}
\end{equation}
Here and below $g(\sigma)=\bar g(t_0-\sigma)$.  If
$\Omega\subset M\times I$ is open, the path is $\Omega$-admissible when
\begin{equation}
                   (\gamma(\sigma),t_0-\sigma)\in\Omega
                   \quad(a\leq\sigma\leq b).             \label{eq:red-domain-admissible}
\end{equation}
For endpoints in the corresponding time slices, put
\begin{equation}
 \mathbf L^{t_0,g,\Omega}_{a,b}(x,y)
 :=\inf\{\mathcal A^{t_0,g}_{a,b}(\gamma):
       \gamma(a)=x,\ \gamma(b)=y,\ \gamma\text{ is }\Omega
       \text{-admissible}\}.                            \label{eq:red-restricted-value}
\end{equation}
The empty infimum is $+\infty$.  When no domain is displayed,
$\Omega=M\times I$.
For a time $t\in I$, write
\begin{equation}
                    \Omega_t:=\{x\in M:(x,t)\in\Omega\}.       \label{eq:red-domain-slice}
\end{equation}
\end{definition}

The action of an individual finite-segment-energy curve is real: its graph is
compact, so the scalar term is bounded, and its kinetic term is finite by
hypothesis.  The infimum in \eqref{eq:red-restricted-value} initially takes
values in the extended real line and can in principle be $-\infty$.  The
following lower bound excludes that value.  After it is imposed, the value
lives in $\mathbb R\cup\{+\infty\}$, the mathematical counterpart of
\texttt{WithTop Real}; its addition is total in the dynamic-programming
formula.

\begin{lemma}[Scalar lower bound for the segment action]
\label{lem:red-action-lower-bound}
If $R\geq-K$ along every admissible curve graph, then
\begin{equation}
 \mathcal A^{t_0,g}_{a,b}(\gamma)
 \geq-\frac{2K}{3}\bigl(b^{3/2}-a^{3/2}\bigr).          \label{eq:red-action-scalar-lower}
\end{equation}
Consequently $\mathbf L_{a,b}$ belongs to
$\mathbb R\cup\{+\infty\}$.  It is a real number whenever at least one
finite-segment-energy admissible curve joins the endpoints.
\end{lemma}

\begin{proof}
Discard the nonnegative kinetic term and integrate
$-K\sqrt\sigma$.  A finite-segment-energy curve has a finite scalar term on a
compact curve graph, so its action is finite.
\end{proof}

\begin{proposition}[Exact concatenation and dynamic programming]
\label{prop:red-dynamic-programming}
Let $0\leq a<b<c$.  If $\gamma_1(a)=x$, $\gamma_1(b)=y=\gamma_2(b)$, and
$\gamma_2(c)=z$, then
\begin{equation}
 \mathcal A^{t_0,g}_{a,c}(\gamma_1*\gamma_2)
 =\mathcal A^{t_0,g}_{a,b}(\gamma_1)
  +\mathcal A^{t_0,g}_{b,c}(\gamma_2).                  \label{eq:red-action-additivity}
\end{equation}
If the scalar lower bound of Lemma~\ref{lem:red-action-lower-bound} holds,
then, in $\mathbb R\cup\{+\infty\}$,
\begin{equation}
 \mathbf L^{t_0,g,\Omega}_{a,c}(x,z)
 =\inf_{y\in\Omega_{t_0-b}}
   \left(\mathbf L^{t_0,g,\Omega}_{a,b}(x,y)
        +\mathbf L^{t_0,g,\Omega}_{b,c}(y,z)\right).    \label{eq:red-dynamic-programming}
\end{equation}
\end{proposition}

\begin{proof}
Equation~\eqref{eq:red-action-additivity} is the decomposition of one
integral over $[a,b]\cup[b,c]$.  Every full admissible path restricts to
two admissible segments and supplies the ``$\geq$'' inequality in
\eqref{eq:red-dynamic-programming}.  For the reverse inequality, first
suppose that the right side of \eqref{eq:red-dynamic-programming} is finite.
Choose $y$ for which the sum is within $\varepsilon$ of that infimum.  The
two summands are then finite, because both have a common finite lower bound.
Choose a path within $\varepsilon$ of each summand, concatenate, and let
$\varepsilon\downarrow0$.  If the right side is $+\infty$ and the left side
were finite, a finite-action $\varepsilon$-competitor for the left side
would have a midpoint and two finite-action restrictions, contradicting the
right side.  Thus both sides are $+\infty$.  This case split avoids an
illicit ``$\varepsilon$-minimizer of $+\infty$''; paths of infinite action
need not be absent.
\end{proof}

\begin{remark}[Do not restart the clock]
\label{rem:red-do-not-restart-clock}
The second summand in \eqref{eq:red-action-additivity} retains the weight
$\sqrt\sigma$.  The reduced action newly based at $(y,t_0-b)$ would instead
have weight $\sqrt{\sigma-b}$.  It is a different functional.  Thus
\eqref{eq:red-dynamic-programming}, not a triangle inequality for re-centered
reduced distances, is the correct interface for surgery.
\end{remark}

\begin{definition}[Finite reduced length and reduced distance from a pole]
\label{def:red-distance}
For $\tau>0$ with $[t_0-\tau,t_0]\subset I$, first define the generalized
reduced length
\begin{equation}
 \overline L_{(p,t_0)}(q,\tau)
 :=\mathbf L^{t_0,g}_{0,\tau}(p,q)
 \quad\text{in }[-\infty,+\infty].                    \label{eq:red-generalized-L}
\end{equation}
If a scalar lower bound excludes $-\infty$ and a finite-action competitor
excludes $+\infty$, let $L_{(p,t_0)}(q,\tau)\in\mathbb R$ be the unique real
number represented by \eqref{eq:red-generalized-L}, and define
\begin{equation}
 \ell_{(p,t_0)}(q,\tau)
 :=\frac{L_{(p,t_0)}(q,\tau)}{2\sqrt\tau}.              \label{eq:red-L-ell}
\end{equation}
Equivalently, for $t<t_0$ we write
$L_{\bar g}((p,t_0),(q,t))$ and
$\ell_{\bar g}((p,t_0),(q,t))$, with $\tau=t_0-t$.
When the pole is fixed, its subscript is suppressed.  Every use of the
real-valued symbols $L$ and $\ell$ includes these two finiteness proofs; an
implementation may package them as hypotheses for an extended-real
\texttt{toReal} operation.  Under \eqref{eq:red-standing-hypothesis} below,
Lemma~\ref{lem:red-action-lower-bound} and a smooth comparison curve supply
the proofs for every endpoint.
\end{definition}

The functional $\mathcal A$, and hence $L$ and $\ell$, may be negative.
They are spacetime action values, not metrics.

\section{Square-root time, coercivity, and minimizers}
\label{sec:red-minimizers}

An ordinary $C^1$ condition on $\gamma$ at $\tau=0$ is wrong: even a
nonconstant Euclidean minimizer has velocity of order $\tau^{-1/2}$.  The
following change of variables supplies the correct path space.

\begin{definition}[Finite energy and square-root smoothness]
\label{def:red-square-root-path}
For a path $\gamma:[0,\bar\tau]\to M$, put
\begin{equation}
       \bar r=2\sqrt{\bar\tau},
       \qquad \beta(r)=\gamma(r^2/4),
       \qquad 0\leq r\leq\bar r.                       \label{eq:red-square-root-path}
\end{equation}
The path $\gamma$ has \emph{finite energy} when $\beta$ is absolutely
continuous and
\begin{equation}
        \int_0^{\bar r}|\beta'(r)|_{g(r^2/4)}^2dr<\infty.
                                                               \label{eq:red-square-root-energy}
\end{equation}
It is \emph{square-root smooth} when $\beta$ is smooth on
$[0,\bar r]$.
\end{definition}

For a pole segment $[0,\bar\tau]$, the two path descriptions preserve the
required absolute continuity.  If $\gamma$ is absolutely continuous, then
$\beta(r)=\gamma(r^2/4)$ is absolutely continuous by composition.  Conversely,
if $\beta$ is absolutely continuous with $\beta'\in L^2$, then
$\gamma(\tau)=\beta(2\sqrt\tau)$ is absolutely continuous because
\begin{equation}
 \int_0^{\bar\tau}|\dot\gamma(\tau)|_{g(\tau)}\,d\tau
       =\int_0^{2\sqrt{\bar\tau}}
          |\beta'(r)|_{g(r^2/4)}\,dr<\infty.            \label{eq:red-AC-reparametrization}
\end{equation}
The substitution also shows that
\eqref{eq:red-square-root-energy} is equivalent to
\eqref{eq:red-finite-segment-energy}.  Thus ``finite energy'' has one
meaning in the singular and clocked descriptions.

\begin{lemma}[Endpoint regularization]
\label{lem:red-square-root-action}
For every finite-energy path,
\begin{equation}
 \mathcal A^{t_0,g}_{0,\bar\tau}(\gamma)
 =\int_0^{2\sqrt{\bar\tau}}
   \left\{|\beta'(r)|_{g(r^2/4)}^2
       +\frac{r^2}{4}R(\beta(r),r^2/4)\right\}dr.       \label{eq:red-square-root-action}
\end{equation}
\end{lemma}

\begin{proof}
Use $\tau=r^2/4$, $d\tau=(r/2)dr$, and
$\dot\gamma=(2/r)\beta'$.  The two factors of $r$ cancel in the kinetic
term.  Formula~\eqref{eq:red-square-root-action}, rather than the singular
formula \eqref{eq:red-clocked-action}, is the foundational integral at the
pole.
\end{proof}

For the remainder of Sections~\ref{sec:red-minimizers}--
\ref{sec:red-regular-rigidity}, we impose the following reusable
hypothesis:
\begin{equation}
 \begin{gathered}
 g(\tau),\quad 0\leq\tau<T,\quad\text{is a smooth complete backward
 Ricci flow},\\
 \text{and for every }T_0<T,\qquad
 \sup_{M\times[0,T_0]}|\Rm|<\infty .
 \end{gathered}                                           \label{eq:red-standing-hypothesis}
\end{equation}
The pole is $(p,0)$ in backward notation.  A pole $(p,t_0)$ in forward
notation is converted by \eqref{eq:red-backward-flow}.  On every compact
backward-time slab strictly inside the domain, Chapter~8 supplies the
curvature-derivative bounds used in the ordinary differential equations
below.

\begin{lemma}[Metric comparison and quadratic action bounds]
\label{lem:red-quadratic-bounds}
Fix $T_0<T$, and suppose
\begin{equation}
       |\Ric|_{\mathrm{op}}\leq\Lambda,
       \qquad |R|\leq C_R
       \quad\text{on }M\times[0,T_0].                  \label{eq:red-curvature-bounds}
\end{equation}
Then, for $0\leq s,\tau\leq T_0$,
\begin{equation}
 e^{-2\Lambda|\tau-s|}g(s)\leq g(\tau)
       \leq e^{2\Lambda|\tau-s|}g(s).                 \label{eq:red-metric-comparison}
\end{equation}
For $0<\tau\leq T_0$ and $r(q)=d_{g(0)}(p,q)$,
\begin{align}
 L(q,\tau)&\geq
 \frac{e^{-2\Lambda\tau}}{2\sqrt\tau}r(q)^2
       -\frac23C_R\tau^{3/2},                          \label{eq:red-L-lower}\\
 L(q,\tau)&\leq
 \frac{e^{2\Lambda\tau}}{2\sqrt\tau}r(q)^2
       +\frac23C_R\tau^{3/2}.                          \label{eq:red-L-upper}
\end{align}
Consequently,
\begin{equation}
 e^{-2\Lambda\tau}\frac{r(q)^2}{4\tau}-\frac{C_R}{3}\tau
 \leq\ell(q,\tau)\leq
 e^{2\Lambda\tau}\frac{r(q)^2}{4\tau}+\frac{C_R}{3}\tau.
                                                               \label{eq:red-ell-quadratic}
\end{equation}
\end{lemma}

\begin{proof}
Integrating
$|\partial_\tau\log g(\tau)(v,v)|\leq2\Lambda$ gives
\eqref{eq:red-metric-comparison}.  Write the action in the regular variable
$r\in[0,2\sqrt\tau]$.  Cauchy--Schwarz gives
\begin{equation}
 d_{g(0)}(p,q)^2\leq2\sqrt\tau
       \int_0^{2\sqrt\tau}|\beta'|_{g(0)}^2dr.          \label{eq:red-energy-distance}
\end{equation}
Metric comparison bounds the kinetic term below, while
\begin{equation}
 \left|\int_0^{2\sqrt\tau}\frac{r^2}{4}R\,dr\right|
       \leq\frac23C_R\tau^{3/2}.                       \label{eq:red-potential-bound}
\end{equation}
This proves \eqref{eq:red-L-lower}.  For the upper bound, traverse a
constant-speed minimizing $g(0)$-geodesic from $p$ to $q$ during the
$r$-interval and apply the other side of metric comparison.  Division by
$2\sqrt\tau$ proves \eqref{eq:red-ell-quadratic}.
\end{proof}

\begin{theorem}[Minimizing $\mathcal L$-curves]
\label{thm:red-minimizer-existence}
Under \eqref{eq:red-standing-hypothesis}, for every $q\in M$ and
$\bar\tau\in(0,T)$ there is a finite-energy path from $(p,0)$ to
$(q,\bar\tau)$ whose action equals $L(q,\bar\tau)$.  Every minimizer is
square-root smooth.  Thus it is smooth for positive backward time, although
it need not be differentiable as a $\tau$-parametrized path at the pole.
\end{theorem}

\begin{proof}
Choose a minimizing sequence $\beta_j$ in the regularized path space of
Definition~\ref{def:red-square-root-path}.  The scalar lower bound and a
fixed finite-action competitor give a uniform kinetic-energy bound.  Hence,
by Cauchy--Schwarz and metric comparison,
\begin{equation}
 d_{g(0)}\bigl(\beta_j(r_1),\beta_j(r_2)\bigr)
       \leq C|r_2-r_1|^{1/2}
 \quad(0\leq r_1,r_2\leq2\sqrt{\bar\tau}),             \label{eq:red-minimizer-holder}
\end{equation}
with one constant $C$.  Taking $r_1=0$ confines all images to one closed
bounded $g(0)$-ball.  Completeness and Hopf--Rinow make this ball compact,
while the full two-point estimate \eqref{eq:red-minimizer-holder} gives
equicontinuity.  Arzel\`a--Ascoli
produces a uniformly convergent subsequence $\beta_j\to\beta$ with the
prescribed endpoints.

Cover the compact image of $\beta$ by finitely many coordinate charts and
subdivide the parameter interval so that each subarc lies in one chart.
The velocities converge weakly in $L^2$ after a further subsequence.  The
metric coefficients converge uniformly along the curves, and the kinetic
integrand is uniformly positive definite and convex in velocity.  Weak
lower semicontinuity therefore gives
\begin{equation}
 \int|\beta'|^2_{g(r^2/4)}dr
 \leq\liminf_j\int|\beta_j'|^2_{g(r^2/4)}dr.             \label{eq:red-weak-lsc}
\end{equation}
The scalar term converges by uniform convergence.  Thus $\beta$ minimizes.

To make the endpoint regularity explicit, in coordinates put
\begin{equation}
 h_{ij}(r,x)=g_{ij}(x,r^2/4),
 \qquad W(r,x)=\frac{r^2}{4}R(x,r^2/4).                 \label{eq:red-coordinate-Lagrangian-data}
\end{equation}
The weak Euler--Lagrange equation is
\begin{equation}
 \frac d{dr}\left(2h_{ij}(r,\beta)\beta'^j\right)
 =\partial_i h_{jk}(r,\beta)\beta'^j\beta'^k
       +\partial_iW(r,\beta)                            \label{eq:red-coordinate-momentum-equation}
\end{equation}
in distributions.  Its right side is $L^1$, so the momentum is
$W^{1,1}$.  Uniform positive definiteness of $h$ first gives
$\beta'\in C^0$; substituting again and bootstrapping gives smoothness
through $r=0$.  This proves every assertion without invoking an
$\mathcal L$-cut-locus theorem.
\end{proof}

\section{Covariance, domains, and convergence interfaces}
\label{sec:red-covariance}

The next propositions state exactly what information about the ambient flow
the action uses.  These facts are elementary, but later surgery arguments
depend on their quantifiers.

\begin{proposition}[Domain calculus]
\label{prop:red-domain-calculus}
Assume the scalar lower bound of Lemma~\ref{lem:red-action-lower-bound} on
the ambient competitor domain, hence also on every restricted domain below.
\begin{enumerate}
\item If $\Omega_1\subset\Omega_2$, then
\begin{equation}
       \mathbf L^{\Omega_2}_{a,b}(x,y)
       \leq\mathbf L^{\Omega_1}_{a,b}(x,y).             \label{eq:red-domain-monotonicity}
\end{equation}
\item Suppose $\mathbf L^\Omega_{a,b}(x,y)$ is finite.  Let
      $\Gamma_{\rm in}$ be the ambient paths whose spacetime graphs stay
      in $\Omega$, and let $\Gamma_{\rm out}$ be the remaining paths.  If
\begin{equation}
 \inf_{\Gamma_{\rm out}}\mathcal A
       \geq \mathbf L^\Omega_{a,b}(x,y)+\eta             \label{eq:red-domain-gap}
\end{equation}
for some $\eta>0$, then the ambient and restricted values agree, and every
      ambient $\varepsilon$-minimizer with $\varepsilon<\eta$ remains in
      $\Omega$.  Here this means an ambient admissible path $\gamma$ with
      $\mathcal A(\gamma)\leq\mathbf L^{\rm ambient}_{a,b}(x,y)+\varepsilon$.
\item If $\Omega_j\nearrow\Omega$ and every compact admissible curve graph
      in $\Omega$ is contained in some $\Omega_j$, then
\begin{equation}
             \mathbf L^{\Omega_j}_{a,b}(x,y)
             \downarrow\mathbf L^\Omega_{a,b}(x,y).     \label{eq:red-domain-exhaustion}
\end{equation}
\end{enumerate}
\end{proposition}

\begin{proof}
The first assertion is inclusion of competitor sets.  For the second,
partition the ambient competitor set into $\Gamma_{\rm in}$ and
$\Gamma_{\rm out}$ and use
\begin{equation}
 \mathbf L^{\rm ambient}
 =\min\left\{\inf_{\Gamma_{\rm in}}\mathcal A,
              \inf_{\Gamma_{\rm out}}\mathcal A\right\}.
                                                               \label{eq:red-domain-partition}
\end{equation}
For the third, monotonicity gives a decreasing limit no smaller than
$\mathbf L^\Omega$.  If $\mathbf L^\Omega=+\infty$, all the restricted
values are $+\infty$ and there is nothing to prove.  Otherwise a compact
$\varepsilon$-minimizing curve is admitted by some $\Omega_j$, which gives
the opposite inequality after $\varepsilon\downarrow0$.
\end{proof}

For a pole value, assume the scalar lower bound and define the
extended-valued restricted reduced distance and its density by
\begin{equation}
 \overline\ell^{\,\Omega}(q,\tau)
   :=\frac{\mathbf L^{\Omega}_{0,\tau}(p,q)}{2\sqrt\tau}
       \quad\text{in }\mathbb R\cup\{+\infty\},
 \qquad
 u^\Omega(q,\tau)
   :=(4\pi\tau)^{-n/2}\operatorname{expNeg}\!\left(
                    \overline\ell^{\,\Omega}(q,\tau)\right),  \label{eq:red-restricted-density}
\end{equation}
where $\operatorname{expNeg}(s)=e^{-s}$ for $s\in\mathbb R$ and
$\operatorname{expNeg}(+\infty)=0$.  When the restricted value is finite we omit the bar
and recover the ordinary real-valued $\ell^\Omega$.  Domain enlargement
decreases $\overline\ell^{\,\Omega}$ and increases $u^\Omega$.  We make no
time-monotonicity claim for the integral of a restricted density: a
boundary-flux theorem would be needed.

\begin{proposition}[Locality and covariance]
\label{prop:red-covariance}
The clocked action has the following properties.
\begin{enumerate}
\item It uses the metric history only on $[t_0-b,t_0-a]$ and on an open
      spacetime neighborhood of the curve graph.
\item Translating forward time and translating $t_0$ by the same amount
      leaves $\mathcal A$, $L$, and $\ell$ unchanged.
\item If $F:N\to M$ is a time-independent diffeomorphism and
      $h(t)=F^*\bar g(t)$, then
\begin{equation}
 \mathcal A^{t_0,h}_{a,b}(\gamma)
  =\mathcal A^{t_0,\bar g}_{a,b}(F\circ\gamma),
 \qquad
 L_h((x,t_0),(y,t))=L_{\bar g}((Fx,t_0),(Fy,t)).         \label{eq:red-pullback-covariance}
\end{equation}
If $F$ is a same-dimensional open embedding (equivalently, a local
diffeomorphism onto an open set) and $h(t)=F^*\bar g(t)$ is the pulled-back
open subflow, the equality is with the $F(N)$-restricted ambient value; the
unrestricted ambient value can only be smaller.  No such assertion is made
for a positive-codimension embedding, because its intrinsic scalar
curvature need not equal the ambient scalar curvature.
\item For $c\in\mathbb R$ and $\lambda>0$, set
\begin{equation}
 \widehat g(s)=\lambda\bar g(c+s/\lambda),
 \qquad \widehat t_0=\lambda(t_0-c),
 \qquad \widehat\gamma(\widehat\sigma)
       =\gamma(\widehat\sigma/\lambda).                 \label{eq:red-parabolic-rescaling}
\end{equation}
Then
\begin{equation}
 \mathcal A^{\widehat t_0,\widehat g}_{\lambda a,\lambda b}
       (\widehat\gamma)=\sqrt\lambda\,
       \mathcal A^{t_0,\bar g}_{a,b}(\gamma),
 \quad \widehat L=\sqrt\lambda L,
 \quad \widehat\ell=\ell.                             \label{eq:red-scaling-action}
\end{equation}
Moreover $\widehat u=\lambda^{-n/2}u$,
$d\mu_{\widehat g}=\lambda^{n/2}d\mu_g$, and reduced volume is invariant.
\end{enumerate}
\end{proposition}

\begin{proof}
Locality follows from the definition of scalar curvature and the metric norm;
agreement merely at the points of a curve would not suffice, because scalar
curvature uses spatial derivatives.  Time translation and pullback
covariance follow by substitution.  Under \eqref{eq:red-parabolic-rescaling},
scalar curvature and squared speed both acquire the factor $\lambda^{-1}$,
elapsed time and its differential acquire factors $\lambda$ and $\lambda$,
and the square-root weight acquires $\sqrt\lambda$.  This gives
\eqref{eq:red-scaling-action}; the remaining scaling laws are immediate.
\end{proof}

\begin{remark}[Time-dependent pullbacks]
\label{rem:red-time-dependent-pullback}
The standard action is not invariant under an arbitrary time-dependent
diffeomorphism.  If the pullback map has spatial generator $W$, the physical
velocity contains $\dot\gamma+W$, so the energy changes to
$|\dot\gamma+W|^2$.  Proposition~\ref{prop:red-covariance} intentionally
asserts covariance only for time-independent pullbacks or genuine spacetime
isometries.
\end{remark}

\begin{proposition}[Stability under smooth pointed convergence]
\label{prop:red-action-convergence}
Fix $\tau>0$.  Suppose backward Ricci flows
$(M_i,g_i(s),p_i)$ on $0\leq s\leq\tau$ converge smoothly on compact
spacetime subsets to $(M_\infty,g_\infty(s),p_\infty)$ through
time-independent open embeddings $\Phi_i:U_i\to M_i$ such that
$\Phi_i(p_\infty)=p_i$, the $U_i$ exhaust $M_\infty$, and
$\Phi_i^*g_i(s)\to g_\infty(s)$ smoothly on every compact subset of
$M_\infty\times[0,\tau]$.  Fix $q\in M_\infty$ and put
$q_i=\Phi_i(q)$ for all sufficiently large $i$.
\begin{enumerate}
\item The action of every fixed finite-energy curve whose graph is compactly
      contained in the convergence region converges to its limiting action.
\item Transporting an $\varepsilon$-minimizer gives
\begin{equation}
 \limsup_{i\to\infty}L_i(q_i,\tau)
       \leq L_\infty(q,\tau).                           \label{eq:red-action-limsup}
\end{equation}
\item If, in addition, the spacetime graph of each chosen
      $i^{-1}$-minimizer is contained in
      $(\Phi_i\times\operatorname{id})(K)$ for one
      compact $K\subset M_\infty\times[0,\tau]$ whose spatial projection
      lies in every sufficiently large $U_i$, then
\begin{equation}
                 L_i(q_i,\tau)\longrightarrow
                 L_\infty(q,\tau).                      \label{eq:red-action-full-convergence}
\end{equation}
\item More generally, fix a time $\theta$ in a common convergence interval,
      suppose $\ell_i(\Phi_i x,\theta)\to\ell_\infty(x,\theta)$ locally
      uniformly, fix $z\in M_\infty$, and put $z_i=\Phi_i(z)$.  Then
\begin{equation}
 \Phi_i^*\bigl(u_i(\cdot,\theta)d\mu_{g_i(\theta)}\bigr)
 \rightharpoonup
 u_\infty(\cdot,\theta)d\mu_{g_\infty(\theta)}          \label{eq:red-local-measure-convergence}
\end{equation}
      weakly on compact sets.  The full masses
      converge under the following chart-tightness condition: for every
      $\varepsilon>0$ there are a compact $K\Subset M_\infty$ and $i_0$,
      with $K\subset U_i$ for $i\geq i_0$, such that
\begin{equation}
 \sup_{i\geq i_0}\int_{M_i\setminus\Phi_i(K)}
             u_i(x,\theta)\,d\mu_{g_i(\theta)}(x)
       <\varepsilon.                                    \label{eq:red-chart-tightness}
\end{equation}
A convenient sufficient condition consists of both radial tightness
\begin{equation}
 \lim_{A\to\infty}\sup_i
 \int_{M_i\setminus B_{g_i(\theta)}(z_i,A)}
             u_i(x,\theta)\,d\mu_{g_i(\theta)}(x)=0     \label{eq:red-tightness-definition}
\end{equation}
and eventual compact ball capture: for every $A<\infty$ there is one compact
$K_A\Subset M_\infty$ such that
$B_{g_i(\theta)}(z_i,A)\subset\Phi_i(K_A)$ for all sufficiently large $i$.
Without ball capture, a disappearing region inside a bounded ambient ball
could carry mass, so radial tightness alone is not sufficient.
Without tightness, only
\begin{equation}
 \int_{M_\infty}u_\infty(x,\theta)d\mu_{g_\infty(\theta)}(x)
       \leq\liminf_i\int_{M_i}u_i(x,\theta)d\mu_{g_i(\theta)}(x)
                                                               \label{eq:red-possible-mass-loss}
\end{equation}
is automatic.
If mass convergence is asserted uniformly for $\theta$ in a compact
interval $J$, local convergence, chart tightness, and any ball-capture
criterion used to prove it must all be uniform in $\theta\in J$.
\end{enumerate}
\end{proposition}

\begin{proof}
Smooth convergence of coefficients and dominated convergence prove the
fixed-curve assertion.  Transport a limiting $\varepsilon$-minimizer to
obtain \eqref{eq:red-action-limsup}.  Under confinement, the regularized
formula \eqref{eq:red-square-root-action} and coercivity give a uniform
$H^1$ bound for minimizing curves.  One-dimensional compactness and weak
lower semicontinuity give the matching liminf.  For the mass statement,
insert a compactly supported cutoff, pass to the limit there, and then use
\eqref{eq:red-chart-tightness}.  Under the sufficient radial criterion,
ball capture permits a compact limiting set to cover the inner ambient ball,
and \eqref{eq:red-tightness-definition} controls its complement.  Exhaustion
without uniform tails gives
only Fatou's inequality \eqref{eq:red-possible-mass-loss}.
\end{proof}

The old pole need not belong to the pointed charts in an ancient blow-down.
In that situation Proposition~\ref{prop:red-action-convergence} is not
misapplied: we extract the functions $\ell_i$ from local estimates and use a
separate two-point bound to prove tightness in
Section~\ref{sec:red-asymptotic}.

\section{First variation and the Hamilton--Jacobi identity}
\label{sec:red-first-variation}

Let $\gamma_u$ be a square-root-smooth variation of a curve, and put
\begin{equation}
        X=\partial_\tau\gamma,
        \qquad Y=\left.\partial_u\gamma_u\right|_{u=0}.
                                                               \label{eq:red-variation-fields}
\end{equation}
The covariant derivatives below use the spacetime family of Levi--Civita
connections in the usual fixed-time way.

\begin{proposition}[First variation]
\label{prop:red-first-variation}
For $0<a<b$,
\begin{align}
 \delta_Y\mathcal A_{a,b}
={}&2\sqrt\tau\,\langle X,Y\rangle\big|_a^b             \label{eq:red-first-variation}\\
 &-2\int_a^b\sqrt\tau\left\langle
   \nabla_X X-\frac12\nabla R+2\Ric^\sharp(X)
       +\frac{1}{2\tau}X,Y\right\rangle d\tau .\nonumber
\end{align}
The same formula holds for $a=0$ and fixed initial point, with zero lower
boundary term.
\end{proposition}

\begin{proof}
Differentiate under the integral.  Metric compatibility in the spatial
variables gives
\begin{equation}
 \delta_Y|X|^2=2\langle\nabla_XY,X\rangle,
 \qquad \delta_YR=\langle\nabla R,Y\rangle.             \label{eq:red-first-variation-raw}
\end{equation}
Since the connection is torsion free, $\nabla_YX=\nabla_XY$.  Along the
two-parameter map,
\begin{equation}
 \frac d{d\tau}\langle Y,X\rangle
 =\langle\nabla_XY,X\rangle+\langle Y,\nabla_XX\rangle
       +2\Ric(Y,X),                                     \label{eq:red-time-pairing}
\end{equation}
where the last term is $\partial_\tau g(Y,X)$.  Integrating
$2\sqrt\tau\langle\nabla_XY,X\rangle$ by parts gives
\eqref{eq:red-first-variation}.  At the pole, a square-root-smooth fixed-end
variation has $Y=O(\sqrt\tau)$, while $\sqrt\tau X=O(1)$, so the lower
boundary term tends to zero.
\end{proof}

\begin{corollary}[$\mathcal L$-geodesic equation]
\label{cor:red-L-geodesic-equation}
A critical path with fixed endpoints satisfies
\begin{equation}
 \nabla_X X-\frac12\nabla R+2\Ric^\sharp(X)
       +\frac{1}{2\tau}X=0.                             \label{eq:red-L-geodesic-equation}
\end{equation}
In square-root time this is the nonsingular equation
\begin{equation}
 \nabla_{\beta'}\beta'+r\Ric^\sharp(\beta')
       -\frac{r^2}{8}\nabla R=0,                       \label{eq:red-L-geodesic-square-root}
\end{equation}
with all tensors evaluated at $(\beta(r),r^2/4)$.
\end{corollary}

\begin{proof}
The fundamental lemma of the calculus of variations applied to
\eqref{eq:red-first-variation} gives
\eqref{eq:red-L-geodesic-equation}.  Substitute
$X=(2/r)\beta'$ and use the time-dependent connection to obtain
\eqref{eq:red-L-geodesic-square-root}.  Alternatively it is the ordinary
Euler--Lagrange equation of \eqref{eq:red-square-root-action}.
\end{proof}

For every $V\in T_pM$, equation
\eqref{eq:red-L-geodesic-square-root} has a unique solution with
\begin{equation}
                  \beta(0)=p,
                  \qquad\beta'(0)=V.                   \label{eq:red-L-initial-data}
\end{equation}
On compact time subintervals, curvature-derivative bounds and
\begin{equation}
 \frac d{dr}|\beta'|^2
 =-r\Ric(\beta',\beta')
       +\frac{r^2}{4}\langle\nabla R,\beta'\rangle      \label{eq:red-L-speed-evolution}
\end{equation}
give a speed bound by Young's inequality and Gronwall.  Completeness then
prevents finite-parameter escape.  Thus the solution exists as long as the
backward flow does.  We define
\begin{equation}
             \mathcal L\exp_\tau(V)=\beta_V(2\sqrt\tau),
 \qquad
             \lim_{\tau\downarrow0}\sqrt\tau\,
                    \dot\gamma_V(\tau)=V.               \label{eq:red-L-exponential}
\end{equation}
We use this map for normalization checks, but the weak monotonicity proof
will not need its cut locus or Jacobian.

\begin{lemma}[Terminal speed of a minimizing curve]
\label{lem:red-terminal-speed}
Fix $0<a<b<T$.  There is $C<\infty$, depending only on $a,b$ and bounded
curvature and curvature-derivative data on a slightly larger slab, such
that every minimizing curve from $(p,0)$ to $(q,\tau)$, $a\leq\tau\leq b$,
satisfies
\begin{equation}
 |\beta'(2\sqrt\tau)|\leq C\bigl(1+d_{g(0)}(p,q)\bigr),
 \qquad
 |X(\tau)|\leq C\bigl(1+d_{g(0)}(p,q)\bigr).           \label{eq:red-terminal-speed}
\end{equation}
The estimate is uniform over all minimizing curves to the same endpoint.
\end{lemma}

\begin{proof}
Choose $b_1\in(b,T)$.  The standing curvature bound and the Bando--Shi
estimates, applied after reversing the backward clock on $[0,b_1]$, give
global bounds for $|\Ric|$ and $|\nabla R|$ on $[0,b]$.  If
$v(r)=|\beta'(r)|$, equation \eqref{eq:red-L-speed-evolution}, interpreted
at zeros by replacing $v$ with $(v^2+\varepsilon)^{1/2}$ and then letting
$\varepsilon\downarrow0$, gives
\begin{equation}
                     |v'(r)|\leq C r v(r)+Cr^2
                     \quad\text{a.e.}                  \label{eq:red-speed-ode-bound}
\end{equation}
for $0\leq r\leq2\sqrt b$.

The upper action bound \eqref{eq:red-L-upper} and the absolute bound on the
scalar term in \eqref{eq:red-square-root-action} give
\begin{equation}
 \int_0^{2\sqrt\tau}v(r)^2dr
       \leq C\bigl(1+d_{g(0)}(p,q)^2\bigr).             \label{eq:red-minimizer-energy-growth}
\end{equation}
Put $\delta=\sqrt a$.  Since $2\sqrt\tau\geq2\sqrt a$, the terminal interval
$[2\sqrt\tau-\delta,2\sqrt\tau]$ has length $\delta$.  Its $L^2$ average
contains a point $r_*$ with
\begin{equation}
                   v(r_*)\leq C\bigl(1+d_{g(0)}(p,q)\bigr).  \label{eq:red-speed-average-point}
\end{equation}
Integrating \eqref{eq:red-speed-ode-bound} from $r_*$ to $2\sqrt\tau$ and
using Gronwall proves the first inequality in
\eqref{eq:red-terminal-speed}.  The identity
$X(\tau)=\beta'(2\sqrt\tau)/\sqrt\tau$ proves the second because
$\tau\geq a$.
\end{proof}

\begin{example}[Gaussian normalization]
\label{ex:red-Gaussian}
For the static Euclidean flow with pole $p$, the solution with initial vector
$V$ is
\begin{equation}
 \beta_V(r)=p+rV,
 \qquad \gamma_V(\tau)=p+2\sqrt\tau V,
 \qquad \ell(\gamma_V(\tau),\tau)=|V|^2.                \label{eq:red-Gaussian-initial-vector}
\end{equation}
Equivalently,
\begin{equation}
 L(q,\tau)=\frac{|q-p|^2}{2\sqrt\tau},
 \qquad \ell(q,\tau)=\frac{|q-p|^2}{4\tau}.             \label{eq:red-Gaussian-distance}
\end{equation}
These formulas fix every normalization in the chapter.
\end{example}

\begin{proposition}[Endpoint upper supports]
\label{prop:red-endpoint-supports}
Let $\gamma$ minimize from $(p,0)$ to $(q,\bar\tau)$, and put
$X=\dot\gamma$.  There is a smooth local upper support $\widehat L\geq L$
touching $L$ at $(q,\bar\tau)$ such that
\begin{equation}
 \nabla\widehat L(q,\bar\tau)=2\sqrt{\bar\tau}\,X(\bar\tau),
 \qquad
 \partial_\tau\widehat L(q,\bar\tau)
   =\sqrt{\bar\tau}\bigl(R-|X|^2\bigr)(q,\bar\tau).    \label{eq:red-endpoint-support-formulas}
\end{equation}
The functions $L$ and $\ell$ are locally Lipschitz on
$M\times(0,T)$.
\end{proposition}

\begin{proof}
Keep the minimizing curve fixed up to a time just before $\bar\tau$ and
join its last point to nearby spacetime endpoints by the unique short
regularized Euler--Lagrange segment in a coordinate neighborhood.  The
action of this broken family is smooth, is at least the infimum $L$, and
agrees with $L$ at the chosen endpoint.  Its spatial derivative is the
upper boundary term in \eqref{eq:red-first-variation}.  To vary the final
time while keeping $q$ fixed, the change of the upper integration limit
contributes $\sqrt{\bar\tau}(R+|X|^2)$ and the endpoint correction contributes
$-2\sqrt{\bar\tau}|X|^2$, giving the second formula.

On a compact spacetime cylinder, Lemma~\ref{lem:red-terminal-speed} gives a
common terminal-speed bound for every minimizing curve, and the construction
uses uniformly bounded short segments.  Append such a segment from one
nearby endpoint to the other and then reverse the comparison.  The two
action inequalities give uniform difference-quotient bounds in space and
time.  Hence $L$, and therefore $\ell$, is locally Lipschitz.
\end{proof}

At a differentiability point, an upper support touching a locally Lipschitz
function has the same differential.  We therefore obtain the exact
almost-everywhere statement below.

\begin{corollary}[Endpoint identities and Hamilton--Jacobi equation]
\label{cor:red-Hamilton-Jacobi}
At every differentiability point $(q,\tau)$, the terminal velocity of every
minimizer is the same and
\begin{equation}
 \nabla L=2\sqrt\tau X,
 \qquad
 \partial_\tau L=\sqrt\tau(R-|X|^2),
 \qquad
 \nabla\ell=X.                                         \label{eq:red-endpoint-identities}
\end{equation}
Consequently
\begin{equation}
 2\partial_\tau\ell+|\nabla\ell|^2-R+\frac\ell\tau=0
                                                               \label{eq:red-Hamilton-Jacobi}
\end{equation}
almost everywhere on $M\times(0,T)$.
\end{corollary}

\begin{proof}
The first two identities are
\eqref{eq:red-endpoint-support-formulas}; division of $L$ by
$2\sqrt\tau$ gives
\begin{equation}
 \partial_\tau\ell
 =\frac12(R-|X|^2)-\frac\ell{2\tau}.                   \label{eq:red-ell-time-endpoint}
\end{equation}
Substitute $X=\nabla\ell$.  If two minimizing curves ended with different
velocities, the first identity would give two values for $\nabla L$.
\end{proof}

\begin{corollary}[Global first-order growth on a positive slab]
\label{cor:red-first-order-growth}
Fix $0<a<b<T$ and put $r(q)=d_{g(0)}(p,q)$.  There is $C<\infty$ such that,
at almost every $(q,\tau)\in M\times[a,b]$,
\begin{equation}
 |\nabla\ell(q,\tau)|\leq C(1+r(q)),
 \qquad
 |\partial_\tau\ell(q,\tau)|\leq C(1+r(q)^2).         \label{eq:red-ell-first-growth}
\end{equation}
\end{corollary}

\begin{proof}
At a differentiability point, \eqref{eq:red-endpoint-identities} and
Lemma~\ref{lem:red-terminal-speed} give the gradient estimate.  Formula
\eqref{eq:red-ell-time-endpoint}, the bounded scalar curvature on the slab,
the gradient estimate, and the quadratic bound
\eqref{eq:red-ell-quadratic} give the time estimate.
\end{proof}

\section{Second variation and the global weak inequalities}
\label{sec:red-weak}

The purpose of the second variation is a trace inequality, not a complete
description of the $\mathcal L$-cut locus.  We expose the finite tensor
calculation needed to obtain it.

\begin{proposition}[$\mathcal L$-index form]
\label{prop:red-index-form}
Let $\gamma$ be an $\mathcal L$-geodesic and let $Y$ be a square-root-smooth
variation field with $Y(0)=0$.  Define
\begin{equation}
 I_{\mathcal L}(Y,Y)
 :=\delta_Y^2\mathcal A^{t_0,g}_{0,\bar\tau}
   -2\sqrt{\bar\tau}\langle\nabla_YY,X\rangle_{\bar\tau}.
                                                               \label{eq:red-index-definition}
\end{equation}
Then
\begin{align}
 I_{\mathcal L}(Y,Y)
={}&\int_0^{\bar\tau}\sqrt\tau\bigl\{
 2|\nabla_XY|^2+\nabla^2R(Y,Y)
 +2\langle\Rm(Y,X)Y,X\rangle                         \label{eq:red-index-form}\\
 &\hspace{5em}-4(\nabla_Y\Ric)(Y,X)
 +2(\nabla_X\Ric)(Y,Y)\bigr\}\,d\tau .\nonumber
\end{align}
\end{proposition}

\begin{proof}
Differentiate \eqref{eq:red-first-variation}.  Commute the two spatial
covariant derivatives using the curvature convention of Chapter~4, and use
\begin{equation}
 (\partial_\tau\nabla)_UV
 =(\nabla_U\Ric^\sharp)(V)+(\nabla_V\Ric^\sharp)(U)
       -(\nabla\Ric)^\sharp(U,V).                       \label{eq:red-connection-variation}
\end{equation}
The endpoint acceleration is precisely the term subtracted in
\eqref{eq:red-index-definition}.  Pairing the remaining connection-variation
terms and using the symmetry of $\Ric$ gives respectively the coefficients
$-4$ and $+2$ in \eqref{eq:red-index-form}.  The pole boundary term vanishes
because $Y=O(\sqrt\tau)$.  This is a finite differentiation-and-commutation
calculation; no minimizing or positivity assertion is used in it.
\end{proof}

For a minimizing $\gamma$ and a terminal vector $E\in T_qM$, take the
comparison field determined by
\begin{equation}
 \nabla_XY=-\Ric^\sharp(Y)+\frac{1}{2\tau}Y,
 \qquad Y(\bar\tau)=E.                                 \label{eq:red-comparison-field}
\end{equation}
Equivalently $Z=Y/\sqrt\tau$ satisfies
$\nabla_XZ=-\Ric^\sharp(Z)$, which is regular in square-root time.  If
$E_1,\ldots,E_n$ is orthonormal at the endpoint, the corresponding fields
satisfy
\begin{equation}
                 \langle Y_i(\tau),Y_j(\tau)\rangle
                 =\frac\tau{\bar\tau}\delta_{ij}.       \label{eq:red-comparison-orthogonality}
\end{equation}

\begin{definition}[Backward Harnack integral]
\label{def:red-Harnack-integral}
Along an $\mathcal L$-geodesic define the matrix expression
\begin{align}
 \mathfrak H(X,Y):={}&-2(\partial_\tau\Ric)(Y,Y)
       -\nabla^2R(Y,Y)+2|\Ric^\sharp(Y)|^2
       -\frac1\tau\Ric(Y,Y)                            \label{eq:red-matrix-Harnack-integrand}\\
 &-2\langle\Rm(Y,X)Y,X\rangle
       -4(\nabla_X\Ric)(Y,Y)+4(\nabla_Y\Ric)(Y,X).\nonumber
\end{align}
Here $\partial_\tau\Ric$ is the fixed-vector time derivative of the
covariant Ricci tensor.  Define also
\begin{align}
 \mathcal H(X)&=-\partial_\tau R-2\langle\nabla R,X\rangle
       +2\Ric(X,X)-\frac R\tau,                         \label{eq:red-Harnack-integrand}\\
 K(\gamma,\bar\tau)&=
       \int_0^{\bar\tau}\tau^{3/2}\mathcal H(X)d\tau. \label{eq:red-K-definition}
\end{align}
\end{definition}

\begin{lemma}[Completed square and traced comparison identity]
\label{lem:red-traced-index-identity}
For every square-root-smooth $Y$ along an $\mathcal L$-geodesic with
$Y(0)=0$,
\begin{align}
 I_{\mathcal L}(Y,Y)+2\sqrt{\bar\tau}\Ric(Y,Y)\big|_{\bar\tau}
={}&\frac{|Y(\bar\tau)|^2}{\sqrt{\bar\tau}}
 -\int_0^{\bar\tau}\sqrt\tau\,\mathfrak H(X,Y)\,d\tau       \label{eq:red-index-completed-square}\\
 &+2\int_0^{\bar\tau}\sqrt\tau\,
   \left|\nabla_XY+\Ric^\sharp(Y)-\frac1{2\tau}Y\right|^2d\tau.\nonumber
\end{align}
Let $E_1,\ldots,E_n$ be orthonormal at $(q,\bar\tau)$ and let $Y_i$ solve
\eqref{eq:red-comparison-field} with $Y_i(\bar\tau)=E_i$.  Then
\begin{equation}
 \sum_{i=1}^n\mathfrak H(X,Y_i)(\tau)
       =\frac\tau{\bar\tau}\mathcal H(X)(\tau),        \label{eq:red-matrix-trace-contraction}
\end{equation}
and consequently
\begin{equation}
 \sum_{i=1}^n I_{\mathcal L}(Y_i,Y_i)
       =-\frac K{\bar\tau}+\frac n{\sqrt{\bar\tau}}
              -2\sqrt{\bar\tau}R(q,\bar\tau).         \label{eq:red-traced-index-identity}
\end{equation}
\end{lemma}

\begin{proof}
Expand the square in \eqref{eq:red-index-completed-square}, use
\begin{equation}
 \frac d{d\tau}|Y|^2
       =2\langle\nabla_XY,Y\rangle+2\Ric(Y,Y),         \label{eq:red-Y-norm-derivative}
\end{equation}
integrate the resulting weighted total derivative, and substitute the index
form \eqref{eq:red-index-form}.  The lower boundary is zero because
$Y=O(\sqrt\tau)$.  After grouping, the remaining terms are exactly the seven
displayed contractions in \eqref{eq:red-matrix-Harnack-integrand}.  This is
a term-by-term polynomial identity in the listed tensors.

For the contraction, \eqref{eq:red-comparison-orthogonality} says that
$\sqrt{\bar\tau/\tau}\,Y_1,\ldots,
\sqrt{\bar\tau/\tau}\,Y_n$ is orthonormal.  Trace each term in
\eqref{eq:red-matrix-Harnack-integrand}; use
$\div\Ric=\frac12\nabla R$ and
$\partial_\tau R=-\Delta R-2|\Ric|^2$.  The result is precisely
\eqref{eq:red-matrix-trace-contraction}.  For the comparison fields the
square in \eqref{eq:red-index-completed-square} vanishes.  Sum over $i$,
use $\sum_i\Ric(E_i,E_i)=R$, and note that
\begin{equation}
 \int_0^{\bar\tau}\sqrt\tau\,
        \frac\tau{\bar\tau}\mathcal H(X)d\tau
       =\frac K{\bar\tau}.                              \label{eq:red-K-trace-integral}
\end{equation}
This proves \eqref{eq:red-traced-index-identity}.
\end{proof}

\begin{lemma}[Full-path traced upper support]
\label{lem:red-full-path-support}
Let $\gamma$ be a minimizing $\mathcal L$-geodesic from the pole to
$(q,\bar\tau)$.  There is a smooth spacetime upper support
$\widehat L_\gamma\geq L$, touching at $(q,\bar\tau)$, whose first
derivatives are \eqref{eq:red-endpoint-support-formulas} and whose spatial
Laplacian satisfies
\begin{equation}
 \Delta\widehat L_\gamma(q,\bar\tau)
 =-\frac K{\bar\tau}+\frac n{\sqrt{\bar\tau}}
                    -2\sqrt{\bar\tau}R(q,\bar\tau).    \label{eq:red-full-support-trace}
\end{equation}
Here $K$ is computed along this chosen minimizer.
\end{lemma}

\begin{proof}
Let $Y_i$ be the full comparison fields from
\eqref{eq:red-comparison-field}.  In endpoint normal coordinates, vary the
whole path by
\begin{equation}
 \gamma_s(\tau)=\exp_{\gamma(\tau)}
                 \left(\sum_{i=1}^n s_iY_i(\tau)\right).  \label{eq:red-full-path-support-variation}
\end{equation}
For sufficiently small $s$, compactness of the curve graph places this
family in finitely many normal neighborhoods.  Since $Y_i(0)=0$ and
$Y_i=O(\sqrt\tau)$, it is a square-root-smooth fixed-pole variation.  The
action of the varied path is a smooth spatial upper support for $L$, and
endpoint normal coordinates remove the acceleration term.  Its Hessian
trace is therefore $\sum_iI_{\mathcal L}(Y_i,Y_i)$, which is
\eqref{eq:red-full-support-trace} by
Lemma~\ref{lem:red-traced-index-identity}.  Extend the family smoothly in
its final time as in Proposition~\ref{prop:red-endpoint-supports}; this gives
the spacetime support and its stated first derivatives without changing the
spatial trace at the touching point.
\end{proof}

\begin{proposition}[Trace comparison and elimination of $K$]
\label{prop:red-trace-comparison}
Let $\gamma$ be a minimizing $\mathcal L$-geodesic ending at
$(q,\tau)$.  Lemma~\ref{lem:red-full-path-support} supplies a smooth
spacetime upper support $\widehat L_\gamma$ for which
\begin{equation}
 \Delta\widehat L_\gamma(q,\tau)
       \leq-\frac K\tau+\frac n{\sqrt\tau}
                    -2\sqrt\tau R.                     \label{eq:red-L-Laplacian}
\end{equation}
At every point where $\ell$ is spacetime differentiable and has a spatial
Alexandrov Hessian, choose any minimizing curve.  The support has the same
first derivatives as $\ell$, and
$\Delta^{\rm A}L\leq\Delta\widehat L_\gamma$.  Consequently
\begin{align}
 \partial_\tau\ell
   &=R-\frac\ell\tau+\frac K{2\tau^{3/2}},              \label{eq:red-ell-time-K}\\
 |\nabla\ell|^2
   &=-R+\frac\ell\tau-\frac K{\tau^{3/2}},             \label{eq:red-ell-gradient-K}\\
 \Delta\ell
   &\leq-R+\frac n{2\tau}-\frac K{2\tau^{3/2}}.        \label{eq:red-ell-Laplacian-K}
\end{align}
where $\Delta\ell$ in these pointwise formulas denotes the Alexandrov
Laplacian.  The same three relations hold for the corresponding upper
support at an arbitrary endpoint, with equality in the first two and an
inequality in the third.
Consequently
\begin{align}
 \partial_\tau\ell-\Delta\ell+|\nabla\ell|^2-R
       +\frac n{2\tau}&\geq0,                           \label{eq:red-ell-main-inequality}\\
 2\Delta\ell-|\nabla\ell|^2+R+\frac{\ell-n}{\tau}
       &\leq0,                                           \label{eq:red-ell-second-inequality}\\
 \partial_\tau\ell+\Delta\ell+\frac\ell\tau
       -\frac n{2\tau}&\leq0.                          \label{eq:red-ell-sum-inequality}
\end{align}
\end{proposition}

\begin{proof}
Use endpoint variations with $(\nabla_{E_i}E_i)(q)=0$.  The minimizing
property bounds the endpoint Hessian of $L$ in direction $E_i$ by
$I_{\mathcal L}(Y_i,Y_i)$.  Trace and apply the already isolated finite
identity \eqref{eq:red-traced-index-identity}.  This is exactly
\eqref{eq:red-L-Laplacian}; no curvature contraction remains hidden in this
step.

Along the geodesic equation \eqref{eq:red-L-geodesic-equation}, direct
differentiation gives
\begin{equation}
 \frac d{d\tau}(R+|X|^2)
       =-\mathcal H(X)-\frac1\tau(R+|X|^2).              \label{eq:red-energy-Harnack-derivative}
\end{equation}
Multiplying by $\tau^{3/2}$ and integrating from the pole yields
\begin{equation}
 \tau^{3/2}(R+|X|^2)=-K+\frac12L.                       \label{eq:red-energy-K-L}
\end{equation}
The lower boundary is zero: bounded scalar curvature gives
$\tau^{3/2}R\to0$, while square-root smoothness gives
$\sqrt\tau X\to V$ and hence
$\tau^{3/2}|X|^2=\sqrt\tau\,|\sqrt\tau X|^2\to0$.
Combine this with \eqref{eq:red-endpoint-identities} and
\eqref{eq:red-L-Laplacian}, and divide by $2\sqrt\tau$, to obtain
\eqref{eq:red-ell-time-K}--\eqref{eq:red-ell-Laplacian-K}.  Eliminating
$K$ gives the first two inequalities.  The third follows by combining the
first with the Hamilton--Jacobi equality
\eqref{eq:red-Hamilton-Jacobi}.
\end{proof}

All three inequalities are equalities for
\eqref{eq:red-Gaussian-distance}.  At a cut point the preceding Hessian
calculation is not classical.  The next theorem records the exact global
replacement.

\begin{theorem}[Supports, semiconcavity, and weak reduced geometry]
\label{thm:red-weak-reduced-geometry}
On every compact cylinder $K\times[a,b]\subset M\times(0,T)$, the function
$\ell$ is locally Lipschitz in spacetime and uniformly semiconcave in the
spatial variables.  Each minimizing $\mathcal L$-geodesic supplies an upper
support at its endpoint for which
\eqref{eq:red-ell-Laplacian-K} holds.  Consequently
\eqref{eq:red-ell-main-inequality} holds weakly in the precise sense that,
for every nonnegative $\psi\in C_c^\infty(M\times(0,T))$,
\begin{equation}
 \int_0^T\!\!\int_M\left\{
  \langle\nabla\ell,\nabla\psi\rangle
 +\left(\partial_\tau\ell+|\nabla\ell|^2-R
          +\frac n{2\tau}\right)\psi\right\}
 d\mu_{g(\tau)}d\tau\geq0.                            \label{eq:red-ell-weak-test}
\end{equation}
\end{theorem}

\begin{proof}
For an endpoint near $(q,\tau)$, vary a fixed minimizing curve and use the
index form on the last short segment.  On a compact cylinder the metric,
curvature, and required derivatives are uniformly bounded, so the resulting
upper supports have a common spatial Hessian upper bound.  In coordinates,
subtracting one fixed quadratic function makes every support concave.  The
standard support characterization of concavity then makes $\ell$ locally
semiconcave.

Lemma~\ref{lem:red-full-path-support} constructs the sharp support carrying
the chosen minimizer's $K$.  It gives
\eqref{eq:red-ell-Laplacian-K} with the $K$ belonging to the chosen
minimizer.  Thus the short-terminal construction supplies uniform
semiconcavity, while the full-path construction supplies the sharp traced
barrier; the two roles are not conflated.

We next remove the minimizer-dependent quantity $K$ before making any
distributional argument.  By Rademacher's and Alexandrov's theorems and
Fubini, at almost every spacetime point $\ell$ has a time derivative, a
spatial gradient, and an Alexandrov spatial Hessian.  Choose any minimizing
$\mathcal L$-geodesic to such a point.  Its upper support has the same first
derivatives as $\ell$, while the Alexandrov Hessian of $\ell$ is bounded
above by the Hessian of the support.  Therefore
\eqref{eq:red-ell-time-K}, \eqref{eq:red-ell-gradient-K}, and
\eqref{eq:red-ell-Laplacian-K} hold there, with the same $K$ attached to the
chosen minimizer.  Eliminating that $K$ at this point gives the
endpoint-independent almost-everywhere inequality
\begin{equation}
 \Delta^{\rm ac}\ell\leq
 \partial_\tau\ell+|\nabla\ell|^2-R+\frac n{2\tau}.
                                                               \label{eq:red-main-ae}
\end{equation}

For completeness, the remaining analytic transfer is the following.  A
locally semiconcave function has a distributional Hessian which is a
matrix-valued Radon measure.  Its singular part is negative semidefinite,
and its absolutely continuous part agrees almost everywhere with the
Alexandrov Hessian.  To prove this, subtract the coordinate quadratic,
mollify the resulting concave function, observe that every mollified Hessian
is negative semidefinite, and take weak-* limits of the Hessian measures.
In a coordinate chart the covariant Hessian differs from the coordinate
Hessian by Christoffel symbols times $df$; since $f$ is locally Lipschitz,
that correction is absolutely continuous.  Thus the singular part is
unchanged, and contraction with the positive-definite tensor $g^{-1}$
preserves its nonpositive sign.
Thus an almost-everywhere upper bound $\Delta^{\rm ac} f\leq F$ passes to
distributions:
\begin{equation}
       -\int\langle\nabla f,\nabla\varphi\rangle d\mu
       \leq\int F\varphi\,d\mu
       \quad(\varphi\geq0).                             \label{eq:red-semiconcavity-transfer}
\end{equation}
Apply this slice by slice to \eqref{eq:red-main-ae}.  Local Lipschitz
regularity supplies the weak time derivative, and Fubini's theorem gives
\eqref{eq:red-ell-weak-test}.  In particular, no measurable selection of a
minimizer and no globally defined version of $K$ is used.
\end{proof}

\begin{definition}[Reduced density]
\label{def:red-density}
The reduced density based at the regular pole is
\begin{equation}
                     u(q,\tau)
       =(4\pi\tau)^{-n/2}e^{-\ell(q,\tau)}.              \label{eq:red-density}
\end{equation}
\end{definition}

\begin{corollary}[Weak conjugate-heat subsolution]
\label{cor:red-density-subsolution}
The reduced density satisfies
\begin{equation}
                  (\partial_\tau-\Delta+R)u\leq0        \label{eq:red-conjugate-subsolution}
\end{equation}
in the following exact distributional sense: for every nonnegative
$\phi\in C_c^\infty(M\times(0,T))$,
\begin{equation}
 \int_0^T\!\!\int_Mu(\partial_\tau\phi+\Delta\phi)
       \,d\mu_{g(\tau)}d\tau\geq0.                    \label{eq:red-density-weak-test}
\end{equation}
Equivalently,
\begin{equation}
 \int_0^T\!\!\int_Mu\left(\partial_\tau\phi
       +\langle\nabla\ell,\nabla\phi\rangle\right)
       d\mu_{g(\tau)}d\tau\geq0.                       \label{eq:red-density-weak-first-order}
\end{equation}
For $0<a<b<T$ and a nonnegative test function compactly supported in space
but allowed to meet the temporal endpoints,
\begin{align}
 &\int_Mu(\cdot,b)\phi(\cdot,b)d\mu_{g(b)}
 -\int_Mu(\cdot,a)\phi(\cdot,a)d\mu_{g(a)}             \label{eq:red-density-endpoint-test}\\
 &\hspace{4em}\leq
 \int_a^b\!\!\int_Mu(\partial_\tau\phi+\Delta\phi)
       d\mu_{g(\tau)}d\tau.\nonumber
\end{align}
The right side of \eqref{eq:red-density-endpoint-test} can equivalently be
written with $\Delta\phi$ replaced by
$\langle\nabla\ell,\nabla\phi\rangle$.
\end{corollary}

\begin{proof}
Use $\psi=u\phi$ in \eqref{eq:red-ell-weak-test}, justified by locally
Lipschitz approximation.  Since
\begin{equation}
 \nabla u=-u\nabla\ell,
 \qquad
 \partial_\tau u=-u\left(\partial_\tau\ell+\frac n{2\tau}\right)
                                                               \label{eq:red-density-derivatives}
\end{equation}
almost everywhere, the two $|\nabla\ell|^2$ terms cancel.  Integrating in
time and using $\partial_\tau d\mu=R\,d\mu$ cancels the scalar-curvature
term.  Spatial integration by parts converts
$u\langle\nabla\ell,\nabla\phi\rangle$ to $u\Delta\phi$.  This proves
\eqref{eq:red-density-weak-test}; retaining the temporal boundary terms
gives \eqref{eq:red-density-endpoint-test}.  On the Gaussian, equality holds,
which checks the sign.
\end{proof}

\section{Reduced volume: normalization and monotonicity}
\label{sec:red-volume}

\begin{definition}[Reduced volume]
\label{def:red-volume}
The reduced volume based at the regular pole $(p,0)$ is
\begin{equation}
             \widetilde V_{(p,0)}(\tau)
       :=\int_M(4\pi\tau)^{-n/2}e^{-\ell(q,\tau)}
                    d\mu_{g(\tau)}(q).                  \label{eq:red-volume}
\end{equation}
\end{definition}

The normalization $(4\pi)^{-n/2}$ makes the Gaussian value equal to one.
We first justify that the integral is finite and that the regular pole has
unit mass at infinitesimal time.

\begin{lemma}[Gaussian integrability on positive time slabs]
\label{lem:red-Gaussian-integrability}
Fix $0<a<b<T$.  There are constants $c,C>0$, depending on the slab and the
curvature bound, such that, with $r(q)=d_{g(0)}(p,q)$,
\begin{align}
 u(q,\tau)&\leq Ce^{-cr(q)^2},                           \label{eq:red-density-Gaussian-tail}\\
 u(q,\tau)|\nabla\ell(q,\tau)|
       &\leq C(1+r(q))e^{-cr(q)^2}                       \label{eq:red-gradient-Gaussian-tail}
\end{align}
for almost every $(q,\tau)\in M\times[a,b]$.  Both right sides are
integrable with respect to $d\mu_{g(\tau)}$, uniformly in $\tau\in[a,b]$.
In particular, $0<\widetilde V(\tau)<\infty$.
\end{lemma}

\begin{proof}
The lower bound in \eqref{eq:red-ell-quadratic} gives
\eqref{eq:red-density-Gaussian-tail}.  Since $\Ric_{g(0)}$ is bounded below,
Bishop comparison bounds the volume of the $k$th unit annulus about $p$ by
$Ce^{Ck}$.  Summing $e^{-ck^2+Ck}$ proves uniform integrability.

Corollary~\ref{cor:red-first-order-growth} gives the required linear growth
of $|\nabla\ell|$ (and, for later differentiation, quadratic growth of
$|\partial_\tau\ell|$).  Multiplication by the Gaussian in
\eqref{eq:red-density-Gaussian-tail} proves
\eqref{eq:red-gradient-Gaussian-tail}.  The same annular sum is integrable.
\end{proof}

\begin{lemma}[$L^1$ differentiation on a moving measure]
\label{lem:red-L1-moving-measure}
Fix $0<a<b<T$, and write
$d\mu_{g(\tau)}=\rho(\cdot,\tau)d\mu_{g(0)}$.  Suppose that for almost every
$q$, the function $\tau\mapsto F(q,\tau)\rho(q,\tau)$ is absolutely
continuous, and that one $H\in L^1(M,d\mu_{g(0)})$ satisfies
\begin{equation}
 |F\rho|+|\partial_\tau(F\rho)|\leq H
       \quad\text{a.e. on }M\times[a,b].               \label{eq:red-L1-AC-domination}
\end{equation}
Then $\tau\mapsto\int_MF(\cdot,\tau)d\mu_{g(\tau)}$ is absolutely continuous
and, for almost every $\tau$,
\begin{equation}
 \frac d{d\tau}\int_MF\,d\mu_{g(\tau)}
       =\int_M(\partial_\tau F+RF)\,d\mu_{g(\tau)}.    \label{eq:red-moving-measure-derivative}
\end{equation}
\end{lemma}

\begin{proof}
Apply the scalar fundamental theorem of calculus to $F(q,\cdot)\rho(q,\cdot)$,
then use Fubini and the domination
\eqref{eq:red-L1-AC-domination}.  This proves the $L^1$-valued fundamental
theorem and absolute continuity of the integral.  Finally
$\partial_\tau\rho=R\rho$ gives
\eqref{eq:red-moving-measure-derivative}.
\end{proof}

\begin{theorem}[Regular-pole normalization]
\label{thm:red-volume-normalization}
As $\tau\downarrow0$,
\begin{equation}
                         \widetilde V(\tau)\longrightarrow1.
                                                               \label{eq:red-volume-normalization}
\end{equation}
More precisely, the measures $u(\cdot,\tau)d\mu_{g(\tau)}$ converge weakly
to the unit point mass at $p$: for every bounded continuous
$\varphi:M\to\mathbb R$,
\begin{equation}
 \lim_{\tau\downarrow0}\int_M\varphi u\,d\mu_{g(\tau)}
                         =\varphi(p).                    \label{eq:red-density-delta-limit}
\end{equation}
\end{theorem}

\begin{proof}
Let $|\Ric|_{\mathrm{op}}\leq\Lambda$ on a short slab.  Then
$|R|\leq n\Lambda$, and
\begin{equation}
 e^{-n\Lambda\tau}d\mu_{g(0)}
 \leq d\mu_{g(\tau)}
 \leq e^{n\Lambda\tau}d\mu_{g(0)}.                    \label{eq:red-measure-comparison}
\end{equation}
The two sides of \eqref{eq:red-ell-quadratic} sandwich
$\widetilde V(\tau)$ between
\begin{align}
 e^{-4n\Lambda\tau/3}
 \int_M(4\pi\tau)^{-n/2}
   e^{-e^{2\Lambda\tau}r^2/(4\tau)}d\mu_{g(0)}          \label{eq:red-volume-sandwich-lower}
\end{align}
and
\begin{align}
 e^{4n\Lambda\tau/3}
 \int_M(4\pi\tau)^{-n/2}
   e^{-e^{-2\Lambda\tau}r^2/(4\tau)}d\mu_{g(0)}.       \label{eq:red-volume-sandwich-upper}
\end{align}
Each integral is a Riemannian Gaussian approximate identity.  Indeed, on a
fixed normal ball at $p$, put $x=\sqrt\tau z$; the metric density tends
uniformly on bounded $z$-sets to the Euclidean density, so the local limit is
\begin{equation}
 (4\pi)^{-n/2}\int_{\mathbb R^n}e^{-|z|^2/4}dz=1.        \label{eq:red-Euclidean-Gaussian-mass}
\end{equation}
The annular Bishop bound used in
Lemma~\ref{lem:red-Gaussian-integrability} makes the complement of the
normal ball tend uniformly to zero.  Both sandwich bounds therefore tend to
one.  Inserting $\varphi$, using its continuity on a small ball and its
boundedness on the Gaussian tail, proves
\eqref{eq:red-density-delta-limit}.
\end{proof}

\begin{theorem}[Reduced-volume monotonicity]
\label{thm:red-volume-monotonicity}
Under \eqref{eq:red-standing-hypothesis}, $\widetilde V$ is locally
absolutely continuous and nonincreasing on $(0,T)$.  Thus
\begin{equation}
 0<\widetilde V(\tau_2)\leq\widetilde V(\tau_1)\leq1
 \qquad(0<\tau_1<\tau_2<T).                            \label{eq:red-volume-monotone}
\end{equation}
At almost every $\tau$,
\begin{equation}
 \frac d{d\tau}\widetilde V(\tau)
 =\int_M\left(-\frac n{2\tau}-\partial_\tau\ell+R\right)
               u\,d\mu_{g(\tau)}\leq0.                 \label{eq:red-volume-derivative}
\end{equation}
\end{theorem}

\begin{proof}
Choose a nonincreasing Lipschitz function $\eta:[0,\infty)\to[0,1]$ equal
to one on $[0,1]$ and zero on $[2,\infty)$, and set
\begin{equation}
              \chi_A(q)=\eta(r(q)/A).                  \label{eq:red-static-cutoff}
\end{equation}
Hopf--Rinow makes its support compact, and metric comparison gives
$|\nabla\chi_A|_{g(\tau)}\leq C/A$ on a fixed slab.  Smooth approximation
allows it as a test in the first-order form of
Corollary~\ref{cor:red-density-subsolution}.  Therefore, for $0<a<b<T$,
\begin{align}
 &\int_M\chi_Au(\cdot,b)d\mu_{g(b)}
 -\int_M\chi_Au(\cdot,a)d\mu_{g(a)}                   \label{eq:red-volume-cutoff-inequality}\\
 &\hspace{3em}\leq
 \int_a^b\!\!\int_Mu\langle\nabla\ell,\nabla\chi_A\rangle
                  d\mu_{g(\tau)}d\tau.\nonumber
\end{align}
By \eqref{eq:red-gradient-Gaussian-tail}, the absolute value of the right
side is at most
\begin{equation}
 \frac CA\int_a^b\!\!\int_{\{A\leq r\leq2A\}}
          (1+r)e^{-cr^2}d\mu_{g(\tau)}d\tau,             \label{eq:red-volume-flux-error}
\end{equation}
which tends to zero.  Dominated convergence on the left gives
$\widetilde V(b)\leq\widetilde V(a)$.  Theorem
\ref{thm:red-volume-normalization} gives the upper bound one.

On a fixed positive slab, metric comparison bounds the measure density
$\rho=d\mu_{g(\tau)}/d\mu_{g(0)}$.  Equations
\eqref{eq:red-density-derivatives},
\eqref{eq:red-density-Gaussian-tail}, and
\eqref{eq:red-ell-first-growth} give one integrable majorant of the form
$C(1+r^2)e^{-cr^2}$ for both $u\rho$ and
$\partial_\tau(u\rho)$.  Apply
Lemma~\ref{lem:red-L1-moving-measure} with $F=u$.  It gives
$\widetilde V\in AC_{\mathrm{loc}}$ and exactly
\eqref{eq:red-volume-derivative}.  We claim differentiability only almost
everywhere, which is the exact consequence of absolute continuity.
\end{proof}

\section{The local defect identity and regular-pole rigidity}
\label{sec:red-regular-rigidity}

The equality argument below will be used twice.  At an ordinary pole it
forces the Gaussian.  On a blow-down whose pole has escaped, it forces only
the gradient-shrinker equation; this distinction is essential.

\begin{lemma}[Perelman's local $v$-identity]
\label{lem:red-v-identity}
Let $u>0$ be a smooth solution of
\begin{equation}
                    (\partial_\tau-\Delta+R)u=0         \label{eq:red-conjugate-heat-equality}
\end{equation}
on an open spacetime set, and write
\begin{equation}
                    u=(4\pi\tau)^{-n/2}e^{-f}.           \label{eq:red-density-potential}
\end{equation}
Define
\begin{equation}
 v=\left[\tau(2\Delta f-|\nabla f|^2+R)+f-n\right]u.    \label{eq:red-v-definition}
\end{equation}
Then
\begin{equation}
 (\partial_\tau-\Delta+R)v
 =-2\tau\left|\Ric+\nabla^2f-\frac1{2\tau}g\right|^2u.
                                                               \label{eq:red-v-identity}
\end{equation}
\end{lemma}

\begin{proof}
Equation~\eqref{eq:red-conjugate-heat-equality} is equivalent to
\begin{equation}
 \partial_\tau f=\Delta f-|\nabla f|^2+R-\frac n{2\tau}.
                                                               \label{eq:red-f-evolution}
\end{equation}
For any scalar $A$,
\begin{equation}
 (\partial_\tau-\Delta+R)(Au)
 =u\{(\partial_\tau-\Delta)A+2\langle\nabla f,\nabla A\rangle\}.
                                                               \label{eq:red-product-conjugate}
\end{equation}
Apply this with the bracket $A$ in \eqref{eq:red-v-definition}.  The local
calculation uses exactly
\begin{align}
 \partial_\tau\Delta f
   &=\Delta\partial_\tau f-2\langle\Ric,\nabla^2f\rangle,
                                                               \label{eq:red-laplacian-evolution-f}\\
 \partial_\tau|\nabla f|^2
   &=-2\Ric(\nabla f,\nabla f)
       +2\langle\nabla f,\nabla\partial_\tau f\rangle,
                                                               \label{eq:red-gradient-evolution-f}\\
 \Delta|\nabla f|^2
   &=2|\nabla^2f|^2+2\langle\nabla f,\nabla\Delta f\rangle
       +2\Ric(\nabla f,\nabla f),                       \label{eq:red-Bochner-f}\\
 \partial_\tau R&=-\Delta R-2|\Ric|^2.                 \label{eq:red-scalar-backward-repeat}
\end{align}
Substitute \eqref{eq:red-f-evolution}, cancel the gradient terms, and
complete the square.  The remainder is
\begin{equation}
 -2\tau|\Ric+\nabla^2f|^2
 +2\langle\Ric+\nabla^2f,g\rangle-\frac n{2\tau},       \label{eq:red-v-pre-square}
\end{equation}
which is exactly the right side of \eqref{eq:red-v-identity} after expanding
the square.  The calculation is local: it uses neither integration by parts
nor compactness.  The factor $u$ on the right is indispensable.
\end{proof}

\begin{theorem}[Rigidity at a regular spacetime pole]
\label{thm:red-regular-pole-rigidity}
Suppose
\begin{equation}
          \widetilde V(\tau_1)=\widetilde V(\tau_2)
          \qquad(0<\tau_1<\tau_2<T).                   \label{eq:red-volume-equality-times}
\end{equation}
Then $(M,g(\tau),p)$ is the static pointed Euclidean flow for
$0\leq\tau\leq\tau_2$.  In particular, equality
$\widetilde V(\tau_0)=1$ at one positive time has the same conclusion on
$[0,\tau_0]$.
\end{theorem}

\begin{proof}
Set
\begin{equation}
                 \mathcal D
       :=-(\partial_\tau-\Delta+R)u.                    \label{eq:red-defect-measure}
\end{equation}
Corollary~\ref{cor:red-density-subsolution} makes $\mathcal D$ a positive
distribution and hence a nonnegative Radon measure.  Monotonicity makes
$\widetilde V$ constant on $[\tau_1,\tau_2]$.  For
$\tau_1<a<b<\tau_2$, test $\mathcal D$ with the spatial cutoff $\chi_A$ and
smooth temporal approximations to $\mathbf1_{[a,b]}$.  The exact result is
\begin{align}
 \langle\mathcal D,\chi_A\mathbf1_{[a,b]}\rangle
={}&\int_M\chi_Au(\cdot,a)d\mu_{g(a)}
   -\int_M\chi_Au(\cdot,b)d\mu_{g(b)}                  \label{eq:red-defect-cutoff}\\
 &+\int_a^b\!\!\int_M
     u\langle\nabla\ell,\nabla\chi_A\rangle d\mu d\tau.\nonumber
\end{align}
The first line tends to
$\widetilde V(a)-\widetilde V(b)=0$, and
\eqref{eq:red-volume-flux-error} sends the last line to zero.  Positivity
then implies $\mathcal D=0$ on
$M\times(\tau_1,\tau_2)$.

Interior regularity for a weak solution of the smooth linear equation makes
$u$, and hence $\ell$, smooth there.  The almost-everywhere Hamilton--Jacobi
identity now holds everywhere, while the conjugate-heat equality gives
\begin{equation}
 \partial_\tau\ell-\Delta\ell+|\nabla\ell|^2-R
                    +\frac n{2\tau}=0.                 \label{eq:red-ell-conjugate-equality}
\end{equation}
Combining it with \eqref{eq:red-Hamilton-Jacobi} yields
\begin{equation}
 \tau(2\Delta\ell-|\nabla\ell|^2+R)+\ell-n=0.          \label{eq:red-v-zero-bracket}
\end{equation}
Thus $v\equiv0$, and Lemma~\ref{lem:red-v-identity} gives
\begin{equation}
                   \Ric+\nabla^2\ell=\frac1{2\tau}g
       \quad\text{on }M\times(\tau_1,\tau_2).           \label{eq:red-regular-pole-soliton}
\end{equation}

Fix $\tau_*\in(\tau_1,\tau_2)$.  The complete bounded-curvature shrinker
in \eqref{eq:red-regular-pole-soliton} generates its canonical self-similar
backward flow.  On $[\varepsilon,\tau_*]$, forward uniqueness after reversing
the backward clock identifies it with the given flow.  Hence
\begin{equation}
 g(\tau)=\frac\tau{\tau_*}\Phi_\tau^*g(\tau_*),
 \qquad
 \sup_M|\Rm(g(\tau))|
  =\frac{\tau_*}{\tau}\sup_M|\Rm(g(\tau_*))|.           \label{eq:red-regular-pole-scaling}
\end{equation}
The original flow has uniformly bounded curvature as $\tau\downarrow0$.
Therefore $\Rm(g(\tau_*))=0$.  Varying $\tau_*$, using smoothness at the
endpoints, and using forward uniqueness shows that $g$ is flat and static on
$[0,\tau_2]$.

For a static complete flat flow the action can be minimized directly in
square-root time, as in Example~\ref{ex:red-Gaussian}; hence
\begin{equation}
                         \ell(q,\tau)=\frac{d_g(p,q)^2}{4\tau}.
                                                               \label{eq:red-flat-reduced-distance}
\end{equation}
Equation~\eqref{eq:red-regular-pole-soliton} says this squared distance is
smooth away from $p$ and has Hessian $2g$ there.  Thus $p$ has empty cut
locus.  The exponential map at $p$ is a global diffeomorphism and local
Euclidean isometry, hence a pointed global isometry
$(M,g,p)\cong(\mathbb R^n,g_{\rm E},0)$.

If $\widetilde V(\tau_0)=1$, Theorem~\ref{thm:red-volume-normalization} and
monotonicity make it constant between zero and $\tau_0$; apply the result on
interior pairs and take their union.
\end{proof}

\begin{remark}[Two equality regimes]
\label{rem:red-two-rigidity-regimes}
The proof used bounded curvature through the regular time $\tau=0$ to turn
the self-similar scaling into flatness.  In an ancient blow-down, the limiting
flow is defined only for $\theta>0$, the old pole may escape, and this
argument is unavailable.  Constant limiting mass then gives a general
gradient shrinker, not necessarily the Gaussian.
\end{remark}

\section[Ancient kappa-solutions and reduced distance]
{Ancient \texorpdfstring{$\kappa$}{kappa}-solutions and reduced-distance
estimates}
\label{sec:red-kappa}

We now pass to the ancient setting.  The forward flow is denoted
$\bar g(t)$, $-\infty<t\leq0$, and the pole is $(p,0)$.  Thus
\begin{equation}
                         g(\tau)=\bar g(-\tau),
                         \qquad 0\leq\tau<\infty         \label{eq:red-ancient-backward-flow}
\end{equation}
is a backward Ricci flow.
Throughout the ancient analysis we assume $n\geq2$.  This is automatic for
a nonflat Riemannian manifold with nonnegative curvature operator, but it is
kept explicit because the changing-distance estimate contains $n-1$.

\begin{definition}[$\kappa$-noncollapsing and ancient $\kappa$-solution]
\label{def:red-kappa-solution}
A complete Riemannian metric $h$ is $\kappa$-noncollapsed on all scales when,
for every $x$ and $r>0$, the curvature bound
\begin{equation}
                 |\Rm_h|\leq r^{-2}
                 \quad\text{on }B_h(x,r)                \label{eq:red-noncollapse-curvature-ball}
\end{equation}
implies
\begin{equation}
                  \Vol_hB_h(x,r)\geq\kappa r^n.
                                                               \label{eq:red-kappa-noncollapse}
\end{equation}
A complete Ricci flow is $\kappa$-noncollapsed on all scales when every time
slice has this property with the same $\kappa$.
An \emph{ancient $\kappa$-solution} is a connected, complete, nonflat
ancient Ricci flow with nonnegative curvature operator,
$\kappa$-noncollapsing on all scales, and one global spacetime curvature
bound.  In the scalar formulation used here, the last hypothesis is
\begin{equation}
          \text{there exists }C<\infty\text{ such that }
          0\leq R(x,t)\leq C
          \quad\text{for every }(x,t)\in M\times(-\infty,0].  \label{eq:red-kappa-global-bound}
\end{equation}
Nonnegative curvature converts this to a global bound for $|\Rm|$.  Notice
that ``bounded on each time slice'' has deliberately not been substituted
for \eqref{eq:red-kappa-global-bound}; doing so and then invoking the ancient
Harnack estimate would be circular, because the Harnack theorem itself uses
bounded-curvature slab hypotheses.
\end{definition}

The norm bridge from nonnegative curvature operator to scalar curvature is
dimension dependent.  We use only
\begin{equation}
                  0\leq\Ric\leq Rg,
                  \qquad |\Rm|\leq C_nR.                \label{eq:red-nonnegative-curvature-bridges}
\end{equation}
No dimension-independent identification of $R$ and $|\Rm|$ is intended.

\begin{proposition}[Ancient reduced-distance estimates]
\label{prop:red-ancient-ell-estimates}
Let $\ell$ be based at $(p,0)$ on an ancient $\kappa$-solution of dimension
$n\geq2$.  For every minimizing $\mathcal L$-geodesic $\gamma$ ending at
$(q,\tau)$,
\begin{equation}
       R(q,\tau)+|\dot\gamma(\tau)|^2
       \leq\frac{3\ell(q,\tau)}\tau.                   \label{eq:red-ancient-endpoint-estimate}
\end{equation}
Consequently, almost everywhere,
\begin{align}
 |\nabla\ell|^2+R&\leq\frac{3\ell}\tau,                 \label{eq:red-ancient-gradient-estimate}\\
 -\frac{2\ell}\tau\leq\partial_\tau\ell
       &\leq\frac\ell\tau.                             \label{eq:red-ancient-time-estimate}
\end{align}
In particular,
\begin{equation}
       |\nabla\sqrt\ell|\leq\frac{\sqrt3}{2\sqrt\tau},
                                                               \label{eq:red-sqrt-ell-gradient}
\end{equation}
interpreted by approximation with $\sqrt{\ell+\varepsilon}$ at a possible
zero.  For $0<\tau_1,\tau_2<\infty$,
\begin{equation}
 \ell(q,\tau_2)
 \leq\ell(q,\tau_1)
        \max\left\{(\tau_2/\tau_1)^2,
                    (\tau_1/\tau_2)^2\right\}.          \label{eq:red-ell-time-ratio}
\end{equation}
\end{proposition}

\begin{proof}
The ancient trace Harnack estimate of Chapter~9, converted to backward time,
states
\begin{equation}
 -\partial_\tau R-2\langle\nabla R,X\rangle
       +2\Ric(X,X)\geq0.                                \label{eq:red-ancient-Harnack-backward}
\end{equation}
Thus the integrand \eqref{eq:red-Harnack-integrand} satisfies
$\mathcal H(X)\geq-R/\tau$.  Write
\begin{equation}
 A=\int_0^\tau\sqrt s\,R\,ds,
 \qquad B=\int_0^\tau\sqrt s\,|X|^2ds,
 \qquad L=A+B.                                          \label{eq:red-action-A-B}
\end{equation}
Since $R\geq0$, $A,B\geq0$, and $K\geq-A$.  Formula
\eqref{eq:red-energy-K-L} gives
\begin{equation}
 \tau^{3/2}(R+|X|^2)\leq A+\frac12(A+B)\leq\frac32L,  \label{eq:red-energy-ancient-bound}
\end{equation}
which is \eqref{eq:red-ancient-endpoint-estimate}.  At differentiability
points $X=\nabla\ell$, giving
\eqref{eq:red-ancient-gradient-estimate}.  Substitute this and $R\geq0$ in
the Hamilton--Jacobi identity to get
\eqref{eq:red-ancient-time-estimate}.  The spatial estimate follows by the
chain rule.  For the time-ratio estimate, do not divide by $\ell$, which may
vanish.  The upper inequality in
\eqref{eq:red-ancient-time-estimate} says
$(\tau^{-1}\ell)'\leq0$ almost everywhere, while the lower one says
$(\tau^2\ell)'\geq0$.  The function $\ell(q,\cdot)$ is locally absolutely
continuous on $(0,\infty)$, so integration gives the sharp one-sided bounds
and hence the coarser symmetric bound \eqref{eq:red-ell-time-ratio}.  This
argument includes the zero case.
\end{proof}

\begin{definition}[Upper right Dini derivative]
\label{def:red-upper-Dini}
For a real-valued function $F$ defined to the right of $s$, set
\begin{equation}
 D^+F(s):=\limsup_{h\downarrow0}
              \frac{F(s+h)-F(s)}h
       \quad\text{in }[-\infty,+\infty].               \label{eq:red-upper-Dini-definition}
\end{equation}
An inequality $D^+F(s)\leq c$ with $c\in\mathbb R$ uses the extended-real
order and, in particular, excludes the value $+\infty$.
\end{definition}

\begin{lemma}[Integration of an upper Dini inequality]
\label{lem:red-Dini-integration}
Let $F:[\alpha,\beta]\to\mathbb R$ be continuous, locally absolutely
continuous on $(\alpha,\beta]$, and let
$h:(\alpha,\beta)\to\mathbb R$ be finite-valued with
$\int_\alpha^\beta|h|<\infty$.  If
\begin{equation}
                       D^+F(s)\leq h(s)                 \label{eq:red-Dini-hypothesis}
\end{equation}
for every $s\in(\alpha,\beta)$, then
\begin{equation}
                 F(\beta)-F(\alpha)
                    \leq\int_\alpha^\beta h(s)\,ds.    \label{eq:red-Dini-conclusion}
\end{equation}
The same conclusion holds when \eqref{eq:red-Dini-hypothesis} is known
almost everywhere and $F$ is absolutely continuous on the whole interval.
\end{lemma}

\begin{proof}
On $[\alpha+\varepsilon,\beta]$, the ordinary derivative exists almost
everywhere and is bounded above by the upper right Dini derivative.  Apply
the fundamental theorem for absolutely continuous functions and let
$\varepsilon\downarrow0$, using continuity of $F$ and integrability of $h$.
The final assertion is the same argument without the limiting step.
\end{proof}

\begin{proposition}[Existence of an $\ell$-center]
\label{prop:red-ell-center}
For every $\tau>0$ there is $q_\tau\in M$ such that
\begin{equation}
                         \ell(q_\tau,\tau)\leq\frac n2.
                                                               \label{eq:red-ell-center}
\end{equation}
If $q_i=q_{\tau_i}$ and
\begin{equation}
 g_i(\theta)=\tau_i^{-1}g(\tau_i\theta),
 \qquad \ell_i(x,\theta)=\ell(x,\tau_i\theta),          \label{eq:red-blowdown-definition-preview}
\end{equation}
then, for every $A\geq1$,
\begin{equation}
              \ell_i(q_i,\theta)\leq\frac n2A^2
              \quad(A^{-1}\leq\theta\leq A).           \label{eq:red-center-time-control}
\end{equation}
\end{proposition}

\begin{proof}
The quadratic lower bound makes $\ell(\cdot,\tau)$ proper, so it attains its
minimum $m(\tau)$.  This minimum envelope has the regularity needed below.
Indeed, on every $[a,b]\Subset(0,\infty)$ the stationary competitor gives a
uniform upper bound for $m$, while the quadratic lower bound confines every
minimizing point to one compact ball.  Local spacetime Lipschitzness of
$\ell$ on that ball then makes $m$ Lipschitz on $[a,b]$.

Let $q_\tau$ minimize at time $\tau$ and choose the upper support supplied by
Theorem~\ref{thm:red-weak-reduced-geometry}.  Because the support lies above
$\ell\geq m(\tau)$ and touches at $q_\tau$, it too has a spatial minimum
there; hence its gradient vanishes and its Laplacian is nonnegative.
Comparing $m(\tau+h)$ with the value at the fixed point $q_\tau$ and using
\eqref{eq:red-ell-sum-inequality} for the support gives the upper-Dini
inequality
\begin{equation}
                         D^+m(\tau)+\frac{m(\tau)}\tau
                         \leq\frac n{2\tau}.             \label{eq:red-center-Dini}
\end{equation}
Here the initial datum is explicit.  Nonnegative scalar curvature gives
$m(\tau)\geq0$, and the stationary path $\gamma(s)=p$, together with the
global scalar bound in \eqref{eq:red-kappa-global-bound}, gives
\begin{equation}
             m(\tau)\leq\ell(p,\tau)
             \leq\frac{1}{2\sqrt\tau}
                    \int_0^\tau C\sqrt s\,ds
             =\frac C3\tau.                            \label{eq:red-center-initial-data}
\end{equation}
Thus $F(\tau)=\tau m(\tau)$ extends continuously by $F(0)=0$, and
\eqref{eq:red-center-Dini} says $D^+F\leq n/2$.  Apply
Lemma~\ref{lem:red-Dini-integration} to obtain
$\tau m(\tau)\leq n\tau/2$, hence $m(\tau)\leq n/2$.
Choose a minimizer.  The scaling law
\eqref{eq:red-scaling-action} identifies $\ell_i$ as in
\eqref{eq:red-blowdown-definition-preview}, and
\eqref{eq:red-ell-time-ratio} gives
\eqref{eq:red-center-time-control}.  Notice that the defining bound
$n/2$ holds at $\theta=1$; it is not asserted at every $\theta$.
\end{proof}

The proof that mass does not escape requires a local distance-distortion
estimate.  We state it in the one-sided form valid at cut points.

\begin{lemma}[Perelman's endpoint changing-distance estimate]
\label{lem:red-changing-distance}
Let $g(\tau)$ be a complete backward Ricci flow near $\tau_0$, and let
$x,y\in M$.  Suppose at time $\tau_0$ that
\begin{equation}
 \Ric\leq(n-1)K g
 \quad\text{on }B(x,r_0)\cup B(y,r_0).                  \label{eq:red-distance-Ricci-hypothesis}
\end{equation}
If $d(x,y)\geq2r_0$, then the upper right Dini derivative of distance due to
the changing metric satisfies
\begin{equation}
 D^+_\tau d_{g(\tau)}(x,y)\big|_{\tau_0}
 \leq2(n-1)\left(\frac23Kr_0+\frac1{r_0}\right).        \label{eq:red-distance-long}
\end{equation}
If $d(x,y)<2r_0$, then
\begin{equation}
 D^+_\tau d_{g(\tau)}(x,y)\big|_{\tau_0}
                         \leq2(n-1)Kr_0.                 \label{eq:red-distance-short}
\end{equation}
For absolutely continuous moving endpoints $x(\tau),y(\tau)$, add
$|x'|+|y'|$ to the right side.
\end{lemma}

\begin{proof}
Choose a unit minimizing geodesic at $\tau_0$.  If its length is at least
$2r_0$, use $n-1$ normal comparison fields which increase linearly from zero
to one on the first endpoint segment of length $r_0$, remain one in the
middle, and decrease linearly on the last segment.  The second variation of
length is nonnegative.  Summing the index forms and using
\eqref{eq:red-distance-Ricci-hypothesis} on the endpoint segments gives
\begin{equation}
 \int_\gamma\Ric(\dot\gamma,\dot\gamma)ds
 \leq2(n-1)\left(\frac23Kr_0+\frac1{r_0}\right).        \label{eq:red-distance-index-bound}
\end{equation}
Because $\partial_\tau g=2\Ric$, the right derivative of the length of this
fixed curve is the integral on the left.  It is an upper barrier for the
later distance, proving \eqref{eq:red-distance-long}.  If the length is less
than $2r_0$, the two endpoint balls cover the geodesic and direct integration
gives \eqref{eq:red-distance-short}.  At a cut pair the chosen minimizing
geodesic is still an upper barrier, so the conclusion is a Dini inequality,
not a claim of classical differentiability.  Moving endpoints contribute
their metric speeds by the triangle inequality.
\end{proof}

\begin{lemma}[Two-point coercivity for reduced distance]
\label{lem:red-two-point}
On an ancient $n$-dimensional $\kappa$-solution with $n\geq2$, for every
$\bar\tau>0$ and
$q_1,q_2\in M$,
\begin{equation}
 \ell(q_1,\bar\tau)+\ell(q_2,\bar\tau)
 \geq
 \frac{d_{g(\bar\tau)}(q_1,q_2)^2}
      {576(n+1)^2\bar\tau}-1.                           \label{eq:red-two-point}
\end{equation}
The constant is explicit but not optimized.
\end{lemma}

\begin{proof}
Choose arbitrary minimizing $\mathcal L$-geodesics $\gamma_a$ from $p$ to
$q_a$, and put
\begin{equation}
                 \ell_a=\ell(q_a,\bar\tau),
                 \qquad a=1,2.                          \label{eq:red-two-point-ell-a}
\end{equation}
No endpoint differentiability or uniqueness is needed: the estimate
\eqref{eq:red-ancient-endpoint-estimate} was stated for every chosen
minimizer.  A prefix of a minimizing curve minimizes the same-clock prefix
action.
Since $R\geq0$,
\begin{equation}
 \ell(\gamma_a(s),s)
 \leq\left(\frac{\bar\tau}{s}\right)^{1/2}\ell_a.      \label{eq:red-prefix-ell}
\end{equation}
The endpoint form \eqref{eq:red-ancient-endpoint-estimate} then gives
\begin{equation}
 |\dot\gamma_a(s)|
 \leq\sqrt3\,\bar\tau^{1/4}s^{-3/4}\sqrt{\ell_a}.      \label{eq:red-two-point-endpoint-speed}
\end{equation}

Interchange the endpoints if necessary so that $\ell_1\leq\ell_2$, and set
\begin{equation}
 B(s)=1+\left(\frac{\bar\tau}{s}\right)^{1/4}\sqrt{\ell_2},
 \qquad r(s)=\frac{\sqrt s}{B(s)}.                      \label{eq:red-two-point-radius}
\end{equation}
Integrating \eqref{eq:red-sqrt-ell-gradient} along a minimizing spatial
geodesic shows that, on either endpoint ball
$B_{g(s)}(\gamma_a(s),r(s))$,
\begin{equation}
                         \sqrt{\ell(\cdot,s)}\leq B(s).
                                                               \label{eq:red-two-point-local-ell}
\end{equation}
Therefore, by \eqref{eq:red-ancient-gradient-estimate} and
\eqref{eq:red-nonnegative-curvature-bridges},
\begin{equation}
                  \Ric\leq Rg\leq\frac{3B(s)^2}{s}g    \label{eq:red-two-point-local-Ricci}
\end{equation}
on those balls.

Put $D(s)=d_{g(s)}(\gamma_1(s),\gamma_2(s))$.  On every interval
$[\varepsilon,\bar\tau]$ the endpoint curves are smooth and the metric is
smooth on their compact images, so $D$ is locally absolutely continuous.
The square-root regularity at the pole and smooth convergence $g(s)\to g(0)$
on compact sets give a continuous extension $D(0)=0$.  Apply
Lemma~\ref{lem:red-changing-distance} with
\begin{equation}
                       K(s)=\frac{3B(s)^2}{(n-1)s},
                       \qquad r_0=r(s).                 \label{eq:red-two-point-K}
\end{equation}
In the long-distance case its metric contribution is at most
$2(n+1)B(s)/\sqrt s$; the short-distance bound is no larger.  Adding
\eqref{eq:red-two-point-endpoint-speed} gives the upper-Dini inequality
\begin{align}
 D^+D(s)\leq{}&2(n+1)s^{-1/2}                            \label{eq:red-two-point-D-prime}\\
 &+\left[2(n+1)\sqrt{\ell_2}
       +\sqrt3(\sqrt{\ell_1}+\sqrt{\ell_2})\right]
       \bar\tau^{1/4}s^{-3/4}.\nonumber
\end{align}
The right side is integrable at zero.  Apply
Lemma~\ref{lem:red-Dini-integration} on $[0,\bar\tau]$ and coarsen the
coefficients to obtain
\begin{equation}
 \frac{D(\bar\tau)}{\sqrt{\bar\tau}}
 \leq4(n+1)+12(n+1)(\sqrt{\ell_1}+\sqrt{\ell_2}).       \label{eq:red-two-point-D-integrated}
\end{equation}
Let
\begin{equation}
 x=\frac{D(\bar\tau)}{12(n+1)\sqrt{\bar\tau}},
 \qquad S=\sqrt{\ell_1}+\sqrt{\ell_2}.                 \label{eq:red-two-point-x-S}
\end{equation}
Then $S\geq x-1/3$ and
$\ell_1+\ell_2\geq S^2/2$.  If $x<1/3$, the desired right side is negative.
If $x\geq1/3$, the elementary identity
\begin{equation}
 \frac12(x-1/3)^2
       -\left(\frac{x^2}{4}-\frac1{18}\right)
       =\left(\frac x2-\frac13\right)^2                 \label{eq:red-two-point-elementary}
\end{equation}
gives
$\ell_1+\ell_2\geq x^2/4-1/18\geq x^2/4-1$, which is
\eqref{eq:red-two-point}.
\end{proof}

\begin{remark}[Why changing distances enter]
\label{rem:red-changing-distance-role}
Global metric equivalence over the long interval $[0,\bar\tau]$ would
produce constants depending exponentially on $\bar\tau$ and cannot yield a
scale-invariant two-point estimate.  The endpoint changing-distance lemma
uses curvature only on balls whose radii are chosen from $\ell$ itself.  It
is precisely what turns local Harnack control into the Gaussian coercivity
needed for the asymptotic limit.
\end{remark}

\section{Blow-down compactness and no loss of mass}
\label{sec:red-asymptotic}

Let $(M^n,\bar g(t))$, $t\leq0$, be an ancient $\kappa$-solution with
$n\geq2$, fix
$p\in M$, and let $\ell$ and $\widetilde V$ be based at $(p,0)$.  The
monotone limit
\begin{equation}
                   \widetilde V_\infty
       :=\lim_{\tau\to\infty}\widetilde V(\tau)          \label{eq:red-asymptotic-volume}
\end{equation}
exists in $[0,1]$.

Choose any sequence $\tau_i\to\infty$ and $\ell$-centers $q_i$ satisfying
\eqref{eq:red-ell-center}.  Define the parabolic blow-downs
\begin{equation}
 g_i(\theta)=\tau_i^{-1}g(\tau_i\theta),
 \qquad
 \ell_i(x,\theta)=\ell(x,\tau_i\theta),
 \qquad 0<\theta<\infty.                                \label{eq:red-blowdown-definition}
\end{equation}
Then $\partial_\theta g_i=2\Ric_{g_i}$, and parabolic covariance gives
\begin{equation}
 \widetilde V_i(\theta)=\widetilde V(\tau_i\theta).      \label{eq:red-blowdown-volume}
\end{equation}

\begin{lemma}[Local compactness about an $\ell$-center]
\label{lem:red-blowdown-local-compactness}
For $A\geq1$ and $D>0$, define the fixed-base-time cylinders
\begin{align}
 Q_i(A,D)&:=B_{g_i(1)}(q_i,D)\times[A^{-1},A],
                                                               \label{eq:red-inner-cylinder}\\
 Q_i^+(A,D)&:=B_{g_i(1)}(q_i,D+2)
                  \times[(2A)^{-1},2A].                \label{eq:red-outer-cylinder}
\end{align}
There are constants $C_m=C_m(n,A,D)$, independent of $i$, such that on
$Q_i(A,D)$,
\begin{equation}
 |\nabla^m\Rm_{g_i}|\leq C_m,
 \qquad
 |\ell_i|+|\nabla\ell_i|+|\partial_\theta\ell_i|\leq C_0
                                                               \label{eq:red-blowdown-local-bounds}
\end{equation}
for every $m\geq0$; the two first-derivative bounds for $\ell_i$ are
almost-everywhere bounds.  The functions $\ell_i$ also have a common spatial
semiconcavity constant on $Q_i(A,D)$.  The injectivity
radii at $(q_i,1)$ have a positive lower bound depending only on
$n,\kappa$.

After passage to a subsequence and a diagonal argument, there are a complete
pointed backward Ricci flow
\begin{equation}
 (M_\infty,g_\infty(\theta),q_\infty),
 \qquad 0<\theta<\infty,                                \label{eq:red-limit-flow}
\end{equation}
and a locally Lipschitz function $\ell_\infty$ such that the pointed flows
converge smoothly on compact spacetime sets and the pulled-back $\ell_i$
converge locally uniformly to $\ell_\infty$.  On compact sets the functions
are uniformly spatially semiconcave.  The limit is complete, has
nonnegative curvature operator, and is $\kappa$-noncollapsed on all scales.
The convergence embeddings may be chosen in the ball-capturing form: for
every finite $A,D$ there is a compact $K_{A,D}\Subset M_\infty$ whose
spacetime image contains $Q_i(A,D)$ for all sufficiently large $i$.
\end{lemma}

\begin{proof}
At the base time, \eqref{eq:red-sqrt-ell-gradient} and
$\ell_i(q_i,1)\leq n/2$ give the explicit estimate
\begin{equation}
 \sqrt{\ell_i(x,1)}
 \leq\sqrt{n/2}+\frac{\sqrt3}{2}
              d_{g_i(1)}(x,q_i).                       \label{eq:red-base-ball-ell}
\end{equation}
For each fixed $x$, the time-ratio estimate
\eqref{eq:red-ell-time-ratio} then bounds $\ell_i(x,\theta)$ on the entire
outer cylinder $Q_i^+(A,D)$ by a constant depending only on $n,A,D$.
Consequently
\begin{equation}
                         R_{g_i}\leq\frac{3\ell_i}\theta
                                                               \label{eq:red-blowdown-R-bound}
\end{equation}
and \eqref{eq:red-nonnegative-curvature-bridges} give a uniform curvature
bound on $Q_i^+(A,D)$.  Integrating the resulting Ricci bound gives uniform
equivalence of $g_i(\theta)$ and $g_i(1)$ there.  Hence the difference of two
units between the spatial radii in \eqref{eq:red-inner-cylinder} and
\eqref{eq:red-outer-cylinder} contains a positive metric collar, uniformly
in $i$ and $\theta$; the time intervals likewise have a positive buffer at
both ends.

Reverse the backward clock on $[(2A)^{-1},2A]$ and apply the local
Bando--Shi theorem of Chapter~8 on $Q_i^+(A,D)$.  Its interior conclusion on
$Q_i(A,D)$ gives every curvature-derivative bound in
\eqref{eq:red-blowdown-local-bounds}; if the theorem is stated using one
intermediate cylinder, insert the radii $D+1$ and the corresponding midpoint
time interval.  Equations \eqref{eq:red-ancient-gradient-estimate} and
\eqref{eq:red-ancient-time-estimate} give the almost-everywhere bounds for
$\ell_i$.  The comparison-field construction in
Theorem~\ref{thm:red-weak-reduced-geometry}, applied with the curvature
derivative bounds on the outer cylinder, gives common spatial
semiconcavity on $Q_i(A,D)$.

For injectivity, let $C_*(n)$ be the curvature bound obtained from
\eqref{eq:red-base-ball-ell} on
$B_{g_i(1)}(q_i,2)$ and set
\begin{equation}
                 \rho:=\min\{1,C_*(n)^{-1/2}\}.        \label{eq:red-injectivity-scale}
\end{equation}
Then $|\Rm_{g_i(1)}|\leq\rho^{-2}$ on
$B_{g_i(1)}(q_i,\rho)$.
$\kappa$-noncollapsing gives volume at least $\kappa\rho^n$, and the
Cheeger--Gromov--Taylor estimate gives
$\inj_{g_i(1)}(q_i)\geq\iota(n,\kappa)\rho$.
Apply Hamilton's
pointed compactness theorem in its ball-capturing formulation and diagonalize
over integer $A,D$.  Completeness of the approximating slices and the
uniform metric and injectivity-radius bounds are the hypotheses which give
the compact sets $K_{A,D}$ in the last assertion.  The
first-derivative bounds for $\ell_i$ give local equicontinuity, while the
support construction in Theorem~\ref{thm:red-weak-reduced-geometry} gives
uniform semiconcavity.  Arzel\`a--Ascoli gives $\ell_\infty$.  Completeness,
curvature nonnegativity, and $\kappa$-noncollapsing pass to a smooth pointed
limit.
\end{proof}

Hamilton compactness is an imported project interface here.  Before the
formalization blueprint is frozen, its already developed Lean theorem must
be matched to the exact pointed, time-interval, basepoint, completeness,
curvature-derivative, and injectivity-radius data used in the preceding
proof; no declaration name or status is inferred merely from this prose.

\begin{lemma}[Uniform Gaussian tails about the moving centers]
\label{lem:red-blowdown-Gaussian-tail}
Let $A\geq1$.  With
\begin{equation}
                       c_n=\frac1{576(n+1)^2},           \label{eq:red-cn}
\end{equation}
for $A^{-1}\leq\theta\leq A$ one has
\begin{equation}
 \ell_i(x,\theta)
 \geq\frac{c_n}{\theta}d_{g_i(\theta)}(x,q_i)^2
       -1-\frac n2A^2.                                  \label{eq:red-blowdown-ell-tail}
\end{equation}
Consequently there is a function $\Psi_{n,A}(D)\downarrow0$ as
$D\to\infty$ such that
\begin{equation}
 \sup_i\sup_{A^{-1}\leq\theta\leq A}
 \int_{M\setminus B_{g_i(\theta)}(q_i,D)}
       (4\pi\theta)^{-n/2}e^{-\ell_i}\,d\mu_{g_i(\theta)}
 \leq\Psi_{n,A}(D).                                    \label{eq:red-blowdown-uniform-tail}
\end{equation}
The same coercive lower bound holds for $\ell_\infty$, centered at
$q_\infty$.
\end{lemma}

\begin{proof}
Apply Lemma~\ref{lem:red-two-point} at rescaled time $\theta$ to $x$ and
$q_i$, and use \eqref{eq:red-center-time-control}.  This gives
\eqref{eq:red-blowdown-ell-tail}.  Since $\Ric_{g_i(\theta)}\geq0$,
Bishop--Gromov comparison gives
\begin{equation}
        \Vol B_{g_i(\theta)}(q_i,k+1)\leq\omega_n(k+1)^n.
                                                               \label{eq:red-Bishop-Euclidean-upper}
\end{equation}
Split the complement of the $D$-ball into unit annuli.  The density there is
bounded by $C(n,A)e^{-(c_n/A)k^2}$, so the tail is bounded by
\begin{equation}
 C(n,A)\sum_{k\geq\lfloor D\rfloor}(k+1)^n
                   e^{-(c_n/A)k^2},                    \label{eq:red-Gaussian-annular-sum}
\end{equation}
which tends to zero independently of $i$ and $\theta$.  Passing the pointwise
lower bound to the smooth pointed limit gives the last assertion.  This
argument is a quantitative tightness proof, not an invocation of dominated
convergence on varying manifolds.
\end{proof}

\begin{theorem}[No loss of reduced-volume mass]
\label{thm:red-no-mass-loss}
For every $\theta>0$, set
\begin{equation}
 u_\infty(x,\theta)
       =(4\pi\theta)^{-n/2}e^{-\ell_\infty(x,\theta)}.   \label{eq:red-limit-density}
\end{equation}
Then
\begin{equation}
 \int_{M_\infty}u_\infty(\cdot,\theta)d\mu_{g_\infty(\theta)}
 =\lim_{i\to\infty}\widetilde V(\tau_i\theta)
 =\widetilde V_\infty                                  \label{eq:red-no-mass-loss}
\end{equation}
for every $\theta>0$.  In particular, the limiting mass is independent of
$\theta$.
\end{theorem}

\begin{proof}
Fix $\theta$ and a smooth cutoff equal to one on the pointed $D$-ball and
supported in the $2D$-ball.  Smooth metric convergence and local uniform
convergence of $\ell_i$ give convergence of the cutoff integrals.  The
ball-capturing conclusion of
Lemma~\ref{lem:red-blowdown-local-compactness} ensures that every inner
$g_i(\theta)$-ball is represented in one limiting compact chart; this is
the additional hypothesis required in
Proposition~\ref{prop:red-action-convergence}.  The
uniform tail \eqref{eq:red-blowdown-uniform-tail} allows $D\to\infty$ on both
the approximating and limiting manifolds, proving the first equality.  The
second follows from \eqref{eq:red-blowdown-volume} and the definition
\eqref{eq:red-asymptotic-volume}.  This is exactly the tightness step which
cannot be replaced by pointed convergence alone.
\end{proof}

\section{The asymptotic-shrinker theorem}
\label{sec:red-asymptotic-shrinker}

\begin{theorem}[Existence of an asymptotic shrinker]
\label{thm:red-asymptotic-shrinker}
Let $(M^n,\bar g(t))$, $-\infty<t\leq0$, be an ancient $\kappa$-solution
with $n\geq2$,
and fix $p\in M$.  For every sequence $\tau_i\to\infty$, choose
$q_i$ with $\ell(q_i,\tau_i)\leq n/2$ and form
\eqref{eq:red-blowdown-definition}.  After passing to a subsequence, the
following hold.
\begin{enumerate}
\item The pointed flows $(M,g_i(\theta),q_i)$ converge smoothly on compact
      subsets of $M_\infty\times(0,\infty)$ to the complete pointed flow
      \eqref{eq:red-limit-flow}, and $\ell_i\to\ell_\infty$ locally
      uniformly.
\item The function $\ell_\infty$ is smooth and
\begin{equation}
 \Ric_{g_\infty(\theta)}+\nabla^2\ell_\infty(\cdot,\theta)
              =\frac1{2\theta}g_\infty(\theta).         \label{eq:red-limit-shrinker-equation}
\end{equation}
At every $\theta>0$,
\begin{equation}
 R+|\nabla\ell_\infty|^2=\frac{\ell_\infty}\theta.      \label{eq:red-limit-Hamilton-normalization}
\end{equation}
\item The normalized shrinker density has the constant mass
\begin{equation}
 0<\int_{M_\infty}(4\pi\theta)^{-n/2}e^{-\ell_\infty}
                   d\mu_{g_\infty(\theta)}
       =\widetilde V_\infty<1.                          \label{eq:red-limit-mass-strict}
\end{equation}
\item The limiting shrinker is nonflat, has nonnegative curvature operator,
      is $\kappa$-noncollapsed on all scales, and satisfies the ancient trace
      Harnack inequality.
\end{enumerate}
No uniqueness of the asymptotic shrinker, and no independence of the
sequence $\tau_i$, is asserted.
\end{theorem}

\begin{proof}
The first assertion and the inherited geometric properties, except
nonflatness, are Lemma~\ref{lem:red-blowdown-local-compactness}.  The weak
inequality \eqref{eq:red-density-weak-test} passes through smooth local
convergence, so
\begin{equation}
 P_\infty^*u_\infty
 :=(\partial_\theta-\Delta_{g_\infty(\theta)}
                 +R_{g_\infty(\theta)})u_\infty\leq0.   \label{eq:red-limit-weak-subsolution}
\end{equation}
Let
\begin{equation}
                         \nu=-P_\infty^*u_\infty.        \label{eq:red-limit-defect}
\end{equation}
This is a nonnegative Radon measure.

Fix $0<a<b<\infty$ and a nonnegative
$\eta\in C_c^\infty((a,b))$.  Since $\Ric\geq0$ and
$\partial_\theta g_\infty=2\Ric$, a cutoff $\chi_D$ based on
$g_\infty(a)$ can be chosen with
\begin{align}
 0\leq\chi_D\leq1,\qquad
 &\chi_D=1\text{ on }B_{g_\infty(a)}(q_\infty,D),      \label{eq:red-limit-cutoff}\\
 &\supp\chi_D\subset B_{g_\infty(a)}(q_\infty,2D),
 \qquad |\nabla\chi_D|_{g_\infty(\theta)}\leq C/D.\nonumber
\end{align}
Indeed $g_\infty(\theta)\geq g_\infty(a)$ on $[a,b]$, so a standard
$g_\infty(a)$-cutoff has the displayed gradient bound for every such
$\theta$.  Its transition annulus is outside the radius-$D$ ball for every
$g_\infty(\theta)$ as well.  In particular, $\eta\chi_D$ is a legitimate
compactly supported test for the Radon measure $\nu$.
By the definition of the weak operator and spatial integration by parts,
\begin{align}
 \nu(\eta\chi_D)
={}&\int_a^b\!\!\int_{M_\infty}
  \left(u_\infty\eta'\chi_D
       -\eta\langle\nabla u_\infty,\nabla\chi_D\rangle\right)
                  d\mu d\theta.                        \label{eq:red-limit-defect-cutoff}
\end{align}
The first term tends, by Theorem~\ref{thm:red-no-mass-loss}, to
\begin{equation}
                      \widetilde V_\infty\int_a^b\eta'=0.
                                                               \label{eq:red-limit-defect-time-zero}
\end{equation}
The gradient estimate passes to the limit and gives
\begin{equation}
 |\nabla u_\infty|
 =u_\infty|\nabla\ell_\infty|
 \leq C(a,b)\sqrt{\ell_\infty}\,e^{-\ell_\infty}.      \label{eq:red-limit-gradient-density}
\end{equation}
The coercive bound in Lemma~\ref{lem:red-blowdown-Gaussian-tail} makes this
uniformly integrable, so the second term in
\eqref{eq:red-limit-defect-cutoff} is $O(D^{-1})$.  Hence
$\nu(\eta\chi_D)\to0$.  Positivity of $\nu$, followed by exhaustion in
space and time, proves
\begin{equation}
                         P_\infty^*u_\infty=0.           \label{eq:red-limit-conjugate-equality}
\end{equation}
Interior parabolic regularity makes $u_\infty$ smooth and positive, and
therefore makes $\ell_\infty$ smooth.

We next pass the Hamilton--Jacobi identity.  On each compact spacetime set,
the functions $\ell_i$ are uniformly Lipschitz and uniformly semiconcave in
space.  Local uniform convergence of semiconcave functions implies
\begin{equation}
          \nabla\ell_i\longrightarrow\nabla\ell_\infty
          \quad\text{almost everywhere and in }L^p_{\rm loc}
          \quad(1\leq p<\infty).                        \label{eq:red-limit-gradient-convergence}
\end{equation}
To see this without a black-box differentiability claim, subtract the common
coordinate quadratic, mollify the concave functions, use pointwise
convergence of gradients at differentiability points of the limit, and apply
the uniform local Lipschitz bound for dominated convergence.  Insert this
convergence in the weak time-derivative form of
\eqref{eq:red-Hamilton-Jacobi}.  It identifies the time derivative of the
limit and gives
\begin{equation}
 2\partial_\theta\ell_\infty+|\nabla\ell_\infty|^2-R
                     +\frac{\ell_\infty}\theta=0.       \label{eq:red-limit-Hamilton-Jacobi}
\end{equation}
The conjugate-heat equality gives
\begin{equation}
 \partial_\theta\ell_\infty-\Delta\ell_\infty
       +|\nabla\ell_\infty|^2-R+\frac n{2\theta}=0.     \label{eq:red-limit-ell-conjugate}
\end{equation}
Combining these equations yields
\begin{equation}
 \theta(2\Delta\ell_\infty-|\nabla\ell_\infty|^2+R)
       +\ell_\infty-n=0.                                \label{eq:red-limit-v-zero}
\end{equation}
Thus the $v$-quantity of Lemma~\ref{lem:red-v-identity} is identically zero.
Its local identity, which does not require a pole at $\theta=0$, gives
\eqref{eq:red-limit-shrinker-equation}.  Taking the trace there and combining
it with \eqref{eq:red-limit-v-zero} gives
\eqref{eq:red-limit-Hamilton-normalization}.

Theorem~\ref{thm:red-no-mass-loss} gives the equality in
\eqref{eq:red-limit-mass-strict}.  Its value is positive because
$\ell_\infty(q_\infty,1)\leq n/2$ and $u_\infty$ is positive on a small ball.
If it were one, monotonicity and the bound $\widetilde V\leq1$ would make the
original regular-pole reduced volume identically one.  Theorem
\ref{thm:red-regular-pole-rigidity} would make the original ancient solution
Gaussian, contrary to the definition of a $\kappa$-solution.  Hence the
value is strictly less than one.

Finally suppose the limiting shrinker were flat.  At $\theta=1$,
\eqref{eq:red-limit-shrinker-equation} gives
$\nabla^2\ell_\infty=g_\infty/2$.  The Gaussian lower bound makes
$\ell_\infty$ proper, so strict convexity gives a unique minimum.  Here is
the quotient argument explicitly.  The universal Riemannian cover of the
complete flat manifold is $(\mathbb R^n,g_{\rm E})$.  The lifted potential
$\widetilde\ell$ satisfies $\nabla^2\widetilde\ell=g_{\rm E}/2$, hence
\begin{equation}
                 \widetilde\ell(x)=\frac14|x-a|^2+C    \label{eq:red-flat-lift-potential}
\end{equation}
for some $a\in\mathbb R^n$ and $C\in\mathbb R$.  Every deck transformation
preserves $\widetilde\ell$ and therefore fixes its unique minimizer $a$.
The deck action of a covering is free, so the deck group is trivial and the
flat manifold itself is $\mathbb R^n$.  At $\theta=1$,
\eqref{eq:red-limit-Hamilton-normalization} reads
$|\nabla\widetilde\ell|^2=\widetilde\ell$ and forces $C=0$ in
\eqref{eq:red-flat-lift-potential}.  The normalized mass is therefore one,
contradicting the strict inequality.  Thus the limit is nonflat.  The trace
Harnack estimate passes by smooth convergence, completing the proof.
\end{proof}

\begin{remark}[The pole escapes]
\label{rem:red-escaping-pole}
The function $\ell_\infty$ need not be the reduced distance of the limiting
flow from a regular point at $\theta=0$.  Usually the old point $p$ escapes
every compact pointed set, and the limit flow is not even defined at
$\theta=0$.  Properness is supplied by
\eqref{eq:red-blowdown-ell-tail}, centered at $q_\infty$; the shrinker
equation is supplied by the local $v$-identity.  Neither step refers to a
nonexistent limiting pole.
\end{remark}

\section{Three-dimensional asymptotic models}
\label{sec:red-three-models}

\begin{corollary}[Asymptotic shrinkers of three-dimensional
$\kappa$-solutions]
\label{cor:red-three-asymptotic-shrinker}
Let $(M^3,\bar g(t))$ be an ancient $\kappa$-solution.  Every asymptotic
shrinker furnished by Theorem~\ref{thm:red-asymptotic-shrinker} is exactly
one of the following normalized solitons:
\begin{enumerate}
\item a round spherical space form $\Sph^3(2)/\Gamma$;
\item $\Sph^2(\sqrt2)\times\mathbb R$;
\item $\mathbb {RP}^2(\sqrt2)\times\mathbb R$;
\item the diagonal free quotient of $\Sph^2(\sqrt2)\times\mathbb R$ by
      $(y,s)\mapsto(-y,-s)$.
\end{enumerate}
If $M$ is noncompact, the first alternative cannot occur.  If $M$ is
orientable, the third alternative cannot occur.
\end{corollary}

\begin{proof}
At $\theta=1$, equations \eqref{eq:red-limit-shrinker-equation} and
\eqref{eq:red-limit-Hamilton-normalization} make
$(M_\infty,g_\infty(1),\ell_\infty(\cdot,1))$ a complete Hamilton-normalized
gradient shrinker.  It has nonnegative curvature operator and is nonflat.
Theorem~\ref{thm:nonnegative-three-shrinker-classification} and the exact
potential-preserving quotient list in
Chapter~\ref{ch:three-dimensional-shrinkers} give the four alternatives;
the Gaussian branch is excluded by nonflatness.

If a pointed smooth limit of the connected noncompact manifolds $M$ were
compact, an exhausting convergence domain would equal all of $M_\infty$ for
large $i$.  Its same-dimensional embedding into $M$ would have open image by
invariance of domain and closed image by compactness, hence would be
surjective.  This would make $M$ compact, a contradiction.  Orientations
pull back through the convergence embeddings and exhaust the limit, while
$\mathbb {RP}^2\times\mathbb R$ is nonorientable.
\end{proof}

The normalized masses provide useful regression tests:
\begin{equation}
 \begin{array}{c|c}
 \text{model}&(4\pi)^{-3/2}\int e^{-f}d\mu\\ \hline
 \mathbb R^3&1\\
 \Sph^2(\sqrt2)\times\mathbb R&2/e\\
 \text{either order-two cylinder quotient}&1/e\\
 \Sph^3(2)/\Gamma&2\sqrt\pi\,e^{-3/2}/|\Gamma|.
 \end{array}                                             \label{eq:red-three-density-regression}
\end{equation}
For example, the cylinder potential is $f=1+s^2/4$; multiplying the area
$8\pi$ by $\int_\mathbb R e^{-s^2/4}ds=2\sqrt\pi$ and by $e^{-1}$ gives the
second line.  These values distinguish regular-pole Gaussian rigidity from
the non-Gaussian constant masses of blow-down shrinkers.

\section{Interfaces reserved for Ricci flow with surgery}
\label{sec:red-surgery-interface}

The smooth chapter must expose exactly the lemmas which survive when the
ambient spacetime later has discrete surgery times.  It must not silently
define a path across a surgery before the survivor and cap identifications
have been fixed.

\begin{lemma}[A reusable crossing-cost bound]
\label{lem:red-crossing-cost}
Let $0<a<b$.  Suppose on a spacetime region that
\begin{equation}
                        R\geq Q\geq0,
                        \qquad g(\sigma)\geq c h         \label{eq:red-crossing-hypotheses}
\end{equation}
for a fixed Riemannian metric $h$ and $c>0$.  If a path segment in that
region joins two sets at $h$-distance at least $d$, then
\begin{equation}
                 \mathcal A^{t_0,g}_{a,b}(\gamma)
                 \geq2\sqrt{acQ}\,d.                   \label{eq:red-crossing-cost}
\end{equation}
\end{lemma}

\begin{proof}
For every velocity $v$,
\begin{equation}
 \sqrt\sigma(R+|v|_g^2)
 \geq\sqrt a\,(Q+c|v|_h^2)
 \geq2\sqrt{acQ}\,|v|_h.                               \label{eq:red-crossing-pointwise}
\end{equation}
Integrate and use that the $h$-length is at least $d$.
\end{proof}

Together with the scalar lower bound
\eqref{eq:red-action-scalar-lower} outside the crossing interval, this lemma
will show that a curve entering a sufficiently curved new cap pays more than
an unscathed comparison curve.  Proposition~\ref{prop:red-domain-calculus}
then converts that strict action gap into confinement.

The public smooth-to-surgery interfaces supplied by this chapter are:
\begin{enumerate}
\item the arbitrary-pole clock and the additive same-clock segment action,
      Definitions~\ref{def:red-clocked-action} and \ref{def:red-distance};
\item exact restriction, enlargement, exhaustion, and leaving-curve gap
      rules, Proposition~\ref{prop:red-domain-calculus};
\item locality under agreement on an open spacetime neighborhood and
      covariance under fixed pullback, time translation, and parabolic
      scaling, Proposition~\ref{prop:red-covariance};
\item action convergence under confinement and the separate tightness
      hypothesis for total masses, Proposition~\ref{prop:red-action-convergence};
\item the weak endpoint inequality
      \eqref{eq:red-density-endpoint-test}, whose temporal boundary terms
      can later record errors at surgery times;
\item endpoint changing-distance control,
      Lemma~\ref{lem:red-changing-distance}; and
\item the cap- or neck-crossing estimate
      \eqref{eq:red-crossing-cost}.
\end{enumerate}

The later surgery chapter must still define: surviving points and
worldlines; pre- and post-surgery comparison maps; unscathed, admissible, and
barely admissible curves; the status of a curve ending in a newly inserted
cap; action at a discrete surgery time; quantitative cap-encounter costs in
terms of the surgery parameters; reduced-volume jump errors; and the
surgery-uniform no-local-collapsing theorem.  None of those notions is
totalized or assumed here.  In particular, the present chapter proves smooth
noncollapsing consequences only through the explicitly imported
$\kappa$-hypothesis; it does not yet prove Perelman's noncollapsing theorem
for a flow with surgery.

\section{Formalization units and dependency boundary}
\label{sec:red-formalization}

The mathematical proof above is intentionally factored into the following
fine nodes.  The identifiers are stable semantic suggestions for the
blueprint, not declarations asserted to exist in Lean.

\begin{enumerate}
\item \nodeid{PC-F-RED-BACKWARD-POLE} is
      \eqref{eq:red-backward-flow}, including the proof that elapsed time is
      positive.  \nodeid{PC-F-RED-CLOCKED-SEGMENT-ACTION} is
      Definition~\ref{def:red-clocked-action}, with values first in the
      extended reals.  \nodeid{PC-F-RED-SCALAR-LOWER} excludes $-\infty$.
      \nodeid{PC-F-RED-VALUE-FINITE} combines that exclusion with a finite
      competitor and extracts the unique real representative used by $L$
      and $\ell$.

\item \nodeid{PC-F-RED-ACTION-CONCATENATION} and
      \nodeid{PC-F-RED-DYNAMIC-PROGRAMMING} are the two assertions of
      Proposition~\ref{prop:red-dynamic-programming}.  Their common clock is
      part of the interface and must not be erased by a simplifier.

\item \nodeid{PC-F-RED-SQRT-TIME-PATH} and
      \nodeid{PC-F-RED-SQRT-TIME-ACTION} are
      \eqref{eq:red-square-root-path} and
      \eqref{eq:red-square-root-action}.
      \nodeid{PC-F-RED-AC-REPARAMETRIZATION} is
      \eqref{eq:red-AC-reparametrization}, and
      \nodeid{PC-F-RED-COORDINATE-MOMENTUM} is
      \eqref{eq:red-coordinate-momentum-equation}.
      \nodeid{PC-F-RED-MINIMIZER-DIRECT-METHOD} is
      Theorem~\ref{thm:red-minimizer-existence}, internally split into
      energy boundedness, confinement, uniform convergence, weak $H^1$
      convergence, lower semicontinuity, and one-dimensional regularity.

\item \nodeid{PC-F-RED-FIRST-VARIATION} is
      \eqref{eq:red-first-variation};
      \nodeid{PC-F-RED-GEODESIC-ODE} is the regular equation
      \eqref{eq:red-L-geodesic-square-root};
      \nodeid{PC-F-RED-TERMINAL-SPEED} is
      Lemma~\ref{lem:red-terminal-speed} and
      Corollary~\ref{cor:red-first-order-growth};
      \nodeid{PC-F-RED-ENDPOINT-SUPPORT} is
      Proposition~\ref{prop:red-endpoint-supports}; and
      \nodeid{PC-F-RED-HAMILTON-JACOBI} is the almost-everywhere identity
      \eqref{eq:red-Hamilton-Jacobi}.

\item \nodeid{PC-F-RED-INDEX-FORM} is
      \eqref{eq:red-index-form};
      \nodeid{PC-F-RED-FULL-PATH-K-SUPPORT} is
      Lemma~\ref{lem:red-full-path-support};
      \nodeid{PC-F-RED-TRACE-COMPARISON} is the finite trace calculation
      \eqref{eq:red-matrix-trace-contraction}--
      \eqref{eq:red-L-Laplacian}, using that full-path support; and
      \nodeid{PC-F-RED-K-ELIMINATION} is
      \eqref{eq:red-ell-time-K}--\eqref{eq:red-ell-Laplacian-K}.

\item \nodeid{PC-F-RED-UPPER-SUPPORT-SEMICONCAVITY} is the first assertion
      of Theorem~\ref{thm:red-weak-reduced-geometry};
      \nodeid{PC-F-RED-K-ELIMINATION-AE} is
      \eqref{eq:red-main-ae};
      \nodeid{PC-F-RED-SEMICONCAVITY-DISTRIBUTION} is
      \eqref{eq:red-semiconcavity-transfer}; and
      \nodeid{PC-F-RED-WEAK-DENSITY} is the exact test inequality
      \eqref{eq:red-density-weak-test}.

\item \nodeid{PC-F-RED-REGULAR-GAUSSIAN-BOUNDS} is
      Lemma~\ref{lem:red-quadratic-bounds}.
      \par\noindent
      \nodeid{PC-F-RED-L1-MOVING-MEASURE} is
      Lemma~\ref{lem:red-L1-moving-measure}.
      \par\noindent
      \nodeid{PC-F-RED-POLE-NORMALIZATION} is
      Theorem~\ref{thm:red-volume-normalization}.
      \par\noindent
      \nodeid{PC-F-RED-VOLUME-MONOTONICITY} is the cutoff proof of
      Theorem~\ref{thm:red-volume-monotonicity}.

\item \nodeid{PC-F-RED-V-LOCAL-IDENTITY} is
      Lemma~\ref{lem:red-v-identity};
      \nodeid{PC-F-RED-DEFECT-MEASURE} converts the positive distribution
      to a Radon measure and exhausts compactly supported tests;
      \nodeid{PC-F-RED-DEFECT-ZERO} is the positive-measure argument
      \eqref{eq:red-defect-cutoff}; and
      \nodeid{PC-F-RED-REGULAR-POLE-RIGIDITY} is
      Theorem~\ref{thm:red-regular-pole-rigidity}.

\item \nodeid{PC-F-RED-KAPPA-CENTER} is
      Proposition~\ref{prop:red-ell-center}.
      \par\noindent
      \nodeid{PC-F-RED-DINI-INTEGRATION} is
      Lemma~\ref{lem:red-Dini-integration}.
      \par\noindent
      \nodeid{PC-F-RED-ANCIENT-HARNACK-ELL} is
      Proposition~\ref{prop:red-ancient-ell-estimates}.
      \par\noindent
      \nodeid{PC-F-RED-DISTANCE-DINI} is
      Lemma~\ref{lem:red-changing-distance}.
      \par\noindent
      \nodeid{PC-F-RED-TWO-POINT-COERCIVITY} is
      Lemma~\ref{lem:red-two-point}.

\item \nodeid{PC-F-RED-BLOWDOWN-LOCAL-COMPACTNESS} is
      Lemma~\ref{lem:red-blowdown-local-compactness};
      \nodeid{PC-F-RED-GAUSSIAN-TIGHTNESS} is
      Lemma~\ref{lem:red-blowdown-Gaussian-tail};
      \nodeid{PC-F-RED-NO-MASS-LOSS} is
      Theorem~\ref{thm:red-no-mass-loss}; and
      \nodeid{PC-F-RED-HJ-LIMIT} is semiconcave-gradient convergence
      \eqref{eq:red-limit-gradient-convergence} and passage to the weak
      Hamilton--Jacobi identity;
      \nodeid{PC-F-RED-ASYMPTOTIC-SHRINKER} is
      Theorem~\ref{thm:red-asymptotic-shrinker}.

\item \nodeid{PC-B-RED-DOMAIN-RESTRICTION} collects
      Proposition~\ref{prop:red-domain-calculus};
      \nodeid{PC-B-RED-POINTED-STABILITY} is
      Proposition~\ref{prop:red-action-convergence}; and
      \nodeid{PC-B-RED-SURGERY-INTERFACE} is the explicit producer--consumer
      list in Section~\ref{sec:red-surgery-interface}.
\end{enumerate}

The coarse dependency forest is
\begin{equation}
\begin{gathered}
 \text{square-root action}\Longrightarrow
 \text{minimizer}\Longrightarrow\text{variations},\\
 \text{variations}\Longrightarrow\text{weak density},\\
 \text{weak density}+\text{Gaussian tails}\Longrightarrow
 \text{reduced-volume monotonicity},\\
 \text{ancient Harnack}+\text{changing distance}
 \Longrightarrow\text{two-point coercivity},\\
 \text{two-point coercivity}
 \Longrightarrow\text{tight blow-down mass},\\
 \text{Hamilton compactness}+\text{center estimates}
 \Longrightarrow\text{pointed limit},\\
 \text{tight limiting mass}+v\text{-identity}
 \Longrightarrow\text{asymptotic shrinker}.
\end{gathered}                                             \label{eq:red-coarse-dependency-forest}
\end{equation}

The imported mathematical inputs are: Hopf--Rinow and Arzel\`a--Ascoli;
weak compactness and lower semicontinuity in one-dimensional $H^1$; the
support characterization and Hessian-measure theorem for semiconcave
functions; Bishop--Gromov comparison; parabolic regularity; the injectivity
radius estimate; completeness of the shrinker vector field and its canonical
flow; forward uniqueness for complete bounded-curvature Ricci flows;
Hamilton compactness; Chapter~8's derivative estimates;
Chapter~9's ancient trace Harnack estimate; and the classification theorem
of Chapter~\ref{ch:three-dimensional-shrinkers}.  Each imported interface
must be matched to an exact statement in the blueprint.  In particular,
although collaborators have separately developed Hamilton compactness and
$\mathcal W$-entropy formalizations, this chapter makes no compilation,
kernel, axiom-footprint, or semantic-correspondence claim about them without
an exact declaration audit.  The $\mathcal W$-entropy is not used in the
proof given here.

\subsection*{Lean status}

Every new node in this chapter currently has Lean statement and Lean proof
status \emph{not yet implemented} unless and until the project status ledger
records contrary declaration-level evidence.  The displayed statements and
proofs are the human mathematical source; they do not themselves constitute
typechecking or kernel-checking evidence.

\section{Sources, corrections, and deliberate omissions}
\label{sec:red-sources}

The primary source is G.~Perelman,
\emph{The entropy formula for the Ricci flow and its geometric applications},
\href{https://arxiv.org/abs/math/0211159}{arXiv:math/0211159}, especially
Sections~7.1--7.2 and Proposition~11.2.  The detailed foundational spine is
Chapter~7 of MSM~135, especially its square-root-time formulas, first and
second variations, endpoint identities, reduced-distance inequalities, and
weak formulations.  The reduced-volume argument is cross-checked against
Chapter~8 of MSM~135 and Sections~23--24 and~29 of Kleiner--Lott's
\emph{Notes on Perelman's papers} \cite{MR2460872}.  MSM~144, Chapter~16,
provides an independent check of the weak conjugate-heat subsolution and the
$v$-identity.  GSM~77, Chapter~11, Section~4, supplies the large-dimensional
spacetime motivation, but is not used as a proof source.

The asymptotic-shrinker proof uses the correction in MSM~163, Chapter~19,
Section~5.  Its critical two-point estimate is also found in R.~Ye,
\emph{On the $l$-function and the reduced volume of Perelman I}, Lemma~2.2;
Kleiner--Lott, equation~(39.6); and J.~Morgan--G.~Tian,
\emph{Ricci Flow and the Poincar\'e Conjecture}, Lemma~9.25.  MSM~206,
Theorem~28.24, gives a later concise statement of the asymptotic-shrinker
theorem.

The following source issues have been repaired explicitly.
\begin{enumerate}
\item The direct-method proof of a minimizing curve at the singular endpoint
      is only sketched in MSM~135.  Theorem~\ref{thm:red-minimizer-existence}
      supplies confinement, $H^1$ compactness, lower semicontinuity, endpoint
      preservation, and regularity in the nonsingular $r=2\sqrt\tau$
      variable.

\item Pointed convergence about $q_i$ does not by itself imply convergence
      of the full reduced-volume integrals.  The original pole $p$ can escape
      every pointed ball.  This is the gap in MSM~135, Lemma~8.38.  Theorem
      \ref{thm:red-no-mass-loss} inserts the missing uniform tightness proof,
      following the Ye--Morgan--Tian--Kleiner--Lott estimate incorporated in
      MSM~163.

\item The properness estimate on the next page of the older proof refers to
      the old point $p$ and a metric $g_\infty(0)$.  Neither need exist in
      the pointed blow-down.  Formula~\eqref{eq:red-blowdown-ell-tail} uses
      the actual limiting center $q_\infty$ and a positive time $\theta$.

\item The selected center satisfies $\ell_i(q_i,1)\leq n/2$, not
      $\ell_i(q_i,\theta)\leq n/2$ for arbitrary $\theta$.  Equation
      \eqref{eq:red-center-time-control} replaces the latter by the correct
      uniform bound $(n/2)A^2$ on $A^{-1}\leq\theta\leq A$.  This repairs the
      parameter slip in MSM~163, equation~(19.66), without weakening the
      Gaussian tail conclusion.

\item Perelman's original equality-to-shrinker passage used a formula whose
      stated hypotheses implicitly require a flow down to $\theta=0$.
      A blow-down has no such regular pole.  Equations
      \eqref{eq:red-limit-conjugate-equality}--
      \eqref{eq:red-limit-shrinker-equation} instead use the local
      $v$-identity, as in the corrected Kleiner--Lott and MSM~135 treatments.
      The factor $u$ on the right side of \eqref{eq:red-v-identity}, omitted
      from one original displayed formula, is retained.

\item In the finite-terminal-time clause of MSM~135, Lemma~7.64, the
      preceding calculation gives $1+2T/(T-\tau)$ and hence the uniform
      constant $1+2/c$ under $\tau<(1-c)T$, not the printed $1+2c$.  The same
      correction propagates to the corresponding finite-time clause of the
      next lemma.  Our ancient estimate
      \eqref{eq:red-ancient-gradient-estimate} has the correct constant $3$
      and does not use the erroneous finite-time display.

\item Reduced volume is nonincreasing in backward time $\tau$.  Statements
      with the opposite monotonicity use forward time.  Kleiner--Lott also
      uses a convention without the factor $(4\pi)^{-n/2}$; this chapter uses
      the unit-Gaussian normalization throughout.
\end{enumerate}

We deliberately omit the full $\mathcal L$-Jacobi-field, cut-time,
cut-locus, and Jacobian change-of-variables package from the dependency path
of the main theorem.  Classically, on the minimizing domain of
$\mathcal L\exp$, the density
\begin{equation}
 (4\pi\tau)^{-n/2}e^{-\ell(\gamma_V(\tau),\tau)}
                J_V^{\mathcal L}(\tau)                  \label{eq:red-classical-Jacobian-density}
\end{equation}
is pointwise nonincreasing and bounded by
$\pi^{-n/2}e^{-|V|^2}$; integrating in $V$ gives another beautiful proof of
monotonicity.  Making that route foundational would additionally require the
cut-time function, measure zero of the $\mathcal L$-cut locus, the
diffeomorphism domain of $\mathcal L\exp$, and the change-of-variables
theorem.  The weak proof above proves the required result with fewer and more
surgery-compatible interfaces.

Forward reduced distance, large-dimensional thermostat constructions,
reduced distance based at a singular Type~I time, uniqueness of asymptotic
shrinkers, and the definitions specific to a surgery spacetime are also
deferred.  They are not used in the proof of
Theorem~\ref{thm:red-asymptotic-shrinker}; the last group is explicitly
reserved for the next surgery-facing stages listed in
Section~\ref{sec:red-surgery-interface}.

The complete human proofs in this chapter are conditional only on the named
classical analytic and compactness inputs recorded in
Section~\ref{sec:red-formalization}.  No assertion is made that the new
chapter has already been formalized or independently certified in Lean.

\chapter[Compactness of Three-Dimensional Ancient Kappa-Solutions]
{Curvature-Normalized Compactness and Uniform Geometry of
Three-Dimensional Ancient $\kappa$-Solutions}
\label{ch:kappa-compactness}
\index{$\kappa$-solution!compactness}
\index{curvature normalization}
\index{asymptotic scalar curvature ratio}
\index{asymptotic volume ratio}
\index{strong neck}

\section{Purpose, scope, and the closure issue}
\label{sec:ksol-purpose}

The purpose of this chapter is to turn the definition of an ancient
$\kappa$-solution into a uniform geometric object which can be rescaled at
an arbitrary spacetime point.  This is the compactness package needed before
one can prove that distant regions are necks, organize the complementary
regions into caps, and classify three-dimensional ancient $\kappa$-solutions.
The chapter stops at the precise strong-neck detection interface.  The
global neck--cap decomposition and neck stability belong to the next stage.

An ancient $\kappa$-solution always means Definition
\ref{def:red-kappa-solution}.  In particular, it has one global spacetime
curvature bound, not merely a bound on each compact set.  This convention
creates a real closure issue.  Pointed smooth convergence initially provides
curvature bounds only on bounded-distance spacetime cylinders.  The limit
plainly retains nonnegative curvature, noncollapsing, and the trace Harnack
inequality, but it need not plainly retain a global bound at spatial
infinity.  We shall prove the missing boundedness upgrade in dimensions two
and three before calling the limit a $\kappa$-solution.

This explains the useful part of the notion historically called a
``$\kappa$-solution with Harnack.''  We do not create a second public meaning
of $\kappa$-solution.  Instead, the following predicate is used only as a
temporary carrier inside limit arguments.

\begin{definition}[Harnack limit package]
\label{def:ksol-Harnack-limit-package}
For $n\geq2$ and $\kappa>0$, a smooth Ricci flow $g(t)$ on $M^n$,
$t\leq0$, satisfies
$\mathsf{KLim}_{n,\kappa}$ when all of the following hold.
\begin{enumerate}
\item $M$ is nonempty and connected, every $(M,g(t))$ is complete, and the
      flow is nonflat.
\item $\Rm_{g(t)}\geq0$ for every $t\leq0$.
\item The same $\kappa$ gives the all-scale, every-time noncollapsing
      condition of Definition~\ref{def:red-kappa-solution}.
\item For every $(x,t)\in M\times(-\infty,0]$ and every
      $X\in T_xM$,
\begin{equation}
 \partial_tR+2\langle\nabla R,X\rangle+2\Ric(X,X)\geq0
 \quad\hbox{at }(x,t).                              \label{eq:ksol-KLim-Harnack}
\end{equation}
      At the terminal time $t=0$, the symbol $\partial_tR$ denotes the
      left derivative.  All spatial tensors in this inequality are those
      of $g(t)$.
\end{enumerate}
No global curvature bound is included in this predicate.
\end{definition}

Chapter~\ref{ch:matrix-harnack} and time translation give
\begin{equation}
 \text{ancient $\kappa$-solution}
       \Longrightarrow\mathsf{KLim}_{n,\kappa}.
                                                               \label{eq:ksol-standard-to-KLim}
\end{equation}
The reverse implication will be proved for $n=2,3$ in
Theorems~\ref{thm:ksol-two-dimensional-KLim-bounded} and
\ref{thm:ksol-three-dimensional-KLim-bounded}.  Thus the predicate is an
internal closure interface, not an extra geometric species.

\begin{roadmap}{Sections~\ref{sec:ksol-normalization}--
\ref{sec:ksol-split-branch} establish curvature normalization, the
two-dimensional theorem, and the complete rank-one quotient branch.
Sections~\ref{sec:ksol-spatial-infinity}--
\ref{sec:ksol-fixed-compactness} prove the spatial-infinity, almost-ancient,
and fixed-$\kappa$ compactness packages.  Section~\ref{sec:ksol-Harnack-closure}
closes pointed limits inside the Chapter~10 definition.  Sections
\ref{sec:ksol-universal-gap}--\ref{sec:ksol-neck-detection} prove the
universal noncollapsing gap, uniform derivative estimates, and the local
strong-neck detection theorem.}
\end{roadmap}

All constants are positive unless explicitly allowed to vanish.  A
``universal'' constant in this chapter may depend on the dimension and on
the fixed tensor-norm convention, but not on a particular solution or on
its input noncollapsing constant.  We work in dimension three except in
lemmas whose proofs are no harder in general dimension.

\section{Scalar positivity, normalization, and pointed convergence}
\label{sec:ksol-normalization}

\begin{lemma}[Scalar positivity for genuine ancient solutions]
\label{lem:ksol-scalar-positive}
Let $g(t)$, $t\leq0$, be an ancient $\kappa$-solution.  Then
\begin{equation}
                         R(x,t)>0
       \quad\text{for every }(x,t)\in M\times(-\infty,0].       \label{eq:ksol-R-positive}
\end{equation}
\end{lemma}

\begin{proof}
Nonnegative curvature operator gives $R\geq0$.  If $R(x_0,t_0)=0$, then all
eigenvalues of the curvature operator vanish at $(x_0,t_0)$.  The scalar
evolution
\begin{equation}
                  \partial_tR=\Delta R+2|\Ric|^2              \label{eq:ksol-scalar-evolution}
\end{equation}
and Harnack at vector zero imply $R(x_0,t)=0$ for every $t\leq t_0$.
Apply the scalar strong maximum principle on every compact time slab ending
at $t_0$: all slices in the past of $t_0$ are flat.  If $t_0<0$, complete
bounded-curvature forward uniqueness propagates this flatness to time zero;
if $t_0=0$, continuity already includes the terminal slice.  Either case
contradicts nonflatness.
\end{proof}

\begin{definition}[Curvature normalization]
\label{def:ksol-curvature-normalization}
Let $(x_0,t_0)$ be a point of an ancient flow covered by
Lemma~\ref{lem:ksol-scalar-positive}, and put $Q=R_g(x_0,t_0)$.  Its
curvature-normalized pointed flow is
\begin{equation}
 g^{(x_0,t_0)}(s)=Q\,g\!\left(t_0+\frac{s}{Q}\right),
 \qquad -\infty<s\leq0.                              \label{eq:ksol-curvature-normalization}
\end{equation}
Then
\begin{equation}
                  R_{g^{(x_0,t_0)}}(x_0,0)=1.         \label{eq:ksol-normalized-base-curvature}
\end{equation}
\end{definition}

Completeness, nonnegative curvature operator, the Harnack inequality, and
the same noncollapsing constant are invariant under
\eqref{eq:ksol-curvature-normalization}.  Distances and volumes are
multiplied by $Q^{1/2}$ and $Q^{n/2}$; scalar curvature and $|\Rm|$ are
multiplied by $Q^{-1}$.  Mixed derivatives have the following exact
scaling law.  If
\[
       \widehat g(s)=Q\,g(t_0+s/Q),\qquad t=t_0+s/Q,
\]
then
\begin{equation}
 \left|\nabla_s^b\nabla_{\widehat g(s)}^a
            \Rm_{\widehat g(s)}\right|_{\widehat g(s)}
 =Q^{-1-a/2-b}
   \left|\nabla_t^b\nabla_{g(t)}^a
            \Rm_{g(t)}\right|_{g(t)}.                \label{eq:ksol-mixed-scaling}
\end{equation}
Here $\nabla_t$ and $\nabla_s$ are the metric-compatible time derivatives
defined in \eqref{eq:shi-metric-time-derivative}.  Formula
\eqref{eq:ksol-mixed-scaling} includes both metric and time rescaling; the
factor $Q^{-b}$ would be false for a mere constant rescaling at fixed time.

We use pointed smooth convergence of flows in the following exact sense.
For every $A,T<\infty$, an exhaustion embedding from a neighborhood of
$\overline B_{g_\infty(0)}(x_\infty,A)$ into $M_i$ is defined for all large
$i$, takes $x_\infty$ to $x_i$, and pulls $g_i(t)$ back to $g_\infty(t)$ in
$C^\infty$ on that neighborhood times $[-T,0]$.  The exhaustion is chosen
with source-ball capture.  These data, rather than a bare convergence
symbol, are what imply completeness of every limit slice.

\subsection{Two point-selection lemmas}

\begin{lemma}[Spatial point selection]
\label{lem:ksol-spatial-point-selection}
Let $(N,h)$ be complete, let $R_h\geq0$ be smooth and unbounded, and fix
$p\in N$.  There are points $x_i$ and radii $r_i\in(0,1/2]$ such that, with
$Q_i=R_h(x_i)$,
\begin{align}
 d_h(p,x_i)&\longrightarrow\infty,& Q_i&\longrightarrow\infty,
                                                        \label{eq:ksol-pointpick-escape}\\
 Q_ir_i^2&\longrightarrow\infty,&
 R_h&\leq4Q_i\quad\hbox{on }B_h(x_i,r_i).              \label{eq:ksol-pointpick-control}
\end{align}
In addition $d_h(p,x_i)/r_i\to\infty$.
\end{lemma}

\begin{proof}
Choose $y_i$ escaping every compact set with $R_h(y_i)\to\infty$.  The
closed unit ball about $y_i$ is compact.  On it maximize
\[
       F_i(z)=R_h(z)d_h\bigl(z,\partial B_h(y_i,1)\bigr)^2.
\]
Let $x_i$ be a maximizer, let
$\sigma_i=d_h(x_i,\partial B_h(y_i,1))$, and put $r_i=\sigma_i/2$.
Then $F_i(x_i)\geq F_i(y_i)=R_h(y_i)$ and hence
$Q_ir_i^2=F_i(x_i)/4\to\infty$.  If $z\in B_h(x_i,r_i)$, its distance from
the boundary is at least $\sigma_i/2$, and maximality gives
$R_h(z)\leq4Q_i$.  The remaining conclusions follow from
$d_h(p,x_i)\geq d_h(p,y_i)-1$ and $r_i\leq1/2$.
\end{proof}

\begin{lemma}[Line in an escaping pointed Riemannian limit]
\label{lem:ksol-Riemannian-line-limit}
Let $(N,h)$ be complete, noncompact, and have nonnegative sectional
curvature, and fix $p\in N$.  Let $x_i\in N$, let $\lambda_i>0$, and put
$D_i=d_h(p,x_i)$.  Suppose
\begin{equation}
                    D_i\longrightarrow\infty,\qquad
                    \lambda_iD_i\longrightarrow\infty.       \label{eq:ksol-Riemannian-line-hypotheses}
\end{equation}
If $(N,\lambda_i^2h,x_i)$ converges smoothly and pointedly to a complete
Riemannian manifold $(N_\infty,h_\infty,x_\infty)$, then
$(N_\infty,h_\infty)$ contains a line through $x_\infty$.
\end{lemma}

\begin{proof}
Pass to a subsequence so that $D_{i+1}/D_i\to\infty$.  Choose unit-speed
minimizing segments $\gamma_i:[0,D_i]\to N$ from $p$ to $x_i$ and pass
again so that their initial directions converge.  The segments converge on
compact parameter intervals to a ray from $p$, so
$\angle x_i p x_{i+1}\to0$.  Let $\alpha_i$ be a minimizing segment from
$x_i$ to $x_{i+1}$.  Toponogov says that the Euclidean comparison angle at
$p$ is no larger than the actual angle.  The Euclidean law of cosines and
$D_{i+1}/D_i\to\infty$ then give
\begin{equation}
                  \widetilde\angle p x_i x_{i+1}\longrightarrow\pi.
                                                               \label{eq:ksol-line-comparison-angle}
\end{equation}

Fix $\rho>0$.  By \eqref{eq:ksol-Riemannian-line-hypotheses} and the
choice of subsequence, there are points $p_i\in\gamma_i$ and
$q_i\in\alpha_i$ with
\[
              d_h(p_i,x_i)=d_h(q_i,x_i)=\rho/\lambda_i.
\]
Toponogov comparison-angle monotonicity and
\eqref{eq:ksol-line-comparison-angle} imply that the comparison angle
$\widetilde\angle p_i x_i q_i$ tends to $\pi$.  Consequently,
\begin{equation}
                         \lambda_i d_h(p_i,q_i)\longrightarrow2\rho.
                                                               \label{eq:ksol-line-distance-equality}
\end{equation}
Pointed convergence captures these points and their minimizing connectors.
Diagonalizing over integer $\rho$ produces two unit-speed limiting rays
$\gamma_-,\gamma_+$ from $x_\infty$ with
$d(\gamma_-(\rho),\gamma_+(\rho))=2\rho$ for every $\rho>0$.

For $0<a\leq b$, the equal-radius identity at $b$ and the triangle
inequality give
\[
 2b\leq d(\gamma_-(b),\gamma_-(a))
       +d(\gamma_-(a),\gamma_+(b))
 =(b-a)+d(\gamma_-(a),\gamma_+(b)).
\]
The reverse bound
$d(\gamma_-(a),\gamma_+(b))\leq a+b$ follows through $x_\infty$, so
equality holds.  The case $b\leq a$ is symmetric.  Thus the concatenation
of the two rays is a line.
\end{proof}

\begin{lemma}[Line in an escaping curvature-normalized limit]
\label{lem:ksol-dimension-reduction-line}
Let $(M,g(t))$, $t\leq0$, be a complete smooth flow with nonnegative
curvature operator and $p\in M$.  Suppose $x_i,r_i,Q_i$ satisfy
\begin{equation}
 Q_i=R(x_i,0),\quad
 Q_i d_{g(0)}(p,x_i)^2\to\infty,\quad Q_ir_i^2\to\infty,       \label{eq:ksol-line-hypotheses}
\end{equation}
and
\begin{equation}
 R(\cdot,t)\leq C Q_i
 \quad\hbox{on }B_{g(0)}(x_i,r_i)\times(-\infty,0].            \label{eq:ksol-line-cylinder-bound}
\end{equation}
If the locally bounded rescaled flows $Q_i g(t/Q_i)$ have a complete
pointed smooth limit
with $R(x_\infty,0)=1$, then the time-zero limit contains a line through
$x_\infty$.  The whole limiting flow splits isometrically as
\begin{equation}
       (M_\infty,g_\infty(t))
       =(W^{n-1},h(t))\times(\R,ds^2).                         \label{eq:ksol-dimension-reduction-product}
\end{equation}
\end{lemma}

\begin{proof}
Put $D_i=d_{g(0)}(p,x_i)$ and $\lambda_i=\sqrt{Q_i}$.  Properness and
continuity of $R(\cdot,0)$ show that $D_i\to\infty$: otherwise $Q_i$ would
be bounded on one compact ball, contrary to $Q_iD_i^2\to\infty$.  The same
hypothesis gives $\lambda_iD_i\to\infty$.  Apply
Lemma~\ref{lem:ksol-Riemannian-line-limit} to the convergent time-zero
slices.  The limiting slice contains a line, so Cheeger--Gromoll gives the
product at time zero.  The expanding local bounds
\eqref{eq:ksol-line-cylinder-bound} give one global curvature bound on the
limiting flow.  The
curvature-null two-forms are parallel by the strong maximum principle of
Chapter~\ref{ch:strong-curvature}.  To make the terminal time an interior
time, use complete bounded-curvature short-time existence from
$g_\infty(0)$ and bounded-curvature uniqueness to extend and glue the flow
slightly past zero.  The whole-flow nullspace clause then gives one fixed
distribution and the same product for every $t\leq0$.
\end{proof}

\section{The two-dimensional ancient theorem}
\label{sec:ksol-surface}

The surface theorem is proved from Chapter~\ref{ch:reduced-geometry} and a
short compact entropy identity.  This avoids the longer Type~I/Type~II and
cigar route.

\begin{lemma}[Hamilton's compact surface entropy]
\label{lem:ksol-surface-entropy}
Let $(\Sigma^2,h(t))$ be a Ricci flow on a closed connected surface with
$R>0$.  Put
\begin{equation}
 A(t)=\int_\Sigma d\mu,
 \qquad C=\int_\Sigma R\,d\mu=4\pi\chi(\Sigma),
 \qquad r(t)=\frac{C}{A(t)},                             \label{eq:ksol-surface-entropy-data}
\end{equation}
and
\begin{equation}
             \mathcal N(h)=\int_\Sigma R\log(RA)\,d\mu. \label{eq:ksol-surface-entropy}
\end{equation}
Then $\mathcal N$ is invariant under diffeomorphism and constant rescaling,
and
\begin{equation}
 \frac{d}{dt}\mathcal N
 =-\int_\Sigma R|\nabla\log R-\nabla f|^2d\mu
  -2\int_\Sigma\left|\nabla^2f-\frac{\Delta f}{2}h\right|^2d\mu\leq0,       \label{eq:ksol-surface-entropy-square}
\end{equation}
where $f$ is the unique mean-zero solution of $\Delta f=r-R$.
Moreover,
\begin{equation}
                      \mathcal N(h)\geq C\log C,        \label{eq:ksol-surface-entropy-lower}
\end{equation}
with equality if and only if $R$ is spatially constant.
\end{lemma}

\begin{proof}
Under surface Ricci flow,
\[
 \partial_t h=-Rh,\quad \partial_tR=\Delta R+R^2,
 \quad\partial_t d\mu=-R\,d\mu,\quad A'=-C.
\]
Direct differentiation and integration by parts first give
\begin{equation}
 \mathcal N'=-\int\frac{|\nabla R|^2}{R}\,d\mu
                 +\int(R-r)^2d\mu.                    \label{eq:ksol-surface-entropy-first}
\end{equation}
Solvability and uniqueness of $f$ follow because $\int(r-R)d\mu=0$.
Integration by parts gives
\(
 \int\langle\nabla R,\nabla f\rangle d\mu
       =\int(R-r)^2d\mu.
\)
The integrated Bochner identity in dimension two gives
\begin{equation}
 \int(R-r)^2d\mu
 =\int|\nabla^2f|^2d\mu+\frac12\int R|\nabla f|^2d\mu.   \label{eq:ksol-surface-Bochner}
\end{equation}
Expanding the two squares on the right of
\eqref{eq:ksol-surface-entropy-square}, using
$|\nabla^2f-(\Delta f/2)h|^2=|\nabla^2f|^2-(\Delta f)^2/2$, gives exactly
\eqref{eq:ksol-surface-entropy-first}.

For the probability measure $A^{-1}d\mu$, the random variable $Z=AR$
has mean $C$, and $\mathcal N=\mathbb E[Z\log Z]$.  Jensen's inequality for
$z\mapsto z\log z$ proves \eqref{eq:ksol-surface-entropy-lower}; strict
convexity gives its equality clause.
\end{proof}

\begin{theorem}[Two-dimensional ancient $\kappa$-solutions]
\label{thm:ksol-two-dimensional-classification}
Let $(\Sigma^2,h(t))$, $t\leq0$, be an ancient $\kappa$-solution.  There
are $T>0$ and a fixed Riemannian covering $\pi:\Sph^2\to\Sigma$ such that
\begin{equation}
             \pi^*h(t)=2(T-t)g_{\Sph^2(1)}
             \quad\hbox{for every }t\leq0.             \label{eq:ksol-surface-round-flow}
\end{equation}
The deck group is trivial or is generated by the antipodal map.  Hence the
only models are the shrinking round $\Sph^2$ and $\mathbb{RP}^2$; the
oriented model is $\Sph^2$.
\end{theorem}

\begin{proof}
The scalar strong maximum principle gives $R>0$.  Apply the asymptotic
shrinker theorem of Chapter~\ref{ch:reduced-geometry}.  The limit is a
complete nonflat two-dimensional shrinker, so
Theorem~\ref{thm:complete-two-dimensional-shrinker-classification} makes it
a round $\Sph^2$ or $\mathbb{RP}^2$.  In particular the limit is compact.
For large indices, a convergence embedding of the whole compact limit into
$\Sigma$ has image open by invariance of domain and closed by compactness.
Connectedness makes it onto.  Thus $\Sigma$ is compact and the blow-down
convergence is global.

Use the notation of Lemma~\ref{lem:ksol-surface-entropy}.  The entropy is
scale and diffeomorphism invariant, so global convergence of
$\tau_i^{-1}h(-\tau_i)$ to the round limit gives
\begin{equation}
                    \mathcal N(h(-\tau_i))\longrightarrow C\log C.       \label{eq:ksol-surface-entropy-backward-limit}
\end{equation}
Because $\mathcal N$ is nonincreasing forward and bounded below by
$C\log C$, its limit as $t\to-\infty$ exists; the displayed subsequence
identifies it with $C\log C$.  Hence, for every finite $t$,
\[
 C\log C\leq\mathcal N(h(t))
       \leq\lim_{s\to-\infty}\mathcal N(h(s))=C\log C.
\]
The equality clause gives spatially constant $R$ at every time.  Since
$A'=-C$, there is $T=A(0)/C>0$ with $A(t)=C(T-t)$ and
$R(t)=1/(T-t)$.  Integrating $\partial_t h=-Rh$ gives
$h(t)=((T-t)/T)h(0)$, and uniformization gives
\eqref{eq:ksol-surface-round-flow}.  Finally
$\chi(\Sph^2)=|\Gamma|\chi(\Sigma)$ forces a free round deck group to have
order one or two; its nontrivial element is conjugate to the antipodal map.
\end{proof}

\begin{theorem}[Two-dimensional boundedness upgrade]
\label{thm:ksol-two-dimensional-KLim-bounded}
If $g(t)$ satisfies $\mathsf{KLim}_{2,\kappa}$, then
\begin{equation}
              \sup_{\Sigma\times(-\infty,0]}R<\infty.  \label{eq:ksol-KLim2-global-bound}
\end{equation}
Consequently it is an ancient $\kappa$-solution in the sense of
Definition~\ref{def:red-kappa-solution} and is one of the two models in
Theorem~\ref{thm:ksol-two-dimensional-classification}.
\end{theorem}

\begin{proof}
Taking $X=0$ in \eqref{eq:ksol-KLim-Harnack} gives $\partial_tR\geq0$.
If the conclusion failed, $R(\cdot,0)$ would be unbounded.  Apply
Lemma~\ref{lem:ksol-spatial-point-selection}, rescale by
$Q_i=R(x_i,0)$, and use the Harnack inequality to extend the bound
$R\leq4Q_i$ backward in time on the selected balls.  Noncollapsing and the
local curvature bound give a uniform injectivity-radius lower bound at the
rescaled basepoints.  Local Hamilton compactness and
$r_i\sqrt{Q_i}\to\infty$ give a complete ancient limit with
$R(x_\infty,0)=1$ and $R\leq4$.

Lemma~\ref{lem:ksol-dimension-reduction-line} makes the limit split as a
line times a one-dimensional flow.  Every one-dimensional Riemannian
manifold is flat, contradicting $R(x_\infty,0)=1$.  Thus
$\sup_\Sigma R(\cdot,0)<\infty$; monotonicity gives the same bound for all
$t\leq0$.  Nonnegative curvature converts it to the required $|\Rm|$ bound,
and Theorem~\ref{thm:ksol-two-dimensional-classification} applies.
\end{proof}

\section{The rank-one branch and its quotients}
\label{sec:ksol-split-branch}

Here ``rank one'' refers to the curvature operator on $\Lambda^2$; on the
round cylinder the Ricci tensor has rank two.

\begin{theorem}[Complete three-dimensional splitting branch]
\label{thm:ksol-three-dimensional-split-branch}
Let $(M^3,g(t))$, $t\leq0$, be an ancient $\kappa$-solution.  Exactly one of
the following holds.
\begin{enumerate}
\item The curvature operator is positive at every spacetime point.
\item For some $T>0$, the universal spacetime cover is
\begin{equation}
 \mathcal C_T(t)=
 \bigl(\Sph^2\times\R,
       2(T-t)g_{\Sph^2(1)}+ds^2\bigr),                \label{eq:ksol-cylinder-cover}
\end{equation}
and, up to conjugacy, the deck group is one of
\begin{equation}
 \{1\},\qquad
 \langle a\rangle,\ a(y,s)=(-y,s),\qquad
 \langle b\rangle,\ b(y,s)=(-y,-s).                  \label{eq:ksol-cylinder-deck-list}
\end{equation}
Thus the quotient is respectively $\Sph^2\times\R$,
$\mathbb{RP}^2\times\R$, or the diagonal $\mathbb Z_2$ quotient.  If $M$
is orientable, the middle model is excluded.
\end{enumerate}
\end{theorem}

\begin{proof}
If a curvature-operator eigenvalue vanishes anywhere, first extend the
complete bounded-curvature flow slightly beyond zero and glue by
uniqueness if the vanishing point is terminal.  The whole-flow strong
maximum principle of Chapter~\ref{ch:strong-curvature} then gives a fixed
universal cover
\[
             (\widetilde M,\widetilde g(t))
             =(\Sigma^2,h(t))\times(\R,ds^2).
\]
The cover preserves noncollapsing.  If $|\Rm_h|\leq r^{-2}$ on
$B_h(y,r)$, then the product ball obeys the same curvature bound and
\[
 \kappa r^3
 \leq\Vol_{h+ds^2}B((y,0),r)
 \leq2r\,\operatorname{Area}_hB_h(y,r).
\]
Thus the surface factor is $(\kappa/2)$-noncollapsed.  It is nonflat and
inherits the global curvature bound, so
Theorem~\ref{thm:ksol-two-dimensional-classification} proves
\eqref{eq:ksol-cylinder-cover}.

Every deck transformation preserves the positive and zero Ricci
eigenspaces, hence has the form
\begin{equation}
                      \gamma(y,s)=(Ay,\epsilon s+c),
       \qquad A\in O(3),\quad\epsilon\in\{1,-1\}.       \label{eq:ksol-deck-product-form}
\end{equation}
There can be no nonzero translation in the induced action on $\R$.  Indeed,
after taking a power, such an element has translation length $L>0$ and its
sphere part fixes some $y_0$.  At a time tending to $-\infty$, put
$\rho=\sqrt{T-t}$ and $r=\sqrt{L\rho}$.  For $\rho$ large, the curvature on
the radius-$r$ ball is at most $r^{-2}$, whereas a fundamental-slab estimate
in the translation quotient gives
\begin{equation}
       \Vol B([y_0,0],r)\leq \pi Lr^2,
       \qquad \frac{\Vol B([y_0,0],r)}{r^3}
                    \leq\pi\sqrt{\frac L\rho}\longrightarrow0.             \label{eq:ksol-translation-collapse}
\end{equation}
A further quotient only decreases ball volume.  This contradicts
all-scale $\kappa$-noncollapsing.

The line action therefore has trivial image or is generated by one
reflection; two distinct reflections would multiply to a translation.
Translate the coordinate so the reflection is $s\mapsto-s$.  The kernel
acts freely on each round sphere.  Euler characteristic shows that it is
trivial or generated by the antipodal map $a$.  If the line image is
trivial, this gives the first two groups in
\eqref{eq:ksol-cylinder-deck-list}.  Otherwise choose
$\delta=(A,-s)$.  If the kernel is trivial, then $A^2=I$ and freeness makes
$A=-I$, giving $\langle b\rangle$.  If the kernel is $\langle a\rangle$,
the alternative $A^2=-I$ is impossible because
$\det(A^2)=1\neq-1=\det(-I)$; hence $A^2=I$ and again $A=-I$.  But then
$a\delta=(I,-s)$ fixes the central sphere, a contradiction.  The list is
complete.  The antipodal map reverses orientation on $\Sph^2$; hence $a$
reverses product orientation while $b$ preserves it.
\end{proof}

\section{Spatial invariants and the geometry at infinity}
\label{sec:ksol-spatial-infinity}

The extended-real formulation is important: the value $+\infty$ must not be
silently coerced to a real number.

\begin{definition}[AVR and ASCR]
\label{def:ksol-AVR-ASCR}
Let $(N^n,h)$ be complete and noncompact, let $p\in N$, and suppose
$\Ric_h\geq0$ when defining $\AVR$ and $R_h\geq0$ when defining $\ASCR$.
Put
\begin{align}
 \AVR_p(h)
   &=\inf_{r>0}\frac{\Vol_hB_h(p,r)}{\omega_n r^n}
     \in[0,1],                                          \label{eq:ksol-AVR-definition}\\
 \ASCR_p(h)
   &=\inf_{\rho>0}\ \sup_{d_h(p,x)\geq\rho}
           R_h(x)d_h(p,x)^2
     \in[0,+\infty].                                   \label{eq:ksol-ASCR-definition}
\end{align}
The supremum and infimum in \eqref{eq:ksol-ASCR-definition} are taken in
the extended nonnegative reals.
\end{definition}

Bishop--Gromov says that the volume ratio in
\eqref{eq:ksol-AVR-definition} is nonincreasing in $r$, so its infimum is
its limit at infinity.  Formula \eqref{eq:ksol-ASCR-definition} is the usual
limsup as $d(p,x)\to\infty$.  Directly from the tail infimum, the quantifier
form used later is
\begin{equation}
 \ASCR(h)=+\infty
 \Longleftrightarrow
 \bigl(\forall A,D>0\ \exists x:\
       d_h(p,x)\geq D\ \hbox{ and }R_h(x)d_h(p,x)^2>A\bigr).    \label{eq:ksol-ASCR-quantifiers}
\end{equation}

\begin{lemma}[Basepoint and scaling invariance]
\label{lem:ksol-spatial-invariants-well-defined}
For any $p,q\in N$,
\begin{equation}
             \AVR_p(h)=\AVR_q(h),\qquad
             \ASCR_p(h)=\ASCR_q(h).                  \label{eq:ksol-basepoint-invariance}
\end{equation}
If $\widehat h=\lambda h$ for $\lambda>0$, then
\begin{equation}
             \AVR_p(\widehat h)=\AVR_p(h),\qquad
             \ASCR_p(\widehat h)=\ASCR_p(h).         \label{eq:ksol-invariant-scaling}
\end{equation}
All equalities for ASCR are equalities in the extended nonnegative reals.
\end{lemma}

\begin{proof}
Put $d=d_h(p,q)$.  For $r>d$, the triangle inequality gives
\begin{equation}
 B_h(p,r-d)\subset B_h(q,r)\subset B_h(p,r+d).        \label{eq:ksol-basepoint-ball-sandwich}
\end{equation}
Divide the corresponding volume inequalities by $\omega_nr^n$ and let
$r\to\infty$.  The factors $((r\pm d)/r)^n$ tend to one, proving the AVR
equality.

For ASCR, fix $\delta\in(0,1)$.  Outside the ball
$B_h(p,d/\delta)$ one has
\begin{equation}
 (1-\delta)d_h(p,x)\leq d_h(q,x)
       \leq(1+\delta)d_h(p,x).                       \label{eq:ksol-distance-tail-comparison}
\end{equation}
Escaping the $p$-balls is equivalent to escaping the $q$-balls.  If either
ASCR is infinite, \eqref{eq:ksol-ASCR-quantifiers} and
\eqref{eq:ksol-distance-tail-comparison} show that both are infinite.  If
both are finite, take tail suprema in
\eqref{eq:ksol-distance-tail-comparison}.  More explicitly, for all
sufficiently large $\rho$,
\[
 \sup_{d_h(q,x)\geq\rho}R(x)d_h(q,x)^2
 \leq(1+\delta)^2
     \sup_{d_h(p,x)\geq\rho-d}R(x)d_h(p,x)^2.
\]
Taking decreasing infima gives
$\ASCR_q(h)\leq(1+\delta)^2\ASCR_p(h)$.  Interchanging $p$ and $q$ gives
the reverse inequality with the same factor.  Letting
$\delta\downarrow0$ proves equality.  No finite real representative is
chosen until finiteness has been established.

For $\widehat h=\lambda h$, distances multiply by $\sqrt\lambda$, volumes
by $\lambda^{n/2}$, and scalar curvature by $\lambda^{-1}$.  Thus the ball
volume ratio and the product $R(x)d(p,x)^2$ are unchanged after the
corresponding change of radius.  This proves
\eqref{eq:ksol-invariant-scaling}.
\end{proof}

We henceforth omit the basepoint.

\begin{theorem}[AVR invariance under decay at the later slice]
\label{thm:ksol-invariant-in-time}
Let $(N^n,g(t))$ be a complete noncompact ancient Ricci flow with
nonnegative curvature operator, and suppose it satisfies the ancient trace
Harnack inequality \eqref{eq:ksol-KLim-Harnack}.  Fix $p\in N$ and
$t_1<t_2$, and assume
\begin{equation}
             \sup_{N\times[t_1,t_2]}|\Rm|<\infty.    \label{eq:ksol-AVR-slab-bound}
\end{equation}
If
\begin{equation}
 |\Rm|(x,t_2)\longrightarrow0
 \quad\hbox{as }d_{g(t_2)}(p,x)\longrightarrow\infty,  \label{eq:ksol-final-slice-curvature-decay}
\end{equation}
then
\begin{equation}
                         \AVR(g(t_1))=\AVR(g(t_2)).    \label{eq:ksol-AVR-time-invariance}
\end{equation}
\end{theorem}

\begin{proof}
Write $\Delta t=t_2-t_1$.  For every $\varepsilon>0$,
\eqref{eq:ksol-final-slice-curvature-decay} and the norm bridge
\eqref{eq:red-nonnegative-curvature-bridges} give a compact set $K$ such
that $R(x,t_2)\leq\varepsilon$ off $K$.  Harnack gives
$0\leq R(x,t)\leq R(x,t_2)$ on $[t_1,t_2]$, so the same $K$ works
uniformly on the whole slab.  Since $g(t_1)\geq g(t_2)$,
\[
                         B_{t_1}(p,r)\subset B_{t_2}(p,r).
\]
Moreover
\[
 d\mu_{t_2}(x)
 =\exp\!\left(-\int_{t_1}^{t_2}R(x,t)\,dt\right)d\mu_{t_1}(x)
 \geq e^{-\varepsilon\Delta t}d\mu_{t_1}(x)
\]
off $K$.  Hence
\[
 \Vol_{t_2}B_{t_2}(p,r)
 \geq e^{-\varepsilon\Delta t}
      \bigl(\Vol_{t_1}B_{t_1}(p,r)-\Vol_{t_1}K\bigr).
\]
Divide by $\omega_nr^n$, let $r\to\infty$, and then let
$\varepsilon\downarrow0$.  This gives
$\AVR(g(t_2))\geq\AVR(g(t_1))$.

For the opposite inequality, bounded curvature on the slab and Hamilton's
additive distance estimate give a constant $C<\infty$ such that
\[
 d_{t_2}(p,x)\geq d_{t_1}(p,x)-C.
\]
Thus $B_{t_2}(p,r)\subset B_{t_1}(p,r+C)$.  Since
$d\mu_{t_2}\leq d\mu_{t_1}$,
\[
 \Vol_{t_2}B_{t_2}(p,r)\leq\Vol_{t_1}B_{t_1}(p,r+C).
\]
After division by $\omega_nr^n$ and passage to infinity, this gives the
reverse inequality.  The final-slice decay hypothesis is essential.
\end{proof}

\begin{lemma}[Finite ASCR forces positive AVR]
\label{lem:ksol-finite-ASCR-positive-AVR}
Let $(N^n,h)$ be complete and noncompact with $\Rm\geq0$, suppose it is
$\kappa$-noncollapsed on all scales, and suppose $\ASCR(h)<+\infty$.
Let $A\in[0,\infty)$ be the unique real number whose image in the extended
nonnegative reals is $\ASCR(h)$.  Then $\AVR(h)>0$.  More precisely, if
$|\Rm|\leq c_nR$, then, with
\begin{equation}
 \alpha=\min\left\{\frac12,
       \frac1{2\sqrt{c_n(A+1)}}\right\},              \label{eq:ksol-ASCR-alpha}
\end{equation}
one has
\begin{equation}
       \AVR(h)\geq\frac{\kappa}{\omega_n}
          \left(\frac{\alpha}{1+\alpha}\right)^n>0.  \label{eq:ksol-finite-ASCR-AVR-bound}
\end{equation}
\end{lemma}

\begin{proof}
For some $\rho$, the definition of $A$ gives
$R(y)d(p,y)^2\leq A+1$ outside $B(p,\rho)$.  Let $x$ lie on a ray from
$p$ with $d(p,x)=D>2\rho$; such a ray follows from Hopf--Rinow by taking a
limit of minimizing segments to points escaping every compact set.  If
$y\in B(x,\alpha D)$, then
$d(p,y)\geq D/2$, and hence
\[
 |\Rm|(y)\leq4c_n(A+1)D^{-2}\leq(\alpha D)^{-2}.
\]
Noncollapsing gives $\Vol B(x,\alpha D)\geq\kappa(\alpha D)^n$.
This ball lies in $B(p,(1+\alpha)D)$.  Divide by
$\omega_n((1+\alpha)D)^n$ and let $D\to\infty$.
\end{proof}

\begin{lemma}[Asymptotic shrinkers of a noncompact three-solution]
\label{lem:ksol-noncompact-asymptotic-shrinker}
Every asymptotic shrinker of a noncompact three-dimensional ancient
$\kappa$-solution is the round cylinder or one of the two order-two
cylindrical quotients in Theorem~\ref{thm:ksol-three-dimensional-split-branch}.
Each has $\AVR=0$.
\end{lemma}

\begin{proof}
Chapter~\ref{ch:reduced-geometry} gives a complete nonflat shrinker, and
Chapter~\ref{ch:three-dimensional-shrinkers} gives its exact model list.
The limit cannot be compact: if it were, a convergence embedding of the
whole compact limit into the connected source manifold would be open and
closed, hence onto, making the source compact.  The Gaussian model is
excluded by nonflatness.  The remaining noncompact models are precisely
the cylindrical quotients.  Their ball volume grows at most linearly, so
their three-dimensional asymptotic volume ratio is zero.
\end{proof}

\begin{lemma}[Point selection from infinite ASCR]
\label{lem:ksol-ASCR-point-selection}
Let $(N,h,p)$ be complete and noncompact with bounded scalar curvature and
$\ASCR(h)=+\infty$.  There are pairwise disjoint balls $B_h(x_i,r_i)$ and
numbers $\varepsilon_i\downarrow0$ such that
\begin{align}
 d_h(p,x_i)&\to\infty,&
 \sup_{B_h(x_i,r_i)}R&\leq(1+\varepsilon_i)R(x_i),      \label{eq:ksol-ASCR-point-control}\\
 R(x_i)r_i^2&\to\infty,&
 d_h(p,x_i)/r_i&\to\infty.                             \label{eq:ksol-ASCR-point-radii}
\end{align}
\end{lemma}

\begin{proof}
Take $\varepsilon_i\downarrow0$ and choose $y_i$ escaping so fast that,
with $d_i=d(p,y_i)$,
\begin{equation}
               \varepsilon_i^2R(y_i)d_i^2\longrightarrow\infty.
                                                               \label{eq:ksol-ASCR-selection-rate}
\end{equation}
On $\overline B(y_i,d_i/4)$ maximize
\[
          F_i(z)=R(z)d\bigl(z,\partial B(y_i,d_i/4)\bigr)^2.
\]
Let $x_i$ be a maximizer, put
$\sigma_i=d(x_i,\partial B(y_i,d_i/4))$,
$\eta_i=1-(1+\varepsilon_i)^{-1/2}$, and
$r_i=\eta_i\sigma_i$.  If $z\in B(x_i,r_i)$, then its boundary distance is
at least $(1-\eta_i)\sigma_i$; maximality gives
$R(z)\leq(1-\eta_i)^{-2}R(x_i)=(1+\varepsilon_i)R(x_i)$.
Moreover
\[
 R(x_i)r_i^2=\eta_i^2F_i(x_i)
 \geq\frac{\eta_i^2}{16}R(y_i)d_i^2\longrightarrow\infty
\]
because $\eta_i/\varepsilon_i\to1/2$.  Since
$r_i/d(p,x_i)\leq C\eta_i\to0$, the last ratio in
\eqref{eq:ksol-ASCR-point-radii} follows.  Passing to a still faster
escaping subsequence makes the balls disjoint.  Every point chosen above is
a maximum on a compact closed ball.
\end{proof}

\begin{theorem}[Spatial infinity of noncompact three-dimensional
$\kappa$-solutions]
\label{thm:ksol-ASCR-AVR}
Let $(M^3,g(t))$, $t\leq0$, be a noncompact ancient $\kappa$-solution.
Then, on every time slice,
\begin{equation}
                 \ASCR(g(t))=+\infty,
                 \qquad \AVR(g(t))=0.                 \label{eq:ksol-ASCR-infty-AVR-zero}
\end{equation}
\end{theorem}

\begin{proof}
Fix a time $t_0$.

\emph{Step 1: $\ASCR=+\infty$.}
Suppose instead that $\ASCR(g(t_0))<+\infty$, and let $A\in[0,\infty)$ be
its unique finite real representative.  The finite tail bound makes
curvature tend to zero at spatial infinity on the
time-$t_0$ slice.  Lemma~\ref{lem:ksol-finite-ASCR-positive-AVR} gives
\[
                         v:=\AVR(g(t_0))>0.
\]
For every $s<t_0$, Theorem~\ref{thm:ksol-invariant-in-time}, applied with
later time $t_0$, gives $\AVR(g(s))=v$.

Translate $t_0$ to zero and take an $\ell$-center blow-down from
Chapter~\ref{ch:reduced-geometry}.
At its time-one slice, every approximating metric is a constant rescaling
of $g(-\tau_i)$ and therefore has AVR equal to $v$.  Bishop--Gromov
then gives, for every fixed $r>0$,
\[
 \Vol B(q_i,r)\geq v\omega_3r^3.
\]
Source-ball capture and smooth convergence pass this inequality to the
limit.  Thus the asymptotic shrinker has AVR at least $v$, contradicting
Lemma~\ref{lem:ksol-noncompact-asymptotic-shrinker}.  Since $t_0$ was
arbitrary, $\ASCR=+\infty$ on all slices.

\emph{Step 2: $\AVR=0$.}
In the rank-one branch this follows directly from the quotient list: every
model has at most linear volume growth.  Suppose therefore that sectional
curvature is positive.  If $\AVR(g(t_0))=v>0$, apply
Lemma~\ref{lem:ksol-ASCR-point-selection} at time $t_0$.  Put
$Q_i=R(x_i,t_0)$ and rescale by $Q_i$, translating $t_0$ to zero.  The
time-zero bounds in \eqref{eq:ksol-ASCR-point-control}, the ancient trace
Harnack inequality, noncollapsing, local Bando--Shi estimates, and Hamilton
compactness give a complete nonflat ancient limit with globally bounded
nonnegative curvature.  Lemma~\ref{lem:ksol-dimension-reduction-line}
makes it split off a line.  The surface factor is a two-dimensional
$\kappa'$-solution, so Theorem
\ref{thm:ksol-two-dimensional-classification} makes its universal cover a
shrinking round sphere.  The limit is a cylindrical quotient and has AVR
zero.

On the other hand, AVR is invariant under scaling and basepoint change on
the fixed slice $g(t_0)$.  Every rescaled time-zero metric therefore has AVR
$v$.  Bishop--Gromov and pointed smooth convergence, exactly as in Step 1,
give AVR at least $v$ for the limit.  This contradiction proves
$\AVR(g(t_0))=0$.  Since $t_0$ was arbitrary, the theorem follows.
\end{proof}

The proof deliberately avoids the older finite-positive-ASCR asymptotic
cone argument and the zero-ASCR Petrunin--Tuschmann gap theorem.  The
already established asymptotic-shrinker theorem makes those detours
unnecessary.

\begin{lemma}[Passing noncollapsing through a smooth pointed limit]
\label{lem:ksol-pass-noncollapse}
Suppose pointed complete $n$-manifolds converge smoothly with source-ball
capture and every source metric is $\kappa$-noncollapsed on all scales.
Then the limit metric is $\kappa$-noncollapsed on all scales.
\end{lemma}

\begin{proof}
Fix a limit point $z$, $r>0$, and suppose
$|\Rm|\leq r^{-2}$ on $B(z,r)$.  Choose
$r'<\rho<r$.  Because $r^{-2}<(r')^{-2}$, smooth convergence on the
compact $\rho$-ball and ball capture imply, for large $i$, that the source
$r'$-ball lies in the convergence chart and satisfies
$|\Rm_i|\leq(r')^{-2}$.  Hence its volume is at least
$\kappa(r')^n$.  Fix $\delta>0$ with $r'+\delta<\rho$.  For all large $i$,
source-ball capture and distance convergence put the pulled-back source
$r'$-ball inside $B_{g_\infty}(z,r'+\delta)$.  Smooth convergence of the
volume forms therefore gives
\[
             \Vol_{g_\infty}B(z,r'+\delta)\geq\kappa(r')^n.
\]
Let $\delta\downarrow0$ using continuity of metric-ball volume, and then
let $r'\uparrow r$.  The two strictly buffered radii are what make both
the non-strict curvature threshold and the ball-volume inequality pass
correctly.
\end{proof}

\section{Almost-ancient collapse at the curvature scale}
\label{sec:ksol-almost-ancient}

The next proposition is the quantitative engine behind compactness.  Its
project-facing form keeps the fixed noncollapsing constant which is present
in every application.  This makes the proof depend only on
Theorem~\ref{thm:ksol-ASCR-AVR}, rather than on the stronger classical
theorem that every bounded nonflat ancient flow with nonnegative curvature
has zero AVR.

\begin{proposition}[Fixed-$\kappa$ almost-ancient collapse]
\label{prop:ksol-almost-ancient-collapse}
Fix $\kappa>0$.  For every $\varepsilon>0$ there exist
$A=A(\kappa,\varepsilon)\geq1$ and
$L=L(\kappa,\varepsilon)\geq A^2$ with the following property.
Suppose $g(t)$ satisfies $\mathsf{KLim}_{3,\kappa}$, $x\in M$,
$Q=R(x,0)>0$, and $r>0$ satisfy
\begin{equation}
 R(y,t)\leq2Q
 \quad\text{for }(y,t)\in B_{g(0)}(x,r)\times(-\infty,0],
 \qquad r^2Q\geq L.                                  \label{eq:ksol-almost-ancient-hyp}
\end{equation}
Then
\begin{equation}
 \frac{\Vol_{g(0)}B_{g(0)}(x,AQ^{-1/2})}
      {(AQ^{-1/2})^3}\leq\varepsilon.                 \label{eq:ksol-almost-ancient-conclusion}
\end{equation}
\end{proposition}

\begin{proof}
If the assertion were false, diagonalize with $A_i=i$ and $L_i=i^2$.
After scaling by $Q_i$, the counterexamples have
\[
 R(x_i,0)=1,\qquad R\leq2
 \quad\hbox{on }B_{g_i(0)}(x_i,r_i\sqrt{Q_i})
                    \times(-\infty,0],
\]
where $r_i\sqrt{Q_i}\to\infty$, and
\begin{equation}
                 \frac{\Vol B(x_i,i)}{i^3}>\varepsilon.        \label{eq:ksol-almost-ancient-countervolume}
\end{equation}
Bishop--Gromov transfers this lower volume ratio to the unit ball.
The local curvature bound, this volume lower bound, and the local
injectivity-radius theorem give a uniform injectivity radius at $x_i$.
Local Bando--Shi estimates and Hamilton compactness produce a complete
ancient pointed limit through time zero.  The increasing spatial radii give
the global bound $R\leq2$ on the limit, while
$R(x_\infty,0)=1$.  Nonnegative curvature and $\kappa$-noncollapsing pass
to the limit coefficientwise and by
Lemma~\ref{lem:ksol-pass-noncollapse}, respectively, so it is a noncompact
ancient $\kappa$-solution.  It is
noncompact because \eqref{eq:ksol-almost-ancient-countervolume}, after
Bishop--Gromov, gives for every fixed $\rho$,
\[
       \Vol_{g_\infty(0)}B(x_\infty,\rho)/\rho^3
                   \geq\varepsilon.
\]
Thus its AVR is at least $\varepsilon/\omega_3$, contradicting
Theorem~\ref{thm:ksol-ASCR-AVR}.
\end{proof}

The classical stronger version allows incomplete finite cylinders, assumes
$R\leq2Q$ on a backward interval, and requires both the scaled time depth
and scaled spatial radius to tend to infinity.  The proof is identical after
diagonalization.  It will be recorded when surgery spacetime domains are
introduced; the ancient version above is the exact consumer needed here.

\section{Anchored volume, bounded-distance curvature, and a seed ball}
\label{sec:ksol-anchored-control}

\begin{lemma}[Bounded curvature at bounded distance from one volume anchor]
\label{lem:ksol-bounded-curvature-bounded-distance}
Fix $\kappa,v>0$ and $D<\infty$.  There is
$C=C(\kappa,v,D)<\infty$ such that every flow satisfying
$\mathsf{KLim}_{3,\kappa}$ and
\begin{equation}
       \Vol_{g(0)}B_{g(0)}(x,1)\geq v                    \label{eq:ksol-anchor-data}
\end{equation}
obeys
\begin{equation}
                 \sup_{B_{g(0)}(x,D)}R(\cdot,0)\leq C. \label{eq:ksol-bounded-distance-bound}
\end{equation}
\end{lemma}

\begin{proof}
Suppose otherwise.  In a counterexample sequence choose
$z_i\in B(x_i,D)$ with $R(z_i,0)\to\infty$.  On the compact ball
$\overline B(x_i,D+1)$ maximize
\[
              F_i(y)=R(y,0)\bigl(D+1-d(x_i,y)\bigr)^2.
\]
Let $w_i$ be a maximizer, put
$\sigma_i=D+1-d(x_i,w_i)$,
$Q_i=R(w_i,0)$, and
$s_i=(1-2^{-1/2})\sigma_i$.  Since $F_i(w_i)\geq F_i(z_i)$,
\begin{equation}
 Q_is_i^2\to\infty,
 \qquad R(\cdot,0)\leq2Q_i\quad\hbox{on }B(w_i,s_i).    \label{eq:ksol-bounded-distance-pointpick}
\end{equation}
Indeed, for $y\in B(w_i,s_i)$ the boundary-distance factor is at least
$2^{-1/2}\sigma_i$, and maximality gives $R(y,0)\leq2Q_i$.
Also $w_i\in B(x_i,D+1)$.

The Harnack inequality at $X=0$ gives $\partial_tR\geq0$, so the same
curvature bound holds on the fixed time-zero ball at every earlier time.
Apply Proposition~\ref{prop:ksol-almost-ancient-collapse} with
\begin{equation}
                    \varepsilon_0=\frac{v}{2(D+2)^3}.   \label{eq:ksol-anchor-epsilon}
\end{equation}
For one fixed $A$ and all large $i$, it gives an $A Q_i^{-1/2}$ ball about
$w_i$ whose volume ratio is at most $\varepsilon_0$.

On the other hand, $B(x_i,1)\subset B(w_i,D+2)$.  Bishop--Gromov at $w_i$
therefore gives, for every $0<\rho\leq D+2$,
\begin{equation}
 \frac{\Vol B(w_i,\rho)}{\rho^3}
 \geq\frac{\Vol B(w_i,D+2)}{(D+2)^3}
 \geq\frac{v}{(D+2)^3}.                               \label{eq:ksol-anchor-BG-lower}
\end{equation}
Since $AQ_i^{-1/2}\to0$, this contradicts
\eqref{eq:ksol-anchor-epsilon}.
\end{proof}

Notice that every ball in the final comparison is at time zero.  No
changing-distance estimate is concealed in this proof.

\begin{proposition}[Uniform seed volume at the scalar scale]
\label{prop:ksol-seed-volume}
For every $\kappa>0$ there is $v_\kappa>0$ such that, if a flow satisfying
$\mathsf{KLim}_{3,\kappa}$ has $R(x,0)=1$, then
\begin{equation}
                     \Vol_{g(0)}B_{g(0)}(x,1)\geq v_\kappa.    \label{eq:ksol-seed-volume}
\end{equation}
\end{proposition}

\begin{proof}
Assume that a sequence has unit-ball volumes tending to zero.  Write
\(\vartheta_i(r)=\Vol B(x_i,r)/r^3\).  Smooth small-ball asymptotics give
$\vartheta_i(r)\to\omega_3$ as $r\downarrow0$, while
$\vartheta_i(1)\to0$.  By continuity choose $\delta_i\in(0,1)$ with
\begin{equation}
                     \vartheta_i(\delta_i)=\omega_3/2. \label{eq:ksol-half-Euclidean-scale}
\end{equation}
The inequality
$\Vol B(x_i,1)\geq(\omega_3/2)\delta_i^3$ shows that
$\delta_i\to0$.

Rescale by
\begin{equation}
                        h_i(t)=\delta_i^{-2}g_i(\delta_i^2t).   \label{eq:ksol-seed-rescaling}
\end{equation}
Then
\begin{equation}
 R_{h_i}(x_i,0)=\delta_i^2\to0,
 \qquad \Vol_{h_i(0)}B_{h_i(0)}(x_i,1)=\omega_3/2.     \label{eq:ksol-seed-rescaled-data}
\end{equation}
Lemma~\ref{lem:ksol-bounded-curvature-bounded-distance} gives scalar
curvature bounds on every fixed time-zero ball.  Choose a fixed small
$\rho>0$ from the bound on the radius-two ball so that
$|\Rm_{h_i}|\leq\rho^{-2}$ on $B_{h_i(0)}(x_i,\rho)$.
Noncollapsing gives volume at least $\kappa\rho^3$, and the local
injectivity-radius estimate gives a positive lower bound for
$\inj_{h_i(0)}(x_i)$.

Harnack bounds curvature on the same spatial sets backward in time;
Bando--Shi and Hamilton compactness now give a complete ancient pointed
limit.  Lemma~\ref{lem:ksol-pass-noncollapse} makes it
$\kappa$-noncollapsed, and it has
$R_{h_\infty}(x_\infty,0)=0$.  Harnack gives
$R(x_\infty,t)=0$ for every $t<0$.  For each fixed $t<0$, apply the scalar
strong maximum principle on a slab having $t$ as an interior time.  It
gives $R(\cdot,t)=0$.  Thus every negative-time slice is flat, and
continuity gives flatness at time zero as well.

There is a particularly short way to identify the complete flat limit.
For every $r>0$, flatness and $\kappa$-noncollapsing give the hypotheses of
the local injectivity-radius theorem on the radius-$r$ ball.  Hence
\begin{equation}
                    \inj_{h_\infty(0)}(x_\infty)
                          \geq c(3,\kappa)r.
\end{equation}
Letting $r\to\infty$ makes the injectivity radius infinite, so the
exponential map is a global Euclidean isometry.  Its unit-ball volume is
$\omega_3$.  Smooth ball convergence, squeezed between radii
$1-\eta$ and $1+\eta$ and followed by $\eta\downarrow0$, instead passes
\eqref{eq:ksol-seed-rescaled-data} to give unit-ball volume $\omega_3/2$.
This contradiction proves the proposition.
\end{proof}

\begin{corollary}[Normalized bounded-distance curvature]
\label{cor:ksol-normalized-bounded-distance}
For every $\kappa>0$ and $D<\infty$ there is $C(\kappa,D)<\infty$ such
that every flow satisfying $\mathsf{KLim}_{3,\kappa}$ with $R(x,0)=1$
satisfies
\begin{equation}
                   R(y,0)\leq C(\kappa,D)
       \quad\text{when }d_{g(0)}(x,y)\leq D.           \label{eq:ksol-normalized-local-R}
\end{equation}
\end{corollary}

\begin{proof}
Combine Proposition~\ref{prop:ksol-seed-volume} with
Lemma~\ref{lem:ksol-bounded-curvature-bounded-distance}.
\end{proof}

\section{Fixed-$\kappa$ preliminary compactness}
\label{sec:ksol-fixed-compactness}

Fix once and for all a dimensional norm-bridge constant
$C_{\mathrm{Rm},3}$ such that $|\Rm|\leq C_{\mathrm{Rm},3}R$ whenever the
three-dimensional curvature operator is nonnegative.

\begin{lemma}[Buffered parabolic derivative bounds]
\label{lem:ksol-buffered-parabolic-bounds}
For fixed $\kappa,D,T>0$ and integers $a,b\geq0$, there is
$C_{a,b}(\kappa,D,T)<\infty$ such that every normalized flow in
Corollary~\ref{cor:ksol-normalized-bounded-distance} satisfies
\begin{equation}
 |\nabla_t^b\nabla^a\Rm|(y,t)\leq C_{a,b}(\kappa,D,T)  \label{eq:ksol-buffered-derivative-bound}
\end{equation}
on $B_{g(0)}(x,D)\times[-T,0]$.
\end{lemma}

\begin{proof}
Apply Corollary~\ref{cor:ksol-normalized-bounded-distance} on the enlarged
time-zero ball $B_{g(0)}(x,D+2)$.  Harnack gives
$R(y,t)\leq R(y,0)$ for every fixed $y$ and $t\leq0$.  Nonnegative
curvature gives $|\Rm|\leq C_{\mathrm{Rm},3}R$.  Moreover
$g(t)\geq g(0)$ for $t\leq0$,
so
\[
 B_{g(t)}(y,1)\subset B_{g(0)}(y,1)
                    \subset B_{g(0)}(x,D+1).
\]
Apply the local Bando--Shi estimates of Chapter~\ref{ch:bando-shi} on
past-buffered intervals $[s-1,s]$ and on smaller balls.  This gives spatial
derivatives also at the terminal time.  The curvature evolution equations
and induction give the intrinsic time and mixed derivatives.  At $t=0$ all
time derivatives mean left derivatives, equivalently limits from $t<0$.
\end{proof}

\begin{lemma}[Normalized base injectivity]
\label{lem:ksol-normalized-injectivity}
For every $\kappa>0$ there is $\iota_\kappa>0$ such that a normalized flow
as above satisfies
\begin{equation}
                         \inj_{g(0)}(x)\geq\iota_\kappa.        \label{eq:ksol-base-injectivity}
\end{equation}
\end{lemma}

\begin{proof}
Use Corollary~\ref{cor:ksol-normalized-bounded-distance} on $B(x,1)$ and
choose
\[
 \rho=\min\{1,(C_{\mathrm{Rm},3}C(\kappa,1))^{-1/2}\}.
\]
Then $|\Rm|\leq\rho^{-2}$ on $B(x,\rho)$, noncollapsing gives volume at
least $\kappa\rho^3$, and the local injectivity-radius theorem gives
\eqref{eq:ksol-base-injectivity}.
\end{proof}

\begin{theorem}[Preliminary curvature-normalized limit]
\label{thm:ksol-preliminary-compactness}
Fix $\kappa>0$.  Let $(M_i^3,g_i(t),x_i)$, $t\leq0$, satisfy
$\mathsf{KLim}_{3,\kappa}$ and $R_{g_i}(x_i,0)=1$.  A subsequence converges
smoothly and pointedly on every compact spacetime set to a connected,
complete ancient flow $(M_\infty,g_\infty(t),x_\infty)$ which satisfies
\begin{equation}
 R_{g_\infty}(x_\infty,0)=1,
 \qquad\Rm_{g_\infty}\geq0,                            \label{eq:ksol-preliminary-limit-data}
\end{equation}
is $\kappa$-noncollapsed on all scales, and satisfies the trace Harnack
inequality.  Thus it satisfies $\mathsf{KLim}_{3,\kappa}$.
\end{theorem}

\begin{proof}
Lemmas~\ref{lem:ksol-buffered-parabolic-bounds} and
\ref{lem:ksol-normalized-injectivity}, followed by Hamilton compactness and
diagonalization over integer $D,T$, give the smooth limit and compatible
exhaustion embeddings.  Source-ball capture proves completeness at time
zero.  Since $\Ric\geq0$, one has $g_\infty(t)\geq g_\infty(0)$ for $t<0$;
therefore every closed $g_\infty(t)$-ball is contained in a compact
$g_\infty(0)$-ball and is compact.  This proves completeness of all slices.
The curvature inequality and normalization pass coefficientwise.
Lemma~\ref{lem:ksol-pass-noncollapse} passes noncollapsing.  Smooth
spacetime convergence, including $\partial_tR$, $\nabla R$, and $\Ric$,
passes the Harnack quadratic for every limiting vector.  The normalized
scalar makes the limit nonflat.
\end{proof}

The theorem has deliberately not yet called the limit an ancient
$\kappa$-solution.  The constants in
\eqref{eq:ksol-normalized-local-R} may grow with $D$, so they do not furnish
one global spatial bound.  The next section supplies the indispensable
closure theorem.

\section{Cylindrical slices and the boundedness upgrade}
\label{sec:ksol-Harnack-closure}

We first isolate the one global geometric input needed to convert a
locally controlled Harnack limit into a globally bounded solution.

\begin{definition}[Spatial $\varepsilon$-neck witness]
\label{def:ksol-spatial-neck}
Fix $\varepsilon>0$.  Let $(N^3,h)$ be complete, let $p\in N$ have
$R_h(p)>0$, and set
$r_p=\sqrt2\,R_h(p)^{-1/2}$.  A \emph{spatial $\varepsilon$-neck witness
centered at $p$} is an embedding
\begin{equation}
 \Phi:\Sph^2\times(-\varepsilon^{-1}-1,\varepsilon^{-1}+1)\longrightarrow N,
 \qquad \Phi(y_*,0)=p,                                 \label{eq:ksol-spatial-neck-embedding}
\end{equation}
such that $r_p^{-2}\Phi^*h$ is within $\varepsilon$ in
$C^{\lceil\varepsilon^{-1}\rceil}$ of
$g_{\Sph^2(1)}+dz^2$ on the subcylinder with
$|z|\leq\varepsilon^{-1}$.  For a chosen witness, write
\begin{equation}
 \mathcal N_\Phi=\operatorname{image}(\Phi),\qquad
 \mathcal C_\Phi=\Phi(\Sph^2\times[-\varepsilon^{-1},\varepsilon^{-1}]),
 \qquad S_\Phi=\Phi(\Sph^2\times\{0\}).               \label{eq:ksol-spatial-neck-witness-data}
\end{equation}
The point $p$ is a \emph{center of a spatial $\varepsilon$-neck} when such
a witness exists.  Later separation arguments always quantify over a
chosen witness; they never recover its embedding from this existential
proposition.
\end{definition}

\begin{lemma}[A smooth cylinder limit detects spatial necks]
\label{lem:ksol-spatial-neck-detection}
Let $(N_i^3,h_i,x_i)$ be pointed complete Riemannian manifolds with
$R_{h_i}(x_i)=1$.  Suppose they converge smoothly and pointedly to
\begin{equation}
 \bigl(\Sph^2\times\R,
       2g_{\Sph^2(1)}+ds^2,(y_*,0)\bigr).             \label{eq:ksol-spatial-cylinder-limit}
\end{equation}
Then, for every $\varepsilon>0$, the point $x_i$ is the center of a spatial
$\varepsilon$-neck for all sufficiently large $i$.
\end{lemma}

\begin{proof}
Let $\Psi(y,z)=(y,\sqrt2z)$.  Since $R_{h_i}(x_i)=1$, the scale in
Definition~\ref{def:ksol-spatial-neck} is exactly $r_{x_i}=\sqrt2$, and
\[
 r_{x_i}^{-2}\Psi^*
       \bigl(2g_{\Sph^2(1)}+ds^2\bigr)
       =g_{\Sph^2(1)}+dz^2.
\]
Use one convergence embedding on the image under $\Psi$ of the fixed
compact buffered cylinder in
\eqref{eq:ksol-spatial-neck-embedding}.  Composing it with $\Psi$ gives
the chosen witness, and smooth convergence controls the finite required
norm on its core.
\end{proof}

\begin{lemma}[The two sides of a chosen spatial neck]
\label{lem:ksol-spatial-neck-two-sides}
Let $(N^3,h)$ be complete, noncompact, and have positive sectional
curvature, and let $\Phi$ be a chosen spatial $\varepsilon$-neck witness.
Then $N\setminus S_\Phi$ has exactly two components.  Exactly one has
compact closure; denote that closure by $K_\Phi$ and denote the other
component by $U_\Phi$.  After replacing $\Phi(y,z)$ by $\Phi(y,-z)$ if
necessary, the negative end of $\mathcal N_\Phi$ faces
$\operatorname{int}K_\Phi$ and the positive end faces $U_\Phi$.
\end{lemma}

\begin{proof}
The positive-curvature soul theorem makes $N$ diffeomorphic to $\R^3$.
The smooth embedded sphere $S_\Phi$ therefore separates $N$ into one
bounded and one unbounded component, and smooth Schoenflies identifies the
closure of the bounded component with a three-ball.  The two ends of the
product neighborhood $\mathcal N_\Phi\setminus S_\Phi$ lie in different
components.  Reversing the product coordinate gives the asserted
orientation.
\end{proof}

\begin{lemma}[Busemann-level diameter through a good neck]
\label{lem:ksol-neck-Busemann-diameter}
There is a dimensional $\bar\varepsilon$ with
$0<\bar\varepsilon<(5\pi)^{-1}$ and the following property for every
$\varepsilon\in(0,\bar\varepsilon]$.
Let $(N^3,h)$ be complete, noncompact, and have positive sectional
curvature.  Let $\gamma$ be a unit-speed ray and use the convex-sign
Busemann function
\begin{equation}
             b_\gamma(x)=\lim_{s\to\infty}
                    \bigl(s-d_h(\gamma(s),x)\bigr).    \label{eq:ksol-convex-Busemann}
\end{equation}
Let $\Phi$ be a chosen spatial $\varepsilon$-neck witness centered at
$y$, oriented as in Lemma~\ref{lem:ksol-spatial-neck-two-sides}.  Then,
with $b_y=b_\gamma(y)$,
\begin{equation}
 \frac{11}{12}\pi\sqrt2\,R_h(y)^{-1/2}
 \leq\operatorname{diam}b_\gamma^{-1}(b_y)
 \leq11\pi\sqrt2\,R_h(y)^{-1/2}.                      \label{eq:ksol-neck-level-diameter-bound}
\end{equation}
Moreover $b_\gamma$ is bounded below and attains its minimum.
\end{lemma}

\begin{proof}
We record the finite comparison argument used later.  Use the oriented
cylindrical coordinate $u$ of the chosen witness.  For a slice
$\Phi(\Sph^2\times\{c\})$, ``inward'' means the component on its
$K_\Phi$ side and ``outward'' means the component on its $U_\Phi$ side.
Toponogov comparison and the
$C^{\lceil\varepsilon^{-1}\rceil}$ closeness to the round cylinder gives,
after decreasing $\bar\varepsilon$ once and for all, the following two
implications for
$A=5\pi$:
\begin{align}
 q\text{ inward of }\{u=-A\}
     &\Longrightarrow b_\gamma(q)<b_y,                 \label{eq:ksol-Busemann-inward}\\
 q\text{ outward of }\{u=A\}
     &\Longrightarrow b_\gamma(q)>b_y.                 \label{eq:ksol-Busemann-outward}
\end{align}
Put $a=\sqrt2\,R(y)^{-1/2}$.  For the first implication, join $q$ to
$\gamma(s)$ and let $z_s$ be its intersection with the center sphere.
Cylindrical closeness gives
$d(z_s,y)\leq(6/5)\pi a$ and
$d(z_s,q)>(5/6)Aa$.  Therefore
\begin{align*}
 s-d(\gamma(s),q)
 &\leq s-d(\gamma(s),y)+d(z_s,y)-d(z_s,q)\\
 &\leq s-d(\gamma(s),y)
       +\frac65\pi a-\frac56Aa
 <s-d(\gamma(s),y).
\end{align*}
Let $s\to\infty$.  For the second implication, compare the triangle with
vertices $y,q,\gamma(s)$ after both long sides cross the slice $u=A$.
Cylindrical closeness makes the relevant comparison angle at $y$ at most
$\pi/6$.  Toponogov and the Euclidean law of cosines give
\[
 d(q,\gamma(s))^2
 \leq d(y,\gamma(s))^2+d(y,q)^2
       -d(y,\gamma(s))d(y,q).
\]
For large $s$, elementary rearrangement gives
\[
 d(y,\gamma(s))
 \geq d(q,\gamma(s))+\frac14d(y,q).
\]
Hence
$b_\gamma(q)\geq b_\gamma(y)+d(y,q)/4>b_\gamma(y)$.
This proves the two displayed implications.

Thus $b_\gamma^{-1}(b_y)$ lies in the compact neck band $|u|\leq5\pi$.
Cylindrical closeness bounds its diameter above by the right side of
\eqref{eq:ksol-neck-level-diameter-bound}.  At $u=-5\pi$ and $u=5\pi$,
continuity from points strictly beyond the two slices turns
\eqref{eq:ksol-Busemann-inward}--\eqref{eq:ksol-Busemann-outward} into the
weak inequalities $b_\gamma(\Phi(z,-5\pi))\leq b_y$ and
$b_\gamma(\Phi(z,5\pi))\geq b_y$.  Consequently, for every
$z\in\Sph^2$, the intermediate value theorem provides
$u(z)\in[-5\pi,5\pi]$ with
$b_\gamma(\Phi(z,u(z)))=b_y$.  Choosing $z$ antipodal to the sphere
coordinate of $y$ and using cylindrical closeness gives the lower bound.
Finally \eqref{eq:ksol-Busemann-outward} shows that the infimum of
$b_\gamma$ is attained on the compact set $K_\Phi\cup\mathcal C_\Phi$.
\end{proof}

\begin{lemma}[Outward subsequence of disjoint escaping necks]
\label{lem:ksol-neck-outward-subsequence}
Let $\bar\varepsilon$ be supplied by
Lemma~\ref{lem:ksol-neck-Busemann-diameter}, fix
$\varepsilon\in(0,\bar\varepsilon]$, and retain that lemma's $N,h,\gamma$.
Let
$p_*\in b_\gamma^{-1}(\min_N b_\gamma)$.  Suppose $\Phi_i$ are chosen
spatial $\varepsilon$-neck witnesses whose controlled cores
$\mathcal C_{\Phi_i}$ are pairwise disjoint and whose centers $y_i$ escape
every compact set.  Then, after discarding finitely many terms and passing
to a subsequence,
\begin{equation}
                    K_{\Phi_i}\subset
                    \operatorname{int}K_{\Phi_{i+1}}.          \label{eq:ksol-neck-nested-subsequence}
\end{equation}
\end{lemma}

\begin{proof}
The point $p_*$ belongs to at most one of the pairwise disjoint controlled
cores.  Discard that term.  For every remaining $i$, the point $p_*$
cannot lie in $U_{\Phi_i}$: it would then be beyond the positive boundary
$u=\varepsilon^{-1}>5\pi$ of the controlled core, so
\eqref{eq:ksol-Busemann-outward} would give
$b_\gamma(p_*)>b_\gamma(y_i)$, contrary to minimality.  Hence
$p_*\in\operatorname{int}K_{\Phi_i}$.

Two disjoint embedded spheres in $\R^3$ whose compact sides contain the
same point have nested compact sides.  Fix one index $i$.  Since
$K_{\Phi_i}$ is compact and $y_j$ escapes, some later center $y_j$ lies
outside $K_{\Phi_i}$.  The connected sphere $S_{\Phi_j}$ is disjoint from
$S_{\Phi_i}$, so it lies entirely in $U_{\Phi_i}$; nesting then forces
$K_{\Phi_i}\subset\operatorname{int}K_{\Phi_j}$.  Iteration constructs
the required subsequence.
\end{proof}

\begin{lemma}[Surjective Sharafutdinov level map]
\label{lem:ksol-Sharafutdinov-level-surjective}
Let $(N,h)$ be complete and noncompact with nonnegative sectional curvature,
let $b$ be the Busemann function of a ray in the sign convention
\eqref{eq:ksol-convex-Busemann}, and suppose $b$ is bounded below.  If
$c_1<c_2$ are above its minimum and both levels are nonempty, then the
Sharafutdinov level map
\begin{equation}
             \pi_{c_2,c_1}:b^{-1}(c_2)\longrightarrow b^{-1}(c_1)
                                                               \label{eq:ksol-Sharafutdinov-level-map}
\end{equation}
in the inward direction is surjective and distance nonincreasing.  Hence
\begin{equation}
              \operatorname{diam}b^{-1}(c_1)
                    \leq\operatorname{diam}b^{-1}(c_2).        \label{eq:ksol-Busemann-level-diameter-monotone}
\end{equation}
\end{lemma}

\begin{proof}
Use $f=c_2-b$ on the compact totally convex slab cut off by the level
$b=c_2$.  Sharafutdinov's construction follows the normalized generalized
gradient of the concave function $-b$ and gives a distance-nonincreasing
retraction onto each inner superlevel set.  Its restriction to the outer
level is the map in \eqref{eq:ksol-Sharafutdinov-level-map}.

The surjectivity of this restriction is a separate part of the level-flow
interface and does not follow merely by calling the map a retraction.  Fix
$q\in b^{-1}(c_1)$.  Choose $s_j\to\infty$ and minimizing segments from
$q$ to $\gamma(s_j)$.  Properness and the first variation estimate give a
subsequential limiting coray $\sigma_q:[0,\infty)\to N$ satisfying
\begin{equation}
             b(\sigma_q(s))=c_1+s\qquad(s\geq0).      \label{eq:ksol-coray-linear-busemann}
\end{equation}
Put $p=\sigma_q(c_2-c_1)$.  Equality in the unit Lipschitz bound for $b$
along every subsegment of $\sigma_q$ says precisely that the reversed
segment from $p$ to $q$ is a unit-speed curve of maximal descent for $b$.
The characterization and forward uniqueness of normalized generalized
gradient curves for the concave function $-b$ therefore identify this
reversed segment with the Sharafutdinov trajectory starting at $p$.
Consequently $\pi_{c_2,c_1}(p)=q$.  Thus every $q$ has a preimage; no
backward solution of a possibly noninvertible gradient semigroup was used.
For
$q_1,q_2$ choose preimages $p_1,p_2$; distance nonincrease gives
$d(q_1,q_2)\leq d(p_1,p_2)$.  Taking suprema proves
\eqref{eq:ksol-Busemann-level-diameter-monotone}.
\end{proof}

\begin{proposition}[Rough comparison of outward neck radii]
\label{prop:ksol-outward-neck-comparison}
There is a dimensional $\varepsilon_*$ with
$0<\varepsilon_*\leq\bar\varepsilon<(5\pi)^{-1}$ and the following
property, where $\bar\varepsilon$ is the constant from
Lemma~\ref{lem:ksol-neck-Busemann-diameter}.  Let
$(N^3,h)$ be complete, noncompact, and have positive sectional curvature.
Let $\Phi_1,\Phi_2$ be chosen spatial $\varepsilon_*$-neck witnesses,
centered at $y_1,y_2$, whose controlled cores are disjoint.  Orient both
witnesses as in Lemma~\ref{lem:ksol-spatial-neck-two-sides}.  If
\begin{equation}
              K_{\Phi_1}\subset\operatorname{int}K_{\Phi_2},  \label{eq:ksol-neck-outward-order}
\end{equation}
then
\begin{equation}
                            R_h(y_2)\leq144R_h(y_1).    \label{eq:ksol-neck-radius-comparison}
\end{equation}
\end{proposition}

\begin{proof}
Choose one unit-speed ray $\gamma$ and use its function
\eqref{eq:ksol-convex-Busemann} for both witnesses.
Put $b_i=b_\gamma(y_i)$.  The separation conclusions
\eqref{eq:ksol-Busemann-inward}--\eqref{eq:ksol-Busemann-outward}, the
disjointness of the controlled cores, and
\eqref{eq:ksol-neck-outward-order} place $y_2$ beyond the positive end of
$\Phi_1$ and give $b_1<b_2$.
Lemma~\ref{lem:ksol-Sharafutdinov-level-surjective} gives a surjective
distance-nonincreasing retraction from the farther convex Busemann level to
the nearer one.  Hence
\[
       \operatorname{diam}b_\gamma^{-1}(b_1)
       \leq\operatorname{diam}b_\gamma^{-1}(b_2).
\]
Apply Lemma~\ref{lem:ksol-neck-Busemann-diameter} at $y_1$ and $y_2$:
\[
 \frac{11}{12}\pi\sqrt2\,R(y_1)^{-1/2}
       \leq11\pi\sqrt2\,R(y_2)^{-1/2}.
\]
Squaring gives \eqref{eq:ksol-neck-radius-comparison}.
\end{proof}

This proposition is deliberately much weaker than a neck--cap theorem.  It
is a one-slice obstruction to arbitrarily thinner outward necks, and that is
exactly what the closure proof consumes.

\begin{lemma}[Terminal rank-one product]
\label{lem:ksol-terminal-rank-one-product}
Let $g(t)$ satisfy $\mathsf{KLim}_{3,\kappa}$, and suppose the curvature
operator of $g(0)$ has rank one everywhere.  Then the universal cover of
the terminal slice has an isometric product decomposition
\begin{equation}
       (\widetilde M,\widetilde g(0))
       \cong(\Sigma^2\times\R,h_0+ds^2),              \label{eq:ksol-KLim-terminal-product}
\end{equation}
where $(\Sigma,h_0)$ is complete and simply connected and has positive
Gaussian curvature.  No curvature bound on $M$ is assumed.
\end{lemma}

\begin{proof}
Restrict to the closed slab $[-1,0]$ and shift it to $[0,1]$.  Apply
Theorem~\ref{thm:smp-systems} to the curvature operator in the fixed
Uhlenbeck bundle.  Its positive-time rank statement and the terminal
rank-one hypothesis give $\delta>0$ for which the rank is one on
$M\times(1-\delta,1]$.  On the open interval $(1-\delta,1)$,
Corollary~\ref{cor:curvature-fixed-bundle-rank-interval} gives one smooth
spatially parallel rank-two kernel bundle $K$.

We spell out the one-sided endpoint transfer because $1$ is not an
interior time.  If $v\in K_x$, then $\mathcal M(x,s)v=0$ for $s<1$;
continuity gives $\mathcal M(x,1)v=0$.  Both kernels have dimension two,
so $K_x=\ker\mathcal M(x,1)$.  If $w$ is a local smooth section of $K$,
then $D_Y^sw\in K$ for $s<1$.  Smooth convergence $D^s\to D^1$ and
closedness of the finite-dimensional fibers give $D_Y^1w\in K$.  Thus the
terminal kernel is a smooth parallel bundle.  This is precisely the
left-limit argument of Lemma~\ref{lem:curvature-one-sided-endpoint}, with
no use of a time beyond the terminal endpoint.

Proposition~\ref{prop:ue-nullspace-transfer} transfers $K$ to the parallel
nullspace of the geometric curvature operator of $g(0)$.  After lifting to
the complete simply connected universal cover,
Lemma~\ref{lem:curvature-image-to-parallel-line} gives a parallel tangent
line and Theorem~\ref{thm:parallel-line-splitting} gives
\eqref{eq:ksol-KLim-terminal-product}.  The factor is complete and simply
connected.  Lemma~\ref{lem:rank-one-normalization} identifies its Gaussian
curvature with the positive curvature-operator eigenvalue, so that
curvature is strictly positive.
\end{proof}

\begin{lemma}[Boundedness of a rank-one terminal slice]
\label{lem:ksol-rank-one-terminal-bounded}
Let $g(t)$ satisfy $\mathsf{KLim}_{3,\kappa}$.  Suppose the curvature
operator of $g(0)$ has rank one everywhere.  Then
\begin{equation}
                              \sup_MR(\cdot,0)<\infty.  \label{eq:ksol-rank-one-terminal-bound}
\end{equation}
\end{lemma}

\begin{proof}
Lemma~\ref{lem:ksol-terminal-rank-one-product} gives
\eqref{eq:ksol-KLim-terminal-product}.  Noncollapsing passes to the cover.
The product-ball
calculation from Theorem~\ref{thm:ksol-three-dimensional-split-branch}
therefore makes $(\Sigma,h_0)$ $(\kappa/2)$-noncollapsed on all scales.

Suppose $R_{h_0}$ were unbounded.  Apply
Lemma~\ref{lem:ksol-spatial-point-selection} on $\Sigma$, with a fixed
$p\in\Sigma$, to obtain $y_i,r_i,Q_i$ such that
\begin{equation}
 Q_i=R_{h_0}(y_i),\quad
 Q_i d_{h_0}(p,y_i)^2\to\infty,\quad
 Q_ir_i^2\to\infty,\quad
 R_{h_0}\leq4Q_i\text{ on }B_{h_0}(y_i,r_i).          \label{eq:ksol-rank-one-surface-pointpick}
\end{equation}
At time zero, the product ball about $(y_i,0)$ is contained in
$B_{h_0}(y_i,r_i)\times(-r_i,r_i)$, so its scalar curvature is at most
$4Q_i$.  The lifted Harnack inequality extends this bound backward on the
same points.  Local Bando--Shi estimates on successively smaller balls give
all spatial curvature derivatives at time zero.  The surface
noncollapsing estimate gives a normalized injectivity-radius lower bound at
$y_i$.  Hence the pointed surfaces
\begin{equation}
                       (\Sigma,Q_i h_0,y_i)            \label{eq:ksol-rank-one-surface-rescaling}
\end{equation}
subconverge smoothly on every bounded ball to a complete surface
$(\Sigma_\infty,h_\infty,y_\infty)$ with
\[
        R_{h_\infty}(y_\infty)=1,\qquad
        0\leq R_{h_\infty}\leq4.
\]

The escape conclusion \eqref{eq:ksol-pointpick-escape} in
Lemma~\ref{lem:ksol-spatial-point-selection} gives
$d_{h_0}(p,y_i)\to\infty$, while
\eqref{eq:ksol-rank-one-surface-pointpick} gives
$\sqrt{Q_i}d_{h_0}(p,y_i)\to\infty$.  Therefore
Lemma~\ref{lem:ksol-Riemannian-line-limit} applies directly to
\eqref{eq:ksol-rank-one-surface-rescaling}, and $h_\infty$ contains a
line.  Cheeger--Gromoll splits it as a product of two one-dimensional
manifolds, hence it is flat.  This contradicts
$R_{h_\infty}(y_\infty)=1$.  Therefore $R_{h_0}$ is bounded, and
\eqref{eq:ksol-KLim-terminal-product} gives
\eqref{eq:ksol-rank-one-terminal-bound} downstairs.
\end{proof}

\begin{theorem}[Three-dimensional boundedness upgrade]
\label{thm:ksol-three-dimensional-KLim-bounded}
If $g(t)$ satisfies $\mathsf{KLim}_{3,\kappa}$, then
\begin{equation}
                    \sup_{M\times(-\infty,0]}R<\infty. \label{eq:ksol-KLim3-global-bound}
\end{equation}
Thus it is an ancient $\kappa$-solution in the sense of
Definition~\ref{def:red-kappa-solution}.
\end{theorem}

\begin{proof}
Harnack with $X=0$ gives $R(x,t)\leq R(x,0)$ for $t\leq0$.  It is enough to
bound the time-zero slice.

If the curvature operator of $g(0)$ fails to be positive at some point,
the whole-flow strong maximum principle is not yet available: its
whole-flow form assumes the curvature bound being proved.  Instead apply
the terminal-contact form of
Theorem~\ref{thm:smp-systems} on the closed slab $[-1,0]$, followed by the
one-sided endpoint transfer and the algebraic rank-two exclusion from
Chapter~\ref{ch:strong-curvature}.  The time-zero rank is spatially
constant and belongs to $\{0,1,3\}$.  Since positivity fails somewhere,
the slice is either flat or rank one.  In the flat case Harnack gives
$R(x,t)=0$ for every $t\leq0$,
contrary to nonflatness.  In the rank-one case Lemma
\ref{lem:ksol-rank-one-terminal-bounded} bounds $R(\cdot,0)$, and Harnack
bounds every earlier slice.  This proves the theorem in the nonpositive
branch without circular use of whole-flow uniqueness.

Suppose now that $\Rm_{g(0)}>0$ everywhere and that $R(\cdot,0)$ is
unbounded.  Apply
Lemma~\ref{lem:ksol-spatial-point-selection}.  With
$Q_i=R(x_i,0)$, the selected balls have normalized radii tending to infinity
and scalar curvature at most four.  Harnack supplies this same local bound
for every earlier time.  Noncollapsing gives base injectivity; local
Bando--Shi and Hamilton compactness give a complete ancient limit with
$R(x_\infty,0)=1$ and $R\leq4$.  The repaired local
Lemma~\ref{lem:ksol-dimension-reduction-line} applies: unlike the literal
global formulation sometimes cited here, it uses only the point-picked
local cylinders and does not assume the boundedness conclusion being
proved.  The limit splits off a line.  Its surface factor is a
two-dimensional $\kappa'$-solution and hence is round.

The limit is the unquotiented cylinder.  Indeed, positive sectional
curvature makes the source $M$ diffeomorphic to $\R^3$.  Either cylindrical
order-two quotient contains a compact embedded $\mathbb{RP}^2$ at finite
distance from the basepoint.  A sufficiently large convergence embedding
would put such a surface in $\R^3$, which is impossible.  Fix the
dimensional $\varepsilon_*$ supplied by
Proposition~\ref{prop:ksol-outward-neck-comparison}.  Lemma
\ref{lem:ksol-spatial-neck-detection} therefore supplies, for large $i$, a
chosen spatial $\varepsilon_*$-neck witness $\Phi_i$ centered at $x_i$.

The controlled core has spatial size at most
$C(\varepsilon_*)Q_i^{-1/2}$, so it lies in the selected ball
$B(x_i,r_i)$ for large $i$ because $r_i\sqrt{Q_i}\to\infty$.  Its radius
tends to zero because $Q_i\to\infty$.  Since $x_i$ escapes every compact
set and $r_i\leq1/2$, pass to a subsequence for which the selected balls,
and hence the controlled cores $\mathcal C_{\Phi_i}$, are pairwise
disjoint.  Choose one ray $\gamma$ and use its single function $b_\gamma$
throughout.  Lemma~\ref{lem:ksol-neck-Busemann-diameter}, applied to one
witness, supplies $p_*\in b_\gamma^{-1}(\min b_\gamma)$.  Lemma
\ref{lem:ksol-neck-outward-subsequence} now gives a further subsequence
with $K_{\Phi_i}\subset\operatorname{int}K_{\Phi_{i+1}}$.  Fix its first
witness.  Proposition~\ref{prop:ksol-outward-neck-comparison} gives
\[
                         Q_i\leq144Q_{i_0}
\]
for every later $i$, contradicting $Q_i\to\infty$.  Thus
$\sup_MR(\cdot,0)<\infty$, and Harnack gives
\eqref{eq:ksol-KLim3-global-bound}.
\end{proof}

\begin{theorem}[Curvature-normalized compactness for fixed $\kappa$]
\label{thm:ksol-fixed-kappa-compactness}
Fix $\kappa>0$.  For each $i$, let $(M_i^3,g_i(t))$, $t\leq0$, be an ancient
$\kappa$-solution, choose $(x_i,t_i)$ with $t_i\leq0$, and put
$Q_i=R_{g_i}(x_i,t_i)$.  Then $Q_i>0$, and a subsequence of
\begin{equation}
 \widehat g_i(s)=Q_i g_i\!\left(t_i+\frac{s}{Q_i}\right),
 \qquad s\leq0,                                       \label{eq:ksol-fixed-kappa-normalized-sequence}
\end{equation}
converges smoothly and pointedly on every compact spacetime set to an
ancient $\kappa$-solution
$(M_\infty^3,g_\infty(s),x_\infty)$ with
\begin{equation}
                         R_{g_\infty}(x_\infty,0)=1.   \label{eq:ksol-fixed-limit-normalization}
\end{equation}
If every $M_i$ is noncompact, then $M_\infty$ is noncompact.
\end{theorem}

\begin{proof}
Positivity is Lemma~\ref{lem:ksol-scalar-positive}.  The normalized flows
satisfy $\mathsf{KLim}_{3,\kappa}$ and have base scalar one.
Theorem~\ref{thm:ksol-preliminary-compactness} gives a complete ancient
limit with nonnegative curvature, the same noncollapsing constant, Harnack,
and the normalization.  Theorem
\ref{thm:ksol-three-dimensional-KLim-bounded} supplies the one global
spacetime curvature bound required by Definition
\ref{def:red-kappa-solution}.  Hence the limit is genuinely an ancient
$\kappa$-solution.

If the limit were compact, then for large $i$ a convergence embedding of
all of $M_\infty$ into connected $M_i$ would have open image by invariance
of domain and closed image by compactness.  It would be surjective, making
$M_i$ compact.  This proves the last assertion.
\end{proof}

\section{A universal noncollapsing gap}
\label{sec:ksol-universal-gap}

Fixed-$\kappa$ compactness is sufficient for many contradiction arguments.
The classification program needs more: apart from the collapsing family of
round spherical space forms, the noncollapsing constant can be chosen once
and for all.  We derive this from the exact asymptotic-shrinker masses in
Chapter~\ref{ch:reduced-geometry}.  The only separate rigidity input is the
following elementary pinching argument.

\begin{lemma}[A round backward limit forces the round flow]
\label{lem:ksol-round-backward-rigidity}
Let $(M^3,g(t))$, $t\leq0$, be a connected compact ancient Ricci flow with
bounded nonnegative curvature operator.  Suppose that for some
$\tau_i\to\infty$ the time-one slices of
\[
        (M,\tau_i^{-1}g(-\tau_i\theta),q_i)
\]
converge smoothly and pointedly to a compact round metric of positive
sectional curvature.  Then every $g(t)$ has constant positive sectional
curvature.  Thus the flow is a shrinking round spherical-space-form flow.
\end{lemma}

\begin{proof}
Because the limit is compact, a convergence embedding of the whole limit
into $M$ has open image by invariance of domain and closed image by
compactness.  Connectedness makes it onto.  Hence the convergence at
$\theta=1$ is global.

Let $0\leq\alpha\leq\beta\leq\gamma$ be the curvature-operator eigenvalues,
and put $S=\alpha+\beta+\gamma=R/2$.  Global round convergence supplies
numbers $c_i\to1/3$, with $0<c_i\leq1/3$, such that
\begin{equation}
              \alpha(\cdot,-\tau_i)\geq c_iS(\cdot,-\tau_i).
                                                        \label{eq:ksol-round-backward-pinching}
\end{equation}
We verify the invariant-cone step rather than hide it in a pinching slogan.
For $0\leq c\leq1/3$, set
\[
 \mathcal C_c=\{A\geq0:\lambda_{\min}(A)\geq c\,\operatorname{tr}A\}.
\]
This is closed, convex, and invariant under parallel transport.  Up to a
common positive factor, the three-dimensional curvature-reaction ODE is
\[
 \dot\alpha=\alpha^2+\beta\gamma,\quad
 \dot\beta=\beta^2+\alpha\gamma,\quad
 \dot\gamma=\gamma^2+\alpha\beta.
\]
At a boundary point with $S=0$, all three eigenvalues and the reaction
vector vanish.  Suppose $S>0$ and normalize $S=1$.  Since
$\beta,\gamma\geq c$ and $\beta+\gamma=1-c$,
$\beta\gamma\geq c(1-2c)$.  Therefore
\begin{align}
 \dot\alpha-c\dot S
  &=(1+c)\beta\gamma-c(\beta+\gamma)^2 \notag\\
  &\geq c(1+c)(1-2c)-c(1-c)^2
    =c^2(1-3c)\geq0.                                 \label{eq:ksol-round-cone-boundary}
\end{align}
Chapter~\ref{ch:weak-maximum-principle-systems} preserves
$\mathcal C_c$.

Fix $(x,t)$.  For every sufficiently large $i$, one has $-\tau_i<t$, so
\eqref{eq:ksol-round-backward-pinching} and cone preservation give
$\alpha(x,t)\geq c_iS(x,t)$.  Let $i\to\infty$.  Since $\alpha$ is the
least of three numbers with sum $S$, this forces
$\alpha=\beta=\gamma=S/3$.  Each slice is pointwise isotropic, and Schur's
lemma makes its sectional curvature spatially constant.  It is positive:
otherwise the strong maximum principle would contradict the positive round
backward limit.  The Ricci-flow equation is now the homothetic round ODE.
\end{proof}

\begin{proposition}[Universal lower bound for reduced volume]
\label{prop:ksol-universal-reduced-volume-lower}
Let $(M^3,g(t))$, $t\leq0$, be an ancient $\kappa$-solution which is not a
shrinking round spherical-space-form flow.  For every pole $(p,0)$ and
every $\tau>0$,
\begin{equation}
                         \widetilde V_p(\tau)\geq e^{-1}.       \label{eq:ksol-universal-reduced-volume-lower}
\end{equation}
\end{proposition}

\begin{proof}
Apply Theorem~\ref{thm:red-asymptotic-shrinker} to the reduced geometry
based at $(p,0)$.  If the resulting asymptotic shrinker were spherical,
compact pointed convergence would make $M$ compact by the open-and-closed
argument, and Lemma~\ref{lem:ksol-round-backward-rigidity} would make the
original flow round.  This is excluded.  Corollary
\ref{cor:red-three-asymptotic-shrinker} therefore leaves the round cylinder
or one of its two order-two quotients.  Their normalized masses are,
respectively, $2/e$, $1/e$, and $1/e$ by
\eqref{eq:red-three-density-regression}.  The no-loss theorem
\ref{thm:red-no-mass-loss} identifies this mass with
$\widetilde V_{p,\infty}$.  Since reduced volume is nonincreasing in
backward time,
\[
        \widetilde V_p(\tau)\geq\widetilde V_{p,\infty}\geq e^{-1}.
\]
\end{proof}

\begin{lemma}[A collapsed curvature ball has small reduced volume]
\label{lem:ksol-collapse-small-reduced-volume}
Choose a dimensional constant $a_3$ such that
$R\leq a_3|\Rm|$ in dimension three, and put
\begin{equation}
                    \chi_3=\frac1{576(3+1)^2}.        \label{eq:ksol-collapse-chi3}
\end{equation}
Let $(M^3,g(t))$, $t\leq0$, be an ancient $\kappa$-solution.  Suppose that
for some $p\in M$ and $r,v>0$,
\begin{equation}
 |\Rm|_{g(0)}\leq r^{-2}\quad\hbox{on }B_{g(0)}(p,r),
 \qquad
 \Vol_{g(0)}B_{g(0)}(p,r)<vr^3.                       \label{eq:ksol-collapsed-terminal-ball}
\end{equation}
For $0<\theta\leq1$, define
\begin{equation}
 \mathcal T(\theta)=
 e^{1+a_3\theta/3}(4\pi\theta)^{-3/2}\omega_3
 \sum_{j=1}^{\infty}(j+1)^3e^{-\chi_3j^2/\theta}.     \label{eq:ksol-collapse-tail}
\end{equation}
Then
\begin{equation}
 \widetilde V_p(\theta r^2)
 \leq(4\pi\theta)^{-3/2}e^{a_3\theta}v
       +\mathcal T(\theta),                            \label{eq:ksol-collapse-small-V}
\end{equation}
and $\mathcal T(\theta)\to0$ as $\theta\downarrow0$.
\end{lemma}

\begin{proof}
Write $h(s)=g(-s)$ and $\tau=\theta r^2$.  Since $\Ric\geq0$,
$h(s)\geq h(0)$ as quadratic forms.  The ancient trace Harnack inequality
at vector zero gives $\partial_tR\geq0$.  Hence, for
$x\in B_{g(0)}(p,r)$ and $0\leq s\leq\tau$,
\[
              0\leq R_{h(s)}(x)\leq R_{g(0)}(x)\leq a_3r^{-2}.
\]
It follows from $\partial_s d\mu_{h(s)}=R_{h(s)}d\mu_{h(s)}$ that
\begin{equation}
 d\mu_{h(\tau)}\leq e^{a_3\theta}d\mu_{g(0)}
 \quad\hbox{on }B_{g(0)}(p,r),                         \label{eq:ksol-collapse-measure-comparison}
\end{equation}
and metric monotonicity gives
$B_{h(\tau)}(p,r)\subset B_{g(0)}(p,r)$.

Let $\ell$ be based at $(p,0)$.  The stationary curve at $p$ gives
\begin{equation}
 \ell(p,\tau)
 \leq\frac1{2\sqrt\tau}\int_0^\tau\sqrt s\,a_3r^{-2}\,ds
 =\frac{a_3\theta}{3}.                                 \label{eq:ksol-collapse-pole-ell}
\end{equation}
Apply Lemma~\ref{lem:red-two-point} to $p$ and an arbitrary $q$ at time
$\tau$.  Equations \eqref{eq:ksol-collapse-chi3} and
\eqref{eq:ksol-collapse-pole-ell} give
\begin{equation}
 \ell(q,\tau)\geq
 \chi_3\frac{d_{h(\tau)}(p,q)^2}{\tau}
       -1-\frac{a_3\theta}{3}.                         \label{eq:ksol-collapse-ell-coercive}
\end{equation}

Split the reduced-volume integral into $B_{h(\tau)}(p,r)$ and the annuli
\[
 A_j=B_{h(\tau)}(p,(j+1)r)\setminus B_{h(\tau)}(p,jr),
 \qquad j\geq1.
\]
Because the reduced action has nonnegative integrand, $\ell\geq0$.
Equations \eqref{eq:ksol-collapsed-terminal-ball} and
\eqref{eq:ksol-collapse-measure-comparison} bound the inner integral by
the first term on the right of \eqref{eq:ksol-collapse-small-V}.  On $A_j$,
\eqref{eq:ksol-collapse-ell-coercive} bounds the reduced-volume density by
\[
 (4\pi\tau)^{-3/2}
 e^{1+a_3\theta/3}e^{-\chi_3j^2/\theta}.
\]
Bishop--Gromov at time $-\tau$ gives
$\Vol_{h(\tau)}A_j\leq\omega_3(j+1)^3r^3$.  Summation gives
\eqref{eq:ksol-collapse-small-V}.

For $0<\theta\leq1$,
\[
 e^{-\chi_3j^2/\theta}
 \leq e^{-\chi_3/(2\theta)}e^{-\chi_3j^2/2}.
\]
The series with terms $(j+1)^3e^{-\chi_3j^2/2}$ converges, while
$\theta^{-3/2}e^{-\chi_3/(2\theta)}\to0$.  This proves the last assertion.
\end{proof}

\begin{theorem}[Universal $\kappa_0$-gap]
\label{thm:ksol-universal-kappa-gap}
There is a numerical constant $\kappa_0>0$, depending only on the dimension
and the fixed curvature-norm convention, such that, for every $\kappa>0$,
every three-dimensional ancient $\kappa$-solution which is not a shrinking
round spherical-space-form flow is $\kappa_0$-noncollapsed on all scales.
The conclusion is independent of the input value of $\kappa$.
\end{theorem}

\begin{proof}
Choose $\theta_*\in(0,1]$ so that
$\mathcal T(\theta_*)<1/(4e)$, and set
\begin{equation}
 \kappa_0=
 \frac{(4\pi\theta_*)^{3/2}e^{-a_3\theta_*}}{4e}>0.    \label{eq:ksol-universal-kappa0}
\end{equation}
Fix a time $t_0$, a point $p$, and $r>0$ such that
$|\Rm|\leq r^{-2}$ on $B_{g(t_0)}(p,r)$.  Translate $t_0$ to zero.  If the
ball volume were less than $\kappa_0r^3$, Lemma
\ref{lem:ksol-collapse-small-reduced-volume} would give
\[
 \widetilde V_p(\theta_*r^2)
 <\frac1{4e}+\frac1{4e}<\frac1e,
\]
contrary to Proposition~\ref{prop:ksol-universal-reduced-volume-lower}.
The choices of $t_0,p,r$ were arbitrary.
\end{proof}

\begin{remark}[Why the spherical exception is necessary]
\label{rem:ksol-spherical-gap-exception}
On the round flow on $\Sph^3/\Gamma$, local normalized geometry is
independent of $\Gamma$, while the total volume is divided by $|\Gamma|$.
Finite groups acting freely on $\Sph^3$ have unbounded order.  Thus no
positive noncollapsing constant works simultaneously for all spherical
space forms.
\end{remark}

\begin{theorem}[Universal curvature-normalized precompactness]
\label{thm:ksol-universal-precompactness}
For each $i$, let $(M_i^3,g_i(t))$, $t\leq0$, be an ancient
$\kappa_i$-solution for some $\kappa_i>0$, and assume that it is not a
shrinking spherical-space-form flow.  Choose $(x_i,t_i)$ and put
$Q_i=R_{g_i}(x_i,t_i)$.  A subsequence of
\begin{equation}
 \widehat g_i(s)=Q_i g_i\!\left(t_i+\frac{s}{Q_i}\right),
 \qquad s\leq0,                                        \label{eq:ksol-universal-normalized-sequence}
\end{equation}
converges smoothly and pointedly on every compact spacetime set to a
connected complete ancient $\kappa_0$-solution with base scalar curvature
one.  The limit is allowed to be spherical.  If every $M_i$ is noncompact,
then the limit is noncompact.
\end{theorem}

\begin{proof}
Lemma~\ref{lem:ksol-scalar-positive} gives $Q_i>0$.
Theorem~\ref{thm:ksol-universal-kappa-gap} makes every normalized source
$\kappa_0$-noncollapsed.  Apply
Theorem~\ref{thm:ksol-fixed-kappa-compactness} with this fixed constant.
Nonsphericity is not closed under pointed convergence, so the conclusion
must allow a spherical limit.  Noncompactness is closed here by the
open-and-closed argument in the last paragraph of the fixed-$\kappa$
compactness proof.
\end{proof}

The preceding theorem cannot include all spherical space forms without a
fixed input $\kappa$: a sequence with $|\Gamma_i|\to\infty$ collapses.  This
is why the nonspherical hypothesis occurs in the sources but not in the
conclusion.

\section{Universal curvature-scale derivative estimates}
\label{sec:ksol-universal-derivatives}

\begin{theorem}[All-order curvature-scale estimates]
\label{thm:ksol-universal-derivative-estimates}
For every pair of integers $a,b\geq0$ there is a universal constant
$C_{a,b}<\infty$ such that every three-dimensional ancient
$\kappa$-solution, for every input $\kappa>0$, satisfies
\begin{equation}
 \left|\nabla_t^b\nabla^a\Rm\right|(x,t)
       \leq C_{a,b}R(x,t)^{1+a/2+b}                   \label{eq:ksol-universal-mixed-derivatives}
\end{equation}
at every spacetime point.  Here $\nabla_t$ is the metric-compatible time
derivative in \eqref{eq:shi-metric-time-derivative}; the order in the
display is part of the statement.  Every other fixed word in $\nabla$ and
$\nabla_t$ with
$a$ spatial and $b$ time derivatives satisfies the same scaling power with
a possibly different universal constant.  At the terminal time, every
time derivative is a left derivative.
\end{theorem}

\begin{proof}
Suppose first that the flow is nonspherical.  If
\eqref{eq:ksol-universal-mixed-derivatives} failed for fixed $(a,b)$, there
would be flows and points $(x_i,t_i)$ at which the scale-invariant ratio
tends to infinity.  Put $Q_i=R(x_i,t_i)$ and use
\eqref{eq:ksol-universal-normalized-sequence}.  The scaling identity
\eqref{eq:ksol-mixed-scaling} gives
\begin{equation}
 \left|\nabla_s^b\nabla^a\Rm_{\widehat g_i}\right|(x_i,0)
 =Q_i^{-1-a/2-b}
   \left|\nabla_t^b\nabla^a\Rm_{g_i}\right|(x_i,t_i).  \label{eq:ksol-derivative-normalized-ratio}
\end{equation}
Theorem~\ref{thm:ksol-universal-precompactness} gives a smoothly convergent
subsequence.  The left side converges to the finite corresponding norm of
the smooth limit, contradicting divergence.

For a shrinking spherical space form, lift to its round universal cover.
After scalar normalization the local evolving metric is the single
normalized round-sphere flow, independently of the deck group.  Its mixed
derivatives have universal bounds at every order.  Enlarging the constants
handles this branch.

For any other fixed derivative word, repeat the same contradiction with
that word in place of $\nabla_t^b\nabla^a$.  Under parabolic rescaling each
spatial derivative contributes $Q^{-1/2}$ and each metric-compatible time
derivative contributes $Q^{-1}$, independent of their order.  Pointed
smooth convergence again bounds the normalized value.  This proves the
word assertion directly, without requiring a hidden commutator induction.
Smoothness for $t<0$ and passage from the left gives the endpoint statement.
\end{proof}

\begin{corollary}[Scalar differential estimates and scale comparison]
\label{cor:ksol-scalar-scale-comparison}
There is a universal $\eta\geq1$ such that
\begin{equation}
       |\nabla R|\leq\eta R^{3/2},
       \qquad 0\leq\partial_tR\leq\eta R^2.            \label{eq:ksol-scalar-differential-estimates}
\end{equation}
Consequently, at a fixed time $t$,
\begin{equation}
 d_{g(t)}(x,y)\leq\eta^{-1}R(x,t)^{-1/2}
 \Longrightarrow
 \frac49R(x,t)\leq R(y,t)\leq4R(x,t),                 \label{eq:ksol-spatial-scale-comparison}
\end{equation}
and, whenever
$t-(2\eta R(x,t))^{-1}\leq s\leq t$,
\begin{equation}
                  \frac23R(x,t)\leq R(x,s)\leq R(x,t).\label{eq:ksol-temporal-scale-comparison}
\end{equation}
\end{corollary}

\begin{proof}
Contract the cases $(a,b)=(1,0)$ and $(0,1)$ of
Theorem~\ref{thm:ksol-universal-derivative-estimates}; the lower time bound
is the ancient trace Harnack inequality at vector zero.  Thus
$|\nabla(R^{-1/2})|\leq\eta/2$.  Integrate this along a minimizing
time-$t$ geodesic to obtain
\[
 \frac12R(x,t)^{-1/2}\leq R(y,t)^{-1/2}
       \leq\frac32R(x,t)^{-1/2},
\]
which is \eqref{eq:ksol-spatial-scale-comparison}.  Moreover
$-\eta\leq\partial_t(R^{-1})\leq0$.  Integrating from $s$ to $t$ gives
\[
 R(x,t)^{-1}\leq R(x,s)^{-1}
 \leq R(x,t)^{-1}+\eta(t-s)
 \leq\frac32R(x,t)^{-1},
\]
which proves \eqref{eq:ksol-temporal-scale-comparison}.
\end{proof}

\section{The strong-neck detection interface}
\label{sec:ksol-neck-detection}

Let the scalar-one round cylinder be
\begin{equation}
 g_{\mathrm{cyl}}(s)=2(1-s)g_{\Sph^2(1)}+dz^2,
 \qquad -\infty<s<1.                                  \label{eq:ksol-normalized-cylinder}
\end{equation}
Thus $R_{g_{\mathrm{cyl}}}(\cdot,0)=1$.  For
$0<\varepsilon<1$, put
\begin{equation}
 k_\varepsilon=\lceil\varepsilon^{-1}\rceil,
 \qquad L_\varepsilon=\varepsilon^{-1}.               \label{eq:ksol-neck-parameters}
\end{equation}

\begin{definition}[Strong $\varepsilon$-neck]
\label{def:ksol-strong-neck}
Fix $\varepsilon\in(0,1)$.  Let $(M^3,g(t))$ be a Ricci flow and let
$(x,t)$ satisfy $R(x,t)>0$ and
$[t-R(x,t)^{-1},t]$ lie in the time domain.  Put $Q=R(x,t)$ and
\[
                       g^{(x,t)}(s)=Q\,g(t+s/Q),
                       \qquad -1\leq s\leq0.
\]
The point $(x,t)$ is the center of a \emph{strong $\varepsilon$-neck} if
there is one time-independent embedding
\begin{equation}
 \Phi:\Sph^2\times(-L_\varepsilon-1,L_\varepsilon+1)\longrightarrow M,
 \qquad \Phi(y_*,0)=x,                                 \label{eq:ksol-strong-neck-embedding}
\end{equation}
such that, on
$\Sph^2\times[-L_\varepsilon,L_\varepsilon]\times[-1,0]$,
\begin{equation}
 \max_{\substack{a,b\in\mathbb N\\a+2b\leq k_\varepsilon}}
 \sup\left|
 \nabla_{g_{\mathrm{cyl}}(s)}^a\partial_s^b
 \bigl(\Phi^*g^{(x,t)}(s)-g_{\mathrm{cyl}}(s)\bigr)
 \right|_{g_{\mathrm{cyl}}(s)}
 <\varepsilon.                                        \label{eq:ksol-strong-neck-norm}
\end{equation}
Here $\mathbb N=\{0,1,2,\ldots\}$, so zero spatial or time derivatives are
included.
The extra unit in the spatial domain is a boundary buffer and is not part
of the closeness region.  At the terminal time $s=0$, every
$\partial_s^b$ in \eqref{eq:ksol-strong-neck-norm} is a left derivative.
\end{definition}

The definition fixes the scalar normalization, time orientation, time
interval, embedding, derivative order, background connection, and boundary
buffer.  These data are all needed for a stable formal interface.

\begin{proposition}[Smooth cylindrical convergence detects strong necks]
\label{prop:ksol-cylindrical-convergence-neck}
Suppose curvature-normalized pointed Ricci flows
$(M_i,g_i(s),x_i)$, $s\leq0$, with $R_{g_i}(x_i,0)=1$, converge smoothly
and pointedly on every compact spacetime set to
\[
       (\Sph^2\times\R,g_{\mathrm{cyl}}(s),(y_*,0)).
\]
For every $\varepsilon\in(0,1)$, the points $(x_i,0)$ are centers of strong
$\varepsilon$-necks for all sufficiently large $i$.
\end{proposition}

\begin{proof}
Fix $\varepsilon$ and apply pointed smooth convergence to the compact
spacetime cylinder
\[
 \Sph^2\times[-L_\varepsilon-1,L_\varepsilon+1]\times[-1,0].
\]
For all large $i$, one convergence embedding is defined on its spatial
factor and the pulled-back flows converge in every derivative, hence in the
finite norm \eqref{eq:ksol-strong-neck-norm}.  Restricting the embedding to
the open spatial cylinder gives Definition~\ref{def:ksol-strong-neck}.
\end{proof}

\begin{remark}[The quotient caveat]
\label{rem:ksol-strong-neck-quotients}
The word \emph{unquotiented} is essential.  At the soul of the diagonal
quotient by $(y,z)\mapsto(-y,-z)$ there is no arbitrarily long two-sided
$\Sph^2$-neck, and $\mathbb{RP}^2\times\R$ has the wrong cross-section.
Orientation excludes the latter but not the diagonal quotient.  That
quotient belongs to the cap branch of the next chapter.
\end{remark}

\begin{corollary}[Cylindrical asymptotic shrinkers produce strong necks]
\label{cor:ksol-asymptotic-cylinder-necks}
In the asymptotic-shrinker construction of
Chapter~\ref{ch:reduced-geometry}, suppose the limit is the unquotiented
round cylinder.  Let $\tau_i\to\infty$, put $t_i=-\tau_i$, and let $q_i$
be the selected $\ell$-centers at backward time $\tau_i$.  Then, for every
$\varepsilon\in(0,1)$, $(q_i,t_i)$ is the center of a strong
$\varepsilon$-neck for all sufficiently large $i$.
\end{corollary}

\begin{proof}
The backward rescalings are
$g_i(\theta)=\tau_i^{-1}g(-\tau_i\theta)$.  Smooth convergence at the
centers gives
\[
 a_i:=\tau_iR(q_i,-\tau_i)=R_{g_i(1)}(q_i)\longrightarrow1.
\]
The forward, scalar-normalized flows about $(q_i,t_i)$ are exactly
\begin{equation}
 \widehat g_i(s)
 =R(q_i,t_i)g(t_i+s/R(q_i,t_i))
 =a_i g_i(1-s/a_i).                                   \label{eq:ksol-backward-forward-neck-bridge}
\end{equation}
They converge smoothly to \eqref{eq:ksol-normalized-cylinder}.  Apply
Proposition~\ref{prop:ksol-cylindrical-convergence-neck}.
\end{proof}

\section{Formalization interfaces and dependency forest}
\label{sec:ksol-formalization}

For formalization, the chapter should be read as a sequence of typed
producer--consumer contracts, not as one monolithic compactness proof.
The following are the stable milestone nodes.
\begin{enumerate}
\item \nodeid{PC-M-KSOL-NORMALIZATION} consists of scalar positivity,
      Definition~\ref{def:ksol-curvature-normalization}, and the scaling
      law \eqref{eq:ksol-mixed-scaling}.

\item \nodeid{PC-M-KSOL-TWO-DIMENSIONAL-RIGIDITY} is
      Theorems~\ref{thm:ksol-two-dimensional-classification} and
      \ref{thm:ksol-two-dimensional-KLim-bounded}, with the surface entropy
      identity as a separate fine node.

\item \nodeid{PC-M-KSOL-SPLITTING-BRANCH} is
      Theorem~\ref{thm:ksol-three-dimensional-split-branch}, including the
      exact three deck groups and the translation-collapse argument.

\item \nodeid{PC-M-KSOL-SPATIAL-INFINITY} is
      Theorem~\ref{thm:ksol-ASCR-AVR}; its independent producers are the
      finite-ASCR volume lemma, the later-slice AVR invariance theorem, and
      the curvature point-selection lemma.

\item \nodeid{PC-M-KSOL-ALMOST-ANCIENT-CONTROL} is
      Proposition~\ref{prop:ksol-almost-ancient-collapse}.

\item \nodeid{PC-M-KSOL-POINTED-PRECOMPACTNESS} is the chain
      anchored curvature, seed volume, buffered derivatives, injectivity,
      and Theorem~\ref{thm:ksol-preliminary-compactness}.

\item \nodeid{PC-M-KSOL-HARNACK-CLOSURE} is the two- and
      three-dimensional boundedness upgrade, with
      Lemmas~\ref{lem:ksol-terminal-rank-one-product} and
      \ref{lem:ksol-rank-one-terminal-bounded}, the terminal-time strong
      maximum-principle trichotomy, product-to-factor noncollapsing, the
      Riemannian and spacetime Toponogov line lemmas, and the Busemann
      neck-radius comparison as explicit producers.

\item \nodeid{PC-M-KSOL-FIXED-KAPPA-COMPACTNESS} is
      Theorem~\ref{thm:ksol-fixed-kappa-compactness}.

\item \nodeid{PC-M-KSOL-UNIVERSAL-NONCOLLAPSING} is
      Theorem~\ref{thm:ksol-universal-kappa-gap}; its fine nodes are round
      backward rigidity, asymptotic reduced-volume lower bound, and the
      collapsed-ball reduced-volume upper bound.

\item \nodeid{PC-M-KSOL-UNIVERSAL-DERIVATIVE-BOUNDS} is
      Theorem~\ref{thm:ksol-universal-derivative-estimates} and
      Corollary~\ref{cor:ksol-scalar-scale-comparison}.

\item \nodeid{PC-M-KSOL-CYLINDER-NECK-DETECTION} is
      Definition~\ref{def:ksol-strong-neck}, Proposition
      \ref{prop:ksol-cylindrical-convergence-neck}, and Corollary
      \ref{cor:ksol-asymptotic-cylinder-necks}.
\end{enumerate}

The coarse dependency forest is the following list.
\begin{enumerate}
\item Surface entropy and the asymptotic-shrinker theorem produce
      two-dimensional rigidity.
\item The strong maximum principle and surface rigidity produce the
      rank-one branch.
\item ASCR point selection and dimension reduction produce AVR zero.
\item AVR zero produces almost-ancient collapse, which produces anchored
      curvature control.
\item Anchored curvature and seed volume produce preliminary compactness.
\item The fixed-time trichotomy, factor noncollapsing, and the Riemannian
      line lemma produce rank-one terminal boundedness.
\item Preliminary compactness, rank-one terminal boundedness, the local
      line lemma, and Busemann neck comparison produce Harnack-limit
      closure.
\item Preliminary compactness and Harnack-limit closure produce
      fixed-$\kappa$ compactness.
\item Asymptotic masses and two-point coercivity produce the universal
      $\kappa_0$ gap.
\item The universal $\kappa_0$ gap and fixed-$\kappa$ compactness produce
      universal derivative estimates.
\item Pointed spacetime convergence to the cylinder produces a strong
      neck.
\end{enumerate}

The extended value $+\infty$ in ASCR, finite-input versus universal
constants, source and limit basepoints, time-zero ball capture, strict
radius buffers in noncollapse passage, and left derivatives at a terminal
time are part of these interfaces.  A Lean blueprint should not replace
them with informal coercions or implicit subsequence choices.

The imported mathematical interfaces are: Hopf--Rinow properness;
Toponogov comparison-angle monotonicity; Cheeger--Gromoll splitting and the
positive-curvature soul theorem; smooth Jordan--Schoenflies separation and
its nested-sphere corollary; non-embeddability of $\mathbb{RP}^2$ in
$\R^3$; the curvature-nullspace strong maximum principle; Bishop--Gromov;
the local curvature-plus-volume injectivity estimate;
Chapter~\ref{ch:bando-shi}'s local derivative estimates; Hamilton
compactness with compatible exhaustion embeddings and source-ball capture;
Sharafutdinov retraction for convex Busemann levels; Chapter
\ref{ch:matrix-harnack}'s ancient trace Harnack inequality; and Chapter
\ref{ch:reduced-geometry}'s center, compactness, no-loss, and asymptotic
shrinker theorems.  Each must be matched to an exact blueprint statement.
The collaborator developments of Hamilton compactness and
$\mathcal W$-entropy are not assigned compilation, kernel, axiom-footprint,
or semantic-correspondence status here without a declaration-level audit.
The $\mathcal W$-entropy is not used in this chapter.

\subsection*{Lean status}

Every new node in this chapter has Lean statement and Lean proof status
\emph{not yet implemented}, unless the project status ledger records
contrary declaration-level evidence.  The prose is the intended human
mathematical source; it is not itself typechecking or kernel evidence.

\section{Sources, repairs, and deliberate stopping point}
\label{sec:ksol-sources}

The primary source is G.~Perelman,
\emph{The entropy formula for the Ricci flow and its geometric
applications},
\href{https://arxiv.org/abs/math/0211159}{arXiv:math/0211159}, especially
Section~11.  Hamilton's primary precursor is
\emph{The formation of singularities in the Ricci flow}, especially its
dimension-reduction, spatial-infinity, surface, and neck analysis.  The
detailed classical spine is MSM~163, Chapters~18--20:
the two-dimensional base case, Harnack-limit closure, asymptotic scalar and
volume ratios, almost-ancient collapse, compactness, the universal
noncollapsing gap, derivative estimates, and the Busemann appendices.
The arguments were cross-checked against GSM~77, Chapter~9; MSM~206,
Chapter~28; GSM~235, Section~8.2; and the relevant chapters of MSM~110 and
MSM~144.  Chapter~\ref{ch:reduced-geometry} supplies the corrected
asymptotic-shrinker and no-loss interfaces used here.

The following choices and repairs are mathematically substantive.
\begin{enumerate}
\item The term \emph{$\kappa$-solution} retains the single definition in
      Chapter~\ref{ch:reduced-geometry}.  Lei Ni's useful Harnack closure
      idea appears only as the internal predicate
      $\mathsf{KLim}_{n,\kappa}$; the dimension-two and dimension-three
      boundedness theorems discharge it before any public conclusion.

\item The surface conclusion includes the antipodal quotient
      $\mathbb{RP}^2$.  A formulation concluding only $\Sph^2$ requires
      orientability.  Instead of relying on the long Type~I/Type~II cigar
      route, we combine the already proved shrinker classification with
      Hamilton's exact compact surface entropy identity.

\item A literal use of the usual global dimension-reduction theorem inside
      Harnack-limit closure would assume the curvature boundedness being
      proved.  Lemma~\ref{lem:ksol-dimension-reduction-line} isolates the
      local point-picked hypotheses actually used by the Toponogov proof.

\item The positive-time fixed-slice trichotomy is not cited at the terminal
      time outside its stated domain.  Lemma
      \ref{lem:ksol-terminal-rank-one-product} instead applies the strong
      system theorem on a closed slab and transfers the parallel kernel to
      its terminal endpoint by an explicit one-sided limit.

\item The older sequential almost-ancient statement suppresses the finite
      constants and the available backward time depth.  Proposition
      \ref{prop:ksol-almost-ancient-collapse} states the exact ancient
      quantitative engine consumed here.  Its future surgery-spacetime
      version must add both a spatial buffer and a scaled backward-time
      buffer.

\item The noncollapse threshold is non-strict.  Smooth convergence therefore
      passes it first at radii $r'<\rho<r$ and only then lets $r'\uparrow r$.
      No equality case is passed by an unjustified open condition.

\item Universal $\kappa_0$ is proved directly from the exact reduced masses
      $2/e$ and $1/e$.  Lemma
      \ref{lem:ksol-collapse-small-reduced-volume} keeps the dimensional
      norm bridge between $R$ and $|\Rm|$ and supplies its own Gaussian
      tail.  The spherical exception is necessary, not cosmetic.

\item Nonsphericity is not asserted to be closed under normalized
      convergence.  Theorem~\ref{thm:ksol-universal-precompactness} allows a
      spherical limit.  Noncompactness, by contrast, is proved closed by
      the same-dimensional open-and-closed argument.

\item The derivative theorem uses $\nabla_t$ from
      \eqref{eq:shi-metric-time-derivative} and a fixed operator order.  Raw
      coordinate partial derivatives of an evolving covariant tensor are
      not silently treated as intrinsic tensors.

\item Strong necks are spacetime objects.  Definition
      \ref{def:ksol-strong-neck} fixes one time-independent embedding and a
      parabolic norm; a spatial neck is used only in the Busemann closure
      argument.  The diagonal cylindrical quotient is kept out of the
      unquotiented neck conclusion.
\end{enumerate}

This chapter deliberately stops before the global neck--cap decomposition,
neck stability, the Bryant and ancient-oval branches, and the final
classification of three-dimensional ancient $\kappa$-solutions.  Those
arguments require the uniform package just proved and form the natural next
chapter.  Surgery-spacetime versions of compactness and reduced geometry
will subsequently add survival domains, surgery-time barriers, and cap
crossing estimates; none is silently assumed in the smooth statements here.

\chapter[Qualitative Neck--Cap Structure and Neck Stability]
{Qualitative Neck--Cap Structure and Backward Neck Stability}
\label{ch:neck-cap-stability}
\index{$\kappa$-solution!neck--cap structure}
\index{neck stability}
\index{canonical neighborhood}
\index{cap}
\index{round cylinder}

\section{Purpose, scope, and the model-by-model convention}
\label{sec:ncs-purpose}

This chapter converts the compactness and strong-neck detection package of
Chapter~\ref{ch:kappa-compactness} into the qualitative geometry needed by
Ricci flow with surgery.  There are three outputs.  First, a noncompact
three-dimensional ancient $\kappa$-solution is organized into exact
cylindrical models or a witnessed cap with one cylindrical end.  Second,
compact solutions satisfy a two-cap--tube alternative unless the entire
slice is controlled at one curvature scale.  Third, a sufficiently good
neck remains cylindrical far backward along the \emph{same identity
worldline}.  The last assertion is Kleiner--Lott's backward neck-stability
theorem and is substantially stronger than the existence of convenient
moving neck centers.

No complete isometry classification of ancient $\kappa$-solutions is used.
In particular, the Bryant-soliton uniqueness theorem, ancient-oval
uniqueness, rotational symmetry, constant-mean-curvature neck foliations,
and Brendle's neck-improvement theorem play no role.  This is the
Perelman-level qualitative package needed before the surgery arguments.

An ancient $\kappa$-solution has the meaning fixed in Definition
\ref{def:red-kappa-solution}.  Unless a theorem says otherwise, the
dimension is three and $t\leq0$.  The input $\kappa$ is allowed to vary.
Every noncompact solution is nonspherical, so Theorem
\ref{thm:ksol-universal-kappa-gap} replaces its input constant by the
universal constant $\kappa _0$.  Consequently all constants below can be
chosen universally.  This reduction is made before, not after, applying a
compactness theorem.

The rank-one alternatives from
Theorem~\ref{thm:ksol-three-dimensional-split-branch} must remain separate:
\begin{equation}
 \Sph^2\times\R,\qquad
 \mathbb{RP}^2\times\R,\qquad
 \mathcal C_{\rm diag}:=
 (\Sph^2\times\R)/\langle(y,z)\mapsto(-y,-z)\rangle .
                                                        \label{eq:ncs-three-cylinder-models}
\end{equation}
The first model has two ordinary spherical neck ends.  The second is
nonorientable and has projective cross-sections, so it is not an ordinary
$\Sph^2$-neck.  The third is orientable, one-ended, and has a projective
cap.  Orientation excludes only the middle model.  Treating the last two
quotients as interchangeable is a genuine error.

\begin{roadmap}{Sections~\ref{sec:ncs-witnesses}--
\ref{sec:ncs-cylindrical-end} define witness-bearing objects, prove the
long-range curvature estimate, and identify the geometry at spatial
infinity.  Sections~\ref{sec:ncs-four-alternatives}--
\ref{sec:ncs-compact-structure} produce controlled cores and the global
neck--cap alternatives.  Section~\ref{sec:ncs-canonical-neighborhood}
packages them in the exact form consumed by surgery.  Sections
\ref{sec:ncs-stability-inputs}--\ref{sec:ncs-backward-stability} develop the
reduced-geometric proof of backward neck stability.}
\end{roadmap}

We repeatedly use strict closeness with a reserve.  A target
$\varepsilon$ is proved from $\varepsilon/2$, $\varepsilon/4$, or
$\varepsilon/8$ data.  The fractions have no geometric significance; they
prevent the false inference that a limit of strict
$\varepsilon$-inequalities is still strict at the same threshold.

\section{Witness-bearing necks, caps, and tubes}
\label{sec:ncs-witnesses}

The strong neck of Definition~\ref{def:ksol-strong-neck} is a spacetime
object.  Its terminal slice is a spatial neck witness after a harmless
weakening of the parameter.  All topological gluing in this chapter is
performed on that terminal slice; a spatial product decomposition is never
silently upgraded to parabolic closeness.

\begin{definition}[Good-neck and bad-neck loci]
\label{def:ncs-neck-center-sets}
For $0<\varepsilon<1$ and a time $t$, define
\begin{align}
 \mathscr N_\varepsilon(t)
   &:=\{x\in M:(x,t)\text{ is the center of a strong }
                        \varepsilon\text{-neck}\},             \label{eq:ncs-good-neck-set}\\
 \mathscr B_\varepsilon(t)
   &:=M\setminus\mathscr N_\varepsilon(t).             \label{eq:ncs-bad-neck-set}
\end{align}
Membership in $\mathscr N_\varepsilon(t)$ is existential, but every use of
it below chooses and retains an embedding satisfying
\eqref{eq:ksol-strong-neck-embedding}--
\eqref{eq:ksol-strong-neck-norm}.  The set
$\mathscr B_\varepsilon(t)$ is only a closed set.  It is not asserted to be
a submanifold.
\end{definition}

\begin{lemma}[Openness and buffered recentering]
\label{lem:ncs-neck-openness}
There is $\varepsilon_{\rm rec}>0$ such that, for
$0<\varepsilon\leq\varepsilon_{\rm rec}$ and every ancient
$\kappa$-solution,
\begin{equation}
 \mathscr N_{\varepsilon/4}(t)
       \subset\overline{\mathscr N_{\varepsilon/4}(t)}
       \subset\mathscr N_{\varepsilon/2}(t)
       \subset\mathscr N_\varepsilon(t).              \label{eq:ncs-buffered-neck-closure}
\end{equation}
Moreover every $\mathscr N_\varepsilon(t)$ is open.  In particular,
\begin{equation}
        \partial\mathscr B_{\varepsilon/4}(t)
              \subset\mathscr N_{\varepsilon/2}(t).   \label{eq:ncs-bad-boundary-buffer}
\end{equation}
Every point in the last boundary therefore has an explicitly choosable
$\varepsilon/2$-neck witness, although it need not be an
$\varepsilon/4$-neck center.
\end{lemma}

\begin{proof}
Fix a strong $\varepsilon/4$-neck witness centered at $x$.  Cylindrical
closeness gives a uniform ball about $x$, in scalar units, contained in the
buffered image of the witness.  We first create the small temporal buffer
which recentering needs.  On a spatial subcylinder separated from the
boundary of the chart, cylindrical closeness bounds curvature at the
oldest displayed time.  Hamilton's trace Harnack inequality bounds it on a
uniform additional backward interval, and Theorem
\ref{thm:ksol-universal-derivative-estimates} bounds all spatial and time
derivatives required by the neck norm there.  After integrating
$\partial_tg=-2\Ric$ and weakening $\varepsilon/4$ to
$3\varepsilon/8$, the same spatial embedding is therefore cylindrical on
$[-1-\theta_\varepsilon,0]$ for some
$\theta_\varepsilon>0$ independent of the solution and center.

If $x'$ is sufficiently close to $x$, then $R(x',t)/R(x,t)$ is close to
one by Corollary~\ref{cor:ksol-scalar-scale-comparison}.  Translate the
axial coordinate, rotate the spherical coordinate so that the
distinguished model point maps to $x'$, and restrict the original chart by
half of its spatial buffer.  The normalized interval belonging to
$R(x',t)$ is then contained in $[-1-\theta_\varepsilon,0]$, and pullback
and constant rescaling depend continuously on $x'$.  The finite norm in
\eqref{eq:ksol-strong-neck-norm} is consequently less than
$\varepsilon/2$ after decreasing the neighborhood of $x$.  This proves
openness with a loss of tolerance.

For openness at the same tolerance, start instead with an arbitrary strong
$\varepsilon$-witness and let $m>0$ be the difference between
$\varepsilon$ and its actual strict norm.  Smooth dependence first extends
the chart a witness-dependent amount backward, with norm increase less
than $m/3$.  Recentring at a sufficiently nearby point changes the scalar
normalization, the time interval, and every derivative in the finite norm
by less than another $m/3$.  The recentered norm is therefore still
strictly below $\varepsilon$.  This proves that
$\mathscr N_\varepsilon(t)$ itself is open; no universal radius is claimed
for this same-tolerance assertion.

Now let $x_i\in\mathscr N_{\varepsilon/4}(t)$ tend to $x$.  The preceding
recentring construction can be performed in the witness centered at $x_i$
once $i$ is large, because $d_t(x_i,x)R(x_i,t)^{1/2}\to0$.  It gives a
strong $\varepsilon/2$-neck centered at $x$.  This proves the middle
inclusion in \eqref{eq:ncs-buffered-neck-closure}; the other inclusions are
neck weakening.  Since
$\partial(M\setminus U)\subset\overline U$ for every open $U$, take
$U=\mathscr N_{\varepsilon/4}(t)$ to obtain
\eqref{eq:ncs-bad-boundary-buffer}.
\end{proof}

We next specify the topology that the word \emph{cap} is allowed to carry.
Put
\begin{equation}
 \mathbf B=\overline B^3,\qquad
 \mathbf P=\mathbb{RP}^3\setminus\operatorname{int}\overline B^3.
                                                        \label{eq:ncs-cap-models}
\end{equation}
Both are compact connected smooth three-manifolds with one spherical
boundary component.  Fix smooth parametrizations
$\iota_{\mathbf B}:\Sph^2\to\partial\mathbf B$ and
$\iota_{\mathbf P}:\Sph^2\to\partial\mathbf P$ once and for all.

\begin{definition}[Cap-core and one-cap-end witnesses]
\label{def:ncs-cap-witness}
Fix $t$ and $\varepsilon$.  A \emph{cap-core witness} is the following
tuple.
\begin{enumerate}
\item A compact connected smooth codimension-zero submanifold
      $X\subset M$ and a tag $\mathfrak c\in\{\mathbf B,\mathbf P\}$,
      together with a chosen diffeomorphism $\chi:\mathfrak c\to X$.
\item A point $x\in\operatorname{int}X$ and a chosen strong
      $\varepsilon$-neck witness $\Phi$ centered at a point $y\in\partial X$
      whose terminal-slice spatial embedding is denoted by $\phi_t$.  Its
      center sphere is the defined object
      $S_{\Phi,t}:=\phi_t(\Sph^2\times\{0\})$, and
      \begin{equation}
                     \partial X=S_{\Phi,t}.             \label{eq:ncs-cap-boundary-sphere}
      \end{equation}
\item A collar number $\sigma\in(0,1]$, a choice of sign for the axial
      coordinate, and a boundary reparametrization
      $\beta\in\operatorname{Diff}(\Sph^2)$ such that
      \begin{align}
       \phi_t(\Sph^2\times[-\sigma,0))
                    &\subset\operatorname{int}X,       \label{eq:ncs-cap-inner-collar}\\
       \phi_t(\Sph^2\times(0,\sigma])
                    &\subset M\setminus X,             \label{eq:ncs-cap-outer-collar}\\
       \chi\circ\iota_{\mathfrak c}\circ\beta
                    &=\phi_t(\,\cdot\,,0).            \label{eq:ncs-cap-boundary-parametrization}
      \end{align}
\end{enumerate}
A \emph{one-cap-end witness} additionally contains a diffeomorphism
\begin{equation}
 E:\Sph^2\times[0,\infty)\longrightarrow
                M\setminus\operatorname{int}X              \label{eq:ncs-one-cap-end-map}
\end{equation}
whose boundary parametrization agrees exactly with the recorded
parametrization in \eqref{eq:ncs-cap-boundary-parametrization}, and a proof
that every point in the image of $E$ outside a recorded compact
subcylinder is a strong $\varepsilon$-neck center.  It also stores a
dependent choice function assigning one such strong-neck witness to every
point in that exterior subset.

For $C\geq1$ this witness is \emph{$C$-controlled at $y$} when, with
$Q=R(y,t)$,
\begin{equation}
 \operatorname{diam}_tX\leq C Q^{-1/2},\qquad
 C^{-1}Q\leq R(z,t)\leq CQ\quad(z\in X),               \label{eq:ncs-controlled-cap}
\end{equation}
and
\begin{equation}
                  \Vol_tX\geq C^{-1}Q^{-3/2}.          \label{eq:ncs-controlled-cap-volume}
\end{equation}
The inequalities and all maps are fields of the witness, not consequences
hidden in the noun ``cap.''
\end{definition}

\begin{lemma}[Quantified alignment on a neck overlap]
\label{lem:ncs-neck-overlap-alignment}
There are universal constants $C_{\rm ov}\geq1$ and
$\varepsilon_{\rm ov}>0$ with the following property.  Let
$0<\varepsilon\leq\varepsilon_{\rm ov}$, let $\phi_i$ be the terminal
spatial charts of two chosen strong $\varepsilon$-neck witnesses, and put
\begin{equation}
 Q_i=R(x_i,t),\qquad
 u_i=Q_i^{-1/2}\operatorname{pr}_{\R}\circ\phi_i^{-1}\qquad(i=0,1).
                                                               \label{eq:ncs-physical-axial-functions}
\end{equation}
Let $O'$ be a specified connected compact subdomain of their overlap such
that each inverse image lies at least one axial coordinate unit from the
boundary of the corresponding closeness cylinder, $O'$ contains a full
spherical collar of normalized axial width at least one for each chart,
and every two points of $O'$ can be joined inside $O'$ by a piecewise
smooth curve of $g(t)$-length at most $10Q_0^{-1/2}$.  Then
\begin{equation}
 \left|Q_0/Q_1-1\right|\leq C_{\rm ov}\varepsilon,     \label{eq:ncs-overlap-scale-ratio}
\end{equation}
and there are recorded values $\sigma\in\{-1,1\}$ and $c\in\R$ such that
\begin{equation}
 Q_0^{1/2}\sup_{O'}|u_0-\sigma u_1-c|
 +\sup_{O'}|d(u_0-\sigma u_1)|_{g(t)}
       \leq C_{\rm ov}\varepsilon.                    \label{eq:ncs-overlap-C1-bound}
\end{equation}
If the two collars come with an order, $\sigma$ may be chosen so both
positive axial directions agree with that order.
\end{lemma}

\begin{proof}
On a sufficiently round neck the Ricci tensor has two eigenvalues close to
$1/2$ in scalar normalization and one close to zero.  The least eigenvalue
is therefore simple, and its eigenspace is an intrinsic smooth line field.
At a point of $O'$, both $Q_i^{-1}R(\cdot,t)$ are
$1+O(\varepsilon)$; this gives \eqref{eq:ncs-overlap-scale-ratio}.  The
unit gradients of the two physical axial functions are each
$O(\varepsilon)$-close to the intrinsic least-eigenvalue line.  Choose
$\sigma$ to align them, using the prescribed order when present.  This
gives the derivative term in \eqref{eq:ncs-overlap-C1-bound}.  Subtract the
value at one selected point of $O'$ and integrate along curves of length at
most $10Q_0^{-1/2}$ to obtain its zeroth-order term.  The buffer keeps every
comparison inside the stated neck domains.  No spherical gauge is chosen.
\end{proof}

\begin{definition}[Neck chain, neck tube, and two-cap tube]
\label{def:ncs-neck-chain}
A \emph{finite $\varepsilon$-neck chain} on a time slice consists of an
ordered finite list of chosen terminal-slice neck charts
$\phi_0,\ldots,\phi_m$, closed truncation intervals in their axial
coordinates, a compact connected codimension-zero swept region $W$, and
the following certified data.
\begin{enumerate}
\item There are disjoint extreme center spheres $S_-,S_+$ with
      $\partial W=S_-\sqcup S_+$.  The specified truncated chart images
      cover $W$ and contain full collars of both boundary spheres.
\item The overlap nerve of the truncated images is the linear tree
      $0-1-\cdots-m$.  For each consecutive pair, a compact overlap domain
      $O'_j$ satisfying the three explicit domain conditions of Lemma
      \ref{lem:ncs-neck-overlap-alignment} is recorded.  Nonconsecutive
      truncated images are disjoint.  These fields certify no return and
      absence of lateral boundary.
\item For every $j<m$, the chain stores
      $\sigma_j\in\{-1,1\}$ and $c_j\in\R$ satisfying
      \eqref{eq:ncs-overlap-C1-bound} on $O'_j$.  The signs are compatible
      with the order from $S_-$ to $S_+$.
\end{enumerate}
Thus the region, truncations, linear-nerve proof, overlap domains, signs,
translation constants, and inequalities are fields of the chain.
Boundary-sphere parametrizations remain arbitrary diffeomorphisms; no
global $O(3)$ gauge is inferred from local metric closeness.  Pairwise
intersection by itself is not a neck-chain witness.

A \emph{neck-tube witness} is a compact codimension-zero region $T$ and a
chosen diffeomorphism
\begin{equation}
                 \Psi:\Sph^2\times[0,1]\longrightarrow T              \label{eq:ncs-neck-tube-map}
\end{equation}
together with a finite strong $\varepsilon$-neck cover and explicit
matching of the two boundary parametrizations.  A \emph{two-cap--tube
witness} consists of two cap cores $X_-,X_+$ with disjoint interiors and a
neck tube $T$ such that
\begin{equation}
 M=X_-\cup T\cup X_+,\quad
 X_-\cap X_+=\varnothing,\quad
 X_\pm\cap T=\partial X_\pm.                           \label{eq:ncs-two-cap-union}
\end{equation}
Arbitrary center spheres of overlapping neck charts need not be disjoint;
the ordered-chain data are what produce the tube.
\end{definition}

\begin{proposition}[An ordered finite neck chain is an annulus]
\label{prop:ncs-neck-chain-annulus}
After decreasing $\varepsilon_{\rm ov}$, every finite ordered
$\varepsilon$-neck chain with
$0<\varepsilon\leq\varepsilon_{\rm ov}$ produces a neck-tube witness
between its two recorded boundary spheres.  In particular, its specified
swept region is diffeomorphic to $\Sph^2\times[0,1]$.
\end{proposition}

\begin{proof}
Use the stored signs and constants to regard the physical axial functions
as consistently increasing local functions on $W$.  Choose a subordinate
finite partition of unity $\{\chi_j\}$ adapted to the recorded unit-width
overlap collars and put $u=\sum_j\chi_j u_j$.  The terms $u_jd\chi_j$
occur only through differences $(u_j-u_k)d\chi_j$, because
$\sum_jd\chi_j=0$.  Equation \eqref{eq:ncs-overlap-C1-bound} and the
unit-width collar bound make $du$ nonzero and positive on the aligned
least-Ricci eigenline when $C_{\rm ov}\varepsilon_{\rm ov}$ is chosen
sufficiently small.

The linear nerve, full cross-sectional overlaps, and no-return field imply
that $u$ has no lateral boundary and takes the two boundary components of
$W$ to its two extreme values.  Hence
$u:W\to[u_-,u_+]$ is a proper submersion.  Its fibers are connected
two-spheres: this is true in the first chart and fiber type cannot change
under a proper submersion.  Ehresmann triviality gives
$W\diffeo\Sph^2\times[0,1]$.  Collars and isotopy extension make its two
boundary parametrizations agree with the recorded ones, producing the
neck-tube witness.
\end{proof}

\section{Long-range curvature transfer}
\label{sec:ncs-long-range}

The assertion $\ASCR=+\infty$ in Theorem~\ref{thm:ksol-ASCR-AVR} is a
limsup assertion.  It produces good sequences but says nothing by itself
about an arbitrary escaping sequence.  The next theorem is the missing
uniform implication.

\begin{theorem}[Universal long-range curvature transfer]
\label{thm:ncs-long-range-curvature}
There is a nondecreasing function
$\alpha:[0,\infty)\to[0,\infty)$, depending only on the dimension and the
norm convention, such that
\begin{equation}
                  \lim_{s\to\infty}\alpha(s)=+\infty                 \label{eq:ncs-alpha-proper}
\end{equation}
and every three-dimensional ancient $\kappa$-solution satisfies, at every
fixed time $t$ and for all $x,y\in M$,
\begin{equation}
 R(y,t)d_t(x,y)^2
       \geq\alpha\!\left(R(x,t)d_t(x,y)^2\right).      \label{eq:ncs-long-range-transfer}
\end{equation}
Equivalently, for every $0\leq A<\infty$ there is a universal $B(A)<\infty$
such that
\begin{equation}
 R(x,t)d_t(x,y)^2\geq B(A)
       \Longrightarrow R(y,t)d_t(x,y)^2>A.             \label{eq:ncs-long-range-threshold}
\end{equation}
\end{theorem}

\begin{proof}
We first prove the threshold statement.  If it failed for some $A$, there
would be solutions, times, and points $x_i,y_i$ such that, writing
$d_i=d_{t_i}(x_i,y_i)$,
\begin{equation}
 R(x_i,t_i)d_i^2\longrightarrow\infty,\qquad
 R(y_i,t_i)d_i^2\leq A.                                \label{eq:ncs-transfer-contradiction}
\end{equation}
Round spherical-space-form flows have spatially constant scalar curvature
and cannot furnish this sequence.  All remaining sources are universally
$\kappa_0$-noncollapsed by
Theorem~\ref{thm:ksol-universal-kappa-gap}.

Normalize at $(y_i,t_i)$ by $Q_i=R(y_i,t_i)$.  Then
$d_{Q_ig_i(t_i)}(x_i,y_i)\leq\sqrt A$.  Theorem
\ref{thm:ksol-universal-precompactness}, with source-ball capture, gives a
pointed smooth subsequential limit on a ball of radius $\sqrt A+1$.
Scalar curvature is uniformly bounded on that captured ball.  Hence
\begin{equation}
 \frac{R(x_i,t_i)}{R(y_i,t_i)}\leq C(A).                \label{eq:ncs-transfer-local-ratio}
\end{equation}
Multiplying by $R(y_i,t_i)d_i^2\leq A$ contradicts the first part of
\eqref{eq:ncs-transfer-contradiction}.

For a definition of $\alpha$ involving only integers and finite choices,
choose $B_0=0$ and recursively choose
\begin{equation}
 B_m\geq\max\{m,B_{m-1}+1,B(m)\}\qquad(m\geq1),        \label{eq:ncs-Bm-choice}
\end{equation}
where $B(m)$ is a threshold for $A=m$.  Set
\begin{equation}
 \alpha(s)=\max\bigl(\{m\in\mathbb N:B_m\leq s\}\cup\{0\}\bigr).
                                                               \label{eq:ncs-alpha-discrete}
\end{equation}
Then $\alpha$ is nondecreasing and proper.  If the argument of $\alpha$ in
\eqref{eq:ncs-long-range-transfer} is at least $B_m$, the threshold
statement gives that its left side is greater than $m$.  Taking the largest
such $m$ proves the displayed inequality.  This discrete modulus avoids a
choice of generalized inverse over a class of flows.
\end{proof}

\begin{corollary}[Curvature scale along every escaping sequence]
\label{cor:ncs-escape-curvature-scale}
Fix a noncompact ancient $\kappa$-solution, a time $t$, and a point $p$.
If $x_i$ escapes every compact set and $Q_i=R(x_i,t)$, then
\begin{equation}
             Q_i d_t(p,x_i)^2\longrightarrow\infty.    \label{eq:ncs-escape-QD}
\end{equation}
Consequently $d_t(p,x_i)\sqrt{Q_i}\to\infty$.
\end{corollary}

\begin{proof}
Hopf--Rinow gives $d_t(p,x_i)\to\infty$.  Apply
\eqref{eq:ncs-long-range-transfer} with $x=p$ and $y=x_i$.  Its argument
$R(p,t)d_t(p,x_i)^2$ tends to infinity, and then
\eqref{eq:ncs-alpha-proper} proves \eqref{eq:ncs-escape-QD}.
\end{proof}

\section{Cylindrical geometry at spatial infinity}
\label{sec:ncs-cylindrical-end}

\begin{theorem}[Every escaping normalization in the one-cap branches]
\label{thm:ncs-arbitrary-escape-cylinder}
Fix a time $t$ and let $x_i$ escape every compact subset of a noncompact
three-dimensional ancient $\kappa$-solution.  Put $Q_i=R(x_i,t)$ and
\begin{equation}
 g_i(s)=Q_i g(t+s/Q_i),\qquad s\leq0.                  \label{eq:ncs-escaping-normalization}
\end{equation}
The following model-by-model conclusions hold.
\begin{enumerate}
\item If the curvature operator is positive, every subsequence has a
      further subsequence for which $(M,g_i(s),x_i)$ converges smoothly and
      pointedly to the unquotiented scalar-one round cylinder
      \eqref{eq:ksol-normalized-cylinder}.
\item The same conclusion holds on the diagonal quotient
      $\mathcal C_{\rm diag}$.
\item On the exact product $\Sph^2\times\R$ the normalized flow is already
      that cylinder.  On $\mathbb{RP}^2\times\R$ it is instead the exact
      projective cylinder and does not converge to an ordinary
      $\Sph^2$-neck.
\end{enumerate}
Consequently, in the positive-curvature and diagonal branches, for every
$0<\varepsilon<1$ there is a compact set $K_\varepsilon(t)$ such that
\begin{equation}
              M\setminus K_\varepsilon(t)
                    \subset\mathscr N_\varepsilon(t). \label{eq:ncs-eventual-strong-necks}
\end{equation}
\end{theorem}

\begin{proof}
Consider first the positive-curvature branch and fix $p\in M$.  Theorem
\ref{thm:ksol-universal-precompactness} gives a pointed subsequential limit
of \eqref{eq:ncs-escaping-normalization}.  By
Corollary~\ref{cor:ncs-escape-curvature-scale},
\begin{equation}
        d_t(p,x_i)\sqrt{Q_i}\longrightarrow\infty.      \label{eq:ncs-pole-escapes-normalized}
\end{equation}
The Riemannian argument of
Lemma~\ref{lem:ksol-Riemannian-line-limit} therefore gives a line through
the limiting basepoint.  Cheeger--Gromoll splitting and the whole-flow
curvature-nullspace theorem split the limiting flow as a surface times a
line.  Theorem~\ref{thm:ksol-two-dimensional-classification} makes its
universal cover the shrinking round cylinder.

The positive-curvature soul theorem gives $M\diffeo\R^3$.  A nontrivial
rank-one quotient contains a compact embedded $\mathbb{RP}^2$ in a bounded
pointed ball: it is a factor slice in the projective product and the central
one-sided cross-section in the diagonal quotient.  A convergence embedding
would put that surface into $\R^3$, contradicting the non-embeddability of
$\mathbb{RP}^2$ in $\R^3$.  Thus the limit is unquotiented.  Its base scalar
curvature is one, so its extinction parameter is fixed and it is exactly
\eqref{eq:ksol-normalized-cylinder}.

On $\mathcal C_{\rm diag}$, lift a representative of $x_i$ to
$(y_i,z_i)\in\Sph^2\times\R$.  Escape is equivalent to
$|z_i|\to\infty$.  The distance, in scalar units, from $(y_i,z_i)$ to its
nontrivial deck translate tends to infinity.  Hence the covering map is
injective on every fixed normalized ball about $x_i$ for all large $i$,
and the pointed limit is the unquotiented cylinder.  The two product cases
are immediate from their exact formulas.

If \eqref{eq:ncs-eventual-strong-necks} failed, an escaping sequence of
points outside $\mathscr N_\varepsilon(t)$ would have a cylindrical
subsequence.  Proposition~\ref{prop:ksol-cylindrical-convergence-neck}
would put all sufficiently late terms in
$\mathscr N_\varepsilon(t)$, a contradiction.
\end{proof}

\begin{lemma}[An aligned neck sphere on an escaping ray]
\label{lem:ncs-aligned-ray-cross-section}
There are universal constants $\varepsilon_{\rm ray}>0$ and
$C_{\rm ray}<\infty$ with the following property.  Let
$(N^3,g(s))$ be a noncompact ancient $\kappa$-solution in the positive,
diagonal-quotient, or unquotiented-cylinder branch.  Let
$h=g(t)$ and
$\gamma:[0,\infty)\to N$ be a unit-speed minimizing ray.  For every
$0<\varepsilon\leq\varepsilon_{\rm ray}$ there is
$s_0=s_0(N,h,\gamma,\varepsilon)<\infty$ such that, for every $s\geq s_0$,
$y=\gamma(s)$ has a chosen strong $\varepsilon$-neck at $(y,t)$ whose
terminal center sphere $\Sigma_s$ satisfies
\begin{equation}
 \Sigma_s\subset B_h\bigl(y,C_{\rm ray}R_h(y)^{-1/2}\bigr),
 \qquad \gamma^{-1}(\Sigma_s)=\{s\},                 \label{eq:ncs-aligned-ray-sphere}
\end{equation}
and the intersection is transverse.  If $\Sigma_s$ separates $N$, the
initial and terminal tails of $\gamma$ lie on its two different sides.
In the two exact rank-one branches allowed here, $\Sigma_s$ is chosen as
the standard factor sphere, or as its image in the diagonal quotient.
\end{lemma}

\begin{proof}
If the eventual assertion failed, choose bad parameters $s_i\to\infty$,
put $y_i=\gamma(s_i)$, and set
$Q_i=R_h(y_i)$.  Corollary~\ref{cor:ncs-escape-curvature-scale} gives
$s_i\sqrt{Q_i}\to\infty$.  Theorem
\ref{thm:ncs-arbitrary-escape-cylinder} and strong-neck detection show
that the $Q_i$-normalized pointed slices converge to the round cylinder.
The backward and forward pieces of $\gamma$ through $y_i$ converge to a
complete line.  Every line in $\Sph^2(\sqrt2)\times\mathbb R$ is axial,
because a geodesic with nonzero spherical speed ceases to minimize after
a bounded length.  Choose the detecting product charts so that this line
is the model axis.  Their center spheres then contain $y_i$, have diameter
at most $C_{\rm ray}Q_i^{-1/2}$, and meet $\gamma$ transversely there.
In the exact unquotiented cylinder take the literal factor slice through
$y_i$.  In the diagonal quotient take the image of
$\Sph^2\times\{c_i\}$ through a lift of $y_i$ with $c_i\neq0$.  For large
$i$ this is an embedded ordinary sphere and bounds the standard compact
region $\{|z|\leq|c_i|\}/\mathbb Z_2$.

There can be no second intersection.  If $\gamma(s_i')$ were another
point of the same center sphere, minimizing gives
\begin{equation}
 |s_i'-s_i|=d_h\bigl(\gamma(s_i'),\gamma(s_i)\bigr)
       \leq\operatorname{diam}_h\Sigma_{s_i}
       \leq C_{\rm ray}Q_i^{-1/2}.                   \label{eq:ncs-ray-sphere-no-reentry-bound}
\end{equation}
Thus the normalized intervening segment stays in a fixed subneck, where
the axial coordinate has derivative bounded away from zero for all large
$i$.  It cannot meet the zero level twice.  This constructs the asserted
witness for every sufficiently late term of a further subsequence,
contradicting that all $s_i$ were bad.  The last statement follows from
connectedness of the two ray subarcs and the single transverse crossing.
\end{proof}

\begin{lemma}[All-neck rigidity and nonempty core]
\label{lem:ncs-all-neck-rigidity}
There is a universal $\varepsilon_{\rm top}>0$ such that the following
hold for $0<\varepsilon\leq\varepsilon_{\rm top}$.
\begin{enumerate}
\item If an orientable ancient $\kappa$-solution satisfies
      $\mathscr B_\varepsilon(t)=\varnothing$ on one slice, then it is the
      exact unquotiented shrinking cylinder.
\item On $\mathbb{RP}^2\times\R$ no point is the center of an ordinary
      strong $\varepsilon$-neck.
\item In the positive-curvature and diagonal branches,
      $\mathscr B_\varepsilon(t)$ is nonempty and compact.
\end{enumerate}
\end{lemma}

\begin{proof}
Choose $\varepsilon_{\rm top}$ so small that the least Ricci eigenvalue on
every strong $\varepsilon_{\rm top}$-neck core is simple and its eigenline
is transverse to every central sphere.  We first exclude the
positive-curvature branch from the all-neck hypothesis.

If the positive branch is noncompact, the soul theorem gives
$M\diffeo\mathbb R^3$.  Choose one central neck sphere $S$ and let $K$ be
its compact side.  Smooth Schoenflies gives $K\diffeo\overline B^3$.  The
least-Ricci eigenspaces define a line bundle on $K$; it is trivial because
$K$ is contractible.  Orient a unit section to point outward along the
connected boundary.  It is nowhere zero, contradicting relative
Poincar\'e--Hopf because $\chi(K)=1$.

If the positive branch is compact, pass to its finite simply connected
universal cover and lift one central neck sphere.  The lifted sphere
separates.  Van Kampen applied to the two sides shows that either closed
side $K$ is simply connected, so the lifted least-Ricci line bundle on
$K$ is trivial.  Orient its unit section outward along the boundary.  For
every compact odd-dimensional manifold with boundary,
$2\chi(K)=\chi(\partial K)$; hence $\chi(K)=1$.  Relative
Poincar\'e--Hopf again contradicts the nowhere-zero outward section.

It remains to inspect the exact rank-one models.  The unquotiented round
cylinder is an exact ordinary neck at every point.  At scalar curvature
one,
\begin{equation}
 \operatorname{inj}(\Sph^2\times\mathbb R)=\pi\sqrt2,
 \qquad
 \operatorname{inj}(\mathbb{RP}^2\times\mathbb R)=\pi/\sqrt2. \label{eq:ncs-projective-injectivity-gap}
\end{equation}
A sufficiently strong ordinary spherical-neck chart contains the fixed
metric ball which detects these two injectivity radii and forces the
basepoint injectivity radius to be close to $\pi\sqrt2$.  This is
impossible in the projective product, proving assertion~(2).  The same
smaller injectivity radius occurs at every point of the central
$\mathbb{RP}^2$ in the diagonal quotient, so that surface lies in
$\mathscr B_\varepsilon(t)$ after one universal decrease of
$\varepsilon_{\rm top}$.

The split/positive dichotomy now proves assertion~(1): orientation removes
the projective product, the two positive branches have just been excluded,
and the diagonal model has a nonneck central surface.  Only the
unquotiented cylinder remains.  Finally, compactness of the bad set is
automatic in the compact positive branch and follows from
\eqref{eq:ncs-eventual-strong-necks} in the noncompact positive and
diagonal branches.  The preceding arguments prove nonemptiness in every
one of them, which is assertion~(3).
\end{proof}

\section{The remote-three-point lemma and four alternatives}
\label{sec:ncs-four-alternatives}

The next result is the quantitative localization engine.  It is stated for
the strict strong-neck convention used in this book.  The source version
uses a slightly different neck convention; buffered detection and
Theorem~\ref{thm:ksol-universal-derivative-estimates} make the two forms
equivalent after replacing $\varepsilon$ by $\varepsilon/8$.

The source proof is often phrased using the Tits boundary.  For later
formalization we isolate and prove its finite geometric content instead;
no ideal boundary is needed.

\begin{lemma}[Remote cylindrical separator]
\label{lem:ncs-remote-cylindrical-separator}
Let $(N^3,g(s))$ be an orientable noncompact ancient
$\kappa$-solution and put $h=g(t)$.  Suppose that $a\geq0$, that
$\lambda:[-a,\infty)\to N$ is a unit-speed minimizing geodesic, and that
$\eta:[0,\infty)\to N$ is a unit-speed minimizing ray with
$\lambda(0)=\eta(0)=o$.  Assume
\begin{equation}
 d_h(\eta(s),o)=d_h\bigl(\eta(s),\lambda([-a,\infty))\bigr)
                         \qquad(s\geq0).                \label{eq:ncs-nearest-ray-projection}
\end{equation}
Then $N$ has two ends and its ancient flow is the exact unquotiented round
cylinder.
\end{lemma}

\begin{proof}
Choose $r_j\to\infty$, put $u_j=\lambda(r_j)$, and set
$Q_j=R_h(u_j)$.  Since $R_h(o)r_j^2\to\infty$, long-range curvature
transfer gives
\begin{equation}
                              Q_jr_j^2\longrightarrow\infty.  \label{eq:ncs-separator-scale-remote}
\end{equation}
Apply Lemma~\ref{lem:ncs-aligned-ray-cross-section} to the forward ray
$\lambda|_{[0,\infty)}$.  It supplies, for all large $j$, a full central
sphere $\Sigma_j$ such that
\begin{equation}
 \Sigma_j\subset B_h(u_j,C_{\rm ray}Q_j^{-1/2})        \label{eq:ncs-separator-sphere-size}
\end{equation}
and $\lambda$ meets $\Sigma_j$ transversely exactly once, passing from its
negative to its positive side.

For every $s\geq0$, the nearest-point identity and the triangle inequality
give, respectively,
\begin{equation}
 d_h(u_j,\eta(s))\geq s,
 \qquad d_h(u_j,\eta(s))\geq |r_j-s|.                 \label{eq:ncs-separator-two-distance-bounds}
\end{equation}
Consequently
\begin{equation}
                       d_h(u_j,\eta(s))\geq r_j/2      \label{eq:ncs-separator-ray-avoidance}
\end{equation}
for every $s\geq0$.  Equations
\eqref{eq:ncs-separator-scale-remote} and
\eqref{eq:ncs-separator-sphere-size} show that $\eta$ is disjoint from
$\Sigma_j$ for large $j$.  Since $o$ lies on the negative side and $\eta$
starts at $o$, the whole $\eta$-ray remains on that side; the forward tail
of $\lambda$ lies on the positive side.

It remains only to make ``side'' global.  The noncompact orientable model
dichotomy proved in Chapter~\ref{ch:kappa-compactness} leaves the positive
branch $N\diffeo\mathbb R^3$, the diagonal quotient, and the exact
unquotiented cylinder.  In the first branch smooth Schoenflies and the
exterior product theorem make $\Sigma_j$ separating; in the last two
branches Lemma~\ref{lem:ncs-aligned-ray-cross-section} chose the standard
factor sphere or the image of
$\Sph^2\times\{c_j\}$, which separates directly in the explicit product
or product-quotient model.
Thus the two ray tails lie in distinct unbounded components of
$N\setminus\Sigma_j$, and $N$ has at least two ends.  A complete manifold
with $\Ric\geq0$ has
at most two ends, and in the two-ended case contains a line.
Cheeger--Gromoll splitting, the curvature-nullspace strong maximum
principle, and the two-dimensional classification identify the whole flow
with the unquotiented shrinking round cylinder.
\end{proof}

\begin{lemma}[Remote nonneck three-point localization]
\label{lem:ncs-remote-three-point}
There are $\varepsilon_3>0$ and, for every
$0<\varepsilon\leq\varepsilon_3$, a universal $A=A(\varepsilon)\geq1$
with the following property.  Assume that the flow is not the exact
projective cylinder $\mathbb{RP}^2\times\R$, and let
$x,y\in\mathscr B_\varepsilon(t)$ satisfy
\begin{equation}
                  R(x,t)d_t(x,y)^2>A.                 \label{eq:ncs-remote-pair}
\end{equation}
Choose a minimizing time-$t$ segment $[xy]$.  For every $z\in M$, at least
one of the following holds:
\begin{align}
 R(x,t)d_t(x,z)^2&<A,                                  \label{eq:ncs-three-point-x}\\
 R(y,t)d_t(y,z)^2&<A,                                  \label{eq:ncs-three-point-y}\\
 z&\in\mathscr N_\varepsilon(t)\quad\hbox{and}\quad
 R(z,t)d_t(z,[xy])^2<A.                               \label{eq:ncs-three-point-middle}
\end{align}
\end{lemma}

\begin{proof}
Suppose the result fails for $A_i\to\infty$.  Choose violating triples
$x_i,y_i,z_i$ and violating times $t_i$.  Translate every $t_i$ to zero
and suppress the time from the notation.  A shrinking round
spherical-space-form slice satisfies
\begin{equation}
                 R\,\operatorname{diam}^2\leq6\pi^2. \label{eq:ncs-round-remote-bound}
\end{equation}
Thus the remote-pair quantity tending to infinity excludes round sources
for all large $i$.  The universal $\kappa_0$ gap and pointed compactness
are therefore available.  Choose a minimizing segment
$\gamma_i=[x_iy_i]$, and a point
$z_i'\in\gamma_i$ minimizing the distance from $z_i$ to $\gamma_i$.
There are two separate assertions to prove.

First, $z_i'$ diverges from both endpoints in the curvature scale of that
endpoint.  If, for example,
$R(x_i)d(x_i,z_i')^2$ stayed bounded, normalize at $x_i$.  The segment from
$x_i$ toward $y_i$ has unbounded normalized length, and
$[z_i'z_i]$ also has unbounded normalized length because the violating
triple is remote from $x_i$.  After passing to a pointed limit, the first
segments give a minimizing geodesic
$\lambda:[-a,\infty)\to M_\infty$ with
$\lambda(-a)=x_\infty$ and $\lambda(0)=z'_\infty$; the second segments
give a ray $\eta$ from $z'_\infty$.  More explicitly, with
$\widetilde d_i=R(x_i,0)^{1/2}d_0$, boundedness of
$\widetilde d_i(x_i,z_i')$ and failure of the first and remote-pair
clauses give the two triangle inequalities
\begin{equation}
 \widetilde d_i(z_i',z_i)
   \geq\widetilde d_i(x_i,z_i)-\widetilde d_i(x_i,z_i')
       \longrightarrow\infty,
 \qquad
 \widetilde d_i(z_i',y_i)
   \geq\widetilde d_i(x_i,y_i)-\widetilde d_i(x_i,z_i')
       \longrightarrow\infty.                        \label{eq:ncs-two-rays-have-infinite-length}
\end{equation}
Thus both limiting rays are defined for every finite parameter.

We now verify, rather than assume, persistence of the nearest point.  For
fixed $s\geq0$, let $q_i(s)\in[z_i'z_i]$ have
$\widetilde d_i(z_i',q_i(s))=s$.  If some
$w_i\in[x_iy_i]$ satisfied
$\widetilde d_i(q_i(s),w_i)<s$, then
\begin{equation}
 \widetilde d_i(z_i,w_i)
 \leq \widetilde d_i(z_i,q_i(s))+
      \widetilde d_i(q_i(s),w_i)
 <\widetilde d_i(z_i,z_i'),                           \label{eq:ncs-nearest-point-subsegment-contradiction}
\end{equation}
contradicting the definition of $z_i'$.  Hence $z_i'$ is nearest to every
$q_i(s)$ on the entire segment $[x_iy_i]$, and in fact
\begin{equation}
 \widetilde d_i\bigl(q_i(s),[x_iy_i]\bigr)
      =\widetilde d_i(q_i(s),z_i')=s.                 \label{eq:ncs-nearest-point-persists}
\end{equation}
For every fixed $r\in[-a,\infty)$ choose points
$w_i(r)\in[x_iy_i]$ converging to $\lambda(r)$.  Equation
\eqref{eq:ncs-nearest-point-persists} gives
$d(\eta(s),\lambda(r))\geq s$, while $r=0$ gives equality.  Hence the
limit satisfies exactly
\eqref{eq:ncs-nearest-ray-projection} in the limit.

The flows in the sequence are orientable: after excluding the projective
product, this follows from the rank-one/positive dichotomy, the soul
theorem in the noncompact positive branch, and Synge's theorem in the
compact positive branch.  The pointed limit is noncompact because it
contains a ray.  Pointed smooth limits of
orientable manifolds are orientable.  Indeed, a compact
orientation-reversing loop together with a tubular neighborhood in the
limit would embed in every sufficiently large approximant.  Lemma
\ref{lem:ncs-remote-cylindrical-separator} therefore makes the pointed
limit the unquotiented cylinder.  Its limiting basepoint is a
cylinder center, so strict cylindrical convergence and Proposition
\ref{prop:ksol-cylindrical-convergence-neck} make $x_i$ an
$\varepsilon$-neck center for large $i$, contrary to
$x_i\in\mathscr B_\varepsilon(0)$.  For the other endpoint, long-range
transfer first gives
\begin{equation}
                   R(y_i,0)d_0(x_i,y_i)^2\longrightarrow\infty.  \label{eq:ncs-y-endpoint-remote-scale}
\end{equation}
Failure of the second alternative also gives
$R(y_i,0)d_0(y_i,z_i)^2\to\infty$.  The preceding argument, now normalized
at $y_i$, is therefore literal and proves divergence from that endpoint.

Second, apply long-range transfer twice.  The endpoint divergence just
proved implies
\begin{equation}
 R(z_i',0)d_0(z_i',x_i)^2\to\infty,\qquad
 R(z_i',0)d_0(z_i',y_i)^2\to\infty.                    \label{eq:ncs-zprime-two-long-sides}
\end{equation}
Normalize at $z_i'$.  The two halves of $\gamma_i$ converge to a line, so
the limit splits.  It is orientable by the preceding compact-neighborhood
argument, so the surface classification identifies it with the
unquotiented round cylinder rather than the projective product.  Endpoint
divergence puts $z_i'$ in the interior of $\gamma_i$.  First variation for
the nearest point shows that the initial tangent of $[z_i'z_i]$ is
orthogonal to $\gamma_i$ at $z_i'$.  This passes to orthogonality with the
limiting line.  Their normalized lengths must remain bounded:
otherwise they would converge to a minimizing ray orthogonal to the
$\R$-factor, whereas such a geodesic in
$\Sph^2\times\R$ remains in a compact spherical factor and ceases to
minimize after length $\pi\sqrt2$.  Thus
\begin{equation}
 R(z_i',0)d_0(z_i,z_i')^2\leq C.                       \label{eq:ncs-normal-distance-zprime}
\end{equation}
Smooth cylindrical convergence, scalar comparison over this bounded
distance, and the $\varepsilon/8$ reserve show that $z_i$ is an
$\varepsilon$-neck center and that
$R(z_i,0)d_0(z_i,[x_iy_i])^2$ is uniformly bounded.  This contradicts the
chosen violation of \eqref{eq:ncs-three-point-middle}.

The separating sphere used above is not justified by Jordan separation for
an arbitrary embedded sphere.  It is the model-specific global separator
constructed in Lemma~\ref{lem:ncs-remote-cylindrical-separator}.  The
orientability clause is essential: the lemma is false on
$\mathbb{RP}^2\times\R$.
\end{proof}

\begin{lemma}[Uniform axial exit from a remote terminal neck]
\label{lem:ncs-remote-axial-exit}
There are universal numbers $\varepsilon_{\rm ax}>0$ and
$L_{\rm ax}<\infty$ with the following property.  Let $(N,h)$ be complete
with $\sec_h\geq0$, let $z$ be the center of a chosen spatial
$\varepsilon_{\rm ax}$-neck with terminal chart $\phi$, center sphere
$\Sigma=\phi(\Sph^2\times\{0\})$, and distinguished model point
$z=\phi(y_*,0)$.  Assume that
$\Sigma$ separates $N$, and put $Q=R_h(z)$.  If $x\in N\setminus\Sigma$
and
\begin{equation}
                         Qd_h(x,z)^2\geq L_{\rm ax},    \label{eq:ncs-remote-axial-exit-hypothesis}
\end{equation}
then one may choose the sign so that $x$ lies in the negative component of
$N\setminus\Sigma$, and there is a number $s_+\in(0,1]$ such that the
recorded point $z^+=\phi(y_*,s_+)$ satisfies
\begin{equation}
                              d_h(x,z^+)>d_h(x,z).      \label{eq:ncs-remote-axial-exit-conclusion}
\end{equation}
The chosen sign and the point $z^+$ are outputs.
\end{lemma}

\begin{proof}
If no pair of constants existed, set
$\varepsilon_i=\min\{1/i,1/100\}$ and $L_i=i$.  For every $i$ choose a
counterexample with a spatial $\varepsilon_i$-neck and
$R(z_i)d(x_i,z_i)^2\geq L_i$.  Thus the neck tolerance tends to zero and
the left side of \eqref{eq:ncs-remote-axial-exit-hypothesis} tends to
infinity.  Choose the
axial sign so that $x_i$ lies in the negative component.  Normalize at
$z_i$.  The final portions of minimizing segments from $x_i$ to $z_i$
converge on the round product to a minimizing ray ending at the origin.
Such a ray has zero spherical speed: a nonzero spherical component
produces a conjugate point after a universally bounded length in
$\Sph^2(\sqrt2)$.  Its final direction cannot be the wrong axial
direction.  Otherwise the segment, starting in the negative component,
would cross $\Sigma_i$, enter the positive collar, and return to
$z_i\in\Sigma_i$.  The two crossings are at bounded normalized distance,
but in the limiting product chart the axial coordinate along the limiting
geodesic has derivative bounded away from zero and cannot vanish twice.
Thus every terminal unit tangent $v_i$ of every minimizing segment from
$x_i$ to $z_i$ satisfies, for the physical axial function
$u_i=R(z_i)^{-1/2}\operatorname{pr}_{\mathbb R}\circ\phi_i^{-1}$,
\begin{equation}
                              du_i(v_i)\geq\frac12                 \label{eq:ncs-axial-exit-terminal-tangent}
\end{equation}
for all large $i$.  Indeed, failure would supply a subsequence with a
different limiting terminal direction, which was just excluded.  The
one-sided first-variation formula at a cut point is
\begin{equation}
 D^+_v d_h(x,\cdot)\big|_z
   =\min_{\gamma\in\operatorname{Min}(x,z)}
          \langle v,\dot\gamma(d_h(x,z))\rangle .     \label{eq:ncs-distance-directional-derivative}
\end{equation}
Apply it to the positive axial tangent.  Equation
\eqref{eq:ncs-axial-exit-terminal-tangent}, smooth cylindrical control,
and compactness of the set of terminal minimizing directions give a
positive right derivative uniformly over all minimizing segments.
Therefore some $s_i^+\in(0,1]$ gives $z_i^+=\phi_i(y_*,s_i^+)$ with
\eqref{eq:ncs-remote-axial-exit-conclusion}, contradicting the
counterexample.  This diagonal contradiction supplies finite
$\varepsilon_{\rm ax}$ and $L_{\rm ax}$.  The separation hypothesis is
essential: a long product $\Sph^2\times\Sph^1$ gives a counterexample
without it.
Increase $L_{\rm ax}$ once so that it exceeds the square of the normalized
diameter of every center sphere in a spatial
$\varepsilon_{\rm ax}$-neck.
\end{proof}

\begin{lemma}[A farthest point beyond a nonneck core]
\label{lem:ncs-farthest-point-nonneck}
There are universal constants $\varepsilon_{\rm far}>0$ and, for every
$0<\varepsilon\leq\varepsilon_{\rm far}$, $A_{\rm far}(\varepsilon)<\infty$
with the following property.  Let a compact three-dimensional ancient
$\kappa$-solution be given, let
$x\in\mathscr B_\varepsilon(t)$, and let $z$ maximize $d_t(x,\cdot)$.  If
\begin{equation}
             R(x,t)d_t(x,z)^2>A_{\rm far}(\varepsilon),       \label{eq:ncs-farthest-remote}
\end{equation}
then $z\in\mathscr B_\varepsilon(t)$.
\end{lemma}

\begin{proof}
Choose $\varepsilon_{\rm far}\leq\varepsilon_{\rm ax}$ below the universal
$C^2$ curvature-spectrum gap which prevents a scalar-one
constant-curvature three-metric from being a strong spherical neck.  If the conclusion
failed, then $z\in\mathscr N_\varepsilon(t)$.  Weaken its strong
$\varepsilon$-neck to a strong $\varepsilon_{\rm far}$-neck and restrict
its terminal slice to the spatial witness required by Lemma
\ref{lem:ncs-remote-axial-exit}.  Its center sphere separates.  Indeed, a
compact ancient $\kappa$-solution admitting such a sufficiently good neck
is nonround.  The split-branch theorem excludes a compact rank-one
$\kappa$-solution, and the strong maximum principle therefore makes its
sectional curvature positive.  Synge makes it orientable, and Bonnet--Myers makes its
fundamental group finite.  A nonseparating oriented sphere would define a
nonzero integral first-cohomology class, which is impossible for a finite
fundamental group.  When
$R(x,t)d_t(x,z)^2>A_{\rm far}$, choose
$A_{\rm far}$ so that
$\alpha(A_{\rm far})\geq L_{\rm ax}$.  Long-range transfer gives
$R(z,t)d_t(x,z)^2\geq L_{\rm ax}$.  By the final choice of
$L_{\rm ax}$, this also implies $x\notin\Sigma$.  The lemma then supplies
$z^+$ with $d_t(x,z^+)>d_t(x,z)$, contradicting the farthest-point
property.  Hence $z$ is not a strong $\varepsilon$-neck center.
\end{proof}

\begin{theorem}[Global four-alternative localization]
\label{thm:ncs-four-alternatives}
After decreasing $\varepsilon_3$, for every
$0<\varepsilon\leq\varepsilon_3$ there is a universal finite number
$A_{\rm four}(\varepsilon)\geq1$ such that the following holds.  Put
$A=A_{\rm four}(\varepsilon)$.  At each fixed time, every
three-dimensional ancient $\kappa$-solution lies in one of the following
model cases, selected in the displayed priority order.
\begin{enumerate}
\item[(P)] The flow is the exact projective cylinder
      $\mathbb{RP}^2\times\R$.
\item[(A)] The flow is the exact unquotiented cylinder and
      $\mathscr B_\varepsilon(t)=\varnothing$.
\item[(B)] The manifold is noncompact,
      $\mathscr B_\varepsilon(t)$ is nonempty and compact, and for every
      ordered pair $x,y\in\mathscr B_\varepsilon(t)$,
      \begin{equation}
                         R(x,t)d_t(x,y)^2\leq A.        \label{eq:ncs-alternative-B}
      \end{equation}
\item[(C)] The manifold is compact and there are
      $x_-,x_+\in\mathscr B_\varepsilon(t)$ and a chosen minimizing segment
      $\gamma=[x_-x_+]$ such that
      \begin{equation}
           R(x_-,t)d_t(x_-,x_+)^2>A,                 \label{eq:ncs-C-remote-pair}
      \end{equation}
      and the bad set lies in two endpoint balls,
      \begin{equation}
       \mathscr B_\varepsilon(t)\subset
       B_t(x_-,A R(x_-,t)^{-1/2})
       \cup B_t(x_+,A R(x_+,t)^{-1/2}),                \label{eq:ncs-two-bad-clusters}
      \end{equation}
      while every $z\in\mathscr N_\varepsilon(t)$ outside those balls
      satisfies
      \begin{equation}
            R(z,t)d_t(z,\gamma)^2\leq A.               \label{eq:ncs-neck-close-to-axis}
      \end{equation}
\item[(D)] The manifold is compact and there is
      $x\in\mathscr B_\varepsilon(t)$ such that
      \begin{equation}
                         R(x,t)d_t(x,z)^2\leq A
                         \quad(z\in M).                 \label{eq:ncs-controlled-component-distance}
      \end{equation}
\end{enumerate}
The diagonal quotient occurs in (B).  Exact round spherical-space-form
flows occur in (D), with constants for distance and curvature but not with
a universal total-volume lower bound.
\end{theorem}

\begin{proof}
Let $A_3(\varepsilon)$ be the threshold supplied by Lemma
\ref{lem:ncs-remote-three-point}.  Choose
\begin{equation}
 A_{\rm four}(\varepsilon)
   \geq\max\{A_3(\varepsilon),A_{\rm far}(\varepsilon),
                    6\pi^2,1\}.                              \label{eq:ncs-four-constant-choice}
\end{equation}
This specifies the quantifier order without an instruction to enlarge a
constant during an application.

Remove the exact projective product first.  Every remaining model is
orientable: the two remaining rank-one models are orientable, a
noncompact positive-curvature branch is diffeomorphic to $\R^3$, and a
compact positive-curvature branch is orientable by Synge's theorem.  Thus,
if the bad set is empty, Lemma~\ref{lem:ncs-all-neck-rigidity} gives (A).
Suppose it is nonempty.
If the slice has constant positive sectional curvature, then
$R\operatorname{diam}^2\leq6\pi^2\leq A$ by
\eqref{eq:ncs-round-remote-bound}, so any chosen bad point satisfies
alternative~(D).  Thus round spherical-space-form flows have the placement
asserted in the theorem.
If no ordered pair of bad points violates
\eqref{eq:ncs-alternative-B}, then either the manifold is noncompact and
(B) holds, using Theorem~\ref{thm:ncs-arbitrary-escape-cylinder} for
compactness of the bad set, or the manifold is compact.  In the latter
      case choose $x$ in the bad set and let $z$ maximize $d_t(x,\cdot)$.
By the definition \eqref{eq:ncs-four-constant-choice},
$A\geq A_{\rm far}(\varepsilon)$.  If
$R(x,t)d_t(x,z)^2>A$, Lemma
\ref{lem:ncs-farthest-point-nonneck} makes $z$ a bad point and hence a
remote ordered partner of $x$, a contradiction.  Thus every point is
within the bound in \eqref{eq:ncs-controlled-component-distance}, which is
(D).

It remains to consider a remote bad pair.  Apply
Lemma~\ref{lem:ncs-remote-three-point} with its threshold $A_3$.  Its first two clauses put all bad
points into the two endpoint balls, giving
\eqref{eq:ncs-two-bad-clusters}.  Its third clause puts every remaining
point within a controlled curvature-scale distance of $\gamma$; since
$A\geq A_3$, this is \eqref{eq:ncs-neck-close-to-axis}.

These bounds force compactness.  Otherwise choose $z_j$ escaping every
compact set.  The endpoint clauses eventually fail, so the third clause of
Lemma~\ref{lem:ncs-remote-three-point} gives
\begin{equation}
                  R(z_j,t)d_t(z_j,\gamma)^2\leq A.    \label{eq:ncs-remote-pair-escape-bound}
\end{equation}
Choose $z'_j\in\gamma$ nearest to $z_j$.  The fixed segment $\gamma$ is
compact, hence $\min_\gamma R(\cdot,t)>0$, while
$d_t(z_j,\gamma)\to\infty$.  Consequently
$R(z'_j,t)d_t(z'_j,z_j)^2\to\infty$.  Long-range transfer with
$x=z'_j$ and $y=z_j$ contradicts
\eqref{eq:ncs-remote-pair-escape-bound}.  Thus the remote-pair case is
(C), with \eqref{eq:ncs-C-remote-pair} retained as part of its witness.
The priority convention makes the cases exhaustive without claiming that
their raw predicates are pairwise disjoint.
\end{proof}

\section{Controlled cores and the finite neck-chain theorem}
\label{sec:ncs-controlled-core}

\begin{proposition}[Quantitative control of the nonneck core]
\label{prop:ncs-bad-set-compact}
There is $\varepsilon_{\rm core}>0$ such that the following holds for every
$0<\varepsilon\leq\varepsilon_{\rm core}$.  Let a noncompact ancient
$\kappa$-solution satisfy alternative~(B) of Theorem
\ref{thm:ncs-four-alternatives} at tolerance $\varepsilon/4$; explicitly,
$\mathscr B_{\varepsilon/4}(t)$ is nonempty and compact and its ordered
pairs satisfy \eqref{eq:ncs-alternative-B} with
$A_4=A_{\rm four}(\varepsilon/4)$.  There is a constant
$C_0=C_0(\varepsilon)$ such that, for every
\begin{equation}
       p\in\partial\mathscr B_{\varepsilon/4}(t),
       \qquad Q=R(p,t).                                 \label{eq:ncs-buffered-boundary-point}
\end{equation}
the point $p$ has a chosen strong $\varepsilon/2$-neck witness and
\begin{align}
 \operatorname{diam}_t\mathscr B_{\varepsilon/4}(t)
       &\leq C_0Q^{-1/2},                               \label{eq:ncs-bad-diameter}\\
 C_0^{-1}Q\leq R(z,t)&\leq C_0Q
       \quad\bigl(z\in\mathscr B_{\varepsilon/4}(t)\bigr).  \label{eq:ncs-bad-scalar-comparison}
\end{align}
No smoothness of the set-theoretic boundary is asserted.
\end{proposition}

\begin{proof}
Choose $\varepsilon_{\rm core}$ below the openness, topology, and
four-alternative thresholds, after accounting for the factor $1/4$.
Let $\eta$ be the constant in
\eqref{eq:ksol-scalar-differential-estimates}.  Set
\begin{equation}
                         c_\eta=(4\eta^2)^{-1}.         \label{eq:ncs-scalar-ratio-constant}
\end{equation}
The set is nonempty, compact, and proper by
Lemma~\ref{lem:ncs-all-neck-rigidity} and
Theorem~\ref{thm:ncs-arbitrary-escape-cylinder}; connectedness of $M$
makes its boundary nonempty.  Lemma~\ref{lem:ncs-neck-openness} supplies
the stated witness at $p$.  Alternative (B) at tolerance
$\varepsilon/4$, applied in both orders, gives
\begin{equation}
 Qd_t(p,z)^2\leq A_4,\qquad
 R(z,t)d_t(p,z)^2\leq A_4                              \label{eq:ncs-two-ordered-bad-bounds}
\end{equation}
for every bad $z$.  The first inequality gives
\eqref{eq:ncs-bad-diameter}.

We prove the curvature comparison rather than appealing to a picture of a
cap.  Let $H=R(z,t)/Q$ and $d=d_t(p,z)$.  If $H\geq4$, the Lipschitz
bound for $R^{-1/2}$ gives $Qd^2\geq c_\eta$.  Thus
\begin{equation}
 R(z,t)d^2=H Qd^2\geq c_\eta H.
\end{equation}
Apply \eqref{eq:ncs-long-range-transfer} with $x=z,y=p$ and combine it
with $Qd^2\leq A_4$.  Properness of $\alpha$ bounds $H$ from above in terms
of $A_4$ and $\eta$.

If $H\leq1/4$, the same Lipschitz inequality, now with the two endpoints
reversed, gives
\begin{equation}
 Qd^2\geq c_\eta H^{-1},\qquad R(z,t)d^2\geq c_\eta.
\end{equation}
Apply long-range transfer with $x=p,y=z$.  The second inequality in
\eqref{eq:ncs-two-ordered-bad-bounds} forces
$A_4\geq\alpha(c_\eta H^{-1})$, which bounds $H$ away from zero.
The interval $1/4\leq H\leq4$ needs no argument.  Combining the three
cases gives \eqref{eq:ncs-bad-scalar-comparison}.
\end{proof}

\section{Global noncompact neck--cap structure}
\label{sec:ncs-global-structure}

\begin{lemma}[The diagonal quotient is a projective cap]
\label{lem:ncs-diagonal-projective-cap}
The map
\begin{equation}
 \Theta:[y,s]\longmapsto[y_1:y_2:y_3:s]                \label{eq:ncs-diagonal-projective-map}
\end{equation}
is a diffeomorphism
\begin{equation}
 \mathcal C_{\rm diag}\diffeo
       \mathbb{RP}^3\setminus\{[0:0:0:1]\}.            \label{eq:ncs-diagonal-punctured-RP3}
\end{equation}
For every $a>0$, the image of
$(\Sph^2\times[-a,a])/\mathbb Z_2$ is diffeomorphic to
$\mathbf P$, and
\begin{equation}
 M\setminus\operatorname{int}X_a\diffeo\Sph^2\times[0,\infty),
 \qquad M\setminus X_a\diffeo\Sph^2\times(0,\infty), \label{eq:ncs-diagonal-complements}
\end{equation}
where $X_a=(\Sph^2\times[-a,a])/\mathbb Z_2$.
\end{lemma}

\begin{proof}
The homogeneous point in \eqref{eq:ncs-diagonal-projective-map} is
unchanged when $(y,s)$ is replaced by $(-y,-s)$.  Conversely, normalize the
first three homogeneous coordinates of a point different from
$[0:0:0:1]$ to have Euclidean norm one.  The two choices differ by exactly
the diagonal involution.  These constructions are smooth inverses.

Under this identification, $|s|\leq a$ is the complement of a small open
projective ball about the missing point; hence it is $\mathbf P$.  The two
components $s\geq a$ and $s\leq-a$ are interchanged by the involution, and
their quotient is one copy of $\Sph^2\times[0,\infty)$.  Boundary
parametrizations are supplied by the quotient map.  Notice that the
involution reverses orientation on both factors and therefore preserves
the product orientation.
\end{proof}

\begin{theorem}[Noncompact qualitative neck--cap theorem]
\label{thm:ncs-noncompact-neck-cap}
There is $\varepsilon_{\rm nc}>0$ such that, for every
$0<\varepsilon\leq\varepsilon_{\rm nc}$, the following holds.  At every time slice
of a noncompact three-dimensional ancient $\kappa$-solution exactly one of
the following geometric branches occurs.
\begin{enumerate}
\item The flow is the exact unquotiented round cylinder.  Every point is a
      strong $\varepsilon$-neck center and there are two ends.
\item The flow is the exact projective cylinder
      $\mathbb{RP}^2\times\R$.  It is recorded as a separate two-ended
      branch, not as an ordinary spherical neck.
\item The flow is the diagonal quotient.  There is a one-cap-end witness
      of kind $\mathbf P$.
\item The curvature operator is positive.  Then $M\diffeo\R^3$ and there
      is a one-cap-end witness of kind $\mathbf B$.
\end{enumerate}
In the last two cases the witness can be chosen so that
\begin{equation}
       \mathscr B_{\varepsilon/4}(t)\subset
                     \operatorname{int}X,\qquad
       M\setminus\operatorname{int}X
                      \subset\mathscr N_\varepsilon(t),        \label{eq:ncs-cap-contains-bad-set}
\end{equation}
and its boundary is the center sphere of a chosen strong
$\varepsilon/2$-neck.
\end{theorem}

\begin{proof}
Theorem~\ref{thm:ksol-three-dimensional-split-branch} gives the exact three
rank-one branches and the positive-curvature branch.  The first two items
are immediate.  Lemma~\ref{lem:ncs-diagonal-projective-cap} supplies the
topology in the third.  Scalar curvature is spatially constant on this
model; denote its time-$t$ value by $Q(t)$ and choose
\begin{equation}
                       a=L(\varepsilon)Q(t)^{-1/2},    \label{eq:ncs-diagonal-cutoff-scale}
\end{equation}
where $L(\varepsilon)$ is large enough that the nontrivial deck translate
of every point with $|s|\geq a$ lies beyond the buffered neck domain in
$Q(t)$-units.  Equations \eqref{eq:ncs-diagonal-complements} and the explicit
quotient chart then supply the projective one-cap-end witness.  Notice that
a time-independent axial cutoff would not have the required normalized
meaning.

It remains to construct the positive-curvature cap.  The soul theorem and
strict positivity give $M\diffeo\R^3$.  Proposition
\ref{prop:ncs-bad-set-compact} puts the buffered bad set in a scalar-scale
ball about its boundary point $p$ and compares scalar curvature throughout
that set with $Q=R(p,t)$.  Take a connected compact regular neighborhood
$X_0$ of the bad set in a slightly larger $Q^{-1/2}$-ball.

Choose a minimizing ray from a point of $X_0$.  By
Lemma~\ref{lem:ncs-aligned-ray-cross-section}, its sufficiently remote
points have buffered neck witnesses whose center spheres meet the ray
exactly once.  Corollary~\ref{cor:ncs-escape-curvature-scale} makes the
diameter of such a sphere little-oh of its distance from the initial
point.  Choose one center sphere $S$ disjoint from $X_0$.  Smooth
Schoenflies separates $M\diffeo\mathbb R^3$ into a compact side and an
unbounded side.  The ray tail is in the unbounded side, because it is
unbounded, while its initial point is on the other side by the unique
transverse crossing.  Since $X_0$ is connected and disjoint from $S$, it
lies entirely in that compact side; call the latter $X$.  Choose the axial
sign so the negative collar faces $X$ and the positive collar faces the
ray tail.

Smooth Schoenflies gives $X\diffeo\mathbf B$ and the exterior product
$M\setminus\operatorname{int}X\diffeo\Sph^2\times[0,\infty)$.  Collars
and isotopy extension make its boundary parametrization equal to the
recorded neck parametrization.  Because the entire bad set lies inside
$X$, every exterior point is a strong $\varepsilon/4$-neck center and hence
a strong $\varepsilon$-neck center.  This is the required one-cap-end
witness.  No finite-chain proposition is used to trivialize an infinite
end.
\end{proof}

\begin{remark}[A useful topological alternative to no-backtracking]
\label{rem:ncs-cap-Poincare-Hopf}
In the positive branch, suppose an outer neck sphere bounded a ball
containing no nonneck point.  The simple least-Ricci-eigenline would extend
over the ball and be transverse to its boundary.  Since the ball is
orientable, orient the line and reverse it if necessary to point outward.
Poincar\'e--Hopf would give a nowhere-zero outward vector field on a ball,
contradicting $\chi(B^3)=1$.  This is another finite proof that an outer
neck must enclose the cap core.  The proof above uses the ordered-ray
witness because it also supplies the end parametrization.
\end{remark}

\section{Compact slices: two caps or one controlled component}
\label{sec:ncs-compact-structure}

\begin{lemma}[Two-sided axial crossing in a deep neck]
\label{lem:ncs-two-sided-axial-crossing}
There are universal constants
$\varepsilon_{\rm cross}\in(0,1/100)$ and
$L_{\rm cross}<\infty$ with the following property.  Let $z$ lie on a
unit-speed minimizing segment $\gamma=[ab]$ in a time slice of a
three-dimensional ancient $\kappa$-solution.  Suppose $(z,t)$ is the
center of a chosen strong $\varepsilon_{\rm cross}$-neck, with terminal
map $\phi$, distinguished model point $z=\phi(y_*,0)$, scalar
$Q=R(z,t)$, and physical axial function
\begin{equation}
 u=Q^{-1/2}\operatorname{pr}_{\mathbb R}\circ\phi^{-1}.       \label{eq:ncs-crossing-physical-axis}
\end{equation}
If
\begin{equation}
 Qd_t(z,a)^2\geq L_{\rm cross},
 \qquad Qd_t(z,b)^2\geq L_{\rm cross},                       \label{eq:ncs-two-sided-crossing-depth}
\end{equation}
then one can choose the sign of $u$ so that the following finite data are
certified.
\begin{enumerate}
\item The component of
      $\gamma\cap\phi(\Sph^2\times[-10,10])$ containing $z$ is a single
      segment, and $u\circ\gamma$ has derivative at least $1/2$ on it.
\item This component crosses the recorded negative and positive collars
      \begin{equation}
       \phi(\Sph^2\times[-4,-3]),\qquad
       \phi(\Sph^2\times[3,4])                         \label{eq:ncs-fixed-crossing-collars}
      \end{equation}
      in that order.
\item For every $s\in[-4,4]$, the whole minimizing segment $\gamma$
      meets the slice $\phi(\Sph^2\times\{s\})$ exactly once.
\end{enumerate}
The chosen sign and the restricted component are outputs of the lemma.
\end{lemma}

\begin{proof}
Suppose that no pair $(\varepsilon_{\rm cross},L_{\rm cross})$ works.
For $i\geq100$ put $\varepsilon_i=1/i$ and $L_i=i$.  Choose a
counterexample carrying a strong $\varepsilon_i$-neck and satisfying both
depth inequalities with $L_i$.  Normalize by $Q_i=R(z_i,t_i)$ and
translate $t_i$ to zero.  The two halves of the minimizing segments have
lengths tending to infinity, so on every bounded arclength interval they
subconverge in $C^1$ to a complete minimizing line through the center of
\begin{equation}
                       \Sph^2(\sqrt2)\times\mathbb R.          \label{eq:ncs-crossing-limit-cylinder}
\end{equation}
A complete minimizing line in this product has zero spherical speed: a
nonzero spherical component acquires a conjugate point in bounded length.
Thus it is an axial line.  Choose its increasing axial sign.  Smooth
convergence gives
\begin{equation}
                 \frac{d}{ds}(u_i\circ\gamma_i)\geq\frac12    \label{eq:ncs-crossing-positive-derivative}
\end{equation}
on every portion of $\gamma_i$ lying in the fixed slab
$\phi_i(\Sph^2\times[-10,10])$, for all large $i$.
Here no change-of-scale factor is hidden: differentiation of the physical
coordinate $Q_i^{-1/2}\operatorname{pr}_{\mathbb R}$ with respect to
original arclength equals differentiation of
$\operatorname{pr}_{\mathbb R}$ with respect to $Q_i$-normalized
arclength.  Moreover, any point of the segment in this slab is at distance
from $z_i$ at most the intrinsic diameter of the slightly enlarged
cylindrical slab.  Since every subsegment of $\gamma_i$ is minimizing,
all such portions occur in one bounded arclength interval and are covered
by the preceding $C^1$ convergence.

For completeness, this also rules out a disconnected return to the slab.
If the segment met one slice with coordinate in $[-4,4]$ twice, the
subsegment between the two hits would be minimizing and its length would
not exceed the intrinsic diameter of that slice.  In scalar-one cylinder
units these diameters are uniformly less than $5$ for large $i$.
Consequently the intervening subsegment stays in the slab $[-10,10]$,
where \eqref{eq:ncs-crossing-positive-derivative} makes a repeated axial
value impossible.  The long two sides ensure that the segment reaches
both collars in \eqref{eq:ncs-fixed-crossing-collars}.  Hence every large
$i$ satisfies all three conclusions, contradicting its selection.  This
produces finite universal constants.  Finally decrease
$\varepsilon_{\rm cross}$ so that
\begin{equation}
 \varepsilon_{\rm cross}\leq\varepsilon_{\rm ov},
 \qquad C_{\rm ov}\varepsilon_{\rm cross}<10^{-3};   \label{eq:ncs-crossing-overlap-reserve}
\end{equation}
the conclusions persist under this decrease, and the second inequality
is the graph reserve used in the continuation argument below.
\end{proof}

\begin{lemma}[Remote endpoint caps with a strict buffer]
\label{lem:ncs-remote-endpoint-caps}
There is $\varepsilon_{\rm end}>0$ such that, for every
$0<\varepsilon\leq\varepsilon_{\rm end}$ and every $0\leq H<\infty$, there is
$L=L(\varepsilon,H)<\infty$ with the following property.  Suppose a compact
nonround three-dimensional ancient $\kappa$-solution has points
$x_-,x_+\in\mathscr B_{\varepsilon/8}(t)$ such that
\begin{equation}
 R(x_\pm,t)d_t(x_-,x_+)^2>L,                           \label{eq:ncs-two-sided-remote-endpoints}
\end{equation}
Put
\begin{equation}
 K_\pm=\mathscr B_{\varepsilon/8}(t)\cap
       \overline B_t\bigl(x_\pm,H R(x_\pm,t)^{-1/2}\bigr)       \label{eq:ncs-endpoint-bad-clusters}
\end{equation}
and assume
$\mathscr B_{\varepsilon/8}(t)\subset K_-\cup K_+$.
Then there are cap-core witnesses $X_-,X_+$ with
\begin{equation}
                              X_-\cap X_+=\varnothing,          \label{eq:ncs-endcaps-disjoint}
\end{equation}
whose distinguished points are $x_-,x_+$,
whose boundary necks have tolerance $\varepsilon/2$, and which contain the
respective bad clusters in their interiors:
\begin{equation}
                         K_\pm\subset\operatorname{int}X_\pm.  \label{eq:ncs-clusters-inside-endcaps}
\end{equation}
The chosen segment $[x_-x_+]$ leaves $X_-$
through its positive collar, crosses the middle region, and enters $X_+$
through its positive collar with the latter axial sign reversed.
If $y_\pm$ denotes the distinguished center of the boundary-neck witness,
then the exact crossing output is
\begin{equation}
 \{y_-\}=[x_-x_+]\cap\partial X_-,\qquad
 \{y_+\}=[x_-x_+]\cap\partial X_+.                   \label{eq:ncs-endcap-distinguished-crossings}
\end{equation}
Moreover, every
$z\in[x_-x_+]\setminus(\operatorname{int}X_-\cup
\operatorname{int}X_+)$ satisfies the two explicit axial-depth bounds
\begin{equation}
 R(z,t)d_t(z,x_-)^2\geq H,
 \qquad R(z,t)d_t(z,x_+)^2\geq H.                     \label{eq:ncs-endcap-axial-depth}
\end{equation}
\end{lemma}

\begin{proof}
Choose $\varepsilon_{\rm end}>0$ below the noncompact neck--cap,
aligned-ray, two-sided-crossing, and buffered-detection thresholds, after
accounting for every tolerance down to $\varepsilon/64$.  Fix once and for
all a threshold function $B_{\rm lr}$ in
\eqref{eq:ncs-long-range-threshold}, and put
\begin{equation}
 D_H=10+\max\!\left\{H,\sqrt{B_{\rm lr}(H)},
                    \sqrt{B_{\rm lr}(L_{\rm cross})}\right\}.          \label{eq:ncs-endcap-capture-radius}
\end{equation}

Suppose that no such $L$ exists.  For every integer $i$ choose a
counterexample for which both quantities on the left of
\eqref{eq:ncs-two-sided-remote-endpoints} exceed $i$.  Normalize first at
$x_{-,i}$ and then at $x_{+,i}$.  Universal $\kappa_0$-compactness gives
two noncompact pointed limits.  In each normalization the chosen segment,
oriented away from the basepoint, converges on every compact interval to a
unit-speed ray $\gamma_\pm$.

The strict reserve is essential.  A strong $\varepsilon/16$-neck at a
limiting basepoint would pull back to a strong $\varepsilon/8$-neck,
contradicting the choice of the endpoint.  Thus each limiting basepoint is
in the buffered bad core.  Theorem
\ref{thm:ncs-noncompact-neck-cap}, at tolerance $\varepsilon/64$, gives a
limiting ball or projective cap containing it.  The exact unquotiented
cylinder is excluded by the nonneck basepoint.  The exact projective
cylinder is excluded because a sufficiently large captured product domain
contains a two-sided $\mathbb{RP}^2$, whereas every compact nonround source
is positively curved and orientable.

Do not choose an arbitrary boundary of either limiting cap.  Apply Lemma
\ref{lem:ncs-aligned-ray-cross-section} to $\gamma_\pm$ and choose a
center sphere so far along the ray that its compact side contains
\begin{equation}
             \overline B(x_{\pm,\infty},2D_H)          \label{eq:ncs-limit-endcap-large-ball}
\end{equation}
strictly in its interior, as well as the original buffered bad core.  The
lemma records a unique transverse ray crossing and its outward axial sign.
The center is chosen still farther if necessary so that long-range
transfer makes its basepoint-side depth at least $L_{\rm cross}$.

Source-ball capture now pulls back a finite tuple: the compact cap side,
its topology tag, its distinguished endpoint, its buffered boundary-neck
chart, the exact boundary parametrization, and the signed collar.  The
pulled-back limit sphere need not pass exactly through the corresponding
point of the source minimizing segment.  Recenter the chart there, use the
full product collar to isotope the old boundary to the new center sphere,
and redefine the cap by adjoining or removing the compact collar between
the two spheres.  This finite collar isotopy preserves the topology tag,
the strict ball containment, and the boundary parametrization after
composition with its recorded isotopy.  Strict convergence weakens the
recentered neck tolerance to $\varepsilon/2$.  For all large $i$ the
resulting cap $X_{\pm,i}$
satisfies the source-level inclusion
\begin{equation}
 \overline B_{t_i}\!\left(x_{\pm,i},
          D_H R(x_{\pm,i},t_i)^{-1/2}\right)
       \subset\operatorname{int}X_{\pm,i}.            \label{eq:ncs-source-endcap-large-ball}
\end{equation}
Each limiting cap is contained in some finite ball about its basepoint.
Both normalized endpoint distances tend to infinity; hence the triangle
inequality gives, for all large $i$,
$X_{-,i}\cap X_{+,i}=\varnothing$.  This proves the closed-set
disjointness in \eqref{eq:ncs-endcaps-disjoint}, not merely disjointness of
interiors.

At either pulled-back boundary center, the normalized distance to its own
endpoint is at least $B_{\rm lr}(L_{\rm cross})$ in that endpoint scale;
long-range transfer gives the first depth required by Lemma
\ref{lem:ncs-two-sided-axial-crossing}.  The distance to the opposite
endpoint tends to infinity in the opposite endpoint scale, and a second
application of long-range transfer gives the other depth.  That lemma and
the pulled-back sign therefore show that the source segment crosses each
boundary sphere exactly once: it exits $X_-$ through the positive collar
and enters $X_+$ through the oppositely signed positive collar.  In
particular, the portion between the two crossings is one connected
residual subsegment.

It remains to prove the depth inequalities for every point of this source
subsegment; these are not finite limit data and are not obtained merely by
pullback.  If $z$ lies outside the two cap interiors, then
\eqref{eq:ncs-source-endcap-large-ball} gives, for both signs,
\begin{equation}
 R(x_{\pm,i},t_i)d_{t_i}(x_{\pm,i},z)^2
      \geq D_H^2\geq B_{\rm lr}(H).                  \label{eq:ncs-source-endcap-transfer-input}
\end{equation}
Apply \eqref{eq:ncs-long-range-threshold} in the source, with
$x=x_{\pm,i}$ and $y=z$, to obtain
\eqref{eq:ncs-endcap-axial-depth}.  The same ball inclusion and
$D_H>H$ put the clusters $K_\pm$ strictly inside their caps.  Thus every
sufficiently late counterexample actually has all required witnesses, a
contradiction.  This sequential compactness argument defines the finite
threshold $L(\varepsilon,H)$.
\end{proof}

\begin{lemma}[The ordered middle is a certified neck tube]
\label{lem:ncs-ordered-middle-tube}
Let the data satisfy the conclusions of
Lemma~\ref{lem:ncs-remote-endpoint-caps} with
$0<\varepsilon\leq\varepsilon_{\rm cross}$ and
$H\geq L_{\rm cross}$.  Assume
that every point outside $\operatorname{int}X_-\cup\operatorname{int}X_+$
is a strong $\varepsilon$-neck center.  Let $W$ be the closure of the
component of $M\setminus(X_-\cup X_+)$ which contains the open residual
subsegment between the two points in
\eqref{eq:ncs-endcap-distinguished-crossings}.  Then $W$ admits a finite
$\varepsilon$-neck-chain witness in the sense of Definition
\ref{def:ncs-neck-chain}.  Consequently $W$ is a neck tube and
\begin{equation}
                         M=X_-\cup W\cup X_+           \label{eq:ncs-certified-two-cap-decomposition}
\end{equation}
with the exact intersection identities of
\eqref{eq:ncs-two-cap-union}.
\end{lemma}

\begin{proof}
A compact nonround solution has positive sectional curvature by the strong
maximum principle and the rank-one list.  It is orientable, and
Bonnet--Myers makes its fundamental group finite.  Every two-sided
embedded sphere in the middle therefore separates: a nonseparating
oriented sphere would have nonzero Poincar\'e dual in
$H^1(M;\mathbb Z)$ and induce a homomorphism of $\pi_1(M)$ with infinite
image.

Let $\gamma_0$ be the closed residual subsegment between the unique
boundary crossings furnished by Lemma
\ref{lem:ncs-remote-endpoint-caps}.  Every one of its points is a strong
$\varepsilon$-neck center.  Retain one chosen witness at each point and,
only for the axial-crossing conclusion, weaken it to tolerance
$\varepsilon_{\rm cross}$.  Equation
\eqref{eq:ncs-endcap-axial-depth} and Lemma
\ref{lem:ncs-two-sided-axial-crossing} choose all axial signs in the
direction from $x_-$ to $x_+$.  The endpoint witnesses have the same signs
by construction.

We now construct the linear nerve; it is not inferred from intersections
of arbitrary center spheres.  The following finite continuation procedure
will be used.  Its induction invariant consists of
\begin{enumerate}
\item a separating regular sphere $F_j$ meeting $\gamma_0$ exactly once;
\item a compact swept product $T_j$ from $\partial X_-$ to $F_j$, lying
      on the $x_-$ side of $F_j$; and
\item chosen truncated neck charts covering $T_j$, with full
      cross-sectional consecutive overlaps and with all nonconsecutive
      truncations disjoint.
\end{enumerate}
For $j=0$, choose a regular positive slice $F_0$ in the recorded collar of
$\partial X_-$, and let $T_0$ be the closed product collar from
$\partial X_-$ to $F_0$.  The axial-crossing lemma gives its unique
intersection with $\gamma_0$, so all three invariants hold.

Suppose the data have been constructed and $F_j\ne\partial X_+$.  Let
$p_j=F_j\cap\gamma_0$.  Choose a strong $\varepsilon$-neck chart centered
at $p_j$; for $j=0$ its negative collar is matched with the recorded
endpoint chart.  Quantified overlap
alignment, the reserve
\eqref{eq:ncs-crossing-overlap-reserve}, and the unit collar already stored
at $F_j$ show by the implicit-function theorem that the existing $F_j$ is a graph of norm
$O(\varepsilon)$ over the chart's central product slice.  In particular,
the overlap contains a full spherical collar, not merely a neighborhood
of $p_j$.

Consider the entire closed forward chart slab from the graph $F_j$ to the
positive coordinate-$3$ slice.  If this slab is disjoint from the recorded
collar of $\partial X_+$, declare its terminal regular slice to be
$F_{j+1}$.  The slab between $F_j$ and $F_{j+1}$ is an embedded product.
Its interior is connected,
disjoint from $F_j$, and contains the forward part of $\gamma_0$; hence it
lies on the positive side of the separating sphere $F_j$ and is disjoint
from $T_j$.  It is also disjoint from $X_+$: otherwise a path in this
connected slab from its forward segment point to a point of
$\operatorname{int}X_+$ would meet $\partial X_+$ and hence its recorded
collar, contrary to the branch condition.  With the common collar
identified, define
$T_{j+1}$ to be the union of $T_j$ with this closed slab.  At
$p_{j+1}=F_{j+1}\cap\gamma_0$, choose the next neck chart.  The old chart
has more than ninety units of unused buffer.  We use the following elementary
buffered-capture fact.  In two scalar-normalized
$\varepsilon_{\rm cross}$-necks, if one center lies three axial units from
the other and the first chart has twenty unused units, then the scalar ratio
lies in $[1/2,2]$ and the full unit central collar of the second chart lies
inside the first.  Indeed, either collar lies in a model ball of radius six
about its center; $C^0$ metric closeness gives the two source-ball
containments.  This fact puts a full central spherical collar of the new
chart inside the old chart, and Lemma~\ref{lem:ncs-neck-overlap-alignment} records the
overlap domain, its sign, and its translation constant.  Truncate both
charts inside this common collar.  The new transition sphere is thus
constructed as a regular product slice before any claim of nesting is
made.

If instead that closed forward slab intersects the recorded endpoint
collar, the fixed sphere-diameter bound and the large unused neck buffers
put a full subcollar of each chart in their common domain.  Use overlap
alignment there, truncate the last slab at
$F_{j+1}=\partial X_+$, and retain the exact endpoint boundary
parametrization.  The two collar charts give the final full
cross-sectional overlap; define
$T_{j+1}$ by adjoining the resulting closed final slab to $T_j$.  In either
case, Lemma~\ref{lem:ncs-two-sided-axial-crossing} says that $\gamma_0$
meets $F_{j+1}$ only once.  Separation now implies strict nesting: the
already swept region is on its negative side and the unswept terminal
subsegment is on its positive side.  Thus the induction invariant is
preserved.  Notice the order of the argument: regular disjoint transition
slices are constructed from overlap charts first; separation is used only
afterward.

The procedure terminates after finitely many steps.  Indeed, unless it is
the final step, coordinate $3$ advances arclength along $\gamma_0$ by at
least $cR(p_j,t)^{-1/2}$ for a universal $c>0$.  The compact slice has
$R_{\max}=\max_M R(\cdot,t)<\infty$, so every nonfinal advance is at least
$cR_{\max}^{-1/2}>0$, whereas $\gamma_0$ has finite length.  At
termination, shrink the stored collars once.  The resulting truncated
images have exactly the linear nerve $0-1-\cdots-m$, consecutive overlap
domains satisfying the intrinsic path-length hypothesis of Lemma
\ref{lem:ncs-neck-overlap-alignment}, and disjoint nonconsecutive pieces.
Their final swept region $T_m$ has
\begin{equation}
                       \partial T_m=\partial X_-\sqcup\partial X_+.     \label{eq:ncs-constructed-middle-boundary}
\end{equation}
These are precisely the fields of a finite neck-chain witness.  Only now
do we invoke Proposition~\ref{prop:ncs-neck-chain-annulus}; it equips the
already constructed region $T_m$ with a neck-tube witness.  Put $T=T_m$.

The compact union $X_-\cup T\cup X_+$ is closed in $M$.  It is a smooth
codimension-zero submanifold with empty boundary, because the three
recorded boundary parametrizations cancel in their collars.  A
codimension-zero submanifold without boundary is open.  Hence the union is
both open and closed; connectedness of $M$ makes it all of $M$.  Therefore
$T=W$, and \eqref{eq:ncs-certified-two-cap-decomposition} holds with the
exact intersection identities.
\end{proof}

\begin{theorem}[Compact qualitative alternative]
\label{thm:ncs-compact-alternative}
There is $\varepsilon_{\rm compact}>0$ such that, for every
$0<\varepsilon\leq\varepsilon_{\rm compact}$, there is a universal
$C=C(\varepsilon)$ such that a compact three-dimensional ancient
$\kappa$-solution at a fixed time satisfies one of the following.
\begin{enumerate}
\item The metric has constant positive sectional curvature.
\item There is a point $x$ such that the entire component is controlled at
      the scale $Q=R(x,t)$:
      \begin{equation}
       \operatorname{diam}_tM\leq C Q^{-1/2},\qquad
      C^{-1}Q\leq R(\cdot,t)\leq CQ,\qquad
       \sec_{g(t)}\geq C^{-1}Q.                        \label{eq:ncs-compact-controlled}
      \end{equation}
\item There is a two-cap--tube witness.  Every point of the open middle
      tube is a strong $\varepsilon$-neck center.  Each cap is of kind
      $\mathbf B$ or $\mathbf P$.  At most one cap has kind $\mathbf P$.
\end{enumerate}
No comparison between the two cap scales, no uniform diameter estimate for
a global cap, and no single curvature bound relative to one cap across a
long tube is asserted.
\end{theorem}

\begin{proof}
Choose $\varepsilon_{\rm compact}$ below the four-alternative,
endpoint-cap, axial-crossing, and neck-chain thresholds, with room for the
factor $1/8$ used below.
Apply Theorem~\ref{thm:ncs-four-alternatives} with tolerance
$\varepsilon/8$.  Alternative (D) gives the
diameter estimate in \eqref{eq:ncs-compact-controlled}.  We record the
scalar comparison directly.  Put $Q=R(x,t)$,
$d=d_t(x,z)$, and $J=R(z,t)/Q$.  If $J\geq4$, the Lipschitz estimate for
$R^{-1/2}$ gives $Qd^2\geq c_\eta$ and hence
$R(z,t)d^2\geq c_\eta J$.  Long-range transfer from $z$ to $x$, together
with $Qd^2\leq A_{\rm four}(\varepsilon/8)$, bounds $J$ from above.  If
$J\leq1/4$, the same Lipschitz estimate gives
$Qd^2\geq c_\eta J^{-1}$, while alternative (D) bounds $Qd^2$ from above;
this bounds $J$ away from zero.  Thus scalar curvature is globally
comparable with $Q$ in alternative (D).  If no universal positive lower bound for
$\sec/Q$ existed in the nonround case, a normalized compactness limit would
have a zero curvature-operator eigenvalue.  The strong maximum principle
and Theorem~\ref{thm:ksol-three-dimensional-split-branch} would make its
universal cover a rank-one cylinder.  Universal $\kappa_0$-noncollapsing
passes to the limit, whereas a compact axial quotient of a shrinking
cylinder violates all-scale noncollapsing by
\eqref{eq:ksol-translation-collapse}.  This contradiction proves the last
inequality in \eqref{eq:ncs-compact-controlled}.  Because the argument
starts with a sequence on which every candidate lower bound fails, the
resulting $C(\varepsilon)$ is uniform over all source solutions.

Consider alternative (C), with endpoint bad points $x_-,x_+$.  Put
\begin{equation}
 A_*=A_{\rm four}(\varepsilon/8),
 \qquad H_*=\max\{A_*,L_{\rm cross}\},                \label{eq:ncs-compact-middle-constants}
\end{equation}
and let $L=L(\varepsilon,H_*)$ be supplied by Lemma
\ref{lem:ncs-remote-endpoint-caps}.  If one of the two products in
\eqref{eq:ncs-two-sided-remote-endpoints} is at most $L$, use that endpoint
as base.  Scalar-scale comparison and long-range transfer bound scalar
curvature above and below on the segment in that base scale.  The opposite
endpoint scale is therefore comparable with the base scale, so both
endpoint balls of radius $A_*R(x_\pm,t)^{-1/2}$ have controlled diameter
and scalar curvature in that scale.  Every bad point lies in their union.
For a good point $z$ outside both endpoint balls,
\eqref{eq:ncs-neck-close-to-axis} gives
$R(z)d(z,\gamma)^2\leq A_*$.  If its distance from the segment in the base
scale were unbounded, long-range transfer from a nearest point of
$\gamma$ would contradict this inequality.  These three classes---the two
endpoint balls, all bad points, and good points outside the balls---exhaust
$M$.  Thus the whole component has bounded diameter and scalar ratio in
one curvature scale, and the preceding compactness argument supplies the
sectional lower bound.  This is alternative~(2).

In the remaining case both products exceed $L$.  Since the bad clusters
in (C) use radius $A_*$, they are contained in the radius-$H_*$ clusters.
Lemma
\ref{lem:ncs-remote-endpoint-caps} supplies disjoint endpoint caps
containing both bad clusters.  Hence every point in the residual middle is
a target-$\varepsilon$ neck center.  Lemma
\ref{lem:ncs-ordered-middle-tube} supplies the certified middle tube and
the two-cap--tube witness.

It remains to exclude two projective caps.  Gluing the cap models through a
spherical tube gives the following list.  Absorb the product tube into the
two boundary collars.  Every gluing diffeomorphism of $\Sph^2$ extends over
$B^3$ after composition with a reflection when necessary, so the
diffeomorphism type is independent of the recorded boundary
parametrizations:
\begin{equation}
 (\mathbf B,\mathbf B)\mapsto\Sph^3,\qquad
 (\mathbf B,\mathbf P)\mapsto\mathbb{RP}^3,\qquad
 (\mathbf P,\mathbf P)\mapsto
            \mathbb{RP}^3\#\mathbb{RP}^3.             \label{eq:ncs-cap-gluing-list}
\end{equation}
The last manifold has fundamental group
$\mathbb Z_2*\mathbb Z_2$, which is infinite.  A compact nonround
$\kappa$-solution has positive sectional curvature, so Bonnet--Myers makes
its fundamental group finite.  Hence the last gluing cannot occur.  This
argument, rather than orientation, excludes two projective caps.
\end{proof}

\begin{corollary}[Topology of compact nonround solutions]
\label{cor:ncs-compact-topology}
If a compact three-dimensional ancient $\kappa$-solution is not an exact
round spherical-space-form flow, then its underlying manifold is
diffeomorphic to $\Sph^3$ or $\mathbb{RP}^3$.
\end{corollary}

\begin{proof}
Take an asymptotic shrinker as in
Corollary~\ref{cor:red-three-asymptotic-shrinker}.  A compact spherical
asymptotic shrinker would trigger
Lemma~\ref{lem:ksol-round-backward-rigidity} and make the original flow
round.  Hence the shrinker is cylindrical and the backward rescaled
diameters are unbounded.  More explicitly, choose the $\ell$-centers
$q_i$ and backward times $\tau_i\to\infty$ used in the pointed shrinker
limit.  Smooth convergence to a nonflat cylinder gives constants
$0<c<C<\infty$ such that
\begin{equation}
             c\leq\tau_iR(q_i,-\tau_i)\leq C,         \label{eq:ncs-ell-center-scalar-scale}
\end{equation}
and source-ball capture of arbitrarily large limiting cylinder balls gives
$\tau_i^{-1/2}\operatorname{diam}_{-\tau_i}M\to\infty$.

If the controlled-component branch of Theorem
\ref{thm:ncs-compact-alternative} held along a subsequence, its global
scalar comparability would compare its controlling scale $Q_i$ with
$R(q_i,-\tau_i)$.  Equation
\eqref{eq:ncs-ell-center-scalar-scale} and its diameter estimate would then
bound $\tau_i^{-1/2}\operatorname{diam}_{-\tau_i}M$, a contradiction.
The constant-curvature branch is already excluded.  Thus sufficiently far
along the selected sequence the two-cap--tube branch holds.  One such time
is enough to determine the underlying smooth manifold.  The first two gluings in
\eqref{eq:ncs-cap-gluing-list} give $\Sph^3$ or $\mathbb{RP}^3$.  The
smooth manifold is independent of time.
\end{proof}

\section{The canonical-neighborhood theorem for $\kappa$-solutions}
\label{sec:ncs-canonical-neighborhood}

The following is the exact surgery-facing local output.  A neighborhood is
carried as a chosen codimension-zero submanifold; set inclusions are not
identified with equalities of metric balls.

\begin{definition}[Local canonical neck witness]
\label{def:ncs-local-canonical-neck}
Fix $L_{\rm can}=10$ and $0<\varepsilon<1/11$, so that
$L_\varepsilon>L_{\rm can}+1$.  A local $\varepsilon$-neck witness based at $(x,t)$
consists of a chosen strong $\varepsilon$-neck witness $\Phi$ based at
$(x,t)$, its terminal map $\phi_t$, and the exact compact
codimension-zero submanifold
\begin{equation}
 U=\phi_t\bigl(\Sph^2\times[-L_{\rm can},L_{\rm can}]\bigr).    \label{eq:ncs-local-neck-region}
\end{equation}
The two boundary parametrizations are the restrictions of $\phi_t$ at
$\pm L_{\rm can}$.
\end{definition}

\begin{definition}[Local canonical cap witness]
\label{def:ncs-local-canonical-cap}
At time $t$, a local $\varepsilon$-cap witness based at $x$ consists of a
compact connected codimension-zero submanifold $U$, together with the
following nested and typed data.
\begin{enumerate}
\item There is a compact connected cap core
      $X\subset\operatorname{int}U$, a tag
      $\mathfrak c\in\{\mathbf B,\mathbf P\}$, a diffeomorphism
      $\chi:\mathfrak c\to X$, and a distinguished point
      $x\in\operatorname{int}X$.
\item There is a closed neck tube $V\subset U$ with
      \begin{equation}
       U=X\cup V,\qquad
       X\cap V=\partial X,\qquad
       \partial V=\partial X\sqcup\partial U.          \label{eq:ncs-local-cap-exact-union}
      \end{equation}
      The data store both a finite ordered strong
      $\varepsilon$-neck-chain whose swept region is exactly $V$ and the
      resulting neck-tube witness
      $\Psi:\Sph^2\times[0,1]\to V$, with the exact boundary equations
      \begin{equation}
       \Psi(\Sph^2\times\{0\})=\partial X,
       \qquad
       \Psi(\Sph^2\times\{1\})=\partial U.            \label{eq:ncs-local-cap-tube-boundaries}
      \end{equation}
\item The cap and tube boundary parametrizations agree exactly, and all
      axial signs, collars, and transition constants are retained.
\end{enumerate}
In particular,
\begin{equation}
 \partial U\subset V,\qquad
 \overline{U\setminus V}\subset\operatorname{int}U,   \label{eq:ncs-local-cap-collar-position}
\end{equation}
but these inclusions are consequences of the exact union data rather than
substitutes for them.  The phrase ``the cap is based at $x$'' means equality
with the distinguished-point field, not merely membership of $x$ in $U$.
\end{definition}

\begin{lemma}[Truncating a qualitative cap]
\label{lem:ncs-cap-truncation-sandwich}
There is $\varepsilon_{\rm trunc}>0$ such that the following holds for every
$0<\varepsilon\leq\varepsilon_{\rm trunc}$.  Let a noncompact ancient
$\kappa$-solution lie in the positive-curvature or diagonal-quotient
branch, equip it with the one-cap-end witness supplied by Theorem
\ref{thm:ncs-noncompact-neck-cap} at tolerance $\varepsilon/4$, let $x$
lie in its compact core, and put $Q=R(x,t)$.  For every $0\leq H<\infty$ there
are a local $\varepsilon$-cap witness $(X,U,V)$ based at $x$ and a number
$r$ such that
\begin{equation}
 r\geq Q^{-1/2},\qquad
 B_t(x,r)\subset U\subset B_t(x,2r),\qquad
 d_t(x,V)\geq H Q^{-1/2}.                             \label{eq:ncs-cap-truncation-sandwich}
\end{equation}
For the selected finite truncation, scalar curvature has positive finite
extrema on $U$ and $\Vol_tU>0$.
\end{lemma}

\begin{proof}
Choose $\varepsilon_{\rm trunc}$ so that $\varepsilon/4$ lies below the
noncompact neck--cap threshold and every recentered terminal neck can be
weakened to the target $\varepsilon$.
We first treat the positive-curvature branch.  Choose a unit-speed ray
$\gamma$ issuing from $x$ and normalize its Busemann function so that
$b_\gamma(x)=0$; thus $b_\gamma(\gamma(s))=s$.  Theorem
\ref{thm:ncs-arbitrary-escape-cylinder} permits the choice of
$s_j\to\infty$ such that $y_j=\gamma(s_j)$ has a chosen strong neck
witness, retained with its terminal center sphere, and with tolerance
tending to zero.  Pass to disjoint controlled cores and then to
the nested subsequence of Lemma
\ref{lem:ksol-neck-outward-subsequence}.  Let $U_j$ be the compact cap side
of the center sphere $S_j$ and put
\begin{equation}
 a_j=d_t(x,\partial U_j),\qquad
 b_j=\sup_{z\in U_j}d_t(x,z).                          \label{eq:ncs-cap-inner-outer-radii}
\end{equation}
Write $r_j=R(y_j,t)^{-1/2}$.  Corollary
\ref{cor:ncs-escape-curvature-scale} gives
\begin{equation}
                         r_j/s_j\longrightarrow0.      \label{eq:ncs-cap-neck-radius-little-o}
\end{equation}

We now prove the required radial estimate without summing errors over a
long neck chain.  The comparison in Lemma
\ref{lem:ksol-neck-Busemann-diameter}, applied in a fixed outward
subcollar of the $j$th increasingly round neck, gives one level
$b_\gamma^{-1}(c_j)$ and one dimensional constant $C$ such that
\begin{equation}
 s_j\leq c_j\leq s_j+Cr_j,\qquad
 \operatorname{diam}b_\gamma^{-1}(c_j)\leq Cr_j,       \label{eq:ncs-cap-outer-level-control}
\end{equation}
and $b_\gamma\leq c_j$ on $S_j$.  A minimizing segment from any
$z\in U_j$ to $\gamma(s)$, for large $s$, exits the compact side through
$S_j$.  Letting $s\to\infty$ shows
$b_\gamma(z)\leq\max_{S_j}b_\gamma\leq c_j$.

If $0\leq c=b_\gamma(z)\leq c_j$, then both $z$ and $\gamma(c)$ lie on
$b_\gamma^{-1}(c)$.  The surjective distance-nonincreasing level map of
Lemma~\ref{lem:ksol-Sharafutdinov-level-surjective} and
\eqref{eq:ncs-cap-outer-level-control} give
\begin{equation}
                  d_t(z,\gamma(c))\leq Cr_j.           \label{eq:ncs-cap-same-level-distance}
\end{equation}
The fixed sublevel $\{b_\gamma\leq0\}$ is compact.  Indeed, choose one
remote neck with $b_\gamma(y)>1$; the outward inequality
\eqref{eq:ksol-Busemann-outward} puts the entire region beyond its positive
collar in $\{b_\gamma>0\}$, while the complementary cap side and collar are
compact.  Absorb the sublevel diameter in a constant $D_0$.  Since
$b_\gamma$ is
one-Lipschitz and $|b_\gamma-s_j|\leq Cr_j$ on $S_j$, the preceding
estimates yield
\begin{equation}
 s_j-Cr_j\leq a_j\leq b_j
       \leq s_j+Cr_j+D_0.                              \label{eq:ncs-cap-two-radii-estimate}
\end{equation}
Together with \eqref{eq:ncs-cap-neck-radius-little-o}, this proves
\begin{equation}
                         a_j\longrightarrow\infty,\qquad
                         b_j/a_j\longrightarrow1.      \label{eq:ncs-cap-radius-ratio}
\end{equation}

In the diagonal branch take
$U_j=(\Sph^2\times[-L_j,L_j])/\mathbb Z_2$ with $L_j\to\infty$ and large
enough that $x\in\operatorname{int}U_j$.  The spherical diameter is fixed
at time $t$, whereas the two lifted axial distances from a representative
of $x$ are $L_j+O(1)$.  Thus
\eqref{eq:ncs-cap-radius-ratio} holds directly, with cap type $\mathbf P$.

In either branch, recenter the terminal neck at two fixed axial levels in
its buffered core.  Let $S_j^{\rm in}$ and $S_j^{\rm out}=\partial U_j$
be the resulting ordered center spheres, let $V_j$ be the closed region
between them, and let $X_j$ be the compact side of $S_j^{\rm in}$.  Then
\begin{equation}
 U_j=X_j\cup V_j,\qquad X_j\cap V_j=S_j^{\rm in},\qquad
 \partial V_j=S_j^{\rm in}\sqcup\partial U_j.         \label{eq:ncs-truncated-cap-exact-union}
\end{equation}
For large $j$, the distinguished point $x$ lies in
$\operatorname{int}X_j$; $X_j$ has the same cap type as $U_j$; and the
recentered charts give $V_j$ a finite ordered $\varepsilon$-neck-chain
witness with the exact boundary parametrizations.  Moreover
$d_t(x,V_j)\to\infty$.

Choose $j$ with $a_j>Q^{-1/2}$, $b_j<2a_j$, and
$d_t(x,V_j)>HQ^{-1/2}$.  Select $r$ with
\begin{equation}
             \max\{b_j/2,Q^{-1/2}\}<r<a_j.           \label{eq:ncs-cap-radius-choice}
\end{equation}
The interval is nonempty by the two strict inequalities just imposed.
The definition of $a_j$ gives
$B_t(x,r)\subset U_j$, and the definition of $b_j$ gives
$U_j\subset B_t(x,2r)$.  Take $(X,U,V)=(X_j,U_j,V_j)$.  Compactness,
smoothness, and scalar positivity give the last assertions.
\end{proof}

\begin{theorem}[$\kappa$-solution canonical neighborhoods]
\label{thm:ncs-kappa-canonical-neighborhood}
There is $\varepsilon_{\rm can}\in(0,1/11)$ such that, for every
$0<\varepsilon\leq\varepsilon_{\rm can}$, there are universal constants
$C_1=C_1(\varepsilon),C_2=C_2(\varepsilon)\geq1$ with the following
property.  Let $(x,t)$ be any spacetime point in any three-dimensional
\emph{oriented} ancient $\kappa$-solution and put $Q=R(x,t)$.  There are a
radius $r$ and a chosen time-$t$ neighborhood $U$ such that
\begin{equation}
 Q^{-1/2}\leq r\leq C_1Q^{-1/2},\qquad
 B_t(x,r)\subset U\subset B_t(x,2r),                  \label{eq:ncs-canonical-ball-sandwich}
\end{equation}
and at least one of the following holds.
\begin{enumerate}
\item[(a)] $U$ carries a local canonical $\varepsilon$-neck witness based
      at $(x,t)$ in the sense of Definition
      \ref{def:ncs-local-canonical-neck}.
\item[(b)] $U$ carries a local $\varepsilon$-cap witness of Definition
      \ref{def:ncs-local-canonical-cap}, of kind $\mathbf B$ or $\mathbf P$,
      whose distinguished point is exactly $x$.  Its recorded outward neck
      collar $V\subset U$ satisfies
      \begin{equation}
           d_t(x,V)\geq10^4Q^{-1/2}.                  \label{eq:ncs-deep-in-cap}
      \end{equation}
\item[(c)] $U=M$, the manifold is diffeomorphic to $\Sph^3$ or
      $\mathbb{RP}^3$, and
      \begin{equation}
                           \sec_{g(t)}\geq C_2^{-1}Q.   \label{eq:ncs-canonical-positive-sec}
      \end{equation}
\item[(d)] $U=M$ and $g(t)$ has constant positive sectional curvature.
\end{enumerate}
In every case (a)--(d),
\begin{equation}
             C_2^{-1}Q\leq R(z,t)\leq C_2Q
             \qquad(z\in U).                          \label{eq:ncs-canonical-scalar}
\end{equation}
In cases (a), (b), and (c), but not uniformly in case (d),
\begin{equation}
                         \Vol_tU\geq C_2^{-1}Q^{-3/2}. \label{eq:ncs-canonical-volume}
\end{equation}
\end{theorem}

The orientation hypothesis is the one used by Perelman and by the surgery
chapters.  Without it, the exact projective product
$\mathbb{RP}^2\times\R$ remains the separate model branch in
Theorem~\ref{thm:ncs-noncompact-neck-cap}; it is not mislabeled as an
ordinary spherical neck.

\begin{proof}
Suppose no universal constants exist.  Choose counterexamples with
$R(x_i,t_i)=1$ after curvature normalization and with every candidate
constant at most $i$ failing.  Exact round sources satisfy (d) at a
universal local curvature scale, although their total volumes may tend to
zero as the deck-group order grows.  Thus all nonround sources have the
universal $\kappa_0$.

Apply Theorem~\ref{thm:ksol-universal-precompactness}.  If the pointed limit
is compact, source-ball capture and the same-dimensional open-and-closed
argument identify the whole $M_i$ with the limit for large $i$.  If the
limit is the round quotient $\Sph^3/\Gamma$, topology stability and
Corollary~\ref{cor:ncs-compact-topology} require nearby nonround sources to
have $|\Gamma|=1$ or $2$.  They then satisfy (c).  If $|\Gamma|>2$, the
nearby sources must themselves be round and satisfy (d).  If the limit is
nonround, Theorem
\ref{thm:ncs-compact-alternative} and
Corollary~\ref{cor:ncs-compact-topology} give (c).  Smooth compact
convergence supplies the scalar, sectional, diameter, and positive-volume
bounds.  This contradicts the selected failure.

Suppose the limit is noncompact.  It cannot be the projective cylinder: a
sufficiently large compact limit domain is a nonorientable
codimension-zero manifold, whereas a convergence embedding places it as a
codimension-zero domain in the oriented source.  If its basepoint is the
center of a strong $\varepsilon/2$-neck, smooth pointed convergence and
strict reserve give (a) in the sources: in the limiting flow choose the
assumed strong $\varepsilon/2$-witness and take the exact region
$U_\infty=\phi_t(\Sph^2\times[-L_{\rm can},L_{\rm can}])$ from
Definition~\ref{def:ncs-local-canonical-neck}.  Its boundary distance $a$
and outer radius $b$ satisfy $b<2a$ after the fixed harmless cylinder
estimate.  Choose $\max\{b/2,1\}<r<a$ and pull that entire finite tuple
back.
Otherwise the basepoint belongs to
$\mathscr B_{\varepsilon/2}$.  Since
$\mathscr B_{\varepsilon/2}\subset\mathscr B_{\varepsilon/32}$, Theorem
\ref{thm:ncs-noncompact-neck-cap}, applied at tolerance
$\varepsilon/8$, gives a ball or projective one-cap-end witness whose
compact core contains the limiting basepoint.  Apply Lemma
\ref{lem:ncs-cap-truncation-sandwich} with cap tolerance
$\varepsilon/2$ and $H=2\cdot10^4$.  The resulting fixed finite local
$\varepsilon/2$-cap witness satisfies, with strict reserve, both
\eqref{eq:ncs-deep-in-cap} and the sandwich
\eqref{eq:ncs-canonical-ball-sandwich}.  On this one fixed compact witness,
scalar curvature has positive finite extrema and volume is positive.
Pull the entire tuple of maps, inclusions, collars, and inequalities back by
the convergence embeddings.  Its finite $\varepsilon/2$ neck chain becomes
an $\varepsilon$ neck chain, and the doubled distance margin gives
\eqref{eq:ncs-deep-in-cap}.  For all large $i$ this is alternative (b)
with finite constants independent of $i$, again a contradiction.

The proof used only universal compactness, the qualitative cap theorem,
and compact topology.  It did not use the Bryant or ancient-oval
classification.  The omission of a volume assertion in (d) is forced by
round quotients $\Sph^3/\Gamma$ with $|\Gamma|\to\infty$.
\end{proof}

\section{Reduced-geometric inputs for neck stability}
\label{sec:ncs-stability-inputs}

Fix a pole $(p,0)$ and use backward time
\begin{equation}
 h(\tau)=g(-\tau),\qquad
 \ell_p(q,\tau)=\ell_{(p,0)}(q,\tau).                  \label{eq:ncs-backward-notation}
\end{equation}
For $\rho>0$ define the Type~I rescaling
\begin{equation}
 h^{[\rho]}(\theta)=\rho^{-1}h(\rho\theta),\qquad
                         0<\theta<\infty.              \label{eq:ncs-TypeI-backward-rescaling}
\end{equation}
The extinction-centered backward round cylinder and its normalized
potential are
\begin{equation}
 h_{\rm C}(\theta)=2\theta g_{\Sph^2(1)}+dz^2,\qquad
 \ell_{\rm C}(y,z,\theta)=1+\frac{z^2}{4\theta}.       \label{eq:ncs-extinction-cylinder-triple}
\end{equation}
At $\theta=1$ its scalar curvature is one.  This cylinder is singular at
$\theta=0$; its potential is a blow-down potential, not a reduced distance
from a regular pole at that time.

\begin{definition}[Type~I cylindrical and cylindrical-triple witnesses]
\label{def:ncs-TypeI-cylinder-witness}
Relative to the fixed regular pole $(p,0)$ of
\eqref{eq:ncs-backward-notation}, and for $0<\delta<1$, a point
$(x,-\rho)$ is
\emph{Type~I $\delta$-cylindrical} if one time-independent pointed embedding
\begin{equation}
 \Phi:\Sph^2\times[-\delta^{-1},\delta^{-1}]\longrightarrow M,
 \qquad \Phi(y_*,0)=x,                                \label{eq:ncs-TypeI-embedding}
\end{equation}
makes $h^{[\rho]}(\theta)$ $\delta$-close to
$h_{\rm C}(\theta)$ on
\begin{equation}
 \Sph^2\times[-\delta^{-1},\delta^{-1}]
       \times[\delta,1+\delta^{-2}].                  \label{eq:ncs-TypeI-fixed-window}
\end{equation}
Precisely, with $k_\delta=\lceil\delta^{-1}\rceil$, the sum of the suprema
of the background-cylinder covariant derivatives
$\nabla^j\partial_\theta^m(\Phi^*h^{[\rho]}-h_{\rm C})$ for
$j+2m\leq k_\delta$, measured with $h_{\rm C}$ on the displayed domain,
is less than $\delta$.  Restriction to a smaller spatial cylinder is part
of weakening a witness.

A \emph{$\delta$-cylindrical triple witness} is parameterized by the tuple
$(p,\rho,x)$, retains the same embedding, and also stores a number
$z_0\in\R$ for which the pulled-back reduced distance has
supremum error less than $\delta$ from
\begin{equation}
              1+\frac{(z+z_0)^2}{4\theta}.             \label{eq:ncs-cylinder-offset-potential}
\end{equation}
on the same domain.  Only this $C^0$ function estimate is consumed by the
first-crossing lemma; no unrecorded derivative norm for $\ell$ is assumed.
The offset is part of the witness.  A bounded-$\ell$ moving center need not
converge to the minimum slice $z+z_0=0$.  Convergence embeddings for triples
are required to preserve identity worldlines on their common spacetime
domain.
\end{definition}

\begin{lemma}[Regular-pole convergence of reduced-distance tuples]
\label{lem:ncs-regular-pole-tuple-compactness}
Fix $\kappa>0$.  Let normalized pointed ancient $\kappa$-solutions
$(M_i,g_i(t),p_i)$ converge smoothly and pointedly, with worldline-
preserving embeddings, to $(M_\infty,g_\infty(t),p_\infty)$ and
$R(p_i,0)=1$.  After passing to a subsequence,
\begin{equation}
          \ell_{p_i}\longrightarrow\ell_{p_\infty}    \label{eq:ncs-regular-pole-ell-convergence}
\end{equation}
locally uniformly on $M_\infty\times(0,\infty)$.  For every compact
$J\Subset(0,\infty)$,
\begin{equation}
 \sup_{\rho\in J}
   \left|\widetilde V_{p_i}(\rho)
               -\widetilde V_{p_\infty}(\rho)\right|\longrightarrow0.  \label{eq:ncs-regular-pole-mass-convergence}
\end{equation}
\end{lemma}

\begin{proof}
The stationary path at the pole gives a uniform local upper bound for
$\ell_{p_i}$ on every compact positive time interval.  The spatial
$\sqrt\ell$ estimate, the time inequality, and local semiconcavity then
give equicontinuity, so Arzel\`a--Ascoli produces a locally uniform limit.
To identify it, fix a compact endpoint set and a compact time interval.
Here is the confinement step.  A prefix of an $\epsilon_i$-minimizing
$\mathcal L$-curve is itself $\epsilon_i$-minimizing for its prefix
endpoint, since replacing it by a cheaper prefix would decrease the full
action.  Nonnegativity of the integrand bounds every prefix action by the
full action.  On a time interval bounded away from zero, apply the
two-point coercivity estimate to each prefix, using the stationary-path
bound at the pole, to place its endpoint in one fixed metric ball.  A
finite subdivision gives one common captured spacetime set.  On the
remaining short initial interval, the weighted kinetic-energy bound and a
fixed normal neighborhood of the pole give the same conclusion directly.
Thus every relevant almost-minimizer is confined to one common compact
spacetime set.  We use the standing source-ball-capturing form of pointed
convergence: after the confinement radius and time slab have been fixed,
one compact limit chart has image containing the entire corresponding
source ball for every sufficiently large $i$.  Proposition
\ref{prop:red-action-convergence} therefore gives both the liminf and the
recovery-sequence inequalities for the $\mathcal L$-action.  The limit is
the actual reduced distance from $(p_\infty,0)$, proving
\eqref{eq:ncs-regular-pole-ell-convergence}.

For total mass, two-point coercivity about the pole and Bishop--Gromov give
a Gaussian annular majorant uniform in $i$ and in $\rho\in J$.  Integrate
first on a fixed captured ball, use smooth metric and local function
convergence there, and then let the ball radius tend to infinity.  The
uniform tail tends to zero and proves
\eqref{eq:ncs-regular-pole-mass-convergence}.
\end{proof}

\begin{lemma}[Varying-flow compactness for reduced-distance triples]
\label{lem:ncs-varying-ell-compactness}
Fix $\kappa>0$ and $C<\infty$.  Let
$(M_i,g_i(t),p_i)$ be three-dimensional ancient $\kappa$-solutions with
$R(p_i,0)=1$.  Let $\rho_i\to\infty$, and suppose
\begin{equation}
                         \ell_{p_i}(y_i,\rho_i)\leq C.  \label{eq:ncs-moving-center-ell-bound}
\end{equation}
After passing to a subsequence, the pointed triples
\begin{equation}
 \left(M_i,h_i^{[\rho_i]}(\theta),y_i,
       \ell_{p_i}(\,cdot\,,\rho_i\theta)\right)       \label{eq:ncs-rescaled-ell-triples}
\end{equation}
converge on compact subsets of $0<\theta<\infty$ to a complete pointed
backward Ricci flow and a locally Lipschitz function $\ell_\infty$.
Reduced-volume mass does not escape: for every $\theta>0$,
\begin{equation}
 \int_{M_\infty}(4\pi\theta)^{-3/2}e^{-\ell_\infty(\cdot,\theta)}
                 \,d\mu_{h_\infty(\theta)}
       =\lim_{i\to\infty}\widetilde V_{p_i}(\rho_i\theta).   \label{eq:ncs-varying-no-loss-identity}
\end{equation}
If, for every $0<a<b<\infty$,
\begin{equation}
 \widetilde V_{p_i}(a\rho_i)-\widetilde V_{p_i}(b\rho_i)
                         \longrightarrow0,             \label{eq:ncs-varying-zero-drop}
\end{equation}
then the limit is a gradient shrinking soliton and $\ell_\infty$ is its
smooth normalized potential.
\end{lemma}

\begin{proof}
For each fixed $A>1$, the time-ratio estimate first bounds
$\ell_{p_i}(y_i,\rho_i\theta)$ uniformly for
$\theta\in[A^{-1},A]$.  The spatial estimate for $\sqrt\ell$ propagates
this bound to each fixed ball about $y_i$.  The endpoint estimate
\eqref{eq:red-ancient-endpoint-estimate} then gives
\begin{equation}
 R_{h_i^{[\rho_i]}(\theta)}
       \leq3\ell_{p_i}(\,\cdot\,,\rho_i\theta)/\theta. \label{eq:ncs-varying-endpoint-curvature}
\end{equation}
Thus curvature is bounded on every fixed ball and compact
$\theta$-interval.  Noncollapsing gives a
base injectivity-radius bound, and Chapter~\ref{ch:bando-shi} gives all
interior curvature derivatives.  Hamilton compactness yields a complete
metric limit.  More precisely, for every fixed radius $D$ and compact
$\theta$-interval, the source-ball-capturing compactness theorem supplies
one compact limit chart whose image contains the entire source $D$-ball
throughout that interval.  The local semiconcavity and time estimates from
Proposition~\ref{prop:red-ancient-ell-estimates} give local uniform
convergence of the functions and the derivative convergence required in
the weak reduced-distance identities.

For mass tightness, apply Lemma~\ref{lem:red-two-point} with one endpoint
$y_i$.  The uniform bound \eqref{eq:ncs-moving-center-ell-bound} gives a
quadratic lower bound for $\ell_{p_i}$ away from $y_i$.  Bishop--Gromov and
an annular Gaussian sum then give one tail function tending to zero,
uniformly in $i$ on compact $\theta$-intervals.  This is exactly the
no-loss argument of Theorem~\ref{thm:red-no-mass-loss}, now with varying
source flows, and proves \eqref{eq:ncs-varying-no-loss-identity}.

Assume \eqref{eq:ncs-varying-zero-drop}.  The limiting reduced mass is
constant on every compact time interval.  We retain the equality chain at
the available regularity.  The weak reduced-density subsolution passes to
the limit.  Gaussian cutoffs and constancy of total mass make its total
nonnegative defect zero, so the limiting density solves the conjugate heat
equation.  Interior parabolic regularity makes that density, and hence
$\ell_\infty$, smooth.  The smooth conjugate-heat equation for
$u=(4\pi\theta)^{-3/2}e^{-\ell_\infty}$ is equivalent to
\begin{equation}
 \partial_\theta\ell_\infty
   =\Delta\ell_\infty-|\nabla\ell_\infty|^2+R-\frac{3}{2\theta}.
                                                               \label{eq:ncs-limit-conjugate-ell}
\end{equation}
The Hamilton--Jacobi identity passes by the same mechanism as
\eqref{eq:red-limit-gradient-convergence}: uniform semiconcavity and local
uniform convergence give
$\nabla\ell_i\to\nabla\ell_\infty$ almost everywhere and in every finite
$L^p_{\rm loc}$, and these convergences may be inserted into the weak
time-derivative Hamilton--Jacobi identity.  Combine the resulting smooth
identity with \eqref{eq:ncs-limit-conjugate-ell}.  This gives the pointwise
identity
\begin{equation}
 \theta\bigl(2\Delta\ell_\infty-|\nabla\ell_\infty|^2+R\bigr)
                 +\ell_\infty-3=0.                  \label{eq:ncs-limit-v-bracket-zero}
\end{equation}
Thus the limiting $v$-quantity vanishes identically, not merely after
integration.  The local $v$-identity \eqref{eq:red-v-identity} now makes
its nonnegative square defect vanish.  Therefore
\begin{equation}
        \Ric_{h_\infty(\theta)}+\nabla^2\ell_\infty
                  =\frac1{2\theta}h_\infty(\theta),   \label{eq:ncs-limit-shrinker-equation}
\end{equation}
which is the asserted gradient-shrinker equation.
\end{proof}

\begin{lemma}[Uniform gap below the Gaussian]
\label{lem:ncs-Gaussian-gap}
For every fixed $\kappa>0$ there are
$0<\rho_{\rm E}<\infty$ and $\mu_{\rm E}>0$ such that every normalized
pointed three-dimensional ancient $\kappa$-solution satisfies
\begin{equation}
                     \widetilde V_p(\rho_{\rm E})
                          \leq1-\mu_{\rm E}.            \label{eq:ncs-uniform-Gaussian-gap}
\end{equation}
The number $1$ is the Gaussian value in this book's convention.
\end{lemma}

\begin{proof}
If not, after a diagonal choice there would be normalized pointed solutions
with
\begin{equation}
                         \widetilde V_{p_i}(i)>1-i^{-1}.\label{eq:ncs-Gaussian-gap-failure}
\end{equation}
Fixed-$\kappa$ pointed compactness gives a normalized limit with
$R(p_\infty,0)=1$.  For every fixed $\rho$, monotonicity and
\eqref{eq:ncs-Gaussian-gap-failure} sandwich
$\widetilde V_{p_i}(\rho)$ between $1-i^{-1}$ and $1$ once $i>\rho$.
Lemma~\ref{lem:ncs-regular-pole-tuple-compactness} passes the reduced
distances and their whole integrals to the limit, uniformly on compact
positive time intervals.  Hence the limit has reduced volume one at every
positive backward time.  The regular-pole rigidity theorem
\ref{thm:red-regular-pole-rigidity} makes it Euclidean, contradicting
$R(p_\infty,0)=1$.
\end{proof}

\begin{lemma}[Logarithmic plateau selection]
\label{lem:ncs-logarithmic-plateau}
Let $F_i:[a,\rho_i]\to[0,1]$ be nonincreasing, where $a>0$ and
$\rho_i\to\infty$.  After passing to a subsequence, there are
$\sigma_i\to\infty$ and $A_i\to\infty$ such that
\begin{equation}
 a\leq A_i^{-1}\sigma_i<A_i\sigma_i\leq\rho_i,\qquad
 \frac{\rho_i}{\sigma_i}\longrightarrow\infty,        \label{eq:ncs-plateau-scales}
\end{equation}
and
\begin{equation}
 F_i(A_i^{-1}\sigma_i)-F_i(A_i\sigma_i)\longrightarrow0. \label{eq:ncs-plateau-drop}
\end{equation}
\end{lemma}

\begin{proof}
Put $u=\log\rho$ and consider a central portion of
$[\log a,\log\rho_i]$ whose distance from both endpoints tends to infinity
and whose distance from the right endpoint is still a fixed positive
fraction of the total length.  Partition it into $N_i\to\infty$ disjoint
windows, each of length $2\log A_i$ with $A_i\to\infty$ sufficiently
slowly.  The sum of the drops of the nonincreasing function over these
windows is at most one.  One window has drop at most $1/N_i$.  Let
$\log\sigma_i$ be its midpoint.  The choice of the central portion gives
all three scale conclusions in \eqref{eq:ncs-plateau-scales}, and the
selected drop gives \eqref{eq:ncs-plateau-drop}.  Only the order and bounded
total variation of $F_i$ are used.
\end{proof}

\begin{proposition}[Cylindricality at a bounded-$\ell$ moving center]
\label{prop:ncs-moving-basepoint-cylinder}
Fix $0<\widehat\varepsilon<1$ and $0\leq C<\infty$.  There is a universal
$\rho_*=\rho_*(\widehat\varepsilon,C)<\infty$ with the following property.
Let $(M,g(t))$ be a noncompact orientable ancient $\kappa$-solution,
$R(p,0)=1$, and let
\begin{equation}
             \rho\geq\rho_*,\qquad \ell_p(y,\rho)\leq C.       \label{eq:ncs-moving-prop-hypotheses}
\end{equation}
Then either the flow is the diagonal quotient $\mathcal C_{\rm diag}$, or
the rescaled tuple at $(y,\theta=1)$ has a
$\widehat\varepsilon$-cylindrical triple witness on the expanding window
\eqref{eq:ncs-TypeI-fixed-window}.  In particular $(y,-\rho)$ is Type~I
$\widehat\varepsilon$-cylindrical with no unrecorded conversion factor.
\end{proposition}

\begin{proof}
Assume failure and choose counterexamples with $\rho_i\to\infty$.  The
universal-gap theorem first regards every source as a
$\kappa_0$-solution; hence every compactness and Gaussian-gap constant in
this proof is independent of its originally supplied $\kappa$.  The
exact unquotiented cylinder satisfies the conclusion directly.  To make
this literal, use the normalization
$g(t)=2(1-t)g_{\Sph^2(1)}+dz^2$, for which $R(\cdot,0)=1$.  Its reduced
distance from a regular pole is
\begin{equation}
 \ell_{(y_p,z_p)}(y,z,\tau)
  =1-\frac{\arctan\sqrt\tau}{\sqrt\tau}
   +\frac{(z-z_p)^2}{4\tau}
   +\frac{d_{\Sph^2(1)}(y,y_p)^2}
          {2\sqrt\tau\,\arctan\sqrt\tau}.             \label{eq:ncs-regular-pole-cylinder-ell}
\end{equation}
Write a moving center in product coordinates as $(\bar y_i,z_i)$.  At
scale $\rho_i$, set
$z=z_i+\sqrt{\rho_i}\,w$ and
$z_{0,i}=(z_i-z_p)/\sqrt{\rho_i}$.  Then
\begin{equation}
 h_i^{[\rho_i]}(\theta)
   =2(\theta+\rho_i^{-1})g_{\Sph^2(1)}+dw^2,\qquad
 \ell\longrightarrow1+\frac{(w+z_0)^2}{4\theta}.     \label{eq:ncs-exact-cylinder-triple-limit}
\end{equation}
The bound $\ell(y_i,\rho_i)\leq C$ bounds $z_{0,i}$; a subsequence
therefore converges and supplies the required expanding-window triple.
Discard this model and the exceptional diagonal quotient.  The
split-branch theorem then
makes every source positively curved and hence diffeomorphic to $\R^3$.

Let $\rho_{\rm E},\mu_{\rm E}$ be given by
Lemma~\ref{lem:ncs-Gaussian-gap} for $\kappa_0$.  Apply Lemma
\ref{lem:ncs-logarithmic-plateau} to $F_i=\widetilde V_{p_i}$ with
$a=\rho_{\rm E}$ and choose $\sigma_i,A_i$.  Choose an
$\ell$-center $z_i$ at time $\sigma_i$, so
\begin{equation}
                         \ell_{p_i}(z_i,\sigma_i)\leq3/2.      \label{eq:ncs-plateau-ell-center}
\end{equation}
Rescale by $\sigma_i$ about $z_i$.  Lemma
\ref{lem:ncs-varying-ell-compactness} and
\eqref{eq:ncs-plateau-drop} give a gradient-shrinker limit.  Monotonicity
and the Gaussian gap give
\begin{equation}
 \widetilde V_{p_i}(\sigma_i)\leq
 \widetilde V_{p_i}(\rho_{\rm E})\leq1-\mu_{\rm E}.   \label{eq:ncs-plateau-nonGaussian}
\end{equation}
Together with the no-loss identity, this excludes the Gaussian.  A compact limit would,
by source-ball capture and invariance of domain, make the connected source
compact.  The shrinker classification
\ref{thm:complete-three-dimensional-shrinker-classification} leaves the
round cylinder or an order-two cylindrical quotient.  Either quotient
contains an embedded $\mathbb{RP}^2$ in a fixed pointed ball.  Pushing that
limiting surface forward by a convergence embedding would put
$\mathbb{RP}^2$ in $\R^3$.
Therefore the limit is the unquotiented cylinder, and
\begin{equation}
                   \widetilde V_{p_i}(\sigma_i)\longrightarrow2/e.       \label{eq:ncs-first-plateau-mass}
\end{equation}

Every asymptotic shrinker of each source is also unquotiented cylindrical:
it is noncompact by the same open-and-closed argument, nonflat by
Theorem~\ref{thm:red-asymptotic-shrinker}, and neither order-two quotient
embeds back into $\R^3$.  The no-loss theorem and
\eqref{eq:red-three-density-regression} therefore give the exact identity
\begin{equation}
                    \widetilde V_{p_i,\infty}=2/e.      \label{eq:ncs-source-asymptotic-mass}
\end{equation}
Monotonicity and \eqref{eq:ncs-first-plateau-mass} imply, uniformly for
$s\geq\sigma_i$,
\begin{equation}
             2/e\leq\widetilde V_{p_i}(s)\leq2/e+o(1). \label{eq:ncs-volume-squeezed-after-plateau}
\end{equation}

Now rescale by $\rho_i$ about the originally chosen $y_i$.  Since
$\rho_i/\sigma_i\to\infty$, every fixed rescaled time interval lies beyond
$\sigma_i$ for large $i$.  Equations
\eqref{eq:ncs-moving-prop-hypotheses} and
\eqref{eq:ncs-volume-squeezed-after-plateau}, together with Lemma
\ref{lem:ncs-varying-ell-compactness}, give a shrinker limit of constant
mass $2/e$.  This limit is noncompact: a compact limit would be captured
inside one convergence image and the same open-and-closed argument would
make the connected noncompact source compact.  The classification and the quotient mass table force the
unquotiented cylinder.  The function convergence gives
\eqref{eq:ncs-cylinder-offset-potential} for some offset $z_0$.  Thus the
triples are $\widehat\varepsilon$-close for all large $i$, contradicting
their selection.
\end{proof}

\begin{lemma}[Cylinder threshold at a fixed worldline]
\label{lem:ncs-cylinder-threshold}
Fix $\eta>0$ and $A\in(0,1)$.  There is
$\mu=\mu(\eta,A)>0$ with the following property.  Suppose rescaled
flow--$\ell$ triples have a $\mu$-cylindrical triple witness on
$A\leq\theta\leq1$, the embeddings preserve one identity worldline $x$,
and
\begin{equation}
 \ell(x,\theta)\leq1+\eta\quad(A\leq\theta\leq1),\qquad
 \ell(x,1)=1+\eta.                                    \label{eq:ncs-threshold-crossing-data}
\end{equation}
Then these conditions are inconsistent.
\end{lemma}

\begin{proof}
In model coordinates the preserved worldline has one fixed axial
coordinate, which we take to be zero, so the comparison function is
\begin{equation}
                         1+\frac{z_0^2}{4\theta}.       \label{eq:ncs-threshold-model-worldline}
\end{equation}
The $C^0$ triple error and the equality at $\theta=1$ give
$z_0^2/4\geq\eta-\mu$.  At $\theta=A$ they therefore give
\begin{equation}
       \ell(x,A)\geq1+\frac{\eta-\mu}{A}-\mu.          \label{eq:ncs-threshold-direct-lower}
\end{equation}
Choose
$0<\mu<\min\{A,\eta(1-A)/(1+A)\}$.  The right side of
\eqref{eq:ncs-threshold-direct-lower} is then strictly larger than
$1+\eta$, contradicting \eqref{eq:ncs-threshold-crossing-data}.  The
offset $z_0$ cannot be discarded: it is exactly what makes the strict
contradiction visible.
\end{proof}

\begin{lemma}[Sharp time comparison for reduced distance]
\label{lem:ncs-sharp-ell-time-comparison}
For a fixed spatial point $q$ and $0<\rho_1\leq\rho_2$,
\begin{equation}
 \left(\frac{\rho_1}{\rho_2}\right)^2\ell(q,\rho_1)
 \leq\ell(q,\rho_2)
 \leq\frac{\rho_2}{\rho_1}\ell(q,\rho_1).             \label{eq:ncs-sharp-ell-time-comparison}
\end{equation}
\end{lemma}

\begin{proof}
At almost every time,
\eqref{eq:red-ancient-time-estimate} says
$-2\ell/\rho\leq\partial_\rho\ell\leq\ell/\rho$.
Apply the upper-Dini integration lemma
\ref{lem:red-Dini-integration} to
$\rho^2\ell(q,\rho)$ and to $\rho^{-1}\ell(q,\rho)$, using local absolute
continuity.  The first is nondecreasing and the second nonincreasing, which
is exactly \eqref{eq:ncs-sharp-ell-time-comparison}.
\end{proof}

\begin{proposition}[An initial neck captures the fixed basepoint in $\ell$]
\label{prop:ncs-fixed-worldline-capture}
For every $\eta\in(0,1)$ there are universal constants
$\delta_{\rm bp}=\delta_{\rm bp}(\eta)>0$ and
$\rho_{\rm bp}=\rho_{\rm bp}(\eta)<\infty$ such that the following holds.
If a noncompact ancient $\kappa$-solution satisfies
$R(x,0)=1$ and $(x,0)$ is the center of a strong
$\delta_{\rm bp}$-neck, then
\begin{equation}
                         \ell_x(x,\rho)<1+\eta
                         \qquad(\rho\geq\rho_{\rm bp}).\label{eq:ncs-fixed-worldline-ell-bound}
\end{equation}
\end{proposition}

\begin{proof}
Suppose not.  There are normalized pointed solutions whose basepoints are
centers of strong $i^{-1}$-necks and times
$\bar\rho_i\geq i$ at which
\begin{equation}
                         \ell_{x_i}(x_i,\bar\rho_i)\geq1+\eta. \label{eq:ncs-basepoint-ell-failure}
\end{equation}
Discard the exact rank-one branches before using any oriented theorem.  On
$\Sph^2\times\R$, $\mathbb{RP}^2\times\R$, and
$\mathcal C_{\rm diag}$ scalar curvature is spatially constant.  The
stationary curve at $x_i$ therefore has the same action as the stationary
curve on the unquotiented cylinder, and gives
$\ell_{x_i}(x_i,\rho)<1$ for every $\rho>0$.  None can satisfy
\eqref{eq:ncs-basepoint-ell-failure}.  Thus every counterexample is in the
positive-curvature branch, is diffeomorphic to $\R^3$, and is orientable.

Theorem~\ref{thm:ksol-universal-kappa-gap} now regards every member of the
counterexample sequence as a $\kappa_0$-solution.  This conversion is made
before invoking the fixed-$\kappa$ compactness and reduced-distance tuple
lemmas below.

Universal pointed compactness and the expanding strong-neck witnesses make
the limit an unquotiented cylinder on a spacetime neighborhood ending at
zero.  A zero curvature-operator eigenvalue at an interior time, the
whole-flow strong maximum principle, and the two-dimensional theorem extend
this identification to the entire ancient limit.

Lemma~\ref{lem:ncs-regular-pole-tuple-compactness} gives convergence of the
reduced distances from the regular poles at every fixed finite $\rho$.  On
the scalar-one cylinder
\begin{equation}
          g(t)=2(1-t)g_{\Sph^2(1)}+dz^2,               \label{eq:ncs-regular-pole-cylinder}
\end{equation}
the stationary curve at the basepoint gives the exact reduced distance
\begin{align}
 \ell_{\rm cyl}(x,\rho)
 &=\frac1{2\sqrt\rho}\int_0^\rho\frac{\sqrt s}{1+s}\,ds
   =1-\frac{\arctan\sqrt\rho}{\sqrt\rho}<1.            \label{eq:ncs-regular-cylinder-base-ell}
\end{align}
Choose one large fixed $\rho_2$.  For all large $i$,
\begin{equation}
                  \ell_{x_i}(x_i,\rho_2)<1+\eta/4.     \label{eq:ncs-ell-below-before-crossing}
\end{equation}

Let $\rho_i$ be the first time at or after $\rho_2$ at which the continuous
function $\ell_{x_i}(x_i,\rho)$ reaches $1+\eta$.  Then
\begin{equation}
 \ell_{x_i}(x_i,\rho_i)=1+\eta,\qquad
 \ell_{x_i}(x_i,\rho)\leq1+\eta
       \quad(\rho_2\leq\rho\leq\rho_i),               \label{eq:ncs-first-crossing}
\end{equation}
and local convergence gives $\rho_i\to\infty$.  By
Lemma~\ref{lem:ncs-sharp-ell-time-comparison},
\begin{equation}
 \frac{\rho_2}{\rho_i}
       \leq\frac{1+\eta/4}{1+\eta}=:A_0<1.             \label{eq:ncs-crossing-ratio}
\end{equation}
Since $0<\eta<1$, one has $A_0\geq5/8>1/2$.  Choose
$A\in(A_0,1)$.  After rescaling by $\rho_i$, the first-crossing
inequality holds on $A\leq\theta\leq1$.

The moving-center proposition applies with $y_i=x_i$ and $C=1+\eta$.
Choose its accuracy smaller than both $A$ and the constant in
Lemma~\ref{lem:ncs-cylinder-threshold}.  It gives a cylindrical triple,
and that lemma contradicts \eqref{eq:ncs-first-crossing}.  Its exceptional
diagonal alternative cannot occur because all counterexample sources were
already placed in the positive-curvature branch.
\end{proof}

\begin{lemma}[Type~I scale versus scalar-curvature scale]
\label{lem:ncs-TypeI-scalar-scale-bridge}
For every target $\delta\in(0,1)$ there is
$\varepsilon_{\rm TI}(\delta)>0$ such that the following holds.  If
$(x,-\rho)$ is Type~I $\varepsilon$-cylindrical with
$0<\varepsilon\leq\varepsilon_{\rm TI}(\delta)$, put
\begin{equation}
                  a=\rho R(x,-\rho).                   \label{eq:ncs-TypeI-a}
\end{equation}
Then $|a-1|<\delta$.  If $Q=R(x,-\rho)$, the scalar-normalized forward
flow satisfies the exact identity
\begin{equation}
 Q g\!\left(-\rho+\frac rQ\right)
 =a h^{[\rho]}\!\left(1-\frac r a\right),
                    \qquad -1\leq r\leq0.             \label{eq:ncs-two-normalizations-bridge}
\end{equation}
Moreover $(x,-\rho)$ is the center of a strong $\delta$-neck in the scalar
normalization of Definition~\ref{def:ksol-strong-neck}.
\end{lemma}

\begin{proof}
Formula
\eqref{eq:ncs-two-normalizations-bridge} follows by substituting
$\theta=(\rho-r/Q)/\rho=1-r/a$ in
\eqref{eq:ncs-TypeI-backward-rescaling} and using $Q\rho=a$.
On the exact cylinder its right side is
\begin{equation}
                  2(a-r)g_{\Sph^2(1)}+a dz^2,          \label{eq:ncs-bridge-model-metric}
\end{equation}
and scalar curvature at the model basepoint and $\theta=1$ is one.
For the actual rescaling one has the exact identity
\begin{equation}
 R_{h^{[\rho]}(1)}(x)=\rho R_g(x,-\rho)=a.            \label{eq:ncs-TypeI-rescaled-scalar-a}
\end{equation}
Consequently Type~I error at most $\varepsilon$ gives
$|a-1|\leq\omega(\varepsilon)$ for a universal modulus
$\omega(s)\to0$.  Abbreviate
$\varepsilon_{\rm TI}=\varepsilon_{\rm TI}(\delta)$ and choose it so that
\begin{equation}
 \varepsilon_{\rm TI}\leq\frac12,\qquad
 \omega(\varepsilon_{\rm TI})<\min\{\delta,1/2\},
 \qquad
 \varepsilon_{\rm TI}^{-1}\geq2(\delta^{-1}+1),       \label{eq:ncs-TypeI-bridge-buffer-choice}
\end{equation}
and so that the metric error in
\eqref{eq:ncs-bridge-model-metric} and every derivative required by the
target neck norm are below $\delta$.  The expanding
time window \eqref{eq:ncs-TypeI-fixed-window} contains
$\theta=1-r/a$ for $-1\leq r\leq0$ once
$a\in[1/2,3/2]$: indeed $1\leq\theta\leq3\leq
1+\varepsilon_{\rm TI}^{-2}$.  The last inequality in
\eqref{eq:ncs-TypeI-bridge-buffer-choice} leaves the target spatial buffer
under the explicit map
\begin{equation}
 J_a(y,z)=(y,z/\sqrt a),\qquad
 J_a^*\bigl(2(a-r)g_{\Sph^2(1)}+a\,dz^2\bigr)
        =2(a-r)g_{\Sph^2(1)}+dz^2.                  \label{eq:ncs-bridge-axial-map}
\end{equation}
For $a\in[1/2,3/2]$, the bound
$\varepsilon_{\rm TI}^{-1}\geq2(\delta^{-1}+1)$ ensures that $J_a$ is
defined on a domain containing the required target cylinder and its unit
buffer.  The restricted
pulled-back flow therefore satisfies exactly
Definition~\ref{def:ksol-strong-neck}.
\end{proof}

\section{Backward neck stability}
\label{sec:ncs-backward-stability}

\begin{theorem}[Backward stability of a strong neck]
\label{thm:ncs-backward-neck-stability}
There is a universal $\delta_{\rm neck}>0$ with the following property.  For
every $0<\delta_0,\delta_1\leq\delta_{\rm neck}$ there is a universal
$T=T(\delta_0,\delta_1)<0$ such that the following holds.  Let
$(M^3,g(t))$, $t\leq0$, be a noncompact ancient $\kappa$-solution and let
$x_0\in M$ satisfy
\begin{equation}
             R(x_0,0)=1,\qquad
             (x_0,0)\text{ is the center of a strong }
                         \delta_0\text{-neck}.          \label{eq:ncs-backward-stability-hypotheses}
\end{equation}
Then one of the following conclusions applies.
\begin{enumerate}
\item The flow is the diagonal quotient $\mathcal C_{\rm diag}$.
\item For every $t\leq T$, the point $(x_0,t)$ on the fixed identity
      worldline is the center of a strong $\delta_1$-neck, and it is
      Type~I $\delta_1$-cylindrical in the sense of
      Definition~\ref{def:ncs-TypeI-cylinder-witness}.
\end{enumerate}
The point $x_0$ in the second conclusion is not replaced by a moving
$\ell$-center.
\end{theorem}

\begin{proof}
Apply Theorem~\ref{thm:ksol-universal-kappa-gap}; the noncompact source has
the universal noncollapsing constant $\kappa_0$.  Choose
$\delta_{\rm neck}$ no larger than the topology threshold in
Lemma~\ref{lem:ncs-all-neck-rigidity} and no larger than the basepoint
constant in Proposition~\ref{prop:ncs-fixed-worldline-capture} for
$\eta=1/2$.  A strong $\delta_0$-neck weakens, by restriction of its
buffered chart, to the hypothesis of that proposition.  Hence there is a
universal $\rho_0$ such that
\begin{equation}
                         \ell_{x_0}(x_0,\rho)<3/2
                         \qquad(\rho\geq\rho_0).        \label{eq:ncs-fixed-worldline-three-halves}
\end{equation}

The source is orientable.  Indeed, the positive-curvature branch is
diffeomorphic to $\R^3$; the unquotiented and diagonal products are
orientable; and the only nonorientable split branch is
$\mathbb{RP}^2\times\R$, which cannot satisfy the initial ordinary-neck
hypothesis below the chosen topology threshold.

Choose
\begin{equation}
 0<\widehat\varepsilon\leq
 \min\{\delta_1,\varepsilon_{\rm TI}(\delta_1)\}.      \label{eq:ncs-stability-epsilon-choice}
\end{equation}
Then a $\widehat\varepsilon$-cylindrical triple restricts to a Type~I
$\delta_1$ witness, and Lemma
\ref{lem:ncs-TypeI-scalar-scale-bridge} turns it into a strong
$\delta_1$-neck.  Apply Proposition
\ref{prop:ncs-moving-basepoint-cylinder} with
$C=3/2$, $p=x_0$, $y=x_0$, and every sufficiently large $\rho$.  Its
global exceptional alternative is precisely the diagonal quotient.  If
that alternative does not hold, both desired conclusions hold at
$t=-\rho$ for every
\begin{equation}
 \rho\geq\max\{1,\rho_0,\rho_*(\widehat\varepsilon,3/2)\}.
\end{equation}
Take
\begin{equation}
 T=-\max\{1,\rho_0,\rho_*(\widehat\varepsilon,3/2)\}<0.       \label{eq:ncs-stability-time-choice}
\end{equation}
\end{proof}

\begin{corollary}[Curvature on the stable worldline]
\label{cor:ncs-neck-worldline-TypeI-curvature}
In the nonexceptional branch of
Theorem~\ref{thm:ncs-backward-neck-stability},
\begin{equation}
                       \lim_{t\to-\infty}(-t)R(x_0,t)=1.       \label{eq:ncs-worldline-TypeI-limit}
\end{equation}
In particular, for every $\eta>0$ there is $T_\eta<0$ such that
\begin{equation}
 \frac{1-\eta}{-t}\leq R(x_0,t)\leq\frac{1+\eta}{-t}
                         \qquad(t\leq T_\eta).          \label{eq:ncs-worldline-TypeI-two-sided}
\end{equation}
\end{corollary}

\begin{proof}
Apply Theorem~\ref{thm:ncs-backward-neck-stability} successively with
$\delta_1\downarrow0$.  The Type~I pointed metrics then converge on the
fixed window to $h_{\rm C}$.  Scalar curvature at the basepoint and
$\theta=1$ converges to one.  That scalar is exactly
$(-t)R(x_0,t)$, proving \eqref{eq:ncs-worldline-TypeI-limit}; the two-sided
estimate is its epsilon definition.
\end{proof}

\begin{corollary}[Scale- and time-translated form]
\label{cor:ncs-neck-persistence-past}
Let $0<\delta_0\leq\delta_{\rm neck}$ and let $(x_0,t_0)$ be a strong
$\delta_0$-neck center in a noncompact ancient $\kappa$-solution.  Put
$Q=R(x_0,t_0)$.  Unless the flow is the diagonal quotient, for every target
$0<\delta_1\leq\delta_{\rm neck}$ there is
$A=A(\delta_0,\delta_1)$ such that
\begin{equation}
 t\leq t_0-AQ^{-1}
   \Longrightarrow (x_0,t)\text{ is a strong }\delta_1\text{-neck center}.
                                                               \label{eq:ncs-translated-neck-persistence}
\end{equation}
The same translated Type~I conclusion holds about the fixed worldline.
Explicitly, for $\rho=Q(t_0-t)$ it concerns the family
\begin{equation}
 h_{t_0,Q}^{[\rho]}(\theta)
      =\frac{Q}{\rho}
         g\!\left(t_0-\frac{\rho\theta}{Q}\right),\qquad
 \theta\in[\delta_1,1+\delta_1^{-2}],                 \label{eq:ncs-translated-TypeI-family}
\end{equation}
with pole $(x_0,t_0)$ and the identity worldline through $x_0$.
\end{corollary}

\begin{proof}
Translate $t_0$ to zero and parabolically rescale by $Q$.  The hypotheses
become \eqref{eq:ncs-backward-stability-hypotheses}.  Apply
Theorem~\ref{thm:ncs-backward-neck-stability} and undo the translation and
rescaling.
\end{proof}

\begin{remark}[Why the diagonal exception is necessary]
\label{rem:ncs-diagonal-stability-counterexample}
Given any prescribed $0<\delta<1$, in $\mathcal C_{\rm diag}$ choose
$x_0=[y,z_0]$ with $|z_0|$ sufficiently large.  At time zero it is then the
center of a strong $\delta$ finite two-sided $\Sph^2$-neck.  Along that
chosen fixed worldline, however, its dimensionless
distance from the central projective cap under backward Type~I rescaling is
\begin{equation}
                       \frac{|z_0|}{\sqrt{-t}}\longrightarrow0.          \label{eq:ncs-diagonal-worldline-to-soul}
\end{equation}
The pointed backward limit is the diagonal quotient based at its central
one-sided $\mathbb{RP}^2$, not the unquotiented cylinder.  Thus deleting the
exception would make the theorem false.  The projective product is a
different issue: it never satisfies the initial ordinary spherical-neck
hypothesis when $\delta_0$ is sufficiently small.

There is also no honest repair obtained by merely adjoining a second
``central quotient neck'' constructor and retaining one universal backward
time.  In coordinates centered at a lift of $[y,z_0]$, the diagonal deck
map is
\begin{equation}
                         (y,w)\longmapsto(-y,-w-2z_0).
                         \label{eq:ncs-diagonal-affine-deck-map}
\end{equation}
For any fixed negative time, $z_0$ can be chosen so that the rescaled
offset is neither small enough for a central quotient witness nor large
enough for the prescribed longer ordinary neck.  Thus a future
surgery-spacetime generalized-neck theorem must store the affine offset
and use its own quantified conclusion; it is not a corollary of the fixed-
worldline theorem proved here.
Kleiner--Lott do prove such a smooth generalized-neck corollary with an
arbitrarily pointed quotient model.  We deliberately defer it until the
affine-offset witness and the survival domain of a surgery spacetime have
both been defined.
\end{remark}

\section{Formalization interfaces and dependency forest}
\label{sec:ncs-formalization}

The chapter is intended to be formalized through witness structures and
small producer--consumer contracts.  The following stable semantic nodes
record the proposed interfaces.
\begin{enumerate}
\item \nodeid{PC-F-KSOL-NECK-OPENNESS} is
      Lemma~\ref{lem:ncs-neck-openness}, including the strict
      $\varepsilon/4$--$\varepsilon/2$ closure buffer.

\item \nodeid{PC-F-KSOL-CAP-CORE-WITNESS},
      \nodeid{PC-F-KSOL-ONE-CAP-END-WITNESS}, and
      \nodeid{PC-F-KSOL-CONTROLLED-CAP-WITNESS} are the three nested
      structure layers of Definition~\ref{def:ncs-cap-witness}.  The
      topology tag, distinguished point, end map, dependent exterior-neck
      assignment, and quantitative controls are separate fields.

\item \nodeid{PC-F-KSOL-NECK-CHAIN-WITNESS},
      \nodeid{PC-F-KSOL-NECK-TUBE-WITNESS}, and
      \nodeid{PC-F-KSOL-TWO-CAP-TUBE-WITNESS} are the three structures in
      Definition~\ref{def:ncs-neck-chain}.  The chain stores its linear
      nerve and transition data; the tube additionally stores its global
      diffeomorphism.

\item \nodeid{PC-M-KSOL-LONG-RANGE-CURVATURE} is Theorem
      \ref{thm:ncs-long-range-curvature}, with the finite threshold theorem
      and discrete modulus construction as separate fine nodes.

\item \nodeid{PC-M-KSOL-CYLINDRICAL-END} is Theorem
      \ref{thm:ncs-arbitrary-escape-cylinder}.  Its producers are
      long-range transfer, universal precompactness, the escaping-line
      lemma, splitting, the surface theorem, and quotient exclusion.

\item \nodeid{PC-F-KSOL-ALIGNED-RAY-SPHERE} and
      \nodeid{PC-F-KSOL-ALL-NECK-RIGIDITY} are Lemmas
      \ref{lem:ncs-aligned-ray-cross-section} and
      \ref{lem:ncs-all-neck-rigidity}.  The first returns a strong-neck
      witness, its terminal sphere, a no-reentry equality, and, in exact
      models, the standard separating slice.  The second uses only the
      model dichotomy and finite Poincar\'e--Hopf obstructions.

\item \nodeid{PC-F-KSOL-REMOTE-CYLINDRICAL-SEPARATOR} is Lemma
      \ref{lem:ncs-remote-cylindrical-separator};
      \nodeid{PC-F-KSOL-NEAREST-POINT-PERSISTENCE} is identity
      \eqref{eq:ncs-nearest-point-persists}; and
      \nodeid{PC-F-KSOL-REMOTE-THREE-POINT} is Lemma
      \ref{lem:ncs-remote-three-point}.  Its remaining fine nodes are
      endpoint divergence, pointed orientability, line convergence, and
      the nonexistence of a minimizing ray orthogonal to the cylinder
      axis.  No Tits-boundary node is used.

\item \nodeid{PC-F-KSOL-REMOTE-AXIAL-EXIT},
      \nodeid{PC-F-KSOL-FARTHEST-NONNECK}, and
      \nodeid{PC-M-KSOL-FOUR-ALTERNATIVES} are, respectively, Lemma
      \ref{lem:ncs-remote-axial-exit},
      Lemma~\ref{lem:ncs-farthest-point-nonneck}, and Theorem
      \ref{thm:ncs-four-alternatives}.  Separation of the center sphere is
      an explicit input to axial exit.

\item \nodeid{PC-M-KSOL-CONTROLLED-NONNECK-CORE} is Proposition
      \ref{prop:ncs-bad-set-compact}.  It consumes both orders of the
      long-range estimate; scalar comparability is not hidden in a gradient
      bound.

\item \nodeid{PC-F-KSOL-NECK-OVERLAP} and
      \nodeid{PC-F-KSOL-NECK-CHAIN-ANNULUS} are Lemma
      \ref{lem:ncs-neck-overlap-alignment} and Proposition
      \ref{prop:ncs-neck-chain-annulus}.  Axial signs and additive constants
      are outputs, not quotient choices made silently later.

\item \nodeid{PC-M-KSOL-NONCOMPACT-NECKCAP} is Theorem
      \ref{thm:ncs-noncompact-neck-cap};
      \nodeid{PC-M-KSOL-COMPACT-QUALITATIVE-ALTERNATIVE} is Theorem
      \ref{thm:ncs-compact-alternative}.  Their fine producers include
      Lemmas \ref{lem:ncs-remote-endpoint-caps},
      \ref{lem:ncs-two-sided-axial-crossing}, and
      \ref{lem:ncs-ordered-middle-tube}.  The latter stores its linear
      nerve before invoking the finite-chain annulus theorem.

\item \nodeid{PC-F-KSOL-COMPACT-TOPOLOGY} is Corollary
      \ref{cor:ncs-compact-topology}; its shrinker-sequence witness and
      exclusion of the controlled-component branch are retained.

\item \nodeid{PC-M-KSOL-CANONICAL-NEIGHBORHOODS} is Theorem
      \ref{thm:ncs-kappa-canonical-neighborhood}.
      \nodeid{PC-F-KSOL-LOCAL-CANONICAL-NECK} and
      \nodeid{PC-F-KSOL-LOCAL-CANONICAL-CAP} are Definitions
      \ref{def:ncs-local-canonical-neck} and
      \ref{def:ncs-local-canonical-cap};
      \nodeid{PC-F-KSOL-CAP-TRUNCATION} is Lemma
      \ref{lem:ncs-cap-truncation-sandwich}.  These are separate fine nodes.  The
      theorem is in the oriented category; arbitrary round spherical
      quotients retain their separate constant-curvature constructor.

\item \nodeid{PC-F-KSOL-TYPE-I-CYLINDRICAL-WITNESS} is Definition
      \ref{def:ncs-TypeI-cylinder-witness}.  It is parameterized by the
      regular pole, time scale, point, one spacetime embedding, and axial
      potential offset.

\item \nodeid{PC-F-KSOL-NECKSTAB-REGULAR-POLE},
      \nodeid{PC-F-KSOL-NECKSTAB-BLOWDOWN-TRIPLE},
      \nodeid{PC-F-KSOL-NECKSTAB-GAUSSIAN-GAP}, and
      \nodeid{PC-F-KSOL-NECKSTAB-LOG-PLATEAU} are Lemma
      \ref{lem:ncs-regular-pole-tuple-compactness} and Lemmas
      \ref{lem:ncs-varying-ell-compactness}--
      \ref{lem:ncs-logarithmic-plateau}.

\item \nodeid{PC-F-KSOL-NECKSTAB-MOVING-CENTER},
      \nodeid{PC-F-KSOL-NECKSTAB-CYLINDER-THRESHOLD},
      \nodeid{PC-F-KSOL-NECKSTAB-SHARP-ELL-TIME}, and
      \nodeid{PC-F-KSOL-NECKSTAB-BASEPOINT-ELL} are Proposition
      \ref{prop:ncs-moving-basepoint-cylinder}, Lemma
      \ref{lem:ncs-cylinder-threshold}, Lemma
      \ref{lem:ncs-sharp-ell-time-comparison}, and Proposition
      \ref{prop:ncs-fixed-worldline-capture}.

\item \nodeid{PC-F-KSOL-NECKSTAB-SCALE-BRIDGE} is Lemma
      \ref{lem:ncs-TypeI-scalar-scale-bridge};
      \nodeid{PC-M-KSOL-BACKWARD-NECK-STABILITY} is Theorem
      \ref{thm:ncs-backward-neck-stability}.

\item \nodeid{PC-B-KSOL-SURGERY-CANONICAL-INTERFACE} is the tuple of
      Theorem~\ref{thm:ncs-kappa-canonical-neighborhood}, the cap and neck
      witnesses, and Corollary~\ref{cor:ncs-neck-persistence-past}.  A
      future surgery-spacetime theorem must add survival domains and a
      separately typed affine-offset generalized-neck interface; it cannot
      reuse this smooth interface without those data.
\end{enumerate}

The coarse dependency forest is:
\begin{enumerate}
\item universal compactness $+$ long-range transfer $+$ the rank-one list
      produce cylindrical geometry at arbitrary escaping points;
\item cylindrical infinity $+$ the remote-three-point lemma produce the
      four alternatives and controlled nonneck cores;
\item buffered witnesses $+$ overlap alignment $+$ three-manifold topology
      produce one-cap ends and two-cap tubes;
\item those global alternatives $+$ universal compactness produce the
      canonical-neighborhood theorem;
\item varying-flow reduced compactness $+$ Gaussian gap $+$ plateau
      selection $+$ shrinker classification produce moving-center
      cylindricality;
\item moving-center cylindricality $+$ first crossing $+$ the offset
      cylinder potential produce fixed-worldline $\ell$ capture;
\item fixed-worldline capture $+$ the Type~I/scalar-scale bridge produce
      backward neck stability.
\end{enumerate}

The imported classical topology interfaces are: the positive-curvature
soul theorem; smooth Schoenflies and the exterior product theorem in
$\R^3$; collars and isotopy extension;
irreducibility of $\mathbb{RP}^3$; the explicit punctured-projective model;
van Kampen; Synge orientability; Bonnet--Myers finiteness of the
fundamental group; relative Poincar\'e--Hopf; the odd-dimensional identity
$2\chi(K)=\chi(\partial K)$; nonembeddability of a two-sided
$\mathbb{RP}^2$ in an oriented three-manifold; the smooth simple-eigenline
lemma; and Ehresmann triviality over an interval.  Each is a blueprint
obligation.  So are preservation of orientability under pointed smooth
convergence, Hamilton compactness with worldline-preserving spacetime
embeddings, varying-flow reduced-distance compactness, and convergence of
total reduced mass.

\subsection*{Lean status}

Every new node in this chapter has Lean statement and Lean proof status
\emph{not yet implemented}, unless a project ledger supplies
declaration-level compilation and axiom-footprint evidence.  This prose is
the intended human mathematical source, not kernel-checking evidence.

\section{Sources, repairs, and deliberate stopping point}
\label{sec:ncs-sources}

The qualitative core comes from G.~Perelman,
\emph{The entropy formula for the Ricci flow and its geometric
applications}, especially Sections~11.7--11.8.  The detailed reconstruction
uses B.~Kleiner and J.~Lott, \emph{Notes on Perelman's papers}, especially
Corollaries~42.1, 47.2, and~48.1, Lemma~48.3, and the neck--cap results in
Sections~58--59, especially Lemmas~59.1--59.3 and~59.7.  These results are
stated there under the standing orientation convention.  Backward neck
stability is Theorem~6.1 of Kleiner--Lott, \emph{Singular Ricci Flows I}.
Their literal statement fixes $\kappa$ and allows its neck and time
constants to depend on $\kappa$.  The theorem here is the universalized
corollary obtained by first replacing every noncompact source by the
universal $\kappa_0$ from Theorem
\ref{thm:ksol-universal-kappa-gap}; this is a change in quantifier
presentation, not a stronger geometric input.
In arXiv:1408.2271v3 the supporting chain is Lemma~6.3,
Proposition~6.12, and Lemma~6.13, together with the intervening gap,
plateau, and cylinder lemmas; the published journal version renumbers the
main moving-center and fixed-basepoint inputs as Proposition~6.10 and
Lemma~6.11.

Among the project sources, the main classical expositions are MSM~163,
Chapters~18--20 and Appendix~I; GSM~235, Section~8.2; MSM~206,
Chapter~28; and GSM~77, Chapter~9.  Chapter~\ref{ch:reduced-geometry}
supplies the corrected no-loss asymptotic-shrinker theorem, and Chapters
\ref{ch:three-dimensional-shrinkers} and
\ref{ch:kappa-compactness} supply the exact shrinker and cylinder-quotient
lists.

The following repairs and choices are part of the mathematics.
\begin{enumerate}
\item The concise printed compact-core theorem in GSM~235 occurs under the
      standing positive-curvature noncompact hypothesis, although its
      displayed statement says only $\Rm\geq0$.  Read literally it includes
      the exact cylinder, where the nonneck core is empty.  We separated the
      cylinder before stating Proposition~\ref{prop:ncs-bad-set-compact}.

\item A point of $\partial\mathscr B_\varepsilon$ need not be an
      $\varepsilon$-neck center, and this boundary need not be smooth.
      Lemma~\ref{lem:ncs-neck-openness} uses
      $\partial\mathscr B_{\varepsilon/4}\subset
       \mathscr N_{\varepsilon/2}$ and then chooses the center sphere of an
      actual witness.

\item Recentring changes the scalar normalization and may ask for a
      slightly earlier time than the original interval $[-1,0]$.
      Lemma~\ref{lem:ncs-neck-openness} first creates a genuine temporal
      buffer using Harnack and the universal derivative bounds.

\item $\ASCR=+\infty$ is a limsup statement.  It gives selected sequences,
      not all escaping sequences.  Theorem
      \ref{thm:ncs-long-range-curvature} is the indispensable bridge to
      Theorem~\ref{thm:ncs-arbitrary-escape-cylinder}.

\item Without orientation, $\mathbb{RP}^2\times\R$ is an additional exact
      projective-cylinder branch.  The diagonal quotient is orientable and
      cannot be discarded on orientation grounds; its cap is
      $\mathbb{RP}^3\setminus\operatorname{int}B^3$.

\item The complement of neck centers is not the definition of a cap.
      Definition~\ref{def:ncs-cap-witness} stores a smooth core, a model
      diffeomorphism, a boundary neck, a collar sign, an end map, and all
      quantitative controls.

\item Kleiner--Lott Lemma~59.1 is qualitative.  Accordingly Theorem
      \ref{thm:ncs-noncompact-neck-cap} does not claim that an arbitrary
      far neck bounds a cap uniformly controlled at that neck's scale; the
      Bryant soliton rules out such an inference.  Uniform metric control
      is asserted locally only in Theorem
      \ref{thm:ncs-kappa-canonical-neighborhood}, reconstructed from their
      Lemma~59.7 by the finite cap-truncation lemma.

\item A compact slice need not have a long neck.  A compact ancient oval
      near extinction is the basic regression test.  Theorem
      \ref{thm:ncs-compact-alternative} retains the controlled-whole-
      component alternative and assigns separate scales to two caps.

\item The moving-center proof requires compactness for a varying sequence
      of flows and no loss of reduced-volume mass.  It is not a formal
      consequence of the single-flow asymptotic-shrinker theorem.  Lemma
      \ref{lem:ncs-varying-ell-compactness} states the required interface.

\item The Type~I witness uses the expanding interval
      $[\delta,1+\delta^{-2}]$.  A fixed interval ending at $2$ would not
      support the exact change of normalization when
      $\rho R(x,-\rho)<1$.

\item The extinction-centered cylinder potential may have the axial offset
      in \eqref{eq:ncs-cylinder-offset-potential}.  Setting it to zero would
      destroy the first-crossing contradiction.  Likewise, a Type~I
      rescaling and a scalar-curvature normalization are not identical;
      their exact conversion is \eqref{eq:ncs-two-normalizations-bridge}.

\item Backward stability concerns the fixed identity worldline.  A theorem
      producing only moving $\ell$-centers is insufficient.  Proposition
      \ref{prop:ncs-fixed-worldline-capture} is the separate bridge, and
      Remark~\ref{rem:ncs-diagonal-stability-counterexample} records the
      sharp quotient obstruction.

\item Appendix~I of MSM~163 states the distance-nonincreasing
      Sharafutdinov level map used in Chapter
      \ref{ch:kappa-compactness}, but the printed diameter argument also
      needs surjectivity of its restriction to the outer level.  Lemma
      \ref{lem:ksol-Sharafutdinov-level-surjective} supplies the omitted
      step by constructing a coray from each prescribed inner-level point;
      it does not assume that the gradient semigroup is invertible.
\end{enumerate}

This chapter deliberately stops before constant-mean-curvature neck
foliations, approximate Killing fields, quantitative neck improvement,
rotational symmetry, Bryant-soliton uniqueness, ancient-oval uniqueness,
and the complete Brendle--Daskalopoulos--\v{S}e\v{s}um classification.
Those results may refine the list of isometry types, but they are not
dependencies of Perelman's qualitative canonical-neighborhood and surgery
theory developed here.

\backmatter

\bibliographystyle{amsalpha}
\bibliography{references}

\printindex

\end{document}
